\section*{Item 4}\label{it:4}

\textbf{Enunciado}: Considerando o circuito apresentado na Lista de Exercícios.  O  transistor  é comandado  por  um  pulso  quadrado  com  largura  $50$\%,  em  $25$ kHz.  Deseja-se obter $10$ V na saída, com um \textit{ripple} de tensão de $1$\%. A corrente nominal de saída é de $5$ A. Os pulsos de comando do transistor devem variar entre $-15$ e $+15$ V, com tempos de subida e de descida de $100$ ns.

\medskip
\textbf{a)} Calcule e use na simulação a indutância para operar no limiar entre o MCC e o MCD, com uma corrente de saída de $1$ A.

\medskip
\textbf{Resolução}: o valor da indutância para operar no limiar pode ser obtida por
\begin{equation*}
	L_{min} = \frac{R_{o,max}T_s}{2}(1 - D),
\end{equation*}
em que $R_{o,max} = V_o/I_{o,min}$. Substituindo-se os valores fornecidos no enunciado, tem-se $L_{min} = 100~\mu$H.

\medskip
\textbf{b)} Calcule o capacitor de filtro para o \textit{ripple} de tensão indicado.

\medskip
\textbf{Resolução}: o valor do capacitor de filtro para que a tensão de saída apresente um \textit{ripple} de $1$\%, ou seja $\Delta v_{c,pico} = 0{,}01V_o = 0{,}1$ V, pode ser obtido por
\begin{equation*}
	C = \Delta i_L\frac{T_s}{4\Delta v_{c,pico}},
\end{equation*}
em que
\begin{equation*}
	\Delta i_L = \frac{E - V_o}{2L_{min}}DT_s.
\end{equation*}

Substituindo-se os valores fornecidos no enunciado, tem-se $C = 100~\mu$F e $\Delta i_L = 1$ A.

\medskip
\textbf{c)} Simule o circuito, pelo menos por $6$ ms, partindo de condições iniciais nulas, tanto no indutor  quanto no capacitor, e verifique se os valores teóricos correspondem aos simulados. Explique eventuais discrepâncias. Deve-se mostrar o resultado correspondentes as cargas com correntes de $5$ e $1$ A (valores médios).

\medskip
\textbf{Resolução}: a Figura~\ref{fig:ex4_1} mostra o resultado da simulação para a carga com corrente de $5$ A.
\begin{figure}[htb!]
	\centering
	%% Creator: Matplotlib, PGF backend
%%
%% To include the figure in your LaTeX document, write
%%   \input{<filename>.pgf}
%%
%% Make sure the required packages are loaded in your preamble
%%   \usepackage{pgf}
%%
%% Also ensure that all the required font packages are loaded; for instance,
%% the lmodern package is sometimes necessary when using math font.
%%   \usepackage{lmodern}
%%
%% Figures using additional raster images can only be included by \input if
%% they are in the same directory as the main LaTeX file. For loading figures
%% from other directories you can use the `import` package
%%   \usepackage{import}
%%
%% and then include the figures with
%%   \import{<path to file>}{<filename>.pgf}
%%
%% Matplotlib used the following preamble
%%   \def\mathdefault#1{#1}
%%   \everymath=\expandafter{\the\everymath\displaystyle}
%%   \IfFileExists{scrextend.sty}{
%%     \usepackage[fontsize=8.000000pt]{scrextend}
%%   }{
%%     \renewcommand{\normalsize}{\fontsize{8.000000}{9.600000}\selectfont}
%%     \normalsize
%%   }
%%   
%%           \usepackage{amsmath}
%%           \usepackage{amssymb}
%%           \usepackage{mathtools}
%%           \usepackage{dsfont}
%%           \providecommand{\mathdefault}[1]{#1}
%%       
%%   \makeatletter\@ifpackageloaded{underscore}{}{\usepackage[strings]{underscore}}\makeatother
%%
\begingroup%
\makeatletter%
\begin{pgfpicture}%
\pgfpathrectangle{\pgfpointorigin}{\pgfqpoint{3.423091in}{3.646235in}}%
\pgfusepath{use as bounding box, clip}%
\begin{pgfscope}%
\pgfsetbuttcap%
\pgfsetmiterjoin%
\definecolor{currentfill}{rgb}{1.000000,1.000000,1.000000}%
\pgfsetfillcolor{currentfill}%
\pgfsetlinewidth{0.000000pt}%
\definecolor{currentstroke}{rgb}{1.000000,1.000000,1.000000}%
\pgfsetstrokecolor{currentstroke}%
\pgfsetdash{}{0pt}%
\pgfpathmoveto{\pgfqpoint{0.000000in}{0.000000in}}%
\pgfpathlineto{\pgfqpoint{3.423091in}{0.000000in}}%
\pgfpathlineto{\pgfqpoint{3.423091in}{3.646235in}}%
\pgfpathlineto{\pgfqpoint{0.000000in}{3.646235in}}%
\pgfpathlineto{\pgfqpoint{0.000000in}{0.000000in}}%
\pgfpathclose%
\pgfusepath{fill}%
\end{pgfscope}%
\begin{pgfscope}%
\pgfsetbuttcap%
\pgfsetmiterjoin%
\definecolor{currentfill}{rgb}{1.000000,1.000000,1.000000}%
\pgfsetfillcolor{currentfill}%
\pgfsetlinewidth{0.000000pt}%
\definecolor{currentstroke}{rgb}{0.000000,0.000000,0.000000}%
\pgfsetstrokecolor{currentstroke}%
\pgfsetstrokeopacity{0.000000}%
\pgfsetdash{}{0pt}%
\pgfpathmoveto{\pgfqpoint{0.573768in}{2.601404in}}%
\pgfpathlineto{\pgfqpoint{3.293768in}{2.601404in}}%
\pgfpathlineto{\pgfqpoint{3.293768in}{3.507654in}}%
\pgfpathlineto{\pgfqpoint{0.573768in}{3.507654in}}%
\pgfpathlineto{\pgfqpoint{0.573768in}{2.601404in}}%
\pgfpathclose%
\pgfusepath{fill}%
\end{pgfscope}%
\begin{pgfscope}%
\pgfpathrectangle{\pgfqpoint{0.573768in}{2.601404in}}{\pgfqpoint{2.720000in}{0.906250in}}%
\pgfusepath{clip}%
\pgfsetbuttcap%
\pgfsetroundjoin%
\pgfsetlinewidth{0.501875pt}%
\definecolor{currentstroke}{rgb}{0.690196,0.690196,0.690196}%
\pgfsetstrokecolor{currentstroke}%
\pgfsetstrokeopacity{0.250000}%
\pgfsetdash{{1.850000pt}{0.800000pt}}{0.000000pt}%
\pgfpathmoveto{\pgfqpoint{0.697405in}{2.601404in}}%
\pgfpathlineto{\pgfqpoint{0.697405in}{3.507654in}}%
\pgfusepath{stroke}%
\end{pgfscope}%
\begin{pgfscope}%
\pgfsetbuttcap%
\pgfsetroundjoin%
\definecolor{currentfill}{rgb}{0.000000,0.000000,0.000000}%
\pgfsetfillcolor{currentfill}%
\pgfsetlinewidth{0.803000pt}%
\definecolor{currentstroke}{rgb}{0.000000,0.000000,0.000000}%
\pgfsetstrokecolor{currentstroke}%
\pgfsetdash{}{0pt}%
\pgfsys@defobject{currentmarker}{\pgfqpoint{0.000000in}{-0.048611in}}{\pgfqpoint{0.000000in}{0.000000in}}{%
\pgfpathmoveto{\pgfqpoint{0.000000in}{0.000000in}}%
\pgfpathlineto{\pgfqpoint{0.000000in}{-0.048611in}}%
\pgfusepath{stroke,fill}%
}%
\begin{pgfscope}%
\pgfsys@transformshift{0.697405in}{2.601404in}%
\pgfsys@useobject{currentmarker}{}%
\end{pgfscope}%
\end{pgfscope}%
\begin{pgfscope}%
\pgfpathrectangle{\pgfqpoint{0.573768in}{2.601404in}}{\pgfqpoint{2.720000in}{0.906250in}}%
\pgfusepath{clip}%
\pgfsetbuttcap%
\pgfsetroundjoin%
\pgfsetlinewidth{0.501875pt}%
\definecolor{currentstroke}{rgb}{0.690196,0.690196,0.690196}%
\pgfsetstrokecolor{currentstroke}%
\pgfsetstrokeopacity{0.250000}%
\pgfsetdash{{1.850000pt}{0.800000pt}}{0.000000pt}%
\pgfpathmoveto{\pgfqpoint{1.109526in}{2.601404in}}%
\pgfpathlineto{\pgfqpoint{1.109526in}{3.507654in}}%
\pgfusepath{stroke}%
\end{pgfscope}%
\begin{pgfscope}%
\pgfsetbuttcap%
\pgfsetroundjoin%
\definecolor{currentfill}{rgb}{0.000000,0.000000,0.000000}%
\pgfsetfillcolor{currentfill}%
\pgfsetlinewidth{0.803000pt}%
\definecolor{currentstroke}{rgb}{0.000000,0.000000,0.000000}%
\pgfsetstrokecolor{currentstroke}%
\pgfsetdash{}{0pt}%
\pgfsys@defobject{currentmarker}{\pgfqpoint{0.000000in}{-0.048611in}}{\pgfqpoint{0.000000in}{0.000000in}}{%
\pgfpathmoveto{\pgfqpoint{0.000000in}{0.000000in}}%
\pgfpathlineto{\pgfqpoint{0.000000in}{-0.048611in}}%
\pgfusepath{stroke,fill}%
}%
\begin{pgfscope}%
\pgfsys@transformshift{1.109526in}{2.601404in}%
\pgfsys@useobject{currentmarker}{}%
\end{pgfscope}%
\end{pgfscope}%
\begin{pgfscope}%
\pgfpathrectangle{\pgfqpoint{0.573768in}{2.601404in}}{\pgfqpoint{2.720000in}{0.906250in}}%
\pgfusepath{clip}%
\pgfsetbuttcap%
\pgfsetroundjoin%
\pgfsetlinewidth{0.501875pt}%
\definecolor{currentstroke}{rgb}{0.690196,0.690196,0.690196}%
\pgfsetstrokecolor{currentstroke}%
\pgfsetstrokeopacity{0.250000}%
\pgfsetdash{{1.850000pt}{0.800000pt}}{0.000000pt}%
\pgfpathmoveto{\pgfqpoint{1.521647in}{2.601404in}}%
\pgfpathlineto{\pgfqpoint{1.521647in}{3.507654in}}%
\pgfusepath{stroke}%
\end{pgfscope}%
\begin{pgfscope}%
\pgfsetbuttcap%
\pgfsetroundjoin%
\definecolor{currentfill}{rgb}{0.000000,0.000000,0.000000}%
\pgfsetfillcolor{currentfill}%
\pgfsetlinewidth{0.803000pt}%
\definecolor{currentstroke}{rgb}{0.000000,0.000000,0.000000}%
\pgfsetstrokecolor{currentstroke}%
\pgfsetdash{}{0pt}%
\pgfsys@defobject{currentmarker}{\pgfqpoint{0.000000in}{-0.048611in}}{\pgfqpoint{0.000000in}{0.000000in}}{%
\pgfpathmoveto{\pgfqpoint{0.000000in}{0.000000in}}%
\pgfpathlineto{\pgfqpoint{0.000000in}{-0.048611in}}%
\pgfusepath{stroke,fill}%
}%
\begin{pgfscope}%
\pgfsys@transformshift{1.521647in}{2.601404in}%
\pgfsys@useobject{currentmarker}{}%
\end{pgfscope}%
\end{pgfscope}%
\begin{pgfscope}%
\pgfpathrectangle{\pgfqpoint{0.573768in}{2.601404in}}{\pgfqpoint{2.720000in}{0.906250in}}%
\pgfusepath{clip}%
\pgfsetbuttcap%
\pgfsetroundjoin%
\pgfsetlinewidth{0.501875pt}%
\definecolor{currentstroke}{rgb}{0.690196,0.690196,0.690196}%
\pgfsetstrokecolor{currentstroke}%
\pgfsetstrokeopacity{0.250000}%
\pgfsetdash{{1.850000pt}{0.800000pt}}{0.000000pt}%
\pgfpathmoveto{\pgfqpoint{1.933768in}{2.601404in}}%
\pgfpathlineto{\pgfqpoint{1.933768in}{3.507654in}}%
\pgfusepath{stroke}%
\end{pgfscope}%
\begin{pgfscope}%
\pgfsetbuttcap%
\pgfsetroundjoin%
\definecolor{currentfill}{rgb}{0.000000,0.000000,0.000000}%
\pgfsetfillcolor{currentfill}%
\pgfsetlinewidth{0.803000pt}%
\definecolor{currentstroke}{rgb}{0.000000,0.000000,0.000000}%
\pgfsetstrokecolor{currentstroke}%
\pgfsetdash{}{0pt}%
\pgfsys@defobject{currentmarker}{\pgfqpoint{0.000000in}{-0.048611in}}{\pgfqpoint{0.000000in}{0.000000in}}{%
\pgfpathmoveto{\pgfqpoint{0.000000in}{0.000000in}}%
\pgfpathlineto{\pgfqpoint{0.000000in}{-0.048611in}}%
\pgfusepath{stroke,fill}%
}%
\begin{pgfscope}%
\pgfsys@transformshift{1.933768in}{2.601404in}%
\pgfsys@useobject{currentmarker}{}%
\end{pgfscope}%
\end{pgfscope}%
\begin{pgfscope}%
\pgfpathrectangle{\pgfqpoint{0.573768in}{2.601404in}}{\pgfqpoint{2.720000in}{0.906250in}}%
\pgfusepath{clip}%
\pgfsetbuttcap%
\pgfsetroundjoin%
\pgfsetlinewidth{0.501875pt}%
\definecolor{currentstroke}{rgb}{0.690196,0.690196,0.690196}%
\pgfsetstrokecolor{currentstroke}%
\pgfsetstrokeopacity{0.250000}%
\pgfsetdash{{1.850000pt}{0.800000pt}}{0.000000pt}%
\pgfpathmoveto{\pgfqpoint{2.345890in}{2.601404in}}%
\pgfpathlineto{\pgfqpoint{2.345890in}{3.507654in}}%
\pgfusepath{stroke}%
\end{pgfscope}%
\begin{pgfscope}%
\pgfsetbuttcap%
\pgfsetroundjoin%
\definecolor{currentfill}{rgb}{0.000000,0.000000,0.000000}%
\pgfsetfillcolor{currentfill}%
\pgfsetlinewidth{0.803000pt}%
\definecolor{currentstroke}{rgb}{0.000000,0.000000,0.000000}%
\pgfsetstrokecolor{currentstroke}%
\pgfsetdash{}{0pt}%
\pgfsys@defobject{currentmarker}{\pgfqpoint{0.000000in}{-0.048611in}}{\pgfqpoint{0.000000in}{0.000000in}}{%
\pgfpathmoveto{\pgfqpoint{0.000000in}{0.000000in}}%
\pgfpathlineto{\pgfqpoint{0.000000in}{-0.048611in}}%
\pgfusepath{stroke,fill}%
}%
\begin{pgfscope}%
\pgfsys@transformshift{2.345890in}{2.601404in}%
\pgfsys@useobject{currentmarker}{}%
\end{pgfscope}%
\end{pgfscope}%
\begin{pgfscope}%
\pgfpathrectangle{\pgfqpoint{0.573768in}{2.601404in}}{\pgfqpoint{2.720000in}{0.906250in}}%
\pgfusepath{clip}%
\pgfsetbuttcap%
\pgfsetroundjoin%
\pgfsetlinewidth{0.501875pt}%
\definecolor{currentstroke}{rgb}{0.690196,0.690196,0.690196}%
\pgfsetstrokecolor{currentstroke}%
\pgfsetstrokeopacity{0.250000}%
\pgfsetdash{{1.850000pt}{0.800000pt}}{0.000000pt}%
\pgfpathmoveto{\pgfqpoint{2.758011in}{2.601404in}}%
\pgfpathlineto{\pgfqpoint{2.758011in}{3.507654in}}%
\pgfusepath{stroke}%
\end{pgfscope}%
\begin{pgfscope}%
\pgfsetbuttcap%
\pgfsetroundjoin%
\definecolor{currentfill}{rgb}{0.000000,0.000000,0.000000}%
\pgfsetfillcolor{currentfill}%
\pgfsetlinewidth{0.803000pt}%
\definecolor{currentstroke}{rgb}{0.000000,0.000000,0.000000}%
\pgfsetstrokecolor{currentstroke}%
\pgfsetdash{}{0pt}%
\pgfsys@defobject{currentmarker}{\pgfqpoint{0.000000in}{-0.048611in}}{\pgfqpoint{0.000000in}{0.000000in}}{%
\pgfpathmoveto{\pgfqpoint{0.000000in}{0.000000in}}%
\pgfpathlineto{\pgfqpoint{0.000000in}{-0.048611in}}%
\pgfusepath{stroke,fill}%
}%
\begin{pgfscope}%
\pgfsys@transformshift{2.758011in}{2.601404in}%
\pgfsys@useobject{currentmarker}{}%
\end{pgfscope}%
\end{pgfscope}%
\begin{pgfscope}%
\pgfpathrectangle{\pgfqpoint{0.573768in}{2.601404in}}{\pgfqpoint{2.720000in}{0.906250in}}%
\pgfusepath{clip}%
\pgfsetbuttcap%
\pgfsetroundjoin%
\pgfsetlinewidth{0.501875pt}%
\definecolor{currentstroke}{rgb}{0.690196,0.690196,0.690196}%
\pgfsetstrokecolor{currentstroke}%
\pgfsetstrokeopacity{0.250000}%
\pgfsetdash{{1.850000pt}{0.800000pt}}{0.000000pt}%
\pgfpathmoveto{\pgfqpoint{3.170132in}{2.601404in}}%
\pgfpathlineto{\pgfqpoint{3.170132in}{3.507654in}}%
\pgfusepath{stroke}%
\end{pgfscope}%
\begin{pgfscope}%
\pgfsetbuttcap%
\pgfsetroundjoin%
\definecolor{currentfill}{rgb}{0.000000,0.000000,0.000000}%
\pgfsetfillcolor{currentfill}%
\pgfsetlinewidth{0.803000pt}%
\definecolor{currentstroke}{rgb}{0.000000,0.000000,0.000000}%
\pgfsetstrokecolor{currentstroke}%
\pgfsetdash{}{0pt}%
\pgfsys@defobject{currentmarker}{\pgfqpoint{0.000000in}{-0.048611in}}{\pgfqpoint{0.000000in}{0.000000in}}{%
\pgfpathmoveto{\pgfqpoint{0.000000in}{0.000000in}}%
\pgfpathlineto{\pgfqpoint{0.000000in}{-0.048611in}}%
\pgfusepath{stroke,fill}%
}%
\begin{pgfscope}%
\pgfsys@transformshift{3.170132in}{2.601404in}%
\pgfsys@useobject{currentmarker}{}%
\end{pgfscope}%
\end{pgfscope}%
\begin{pgfscope}%
\pgfpathrectangle{\pgfqpoint{0.573768in}{2.601404in}}{\pgfqpoint{2.720000in}{0.906250in}}%
\pgfusepath{clip}%
\pgfsetbuttcap%
\pgfsetroundjoin%
\pgfsetlinewidth{0.501875pt}%
\definecolor{currentstroke}{rgb}{0.690196,0.690196,0.690196}%
\pgfsetstrokecolor{currentstroke}%
\pgfsetstrokeopacity{0.250000}%
\pgfsetdash{{1.850000pt}{0.800000pt}}{0.000000pt}%
\pgfpathmoveto{\pgfqpoint{0.573768in}{2.601404in}}%
\pgfpathlineto{\pgfqpoint{3.293768in}{2.601404in}}%
\pgfusepath{stroke}%
\end{pgfscope}%
\begin{pgfscope}%
\pgfsetbuttcap%
\pgfsetroundjoin%
\definecolor{currentfill}{rgb}{0.000000,0.000000,0.000000}%
\pgfsetfillcolor{currentfill}%
\pgfsetlinewidth{0.803000pt}%
\definecolor{currentstroke}{rgb}{0.000000,0.000000,0.000000}%
\pgfsetstrokecolor{currentstroke}%
\pgfsetdash{}{0pt}%
\pgfsys@defobject{currentmarker}{\pgfqpoint{-0.048611in}{0.000000in}}{\pgfqpoint{-0.000000in}{0.000000in}}{%
\pgfpathmoveto{\pgfqpoint{-0.000000in}{0.000000in}}%
\pgfpathlineto{\pgfqpoint{-0.048611in}{0.000000in}}%
\pgfusepath{stroke,fill}%
}%
\begin{pgfscope}%
\pgfsys@transformshift{0.573768in}{2.601404in}%
\pgfsys@useobject{currentmarker}{}%
\end{pgfscope}%
\end{pgfscope}%
\begin{pgfscope}%
\definecolor{textcolor}{rgb}{0.000000,0.000000,0.000000}%
\pgfsetstrokecolor{textcolor}%
\pgfsetfillcolor{textcolor}%
\pgftext[x=0.417518in, y=2.562824in, left, base]{\color{textcolor}{\rmfamily\fontsize{8.000000}{9.600000}\selectfont\catcode`\^=\active\def^{\ifmmode\sp\else\^{}\fi}\catcode`\%=\active\def%{\%}$\mathdefault{0}$}}%
\end{pgfscope}%
\begin{pgfscope}%
\pgfpathrectangle{\pgfqpoint{0.573768in}{2.601404in}}{\pgfqpoint{2.720000in}{0.906250in}}%
\pgfusepath{clip}%
\pgfsetbuttcap%
\pgfsetroundjoin%
\pgfsetlinewidth{0.501875pt}%
\definecolor{currentstroke}{rgb}{0.690196,0.690196,0.690196}%
\pgfsetstrokecolor{currentstroke}%
\pgfsetstrokeopacity{0.250000}%
\pgfsetdash{{1.850000pt}{0.800000pt}}{0.000000pt}%
\pgfpathmoveto{\pgfqpoint{0.573768in}{2.903488in}}%
\pgfpathlineto{\pgfqpoint{3.293768in}{2.903488in}}%
\pgfusepath{stroke}%
\end{pgfscope}%
\begin{pgfscope}%
\pgfsetbuttcap%
\pgfsetroundjoin%
\definecolor{currentfill}{rgb}{0.000000,0.000000,0.000000}%
\pgfsetfillcolor{currentfill}%
\pgfsetlinewidth{0.803000pt}%
\definecolor{currentstroke}{rgb}{0.000000,0.000000,0.000000}%
\pgfsetstrokecolor{currentstroke}%
\pgfsetdash{}{0pt}%
\pgfsys@defobject{currentmarker}{\pgfqpoint{-0.048611in}{0.000000in}}{\pgfqpoint{-0.000000in}{0.000000in}}{%
\pgfpathmoveto{\pgfqpoint{-0.000000in}{0.000000in}}%
\pgfpathlineto{\pgfqpoint{-0.048611in}{0.000000in}}%
\pgfusepath{stroke,fill}%
}%
\begin{pgfscope}%
\pgfsys@transformshift{0.573768in}{2.903488in}%
\pgfsys@useobject{currentmarker}{}%
\end{pgfscope}%
\end{pgfscope}%
\begin{pgfscope}%
\definecolor{textcolor}{rgb}{0.000000,0.000000,0.000000}%
\pgfsetstrokecolor{textcolor}%
\pgfsetfillcolor{textcolor}%
\pgftext[x=0.417518in, y=2.864908in, left, base]{\color{textcolor}{\rmfamily\fontsize{8.000000}{9.600000}\selectfont\catcode`\^=\active\def^{\ifmmode\sp\else\^{}\fi}\catcode`\%=\active\def%{\%}$\mathdefault{5}$}}%
\end{pgfscope}%
\begin{pgfscope}%
\pgfpathrectangle{\pgfqpoint{0.573768in}{2.601404in}}{\pgfqpoint{2.720000in}{0.906250in}}%
\pgfusepath{clip}%
\pgfsetbuttcap%
\pgfsetroundjoin%
\pgfsetlinewidth{0.501875pt}%
\definecolor{currentstroke}{rgb}{0.690196,0.690196,0.690196}%
\pgfsetstrokecolor{currentstroke}%
\pgfsetstrokeopacity{0.250000}%
\pgfsetdash{{1.850000pt}{0.800000pt}}{0.000000pt}%
\pgfpathmoveto{\pgfqpoint{0.573768in}{3.205571in}}%
\pgfpathlineto{\pgfqpoint{3.293768in}{3.205571in}}%
\pgfusepath{stroke}%
\end{pgfscope}%
\begin{pgfscope}%
\pgfsetbuttcap%
\pgfsetroundjoin%
\definecolor{currentfill}{rgb}{0.000000,0.000000,0.000000}%
\pgfsetfillcolor{currentfill}%
\pgfsetlinewidth{0.803000pt}%
\definecolor{currentstroke}{rgb}{0.000000,0.000000,0.000000}%
\pgfsetstrokecolor{currentstroke}%
\pgfsetdash{}{0pt}%
\pgfsys@defobject{currentmarker}{\pgfqpoint{-0.048611in}{0.000000in}}{\pgfqpoint{-0.000000in}{0.000000in}}{%
\pgfpathmoveto{\pgfqpoint{-0.000000in}{0.000000in}}%
\pgfpathlineto{\pgfqpoint{-0.048611in}{0.000000in}}%
\pgfusepath{stroke,fill}%
}%
\begin{pgfscope}%
\pgfsys@transformshift{0.573768in}{3.205571in}%
\pgfsys@useobject{currentmarker}{}%
\end{pgfscope}%
\end{pgfscope}%
\begin{pgfscope}%
\definecolor{textcolor}{rgb}{0.000000,0.000000,0.000000}%
\pgfsetstrokecolor{textcolor}%
\pgfsetfillcolor{textcolor}%
\pgftext[x=0.358489in, y=3.166991in, left, base]{\color{textcolor}{\rmfamily\fontsize{8.000000}{9.600000}\selectfont\catcode`\^=\active\def^{\ifmmode\sp\else\^{}\fi}\catcode`\%=\active\def%{\%}$\mathdefault{10}$}}%
\end{pgfscope}%
\begin{pgfscope}%
\pgfpathrectangle{\pgfqpoint{0.573768in}{2.601404in}}{\pgfqpoint{2.720000in}{0.906250in}}%
\pgfusepath{clip}%
\pgfsetbuttcap%
\pgfsetroundjoin%
\pgfsetlinewidth{0.501875pt}%
\definecolor{currentstroke}{rgb}{0.690196,0.690196,0.690196}%
\pgfsetstrokecolor{currentstroke}%
\pgfsetstrokeopacity{0.250000}%
\pgfsetdash{{1.850000pt}{0.800000pt}}{0.000000pt}%
\pgfpathmoveto{\pgfqpoint{0.573768in}{3.507654in}}%
\pgfpathlineto{\pgfqpoint{3.293768in}{3.507654in}}%
\pgfusepath{stroke}%
\end{pgfscope}%
\begin{pgfscope}%
\pgfsetbuttcap%
\pgfsetroundjoin%
\definecolor{currentfill}{rgb}{0.000000,0.000000,0.000000}%
\pgfsetfillcolor{currentfill}%
\pgfsetlinewidth{0.803000pt}%
\definecolor{currentstroke}{rgb}{0.000000,0.000000,0.000000}%
\pgfsetstrokecolor{currentstroke}%
\pgfsetdash{}{0pt}%
\pgfsys@defobject{currentmarker}{\pgfqpoint{-0.048611in}{0.000000in}}{\pgfqpoint{-0.000000in}{0.000000in}}{%
\pgfpathmoveto{\pgfqpoint{-0.000000in}{0.000000in}}%
\pgfpathlineto{\pgfqpoint{-0.048611in}{0.000000in}}%
\pgfusepath{stroke,fill}%
}%
\begin{pgfscope}%
\pgfsys@transformshift{0.573768in}{3.507654in}%
\pgfsys@useobject{currentmarker}{}%
\end{pgfscope}%
\end{pgfscope}%
\begin{pgfscope}%
\definecolor{textcolor}{rgb}{0.000000,0.000000,0.000000}%
\pgfsetstrokecolor{textcolor}%
\pgfsetfillcolor{textcolor}%
\pgftext[x=0.358489in, y=3.469074in, left, base]{\color{textcolor}{\rmfamily\fontsize{8.000000}{9.600000}\selectfont\catcode`\^=\active\def^{\ifmmode\sp\else\^{}\fi}\catcode`\%=\active\def%{\%}$\mathdefault{15}$}}%
\end{pgfscope}%
\begin{pgfscope}%
\definecolor{textcolor}{rgb}{0.000000,0.000000,0.000000}%
\pgfsetstrokecolor{textcolor}%
\pgfsetfillcolor{textcolor}%
\pgftext[x=0.302933in,y=3.054529in,,bottom,rotate=90.000000]{\color{textcolor}{\rmfamily\fontsize{8.000000}{9.600000}\selectfont\catcode`\^=\active\def^{\ifmmode\sp\else\^{}\fi}\catcode`\%=\active\def%{\%}$v_o$ (V)}}%
\end{pgfscope}%
\begin{pgfscope}%
\pgfpathrectangle{\pgfqpoint{0.573768in}{2.601404in}}{\pgfqpoint{2.720000in}{0.906250in}}%
\pgfusepath{clip}%
\pgfsetrectcap%
\pgfsetroundjoin%
\pgfsetlinewidth{1.505625pt}%
\definecolor{currentstroke}{rgb}{0.000000,0.447059,0.698039}%
\pgfsetstrokecolor{currentstroke}%
\pgfsetdash{}{0pt}%
\pgfpathmoveto{\pgfqpoint{0.697405in}{2.601404in}}%
\pgfpathlineto{\pgfqpoint{0.699132in}{2.602049in}}%
\pgfpathlineto{\pgfqpoint{0.700945in}{2.604919in}}%
\pgfpathlineto{\pgfqpoint{0.703253in}{2.611809in}}%
\pgfpathlineto{\pgfqpoint{0.706355in}{2.626544in}}%
\pgfpathlineto{\pgfqpoint{0.714119in}{2.665109in}}%
\pgfpathlineto{\pgfqpoint{0.717560in}{2.683462in}}%
\pgfpathlineto{\pgfqpoint{0.721681in}{2.714511in}}%
\pgfpathlineto{\pgfqpoint{0.732440in}{2.801619in}}%
\pgfpathlineto{\pgfqpoint{0.737550in}{2.848453in}}%
\pgfpathlineto{\pgfqpoint{0.747840in}{2.947408in}}%
\pgfpathlineto{\pgfqpoint{0.753610in}{3.001800in}}%
\pgfpathlineto{\pgfqpoint{0.763193in}{3.094786in}}%
\pgfpathlineto{\pgfqpoint{0.770581in}{3.157317in}}%
\pgfpathlineto{\pgfqpoint{0.778046in}{3.221516in}}%
\pgfpathlineto{\pgfqpoint{0.794247in}{3.325846in}}%
\pgfpathlineto{\pgfqpoint{0.798609in}{3.341838in}}%
\pgfpathlineto{\pgfqpoint{0.803719in}{3.363380in}}%
\pgfpathlineto{\pgfqpoint{0.809091in}{3.385369in}}%
\pgfpathlineto{\pgfqpoint{0.811894in}{3.390757in}}%
\pgfpathlineto{\pgfqpoint{0.814870in}{3.392943in}}%
\pgfpathlineto{\pgfqpoint{0.817837in}{3.396548in}}%
\pgfpathlineto{\pgfqpoint{0.821480in}{3.403655in}}%
\pgfpathlineto{\pgfqpoint{0.824087in}{3.407854in}}%
\pgfpathlineto{\pgfqpoint{0.826065in}{3.408556in}}%
\pgfpathlineto{\pgfqpoint{0.827713in}{3.407542in}}%
\pgfpathlineto{\pgfqpoint{0.829901in}{3.404034in}}%
\pgfpathlineto{\pgfqpoint{0.833452in}{3.398615in}}%
\pgfpathlineto{\pgfqpoint{0.836420in}{3.396552in}}%
\pgfpathlineto{\pgfqpoint{0.840737in}{3.394461in}}%
\pgfpathlineto{\pgfqpoint{0.843045in}{3.389896in}}%
\pgfpathlineto{\pgfqpoint{0.846098in}{3.379832in}}%
\pgfpathlineto{\pgfqpoint{0.850699in}{3.364334in}}%
\pgfpathlineto{\pgfqpoint{0.854676in}{3.355940in}}%
\pgfpathlineto{\pgfqpoint{0.857232in}{3.350011in}}%
\pgfpathlineto{\pgfqpoint{0.860529in}{3.337885in}}%
\pgfpathlineto{\pgfqpoint{0.870787in}{3.295707in}}%
\pgfpathlineto{\pgfqpoint{0.874065in}{3.285468in}}%
\pgfpathlineto{\pgfqpoint{0.878021in}{3.267184in}}%
\pgfpathlineto{\pgfqpoint{0.884893in}{3.233968in}}%
\pgfpathlineto{\pgfqpoint{0.894988in}{3.195497in}}%
\pgfpathlineto{\pgfqpoint{0.900455in}{3.169885in}}%
\pgfpathlineto{\pgfqpoint{0.904251in}{3.159358in}}%
\pgfpathlineto{\pgfqpoint{0.907548in}{3.150045in}}%
\pgfpathlineto{\pgfqpoint{0.911606in}{3.133505in}}%
\pgfpathlineto{\pgfqpoint{0.916279in}{3.114760in}}%
\pgfpathlineto{\pgfqpoint{0.919576in}{3.107343in}}%
\pgfpathlineto{\pgfqpoint{0.924342in}{3.098971in}}%
\pgfpathlineto{\pgfqpoint{0.927948in}{3.087934in}}%
\pgfpathlineto{\pgfqpoint{0.932394in}{3.074532in}}%
\pgfpathlineto{\pgfqpoint{0.935197in}{3.070659in}}%
\pgfpathlineto{\pgfqpoint{0.937521in}{3.069932in}}%
\pgfpathlineto{\pgfqpoint{0.939664in}{3.068671in}}%
\pgfpathlineto{\pgfqpoint{0.942301in}{3.065081in}}%
\pgfpathlineto{\pgfqpoint{0.945742in}{3.057344in}}%
\pgfpathlineto{\pgfqpoint{0.948544in}{3.052776in}}%
\pgfpathlineto{\pgfqpoint{0.950687in}{3.051712in}}%
\pgfpathlineto{\pgfqpoint{0.952500in}{3.052413in}}%
\pgfpathlineto{\pgfqpoint{0.957087in}{3.055184in}}%
\pgfpathlineto{\pgfqpoint{0.959230in}{3.054231in}}%
\pgfpathlineto{\pgfqpoint{0.961894in}{3.051071in}}%
\pgfpathlineto{\pgfqpoint{0.964202in}{3.049205in}}%
\pgfpathlineto{\pgfqpoint{0.966015in}{3.049426in}}%
\pgfpathlineto{\pgfqpoint{0.967993in}{3.051319in}}%
\pgfpathlineto{\pgfqpoint{0.971091in}{3.057024in}}%
\pgfpathlineto{\pgfqpoint{0.974059in}{3.060591in}}%
\pgfpathlineto{\pgfqpoint{0.976367in}{3.061295in}}%
\pgfpathlineto{\pgfqpoint{0.981642in}{3.061379in}}%
\pgfpathlineto{\pgfqpoint{0.983950in}{3.064748in}}%
\pgfpathlineto{\pgfqpoint{0.987115in}{3.072696in}}%
\pgfpathlineto{\pgfqpoint{0.990576in}{3.079545in}}%
\pgfpathlineto{\pgfqpoint{0.993214in}{3.081870in}}%
\pgfpathlineto{\pgfqpoint{0.997898in}{3.084460in}}%
\pgfpathlineto{\pgfqpoint{1.000536in}{3.089545in}}%
\pgfpathlineto{\pgfqpoint{1.009400in}{3.109716in}}%
\pgfpathlineto{\pgfqpoint{1.015215in}{3.115116in}}%
\pgfpathlineto{\pgfqpoint{1.018182in}{3.122457in}}%
\pgfpathlineto{\pgfqpoint{1.023913in}{3.137565in}}%
\pgfpathlineto{\pgfqpoint{1.026551in}{3.140334in}}%
\pgfpathlineto{\pgfqpoint{1.030664in}{3.143221in}}%
\pgfpathlineto{\pgfqpoint{1.033302in}{3.148219in}}%
\pgfpathlineto{\pgfqpoint{1.042047in}{3.167334in}}%
\pgfpathlineto{\pgfqpoint{1.044276in}{3.167833in}}%
\pgfpathlineto{\pgfqpoint{1.046198in}{3.168543in}}%
\pgfpathlineto{\pgfqpoint{1.048506in}{3.171189in}}%
\pgfpathlineto{\pgfqpoint{1.051308in}{3.176947in}}%
\pgfpathlineto{\pgfqpoint{1.056227in}{3.187522in}}%
\pgfpathlineto{\pgfqpoint{1.058535in}{3.189120in}}%
\pgfpathlineto{\pgfqpoint{1.060576in}{3.188722in}}%
\pgfpathlineto{\pgfqpoint{1.062603in}{3.188451in}}%
\pgfpathlineto{\pgfqpoint{1.064746in}{3.189782in}}%
\pgfpathlineto{\pgfqpoint{1.067219in}{3.193325in}}%
\pgfpathlineto{\pgfqpoint{1.073038in}{3.202739in}}%
\pgfpathlineto{\pgfqpoint{1.075016in}{3.202956in}}%
\pgfpathlineto{\pgfqpoint{1.077106in}{3.201451in}}%
\pgfpathlineto{\pgfqpoint{1.079506in}{3.200082in}}%
\pgfpathlineto{\pgfqpoint{1.081649in}{3.200613in}}%
\pgfpathlineto{\pgfqpoint{1.083957in}{3.202983in}}%
\pgfpathlineto{\pgfqpoint{1.089022in}{3.208983in}}%
\pgfpathlineto{\pgfqpoint{1.091000in}{3.208605in}}%
\pgfpathlineto{\pgfqpoint{1.093172in}{3.206341in}}%
\pgfpathlineto{\pgfqpoint{1.096351in}{3.202945in}}%
\pgfpathlineto{\pgfqpoint{1.098494in}{3.202716in}}%
\pgfpathlineto{\pgfqpoint{1.100638in}{3.204111in}}%
\pgfpathlineto{\pgfqpoint{1.105011in}{3.207674in}}%
\pgfpathlineto{\pgfqpoint{1.106989in}{3.206844in}}%
\pgfpathlineto{\pgfqpoint{1.109297in}{3.203865in}}%
\pgfpathlineto{\pgfqpoint{1.113341in}{3.198146in}}%
\pgfpathlineto{\pgfqpoint{1.115649in}{3.197516in}}%
\pgfpathlineto{\pgfqpoint{1.117822in}{3.198667in}}%
\pgfpathlineto{\pgfqpoint{1.120671in}{3.200333in}}%
\pgfpathlineto{\pgfqpoint{1.122650in}{3.199549in}}%
\pgfpathlineto{\pgfqpoint{1.124957in}{3.196643in}}%
\pgfpathlineto{\pgfqpoint{1.130995in}{3.187461in}}%
\pgfpathlineto{\pgfqpoint{1.133138in}{3.187229in}}%
\pgfpathlineto{\pgfqpoint{1.138590in}{3.188224in}}%
\pgfpathlineto{\pgfqpoint{1.140898in}{3.185435in}}%
\pgfpathlineto{\pgfqpoint{1.148304in}{3.174209in}}%
\pgfpathlineto{\pgfqpoint{1.150447in}{3.174462in}}%
\pgfpathlineto{\pgfqpoint{1.153911in}{3.175526in}}%
\pgfpathlineto{\pgfqpoint{1.155890in}{3.174109in}}%
\pgfpathlineto{\pgfqpoint{1.158362in}{3.170228in}}%
\pgfpathlineto{\pgfqpoint{1.163220in}{3.161778in}}%
\pgfpathlineto{\pgfqpoint{1.165528in}{3.160848in}}%
\pgfpathlineto{\pgfqpoint{1.167573in}{3.161722in}}%
\pgfpathlineto{\pgfqpoint{1.169951in}{3.162674in}}%
\pgfpathlineto{\pgfqpoint{1.171929in}{3.161784in}}%
\pgfpathlineto{\pgfqpoint{1.174237in}{3.158868in}}%
\pgfpathlineto{\pgfqpoint{1.180542in}{3.149500in}}%
\pgfpathlineto{\pgfqpoint{1.182520in}{3.149518in}}%
\pgfpathlineto{\pgfqpoint{1.185125in}{3.151453in}}%
\pgfpathlineto{\pgfqpoint{1.187269in}{3.151890in}}%
\pgfpathlineto{\pgfqpoint{1.189247in}{3.150735in}}%
\pgfpathlineto{\pgfqpoint{1.191719in}{3.147246in}}%
\pgfpathlineto{\pgfqpoint{1.195596in}{3.141478in}}%
\pgfpathlineto{\pgfqpoint{1.197739in}{3.140821in}}%
\pgfpathlineto{\pgfqpoint{1.199717in}{3.141812in}}%
\pgfpathlineto{\pgfqpoint{1.203858in}{3.144670in}}%
\pgfpathlineto{\pgfqpoint{1.205836in}{3.143864in}}%
\pgfpathlineto{\pgfqpoint{1.208144in}{3.141095in}}%
\pgfpathlineto{\pgfqpoint{1.211935in}{3.136186in}}%
\pgfpathlineto{\pgfqpoint{1.214078in}{3.135895in}}%
\pgfpathlineto{\pgfqpoint{1.216056in}{3.137228in}}%
\pgfpathlineto{\pgfqpoint{1.220678in}{3.141237in}}%
\pgfpathlineto{\pgfqpoint{1.222656in}{3.140633in}}%
\pgfpathlineto{\pgfqpoint{1.225004in}{3.138040in}}%
\pgfpathlineto{\pgfqpoint{1.228091in}{3.134638in}}%
\pgfpathlineto{\pgfqpoint{1.230234in}{3.134519in}}%
\pgfpathlineto{\pgfqpoint{1.232377in}{3.136203in}}%
\pgfpathlineto{\pgfqpoint{1.237500in}{3.141412in}}%
\pgfpathlineto{\pgfqpoint{1.239478in}{3.140941in}}%
\pgfpathlineto{\pgfqpoint{1.241788in}{3.138557in}}%
\pgfpathlineto{\pgfqpoint{1.244239in}{3.136255in}}%
\pgfpathlineto{\pgfqpoint{1.246217in}{3.136159in}}%
\pgfpathlineto{\pgfqpoint{1.248360in}{3.137783in}}%
\pgfpathlineto{\pgfqpoint{1.254832in}{3.144530in}}%
\pgfpathlineto{\pgfqpoint{1.256810in}{3.143713in}}%
\pgfpathlineto{\pgfqpoint{1.262201in}{3.140144in}}%
\pgfpathlineto{\pgfqpoint{1.264180in}{3.141421in}}%
\pgfpathlineto{\pgfqpoint{1.266650in}{3.145095in}}%
\pgfpathlineto{\pgfqpoint{1.269530in}{3.148866in}}%
\pgfpathlineto{\pgfqpoint{1.271673in}{3.149556in}}%
\pgfpathlineto{\pgfqpoint{1.273651in}{3.148634in}}%
\pgfpathlineto{\pgfqpoint{1.278153in}{3.145584in}}%
\pgfpathlineto{\pgfqpoint{1.280132in}{3.146558in}}%
\pgfpathlineto{\pgfqpoint{1.282549in}{3.149760in}}%
\pgfpathlineto{\pgfqpoint{1.286180in}{3.154843in}}%
\pgfpathlineto{\pgfqpoint{1.288323in}{3.155444in}}%
\pgfpathlineto{\pgfqpoint{1.290301in}{3.154428in}}%
\pgfpathlineto{\pgfqpoint{1.294531in}{3.151527in}}%
\pgfpathlineto{\pgfqpoint{1.296509in}{3.152406in}}%
\pgfpathlineto{\pgfqpoint{1.298817in}{3.155290in}}%
\pgfpathlineto{\pgfqpoint{1.302829in}{3.160759in}}%
\pgfpathlineto{\pgfqpoint{1.304972in}{3.161149in}}%
\pgfpathlineto{\pgfqpoint{1.306951in}{3.159929in}}%
\pgfpathlineto{\pgfqpoint{1.310850in}{3.156986in}}%
\pgfpathlineto{\pgfqpoint{1.312828in}{3.157628in}}%
\pgfpathlineto{\pgfqpoint{1.315136in}{3.160230in}}%
\pgfpathlineto{\pgfqpoint{1.319479in}{3.165767in}}%
\pgfpathlineto{\pgfqpoint{1.321622in}{3.165857in}}%
\pgfpathlineto{\pgfqpoint{1.323715in}{3.164224in}}%
\pgfpathlineto{\pgfqpoint{1.327006in}{3.161334in}}%
\pgfpathlineto{\pgfqpoint{1.328985in}{3.161556in}}%
\pgfpathlineto{\pgfqpoint{1.331128in}{3.163452in}}%
\pgfpathlineto{\pgfqpoint{1.336623in}{3.169409in}}%
\pgfpathlineto{\pgfqpoint{1.338601in}{3.168849in}}%
\pgfpathlineto{\pgfqpoint{1.340923in}{3.166290in}}%
\pgfpathlineto{\pgfqpoint{1.343431in}{3.164092in}}%
\pgfpathlineto{\pgfqpoint{1.345574in}{3.164132in}}%
\pgfpathlineto{\pgfqpoint{1.347717in}{3.165891in}}%
\pgfpathlineto{\pgfqpoint{1.352778in}{3.171136in}}%
\pgfpathlineto{\pgfqpoint{1.354757in}{3.170631in}}%
\pgfpathlineto{\pgfqpoint{1.357040in}{3.168162in}}%
\pgfpathlineto{\pgfqpoint{1.359954in}{3.165138in}}%
\pgfpathlineto{\pgfqpoint{1.362097in}{3.164998in}}%
\pgfpathlineto{\pgfqpoint{1.364240in}{3.166579in}}%
\pgfpathlineto{\pgfqpoint{1.369263in}{3.171307in}}%
\pgfpathlineto{\pgfqpoint{1.371241in}{3.170623in}}%
\pgfpathlineto{\pgfqpoint{1.373550in}{3.167913in}}%
\pgfpathlineto{\pgfqpoint{1.376442in}{3.164679in}}%
\pgfpathlineto{\pgfqpoint{1.378585in}{3.164367in}}%
\pgfpathlineto{\pgfqpoint{1.380728in}{3.165783in}}%
\pgfpathlineto{\pgfqpoint{1.385748in}{3.170137in}}%
\pgfpathlineto{\pgfqpoint{1.387726in}{3.169318in}}%
\pgfpathlineto{\pgfqpoint{1.390096in}{3.166356in}}%
\pgfpathlineto{\pgfqpoint{1.392930in}{3.163048in}}%
\pgfpathlineto{\pgfqpoint{1.395073in}{3.162618in}}%
\pgfpathlineto{\pgfqpoint{1.397216in}{3.163923in}}%
\pgfpathlineto{\pgfqpoint{1.402218in}{3.168033in}}%
\pgfpathlineto{\pgfqpoint{1.404196in}{3.167141in}}%
\pgfpathlineto{\pgfqpoint{1.406581in}{3.164074in}}%
\pgfpathlineto{\pgfqpoint{1.409417in}{3.160668in}}%
\pgfpathlineto{\pgfqpoint{1.411560in}{3.160176in}}%
\pgfpathlineto{\pgfqpoint{1.413538in}{3.161271in}}%
\pgfpathlineto{\pgfqpoint{1.418853in}{3.165413in}}%
\pgfpathlineto{\pgfqpoint{1.420831in}{3.164375in}}%
\pgfpathlineto{\pgfqpoint{1.423303in}{3.161004in}}%
\pgfpathlineto{\pgfqpoint{1.425910in}{3.157979in}}%
\pgfpathlineto{\pgfqpoint{1.428053in}{3.157488in}}%
\pgfpathlineto{\pgfqpoint{1.430031in}{3.158588in}}%
\pgfpathlineto{\pgfqpoint{1.435334in}{3.162752in}}%
\pgfpathlineto{\pgfqpoint{1.437313in}{3.161738in}}%
\pgfpathlineto{\pgfqpoint{1.439749in}{3.158463in}}%
\pgfpathlineto{\pgfqpoint{1.442394in}{3.155419in}}%
\pgfpathlineto{\pgfqpoint{1.444537in}{3.154970in}}%
\pgfpathlineto{\pgfqpoint{1.446516in}{3.156116in}}%
\pgfpathlineto{\pgfqpoint{1.451817in}{3.160422in}}%
\pgfpathlineto{\pgfqpoint{1.453795in}{3.159470in}}%
\pgfpathlineto{\pgfqpoint{1.456234in}{3.156272in}}%
\pgfpathlineto{\pgfqpoint{1.458880in}{3.153320in}}%
\pgfpathlineto{\pgfqpoint{1.461023in}{3.152950in}}%
\pgfpathlineto{\pgfqpoint{1.463001in}{3.154171in}}%
\pgfpathlineto{\pgfqpoint{1.468468in}{3.158677in}}%
\pgfpathlineto{\pgfqpoint{1.470446in}{3.157684in}}%
\pgfpathlineto{\pgfqpoint{1.472907in}{3.154439in}}%
\pgfpathlineto{\pgfqpoint{1.475530in}{3.151814in}}%
\pgfpathlineto{\pgfqpoint{1.477673in}{3.151679in}}%
\pgfpathlineto{\pgfqpoint{1.479816in}{3.153300in}}%
\pgfpathlineto{\pgfqpoint{1.484464in}{3.157690in}}%
\pgfpathlineto{\pgfqpoint{1.486442in}{3.157160in}}%
\pgfpathlineto{\pgfqpoint{1.488739in}{3.154683in}}%
\pgfpathlineto{\pgfqpoint{1.491841in}{3.151248in}}%
\pgfpathlineto{\pgfqpoint{1.493984in}{3.151066in}}%
\pgfpathlineto{\pgfqpoint{1.496127in}{3.152645in}}%
\pgfpathlineto{\pgfqpoint{1.501117in}{3.157468in}}%
\pgfpathlineto{\pgfqpoint{1.503096in}{3.156896in}}%
\pgfpathlineto{\pgfqpoint{1.505404in}{3.154355in}}%
\pgfpathlineto{\pgfqpoint{1.508160in}{3.151400in}}%
\pgfpathlineto{\pgfqpoint{1.510303in}{3.151166in}}%
\pgfpathlineto{\pgfqpoint{1.512446in}{3.152693in}}%
\pgfpathlineto{\pgfqpoint{1.517940in}{3.157915in}}%
\pgfpathlineto{\pgfqpoint{1.519918in}{3.157157in}}%
\pgfpathlineto{\pgfqpoint{1.522284in}{3.154339in}}%
\pgfpathlineto{\pgfqpoint{1.524808in}{3.151987in}}%
\pgfpathlineto{\pgfqpoint{1.526951in}{3.151953in}}%
\pgfpathlineto{\pgfqpoint{1.529094in}{3.153673in}}%
\pgfpathlineto{\pgfqpoint{1.534270in}{3.158842in}}%
\pgfpathlineto{\pgfqpoint{1.536249in}{3.158249in}}%
\pgfpathlineto{\pgfqpoint{1.538568in}{3.155665in}}%
\pgfpathlineto{\pgfqpoint{1.541127in}{3.153146in}}%
\pgfpathlineto{\pgfqpoint{1.543270in}{3.153013in}}%
\pgfpathlineto{\pgfqpoint{1.545413in}{3.154633in}}%
\pgfpathlineto{\pgfqpoint{1.550926in}{3.160080in}}%
\pgfpathlineto{\pgfqpoint{1.552904in}{3.159378in}}%
\pgfpathlineto{\pgfqpoint{1.555255in}{3.156643in}}%
\pgfpathlineto{\pgfqpoint{1.557775in}{3.154361in}}%
\pgfpathlineto{\pgfqpoint{1.559919in}{3.154374in}}%
\pgfpathlineto{\pgfqpoint{1.562062in}{3.156133in}}%
\pgfpathlineto{\pgfqpoint{1.567081in}{3.161356in}}%
\pgfpathlineto{\pgfqpoint{1.569059in}{3.160902in}}%
\pgfpathlineto{\pgfqpoint{1.571344in}{3.158517in}}%
\pgfpathlineto{\pgfqpoint{1.574095in}{3.155711in}}%
\pgfpathlineto{\pgfqpoint{1.576238in}{3.155571in}}%
\pgfpathlineto{\pgfqpoint{1.578381in}{3.157179in}}%
\pgfpathlineto{\pgfqpoint{1.583896in}{3.162567in}}%
\pgfpathlineto{\pgfqpoint{1.585874in}{3.161834in}}%
\pgfpathlineto{\pgfqpoint{1.588225in}{3.159054in}}%
\pgfpathlineto{\pgfqpoint{1.590745in}{3.156719in}}%
\pgfpathlineto{\pgfqpoint{1.592888in}{3.156680in}}%
\pgfpathlineto{\pgfqpoint{1.595031in}{3.158383in}}%
\pgfpathlineto{\pgfqpoint{1.600051in}{3.163457in}}%
\pgfpathlineto{\pgfqpoint{1.602029in}{3.162937in}}%
\pgfpathlineto{\pgfqpoint{1.604313in}{3.160474in}}%
\pgfpathlineto{\pgfqpoint{1.607065in}{3.157568in}}%
\pgfpathlineto{\pgfqpoint{1.609208in}{3.157348in}}%
\pgfpathlineto{\pgfqpoint{1.611351in}{3.158872in}}%
\pgfpathlineto{\pgfqpoint{1.616865in}{3.164034in}}%
\pgfpathlineto{\pgfqpoint{1.618844in}{3.163217in}}%
\pgfpathlineto{\pgfqpoint{1.621265in}{3.160233in}}%
\pgfpathlineto{\pgfqpoint{1.623880in}{3.157818in}}%
\pgfpathlineto{\pgfqpoint{1.626023in}{3.157821in}}%
\pgfpathlineto{\pgfqpoint{1.628166in}{3.159562in}}%
\pgfpathlineto{\pgfqpoint{1.632856in}{3.164201in}}%
\pgfpathlineto{\pgfqpoint{1.634834in}{3.163721in}}%
\pgfpathlineto{\pgfqpoint{1.637103in}{3.161323in}}%
\pgfpathlineto{\pgfqpoint{1.640201in}{3.157937in}}%
\pgfpathlineto{\pgfqpoint{1.642344in}{3.157763in}}%
\pgfpathlineto{\pgfqpoint{1.644487in}{3.159331in}}%
\pgfpathlineto{\pgfqpoint{1.649503in}{3.164087in}}%
\pgfpathlineto{\pgfqpoint{1.651481in}{3.163446in}}%
\pgfpathlineto{\pgfqpoint{1.653793in}{3.160802in}}%
\pgfpathlineto{\pgfqpoint{1.656522in}{3.157771in}}%
\pgfpathlineto{\pgfqpoint{1.658665in}{3.157428in}}%
\pgfpathlineto{\pgfqpoint{1.660808in}{3.158832in}}%
\pgfpathlineto{\pgfqpoint{1.666313in}{3.163691in}}%
\pgfpathlineto{\pgfqpoint{1.668292in}{3.162776in}}%
\pgfpathlineto{\pgfqpoint{1.670719in}{3.159666in}}%
\pgfpathlineto{\pgfqpoint{1.673337in}{3.157125in}}%
\pgfpathlineto{\pgfqpoint{1.675480in}{3.157033in}}%
\pgfpathlineto{\pgfqpoint{1.677623in}{3.158683in}}%
\pgfpathlineto{\pgfqpoint{1.682299in}{3.163126in}}%
\pgfpathlineto{\pgfqpoint{1.684277in}{3.162579in}}%
\pgfpathlineto{\pgfqpoint{1.686580in}{3.160058in}}%
\pgfpathlineto{\pgfqpoint{1.689657in}{3.156607in}}%
\pgfpathlineto{\pgfqpoint{1.691800in}{3.156371in}}%
\pgfpathlineto{\pgfqpoint{1.693943in}{3.157883in}}%
\pgfpathlineto{\pgfqpoint{1.698946in}{3.162515in}}%
\pgfpathlineto{\pgfqpoint{1.700924in}{3.161843in}}%
\pgfpathlineto{\pgfqpoint{1.703247in}{3.159149in}}%
\pgfpathlineto{\pgfqpoint{1.705978in}{3.156076in}}%
\pgfpathlineto{\pgfqpoint{1.708121in}{3.155708in}}%
\pgfpathlineto{\pgfqpoint{1.710099in}{3.156925in}}%
\pgfpathlineto{\pgfqpoint{1.715924in}{3.161897in}}%
\pgfpathlineto{\pgfqpoint{1.717902in}{3.160858in}}%
\pgfpathlineto{\pgfqpoint{1.720484in}{3.157440in}}%
\pgfpathlineto{\pgfqpoint{1.722957in}{3.155278in}}%
\pgfpathlineto{\pgfqpoint{1.725100in}{3.155333in}}%
\pgfpathlineto{\pgfqpoint{1.727243in}{3.157130in}}%
\pgfpathlineto{\pgfqpoint{1.731584in}{3.161390in}}%
\pgfpathlineto{\pgfqpoint{1.733562in}{3.160997in}}%
\pgfpathlineto{\pgfqpoint{1.735828in}{3.158709in}}%
\pgfpathlineto{\pgfqpoint{1.739277in}{3.154922in}}%
\pgfpathlineto{\pgfqpoint{1.741420in}{3.154862in}}%
\pgfpathlineto{\pgfqpoint{1.743563in}{3.156547in}}%
\pgfpathlineto{\pgfqpoint{1.748234in}{3.161080in}}%
\pgfpathlineto{\pgfqpoint{1.750213in}{3.160581in}}%
\pgfpathlineto{\pgfqpoint{1.752497in}{3.158148in}}%
\pgfpathlineto{\pgfqpoint{1.755597in}{3.154744in}}%
\pgfpathlineto{\pgfqpoint{1.757740in}{3.154570in}}%
\pgfpathlineto{\pgfqpoint{1.759883in}{3.156146in}}%
\pgfpathlineto{\pgfqpoint{1.764886in}{3.160941in}}%
\pgfpathlineto{\pgfqpoint{1.766864in}{3.160336in}}%
\pgfpathlineto{\pgfqpoint{1.769187in}{3.157725in}}%
\pgfpathlineto{\pgfqpoint{1.771916in}{3.154751in}}%
\pgfpathlineto{\pgfqpoint{1.774059in}{3.154463in}}%
\pgfpathlineto{\pgfqpoint{1.776202in}{3.155928in}}%
\pgfpathlineto{\pgfqpoint{1.781703in}{3.160965in}}%
\pgfpathlineto{\pgfqpoint{1.783681in}{3.160127in}}%
\pgfpathlineto{\pgfqpoint{1.786113in}{3.157109in}}%
\pgfpathlineto{\pgfqpoint{1.788730in}{3.154677in}}%
\pgfpathlineto{\pgfqpoint{1.790874in}{3.154680in}}%
\pgfpathlineto{\pgfqpoint{1.793017in}{3.156428in}}%
\pgfpathlineto{\pgfqpoint{1.797696in}{3.161099in}}%
\pgfpathlineto{\pgfqpoint{1.799674in}{3.160651in}}%
\pgfpathlineto{\pgfqpoint{1.801951in}{3.158285in}}%
\pgfpathlineto{\pgfqpoint{1.805050in}{3.154964in}}%
\pgfpathlineto{\pgfqpoint{1.807193in}{3.154843in}}%
\pgfpathlineto{\pgfqpoint{1.809336in}{3.156470in}}%
\pgfpathlineto{\pgfqpoint{1.814348in}{3.161384in}}%
\pgfpathlineto{\pgfqpoint{1.816326in}{3.160817in}}%
\pgfpathlineto{\pgfqpoint{1.818641in}{3.158259in}}%
\pgfpathlineto{\pgfqpoint{1.821370in}{3.155338in}}%
\pgfpathlineto{\pgfqpoint{1.823513in}{3.155087in}}%
\pgfpathlineto{\pgfqpoint{1.825656in}{3.156587in}}%
\pgfpathlineto{\pgfqpoint{1.831164in}{3.161709in}}%
\pgfpathlineto{\pgfqpoint{1.833142in}{3.160892in}}%
\pgfpathlineto{\pgfqpoint{1.835568in}{3.157908in}}%
\pgfpathlineto{\pgfqpoint{1.838184in}{3.155504in}}%
\pgfpathlineto{\pgfqpoint{1.840327in}{3.155525in}}%
\pgfpathlineto{\pgfqpoint{1.842470in}{3.157289in}}%
\pgfpathlineto{\pgfqpoint{1.847155in}{3.161991in}}%
\pgfpathlineto{\pgfqpoint{1.849133in}{3.161548in}}%
\pgfpathlineto{\pgfqpoint{1.851406in}{3.159193in}}%
\pgfpathlineto{\pgfqpoint{1.854504in}{3.155876in}}%
\pgfpathlineto{\pgfqpoint{1.856647in}{3.155755in}}%
\pgfpathlineto{\pgfqpoint{1.858790in}{3.157380in}}%
\pgfpathlineto{\pgfqpoint{1.863805in}{3.162283in}}%
\pgfpathlineto{\pgfqpoint{1.865783in}{3.161707in}}%
\pgfpathlineto{\pgfqpoint{1.868096in}{3.159139in}}%
\pgfpathlineto{\pgfqpoint{1.870824in}{3.156202in}}%
\pgfpathlineto{\pgfqpoint{1.872967in}{3.155935in}}%
\pgfpathlineto{\pgfqpoint{1.875110in}{3.157418in}}%
\pgfpathlineto{\pgfqpoint{1.880620in}{3.162489in}}%
\pgfpathlineto{\pgfqpoint{1.882598in}{3.161650in}}%
\pgfpathlineto{\pgfqpoint{1.885022in}{3.158640in}}%
\pgfpathlineto{\pgfqpoint{1.887639in}{3.156204in}}%
\pgfpathlineto{\pgfqpoint{1.889782in}{3.156198in}}%
\pgfpathlineto{\pgfqpoint{1.891925in}{3.157934in}}%
\pgfpathlineto{\pgfqpoint{1.896609in}{3.162571in}}%
\pgfpathlineto{\pgfqpoint{1.898587in}{3.162100in}}%
\pgfpathlineto{\pgfqpoint{1.900860in}{3.159710in}}%
\pgfpathlineto{\pgfqpoint{1.903959in}{3.156347in}}%
\pgfpathlineto{\pgfqpoint{1.906102in}{3.156193in}}%
\pgfpathlineto{\pgfqpoint{1.908245in}{3.157785in}}%
\pgfpathlineto{\pgfqpoint{1.913258in}{3.162608in}}%
\pgfpathlineto{\pgfqpoint{1.915236in}{3.162001in}}%
\pgfpathlineto{\pgfqpoint{1.917550in}{3.159395in}}%
\pgfpathlineto{\pgfqpoint{1.920279in}{3.156414in}}%
\pgfpathlineto{\pgfqpoint{1.922422in}{3.156114in}}%
\pgfpathlineto{\pgfqpoint{1.924565in}{3.157564in}}%
\pgfpathlineto{\pgfqpoint{1.930071in}{3.162549in}}%
\pgfpathlineto{\pgfqpoint{1.932050in}{3.161683in}}%
\pgfpathlineto{\pgfqpoint{1.934477in}{3.158635in}}%
\pgfpathlineto{\pgfqpoint{1.937094in}{3.156161in}}%
\pgfpathlineto{\pgfqpoint{1.939237in}{3.156126in}}%
\pgfpathlineto{\pgfqpoint{1.941380in}{3.157834in}}%
\pgfpathlineto{\pgfqpoint{1.946061in}{3.162409in}}%
\pgfpathlineto{\pgfqpoint{1.948039in}{3.161916in}}%
\pgfpathlineto{\pgfqpoint{1.950315in}{3.159497in}}%
\pgfpathlineto{\pgfqpoint{1.953414in}{3.156099in}}%
\pgfpathlineto{\pgfqpoint{1.955557in}{3.155924in}}%
\pgfpathlineto{\pgfqpoint{1.957700in}{3.157495in}}%
\pgfpathlineto{\pgfqpoint{1.962709in}{3.162271in}}%
\pgfpathlineto{\pgfqpoint{1.964687in}{3.161651in}}%
\pgfpathlineto{\pgfqpoint{1.967005in}{3.159025in}}%
\pgfpathlineto{\pgfqpoint{1.969734in}{3.156026in}}%
\pgfpathlineto{\pgfqpoint{1.971877in}{3.155713in}}%
\pgfpathlineto{\pgfqpoint{1.974020in}{3.157151in}}%
\pgfpathlineto{\pgfqpoint{1.979523in}{3.162111in}}%
\pgfpathlineto{\pgfqpoint{1.981501in}{3.161240in}}%
\pgfpathlineto{\pgfqpoint{1.983931in}{3.158183in}}%
\pgfpathlineto{\pgfqpoint{1.986549in}{3.155704in}}%
\pgfpathlineto{\pgfqpoint{1.988692in}{3.155666in}}%
\pgfpathlineto{\pgfqpoint{1.990835in}{3.157372in}}%
\pgfpathlineto{\pgfqpoint{1.995513in}{3.161945in}}%
\pgfpathlineto{\pgfqpoint{1.997491in}{3.161455in}}%
\pgfpathlineto{\pgfqpoint{1.999769in}{3.159038in}}%
\pgfpathlineto{\pgfqpoint{2.002869in}{3.155646in}}%
\pgfpathlineto{\pgfqpoint{2.005012in}{3.155476in}}%
\pgfpathlineto{\pgfqpoint{2.007155in}{3.157054in}}%
\pgfpathlineto{\pgfqpoint{2.012163in}{3.161848in}}%
\pgfpathlineto{\pgfqpoint{2.014141in}{3.161237in}}%
\pgfpathlineto{\pgfqpoint{2.016459in}{3.158621in}}%
\pgfpathlineto{\pgfqpoint{2.019189in}{3.155636in}}%
\pgfpathlineto{\pgfqpoint{2.021332in}{3.155336in}}%
\pgfpathlineto{\pgfqpoint{2.023475in}{3.156787in}}%
\pgfpathlineto{\pgfqpoint{2.028978in}{3.161782in}}%
\pgfpathlineto{\pgfqpoint{2.030956in}{3.160924in}}%
\pgfpathlineto{\pgfqpoint{2.033386in}{3.157884in}}%
\pgfpathlineto{\pgfqpoint{2.036003in}{3.155425in}}%
\pgfpathlineto{\pgfqpoint{2.038146in}{3.155403in}}%
\pgfpathlineto{\pgfqpoint{2.040289in}{3.157125in}}%
\pgfpathlineto{\pgfqpoint{2.044969in}{3.161736in}}%
\pgfpathlineto{\pgfqpoint{2.046947in}{3.161261in}}%
\pgfpathlineto{\pgfqpoint{2.049224in}{3.158863in}}%
\pgfpathlineto{\pgfqpoint{2.052323in}{3.155496in}}%
\pgfpathlineto{\pgfqpoint{2.054466in}{3.155343in}}%
\pgfpathlineto{\pgfqpoint{2.056609in}{3.156938in}}%
\pgfpathlineto{\pgfqpoint{2.061619in}{3.161774in}}%
\pgfpathlineto{\pgfqpoint{2.063597in}{3.161177in}}%
\pgfpathlineto{\pgfqpoint{2.065914in}{3.158580in}}%
\pgfpathlineto{\pgfqpoint{2.068643in}{3.155617in}}%
\pgfpathlineto{\pgfqpoint{2.070786in}{3.155332in}}%
\pgfpathlineto{\pgfqpoint{2.072929in}{3.156798in}}%
\pgfpathlineto{\pgfqpoint{2.078434in}{3.161833in}}%
\pgfpathlineto{\pgfqpoint{2.080412in}{3.160987in}}%
\pgfpathlineto{\pgfqpoint{2.082841in}{3.157965in}}%
\pgfpathlineto{\pgfqpoint{2.085458in}{3.155522in}}%
\pgfpathlineto{\pgfqpoint{2.087601in}{3.155512in}}%
\pgfpathlineto{\pgfqpoint{2.089744in}{3.157246in}}%
\pgfpathlineto{\pgfqpoint{2.094425in}{3.161883in}}%
\pgfpathlineto{\pgfqpoint{2.096403in}{3.161416in}}%
\pgfpathlineto{\pgfqpoint{2.098678in}{3.159030in}}%
\pgfpathlineto{\pgfqpoint{2.101777in}{3.155676in}}%
\pgfpathlineto{\pgfqpoint{2.103921in}{3.155532in}}%
\pgfpathlineto{\pgfqpoint{2.106064in}{3.157134in}}%
\pgfpathlineto{\pgfqpoint{2.111075in}{3.161986in}}%
\pgfpathlineto{\pgfqpoint{2.113053in}{3.161393in}}%
\pgfpathlineto{\pgfqpoint{2.115368in}{3.158803in}}%
\pgfpathlineto{\pgfqpoint{2.118097in}{3.155844in}}%
\pgfpathlineto{\pgfqpoint{2.120240in}{3.155563in}}%
\pgfpathlineto{\pgfqpoint{2.122384in}{3.157032in}}%
\pgfpathlineto{\pgfqpoint{2.127890in}{3.162071in}}%
\pgfpathlineto{\pgfqpoint{2.129868in}{3.161225in}}%
\pgfpathlineto{\pgfqpoint{2.132295in}{3.158203in}}%
\pgfpathlineto{\pgfqpoint{2.134912in}{3.155759in}}%
\pgfpathlineto{\pgfqpoint{2.137055in}{3.155747in}}%
\pgfpathlineto{\pgfqpoint{2.139198in}{3.157479in}}%
\pgfpathlineto{\pgfqpoint{2.143880in}{3.162110in}}%
\pgfpathlineto{\pgfqpoint{2.145858in}{3.161640in}}%
\pgfpathlineto{\pgfqpoint{2.148133in}{3.159249in}}%
\pgfpathlineto{\pgfqpoint{2.151232in}{3.155889in}}%
\pgfpathlineto{\pgfqpoint{2.153375in}{3.155739in}}%
\pgfpathlineto{\pgfqpoint{2.155518in}{3.157335in}}%
\pgfpathlineto{\pgfqpoint{2.160530in}{3.162172in}}%
\pgfpathlineto{\pgfqpoint{2.162508in}{3.161573in}}%
\pgfpathlineto{\pgfqpoint{2.164823in}{3.158975in}}%
\pgfpathlineto{\pgfqpoint{2.167552in}{3.156007in}}%
\pgfpathlineto{\pgfqpoint{2.169695in}{3.155718in}}%
\pgfpathlineto{\pgfqpoint{2.171838in}{3.157179in}}%
\pgfpathlineto{\pgfqpoint{2.177344in}{3.162197in}}%
\pgfpathlineto{\pgfqpoint{2.179322in}{3.161343in}}%
\pgfpathlineto{\pgfqpoint{2.181750in}{3.158311in}}%
\pgfpathlineto{\pgfqpoint{2.184367in}{3.155856in}}%
\pgfpathlineto{\pgfqpoint{2.186510in}{3.155836in}}%
\pgfpathlineto{\pgfqpoint{2.188653in}{3.157559in}}%
\pgfpathlineto{\pgfqpoint{2.193334in}{3.162170in}}%
\pgfpathlineto{\pgfqpoint{2.195312in}{3.161692in}}%
\pgfpathlineto{\pgfqpoint{2.197587in}{3.159292in}}%
\pgfpathlineto{\pgfqpoint{2.200687in}{3.155918in}}%
\pgfpathlineto{\pgfqpoint{2.202830in}{3.155760in}}%
\pgfpathlineto{\pgfqpoint{2.204973in}{3.157349in}}%
\pgfpathlineto{\pgfqpoint{2.209983in}{3.162166in}}%
\pgfpathlineto{\pgfqpoint{2.211962in}{3.161560in}}%
\pgfpathlineto{\pgfqpoint{2.214278in}{3.158953in}}%
\pgfpathlineto{\pgfqpoint{2.217007in}{3.155975in}}%
\pgfpathlineto{\pgfqpoint{2.219150in}{3.155679in}}%
\pgfpathlineto{\pgfqpoint{2.221293in}{3.157133in}}%
\pgfpathlineto{\pgfqpoint{2.226798in}{3.162134in}}%
\pgfpathlineto{\pgfqpoint{2.228776in}{3.161276in}}%
\pgfpathlineto{\pgfqpoint{2.231204in}{3.158236in}}%
\pgfpathlineto{\pgfqpoint{2.233821in}{3.155775in}}%
\pgfpathlineto{\pgfqpoint{2.235964in}{3.155750in}}%
\pgfpathlineto{\pgfqpoint{2.238107in}{3.157469in}}%
\pgfpathlineto{\pgfqpoint{2.242788in}{3.162070in}}%
\pgfpathlineto{\pgfqpoint{2.244766in}{3.161590in}}%
\pgfpathlineto{\pgfqpoint{2.247042in}{3.159186in}}%
\pgfpathlineto{\pgfqpoint{2.250141in}{3.155808in}}%
\pgfpathlineto{\pgfqpoint{2.252284in}{3.155648in}}%
\pgfpathlineto{\pgfqpoint{2.254427in}{3.157234in}}%
\pgfpathlineto{\pgfqpoint{2.259437in}{3.162047in}}%
\pgfpathlineto{\pgfqpoint{2.261416in}{3.161441in}}%
\pgfpathlineto{\pgfqpoint{2.263732in}{3.158833in}}%
\pgfpathlineto{\pgfqpoint{2.266461in}{3.155855in}}%
\pgfpathlineto{\pgfqpoint{2.268604in}{3.155559in}}%
\pgfpathlineto{\pgfqpoint{2.270747in}{3.157014in}}%
\pgfpathlineto{\pgfqpoint{2.276252in}{3.162017in}}%
\pgfpathlineto{\pgfqpoint{2.278230in}{3.161159in}}%
\pgfpathlineto{\pgfqpoint{2.280659in}{3.158122in}}%
\pgfpathlineto{\pgfqpoint{2.283276in}{3.155662in}}%
\pgfpathlineto{\pgfqpoint{2.285419in}{3.155640in}}%
\pgfpathlineto{\pgfqpoint{2.287562in}{3.157361in}}%
\pgfpathlineto{\pgfqpoint{2.292242in}{3.161968in}}%
\pgfpathlineto{\pgfqpoint{2.294220in}{3.161491in}}%
\pgfpathlineto{\pgfqpoint{2.296497in}{3.159090in}}%
\pgfpathlineto{\pgfqpoint{2.299596in}{3.155717in}}%
\pgfpathlineto{\pgfqpoint{2.301739in}{3.155560in}}%
\pgfpathlineto{\pgfqpoint{2.303882in}{3.157151in}}%
\pgfpathlineto{\pgfqpoint{2.308892in}{3.161973in}}%
\pgfpathlineto{\pgfqpoint{2.310870in}{3.161370in}}%
\pgfpathlineto{\pgfqpoint{2.313187in}{3.158768in}}%
\pgfpathlineto{\pgfqpoint{2.315916in}{3.155795in}}%
\pgfpathlineto{\pgfqpoint{2.318059in}{3.155503in}}%
\pgfpathlineto{\pgfqpoint{2.320202in}{3.156962in}}%
\pgfpathlineto{\pgfqpoint{2.325707in}{3.161977in}}%
\pgfpathlineto{\pgfqpoint{2.327685in}{3.161123in}}%
\pgfpathlineto{\pgfqpoint{2.330113in}{3.158091in}}%
\pgfpathlineto{\pgfqpoint{2.332730in}{3.155637in}}%
\pgfpathlineto{\pgfqpoint{2.334873in}{3.155619in}}%
\pgfpathlineto{\pgfqpoint{2.337016in}{3.157345in}}%
\pgfpathlineto{\pgfqpoint{2.341697in}{3.161962in}}%
\pgfpathlineto{\pgfqpoint{2.343675in}{3.161487in}}%
\pgfpathlineto{\pgfqpoint{2.345951in}{3.159091in}}%
\pgfpathlineto{\pgfqpoint{2.349050in}{3.155725in}}%
\pgfpathlineto{\pgfqpoint{2.351193in}{3.155572in}}%
\pgfpathlineto{\pgfqpoint{2.353336in}{3.157166in}}%
\pgfpathlineto{\pgfqpoint{2.358347in}{3.161997in}}%
\pgfpathlineto{\pgfqpoint{2.360325in}{3.161397in}}%
\pgfpathlineto{\pgfqpoint{2.362641in}{3.158798in}}%
\pgfpathlineto{\pgfqpoint{2.365370in}{3.155829in}}%
\pgfpathlineto{\pgfqpoint{2.367513in}{3.155540in}}%
\pgfpathlineto{\pgfqpoint{2.369656in}{3.157001in}}%
\pgfpathlineto{\pgfqpoint{2.375162in}{3.162023in}}%
\pgfpathlineto{\pgfqpoint{2.377140in}{3.161171in}}%
\pgfpathlineto{\pgfqpoint{2.379568in}{3.158141in}}%
\pgfpathlineto{\pgfqpoint{2.382185in}{3.155690in}}%
\pgfpathlineto{\pgfqpoint{2.384328in}{3.155673in}}%
\pgfpathlineto{\pgfqpoint{2.386471in}{3.157400in}}%
\pgfpathlineto{\pgfqpoint{2.391152in}{3.162020in}}%
\pgfpathlineto{\pgfqpoint{2.393130in}{3.161546in}}%
\pgfpathlineto{\pgfqpoint{2.395406in}{3.159151in}}%
\pgfpathlineto{\pgfqpoint{2.398505in}{3.155785in}}%
\pgfpathlineto{\pgfqpoint{2.400648in}{3.155632in}}%
\pgfpathlineto{\pgfqpoint{2.402791in}{3.157226in}}%
\pgfpathlineto{\pgfqpoint{2.407802in}{3.162057in}}%
\pgfpathlineto{\pgfqpoint{2.409780in}{3.161457in}}%
\pgfpathlineto{\pgfqpoint{2.412096in}{3.158857in}}%
\pgfpathlineto{\pgfqpoint{2.414825in}{3.155888in}}%
\pgfpathlineto{\pgfqpoint{2.416968in}{3.155598in}}%
\pgfpathlineto{\pgfqpoint{2.419111in}{3.157059in}}%
\pgfpathlineto{\pgfqpoint{2.424616in}{3.162077in}}%
\pgfpathlineto{\pgfqpoint{2.426595in}{3.161224in}}%
\pgfpathlineto{\pgfqpoint{2.429022in}{3.158193in}}%
\pgfpathlineto{\pgfqpoint{2.431639in}{3.155740in}}%
\pgfpathlineto{\pgfqpoint{2.433782in}{3.155721in}}%
\pgfpathlineto{\pgfqpoint{2.435925in}{3.157446in}}%
\pgfpathlineto{\pgfqpoint{2.440607in}{3.162062in}}%
\pgfpathlineto{\pgfqpoint{2.442585in}{3.161587in}}%
\pgfpathlineto{\pgfqpoint{2.444860in}{3.159189in}}%
\pgfpathlineto{\pgfqpoint{2.447959in}{3.155821in}}%
\pgfpathlineto{\pgfqpoint{2.450102in}{3.155666in}}%
\pgfpathlineto{\pgfqpoint{2.452245in}{3.157257in}}%
\pgfpathlineto{\pgfqpoint{2.457256in}{3.162083in}}%
\pgfpathlineto{\pgfqpoint{2.459234in}{3.161481in}}%
\pgfpathlineto{\pgfqpoint{2.461550in}{3.158879in}}%
\pgfpathlineto{\pgfqpoint{2.464279in}{3.155906in}}%
\pgfpathlineto{\pgfqpoint{2.466422in}{3.155614in}}%
\pgfpathlineto{\pgfqpoint{2.468565in}{3.157073in}}%
\pgfpathlineto{\pgfqpoint{2.474071in}{3.162085in}}%
\pgfpathlineto{\pgfqpoint{2.476049in}{3.161230in}}%
\pgfpathlineto{\pgfqpoint{2.478477in}{3.158196in}}%
\pgfpathlineto{\pgfqpoint{2.481094in}{3.155740in}}%
\pgfpathlineto{\pgfqpoint{2.483237in}{3.155720in}}%
\pgfpathlineto{\pgfqpoint{2.485380in}{3.157443in}}%
\pgfpathlineto{\pgfqpoint{2.490061in}{3.162055in}}%
\pgfpathlineto{\pgfqpoint{2.492039in}{3.161577in}}%
\pgfpathlineto{\pgfqpoint{2.494318in}{3.159174in}}%
\pgfpathlineto{\pgfqpoint{2.497427in}{3.155818in}}%
\pgfpathlineto{\pgfqpoint{2.499571in}{3.155685in}}%
\pgfpathlineto{\pgfqpoint{2.501714in}{3.157298in}}%
\pgfpathlineto{\pgfqpoint{2.506694in}{3.162120in}}%
\pgfpathlineto{\pgfqpoint{2.508672in}{3.161534in}}%
\pgfpathlineto{\pgfqpoint{2.510980in}{3.158963in}}%
\pgfpathlineto{\pgfqpoint{2.513734in}{3.155956in}}%
\pgfpathlineto{\pgfqpoint{2.515877in}{3.155667in}}%
\pgfpathlineto{\pgfqpoint{2.518020in}{3.157129in}}%
\pgfpathlineto{\pgfqpoint{2.523526in}{3.162148in}}%
\pgfpathlineto{\pgfqpoint{2.525504in}{3.161294in}}%
\pgfpathlineto{\pgfqpoint{2.527931in}{3.158263in}}%
\pgfpathlineto{\pgfqpoint{2.530548in}{3.155809in}}%
\pgfpathlineto{\pgfqpoint{2.532691in}{3.155790in}}%
\pgfpathlineto{\pgfqpoint{2.534835in}{3.157515in}}%
\pgfpathlineto{\pgfqpoint{2.539516in}{3.162128in}}%
\pgfpathlineto{\pgfqpoint{2.541494in}{3.161651in}}%
\pgfpathlineto{\pgfqpoint{2.543769in}{3.159253in}}%
\pgfpathlineto{\pgfqpoint{2.546868in}{3.155882in}}%
\pgfpathlineto{\pgfqpoint{2.549012in}{3.155725in}}%
\pgfpathlineto{\pgfqpoint{2.551155in}{3.157315in}}%
\pgfpathlineto{\pgfqpoint{2.556165in}{3.162136in}}%
\pgfpathlineto{\pgfqpoint{2.558144in}{3.161532in}}%
\pgfpathlineto{\pgfqpoint{2.560459in}{3.158927in}}%
\pgfpathlineto{\pgfqpoint{2.563189in}{3.155952in}}%
\pgfpathlineto{\pgfqpoint{2.565332in}{3.155658in}}%
\pgfpathlineto{\pgfqpoint{2.567475in}{3.157114in}}%
\pgfpathlineto{\pgfqpoint{2.572980in}{3.162120in}}%
\pgfpathlineto{\pgfqpoint{2.574958in}{3.161262in}}%
\pgfpathlineto{\pgfqpoint{2.577386in}{3.158226in}}%
\pgfpathlineto{\pgfqpoint{2.580003in}{3.155766in}}%
\pgfpathlineto{\pgfqpoint{2.582146in}{3.155744in}}%
\pgfpathlineto{\pgfqpoint{2.584289in}{3.157464in}}%
\pgfpathlineto{\pgfqpoint{2.588970in}{3.162069in}}%
\pgfpathlineto{\pgfqpoint{2.590948in}{3.161590in}}%
\pgfpathlineto{\pgfqpoint{2.593224in}{3.159188in}}%
\pgfpathlineto{\pgfqpoint{2.596323in}{3.155813in}}%
\pgfpathlineto{\pgfqpoint{2.598466in}{3.155654in}}%
\pgfpathlineto{\pgfqpoint{2.600609in}{3.157242in}}%
\pgfpathlineto{\pgfqpoint{2.605619in}{3.162059in}}%
\pgfpathlineto{\pgfqpoint{2.607598in}{3.161453in}}%
\pgfpathlineto{\pgfqpoint{2.609914in}{3.158847in}}%
\pgfpathlineto{\pgfqpoint{2.612643in}{3.155871in}}%
\pgfpathlineto{\pgfqpoint{2.614786in}{3.155576in}}%
\pgfpathlineto{\pgfqpoint{2.616929in}{3.157031in}}%
\pgfpathlineto{\pgfqpoint{2.622434in}{3.162037in}}%
\pgfpathlineto{\pgfqpoint{2.624412in}{3.161180in}}%
\pgfpathlineto{\pgfqpoint{2.626840in}{3.158143in}}%
\pgfpathlineto{\pgfqpoint{2.629458in}{3.155685in}}%
\pgfpathlineto{\pgfqpoint{2.631601in}{3.155663in}}%
\pgfpathlineto{\pgfqpoint{2.633744in}{3.157385in}}%
\pgfpathlineto{\pgfqpoint{2.638424in}{3.161993in}}%
\pgfpathlineto{\pgfqpoint{2.640402in}{3.161515in}}%
\pgfpathlineto{\pgfqpoint{2.642678in}{3.159115in}}%
\pgfpathlineto{\pgfqpoint{2.645778in}{3.155743in}}%
\pgfpathlineto{\pgfqpoint{2.647921in}{3.155586in}}%
\pgfpathlineto{\pgfqpoint{2.650064in}{3.157176in}}%
\pgfpathlineto{\pgfqpoint{2.655074in}{3.161999in}}%
\pgfpathlineto{\pgfqpoint{2.657052in}{3.161396in}}%
\pgfpathlineto{\pgfqpoint{2.659368in}{3.158793in}}%
\pgfpathlineto{\pgfqpoint{2.662098in}{3.155819in}}%
\pgfpathlineto{\pgfqpoint{2.664241in}{3.155527in}}%
\pgfpathlineto{\pgfqpoint{2.666384in}{3.156986in}}%
\pgfpathlineto{\pgfqpoint{2.671889in}{3.161999in}}%
\pgfpathlineto{\pgfqpoint{2.673867in}{3.161145in}}%
\pgfpathlineto{\pgfqpoint{2.676295in}{3.158112in}}%
\pgfpathlineto{\pgfqpoint{2.678912in}{3.155658in}}%
\pgfpathlineto{\pgfqpoint{2.681055in}{3.155639in}}%
\pgfpathlineto{\pgfqpoint{2.683198in}{3.157364in}}%
\pgfpathlineto{\pgfqpoint{2.687879in}{3.161979in}}%
\pgfpathlineto{\pgfqpoint{2.689857in}{3.161504in}}%
\pgfpathlineto{\pgfqpoint{2.692133in}{3.159107in}}%
\pgfpathlineto{\pgfqpoint{2.695232in}{3.155739in}}%
\pgfpathlineto{\pgfqpoint{2.697375in}{3.155585in}}%
\pgfpathlineto{\pgfqpoint{2.699518in}{3.157178in}}%
\pgfpathlineto{\pgfqpoint{2.704529in}{3.162007in}}%
\pgfpathlineto{\pgfqpoint{2.706507in}{3.161406in}}%
\pgfpathlineto{\pgfqpoint{2.708823in}{3.158806in}}%
\pgfpathlineto{\pgfqpoint{2.711552in}{3.155836in}}%
\pgfpathlineto{\pgfqpoint{2.713695in}{3.155546in}}%
\pgfpathlineto{\pgfqpoint{2.715838in}{3.157006in}}%
\pgfpathlineto{\pgfqpoint{2.721343in}{3.162025in}}%
\pgfpathlineto{\pgfqpoint{2.723322in}{3.161173in}}%
\pgfpathlineto{\pgfqpoint{2.725750in}{3.158142in}}%
\pgfpathlineto{\pgfqpoint{2.728367in}{3.155689in}}%
\pgfpathlineto{\pgfqpoint{2.730510in}{3.155672in}}%
\pgfpathlineto{\pgfqpoint{2.732653in}{3.157398in}}%
\pgfpathlineto{\pgfqpoint{2.737334in}{3.162016in}}%
\pgfpathlineto{\pgfqpoint{2.739312in}{3.161541in}}%
\pgfpathlineto{\pgfqpoint{2.741587in}{3.159146in}}%
\pgfpathlineto{\pgfqpoint{2.744687in}{3.155779in}}%
\pgfpathlineto{\pgfqpoint{2.746830in}{3.155625in}}%
\pgfpathlineto{\pgfqpoint{2.748973in}{3.157219in}}%
\pgfpathlineto{\pgfqpoint{2.753984in}{3.162048in}}%
\pgfpathlineto{\pgfqpoint{2.755962in}{3.161447in}}%
\pgfpathlineto{\pgfqpoint{2.758274in}{3.158852in}}%
\pgfpathlineto{\pgfqpoint{2.761020in}{3.155886in}}%
\pgfpathlineto{\pgfqpoint{2.763163in}{3.155620in}}%
\pgfpathlineto{\pgfqpoint{2.765306in}{3.157104in}}%
\pgfpathlineto{\pgfqpoint{2.770781in}{3.162125in}}%
\pgfpathlineto{\pgfqpoint{2.772759in}{3.161290in}}%
\pgfpathlineto{\pgfqpoint{2.775204in}{3.158250in}}%
\pgfpathlineto{\pgfqpoint{2.777821in}{3.155801in}}%
\pgfpathlineto{\pgfqpoint{2.779964in}{3.155787in}}%
\pgfpathlineto{\pgfqpoint{2.782107in}{3.157515in}}%
\pgfpathlineto{\pgfqpoint{2.786789in}{3.162138in}}%
\pgfpathlineto{\pgfqpoint{2.788767in}{3.161665in}}%
\pgfpathlineto{\pgfqpoint{2.791042in}{3.159271in}}%
\pgfpathlineto{\pgfqpoint{2.794141in}{3.155905in}}%
\pgfpathlineto{\pgfqpoint{2.796284in}{3.155752in}}%
\pgfpathlineto{\pgfqpoint{2.798427in}{3.157346in}}%
\pgfpathlineto{\pgfqpoint{2.803438in}{3.162174in}}%
\pgfpathlineto{\pgfqpoint{2.805417in}{3.161572in}}%
\pgfpathlineto{\pgfqpoint{2.807732in}{3.158971in}}%
\pgfpathlineto{\pgfqpoint{2.810461in}{3.155999in}}%
\pgfpathlineto{\pgfqpoint{2.812604in}{3.155707in}}%
\pgfpathlineto{\pgfqpoint{2.814747in}{3.157166in}}%
\pgfpathlineto{\pgfqpoint{2.820253in}{3.162176in}}%
\pgfpathlineto{\pgfqpoint{2.822231in}{3.161321in}}%
\pgfpathlineto{\pgfqpoint{2.824659in}{3.158285in}}%
\pgfpathlineto{\pgfqpoint{2.827276in}{3.155828in}}%
\pgfpathlineto{\pgfqpoint{2.829419in}{3.155806in}}%
\pgfpathlineto{\pgfqpoint{2.831562in}{3.157527in}}%
\pgfpathlineto{\pgfqpoint{2.836243in}{3.162134in}}%
\pgfpathlineto{\pgfqpoint{2.838221in}{3.161655in}}%
\pgfpathlineto{\pgfqpoint{2.840497in}{3.159253in}}%
\pgfpathlineto{\pgfqpoint{2.843596in}{3.155878in}}%
\pgfpathlineto{\pgfqpoint{2.845739in}{3.155718in}}%
\pgfpathlineto{\pgfqpoint{2.847882in}{3.157306in}}%
\pgfpathlineto{\pgfqpoint{2.852892in}{3.162121in}}%
\pgfpathlineto{\pgfqpoint{2.854871in}{3.161514in}}%
\pgfpathlineto{\pgfqpoint{2.857187in}{3.158907in}}%
\pgfpathlineto{\pgfqpoint{2.859916in}{3.155929in}}%
\pgfpathlineto{\pgfqpoint{2.862059in}{3.155633in}}%
\pgfpathlineto{\pgfqpoint{2.864202in}{3.157087in}}%
\pgfpathlineto{\pgfqpoint{2.869707in}{3.162088in}}%
\pgfpathlineto{\pgfqpoint{2.871685in}{3.161230in}}%
\pgfpathlineto{\pgfqpoint{2.874113in}{3.158191in}}%
\pgfpathlineto{\pgfqpoint{2.876730in}{3.155731in}}%
\pgfpathlineto{\pgfqpoint{2.878873in}{3.155707in}}%
\pgfpathlineto{\pgfqpoint{2.881016in}{3.157426in}}%
\pgfpathlineto{\pgfqpoint{2.885697in}{3.162029in}}%
\pgfpathlineto{\pgfqpoint{2.887675in}{3.161550in}}%
\pgfpathlineto{\pgfqpoint{2.889951in}{3.159147in}}%
\pgfpathlineto{\pgfqpoint{2.893050in}{3.155771in}}%
\pgfpathlineto{\pgfqpoint{2.895193in}{3.155612in}}%
\pgfpathlineto{\pgfqpoint{2.897336in}{3.157200in}}%
\pgfpathlineto{\pgfqpoint{2.902347in}{3.162016in}}%
\pgfpathlineto{\pgfqpoint{2.904325in}{3.161411in}}%
\pgfpathlineto{\pgfqpoint{2.906641in}{3.158805in}}%
\pgfpathlineto{\pgfqpoint{2.909370in}{3.155829in}}%
\pgfpathlineto{\pgfqpoint{2.911513in}{3.155535in}}%
\pgfpathlineto{\pgfqpoint{2.913656in}{3.156991in}}%
\pgfpathlineto{\pgfqpoint{2.919161in}{3.161999in}}%
\pgfpathlineto{\pgfqpoint{2.921139in}{3.161143in}}%
\pgfpathlineto{\pgfqpoint{2.923568in}{3.158107in}}%
\pgfpathlineto{\pgfqpoint{2.926185in}{3.155650in}}%
\pgfpathlineto{\pgfqpoint{2.928328in}{3.155629in}}%
\pgfpathlineto{\pgfqpoint{2.930471in}{3.157352in}}%
\pgfpathlineto{\pgfqpoint{2.935152in}{3.161963in}}%
\pgfpathlineto{\pgfqpoint{2.937130in}{3.161487in}}%
\pgfpathlineto{\pgfqpoint{2.939406in}{3.159088in}}%
\pgfpathlineto{\pgfqpoint{2.942505in}{3.155718in}}%
\pgfpathlineto{\pgfqpoint{2.944648in}{3.155562in}}%
\pgfpathlineto{\pgfqpoint{2.946791in}{3.157154in}}%
\pgfpathlineto{\pgfqpoint{2.951801in}{3.161980in}}%
\pgfpathlineto{\pgfqpoint{2.953780in}{3.161378in}}%
\pgfpathlineto{\pgfqpoint{2.956096in}{3.158777in}}%
\pgfpathlineto{\pgfqpoint{2.958825in}{3.155806in}}%
\pgfpathlineto{\pgfqpoint{2.960968in}{3.155516in}}%
\pgfpathlineto{\pgfqpoint{2.963111in}{3.156976in}}%
\pgfpathlineto{\pgfqpoint{2.968616in}{3.161993in}}%
\pgfpathlineto{\pgfqpoint{2.970594in}{3.161140in}}%
\pgfpathlineto{\pgfqpoint{2.973022in}{3.158109in}}%
\pgfpathlineto{\pgfqpoint{2.975639in}{3.155657in}}%
\pgfpathlineto{\pgfqpoint{2.977782in}{3.155639in}}%
\pgfpathlineto{\pgfqpoint{2.979925in}{3.157365in}}%
\pgfpathlineto{\pgfqpoint{2.984606in}{3.161983in}}%
\pgfpathlineto{\pgfqpoint{2.986585in}{3.161509in}}%
\pgfpathlineto{\pgfqpoint{2.988860in}{3.159114in}}%
\pgfpathlineto{\pgfqpoint{2.991959in}{3.155748in}}%
\pgfpathlineto{\pgfqpoint{2.994102in}{3.155595in}}%
\pgfpathlineto{\pgfqpoint{2.996245in}{3.157189in}}%
\pgfpathlineto{\pgfqpoint{3.001256in}{3.162020in}}%
\pgfpathlineto{\pgfqpoint{3.003234in}{3.161420in}}%
\pgfpathlineto{\pgfqpoint{3.005550in}{3.158821in}}%
\pgfpathlineto{\pgfqpoint{3.008279in}{3.155852in}}%
\pgfpathlineto{\pgfqpoint{3.010422in}{3.155563in}}%
\pgfpathlineto{\pgfqpoint{3.012565in}{3.157024in}}%
\pgfpathlineto{\pgfqpoint{3.018071in}{3.162044in}}%
\pgfpathlineto{\pgfqpoint{3.020049in}{3.161192in}}%
\pgfpathlineto{\pgfqpoint{3.022477in}{3.158162in}}%
\pgfpathlineto{\pgfqpoint{3.025094in}{3.155710in}}%
\pgfpathlineto{\pgfqpoint{3.027237in}{3.155693in}}%
\pgfpathlineto{\pgfqpoint{3.029380in}{3.157419in}}%
\pgfpathlineto{\pgfqpoint{3.034061in}{3.162037in}}%
\pgfpathlineto{\pgfqpoint{3.036039in}{3.161563in}}%
\pgfpathlineto{\pgfqpoint{3.038315in}{3.159167in}}%
\pgfpathlineto{\pgfqpoint{3.041414in}{3.155800in}}%
\pgfpathlineto{\pgfqpoint{3.043557in}{3.155646in}}%
\pgfpathlineto{\pgfqpoint{3.045700in}{3.157240in}}%
\pgfpathlineto{\pgfqpoint{3.050711in}{3.162069in}}%
\pgfpathlineto{\pgfqpoint{3.052689in}{3.161467in}}%
\pgfpathlineto{\pgfqpoint{3.055005in}{3.158867in}}%
\pgfpathlineto{\pgfqpoint{3.057734in}{3.155896in}}%
\pgfpathlineto{\pgfqpoint{3.059877in}{3.155606in}}%
\pgfpathlineto{\pgfqpoint{3.062020in}{3.157066in}}%
\pgfpathlineto{\pgfqpoint{3.067525in}{3.162082in}}%
\pgfpathlineto{\pgfqpoint{3.069504in}{3.161228in}}%
\pgfpathlineto{\pgfqpoint{3.071931in}{3.158195in}}%
\pgfpathlineto{\pgfqpoint{3.074548in}{3.155741in}}%
\pgfpathlineto{\pgfqpoint{3.076692in}{3.155722in}}%
\pgfpathlineto{\pgfqpoint{3.078835in}{3.157446in}}%
\pgfpathlineto{\pgfqpoint{3.083516in}{3.162060in}}%
\pgfpathlineto{\pgfqpoint{3.085494in}{3.161584in}}%
\pgfpathlineto{\pgfqpoint{3.087769in}{3.159186in}}%
\pgfpathlineto{\pgfqpoint{3.090868in}{3.155816in}}%
\pgfpathlineto{\pgfqpoint{3.093012in}{3.155661in}}%
\pgfpathlineto{\pgfqpoint{3.095155in}{3.157252in}}%
\pgfpathlineto{\pgfqpoint{3.100165in}{3.162076in}}%
\pgfpathlineto{\pgfqpoint{3.102143in}{3.161473in}}%
\pgfpathlineto{\pgfqpoint{3.104459in}{3.158871in}}%
\pgfpathlineto{\pgfqpoint{3.107189in}{3.155898in}}%
\pgfpathlineto{\pgfqpoint{3.109332in}{3.155605in}}%
\pgfpathlineto{\pgfqpoint{3.111475in}{3.157063in}}%
\pgfpathlineto{\pgfqpoint{3.116980in}{3.162075in}}%
\pgfpathlineto{\pgfqpoint{3.118958in}{3.161220in}}%
\pgfpathlineto{\pgfqpoint{3.121386in}{3.158185in}}%
\pgfpathlineto{\pgfqpoint{3.124003in}{3.155729in}}%
\pgfpathlineto{\pgfqpoint{3.126146in}{3.155709in}}%
\pgfpathlineto{\pgfqpoint{3.128289in}{3.157432in}}%
\pgfpathlineto{\pgfqpoint{3.132970in}{3.162043in}}%
\pgfpathlineto{\pgfqpoint{3.134948in}{3.161566in}}%
\pgfpathlineto{\pgfqpoint{3.137224in}{3.159166in}}%
\pgfpathlineto{\pgfqpoint{3.140323in}{3.155795in}}%
\pgfpathlineto{\pgfqpoint{3.142466in}{3.155639in}}%
\pgfpathlineto{\pgfqpoint{3.144609in}{3.157229in}}%
\pgfpathlineto{\pgfqpoint{3.149620in}{3.162052in}}%
\pgfpathlineto{\pgfqpoint{3.151598in}{3.161448in}}%
\pgfpathlineto{\pgfqpoint{3.153914in}{3.158845in}}%
\pgfpathlineto{\pgfqpoint{3.156643in}{3.155871in}}%
\pgfpathlineto{\pgfqpoint{3.158786in}{3.155579in}}%
\pgfpathlineto{\pgfqpoint{3.160929in}{3.157036in}}%
\pgfpathlineto{\pgfqpoint{3.166434in}{3.162048in}}%
\pgfpathlineto{\pgfqpoint{3.168412in}{3.161193in}}%
\pgfpathlineto{\pgfqpoint{3.170132in}{3.159238in}}%
\pgfpathlineto{\pgfqpoint{3.170132in}{3.159238in}}%
\pgfusepath{stroke}%
\end{pgfscope}%
\begin{pgfscope}%
\pgfsetrectcap%
\pgfsetmiterjoin%
\pgfsetlinewidth{0.803000pt}%
\definecolor{currentstroke}{rgb}{0.000000,0.000000,0.000000}%
\pgfsetstrokecolor{currentstroke}%
\pgfsetdash{}{0pt}%
\pgfpathmoveto{\pgfqpoint{0.573768in}{2.601404in}}%
\pgfpathlineto{\pgfqpoint{0.573768in}{3.507654in}}%
\pgfusepath{stroke}%
\end{pgfscope}%
\begin{pgfscope}%
\pgfsetrectcap%
\pgfsetmiterjoin%
\pgfsetlinewidth{0.803000pt}%
\definecolor{currentstroke}{rgb}{0.000000,0.000000,0.000000}%
\pgfsetstrokecolor{currentstroke}%
\pgfsetdash{}{0pt}%
\pgfpathmoveto{\pgfqpoint{3.293768in}{2.601404in}}%
\pgfpathlineto{\pgfqpoint{3.293768in}{3.507654in}}%
\pgfusepath{stroke}%
\end{pgfscope}%
\begin{pgfscope}%
\pgfsetrectcap%
\pgfsetmiterjoin%
\pgfsetlinewidth{0.803000pt}%
\definecolor{currentstroke}{rgb}{0.000000,0.000000,0.000000}%
\pgfsetstrokecolor{currentstroke}%
\pgfsetdash{}{0pt}%
\pgfpathmoveto{\pgfqpoint{0.573768in}{2.601404in}}%
\pgfpathlineto{\pgfqpoint{3.293768in}{2.601404in}}%
\pgfusepath{stroke}%
\end{pgfscope}%
\begin{pgfscope}%
\pgfsetrectcap%
\pgfsetmiterjoin%
\pgfsetlinewidth{0.803000pt}%
\definecolor{currentstroke}{rgb}{0.000000,0.000000,0.000000}%
\pgfsetstrokecolor{currentstroke}%
\pgfsetdash{}{0pt}%
\pgfpathmoveto{\pgfqpoint{0.573768in}{3.507654in}}%
\pgfpathlineto{\pgfqpoint{3.293768in}{3.507654in}}%
\pgfusepath{stroke}%
\end{pgfscope}%
\begin{pgfscope}%
\definecolor{textcolor}{rgb}{0.000000,0.000000,0.000000}%
\pgfsetstrokecolor{textcolor}%
\pgfsetfillcolor{textcolor}%
\pgftext[x=0.628168in,y=3.444217in,left,top]{\color{textcolor}{\rmfamily\fontsize{8.000000}{9.600000}\selectfont\catcode`\^=\active\def^{\ifmmode\sp\else\^{}\fi}\catcode`\%=\active\def%{\%}(a)}}%
\end{pgfscope}%
\begin{pgfscope}%
\pgfsetbuttcap%
\pgfsetmiterjoin%
\definecolor{currentfill}{rgb}{1.000000,1.000000,1.000000}%
\pgfsetfillcolor{currentfill}%
\pgfsetlinewidth{0.000000pt}%
\definecolor{currentstroke}{rgb}{0.000000,0.000000,0.000000}%
\pgfsetstrokecolor{currentstroke}%
\pgfsetstrokeopacity{0.000000}%
\pgfsetdash{}{0pt}%
\pgfpathmoveto{\pgfqpoint{0.573768in}{1.532029in}}%
\pgfpathlineto{\pgfqpoint{3.293768in}{1.532029in}}%
\pgfpathlineto{\pgfqpoint{3.293768in}{2.438279in}}%
\pgfpathlineto{\pgfqpoint{0.573768in}{2.438279in}}%
\pgfpathlineto{\pgfqpoint{0.573768in}{1.532029in}}%
\pgfpathclose%
\pgfusepath{fill}%
\end{pgfscope}%
\begin{pgfscope}%
\pgfpathrectangle{\pgfqpoint{0.573768in}{1.532029in}}{\pgfqpoint{2.720000in}{0.906250in}}%
\pgfusepath{clip}%
\pgfsetbuttcap%
\pgfsetroundjoin%
\pgfsetlinewidth{0.501875pt}%
\definecolor{currentstroke}{rgb}{0.690196,0.690196,0.690196}%
\pgfsetstrokecolor{currentstroke}%
\pgfsetstrokeopacity{0.250000}%
\pgfsetdash{{1.850000pt}{0.800000pt}}{0.000000pt}%
\pgfpathmoveto{\pgfqpoint{0.697405in}{1.532029in}}%
\pgfpathlineto{\pgfqpoint{0.697405in}{2.438279in}}%
\pgfusepath{stroke}%
\end{pgfscope}%
\begin{pgfscope}%
\pgfsetbuttcap%
\pgfsetroundjoin%
\definecolor{currentfill}{rgb}{0.000000,0.000000,0.000000}%
\pgfsetfillcolor{currentfill}%
\pgfsetlinewidth{0.803000pt}%
\definecolor{currentstroke}{rgb}{0.000000,0.000000,0.000000}%
\pgfsetstrokecolor{currentstroke}%
\pgfsetdash{}{0pt}%
\pgfsys@defobject{currentmarker}{\pgfqpoint{0.000000in}{-0.048611in}}{\pgfqpoint{0.000000in}{0.000000in}}{%
\pgfpathmoveto{\pgfqpoint{0.000000in}{0.000000in}}%
\pgfpathlineto{\pgfqpoint{0.000000in}{-0.048611in}}%
\pgfusepath{stroke,fill}%
}%
\begin{pgfscope}%
\pgfsys@transformshift{0.697405in}{1.532029in}%
\pgfsys@useobject{currentmarker}{}%
\end{pgfscope}%
\end{pgfscope}%
\begin{pgfscope}%
\pgfpathrectangle{\pgfqpoint{0.573768in}{1.532029in}}{\pgfqpoint{2.720000in}{0.906250in}}%
\pgfusepath{clip}%
\pgfsetbuttcap%
\pgfsetroundjoin%
\pgfsetlinewidth{0.501875pt}%
\definecolor{currentstroke}{rgb}{0.690196,0.690196,0.690196}%
\pgfsetstrokecolor{currentstroke}%
\pgfsetstrokeopacity{0.250000}%
\pgfsetdash{{1.850000pt}{0.800000pt}}{0.000000pt}%
\pgfpathmoveto{\pgfqpoint{1.109526in}{1.532029in}}%
\pgfpathlineto{\pgfqpoint{1.109526in}{2.438279in}}%
\pgfusepath{stroke}%
\end{pgfscope}%
\begin{pgfscope}%
\pgfsetbuttcap%
\pgfsetroundjoin%
\definecolor{currentfill}{rgb}{0.000000,0.000000,0.000000}%
\pgfsetfillcolor{currentfill}%
\pgfsetlinewidth{0.803000pt}%
\definecolor{currentstroke}{rgb}{0.000000,0.000000,0.000000}%
\pgfsetstrokecolor{currentstroke}%
\pgfsetdash{}{0pt}%
\pgfsys@defobject{currentmarker}{\pgfqpoint{0.000000in}{-0.048611in}}{\pgfqpoint{0.000000in}{0.000000in}}{%
\pgfpathmoveto{\pgfqpoint{0.000000in}{0.000000in}}%
\pgfpathlineto{\pgfqpoint{0.000000in}{-0.048611in}}%
\pgfusepath{stroke,fill}%
}%
\begin{pgfscope}%
\pgfsys@transformshift{1.109526in}{1.532029in}%
\pgfsys@useobject{currentmarker}{}%
\end{pgfscope}%
\end{pgfscope}%
\begin{pgfscope}%
\pgfpathrectangle{\pgfqpoint{0.573768in}{1.532029in}}{\pgfqpoint{2.720000in}{0.906250in}}%
\pgfusepath{clip}%
\pgfsetbuttcap%
\pgfsetroundjoin%
\pgfsetlinewidth{0.501875pt}%
\definecolor{currentstroke}{rgb}{0.690196,0.690196,0.690196}%
\pgfsetstrokecolor{currentstroke}%
\pgfsetstrokeopacity{0.250000}%
\pgfsetdash{{1.850000pt}{0.800000pt}}{0.000000pt}%
\pgfpathmoveto{\pgfqpoint{1.521647in}{1.532029in}}%
\pgfpathlineto{\pgfqpoint{1.521647in}{2.438279in}}%
\pgfusepath{stroke}%
\end{pgfscope}%
\begin{pgfscope}%
\pgfsetbuttcap%
\pgfsetroundjoin%
\definecolor{currentfill}{rgb}{0.000000,0.000000,0.000000}%
\pgfsetfillcolor{currentfill}%
\pgfsetlinewidth{0.803000pt}%
\definecolor{currentstroke}{rgb}{0.000000,0.000000,0.000000}%
\pgfsetstrokecolor{currentstroke}%
\pgfsetdash{}{0pt}%
\pgfsys@defobject{currentmarker}{\pgfqpoint{0.000000in}{-0.048611in}}{\pgfqpoint{0.000000in}{0.000000in}}{%
\pgfpathmoveto{\pgfqpoint{0.000000in}{0.000000in}}%
\pgfpathlineto{\pgfqpoint{0.000000in}{-0.048611in}}%
\pgfusepath{stroke,fill}%
}%
\begin{pgfscope}%
\pgfsys@transformshift{1.521647in}{1.532029in}%
\pgfsys@useobject{currentmarker}{}%
\end{pgfscope}%
\end{pgfscope}%
\begin{pgfscope}%
\pgfpathrectangle{\pgfqpoint{0.573768in}{1.532029in}}{\pgfqpoint{2.720000in}{0.906250in}}%
\pgfusepath{clip}%
\pgfsetbuttcap%
\pgfsetroundjoin%
\pgfsetlinewidth{0.501875pt}%
\definecolor{currentstroke}{rgb}{0.690196,0.690196,0.690196}%
\pgfsetstrokecolor{currentstroke}%
\pgfsetstrokeopacity{0.250000}%
\pgfsetdash{{1.850000pt}{0.800000pt}}{0.000000pt}%
\pgfpathmoveto{\pgfqpoint{1.933768in}{1.532029in}}%
\pgfpathlineto{\pgfqpoint{1.933768in}{2.438279in}}%
\pgfusepath{stroke}%
\end{pgfscope}%
\begin{pgfscope}%
\pgfsetbuttcap%
\pgfsetroundjoin%
\definecolor{currentfill}{rgb}{0.000000,0.000000,0.000000}%
\pgfsetfillcolor{currentfill}%
\pgfsetlinewidth{0.803000pt}%
\definecolor{currentstroke}{rgb}{0.000000,0.000000,0.000000}%
\pgfsetstrokecolor{currentstroke}%
\pgfsetdash{}{0pt}%
\pgfsys@defobject{currentmarker}{\pgfqpoint{0.000000in}{-0.048611in}}{\pgfqpoint{0.000000in}{0.000000in}}{%
\pgfpathmoveto{\pgfqpoint{0.000000in}{0.000000in}}%
\pgfpathlineto{\pgfqpoint{0.000000in}{-0.048611in}}%
\pgfusepath{stroke,fill}%
}%
\begin{pgfscope}%
\pgfsys@transformshift{1.933768in}{1.532029in}%
\pgfsys@useobject{currentmarker}{}%
\end{pgfscope}%
\end{pgfscope}%
\begin{pgfscope}%
\pgfpathrectangle{\pgfqpoint{0.573768in}{1.532029in}}{\pgfqpoint{2.720000in}{0.906250in}}%
\pgfusepath{clip}%
\pgfsetbuttcap%
\pgfsetroundjoin%
\pgfsetlinewidth{0.501875pt}%
\definecolor{currentstroke}{rgb}{0.690196,0.690196,0.690196}%
\pgfsetstrokecolor{currentstroke}%
\pgfsetstrokeopacity{0.250000}%
\pgfsetdash{{1.850000pt}{0.800000pt}}{0.000000pt}%
\pgfpathmoveto{\pgfqpoint{2.345890in}{1.532029in}}%
\pgfpathlineto{\pgfqpoint{2.345890in}{2.438279in}}%
\pgfusepath{stroke}%
\end{pgfscope}%
\begin{pgfscope}%
\pgfsetbuttcap%
\pgfsetroundjoin%
\definecolor{currentfill}{rgb}{0.000000,0.000000,0.000000}%
\pgfsetfillcolor{currentfill}%
\pgfsetlinewidth{0.803000pt}%
\definecolor{currentstroke}{rgb}{0.000000,0.000000,0.000000}%
\pgfsetstrokecolor{currentstroke}%
\pgfsetdash{}{0pt}%
\pgfsys@defobject{currentmarker}{\pgfqpoint{0.000000in}{-0.048611in}}{\pgfqpoint{0.000000in}{0.000000in}}{%
\pgfpathmoveto{\pgfqpoint{0.000000in}{0.000000in}}%
\pgfpathlineto{\pgfqpoint{0.000000in}{-0.048611in}}%
\pgfusepath{stroke,fill}%
}%
\begin{pgfscope}%
\pgfsys@transformshift{2.345890in}{1.532029in}%
\pgfsys@useobject{currentmarker}{}%
\end{pgfscope}%
\end{pgfscope}%
\begin{pgfscope}%
\pgfpathrectangle{\pgfqpoint{0.573768in}{1.532029in}}{\pgfqpoint{2.720000in}{0.906250in}}%
\pgfusepath{clip}%
\pgfsetbuttcap%
\pgfsetroundjoin%
\pgfsetlinewidth{0.501875pt}%
\definecolor{currentstroke}{rgb}{0.690196,0.690196,0.690196}%
\pgfsetstrokecolor{currentstroke}%
\pgfsetstrokeopacity{0.250000}%
\pgfsetdash{{1.850000pt}{0.800000pt}}{0.000000pt}%
\pgfpathmoveto{\pgfqpoint{2.758011in}{1.532029in}}%
\pgfpathlineto{\pgfqpoint{2.758011in}{2.438279in}}%
\pgfusepath{stroke}%
\end{pgfscope}%
\begin{pgfscope}%
\pgfsetbuttcap%
\pgfsetroundjoin%
\definecolor{currentfill}{rgb}{0.000000,0.000000,0.000000}%
\pgfsetfillcolor{currentfill}%
\pgfsetlinewidth{0.803000pt}%
\definecolor{currentstroke}{rgb}{0.000000,0.000000,0.000000}%
\pgfsetstrokecolor{currentstroke}%
\pgfsetdash{}{0pt}%
\pgfsys@defobject{currentmarker}{\pgfqpoint{0.000000in}{-0.048611in}}{\pgfqpoint{0.000000in}{0.000000in}}{%
\pgfpathmoveto{\pgfqpoint{0.000000in}{0.000000in}}%
\pgfpathlineto{\pgfqpoint{0.000000in}{-0.048611in}}%
\pgfusepath{stroke,fill}%
}%
\begin{pgfscope}%
\pgfsys@transformshift{2.758011in}{1.532029in}%
\pgfsys@useobject{currentmarker}{}%
\end{pgfscope}%
\end{pgfscope}%
\begin{pgfscope}%
\pgfpathrectangle{\pgfqpoint{0.573768in}{1.532029in}}{\pgfqpoint{2.720000in}{0.906250in}}%
\pgfusepath{clip}%
\pgfsetbuttcap%
\pgfsetroundjoin%
\pgfsetlinewidth{0.501875pt}%
\definecolor{currentstroke}{rgb}{0.690196,0.690196,0.690196}%
\pgfsetstrokecolor{currentstroke}%
\pgfsetstrokeopacity{0.250000}%
\pgfsetdash{{1.850000pt}{0.800000pt}}{0.000000pt}%
\pgfpathmoveto{\pgfqpoint{3.170132in}{1.532029in}}%
\pgfpathlineto{\pgfqpoint{3.170132in}{2.438279in}}%
\pgfusepath{stroke}%
\end{pgfscope}%
\begin{pgfscope}%
\pgfsetbuttcap%
\pgfsetroundjoin%
\definecolor{currentfill}{rgb}{0.000000,0.000000,0.000000}%
\pgfsetfillcolor{currentfill}%
\pgfsetlinewidth{0.803000pt}%
\definecolor{currentstroke}{rgb}{0.000000,0.000000,0.000000}%
\pgfsetstrokecolor{currentstroke}%
\pgfsetdash{}{0pt}%
\pgfsys@defobject{currentmarker}{\pgfqpoint{0.000000in}{-0.048611in}}{\pgfqpoint{0.000000in}{0.000000in}}{%
\pgfpathmoveto{\pgfqpoint{0.000000in}{0.000000in}}%
\pgfpathlineto{\pgfqpoint{0.000000in}{-0.048611in}}%
\pgfusepath{stroke,fill}%
}%
\begin{pgfscope}%
\pgfsys@transformshift{3.170132in}{1.532029in}%
\pgfsys@useobject{currentmarker}{}%
\end{pgfscope}%
\end{pgfscope}%
\begin{pgfscope}%
\pgfpathrectangle{\pgfqpoint{0.573768in}{1.532029in}}{\pgfqpoint{2.720000in}{0.906250in}}%
\pgfusepath{clip}%
\pgfsetbuttcap%
\pgfsetroundjoin%
\pgfsetlinewidth{0.501875pt}%
\definecolor{currentstroke}{rgb}{0.690196,0.690196,0.690196}%
\pgfsetstrokecolor{currentstroke}%
\pgfsetstrokeopacity{0.250000}%
\pgfsetdash{{1.850000pt}{0.800000pt}}{0.000000pt}%
\pgfpathmoveto{\pgfqpoint{0.573768in}{1.532029in}}%
\pgfpathlineto{\pgfqpoint{3.293768in}{1.532029in}}%
\pgfusepath{stroke}%
\end{pgfscope}%
\begin{pgfscope}%
\pgfsetbuttcap%
\pgfsetroundjoin%
\definecolor{currentfill}{rgb}{0.000000,0.000000,0.000000}%
\pgfsetfillcolor{currentfill}%
\pgfsetlinewidth{0.803000pt}%
\definecolor{currentstroke}{rgb}{0.000000,0.000000,0.000000}%
\pgfsetstrokecolor{currentstroke}%
\pgfsetdash{}{0pt}%
\pgfsys@defobject{currentmarker}{\pgfqpoint{-0.048611in}{0.000000in}}{\pgfqpoint{-0.000000in}{0.000000in}}{%
\pgfpathmoveto{\pgfqpoint{-0.000000in}{0.000000in}}%
\pgfpathlineto{\pgfqpoint{-0.048611in}{0.000000in}}%
\pgfusepath{stroke,fill}%
}%
\begin{pgfscope}%
\pgfsys@transformshift{0.573768in}{1.532029in}%
\pgfsys@useobject{currentmarker}{}%
\end{pgfscope}%
\end{pgfscope}%
\begin{pgfscope}%
\definecolor{textcolor}{rgb}{0.000000,0.000000,0.000000}%
\pgfsetstrokecolor{textcolor}%
\pgfsetfillcolor{textcolor}%
\pgftext[x=0.417518in, y=1.493449in, left, base]{\color{textcolor}{\rmfamily\fontsize{8.000000}{9.600000}\selectfont\catcode`\^=\active\def^{\ifmmode\sp\else\^{}\fi}\catcode`\%=\active\def%{\%}$\mathdefault{0}$}}%
\end{pgfscope}%
\begin{pgfscope}%
\pgfpathrectangle{\pgfqpoint{0.573768in}{1.532029in}}{\pgfqpoint{2.720000in}{0.906250in}}%
\pgfusepath{clip}%
\pgfsetbuttcap%
\pgfsetroundjoin%
\pgfsetlinewidth{0.501875pt}%
\definecolor{currentstroke}{rgb}{0.690196,0.690196,0.690196}%
\pgfsetstrokecolor{currentstroke}%
\pgfsetstrokeopacity{0.250000}%
\pgfsetdash{{1.850000pt}{0.800000pt}}{0.000000pt}%
\pgfpathmoveto{\pgfqpoint{0.573768in}{1.834113in}}%
\pgfpathlineto{\pgfqpoint{3.293768in}{1.834113in}}%
\pgfusepath{stroke}%
\end{pgfscope}%
\begin{pgfscope}%
\pgfsetbuttcap%
\pgfsetroundjoin%
\definecolor{currentfill}{rgb}{0.000000,0.000000,0.000000}%
\pgfsetfillcolor{currentfill}%
\pgfsetlinewidth{0.803000pt}%
\definecolor{currentstroke}{rgb}{0.000000,0.000000,0.000000}%
\pgfsetstrokecolor{currentstroke}%
\pgfsetdash{}{0pt}%
\pgfsys@defobject{currentmarker}{\pgfqpoint{-0.048611in}{0.000000in}}{\pgfqpoint{-0.000000in}{0.000000in}}{%
\pgfpathmoveto{\pgfqpoint{-0.000000in}{0.000000in}}%
\pgfpathlineto{\pgfqpoint{-0.048611in}{0.000000in}}%
\pgfusepath{stroke,fill}%
}%
\begin{pgfscope}%
\pgfsys@transformshift{0.573768in}{1.834113in}%
\pgfsys@useobject{currentmarker}{}%
\end{pgfscope}%
\end{pgfscope}%
\begin{pgfscope}%
\definecolor{textcolor}{rgb}{0.000000,0.000000,0.000000}%
\pgfsetstrokecolor{textcolor}%
\pgfsetfillcolor{textcolor}%
\pgftext[x=0.417518in, y=1.795533in, left, base]{\color{textcolor}{\rmfamily\fontsize{8.000000}{9.600000}\selectfont\catcode`\^=\active\def^{\ifmmode\sp\else\^{}\fi}\catcode`\%=\active\def%{\%}$\mathdefault{5}$}}%
\end{pgfscope}%
\begin{pgfscope}%
\pgfpathrectangle{\pgfqpoint{0.573768in}{1.532029in}}{\pgfqpoint{2.720000in}{0.906250in}}%
\pgfusepath{clip}%
\pgfsetbuttcap%
\pgfsetroundjoin%
\pgfsetlinewidth{0.501875pt}%
\definecolor{currentstroke}{rgb}{0.690196,0.690196,0.690196}%
\pgfsetstrokecolor{currentstroke}%
\pgfsetstrokeopacity{0.250000}%
\pgfsetdash{{1.850000pt}{0.800000pt}}{0.000000pt}%
\pgfpathmoveto{\pgfqpoint{0.573768in}{2.136196in}}%
\pgfpathlineto{\pgfqpoint{3.293768in}{2.136196in}}%
\pgfusepath{stroke}%
\end{pgfscope}%
\begin{pgfscope}%
\pgfsetbuttcap%
\pgfsetroundjoin%
\definecolor{currentfill}{rgb}{0.000000,0.000000,0.000000}%
\pgfsetfillcolor{currentfill}%
\pgfsetlinewidth{0.803000pt}%
\definecolor{currentstroke}{rgb}{0.000000,0.000000,0.000000}%
\pgfsetstrokecolor{currentstroke}%
\pgfsetdash{}{0pt}%
\pgfsys@defobject{currentmarker}{\pgfqpoint{-0.048611in}{0.000000in}}{\pgfqpoint{-0.000000in}{0.000000in}}{%
\pgfpathmoveto{\pgfqpoint{-0.000000in}{0.000000in}}%
\pgfpathlineto{\pgfqpoint{-0.048611in}{0.000000in}}%
\pgfusepath{stroke,fill}%
}%
\begin{pgfscope}%
\pgfsys@transformshift{0.573768in}{2.136196in}%
\pgfsys@useobject{currentmarker}{}%
\end{pgfscope}%
\end{pgfscope}%
\begin{pgfscope}%
\definecolor{textcolor}{rgb}{0.000000,0.000000,0.000000}%
\pgfsetstrokecolor{textcolor}%
\pgfsetfillcolor{textcolor}%
\pgftext[x=0.358489in, y=2.097616in, left, base]{\color{textcolor}{\rmfamily\fontsize{8.000000}{9.600000}\selectfont\catcode`\^=\active\def^{\ifmmode\sp\else\^{}\fi}\catcode`\%=\active\def%{\%}$\mathdefault{10}$}}%
\end{pgfscope}%
\begin{pgfscope}%
\pgfpathrectangle{\pgfqpoint{0.573768in}{1.532029in}}{\pgfqpoint{2.720000in}{0.906250in}}%
\pgfusepath{clip}%
\pgfsetbuttcap%
\pgfsetroundjoin%
\pgfsetlinewidth{0.501875pt}%
\definecolor{currentstroke}{rgb}{0.690196,0.690196,0.690196}%
\pgfsetstrokecolor{currentstroke}%
\pgfsetstrokeopacity{0.250000}%
\pgfsetdash{{1.850000pt}{0.800000pt}}{0.000000pt}%
\pgfpathmoveto{\pgfqpoint{0.573768in}{2.438279in}}%
\pgfpathlineto{\pgfqpoint{3.293768in}{2.438279in}}%
\pgfusepath{stroke}%
\end{pgfscope}%
\begin{pgfscope}%
\pgfsetbuttcap%
\pgfsetroundjoin%
\definecolor{currentfill}{rgb}{0.000000,0.000000,0.000000}%
\pgfsetfillcolor{currentfill}%
\pgfsetlinewidth{0.803000pt}%
\definecolor{currentstroke}{rgb}{0.000000,0.000000,0.000000}%
\pgfsetstrokecolor{currentstroke}%
\pgfsetdash{}{0pt}%
\pgfsys@defobject{currentmarker}{\pgfqpoint{-0.048611in}{0.000000in}}{\pgfqpoint{-0.000000in}{0.000000in}}{%
\pgfpathmoveto{\pgfqpoint{-0.000000in}{0.000000in}}%
\pgfpathlineto{\pgfqpoint{-0.048611in}{0.000000in}}%
\pgfusepath{stroke,fill}%
}%
\begin{pgfscope}%
\pgfsys@transformshift{0.573768in}{2.438279in}%
\pgfsys@useobject{currentmarker}{}%
\end{pgfscope}%
\end{pgfscope}%
\begin{pgfscope}%
\definecolor{textcolor}{rgb}{0.000000,0.000000,0.000000}%
\pgfsetstrokecolor{textcolor}%
\pgfsetfillcolor{textcolor}%
\pgftext[x=0.358489in, y=2.399699in, left, base]{\color{textcolor}{\rmfamily\fontsize{8.000000}{9.600000}\selectfont\catcode`\^=\active\def^{\ifmmode\sp\else\^{}\fi}\catcode`\%=\active\def%{\%}$\mathdefault{15}$}}%
\end{pgfscope}%
\begin{pgfscope}%
\definecolor{textcolor}{rgb}{0.000000,0.000000,0.000000}%
\pgfsetstrokecolor{textcolor}%
\pgfsetfillcolor{textcolor}%
\pgftext[x=0.302933in,y=1.985154in,,bottom,rotate=90.000000]{\color{textcolor}{\rmfamily\fontsize{8.000000}{9.600000}\selectfont\catcode`\^=\active\def^{\ifmmode\sp\else\^{}\fi}\catcode`\%=\active\def%{\%}$i_L$ (A)}}%
\end{pgfscope}%
\begin{pgfscope}%
\pgfpathrectangle{\pgfqpoint{0.573768in}{1.532029in}}{\pgfqpoint{2.720000in}{0.906250in}}%
\pgfusepath{clip}%
\pgfsetrectcap%
\pgfsetroundjoin%
\pgfsetlinewidth{1.505625pt}%
\definecolor{currentstroke}{rgb}{0.835294,0.368627,0.000000}%
\pgfsetstrokecolor{currentstroke}%
\pgfsetdash{}{0pt}%
\pgfpathmoveto{\pgfqpoint{0.697405in}{1.532029in}}%
\pgfpathlineto{\pgfqpoint{0.697788in}{1.532600in}}%
\pgfpathlineto{\pgfqpoint{0.701605in}{1.643985in}}%
\pgfpathlineto{\pgfqpoint{0.706139in}{1.773317in}}%
\pgfpathlineto{\pgfqpoint{0.706662in}{1.771687in}}%
\pgfpathlineto{\pgfqpoint{0.713091in}{1.748863in}}%
\pgfpathlineto{\pgfqpoint{0.714218in}{1.744384in}}%
\pgfpathlineto{\pgfqpoint{0.714351in}{1.746877in}}%
\pgfpathlineto{\pgfqpoint{0.722635in}{1.969093in}}%
\pgfpathlineto{\pgfqpoint{0.723311in}{1.965344in}}%
\pgfpathlineto{\pgfqpoint{0.729905in}{1.923865in}}%
\pgfpathlineto{\pgfqpoint{0.730704in}{1.918305in}}%
\pgfpathlineto{\pgfqpoint{0.730898in}{1.921670in}}%
\pgfpathlineto{\pgfqpoint{0.739056in}{2.115491in}}%
\pgfpathlineto{\pgfqpoint{0.739468in}{2.112333in}}%
\pgfpathlineto{\pgfqpoint{0.747214in}{2.035434in}}%
\pgfpathlineto{\pgfqpoint{0.747675in}{2.044145in}}%
\pgfpathlineto{\pgfqpoint{0.755541in}{2.200807in}}%
\pgfpathlineto{\pgfqpoint{0.755809in}{2.197912in}}%
\pgfpathlineto{\pgfqpoint{0.763698in}{2.089471in}}%
\pgfpathlineto{\pgfqpoint{0.764152in}{2.096263in}}%
\pgfpathlineto{\pgfqpoint{0.772026in}{2.223878in}}%
\pgfpathlineto{\pgfqpoint{0.772223in}{2.221181in}}%
\pgfpathlineto{\pgfqpoint{0.780195in}{2.084262in}}%
\pgfpathlineto{\pgfqpoint{0.780802in}{2.091854in}}%
\pgfpathlineto{\pgfqpoint{0.788511in}{2.192792in}}%
\pgfpathlineto{\pgfqpoint{0.788651in}{2.190712in}}%
\pgfpathlineto{\pgfqpoint{0.796680in}{2.031467in}}%
\pgfpathlineto{\pgfqpoint{0.797290in}{2.037651in}}%
\pgfpathlineto{\pgfqpoint{0.804996in}{2.121993in}}%
\pgfpathlineto{\pgfqpoint{0.805140in}{2.119555in}}%
\pgfpathlineto{\pgfqpoint{0.813166in}{1.947332in}}%
\pgfpathlineto{\pgfqpoint{0.813716in}{1.952143in}}%
\pgfpathlineto{\pgfqpoint{0.821480in}{2.028852in}}%
\pgfpathlineto{\pgfqpoint{0.821623in}{2.026320in}}%
\pgfpathlineto{\pgfqpoint{0.829654in}{1.849684in}}%
\pgfpathlineto{\pgfqpoint{0.830155in}{1.853965in}}%
\pgfpathlineto{\pgfqpoint{0.837964in}{1.930852in}}%
\pgfpathlineto{\pgfqpoint{0.838186in}{1.926648in}}%
\pgfpathlineto{\pgfqpoint{0.846141in}{1.755066in}}%
\pgfpathlineto{\pgfqpoint{0.846578in}{1.759029in}}%
\pgfpathlineto{\pgfqpoint{0.854449in}{1.842934in}}%
\pgfpathlineto{\pgfqpoint{0.854759in}{1.837145in}}%
\pgfpathlineto{\pgfqpoint{0.862608in}{1.676478in}}%
\pgfpathlineto{\pgfqpoint{0.863020in}{1.680545in}}%
\pgfpathlineto{\pgfqpoint{0.870958in}{1.775931in}}%
\pgfpathlineto{\pgfqpoint{0.871263in}{1.770594in}}%
\pgfpathlineto{\pgfqpoint{0.879086in}{1.622645in}}%
\pgfpathlineto{\pgfqpoint{0.879457in}{1.626824in}}%
\pgfpathlineto{\pgfqpoint{0.887417in}{1.735973in}}%
\pgfpathlineto{\pgfqpoint{0.887735in}{1.731091in}}%
\pgfpathlineto{\pgfqpoint{0.895578in}{1.596157in}}%
\pgfpathlineto{\pgfqpoint{0.895889in}{1.600172in}}%
\pgfpathlineto{\pgfqpoint{0.903927in}{1.723298in}}%
\pgfpathlineto{\pgfqpoint{0.904251in}{1.718720in}}%
\pgfpathlineto{\pgfqpoint{0.912063in}{1.596392in}}%
\pgfpathlineto{\pgfqpoint{0.912373in}{1.600860in}}%
\pgfpathlineto{\pgfqpoint{0.920413in}{1.735191in}}%
\pgfpathlineto{\pgfqpoint{0.920880in}{1.728864in}}%
\pgfpathlineto{\pgfqpoint{0.928533in}{1.617993in}}%
\pgfpathlineto{\pgfqpoint{0.928842in}{1.622542in}}%
\pgfpathlineto{\pgfqpoint{0.936872in}{1.764917in}}%
\pgfpathlineto{\pgfqpoint{0.937356in}{1.758924in}}%
\pgfpathlineto{\pgfqpoint{0.945017in}{1.653636in}}%
\pgfpathlineto{\pgfqpoint{0.945323in}{1.658288in}}%
\pgfpathlineto{\pgfqpoint{0.953337in}{1.804731in}}%
\pgfpathlineto{\pgfqpoint{0.953790in}{1.799386in}}%
\pgfpathlineto{\pgfqpoint{0.961506in}{1.695527in}}%
\pgfpathlineto{\pgfqpoint{0.961811in}{1.700204in}}%
\pgfpathlineto{\pgfqpoint{0.969817in}{1.847044in}}%
\pgfpathlineto{\pgfqpoint{0.970267in}{1.841728in}}%
\pgfpathlineto{\pgfqpoint{0.977995in}{1.736673in}}%
\pgfpathlineto{\pgfqpoint{0.978358in}{1.742402in}}%
\pgfpathlineto{\pgfqpoint{0.986316in}{1.885742in}}%
\pgfpathlineto{\pgfqpoint{0.986623in}{1.882194in}}%
\pgfpathlineto{\pgfqpoint{0.994486in}{1.771756in}}%
\pgfpathlineto{\pgfqpoint{0.994826in}{1.776933in}}%
\pgfpathlineto{\pgfqpoint{1.002813in}{1.916019in}}%
\pgfpathlineto{\pgfqpoint{1.003136in}{1.912018in}}%
\pgfpathlineto{\pgfqpoint{1.010971in}{1.796760in}}%
\pgfpathlineto{\pgfqpoint{1.011317in}{1.801798in}}%
\pgfpathlineto{\pgfqpoint{1.019298in}{1.935076in}}%
\pgfpathlineto{\pgfqpoint{1.019627in}{1.930737in}}%
\pgfpathlineto{\pgfqpoint{1.027456in}{1.809833in}}%
\pgfpathlineto{\pgfqpoint{1.027888in}{1.815995in}}%
\pgfpathlineto{\pgfqpoint{1.035782in}{1.942125in}}%
\pgfpathlineto{\pgfqpoint{1.036112in}{1.937551in}}%
\pgfpathlineto{\pgfqpoint{1.043941in}{1.811287in}}%
\pgfpathlineto{\pgfqpoint{1.044343in}{1.816706in}}%
\pgfpathlineto{\pgfqpoint{1.052292in}{1.938437in}}%
\pgfpathlineto{\pgfqpoint{1.052601in}{1.933862in}}%
\pgfpathlineto{\pgfqpoint{1.060426in}{1.803369in}}%
\pgfpathlineto{\pgfqpoint{1.060839in}{1.808753in}}%
\pgfpathlineto{\pgfqpoint{1.068752in}{1.926612in}}%
\pgfpathlineto{\pgfqpoint{1.069082in}{1.921711in}}%
\pgfpathlineto{\pgfqpoint{1.076911in}{1.788351in}}%
\pgfpathlineto{\pgfqpoint{1.077363in}{1.794192in}}%
\pgfpathlineto{\pgfqpoint{1.085237in}{1.909466in}}%
\pgfpathlineto{\pgfqpoint{1.085560in}{1.904616in}}%
\pgfpathlineto{\pgfqpoint{1.093395in}{1.769963in}}%
\pgfpathlineto{\pgfqpoint{1.093751in}{1.774363in}}%
\pgfpathlineto{\pgfqpoint{1.101722in}{1.890641in}}%
\pgfpathlineto{\pgfqpoint{1.102044in}{1.885800in}}%
\pgfpathlineto{\pgfqpoint{1.109876in}{1.751457in}}%
\pgfpathlineto{\pgfqpoint{1.110235in}{1.755902in}}%
\pgfpathlineto{\pgfqpoint{1.118206in}{1.873157in}}%
\pgfpathlineto{\pgfqpoint{1.118528in}{1.868376in}}%
\pgfpathlineto{\pgfqpoint{1.126359in}{1.735524in}}%
\pgfpathlineto{\pgfqpoint{1.126721in}{1.740076in}}%
\pgfpathlineto{\pgfqpoint{1.134672in}{1.859206in}}%
\pgfpathlineto{\pgfqpoint{1.134977in}{1.854839in}}%
\pgfpathlineto{\pgfqpoint{1.142839in}{1.723828in}}%
\pgfpathlineto{\pgfqpoint{1.143204in}{1.728517in}}%
\pgfpathlineto{\pgfqpoint{1.151152in}{1.850116in}}%
\pgfpathlineto{\pgfqpoint{1.151453in}{1.845931in}}%
\pgfpathlineto{\pgfqpoint{1.159327in}{1.717136in}}%
\pgfpathlineto{\pgfqpoint{1.159664in}{1.721516in}}%
\pgfpathlineto{\pgfqpoint{1.167657in}{1.846188in}}%
\pgfpathlineto{\pgfqpoint{1.167973in}{1.841805in}}%
\pgfpathlineto{\pgfqpoint{1.175811in}{1.716018in}}%
\pgfpathlineto{\pgfqpoint{1.176158in}{1.720644in}}%
\pgfpathlineto{\pgfqpoint{1.184145in}{1.847461in}}%
\pgfpathlineto{\pgfqpoint{1.184466in}{1.843081in}}%
\pgfpathlineto{\pgfqpoint{1.192297in}{1.719394in}}%
\pgfpathlineto{\pgfqpoint{1.192650in}{1.724224in}}%
\pgfpathlineto{\pgfqpoint{1.200604in}{1.852460in}}%
\pgfpathlineto{\pgfqpoint{1.200905in}{1.848534in}}%
\pgfpathlineto{\pgfqpoint{1.208778in}{1.725636in}}%
\pgfpathlineto{\pgfqpoint{1.209142in}{1.730673in}}%
\pgfpathlineto{\pgfqpoint{1.217092in}{1.859807in}}%
\pgfpathlineto{\pgfqpoint{1.217395in}{1.855862in}}%
\pgfpathlineto{\pgfqpoint{1.225268in}{1.733460in}}%
\pgfpathlineto{\pgfqpoint{1.225629in}{1.738494in}}%
\pgfpathlineto{\pgfqpoint{1.233581in}{1.867902in}}%
\pgfpathlineto{\pgfqpoint{1.233887in}{1.863894in}}%
\pgfpathlineto{\pgfqpoint{1.241754in}{1.741548in}}%
\pgfpathlineto{\pgfqpoint{1.242113in}{1.746542in}}%
\pgfpathlineto{\pgfqpoint{1.250074in}{1.875667in}}%
\pgfpathlineto{\pgfqpoint{1.250383in}{1.871577in}}%
\pgfpathlineto{\pgfqpoint{1.258239in}{1.748813in}}%
\pgfpathlineto{\pgfqpoint{1.258599in}{1.753781in}}%
\pgfpathlineto{\pgfqpoint{1.266570in}{1.882128in}}%
\pgfpathlineto{\pgfqpoint{1.266892in}{1.877764in}}%
\pgfpathlineto{\pgfqpoint{1.274729in}{1.754433in}}%
\pgfpathlineto{\pgfqpoint{1.275073in}{1.759152in}}%
\pgfpathlineto{\pgfqpoint{1.283055in}{1.886592in}}%
\pgfpathlineto{\pgfqpoint{1.283377in}{1.882180in}}%
\pgfpathlineto{\pgfqpoint{1.291214in}{1.757727in}}%
\pgfpathlineto{\pgfqpoint{1.291590in}{1.762924in}}%
\pgfpathlineto{\pgfqpoint{1.299540in}{1.888687in}}%
\pgfpathlineto{\pgfqpoint{1.299862in}{1.884228in}}%
\pgfpathlineto{\pgfqpoint{1.307698in}{1.758678in}}%
\pgfpathlineto{\pgfqpoint{1.308075in}{1.763814in}}%
\pgfpathlineto{\pgfqpoint{1.316025in}{1.888546in}}%
\pgfpathlineto{\pgfqpoint{1.316347in}{1.884046in}}%
\pgfpathlineto{\pgfqpoint{1.324183in}{1.757575in}}%
\pgfpathlineto{\pgfqpoint{1.324560in}{1.762681in}}%
\pgfpathlineto{\pgfqpoint{1.332509in}{1.886599in}}%
\pgfpathlineto{\pgfqpoint{1.332832in}{1.882071in}}%
\pgfpathlineto{\pgfqpoint{1.340668in}{1.754955in}}%
\pgfpathlineto{\pgfqpoint{1.341011in}{1.759498in}}%
\pgfpathlineto{\pgfqpoint{1.348994in}{1.883453in}}%
\pgfpathlineto{\pgfqpoint{1.349317in}{1.878910in}}%
\pgfpathlineto{\pgfqpoint{1.357149in}{1.751461in}}%
\pgfpathlineto{\pgfqpoint{1.357508in}{1.756208in}}%
\pgfpathlineto{\pgfqpoint{1.365479in}{1.879800in}}%
\pgfpathlineto{\pgfqpoint{1.365801in}{1.875255in}}%
\pgfpathlineto{\pgfqpoint{1.373633in}{1.747785in}}%
\pgfpathlineto{\pgfqpoint{1.373992in}{1.752529in}}%
\pgfpathlineto{\pgfqpoint{1.381964in}{1.876244in}}%
\pgfpathlineto{\pgfqpoint{1.382286in}{1.871709in}}%
\pgfpathlineto{\pgfqpoint{1.390118in}{1.744469in}}%
\pgfpathlineto{\pgfqpoint{1.390477in}{1.749223in}}%
\pgfpathlineto{\pgfqpoint{1.398442in}{1.873257in}}%
\pgfpathlineto{\pgfqpoint{1.398756in}{1.868889in}}%
\pgfpathlineto{\pgfqpoint{1.406602in}{1.741815in}}%
\pgfpathlineto{\pgfqpoint{1.406961in}{1.746588in}}%
\pgfpathlineto{\pgfqpoint{1.414922in}{1.871077in}}%
\pgfpathlineto{\pgfqpoint{1.415229in}{1.866841in}}%
\pgfpathlineto{\pgfqpoint{1.423087in}{1.740081in}}%
\pgfpathlineto{\pgfqpoint{1.423451in}{1.744964in}}%
\pgfpathlineto{\pgfqpoint{1.431406in}{1.869889in}}%
\pgfpathlineto{\pgfqpoint{1.431712in}{1.865698in}}%
\pgfpathlineto{\pgfqpoint{1.439572in}{1.739405in}}%
\pgfpathlineto{\pgfqpoint{1.439935in}{1.744300in}}%
\pgfpathlineto{\pgfqpoint{1.447890in}{1.869721in}}%
\pgfpathlineto{\pgfqpoint{1.448195in}{1.865567in}}%
\pgfpathlineto{\pgfqpoint{1.456056in}{1.739685in}}%
\pgfpathlineto{\pgfqpoint{1.456421in}{1.744603in}}%
\pgfpathlineto{\pgfqpoint{1.464375in}{1.870416in}}%
\pgfpathlineto{\pgfqpoint{1.464681in}{1.866264in}}%
\pgfpathlineto{\pgfqpoint{1.472541in}{1.740726in}}%
\pgfpathlineto{\pgfqpoint{1.472907in}{1.745677in}}%
\pgfpathlineto{\pgfqpoint{1.480862in}{1.871736in}}%
\pgfpathlineto{\pgfqpoint{1.481169in}{1.867555in}}%
\pgfpathlineto{\pgfqpoint{1.489026in}{1.742263in}}%
\pgfpathlineto{\pgfqpoint{1.489386in}{1.747121in}}%
\pgfpathlineto{\pgfqpoint{1.497348in}{1.873389in}}%
\pgfpathlineto{\pgfqpoint{1.497657in}{1.869185in}}%
\pgfpathlineto{\pgfqpoint{1.505512in}{1.743973in}}%
\pgfpathlineto{\pgfqpoint{1.505871in}{1.748833in}}%
\pgfpathlineto{\pgfqpoint{1.513836in}{1.875083in}}%
\pgfpathlineto{\pgfqpoint{1.514149in}{1.870808in}}%
\pgfpathlineto{\pgfqpoint{1.521997in}{1.745624in}}%
\pgfpathlineto{\pgfqpoint{1.522356in}{1.750479in}}%
\pgfpathlineto{\pgfqpoint{1.530324in}{1.876602in}}%
\pgfpathlineto{\pgfqpoint{1.530644in}{1.872216in}}%
\pgfpathlineto{\pgfqpoint{1.538482in}{1.747013in}}%
\pgfpathlineto{\pgfqpoint{1.538841in}{1.751861in}}%
\pgfpathlineto{\pgfqpoint{1.546812in}{1.877777in}}%
\pgfpathlineto{\pgfqpoint{1.547135in}{1.873324in}}%
\pgfpathlineto{\pgfqpoint{1.554967in}{1.747980in}}%
\pgfpathlineto{\pgfqpoint{1.555326in}{1.752819in}}%
\pgfpathlineto{\pgfqpoint{1.563297in}{1.878480in}}%
\pgfpathlineto{\pgfqpoint{1.563619in}{1.874017in}}%
\pgfpathlineto{\pgfqpoint{1.571452in}{1.748425in}}%
\pgfpathlineto{\pgfqpoint{1.571811in}{1.753253in}}%
\pgfpathlineto{\pgfqpoint{1.579782in}{1.878669in}}%
\pgfpathlineto{\pgfqpoint{1.580104in}{1.874196in}}%
\pgfpathlineto{\pgfqpoint{1.587936in}{1.748380in}}%
\pgfpathlineto{\pgfqpoint{1.588295in}{1.753198in}}%
\pgfpathlineto{\pgfqpoint{1.596267in}{1.878411in}}%
\pgfpathlineto{\pgfqpoint{1.596589in}{1.873930in}}%
\pgfpathlineto{\pgfqpoint{1.604421in}{1.747943in}}%
\pgfpathlineto{\pgfqpoint{1.604780in}{1.752753in}}%
\pgfpathlineto{\pgfqpoint{1.612752in}{1.877824in}}%
\pgfpathlineto{\pgfqpoint{1.613074in}{1.873339in}}%
\pgfpathlineto{\pgfqpoint{1.620906in}{1.747247in}}%
\pgfpathlineto{\pgfqpoint{1.621265in}{1.752052in}}%
\pgfpathlineto{\pgfqpoint{1.629237in}{1.877055in}}%
\pgfpathlineto{\pgfqpoint{1.629559in}{1.872567in}}%
\pgfpathlineto{\pgfqpoint{1.637391in}{1.746440in}}%
\pgfpathlineto{\pgfqpoint{1.637750in}{1.751244in}}%
\pgfpathlineto{\pgfqpoint{1.645720in}{1.876241in}}%
\pgfpathlineto{\pgfqpoint{1.646041in}{1.871778in}}%
\pgfpathlineto{\pgfqpoint{1.653875in}{1.745640in}}%
\pgfpathlineto{\pgfqpoint{1.654234in}{1.750445in}}%
\pgfpathlineto{\pgfqpoint{1.662203in}{1.875491in}}%
\pgfpathlineto{\pgfqpoint{1.662522in}{1.871071in}}%
\pgfpathlineto{\pgfqpoint{1.670360in}{1.744943in}}%
\pgfpathlineto{\pgfqpoint{1.670719in}{1.749751in}}%
\pgfpathlineto{\pgfqpoint{1.678686in}{1.874888in}}%
\pgfpathlineto{\pgfqpoint{1.679002in}{1.870520in}}%
\pgfpathlineto{\pgfqpoint{1.686845in}{1.744427in}}%
\pgfpathlineto{\pgfqpoint{1.687204in}{1.749239in}}%
\pgfpathlineto{\pgfqpoint{1.695170in}{1.874490in}}%
\pgfpathlineto{\pgfqpoint{1.695484in}{1.870158in}}%
\pgfpathlineto{\pgfqpoint{1.703330in}{1.744134in}}%
\pgfpathlineto{\pgfqpoint{1.703689in}{1.748951in}}%
\pgfpathlineto{\pgfqpoint{1.711654in}{1.874320in}}%
\pgfpathlineto{\pgfqpoint{1.711967in}{1.870010in}}%
\pgfpathlineto{\pgfqpoint{1.719815in}{1.744072in}}%
\pgfpathlineto{\pgfqpoint{1.720174in}{1.748894in}}%
\pgfpathlineto{\pgfqpoint{1.728139in}{1.874368in}}%
\pgfpathlineto{\pgfqpoint{1.728452in}{1.870065in}}%
\pgfpathlineto{\pgfqpoint{1.736299in}{1.744215in}}%
\pgfpathlineto{\pgfqpoint{1.736659in}{1.749041in}}%
\pgfpathlineto{\pgfqpoint{1.744624in}{1.874596in}}%
\pgfpathlineto{\pgfqpoint{1.744937in}{1.870286in}}%
\pgfpathlineto{\pgfqpoint{1.752784in}{1.744513in}}%
\pgfpathlineto{\pgfqpoint{1.753143in}{1.749342in}}%
\pgfpathlineto{\pgfqpoint{1.761109in}{1.874945in}}%
\pgfpathlineto{\pgfqpoint{1.761424in}{1.870618in}}%
\pgfpathlineto{\pgfqpoint{1.769269in}{1.744902in}}%
\pgfpathlineto{\pgfqpoint{1.769628in}{1.749732in}}%
\pgfpathlineto{\pgfqpoint{1.777595in}{1.875349in}}%
\pgfpathlineto{\pgfqpoint{1.777911in}{1.870997in}}%
\pgfpathlineto{\pgfqpoint{1.785754in}{1.745316in}}%
\pgfpathlineto{\pgfqpoint{1.786113in}{1.750146in}}%
\pgfpathlineto{\pgfqpoint{1.794081in}{1.875747in}}%
\pgfpathlineto{\pgfqpoint{1.794399in}{1.871363in}}%
\pgfpathlineto{\pgfqpoint{1.802239in}{1.745696in}}%
\pgfpathlineto{\pgfqpoint{1.802598in}{1.750525in}}%
\pgfpathlineto{\pgfqpoint{1.810567in}{1.876086in}}%
\pgfpathlineto{\pgfqpoint{1.810886in}{1.871678in}}%
\pgfpathlineto{\pgfqpoint{1.818724in}{1.745994in}}%
\pgfpathlineto{\pgfqpoint{1.819083in}{1.750821in}}%
\pgfpathlineto{\pgfqpoint{1.827052in}{1.876326in}}%
\pgfpathlineto{\pgfqpoint{1.827373in}{1.871901in}}%
\pgfpathlineto{\pgfqpoint{1.835209in}{1.746182in}}%
\pgfpathlineto{\pgfqpoint{1.835568in}{1.751006in}}%
\pgfpathlineto{\pgfqpoint{1.843538in}{1.876450in}}%
\pgfpathlineto{\pgfqpoint{1.843858in}{1.872014in}}%
\pgfpathlineto{\pgfqpoint{1.851694in}{1.746250in}}%
\pgfpathlineto{\pgfqpoint{1.852053in}{1.751071in}}%
\pgfpathlineto{\pgfqpoint{1.860023in}{1.876459in}}%
\pgfpathlineto{\pgfqpoint{1.860343in}{1.872017in}}%
\pgfpathlineto{\pgfqpoint{1.868178in}{1.746207in}}%
\pgfpathlineto{\pgfqpoint{1.868538in}{1.751026in}}%
\pgfpathlineto{\pgfqpoint{1.876508in}{1.876368in}}%
\pgfpathlineto{\pgfqpoint{1.876828in}{1.871927in}}%
\pgfpathlineto{\pgfqpoint{1.884663in}{1.746076in}}%
\pgfpathlineto{\pgfqpoint{1.885022in}{1.750894in}}%
\pgfpathlineto{\pgfqpoint{1.892992in}{1.876205in}}%
\pgfpathlineto{\pgfqpoint{1.893312in}{1.871770in}}%
\pgfpathlineto{\pgfqpoint{1.901148in}{1.745888in}}%
\pgfpathlineto{\pgfqpoint{1.901507in}{1.750705in}}%
\pgfpathlineto{\pgfqpoint{1.909476in}{1.876003in}}%
\pgfpathlineto{\pgfqpoint{1.909796in}{1.871578in}}%
\pgfpathlineto{\pgfqpoint{1.917633in}{1.745677in}}%
\pgfpathlineto{\pgfqpoint{1.917992in}{1.750493in}}%
\pgfpathlineto{\pgfqpoint{1.925961in}{1.875796in}}%
\pgfpathlineto{\pgfqpoint{1.926280in}{1.871381in}}%
\pgfpathlineto{\pgfqpoint{1.934118in}{1.745475in}}%
\pgfpathlineto{\pgfqpoint{1.934477in}{1.750292in}}%
\pgfpathlineto{\pgfqpoint{1.942445in}{1.875612in}}%
\pgfpathlineto{\pgfqpoint{1.942764in}{1.871208in}}%
\pgfpathlineto{\pgfqpoint{1.950603in}{1.745309in}}%
\pgfpathlineto{\pgfqpoint{1.950962in}{1.750127in}}%
\pgfpathlineto{\pgfqpoint{1.958929in}{1.875473in}}%
\pgfpathlineto{\pgfqpoint{1.959247in}{1.871083in}}%
\pgfpathlineto{\pgfqpoint{1.967087in}{1.745196in}}%
\pgfpathlineto{\pgfqpoint{1.967447in}{1.750015in}}%
\pgfpathlineto{\pgfqpoint{1.975414in}{1.875391in}}%
\pgfpathlineto{\pgfqpoint{1.975731in}{1.871011in}}%
\pgfpathlineto{\pgfqpoint{1.983572in}{1.745141in}}%
\pgfpathlineto{\pgfqpoint{1.983931in}{1.749962in}}%
\pgfpathlineto{\pgfqpoint{1.991899in}{1.875368in}}%
\pgfpathlineto{\pgfqpoint{1.992216in}{1.870993in}}%
\pgfpathlineto{\pgfqpoint{2.000057in}{1.745145in}}%
\pgfpathlineto{\pgfqpoint{2.000416in}{1.749967in}}%
\pgfpathlineto{\pgfqpoint{2.008384in}{1.875399in}}%
\pgfpathlineto{\pgfqpoint{2.008701in}{1.871024in}}%
\pgfpathlineto{\pgfqpoint{2.016542in}{1.745199in}}%
\pgfpathlineto{\pgfqpoint{2.016901in}{1.750022in}}%
\pgfpathlineto{\pgfqpoint{2.024869in}{1.875471in}}%
\pgfpathlineto{\pgfqpoint{2.025186in}{1.871093in}}%
\pgfpathlineto{\pgfqpoint{2.033027in}{1.745287in}}%
\pgfpathlineto{\pgfqpoint{2.033386in}{1.750110in}}%
\pgfpathlineto{\pgfqpoint{2.041354in}{1.875569in}}%
\pgfpathlineto{\pgfqpoint{2.041672in}{1.871184in}}%
\pgfpathlineto{\pgfqpoint{2.049512in}{1.745393in}}%
\pgfpathlineto{\pgfqpoint{2.049871in}{1.750216in}}%
\pgfpathlineto{\pgfqpoint{2.057839in}{1.875676in}}%
\pgfpathlineto{\pgfqpoint{2.058157in}{1.871282in}}%
\pgfpathlineto{\pgfqpoint{2.065997in}{1.745499in}}%
\pgfpathlineto{\pgfqpoint{2.066356in}{1.750323in}}%
\pgfpathlineto{\pgfqpoint{2.074324in}{1.875776in}}%
\pgfpathlineto{\pgfqpoint{2.074643in}{1.871375in}}%
\pgfpathlineto{\pgfqpoint{2.082481in}{1.745592in}}%
\pgfpathlineto{\pgfqpoint{2.082841in}{1.750415in}}%
\pgfpathlineto{\pgfqpoint{2.090809in}{1.875856in}}%
\pgfpathlineto{\pgfqpoint{2.091128in}{1.871450in}}%
\pgfpathlineto{\pgfqpoint{2.098966in}{1.745660in}}%
\pgfpathlineto{\pgfqpoint{2.099325in}{1.750482in}}%
\pgfpathlineto{\pgfqpoint{2.107294in}{1.875908in}}%
\pgfpathlineto{\pgfqpoint{2.107613in}{1.871498in}}%
\pgfpathlineto{\pgfqpoint{2.115451in}{1.745697in}}%
\pgfpathlineto{\pgfqpoint{2.115810in}{1.750519in}}%
\pgfpathlineto{\pgfqpoint{2.123779in}{1.875929in}}%
\pgfpathlineto{\pgfqpoint{2.124098in}{1.871516in}}%
\pgfpathlineto{\pgfqpoint{2.131936in}{1.745704in}}%
\pgfpathlineto{\pgfqpoint{2.132295in}{1.750525in}}%
\pgfpathlineto{\pgfqpoint{2.140264in}{1.875921in}}%
\pgfpathlineto{\pgfqpoint{2.140583in}{1.871508in}}%
\pgfpathlineto{\pgfqpoint{2.148421in}{1.745683in}}%
\pgfpathlineto{\pgfqpoint{2.148780in}{1.750504in}}%
\pgfpathlineto{\pgfqpoint{2.156749in}{1.875889in}}%
\pgfpathlineto{\pgfqpoint{2.157068in}{1.871477in}}%
\pgfpathlineto{\pgfqpoint{2.164906in}{1.745642in}}%
\pgfpathlineto{\pgfqpoint{2.165265in}{1.750463in}}%
\pgfpathlineto{\pgfqpoint{2.173233in}{1.875842in}}%
\pgfpathlineto{\pgfqpoint{2.173552in}{1.871431in}}%
\pgfpathlineto{\pgfqpoint{2.181391in}{1.745590in}}%
\pgfpathlineto{\pgfqpoint{2.181750in}{1.750410in}}%
\pgfpathlineto{\pgfqpoint{2.189718in}{1.875788in}}%
\pgfpathlineto{\pgfqpoint{2.190037in}{1.871380in}}%
\pgfpathlineto{\pgfqpoint{2.197875in}{1.745535in}}%
\pgfpathlineto{\pgfqpoint{2.198234in}{1.750355in}}%
\pgfpathlineto{\pgfqpoint{2.206203in}{1.875735in}}%
\pgfpathlineto{\pgfqpoint{2.206522in}{1.871330in}}%
\pgfpathlineto{\pgfqpoint{2.214360in}{1.745486in}}%
\pgfpathlineto{\pgfqpoint{2.214719in}{1.750306in}}%
\pgfpathlineto{\pgfqpoint{2.222687in}{1.875691in}}%
\pgfpathlineto{\pgfqpoint{2.223006in}{1.871289in}}%
\pgfpathlineto{\pgfqpoint{2.230845in}{1.745447in}}%
\pgfpathlineto{\pgfqpoint{2.231204in}{1.750268in}}%
\pgfpathlineto{\pgfqpoint{2.239172in}{1.875660in}}%
\pgfpathlineto{\pgfqpoint{2.239491in}{1.871260in}}%
\pgfpathlineto{\pgfqpoint{2.247330in}{1.745423in}}%
\pgfpathlineto{\pgfqpoint{2.247689in}{1.750244in}}%
\pgfpathlineto{\pgfqpoint{2.255657in}{1.875645in}}%
\pgfpathlineto{\pgfqpoint{2.255976in}{1.871246in}}%
\pgfpathlineto{\pgfqpoint{2.263815in}{1.745415in}}%
\pgfpathlineto{\pgfqpoint{2.264174in}{1.750237in}}%
\pgfpathlineto{\pgfqpoint{2.272142in}{1.875645in}}%
\pgfpathlineto{\pgfqpoint{2.272460in}{1.871247in}}%
\pgfpathlineto{\pgfqpoint{2.280300in}{1.745422in}}%
\pgfpathlineto{\pgfqpoint{2.280659in}{1.750243in}}%
\pgfpathlineto{\pgfqpoint{2.288627in}{1.875658in}}%
\pgfpathlineto{\pgfqpoint{2.288945in}{1.871259in}}%
\pgfpathlineto{\pgfqpoint{2.296784in}{1.745440in}}%
\pgfpathlineto{\pgfqpoint{2.297144in}{1.750262in}}%
\pgfpathlineto{\pgfqpoint{2.305112in}{1.875680in}}%
\pgfpathlineto{\pgfqpoint{2.305430in}{1.871281in}}%
\pgfpathlineto{\pgfqpoint{2.313269in}{1.745465in}}%
\pgfpathlineto{\pgfqpoint{2.313628in}{1.750287in}}%
\pgfpathlineto{\pgfqpoint{2.321597in}{1.875706in}}%
\pgfpathlineto{\pgfqpoint{2.321915in}{1.871306in}}%
\pgfpathlineto{\pgfqpoint{2.329754in}{1.745493in}}%
\pgfpathlineto{\pgfqpoint{2.330113in}{1.750315in}}%
\pgfpathlineto{\pgfqpoint{2.338081in}{1.875734in}}%
\pgfpathlineto{\pgfqpoint{2.338400in}{1.871332in}}%
\pgfpathlineto{\pgfqpoint{2.346239in}{1.745519in}}%
\pgfpathlineto{\pgfqpoint{2.346598in}{1.750341in}}%
\pgfpathlineto{\pgfqpoint{2.354566in}{1.875757in}}%
\pgfpathlineto{\pgfqpoint{2.354885in}{1.871354in}}%
\pgfpathlineto{\pgfqpoint{2.362724in}{1.745541in}}%
\pgfpathlineto{\pgfqpoint{2.363083in}{1.750362in}}%
\pgfpathlineto{\pgfqpoint{2.371051in}{1.875775in}}%
\pgfpathlineto{\pgfqpoint{2.371370in}{1.871371in}}%
\pgfpathlineto{\pgfqpoint{2.379209in}{1.745555in}}%
\pgfpathlineto{\pgfqpoint{2.379568in}{1.750377in}}%
\pgfpathlineto{\pgfqpoint{2.387536in}{1.875785in}}%
\pgfpathlineto{\pgfqpoint{2.387855in}{1.871380in}}%
\pgfpathlineto{\pgfqpoint{2.395694in}{1.745561in}}%
\pgfpathlineto{\pgfqpoint{2.396053in}{1.750383in}}%
\pgfpathlineto{\pgfqpoint{2.404021in}{1.875788in}}%
\pgfpathlineto{\pgfqpoint{2.404340in}{1.871382in}}%
\pgfpathlineto{\pgfqpoint{2.412178in}{1.745560in}}%
\pgfpathlineto{\pgfqpoint{2.412537in}{1.750381in}}%
\pgfpathlineto{\pgfqpoint{2.420506in}{1.875783in}}%
\pgfpathlineto{\pgfqpoint{2.420825in}{1.871377in}}%
\pgfpathlineto{\pgfqpoint{2.428663in}{1.745553in}}%
\pgfpathlineto{\pgfqpoint{2.429022in}{1.750374in}}%
\pgfpathlineto{\pgfqpoint{2.436991in}{1.875773in}}%
\pgfpathlineto{\pgfqpoint{2.437310in}{1.871368in}}%
\pgfpathlineto{\pgfqpoint{2.445148in}{1.745541in}}%
\pgfpathlineto{\pgfqpoint{2.445507in}{1.750362in}}%
\pgfpathlineto{\pgfqpoint{2.453475in}{1.875760in}}%
\pgfpathlineto{\pgfqpoint{2.453794in}{1.871355in}}%
\pgfpathlineto{\pgfqpoint{2.461633in}{1.745527in}}%
\pgfpathlineto{\pgfqpoint{2.461992in}{1.750348in}}%
\pgfpathlineto{\pgfqpoint{2.469960in}{1.875746in}}%
\pgfpathlineto{\pgfqpoint{2.470279in}{1.871342in}}%
\pgfpathlineto{\pgfqpoint{2.478118in}{1.745513in}}%
\pgfpathlineto{\pgfqpoint{2.478477in}{1.750334in}}%
\pgfpathlineto{\pgfqpoint{2.486445in}{1.875734in}}%
\pgfpathlineto{\pgfqpoint{2.486764in}{1.871330in}}%
\pgfpathlineto{\pgfqpoint{2.494597in}{1.745702in}}%
\pgfpathlineto{\pgfqpoint{2.494963in}{1.750606in}}%
\pgfpathlineto{\pgfqpoint{2.502924in}{1.875952in}}%
\pgfpathlineto{\pgfqpoint{2.503233in}{1.871708in}}%
\pgfpathlineto{\pgfqpoint{2.511088in}{1.745615in}}%
\pgfpathlineto{\pgfqpoint{2.511447in}{1.750436in}}%
\pgfpathlineto{\pgfqpoint{2.519415in}{1.875825in}}%
\pgfpathlineto{\pgfqpoint{2.519734in}{1.871416in}}%
\pgfpathlineto{\pgfqpoint{2.527572in}{1.745582in}}%
\pgfpathlineto{\pgfqpoint{2.527931in}{1.750403in}}%
\pgfpathlineto{\pgfqpoint{2.535900in}{1.875789in}}%
\pgfpathlineto{\pgfqpoint{2.536219in}{1.871382in}}%
\pgfpathlineto{\pgfqpoint{2.544057in}{1.745544in}}%
\pgfpathlineto{\pgfqpoint{2.544416in}{1.750365in}}%
\pgfpathlineto{\pgfqpoint{2.552385in}{1.875752in}}%
\pgfpathlineto{\pgfqpoint{2.552704in}{1.871346in}}%
\pgfpathlineto{\pgfqpoint{2.560542in}{1.745508in}}%
\pgfpathlineto{\pgfqpoint{2.560901in}{1.750328in}}%
\pgfpathlineto{\pgfqpoint{2.568869in}{1.875718in}}%
\pgfpathlineto{\pgfqpoint{2.569188in}{1.871315in}}%
\pgfpathlineto{\pgfqpoint{2.577027in}{1.745478in}}%
\pgfpathlineto{\pgfqpoint{2.577386in}{1.750298in}}%
\pgfpathlineto{\pgfqpoint{2.585354in}{1.875693in}}%
\pgfpathlineto{\pgfqpoint{2.585673in}{1.871291in}}%
\pgfpathlineto{\pgfqpoint{2.593512in}{1.745457in}}%
\pgfpathlineto{\pgfqpoint{2.593871in}{1.750278in}}%
\pgfpathlineto{\pgfqpoint{2.601839in}{1.875678in}}%
\pgfpathlineto{\pgfqpoint{2.602158in}{1.871277in}}%
\pgfpathlineto{\pgfqpoint{2.609997in}{1.745447in}}%
\pgfpathlineto{\pgfqpoint{2.610356in}{1.750268in}}%
\pgfpathlineto{\pgfqpoint{2.618324in}{1.875673in}}%
\pgfpathlineto{\pgfqpoint{2.618642in}{1.871273in}}%
\pgfpathlineto{\pgfqpoint{2.626481in}{1.745447in}}%
\pgfpathlineto{\pgfqpoint{2.626840in}{1.750269in}}%
\pgfpathlineto{\pgfqpoint{2.634809in}{1.875678in}}%
\pgfpathlineto{\pgfqpoint{2.635127in}{1.871279in}}%
\pgfpathlineto{\pgfqpoint{2.642966in}{1.745457in}}%
\pgfpathlineto{\pgfqpoint{2.643325in}{1.750278in}}%
\pgfpathlineto{\pgfqpoint{2.651293in}{1.875691in}}%
\pgfpathlineto{\pgfqpoint{2.651612in}{1.871291in}}%
\pgfpathlineto{\pgfqpoint{2.659451in}{1.745472in}}%
\pgfpathlineto{\pgfqpoint{2.659810in}{1.750294in}}%
\pgfpathlineto{\pgfqpoint{2.667778in}{1.875709in}}%
\pgfpathlineto{\pgfqpoint{2.668097in}{1.871308in}}%
\pgfpathlineto{\pgfqpoint{2.675936in}{1.745491in}}%
\pgfpathlineto{\pgfqpoint{2.676295in}{1.750313in}}%
\pgfpathlineto{\pgfqpoint{2.684263in}{1.875728in}}%
\pgfpathlineto{\pgfqpoint{2.684582in}{1.871326in}}%
\pgfpathlineto{\pgfqpoint{2.692421in}{1.745510in}}%
\pgfpathlineto{\pgfqpoint{2.692780in}{1.750331in}}%
\pgfpathlineto{\pgfqpoint{2.700748in}{1.875745in}}%
\pgfpathlineto{\pgfqpoint{2.701067in}{1.871342in}}%
\pgfpathlineto{\pgfqpoint{2.708906in}{1.745526in}}%
\pgfpathlineto{\pgfqpoint{2.709265in}{1.750348in}}%
\pgfpathlineto{\pgfqpoint{2.717233in}{1.875759in}}%
\pgfpathlineto{\pgfqpoint{2.717552in}{1.871356in}}%
\pgfpathlineto{\pgfqpoint{2.725391in}{1.745538in}}%
\pgfpathlineto{\pgfqpoint{2.725750in}{1.750360in}}%
\pgfpathlineto{\pgfqpoint{2.733718in}{1.875769in}}%
\pgfpathlineto{\pgfqpoint{2.734037in}{1.871364in}}%
\pgfpathlineto{\pgfqpoint{2.741875in}{1.745545in}}%
\pgfpathlineto{\pgfqpoint{2.742234in}{1.750366in}}%
\pgfpathlineto{\pgfqpoint{2.750203in}{1.875773in}}%
\pgfpathlineto{\pgfqpoint{2.750522in}{1.871368in}}%
\pgfpathlineto{\pgfqpoint{2.758354in}{1.745747in}}%
\pgfpathlineto{\pgfqpoint{2.758721in}{1.750651in}}%
\pgfpathlineto{\pgfqpoint{2.766681in}{1.875999in}}%
\pgfpathlineto{\pgfqpoint{2.766991in}{1.871755in}}%
\pgfpathlineto{\pgfqpoint{2.774845in}{1.745664in}}%
\pgfpathlineto{\pgfqpoint{2.775204in}{1.750484in}}%
\pgfpathlineto{\pgfqpoint{2.783173in}{1.875872in}}%
\pgfpathlineto{\pgfqpoint{2.783492in}{1.871461in}}%
\pgfpathlineto{\pgfqpoint{2.791330in}{1.745628in}}%
\pgfpathlineto{\pgfqpoint{2.791689in}{1.750448in}}%
\pgfpathlineto{\pgfqpoint{2.799658in}{1.875831in}}%
\pgfpathlineto{\pgfqpoint{2.799977in}{1.871421in}}%
\pgfpathlineto{\pgfqpoint{2.807815in}{1.745582in}}%
\pgfpathlineto{\pgfqpoint{2.808174in}{1.750402in}}%
\pgfpathlineto{\pgfqpoint{2.816142in}{1.875783in}}%
\pgfpathlineto{\pgfqpoint{2.816461in}{1.871375in}}%
\pgfpathlineto{\pgfqpoint{2.824300in}{1.745533in}}%
\pgfpathlineto{\pgfqpoint{2.824659in}{1.750354in}}%
\pgfpathlineto{\pgfqpoint{2.832627in}{1.875736in}}%
\pgfpathlineto{\pgfqpoint{2.832946in}{1.871331in}}%
\pgfpathlineto{\pgfqpoint{2.840784in}{1.745489in}}%
\pgfpathlineto{\pgfqpoint{2.841144in}{1.750310in}}%
\pgfpathlineto{\pgfqpoint{2.849112in}{1.875697in}}%
\pgfpathlineto{\pgfqpoint{2.849431in}{1.871295in}}%
\pgfpathlineto{\pgfqpoint{2.857269in}{1.745455in}}%
\pgfpathlineto{\pgfqpoint{2.857628in}{1.750276in}}%
\pgfpathlineto{\pgfqpoint{2.865596in}{1.875670in}}%
\pgfpathlineto{\pgfqpoint{2.865915in}{1.871269in}}%
\pgfpathlineto{\pgfqpoint{2.873754in}{1.745434in}}%
\pgfpathlineto{\pgfqpoint{2.874113in}{1.750255in}}%
\pgfpathlineto{\pgfqpoint{2.882081in}{1.875656in}}%
\pgfpathlineto{\pgfqpoint{2.882400in}{1.871257in}}%
\pgfpathlineto{\pgfqpoint{2.890239in}{1.745427in}}%
\pgfpathlineto{\pgfqpoint{2.890598in}{1.750248in}}%
\pgfpathlineto{\pgfqpoint{2.898566in}{1.875656in}}%
\pgfpathlineto{\pgfqpoint{2.898885in}{1.871257in}}%
\pgfpathlineto{\pgfqpoint{2.906724in}{1.745432in}}%
\pgfpathlineto{\pgfqpoint{2.907083in}{1.750254in}}%
\pgfpathlineto{\pgfqpoint{2.915051in}{1.875667in}}%
\pgfpathlineto{\pgfqpoint{2.915370in}{1.871268in}}%
\pgfpathlineto{\pgfqpoint{2.923209in}{1.745448in}}%
\pgfpathlineto{\pgfqpoint{2.923568in}{1.750270in}}%
\pgfpathlineto{\pgfqpoint{2.931536in}{1.875686in}}%
\pgfpathlineto{\pgfqpoint{2.931855in}{1.871287in}}%
\pgfpathlineto{\pgfqpoint{2.939694in}{1.745470in}}%
\pgfpathlineto{\pgfqpoint{2.940053in}{1.750292in}}%
\pgfpathlineto{\pgfqpoint{2.948021in}{1.875710in}}%
\pgfpathlineto{\pgfqpoint{2.948340in}{1.871309in}}%
\pgfpathlineto{\pgfqpoint{2.956178in}{1.745495in}}%
\pgfpathlineto{\pgfqpoint{2.956537in}{1.750317in}}%
\pgfpathlineto{\pgfqpoint{2.964506in}{1.875734in}}%
\pgfpathlineto{\pgfqpoint{2.964825in}{1.871332in}}%
\pgfpathlineto{\pgfqpoint{2.972663in}{1.745518in}}%
\pgfpathlineto{\pgfqpoint{2.973022in}{1.750340in}}%
\pgfpathlineto{\pgfqpoint{2.980991in}{1.875755in}}%
\pgfpathlineto{\pgfqpoint{2.981309in}{1.871352in}}%
\pgfpathlineto{\pgfqpoint{2.989148in}{1.745537in}}%
\pgfpathlineto{\pgfqpoint{2.989507in}{1.750359in}}%
\pgfpathlineto{\pgfqpoint{2.997475in}{1.875771in}}%
\pgfpathlineto{\pgfqpoint{2.997794in}{1.871367in}}%
\pgfpathlineto{\pgfqpoint{3.005633in}{1.745550in}}%
\pgfpathlineto{\pgfqpoint{3.005992in}{1.750371in}}%
\pgfpathlineto{\pgfqpoint{3.013960in}{1.875780in}}%
\pgfpathlineto{\pgfqpoint{3.014279in}{1.871375in}}%
\pgfpathlineto{\pgfqpoint{3.022118in}{1.745556in}}%
\pgfpathlineto{\pgfqpoint{3.022477in}{1.750377in}}%
\pgfpathlineto{\pgfqpoint{3.030445in}{1.875782in}}%
\pgfpathlineto{\pgfqpoint{3.030764in}{1.871377in}}%
\pgfpathlineto{\pgfqpoint{3.038603in}{1.745555in}}%
\pgfpathlineto{\pgfqpoint{3.038962in}{1.750376in}}%
\pgfpathlineto{\pgfqpoint{3.046930in}{1.875778in}}%
\pgfpathlineto{\pgfqpoint{3.047249in}{1.871373in}}%
\pgfpathlineto{\pgfqpoint{3.055087in}{1.745548in}}%
\pgfpathlineto{\pgfqpoint{3.055447in}{1.750369in}}%
\pgfpathlineto{\pgfqpoint{3.063415in}{1.875769in}}%
\pgfpathlineto{\pgfqpoint{3.063734in}{1.871364in}}%
\pgfpathlineto{\pgfqpoint{3.071572in}{1.745538in}}%
\pgfpathlineto{\pgfqpoint{3.071931in}{1.750359in}}%
\pgfpathlineto{\pgfqpoint{3.079900in}{1.875758in}}%
\pgfpathlineto{\pgfqpoint{3.080219in}{1.871353in}}%
\pgfpathlineto{\pgfqpoint{3.088057in}{1.745525in}}%
\pgfpathlineto{\pgfqpoint{3.088416in}{1.750346in}}%
\pgfpathlineto{\pgfqpoint{3.096385in}{1.875746in}}%
\pgfpathlineto{\pgfqpoint{3.096703in}{1.871342in}}%
\pgfpathlineto{\pgfqpoint{3.104542in}{1.745513in}}%
\pgfpathlineto{\pgfqpoint{3.104901in}{1.750334in}}%
\pgfpathlineto{\pgfqpoint{3.112869in}{1.875734in}}%
\pgfpathlineto{\pgfqpoint{3.113188in}{1.871331in}}%
\pgfpathlineto{\pgfqpoint{3.121027in}{1.745503in}}%
\pgfpathlineto{\pgfqpoint{3.121386in}{1.750324in}}%
\pgfpathlineto{\pgfqpoint{3.129354in}{1.875725in}}%
\pgfpathlineto{\pgfqpoint{3.129673in}{1.871323in}}%
\pgfpathlineto{\pgfqpoint{3.137512in}{1.745495in}}%
\pgfpathlineto{\pgfqpoint{3.137871in}{1.750317in}}%
\pgfpathlineto{\pgfqpoint{3.145839in}{1.875720in}}%
\pgfpathlineto{\pgfqpoint{3.146158in}{1.871317in}}%
\pgfpathlineto{\pgfqpoint{3.153997in}{1.745491in}}%
\pgfpathlineto{\pgfqpoint{3.154356in}{1.750313in}}%
\pgfpathlineto{\pgfqpoint{3.162324in}{1.875718in}}%
\pgfpathlineto{\pgfqpoint{3.162643in}{1.871316in}}%
\pgfpathlineto{\pgfqpoint{3.170132in}{1.750854in}}%
\pgfpathlineto{\pgfqpoint{3.170132in}{1.750854in}}%
\pgfusepath{stroke}%
\end{pgfscope}%
\begin{pgfscope}%
\pgfsetrectcap%
\pgfsetmiterjoin%
\pgfsetlinewidth{0.803000pt}%
\definecolor{currentstroke}{rgb}{0.000000,0.000000,0.000000}%
\pgfsetstrokecolor{currentstroke}%
\pgfsetdash{}{0pt}%
\pgfpathmoveto{\pgfqpoint{0.573768in}{1.532029in}}%
\pgfpathlineto{\pgfqpoint{0.573768in}{2.438279in}}%
\pgfusepath{stroke}%
\end{pgfscope}%
\begin{pgfscope}%
\pgfsetrectcap%
\pgfsetmiterjoin%
\pgfsetlinewidth{0.803000pt}%
\definecolor{currentstroke}{rgb}{0.000000,0.000000,0.000000}%
\pgfsetstrokecolor{currentstroke}%
\pgfsetdash{}{0pt}%
\pgfpathmoveto{\pgfqpoint{3.293768in}{1.532029in}}%
\pgfpathlineto{\pgfqpoint{3.293768in}{2.438279in}}%
\pgfusepath{stroke}%
\end{pgfscope}%
\begin{pgfscope}%
\pgfsetrectcap%
\pgfsetmiterjoin%
\pgfsetlinewidth{0.803000pt}%
\definecolor{currentstroke}{rgb}{0.000000,0.000000,0.000000}%
\pgfsetstrokecolor{currentstroke}%
\pgfsetdash{}{0pt}%
\pgfpathmoveto{\pgfqpoint{0.573768in}{1.532029in}}%
\pgfpathlineto{\pgfqpoint{3.293768in}{1.532029in}}%
\pgfusepath{stroke}%
\end{pgfscope}%
\begin{pgfscope}%
\pgfsetrectcap%
\pgfsetmiterjoin%
\pgfsetlinewidth{0.803000pt}%
\definecolor{currentstroke}{rgb}{0.000000,0.000000,0.000000}%
\pgfsetstrokecolor{currentstroke}%
\pgfsetdash{}{0pt}%
\pgfpathmoveto{\pgfqpoint{0.573768in}{2.438279in}}%
\pgfpathlineto{\pgfqpoint{3.293768in}{2.438279in}}%
\pgfusepath{stroke}%
\end{pgfscope}%
\begin{pgfscope}%
\definecolor{textcolor}{rgb}{0.000000,0.000000,0.000000}%
\pgfsetstrokecolor{textcolor}%
\pgfsetfillcolor{textcolor}%
\pgftext[x=0.628168in,y=2.374842in,left,top]{\color{textcolor}{\rmfamily\fontsize{8.000000}{9.600000}\selectfont\catcode`\^=\active\def^{\ifmmode\sp\else\^{}\fi}\catcode`\%=\active\def%{\%}(b)}}%
\end{pgfscope}%
\begin{pgfscope}%
\pgfsetbuttcap%
\pgfsetmiterjoin%
\definecolor{currentfill}{rgb}{1.000000,1.000000,1.000000}%
\pgfsetfillcolor{currentfill}%
\pgfsetlinewidth{0.000000pt}%
\definecolor{currentstroke}{rgb}{0.000000,0.000000,0.000000}%
\pgfsetstrokecolor{currentstroke}%
\pgfsetstrokeopacity{0.000000}%
\pgfsetdash{}{0pt}%
\pgfpathmoveto{\pgfqpoint{0.573768in}{0.462654in}}%
\pgfpathlineto{\pgfqpoint{3.293768in}{0.462654in}}%
\pgfpathlineto{\pgfqpoint{3.293768in}{1.368904in}}%
\pgfpathlineto{\pgfqpoint{0.573768in}{1.368904in}}%
\pgfpathlineto{\pgfqpoint{0.573768in}{0.462654in}}%
\pgfpathclose%
\pgfusepath{fill}%
\end{pgfscope}%
\begin{pgfscope}%
\pgfpathrectangle{\pgfqpoint{0.573768in}{0.462654in}}{\pgfqpoint{2.720000in}{0.906250in}}%
\pgfusepath{clip}%
\pgfsetbuttcap%
\pgfsetroundjoin%
\pgfsetlinewidth{0.501875pt}%
\definecolor{currentstroke}{rgb}{0.690196,0.690196,0.690196}%
\pgfsetstrokecolor{currentstroke}%
\pgfsetstrokeopacity{0.250000}%
\pgfsetdash{{1.850000pt}{0.800000pt}}{0.000000pt}%
\pgfpathmoveto{\pgfqpoint{0.697405in}{0.462654in}}%
\pgfpathlineto{\pgfqpoint{0.697405in}{1.368904in}}%
\pgfusepath{stroke}%
\end{pgfscope}%
\begin{pgfscope}%
\pgfsetbuttcap%
\pgfsetroundjoin%
\definecolor{currentfill}{rgb}{0.000000,0.000000,0.000000}%
\pgfsetfillcolor{currentfill}%
\pgfsetlinewidth{0.803000pt}%
\definecolor{currentstroke}{rgb}{0.000000,0.000000,0.000000}%
\pgfsetstrokecolor{currentstroke}%
\pgfsetdash{}{0pt}%
\pgfsys@defobject{currentmarker}{\pgfqpoint{0.000000in}{-0.048611in}}{\pgfqpoint{0.000000in}{0.000000in}}{%
\pgfpathmoveto{\pgfqpoint{0.000000in}{0.000000in}}%
\pgfpathlineto{\pgfqpoint{0.000000in}{-0.048611in}}%
\pgfusepath{stroke,fill}%
}%
\begin{pgfscope}%
\pgfsys@transformshift{0.697405in}{0.462654in}%
\pgfsys@useobject{currentmarker}{}%
\end{pgfscope}%
\end{pgfscope}%
\begin{pgfscope}%
\definecolor{textcolor}{rgb}{0.000000,0.000000,0.000000}%
\pgfsetstrokecolor{textcolor}%
\pgfsetfillcolor{textcolor}%
\pgftext[x=0.697405in,y=0.365432in,,top]{\color{textcolor}{\rmfamily\fontsize{8.000000}{9.600000}\selectfont\catcode`\^=\active\def^{\ifmmode\sp\else\^{}\fi}\catcode`\%=\active\def%{\%}$\mathdefault{0}$}}%
\end{pgfscope}%
\begin{pgfscope}%
\pgfpathrectangle{\pgfqpoint{0.573768in}{0.462654in}}{\pgfqpoint{2.720000in}{0.906250in}}%
\pgfusepath{clip}%
\pgfsetbuttcap%
\pgfsetroundjoin%
\pgfsetlinewidth{0.501875pt}%
\definecolor{currentstroke}{rgb}{0.690196,0.690196,0.690196}%
\pgfsetstrokecolor{currentstroke}%
\pgfsetstrokeopacity{0.250000}%
\pgfsetdash{{1.850000pt}{0.800000pt}}{0.000000pt}%
\pgfpathmoveto{\pgfqpoint{1.109526in}{0.462654in}}%
\pgfpathlineto{\pgfqpoint{1.109526in}{1.368904in}}%
\pgfusepath{stroke}%
\end{pgfscope}%
\begin{pgfscope}%
\pgfsetbuttcap%
\pgfsetroundjoin%
\definecolor{currentfill}{rgb}{0.000000,0.000000,0.000000}%
\pgfsetfillcolor{currentfill}%
\pgfsetlinewidth{0.803000pt}%
\definecolor{currentstroke}{rgb}{0.000000,0.000000,0.000000}%
\pgfsetstrokecolor{currentstroke}%
\pgfsetdash{}{0pt}%
\pgfsys@defobject{currentmarker}{\pgfqpoint{0.000000in}{-0.048611in}}{\pgfqpoint{0.000000in}{0.000000in}}{%
\pgfpathmoveto{\pgfqpoint{0.000000in}{0.000000in}}%
\pgfpathlineto{\pgfqpoint{0.000000in}{-0.048611in}}%
\pgfusepath{stroke,fill}%
}%
\begin{pgfscope}%
\pgfsys@transformshift{1.109526in}{0.462654in}%
\pgfsys@useobject{currentmarker}{}%
\end{pgfscope}%
\end{pgfscope}%
\begin{pgfscope}%
\definecolor{textcolor}{rgb}{0.000000,0.000000,0.000000}%
\pgfsetstrokecolor{textcolor}%
\pgfsetfillcolor{textcolor}%
\pgftext[x=1.109526in,y=0.365432in,,top]{\color{textcolor}{\rmfamily\fontsize{8.000000}{9.600000}\selectfont\catcode`\^=\active\def^{\ifmmode\sp\else\^{}\fi}\catcode`\%=\active\def%{\%}$\mathdefault{1}$}}%
\end{pgfscope}%
\begin{pgfscope}%
\pgfpathrectangle{\pgfqpoint{0.573768in}{0.462654in}}{\pgfqpoint{2.720000in}{0.906250in}}%
\pgfusepath{clip}%
\pgfsetbuttcap%
\pgfsetroundjoin%
\pgfsetlinewidth{0.501875pt}%
\definecolor{currentstroke}{rgb}{0.690196,0.690196,0.690196}%
\pgfsetstrokecolor{currentstroke}%
\pgfsetstrokeopacity{0.250000}%
\pgfsetdash{{1.850000pt}{0.800000pt}}{0.000000pt}%
\pgfpathmoveto{\pgfqpoint{1.521647in}{0.462654in}}%
\pgfpathlineto{\pgfqpoint{1.521647in}{1.368904in}}%
\pgfusepath{stroke}%
\end{pgfscope}%
\begin{pgfscope}%
\pgfsetbuttcap%
\pgfsetroundjoin%
\definecolor{currentfill}{rgb}{0.000000,0.000000,0.000000}%
\pgfsetfillcolor{currentfill}%
\pgfsetlinewidth{0.803000pt}%
\definecolor{currentstroke}{rgb}{0.000000,0.000000,0.000000}%
\pgfsetstrokecolor{currentstroke}%
\pgfsetdash{}{0pt}%
\pgfsys@defobject{currentmarker}{\pgfqpoint{0.000000in}{-0.048611in}}{\pgfqpoint{0.000000in}{0.000000in}}{%
\pgfpathmoveto{\pgfqpoint{0.000000in}{0.000000in}}%
\pgfpathlineto{\pgfqpoint{0.000000in}{-0.048611in}}%
\pgfusepath{stroke,fill}%
}%
\begin{pgfscope}%
\pgfsys@transformshift{1.521647in}{0.462654in}%
\pgfsys@useobject{currentmarker}{}%
\end{pgfscope}%
\end{pgfscope}%
\begin{pgfscope}%
\definecolor{textcolor}{rgb}{0.000000,0.000000,0.000000}%
\pgfsetstrokecolor{textcolor}%
\pgfsetfillcolor{textcolor}%
\pgftext[x=1.521647in,y=0.365432in,,top]{\color{textcolor}{\rmfamily\fontsize{8.000000}{9.600000}\selectfont\catcode`\^=\active\def^{\ifmmode\sp\else\^{}\fi}\catcode`\%=\active\def%{\%}$\mathdefault{2}$}}%
\end{pgfscope}%
\begin{pgfscope}%
\pgfpathrectangle{\pgfqpoint{0.573768in}{0.462654in}}{\pgfqpoint{2.720000in}{0.906250in}}%
\pgfusepath{clip}%
\pgfsetbuttcap%
\pgfsetroundjoin%
\pgfsetlinewidth{0.501875pt}%
\definecolor{currentstroke}{rgb}{0.690196,0.690196,0.690196}%
\pgfsetstrokecolor{currentstroke}%
\pgfsetstrokeopacity{0.250000}%
\pgfsetdash{{1.850000pt}{0.800000pt}}{0.000000pt}%
\pgfpathmoveto{\pgfqpoint{1.933768in}{0.462654in}}%
\pgfpathlineto{\pgfqpoint{1.933768in}{1.368904in}}%
\pgfusepath{stroke}%
\end{pgfscope}%
\begin{pgfscope}%
\pgfsetbuttcap%
\pgfsetroundjoin%
\definecolor{currentfill}{rgb}{0.000000,0.000000,0.000000}%
\pgfsetfillcolor{currentfill}%
\pgfsetlinewidth{0.803000pt}%
\definecolor{currentstroke}{rgb}{0.000000,0.000000,0.000000}%
\pgfsetstrokecolor{currentstroke}%
\pgfsetdash{}{0pt}%
\pgfsys@defobject{currentmarker}{\pgfqpoint{0.000000in}{-0.048611in}}{\pgfqpoint{0.000000in}{0.000000in}}{%
\pgfpathmoveto{\pgfqpoint{0.000000in}{0.000000in}}%
\pgfpathlineto{\pgfqpoint{0.000000in}{-0.048611in}}%
\pgfusepath{stroke,fill}%
}%
\begin{pgfscope}%
\pgfsys@transformshift{1.933768in}{0.462654in}%
\pgfsys@useobject{currentmarker}{}%
\end{pgfscope}%
\end{pgfscope}%
\begin{pgfscope}%
\definecolor{textcolor}{rgb}{0.000000,0.000000,0.000000}%
\pgfsetstrokecolor{textcolor}%
\pgfsetfillcolor{textcolor}%
\pgftext[x=1.933768in,y=0.365432in,,top]{\color{textcolor}{\rmfamily\fontsize{8.000000}{9.600000}\selectfont\catcode`\^=\active\def^{\ifmmode\sp\else\^{}\fi}\catcode`\%=\active\def%{\%}$\mathdefault{3}$}}%
\end{pgfscope}%
\begin{pgfscope}%
\pgfpathrectangle{\pgfqpoint{0.573768in}{0.462654in}}{\pgfqpoint{2.720000in}{0.906250in}}%
\pgfusepath{clip}%
\pgfsetbuttcap%
\pgfsetroundjoin%
\pgfsetlinewidth{0.501875pt}%
\definecolor{currentstroke}{rgb}{0.690196,0.690196,0.690196}%
\pgfsetstrokecolor{currentstroke}%
\pgfsetstrokeopacity{0.250000}%
\pgfsetdash{{1.850000pt}{0.800000pt}}{0.000000pt}%
\pgfpathmoveto{\pgfqpoint{2.345890in}{0.462654in}}%
\pgfpathlineto{\pgfqpoint{2.345890in}{1.368904in}}%
\pgfusepath{stroke}%
\end{pgfscope}%
\begin{pgfscope}%
\pgfsetbuttcap%
\pgfsetroundjoin%
\definecolor{currentfill}{rgb}{0.000000,0.000000,0.000000}%
\pgfsetfillcolor{currentfill}%
\pgfsetlinewidth{0.803000pt}%
\definecolor{currentstroke}{rgb}{0.000000,0.000000,0.000000}%
\pgfsetstrokecolor{currentstroke}%
\pgfsetdash{}{0pt}%
\pgfsys@defobject{currentmarker}{\pgfqpoint{0.000000in}{-0.048611in}}{\pgfqpoint{0.000000in}{0.000000in}}{%
\pgfpathmoveto{\pgfqpoint{0.000000in}{0.000000in}}%
\pgfpathlineto{\pgfqpoint{0.000000in}{-0.048611in}}%
\pgfusepath{stroke,fill}%
}%
\begin{pgfscope}%
\pgfsys@transformshift{2.345890in}{0.462654in}%
\pgfsys@useobject{currentmarker}{}%
\end{pgfscope}%
\end{pgfscope}%
\begin{pgfscope}%
\definecolor{textcolor}{rgb}{0.000000,0.000000,0.000000}%
\pgfsetstrokecolor{textcolor}%
\pgfsetfillcolor{textcolor}%
\pgftext[x=2.345890in,y=0.365432in,,top]{\color{textcolor}{\rmfamily\fontsize{8.000000}{9.600000}\selectfont\catcode`\^=\active\def^{\ifmmode\sp\else\^{}\fi}\catcode`\%=\active\def%{\%}$\mathdefault{4}$}}%
\end{pgfscope}%
\begin{pgfscope}%
\pgfpathrectangle{\pgfqpoint{0.573768in}{0.462654in}}{\pgfqpoint{2.720000in}{0.906250in}}%
\pgfusepath{clip}%
\pgfsetbuttcap%
\pgfsetroundjoin%
\pgfsetlinewidth{0.501875pt}%
\definecolor{currentstroke}{rgb}{0.690196,0.690196,0.690196}%
\pgfsetstrokecolor{currentstroke}%
\pgfsetstrokeopacity{0.250000}%
\pgfsetdash{{1.850000pt}{0.800000pt}}{0.000000pt}%
\pgfpathmoveto{\pgfqpoint{2.758011in}{0.462654in}}%
\pgfpathlineto{\pgfqpoint{2.758011in}{1.368904in}}%
\pgfusepath{stroke}%
\end{pgfscope}%
\begin{pgfscope}%
\pgfsetbuttcap%
\pgfsetroundjoin%
\definecolor{currentfill}{rgb}{0.000000,0.000000,0.000000}%
\pgfsetfillcolor{currentfill}%
\pgfsetlinewidth{0.803000pt}%
\definecolor{currentstroke}{rgb}{0.000000,0.000000,0.000000}%
\pgfsetstrokecolor{currentstroke}%
\pgfsetdash{}{0pt}%
\pgfsys@defobject{currentmarker}{\pgfqpoint{0.000000in}{-0.048611in}}{\pgfqpoint{0.000000in}{0.000000in}}{%
\pgfpathmoveto{\pgfqpoint{0.000000in}{0.000000in}}%
\pgfpathlineto{\pgfqpoint{0.000000in}{-0.048611in}}%
\pgfusepath{stroke,fill}%
}%
\begin{pgfscope}%
\pgfsys@transformshift{2.758011in}{0.462654in}%
\pgfsys@useobject{currentmarker}{}%
\end{pgfscope}%
\end{pgfscope}%
\begin{pgfscope}%
\definecolor{textcolor}{rgb}{0.000000,0.000000,0.000000}%
\pgfsetstrokecolor{textcolor}%
\pgfsetfillcolor{textcolor}%
\pgftext[x=2.758011in,y=0.365432in,,top]{\color{textcolor}{\rmfamily\fontsize{8.000000}{9.600000}\selectfont\catcode`\^=\active\def^{\ifmmode\sp\else\^{}\fi}\catcode`\%=\active\def%{\%}$\mathdefault{5}$}}%
\end{pgfscope}%
\begin{pgfscope}%
\pgfpathrectangle{\pgfqpoint{0.573768in}{0.462654in}}{\pgfqpoint{2.720000in}{0.906250in}}%
\pgfusepath{clip}%
\pgfsetbuttcap%
\pgfsetroundjoin%
\pgfsetlinewidth{0.501875pt}%
\definecolor{currentstroke}{rgb}{0.690196,0.690196,0.690196}%
\pgfsetstrokecolor{currentstroke}%
\pgfsetstrokeopacity{0.250000}%
\pgfsetdash{{1.850000pt}{0.800000pt}}{0.000000pt}%
\pgfpathmoveto{\pgfqpoint{3.170132in}{0.462654in}}%
\pgfpathlineto{\pgfqpoint{3.170132in}{1.368904in}}%
\pgfusepath{stroke}%
\end{pgfscope}%
\begin{pgfscope}%
\pgfsetbuttcap%
\pgfsetroundjoin%
\definecolor{currentfill}{rgb}{0.000000,0.000000,0.000000}%
\pgfsetfillcolor{currentfill}%
\pgfsetlinewidth{0.803000pt}%
\definecolor{currentstroke}{rgb}{0.000000,0.000000,0.000000}%
\pgfsetstrokecolor{currentstroke}%
\pgfsetdash{}{0pt}%
\pgfsys@defobject{currentmarker}{\pgfqpoint{0.000000in}{-0.048611in}}{\pgfqpoint{0.000000in}{0.000000in}}{%
\pgfpathmoveto{\pgfqpoint{0.000000in}{0.000000in}}%
\pgfpathlineto{\pgfqpoint{0.000000in}{-0.048611in}}%
\pgfusepath{stroke,fill}%
}%
\begin{pgfscope}%
\pgfsys@transformshift{3.170132in}{0.462654in}%
\pgfsys@useobject{currentmarker}{}%
\end{pgfscope}%
\end{pgfscope}%
\begin{pgfscope}%
\definecolor{textcolor}{rgb}{0.000000,0.000000,0.000000}%
\pgfsetstrokecolor{textcolor}%
\pgfsetfillcolor{textcolor}%
\pgftext[x=3.170132in,y=0.365432in,,top]{\color{textcolor}{\rmfamily\fontsize{8.000000}{9.600000}\selectfont\catcode`\^=\active\def^{\ifmmode\sp\else\^{}\fi}\catcode`\%=\active\def%{\%}$\mathdefault{6}$}}%
\end{pgfscope}%
\begin{pgfscope}%
\definecolor{textcolor}{rgb}{0.000000,0.000000,0.000000}%
\pgfsetstrokecolor{textcolor}%
\pgfsetfillcolor{textcolor}%
\pgftext[x=1.933768in,y=0.211111in,,top]{\color{textcolor}{\rmfamily\fontsize{8.000000}{9.600000}\selectfont\catcode`\^=\active\def^{\ifmmode\sp\else\^{}\fi}\catcode`\%=\active\def%{\%}Time (ms)}}%
\end{pgfscope}%
\begin{pgfscope}%
\pgfpathrectangle{\pgfqpoint{0.573768in}{0.462654in}}{\pgfqpoint{2.720000in}{0.906250in}}%
\pgfusepath{clip}%
\pgfsetbuttcap%
\pgfsetroundjoin%
\pgfsetlinewidth{0.501875pt}%
\definecolor{currentstroke}{rgb}{0.690196,0.690196,0.690196}%
\pgfsetstrokecolor{currentstroke}%
\pgfsetstrokeopacity{0.250000}%
\pgfsetdash{{1.850000pt}{0.800000pt}}{0.000000pt}%
\pgfpathmoveto{\pgfqpoint{0.573768in}{0.503801in}}%
\pgfpathlineto{\pgfqpoint{3.293768in}{0.503801in}}%
\pgfusepath{stroke}%
\end{pgfscope}%
\begin{pgfscope}%
\pgfsetbuttcap%
\pgfsetroundjoin%
\definecolor{currentfill}{rgb}{0.000000,0.000000,0.000000}%
\pgfsetfillcolor{currentfill}%
\pgfsetlinewidth{0.803000pt}%
\definecolor{currentstroke}{rgb}{0.000000,0.000000,0.000000}%
\pgfsetstrokecolor{currentstroke}%
\pgfsetdash{}{0pt}%
\pgfsys@defobject{currentmarker}{\pgfqpoint{-0.048611in}{0.000000in}}{\pgfqpoint{-0.000000in}{0.000000in}}{%
\pgfpathmoveto{\pgfqpoint{-0.000000in}{0.000000in}}%
\pgfpathlineto{\pgfqpoint{-0.048611in}{0.000000in}}%
\pgfusepath{stroke,fill}%
}%
\begin{pgfscope}%
\pgfsys@transformshift{0.573768in}{0.503801in}%
\pgfsys@useobject{currentmarker}{}%
\end{pgfscope}%
\end{pgfscope}%
\begin{pgfscope}%
\definecolor{textcolor}{rgb}{0.000000,0.000000,0.000000}%
\pgfsetstrokecolor{textcolor}%
\pgfsetfillcolor{textcolor}%
\pgftext[x=0.266667in, y=0.465220in, left, base]{\color{textcolor}{\rmfamily\fontsize{8.000000}{9.600000}\selectfont\catcode`\^=\active\def^{\ifmmode\sp\else\^{}\fi}\catcode`\%=\active\def%{\%}$\mathdefault{\ensuremath{-}20}$}}%
\end{pgfscope}%
\begin{pgfscope}%
\pgfpathrectangle{\pgfqpoint{0.573768in}{0.462654in}}{\pgfqpoint{2.720000in}{0.906250in}}%
\pgfusepath{clip}%
\pgfsetbuttcap%
\pgfsetroundjoin%
\pgfsetlinewidth{0.501875pt}%
\definecolor{currentstroke}{rgb}{0.690196,0.690196,0.690196}%
\pgfsetstrokecolor{currentstroke}%
\pgfsetstrokeopacity{0.250000}%
\pgfsetdash{{1.850000pt}{0.800000pt}}{0.000000pt}%
\pgfpathmoveto{\pgfqpoint{0.573768in}{0.882945in}}%
\pgfpathlineto{\pgfqpoint{3.293768in}{0.882945in}}%
\pgfusepath{stroke}%
\end{pgfscope}%
\begin{pgfscope}%
\pgfsetbuttcap%
\pgfsetroundjoin%
\definecolor{currentfill}{rgb}{0.000000,0.000000,0.000000}%
\pgfsetfillcolor{currentfill}%
\pgfsetlinewidth{0.803000pt}%
\definecolor{currentstroke}{rgb}{0.000000,0.000000,0.000000}%
\pgfsetstrokecolor{currentstroke}%
\pgfsetdash{}{0pt}%
\pgfsys@defobject{currentmarker}{\pgfqpoint{-0.048611in}{0.000000in}}{\pgfqpoint{-0.000000in}{0.000000in}}{%
\pgfpathmoveto{\pgfqpoint{-0.000000in}{0.000000in}}%
\pgfpathlineto{\pgfqpoint{-0.048611in}{0.000000in}}%
\pgfusepath{stroke,fill}%
}%
\begin{pgfscope}%
\pgfsys@transformshift{0.573768in}{0.882945in}%
\pgfsys@useobject{currentmarker}{}%
\end{pgfscope}%
\end{pgfscope}%
\begin{pgfscope}%
\definecolor{textcolor}{rgb}{0.000000,0.000000,0.000000}%
\pgfsetstrokecolor{textcolor}%
\pgfsetfillcolor{textcolor}%
\pgftext[x=0.266667in, y=0.844365in, left, base]{\color{textcolor}{\rmfamily\fontsize{8.000000}{9.600000}\selectfont\catcode`\^=\active\def^{\ifmmode\sp\else\^{}\fi}\catcode`\%=\active\def%{\%}$\mathdefault{\ensuremath{-}10}$}}%
\end{pgfscope}%
\begin{pgfscope}%
\pgfpathrectangle{\pgfqpoint{0.573768in}{0.462654in}}{\pgfqpoint{2.720000in}{0.906250in}}%
\pgfusepath{clip}%
\pgfsetbuttcap%
\pgfsetroundjoin%
\pgfsetlinewidth{0.501875pt}%
\definecolor{currentstroke}{rgb}{0.690196,0.690196,0.690196}%
\pgfsetstrokecolor{currentstroke}%
\pgfsetstrokeopacity{0.250000}%
\pgfsetdash{{1.850000pt}{0.800000pt}}{0.000000pt}%
\pgfpathmoveto{\pgfqpoint{0.573768in}{1.262090in}}%
\pgfpathlineto{\pgfqpoint{3.293768in}{1.262090in}}%
\pgfusepath{stroke}%
\end{pgfscope}%
\begin{pgfscope}%
\pgfsetbuttcap%
\pgfsetroundjoin%
\definecolor{currentfill}{rgb}{0.000000,0.000000,0.000000}%
\pgfsetfillcolor{currentfill}%
\pgfsetlinewidth{0.803000pt}%
\definecolor{currentstroke}{rgb}{0.000000,0.000000,0.000000}%
\pgfsetstrokecolor{currentstroke}%
\pgfsetdash{}{0pt}%
\pgfsys@defobject{currentmarker}{\pgfqpoint{-0.048611in}{0.000000in}}{\pgfqpoint{-0.000000in}{0.000000in}}{%
\pgfpathmoveto{\pgfqpoint{-0.000000in}{0.000000in}}%
\pgfpathlineto{\pgfqpoint{-0.048611in}{0.000000in}}%
\pgfusepath{stroke,fill}%
}%
\begin{pgfscope}%
\pgfsys@transformshift{0.573768in}{1.262090in}%
\pgfsys@useobject{currentmarker}{}%
\end{pgfscope}%
\end{pgfscope}%
\begin{pgfscope}%
\definecolor{textcolor}{rgb}{0.000000,0.000000,0.000000}%
\pgfsetstrokecolor{textcolor}%
\pgfsetfillcolor{textcolor}%
\pgftext[x=0.417518in, y=1.223510in, left, base]{\color{textcolor}{\rmfamily\fontsize{8.000000}{9.600000}\selectfont\catcode`\^=\active\def^{\ifmmode\sp\else\^{}\fi}\catcode`\%=\active\def%{\%}$\mathdefault{0}$}}%
\end{pgfscope}%
\begin{pgfscope}%
\definecolor{textcolor}{rgb}{0.000000,0.000000,0.000000}%
\pgfsetstrokecolor{textcolor}%
\pgfsetfillcolor{textcolor}%
\pgftext[x=0.211111in,y=0.915779in,,bottom,rotate=90.000000]{\color{textcolor}{\rmfamily\fontsize{8.000000}{9.600000}\selectfont\catcode`\^=\active\def^{\ifmmode\sp\else\^{}\fi}\catcode`\%=\active\def%{\%}$v_D$ (V)}}%
\end{pgfscope}%
\begin{pgfscope}%
\pgfpathrectangle{\pgfqpoint{0.573768in}{0.462654in}}{\pgfqpoint{2.720000in}{0.906250in}}%
\pgfusepath{clip}%
\pgfsetrectcap%
\pgfsetroundjoin%
\pgfsetlinewidth{1.505625pt}%
\definecolor{currentstroke}{rgb}{0.000000,0.619608,0.450980}%
\pgfsetstrokecolor{currentstroke}%
\pgfsetdash{}{0pt}%
\pgfpathmoveto{\pgfqpoint{0.697405in}{1.262090in}}%
\pgfpathlineto{\pgfqpoint{0.697428in}{1.272694in}}%
\pgfpathlineto{\pgfqpoint{0.697481in}{1.321636in}}%
\pgfpathlineto{\pgfqpoint{0.697749in}{0.850273in}}%
\pgfpathlineto{\pgfqpoint{0.698145in}{0.503848in}}%
\pgfpathlineto{\pgfqpoint{0.698473in}{0.503855in}}%
\pgfpathlineto{\pgfqpoint{0.705897in}{0.504801in}}%
\pgfpathlineto{\pgfqpoint{0.705987in}{0.513810in}}%
\pgfpathlineto{\pgfqpoint{0.706139in}{1.326702in}}%
\pgfpathlineto{\pgfqpoint{0.706827in}{1.326700in}}%
\pgfpathlineto{\pgfqpoint{0.709300in}{1.326600in}}%
\pgfpathlineto{\pgfqpoint{0.711443in}{1.326627in}}%
\pgfpathlineto{\pgfqpoint{0.713586in}{1.326527in}}%
\pgfpathlineto{\pgfqpoint{0.714156in}{1.326552in}}%
\pgfpathlineto{\pgfqpoint{0.714206in}{1.326384in}}%
\pgfpathlineto{\pgfqpoint{0.714218in}{1.240088in}}%
\pgfpathlineto{\pgfqpoint{0.714923in}{0.504442in}}%
\pgfpathlineto{\pgfqpoint{0.717231in}{0.504404in}}%
\pgfpathlineto{\pgfqpoint{0.722304in}{0.504974in}}%
\pgfpathlineto{\pgfqpoint{0.722396in}{0.505919in}}%
\pgfpathlineto{\pgfqpoint{0.722570in}{0.881077in}}%
\pgfpathlineto{\pgfqpoint{0.722635in}{1.327265in}}%
\pgfpathlineto{\pgfqpoint{0.723311in}{1.327259in}}%
\pgfpathlineto{\pgfqpoint{0.729410in}{1.327145in}}%
\pgfpathlineto{\pgfqpoint{0.730691in}{1.327091in}}%
\pgfpathlineto{\pgfqpoint{0.731616in}{0.504793in}}%
\pgfpathlineto{\pgfqpoint{0.735572in}{0.504815in}}%
\pgfpathlineto{\pgfqpoint{0.738789in}{0.505385in}}%
\pgfpathlineto{\pgfqpoint{0.738881in}{0.506655in}}%
\pgfpathlineto{\pgfqpoint{0.739056in}{0.892282in}}%
\pgfpathlineto{\pgfqpoint{0.739120in}{1.327546in}}%
\pgfpathlineto{\pgfqpoint{0.739798in}{1.327537in}}%
\pgfpathlineto{\pgfqpoint{0.747176in}{1.327354in}}%
\pgfpathlineto{\pgfqpoint{0.748664in}{0.504958in}}%
\pgfpathlineto{\pgfqpoint{0.753445in}{0.505071in}}%
\pgfpathlineto{\pgfqpoint{0.755274in}{0.505626in}}%
\pgfpathlineto{\pgfqpoint{0.755366in}{0.507086in}}%
\pgfpathlineto{\pgfqpoint{0.755541in}{0.898257in}}%
\pgfpathlineto{\pgfqpoint{0.755605in}{1.327679in}}%
\pgfpathlineto{\pgfqpoint{0.756283in}{1.327667in}}%
\pgfpathlineto{\pgfqpoint{0.763661in}{1.327456in}}%
\pgfpathlineto{\pgfqpoint{0.763673in}{1.324316in}}%
\pgfpathlineto{\pgfqpoint{0.764482in}{0.505279in}}%
\pgfpathlineto{\pgfqpoint{0.766130in}{0.505033in}}%
\pgfpathlineto{\pgfqpoint{0.771712in}{0.505411in}}%
\pgfpathlineto{\pgfqpoint{0.771851in}{0.507205in}}%
\pgfpathlineto{\pgfqpoint{0.772026in}{0.899850in}}%
\pgfpathlineto{\pgfqpoint{0.772090in}{1.327711in}}%
\pgfpathlineto{\pgfqpoint{0.772770in}{1.327697in}}%
\pgfpathlineto{\pgfqpoint{0.780146in}{1.327447in}}%
\pgfpathlineto{\pgfqpoint{0.780158in}{1.324103in}}%
\pgfpathlineto{\pgfqpoint{0.780967in}{0.505259in}}%
\pgfpathlineto{\pgfqpoint{0.782615in}{0.505006in}}%
\pgfpathlineto{\pgfqpoint{0.788197in}{0.505339in}}%
\pgfpathlineto{\pgfqpoint{0.788336in}{0.507051in}}%
\pgfpathlineto{\pgfqpoint{0.788511in}{0.897804in}}%
\pgfpathlineto{\pgfqpoint{0.788584in}{1.327681in}}%
\pgfpathlineto{\pgfqpoint{0.789302in}{1.327649in}}%
\pgfpathlineto{\pgfqpoint{0.796445in}{1.327397in}}%
\pgfpathlineto{\pgfqpoint{0.796631in}{1.327347in}}%
\pgfpathlineto{\pgfqpoint{0.798114in}{0.504920in}}%
\pgfpathlineto{\pgfqpoint{0.802895in}{0.504947in}}%
\pgfpathlineto{\pgfqpoint{0.804728in}{0.505413in}}%
\pgfpathlineto{\pgfqpoint{0.804821in}{0.506698in}}%
\pgfpathlineto{\pgfqpoint{0.804996in}{0.892929in}}%
\pgfpathlineto{\pgfqpoint{0.805068in}{1.327571in}}%
\pgfpathlineto{\pgfqpoint{0.805794in}{1.327534in}}%
\pgfpathlineto{\pgfqpoint{0.812224in}{1.327238in}}%
\pgfpathlineto{\pgfqpoint{0.813115in}{1.327164in}}%
\pgfpathlineto{\pgfqpoint{0.814045in}{0.504832in}}%
\pgfpathlineto{\pgfqpoint{0.817672in}{0.504729in}}%
\pgfpathlineto{\pgfqpoint{0.821213in}{0.505154in}}%
\pgfpathlineto{\pgfqpoint{0.821305in}{0.506233in}}%
\pgfpathlineto{\pgfqpoint{0.821480in}{0.886101in}}%
\pgfpathlineto{\pgfqpoint{0.821552in}{1.327406in}}%
\pgfpathlineto{\pgfqpoint{0.822273in}{1.327361in}}%
\pgfpathlineto{\pgfqpoint{0.827713in}{1.327048in}}%
\pgfpathlineto{\pgfqpoint{0.829600in}{1.326834in}}%
\pgfpathlineto{\pgfqpoint{0.830485in}{0.504608in}}%
\pgfpathlineto{\pgfqpoint{0.834442in}{0.504536in}}%
\pgfpathlineto{\pgfqpoint{0.837698in}{0.504880in}}%
\pgfpathlineto{\pgfqpoint{0.837790in}{0.505743in}}%
\pgfpathlineto{\pgfqpoint{0.837964in}{0.878258in}}%
\pgfpathlineto{\pgfqpoint{0.838037in}{1.327178in}}%
\pgfpathlineto{\pgfqpoint{0.838759in}{1.327142in}}%
\pgfpathlineto{\pgfqpoint{0.842880in}{1.326841in}}%
\pgfpathlineto{\pgfqpoint{0.845853in}{1.326624in}}%
\pgfpathlineto{\pgfqpoint{0.846085in}{1.326442in}}%
\pgfpathlineto{\pgfqpoint{0.846908in}{0.504393in}}%
\pgfpathlineto{\pgfqpoint{0.850864in}{0.504352in}}%
\pgfpathlineto{\pgfqpoint{0.854183in}{0.504635in}}%
\pgfpathlineto{\pgfqpoint{0.854275in}{0.505304in}}%
\pgfpathlineto{\pgfqpoint{0.854449in}{0.870358in}}%
\pgfpathlineto{\pgfqpoint{0.854525in}{1.326939in}}%
\pgfpathlineto{\pgfqpoint{0.855253in}{1.326897in}}%
\pgfpathlineto{\pgfqpoint{0.858386in}{1.326597in}}%
\pgfpathlineto{\pgfqpoint{0.860529in}{1.326455in}}%
\pgfpathlineto{\pgfqpoint{0.862570in}{1.325947in}}%
\pgfpathlineto{\pgfqpoint{0.863349in}{0.504213in}}%
\pgfpathlineto{\pgfqpoint{0.866481in}{0.504179in}}%
\pgfpathlineto{\pgfqpoint{0.870700in}{0.504554in}}%
\pgfpathlineto{\pgfqpoint{0.870787in}{0.505456in}}%
\pgfpathlineto{\pgfqpoint{0.871030in}{1.326694in}}%
\pgfpathlineto{\pgfqpoint{0.872252in}{1.326566in}}%
\pgfpathlineto{\pgfqpoint{0.875714in}{1.326252in}}%
\pgfpathlineto{\pgfqpoint{0.879055in}{1.325267in}}%
\pgfpathlineto{\pgfqpoint{0.879782in}{0.504071in}}%
\pgfpathlineto{\pgfqpoint{0.887063in}{0.504202in}}%
\pgfpathlineto{\pgfqpoint{0.887245in}{0.504768in}}%
\pgfpathlineto{\pgfqpoint{0.887417in}{0.858915in}}%
\pgfpathlineto{\pgfqpoint{0.887497in}{1.326540in}}%
\pgfpathlineto{\pgfqpoint{0.888229in}{1.326491in}}%
\pgfpathlineto{\pgfqpoint{0.890372in}{1.326176in}}%
\pgfpathlineto{\pgfqpoint{0.891856in}{1.326101in}}%
\pgfpathlineto{\pgfqpoint{0.893340in}{1.325750in}}%
\pgfpathlineto{\pgfqpoint{0.894494in}{1.325701in}}%
\pgfpathlineto{\pgfqpoint{0.895540in}{1.324641in}}%
\pgfpathlineto{\pgfqpoint{0.896334in}{0.503991in}}%
\pgfpathlineto{\pgfqpoint{0.903668in}{0.504376in}}%
\pgfpathlineto{\pgfqpoint{0.903755in}{0.505048in}}%
\pgfpathlineto{\pgfqpoint{0.904006in}{1.326465in}}%
\pgfpathlineto{\pgfqpoint{0.905240in}{1.326313in}}%
\pgfpathlineto{\pgfqpoint{0.907713in}{1.326114in}}%
\pgfpathlineto{\pgfqpoint{0.909856in}{1.325742in}}%
\pgfpathlineto{\pgfqpoint{0.911339in}{1.325588in}}%
\pgfpathlineto{\pgfqpoint{0.912025in}{1.324647in}}%
\pgfpathlineto{\pgfqpoint{0.912818in}{0.503995in}}%
\pgfpathlineto{\pgfqpoint{0.920157in}{0.504426in}}%
\pgfpathlineto{\pgfqpoint{0.920245in}{0.505226in}}%
\pgfpathlineto{\pgfqpoint{0.920484in}{1.326517in}}%
\pgfpathlineto{\pgfqpoint{0.921704in}{1.326396in}}%
\pgfpathlineto{\pgfqpoint{0.925166in}{1.326122in}}%
\pgfpathlineto{\pgfqpoint{0.928205in}{1.325715in}}%
\pgfpathlineto{\pgfqpoint{0.928459in}{1.325677in}}%
\pgfpathlineto{\pgfqpoint{0.928509in}{1.325187in}}%
\pgfpathlineto{\pgfqpoint{0.929262in}{0.504062in}}%
\pgfpathlineto{\pgfqpoint{0.936561in}{0.504319in}}%
\pgfpathlineto{\pgfqpoint{0.936699in}{0.504909in}}%
\pgfpathlineto{\pgfqpoint{0.936872in}{0.862131in}}%
\pgfpathlineto{\pgfqpoint{0.936952in}{1.326667in}}%
\pgfpathlineto{\pgfqpoint{0.937686in}{1.326635in}}%
\pgfpathlineto{\pgfqpoint{0.939829in}{1.326414in}}%
\pgfpathlineto{\pgfqpoint{0.941642in}{1.326359in}}%
\pgfpathlineto{\pgfqpoint{0.943455in}{1.326106in}}%
\pgfpathlineto{\pgfqpoint{0.944919in}{1.326016in}}%
\pgfpathlineto{\pgfqpoint{0.944994in}{1.325678in}}%
\pgfpathlineto{\pgfqpoint{0.945742in}{0.504156in}}%
\pgfpathlineto{\pgfqpoint{0.953041in}{0.504405in}}%
\pgfpathlineto{\pgfqpoint{0.953165in}{0.504903in}}%
\pgfpathlineto{\pgfqpoint{0.953255in}{0.511334in}}%
\pgfpathlineto{\pgfqpoint{0.953412in}{1.326821in}}%
\pgfpathlineto{\pgfqpoint{0.954119in}{1.326790in}}%
\pgfpathlineto{\pgfqpoint{0.956592in}{1.326593in}}%
\pgfpathlineto{\pgfqpoint{0.958735in}{1.326527in}}%
\pgfpathlineto{\pgfqpoint{0.960878in}{1.326298in}}%
\pgfpathlineto{\pgfqpoint{0.961429in}{1.326300in}}%
\pgfpathlineto{\pgfqpoint{0.961479in}{1.326087in}}%
\pgfpathlineto{\pgfqpoint{0.962224in}{0.504284in}}%
\pgfpathlineto{\pgfqpoint{0.965026in}{0.504257in}}%
\pgfpathlineto{\pgfqpoint{0.969645in}{0.505026in}}%
\pgfpathlineto{\pgfqpoint{0.969735in}{0.510046in}}%
\pgfpathlineto{\pgfqpoint{0.969892in}{1.326951in}}%
\pgfpathlineto{\pgfqpoint{0.970597in}{1.326926in}}%
\pgfpathlineto{\pgfqpoint{0.974059in}{1.326716in}}%
\pgfpathlineto{\pgfqpoint{0.976861in}{1.326609in}}%
\pgfpathlineto{\pgfqpoint{0.977964in}{1.326343in}}%
\pgfpathlineto{\pgfqpoint{0.978840in}{0.504344in}}%
\pgfpathlineto{\pgfqpoint{0.985958in}{0.504505in}}%
\pgfpathlineto{\pgfqpoint{0.986144in}{0.505351in}}%
\pgfpathlineto{\pgfqpoint{0.986233in}{0.527825in}}%
\pgfpathlineto{\pgfqpoint{0.986392in}{1.327067in}}%
\pgfpathlineto{\pgfqpoint{0.987115in}{1.327045in}}%
\pgfpathlineto{\pgfqpoint{0.990247in}{1.326859in}}%
\pgfpathlineto{\pgfqpoint{0.992719in}{1.326792in}}%
\pgfpathlineto{\pgfqpoint{0.994449in}{1.326522in}}%
\pgfpathlineto{\pgfqpoint{0.995261in}{0.504454in}}%
\pgfpathlineto{\pgfqpoint{0.998723in}{0.504437in}}%
\pgfpathlineto{\pgfqpoint{1.002546in}{0.504834in}}%
\pgfpathlineto{\pgfqpoint{1.002639in}{0.505663in}}%
\pgfpathlineto{\pgfqpoint{1.002813in}{0.876838in}}%
\pgfpathlineto{\pgfqpoint{1.002892in}{1.327145in}}%
\pgfpathlineto{\pgfqpoint{1.003630in}{1.327122in}}%
\pgfpathlineto{\pgfqpoint{1.007092in}{1.326934in}}%
\pgfpathlineto{\pgfqpoint{1.009894in}{1.326846in}}%
\pgfpathlineto{\pgfqpoint{1.010934in}{1.326632in}}%
\pgfpathlineto{\pgfqpoint{1.011753in}{0.504515in}}%
\pgfpathlineto{\pgfqpoint{1.015050in}{0.504475in}}%
\pgfpathlineto{\pgfqpoint{1.019031in}{0.504887in}}%
\pgfpathlineto{\pgfqpoint{1.019124in}{0.505758in}}%
\pgfpathlineto{\pgfqpoint{1.019298in}{0.878475in}}%
\pgfpathlineto{\pgfqpoint{1.019377in}{1.327192in}}%
\pgfpathlineto{\pgfqpoint{1.020122in}{1.327168in}}%
\pgfpathlineto{\pgfqpoint{1.023913in}{1.326967in}}%
\pgfpathlineto{\pgfqpoint{1.026991in}{1.326861in}}%
\pgfpathlineto{\pgfqpoint{1.027418in}{1.326686in}}%
\pgfpathlineto{\pgfqpoint{1.028191in}{0.504566in}}%
\pgfpathlineto{\pgfqpoint{1.030994in}{0.504480in}}%
\pgfpathlineto{\pgfqpoint{1.035516in}{0.504907in}}%
\pgfpathlineto{\pgfqpoint{1.035608in}{0.505794in}}%
\pgfpathlineto{\pgfqpoint{1.035782in}{0.879080in}}%
\pgfpathlineto{\pgfqpoint{1.035862in}{1.327208in}}%
\pgfpathlineto{\pgfqpoint{1.036607in}{1.327183in}}%
\pgfpathlineto{\pgfqpoint{1.040398in}{1.326978in}}%
\pgfpathlineto{\pgfqpoint{1.043476in}{1.326866in}}%
\pgfpathlineto{\pgfqpoint{1.043903in}{1.326692in}}%
\pgfpathlineto{\pgfqpoint{1.044714in}{0.504554in}}%
\pgfpathlineto{\pgfqpoint{1.047682in}{0.504484in}}%
\pgfpathlineto{\pgfqpoint{1.052033in}{0.505074in}}%
\pgfpathlineto{\pgfqpoint{1.052120in}{0.506584in}}%
\pgfpathlineto{\pgfqpoint{1.052365in}{1.327191in}}%
\pgfpathlineto{\pgfqpoint{1.053590in}{1.327124in}}%
\pgfpathlineto{\pgfqpoint{1.059689in}{1.326844in}}%
\pgfpathlineto{\pgfqpoint{1.060388in}{1.326660in}}%
\pgfpathlineto{\pgfqpoint{1.061119in}{0.504566in}}%
\pgfpathlineto{\pgfqpoint{1.063592in}{0.504448in}}%
\pgfpathlineto{\pgfqpoint{1.068486in}{0.504865in}}%
\pgfpathlineto{\pgfqpoint{1.068578in}{0.505718in}}%
\pgfpathlineto{\pgfqpoint{1.068752in}{0.877796in}}%
\pgfpathlineto{\pgfqpoint{1.068832in}{1.327171in}}%
\pgfpathlineto{\pgfqpoint{1.069576in}{1.327143in}}%
\pgfpathlineto{\pgfqpoint{1.073038in}{1.326935in}}%
\pgfpathlineto{\pgfqpoint{1.075840in}{1.326825in}}%
\pgfpathlineto{\pgfqpoint{1.076873in}{1.326597in}}%
\pgfpathlineto{\pgfqpoint{1.077692in}{0.504489in}}%
\pgfpathlineto{\pgfqpoint{1.081154in}{0.504446in}}%
\pgfpathlineto{\pgfqpoint{1.084971in}{0.504817in}}%
\pgfpathlineto{\pgfqpoint{1.085063in}{0.505632in}}%
\pgfpathlineto{\pgfqpoint{1.085237in}{0.876330in}}%
\pgfpathlineto{\pgfqpoint{1.085316in}{1.327128in}}%
\pgfpathlineto{\pgfqpoint{1.086054in}{1.327099in}}%
\pgfpathlineto{\pgfqpoint{1.089516in}{1.326878in}}%
\pgfpathlineto{\pgfqpoint{1.092319in}{1.326759in}}%
\pgfpathlineto{\pgfqpoint{1.093358in}{1.326514in}}%
\pgfpathlineto{\pgfqpoint{1.094208in}{0.504432in}}%
\pgfpathlineto{\pgfqpoint{1.098165in}{0.504422in}}%
\pgfpathlineto{\pgfqpoint{1.101455in}{0.504765in}}%
\pgfpathlineto{\pgfqpoint{1.101548in}{0.505538in}}%
\pgfpathlineto{\pgfqpoint{1.101722in}{0.874679in}}%
\pgfpathlineto{\pgfqpoint{1.101801in}{1.327079in}}%
\pgfpathlineto{\pgfqpoint{1.102539in}{1.327049in}}%
\pgfpathlineto{\pgfqpoint{1.105671in}{1.326833in}}%
\pgfpathlineto{\pgfqpoint{1.108143in}{1.326732in}}%
\pgfpathlineto{\pgfqpoint{1.109843in}{1.326423in}}%
\pgfpathlineto{\pgfqpoint{1.110703in}{0.504379in}}%
\pgfpathlineto{\pgfqpoint{1.116143in}{0.504435in}}%
\pgfpathlineto{\pgfqpoint{1.117940in}{0.504716in}}%
\pgfpathlineto{\pgfqpoint{1.118033in}{0.505451in}}%
\pgfpathlineto{\pgfqpoint{1.118206in}{0.873102in}}%
\pgfpathlineto{\pgfqpoint{1.118286in}{1.327031in}}%
\pgfpathlineto{\pgfqpoint{1.119023in}{1.327000in}}%
\pgfpathlineto{\pgfqpoint{1.122155in}{1.326774in}}%
\pgfpathlineto{\pgfqpoint{1.124628in}{1.326668in}}%
\pgfpathlineto{\pgfqpoint{1.126328in}{1.326338in}}%
\pgfpathlineto{\pgfqpoint{1.127203in}{0.504334in}}%
\pgfpathlineto{\pgfqpoint{1.134338in}{0.504467in}}%
\pgfpathlineto{\pgfqpoint{1.134500in}{0.505148in}}%
\pgfpathlineto{\pgfqpoint{1.134590in}{0.513995in}}%
\pgfpathlineto{\pgfqpoint{1.134747in}{1.326998in}}%
\pgfpathlineto{\pgfqpoint{1.135458in}{1.326963in}}%
\pgfpathlineto{\pgfqpoint{1.138260in}{1.326747in}}%
\pgfpathlineto{\pgfqpoint{1.140403in}{1.326667in}}%
\pgfpathlineto{\pgfqpoint{1.142639in}{1.326467in}}%
\pgfpathlineto{\pgfqpoint{1.142788in}{1.326453in}}%
\pgfpathlineto{\pgfqpoint{1.143524in}{0.504363in}}%
\pgfpathlineto{\pgfqpoint{1.145996in}{0.504282in}}%
\pgfpathlineto{\pgfqpoint{1.150980in}{0.505057in}}%
\pgfpathlineto{\pgfqpoint{1.151070in}{0.510763in}}%
\pgfpathlineto{\pgfqpoint{1.151227in}{1.326962in}}%
\pgfpathlineto{\pgfqpoint{1.151933in}{1.326931in}}%
\pgfpathlineto{\pgfqpoint{1.155065in}{1.326700in}}%
\pgfpathlineto{\pgfqpoint{1.157538in}{1.326586in}}%
\pgfpathlineto{\pgfqpoint{1.159297in}{1.326230in}}%
\pgfpathlineto{\pgfqpoint{1.160088in}{0.504309in}}%
\pgfpathlineto{\pgfqpoint{1.165363in}{0.504362in}}%
\pgfpathlineto{\pgfqpoint{1.167484in}{0.505266in}}%
\pgfpathlineto{\pgfqpoint{1.167657in}{0.858717in}}%
\pgfpathlineto{\pgfqpoint{1.167735in}{1.326953in}}%
\pgfpathlineto{\pgfqpoint{1.168468in}{1.326922in}}%
\pgfpathlineto{\pgfqpoint{1.171270in}{1.326704in}}%
\pgfpathlineto{\pgfqpoint{1.173413in}{1.326625in}}%
\pgfpathlineto{\pgfqpoint{1.175596in}{1.326429in}}%
\pgfpathlineto{\pgfqpoint{1.175782in}{1.326223in}}%
\pgfpathlineto{\pgfqpoint{1.176586in}{0.504304in}}%
\pgfpathlineto{\pgfqpoint{1.182190in}{0.504374in}}%
\pgfpathlineto{\pgfqpoint{1.183880in}{0.504643in}}%
\pgfpathlineto{\pgfqpoint{1.183972in}{0.505317in}}%
\pgfpathlineto{\pgfqpoint{1.184145in}{0.869465in}}%
\pgfpathlineto{\pgfqpoint{1.184225in}{1.326956in}}%
\pgfpathlineto{\pgfqpoint{1.184961in}{1.326926in}}%
\pgfpathlineto{\pgfqpoint{1.187763in}{1.326714in}}%
\pgfpathlineto{\pgfqpoint{1.189906in}{1.326636in}}%
\pgfpathlineto{\pgfqpoint{1.192081in}{1.326446in}}%
\pgfpathlineto{\pgfqpoint{1.192267in}{1.326243in}}%
\pgfpathlineto{\pgfqpoint{1.193123in}{0.504299in}}%
\pgfpathlineto{\pgfqpoint{1.200313in}{0.504516in}}%
\pgfpathlineto{\pgfqpoint{1.200432in}{0.505049in}}%
\pgfpathlineto{\pgfqpoint{1.200522in}{0.510140in}}%
\pgfpathlineto{\pgfqpoint{1.200680in}{1.326967in}}%
\pgfpathlineto{\pgfqpoint{1.201385in}{1.326938in}}%
\pgfpathlineto{\pgfqpoint{1.204847in}{1.326703in}}%
\pgfpathlineto{\pgfqpoint{1.207649in}{1.326569in}}%
\pgfpathlineto{\pgfqpoint{1.208728in}{1.326462in}}%
\pgfpathlineto{\pgfqpoint{1.209462in}{0.504370in}}%
\pgfpathlineto{\pgfqpoint{1.211935in}{0.504292in}}%
\pgfpathlineto{\pgfqpoint{1.216920in}{0.505110in}}%
\pgfpathlineto{\pgfqpoint{1.217010in}{0.511614in}}%
\pgfpathlineto{\pgfqpoint{1.217168in}{1.326993in}}%
\pgfpathlineto{\pgfqpoint{1.217875in}{1.326964in}}%
\pgfpathlineto{\pgfqpoint{1.221007in}{1.326749in}}%
\pgfpathlineto{\pgfqpoint{1.223480in}{1.326652in}}%
\pgfpathlineto{\pgfqpoint{1.225237in}{1.326326in}}%
\pgfpathlineto{\pgfqpoint{1.226113in}{0.504332in}}%
\pgfpathlineto{\pgfqpoint{1.233263in}{0.504512in}}%
\pgfpathlineto{\pgfqpoint{1.233409in}{0.505190in}}%
\pgfpathlineto{\pgfqpoint{1.233499in}{0.514791in}}%
\pgfpathlineto{\pgfqpoint{1.233657in}{1.327024in}}%
\pgfpathlineto{\pgfqpoint{1.234368in}{1.326992in}}%
\pgfpathlineto{\pgfqpoint{1.237170in}{1.326794in}}%
\pgfpathlineto{\pgfqpoint{1.239643in}{1.326708in}}%
\pgfpathlineto{\pgfqpoint{1.241722in}{1.326370in}}%
\pgfpathlineto{\pgfqpoint{1.242591in}{0.504354in}}%
\pgfpathlineto{\pgfqpoint{1.249185in}{0.504464in}}%
\pgfpathlineto{\pgfqpoint{1.249815in}{0.504706in}}%
\pgfpathlineto{\pgfqpoint{1.249902in}{0.505323in}}%
\pgfpathlineto{\pgfqpoint{1.249992in}{0.529832in}}%
\pgfpathlineto{\pgfqpoint{1.250152in}{1.327039in}}%
\pgfpathlineto{\pgfqpoint{1.250876in}{1.327013in}}%
\pgfpathlineto{\pgfqpoint{1.254008in}{1.326802in}}%
\pgfpathlineto{\pgfqpoint{1.256481in}{1.326715in}}%
\pgfpathlineto{\pgfqpoint{1.258206in}{1.326409in}}%
\pgfpathlineto{\pgfqpoint{1.259069in}{0.504374in}}%
\pgfpathlineto{\pgfqpoint{1.264674in}{0.504446in}}%
\pgfpathlineto{\pgfqpoint{1.266304in}{0.504740in}}%
\pgfpathlineto{\pgfqpoint{1.266396in}{0.505495in}}%
\pgfpathlineto{\pgfqpoint{1.266570in}{0.873882in}}%
\pgfpathlineto{\pgfqpoint{1.266650in}{1.327056in}}%
\pgfpathlineto{\pgfqpoint{1.267387in}{1.327028in}}%
\pgfpathlineto{\pgfqpoint{1.270519in}{1.326824in}}%
\pgfpathlineto{\pgfqpoint{1.272992in}{1.326735in}}%
\pgfpathlineto{\pgfqpoint{1.274691in}{1.326438in}}%
\pgfpathlineto{\pgfqpoint{1.275516in}{0.504402in}}%
\pgfpathlineto{\pgfqpoint{1.279472in}{0.504403in}}%
\pgfpathlineto{\pgfqpoint{1.282789in}{0.504753in}}%
\pgfpathlineto{\pgfqpoint{1.282881in}{0.505517in}}%
\pgfpathlineto{\pgfqpoint{1.283055in}{0.874287in}}%
\pgfpathlineto{\pgfqpoint{1.283134in}{1.327068in}}%
\pgfpathlineto{\pgfqpoint{1.283872in}{1.327041in}}%
\pgfpathlineto{\pgfqpoint{1.287004in}{1.326837in}}%
\pgfpathlineto{\pgfqpoint{1.289477in}{1.326748in}}%
\pgfpathlineto{\pgfqpoint{1.291176in}{1.326454in}}%
\pgfpathlineto{\pgfqpoint{1.291894in}{0.504451in}}%
\pgfpathlineto{\pgfqpoint{1.294531in}{0.504366in}}%
\pgfpathlineto{\pgfqpoint{1.299274in}{0.504759in}}%
\pgfpathlineto{\pgfqpoint{1.299366in}{0.505528in}}%
\pgfpathlineto{\pgfqpoint{1.299540in}{0.874478in}}%
\pgfpathlineto{\pgfqpoint{1.299619in}{1.327074in}}%
\pgfpathlineto{\pgfqpoint{1.300357in}{1.327046in}}%
\pgfpathlineto{\pgfqpoint{1.303489in}{1.326842in}}%
\pgfpathlineto{\pgfqpoint{1.305962in}{1.326752in}}%
\pgfpathlineto{\pgfqpoint{1.307661in}{1.326459in}}%
\pgfpathlineto{\pgfqpoint{1.308377in}{0.504454in}}%
\pgfpathlineto{\pgfqpoint{1.310850in}{0.504360in}}%
\pgfpathlineto{\pgfqpoint{1.315759in}{0.504758in}}%
\pgfpathlineto{\pgfqpoint{1.315851in}{0.505527in}}%
\pgfpathlineto{\pgfqpoint{1.316025in}{0.874468in}}%
\pgfpathlineto{\pgfqpoint{1.316104in}{1.327074in}}%
\pgfpathlineto{\pgfqpoint{1.316842in}{1.327045in}}%
\pgfpathlineto{\pgfqpoint{1.319974in}{1.326840in}}%
\pgfpathlineto{\pgfqpoint{1.322446in}{1.326749in}}%
\pgfpathlineto{\pgfqpoint{1.324146in}{1.326454in}}%
\pgfpathlineto{\pgfqpoint{1.324863in}{0.504450in}}%
\pgfpathlineto{\pgfqpoint{1.327501in}{0.504364in}}%
\pgfpathlineto{\pgfqpoint{1.332243in}{0.504753in}}%
\pgfpathlineto{\pgfqpoint{1.332336in}{0.505517in}}%
\pgfpathlineto{\pgfqpoint{1.332509in}{0.874297in}}%
\pgfpathlineto{\pgfqpoint{1.332589in}{1.327068in}}%
\pgfpathlineto{\pgfqpoint{1.333326in}{1.327040in}}%
\pgfpathlineto{\pgfqpoint{1.336458in}{1.326832in}}%
\pgfpathlineto{\pgfqpoint{1.338931in}{1.326740in}}%
\pgfpathlineto{\pgfqpoint{1.340631in}{1.326440in}}%
\pgfpathlineto{\pgfqpoint{1.341453in}{0.504403in}}%
\pgfpathlineto{\pgfqpoint{1.345244in}{0.504397in}}%
\pgfpathlineto{\pgfqpoint{1.348728in}{0.504744in}}%
\pgfpathlineto{\pgfqpoint{1.348821in}{0.505502in}}%
\pgfpathlineto{\pgfqpoint{1.348994in}{0.874015in}}%
\pgfpathlineto{\pgfqpoint{1.349074in}{1.327060in}}%
\pgfpathlineto{\pgfqpoint{1.349811in}{1.327031in}}%
\pgfpathlineto{\pgfqpoint{1.352943in}{1.326821in}}%
\pgfpathlineto{\pgfqpoint{1.355416in}{1.326726in}}%
\pgfpathlineto{\pgfqpoint{1.357115in}{1.326423in}}%
\pgfpathlineto{\pgfqpoint{1.357976in}{0.504381in}}%
\pgfpathlineto{\pgfqpoint{1.363416in}{0.504445in}}%
\pgfpathlineto{\pgfqpoint{1.365213in}{0.504734in}}%
\pgfpathlineto{\pgfqpoint{1.365305in}{0.505483in}}%
\pgfpathlineto{\pgfqpoint{1.365479in}{0.873686in}}%
\pgfpathlineto{\pgfqpoint{1.365559in}{1.327050in}}%
\pgfpathlineto{\pgfqpoint{1.366296in}{1.327021in}}%
\pgfpathlineto{\pgfqpoint{1.369428in}{1.326808in}}%
\pgfpathlineto{\pgfqpoint{1.371901in}{1.326712in}}%
\pgfpathlineto{\pgfqpoint{1.373600in}{1.326404in}}%
\pgfpathlineto{\pgfqpoint{1.374464in}{0.504370in}}%
\pgfpathlineto{\pgfqpoint{1.380234in}{0.504449in}}%
\pgfpathlineto{\pgfqpoint{1.381698in}{0.504724in}}%
\pgfpathlineto{\pgfqpoint{1.381790in}{0.505466in}}%
\pgfpathlineto{\pgfqpoint{1.381964in}{0.873364in}}%
\pgfpathlineto{\pgfqpoint{1.382044in}{1.327040in}}%
\pgfpathlineto{\pgfqpoint{1.382781in}{1.327011in}}%
\pgfpathlineto{\pgfqpoint{1.385913in}{1.326795in}}%
\pgfpathlineto{\pgfqpoint{1.388385in}{1.326699in}}%
\pgfpathlineto{\pgfqpoint{1.390085in}{1.326386in}}%
\pgfpathlineto{\pgfqpoint{1.390952in}{0.504361in}}%
\pgfpathlineto{\pgfqpoint{1.397216in}{0.504450in}}%
\pgfpathlineto{\pgfqpoint{1.398180in}{0.504706in}}%
\pgfpathlineto{\pgfqpoint{1.398270in}{0.505367in}}%
\pgfpathlineto{\pgfqpoint{1.398359in}{0.537398in}}%
\pgfpathlineto{\pgfqpoint{1.398520in}{1.327032in}}%
\pgfpathlineto{\pgfqpoint{1.399250in}{1.327004in}}%
\pgfpathlineto{\pgfqpoint{1.402383in}{1.326785in}}%
\pgfpathlineto{\pgfqpoint{1.404855in}{1.326690in}}%
\pgfpathlineto{\pgfqpoint{1.406570in}{1.326372in}}%
\pgfpathlineto{\pgfqpoint{1.407439in}{0.504354in}}%
\pgfpathlineto{\pgfqpoint{1.414033in}{0.504456in}}%
\pgfpathlineto{\pgfqpoint{1.414663in}{0.504693in}}%
\pgfpathlineto{\pgfqpoint{1.414750in}{0.505292in}}%
\pgfpathlineto{\pgfqpoint{1.414840in}{0.527922in}}%
\pgfpathlineto{\pgfqpoint{1.414999in}{1.327026in}}%
\pgfpathlineto{\pgfqpoint{1.415721in}{1.326999in}}%
\pgfpathlineto{\pgfqpoint{1.418524in}{1.326797in}}%
\pgfpathlineto{\pgfqpoint{1.420996in}{1.326706in}}%
\pgfpathlineto{\pgfqpoint{1.423055in}{1.326362in}}%
\pgfpathlineto{\pgfqpoint{1.423932in}{0.504347in}}%
\pgfpathlineto{\pgfqpoint{1.430856in}{0.504464in}}%
\pgfpathlineto{\pgfqpoint{1.431234in}{0.505274in}}%
\pgfpathlineto{\pgfqpoint{1.431324in}{0.525724in}}%
\pgfpathlineto{\pgfqpoint{1.431482in}{1.327023in}}%
\pgfpathlineto{\pgfqpoint{1.432202in}{1.326996in}}%
\pgfpathlineto{\pgfqpoint{1.435005in}{1.326794in}}%
\pgfpathlineto{\pgfqpoint{1.437477in}{1.326703in}}%
\pgfpathlineto{\pgfqpoint{1.439540in}{1.326359in}}%
\pgfpathlineto{\pgfqpoint{1.440416in}{0.504346in}}%
\pgfpathlineto{\pgfqpoint{1.447340in}{0.504464in}}%
\pgfpathlineto{\pgfqpoint{1.447718in}{0.505263in}}%
\pgfpathlineto{\pgfqpoint{1.447808in}{0.523979in}}%
\pgfpathlineto{\pgfqpoint{1.447966in}{1.327040in}}%
\pgfpathlineto{\pgfqpoint{1.448684in}{1.326996in}}%
\pgfpathlineto{\pgfqpoint{1.451487in}{1.326794in}}%
\pgfpathlineto{\pgfqpoint{1.453960in}{1.326704in}}%
\pgfpathlineto{\pgfqpoint{1.456025in}{1.326360in}}%
\pgfpathlineto{\pgfqpoint{1.456901in}{0.504347in}}%
\pgfpathlineto{\pgfqpoint{1.463825in}{0.504465in}}%
\pgfpathlineto{\pgfqpoint{1.464203in}{0.505273in}}%
\pgfpathlineto{\pgfqpoint{1.464293in}{0.525058in}}%
\pgfpathlineto{\pgfqpoint{1.464452in}{1.327043in}}%
\pgfpathlineto{\pgfqpoint{1.465171in}{1.326998in}}%
\pgfpathlineto{\pgfqpoint{1.467973in}{1.326797in}}%
\pgfpathlineto{\pgfqpoint{1.470446in}{1.326707in}}%
\pgfpathlineto{\pgfqpoint{1.472509in}{1.326366in}}%
\pgfpathlineto{\pgfqpoint{1.473387in}{0.504349in}}%
\pgfpathlineto{\pgfqpoint{1.480146in}{0.504472in}}%
\pgfpathlineto{\pgfqpoint{1.480602in}{0.504695in}}%
\pgfpathlineto{\pgfqpoint{1.480690in}{0.505298in}}%
\pgfpathlineto{\pgfqpoint{1.480779in}{0.528462in}}%
\pgfpathlineto{\pgfqpoint{1.480939in}{1.327028in}}%
\pgfpathlineto{\pgfqpoint{1.481661in}{1.327001in}}%
\pgfpathlineto{\pgfqpoint{1.484464in}{1.326802in}}%
\pgfpathlineto{\pgfqpoint{1.486937in}{1.326713in}}%
\pgfpathlineto{\pgfqpoint{1.488994in}{1.326374in}}%
\pgfpathlineto{\pgfqpoint{1.489863in}{0.504355in}}%
\pgfpathlineto{\pgfqpoint{1.496292in}{0.504464in}}%
\pgfpathlineto{\pgfqpoint{1.497088in}{0.504701in}}%
\pgfpathlineto{\pgfqpoint{1.497175in}{0.505319in}}%
\pgfpathlineto{\pgfqpoint{1.497265in}{0.530662in}}%
\pgfpathlineto{\pgfqpoint{1.497425in}{1.327033in}}%
\pgfpathlineto{\pgfqpoint{1.498150in}{1.327006in}}%
\pgfpathlineto{\pgfqpoint{1.500953in}{1.326808in}}%
\pgfpathlineto{\pgfqpoint{1.503425in}{1.326719in}}%
\pgfpathlineto{\pgfqpoint{1.505479in}{1.326383in}}%
\pgfpathlineto{\pgfqpoint{1.506346in}{0.504360in}}%
\pgfpathlineto{\pgfqpoint{1.512610in}{0.504453in}}%
\pgfpathlineto{\pgfqpoint{1.513574in}{0.504710in}}%
\pgfpathlineto{\pgfqpoint{1.513663in}{0.505369in}}%
\pgfpathlineto{\pgfqpoint{1.513753in}{0.536644in}}%
\pgfpathlineto{\pgfqpoint{1.513914in}{1.327037in}}%
\pgfpathlineto{\pgfqpoint{1.514643in}{1.327010in}}%
\pgfpathlineto{\pgfqpoint{1.517775in}{1.326795in}}%
\pgfpathlineto{\pgfqpoint{1.520248in}{1.326704in}}%
\pgfpathlineto{\pgfqpoint{1.521964in}{1.326392in}}%
\pgfpathlineto{\pgfqpoint{1.522830in}{0.504365in}}%
\pgfpathlineto{\pgfqpoint{1.528764in}{0.504446in}}%
\pgfpathlineto{\pgfqpoint{1.530060in}{0.504720in}}%
\pgfpathlineto{\pgfqpoint{1.530151in}{0.505426in}}%
\pgfpathlineto{\pgfqpoint{1.530324in}{0.864510in}}%
\pgfpathlineto{\pgfqpoint{1.530404in}{1.327041in}}%
\pgfpathlineto{\pgfqpoint{1.531138in}{1.327013in}}%
\pgfpathlineto{\pgfqpoint{1.534270in}{1.326800in}}%
\pgfpathlineto{\pgfqpoint{1.536743in}{1.326708in}}%
\pgfpathlineto{\pgfqpoint{1.538449in}{1.326399in}}%
\pgfpathlineto{\pgfqpoint{1.539313in}{0.504369in}}%
\pgfpathlineto{\pgfqpoint{1.545248in}{0.504449in}}%
\pgfpathlineto{\pgfqpoint{1.546546in}{0.504728in}}%
\pgfpathlineto{\pgfqpoint{1.546639in}{0.505473in}}%
\pgfpathlineto{\pgfqpoint{1.546812in}{0.873496in}}%
\pgfpathlineto{\pgfqpoint{1.546892in}{1.327044in}}%
\pgfpathlineto{\pgfqpoint{1.547629in}{1.327015in}}%
\pgfpathlineto{\pgfqpoint{1.550761in}{1.326805in}}%
\pgfpathlineto{\pgfqpoint{1.553234in}{1.326711in}}%
\pgfpathlineto{\pgfqpoint{1.554934in}{1.326405in}}%
\pgfpathlineto{\pgfqpoint{1.555797in}{0.504371in}}%
\pgfpathlineto{\pgfqpoint{1.561567in}{0.504452in}}%
\pgfpathlineto{\pgfqpoint{1.563031in}{0.504730in}}%
\pgfpathlineto{\pgfqpoint{1.563124in}{0.505477in}}%
\pgfpathlineto{\pgfqpoint{1.563297in}{0.873561in}}%
\pgfpathlineto{\pgfqpoint{1.563377in}{1.327046in}}%
\pgfpathlineto{\pgfqpoint{1.564114in}{1.327017in}}%
\pgfpathlineto{\pgfqpoint{1.567246in}{1.326807in}}%
\pgfpathlineto{\pgfqpoint{1.569719in}{1.326713in}}%
\pgfpathlineto{\pgfqpoint{1.571418in}{1.326407in}}%
\pgfpathlineto{\pgfqpoint{1.572282in}{0.504372in}}%
\pgfpathlineto{\pgfqpoint{1.578051in}{0.504453in}}%
\pgfpathlineto{\pgfqpoint{1.579516in}{0.504731in}}%
\pgfpathlineto{\pgfqpoint{1.579608in}{0.505478in}}%
\pgfpathlineto{\pgfqpoint{1.579782in}{0.873578in}}%
\pgfpathlineto{\pgfqpoint{1.579862in}{1.327047in}}%
\pgfpathlineto{\pgfqpoint{1.580599in}{1.327018in}}%
\pgfpathlineto{\pgfqpoint{1.583731in}{1.326807in}}%
\pgfpathlineto{\pgfqpoint{1.586204in}{1.326713in}}%
\pgfpathlineto{\pgfqpoint{1.587903in}{1.326407in}}%
\pgfpathlineto{\pgfqpoint{1.588767in}{0.504372in}}%
\pgfpathlineto{\pgfqpoint{1.594536in}{0.504452in}}%
\pgfpathlineto{\pgfqpoint{1.596001in}{0.504730in}}%
\pgfpathlineto{\pgfqpoint{1.596093in}{0.505476in}}%
\pgfpathlineto{\pgfqpoint{1.596267in}{0.873556in}}%
\pgfpathlineto{\pgfqpoint{1.596347in}{1.327046in}}%
\pgfpathlineto{\pgfqpoint{1.597084in}{1.327017in}}%
\pgfpathlineto{\pgfqpoint{1.600216in}{1.326806in}}%
\pgfpathlineto{\pgfqpoint{1.602688in}{1.326712in}}%
\pgfpathlineto{\pgfqpoint{1.604388in}{1.326404in}}%
\pgfpathlineto{\pgfqpoint{1.605252in}{0.504371in}}%
\pgfpathlineto{\pgfqpoint{1.611022in}{0.504451in}}%
\pgfpathlineto{\pgfqpoint{1.612486in}{0.504729in}}%
\pgfpathlineto{\pgfqpoint{1.612578in}{0.505473in}}%
\pgfpathlineto{\pgfqpoint{1.612752in}{0.873503in}}%
\pgfpathlineto{\pgfqpoint{1.612831in}{1.327044in}}%
\pgfpathlineto{\pgfqpoint{1.613568in}{1.327015in}}%
\pgfpathlineto{\pgfqpoint{1.616701in}{1.326803in}}%
\pgfpathlineto{\pgfqpoint{1.619173in}{1.326709in}}%
\pgfpathlineto{\pgfqpoint{1.620873in}{1.326401in}}%
\pgfpathlineto{\pgfqpoint{1.621737in}{0.504369in}}%
\pgfpathlineto{\pgfqpoint{1.627672in}{0.504447in}}%
\pgfpathlineto{\pgfqpoint{1.628971in}{0.504726in}}%
\pgfpathlineto{\pgfqpoint{1.629063in}{0.505470in}}%
\pgfpathlineto{\pgfqpoint{1.629237in}{0.873433in}}%
\pgfpathlineto{\pgfqpoint{1.629316in}{1.327042in}}%
\pgfpathlineto{\pgfqpoint{1.630053in}{1.327013in}}%
\pgfpathlineto{\pgfqpoint{1.633185in}{1.326800in}}%
\pgfpathlineto{\pgfqpoint{1.635658in}{1.326706in}}%
\pgfpathlineto{\pgfqpoint{1.637358in}{1.326397in}}%
\pgfpathlineto{\pgfqpoint{1.638223in}{0.504367in}}%
\pgfpathlineto{\pgfqpoint{1.644157in}{0.504446in}}%
\pgfpathlineto{\pgfqpoint{1.645455in}{0.504722in}}%
\pgfpathlineto{\pgfqpoint{1.645547in}{0.505449in}}%
\pgfpathlineto{\pgfqpoint{1.645720in}{0.869697in}}%
\pgfpathlineto{\pgfqpoint{1.645800in}{1.327040in}}%
\pgfpathlineto{\pgfqpoint{1.646536in}{1.327011in}}%
\pgfpathlineto{\pgfqpoint{1.649668in}{1.326797in}}%
\pgfpathlineto{\pgfqpoint{1.652141in}{1.326703in}}%
\pgfpathlineto{\pgfqpoint{1.653843in}{1.326392in}}%
\pgfpathlineto{\pgfqpoint{1.654708in}{0.504364in}}%
\pgfpathlineto{\pgfqpoint{1.660643in}{0.504444in}}%
\pgfpathlineto{\pgfqpoint{1.661939in}{0.504717in}}%
\pgfpathlineto{\pgfqpoint{1.662030in}{0.505417in}}%
\pgfpathlineto{\pgfqpoint{1.662203in}{0.863571in}}%
\pgfpathlineto{\pgfqpoint{1.662282in}{1.327038in}}%
\pgfpathlineto{\pgfqpoint{1.663017in}{1.327010in}}%
\pgfpathlineto{\pgfqpoint{1.666149in}{1.326794in}}%
\pgfpathlineto{\pgfqpoint{1.668621in}{1.326701in}}%
\pgfpathlineto{\pgfqpoint{1.670328in}{1.326389in}}%
\pgfpathlineto{\pgfqpoint{1.671194in}{0.504363in}}%
\pgfpathlineto{\pgfqpoint{1.677293in}{0.504456in}}%
\pgfpathlineto{\pgfqpoint{1.678423in}{0.504713in}}%
\pgfpathlineto{\pgfqpoint{1.678513in}{0.505393in}}%
\pgfpathlineto{\pgfqpoint{1.678686in}{0.858579in}}%
\pgfpathlineto{\pgfqpoint{1.678765in}{1.327037in}}%
\pgfpathlineto{\pgfqpoint{1.679497in}{1.327008in}}%
\pgfpathlineto{\pgfqpoint{1.682629in}{1.326792in}}%
\pgfpathlineto{\pgfqpoint{1.685102in}{1.326699in}}%
\pgfpathlineto{\pgfqpoint{1.686812in}{1.326386in}}%
\pgfpathlineto{\pgfqpoint{1.687679in}{0.504361in}}%
\pgfpathlineto{\pgfqpoint{1.693943in}{0.504453in}}%
\pgfpathlineto{\pgfqpoint{1.694908in}{0.504710in}}%
\pgfpathlineto{\pgfqpoint{1.694997in}{0.505376in}}%
\pgfpathlineto{\pgfqpoint{1.695087in}{0.537973in}}%
\pgfpathlineto{\pgfqpoint{1.695248in}{1.327036in}}%
\pgfpathlineto{\pgfqpoint{1.695979in}{1.327008in}}%
\pgfpathlineto{\pgfqpoint{1.699111in}{1.326791in}}%
\pgfpathlineto{\pgfqpoint{1.701583in}{1.326698in}}%
\pgfpathlineto{\pgfqpoint{1.703297in}{1.326384in}}%
\pgfpathlineto{\pgfqpoint{1.704164in}{0.504360in}}%
\pgfpathlineto{\pgfqpoint{1.710428in}{0.504452in}}%
\pgfpathlineto{\pgfqpoint{1.711392in}{0.504708in}}%
\pgfpathlineto{\pgfqpoint{1.711482in}{0.505367in}}%
\pgfpathlineto{\pgfqpoint{1.711571in}{0.536778in}}%
\pgfpathlineto{\pgfqpoint{1.711732in}{1.327035in}}%
\pgfpathlineto{\pgfqpoint{1.712462in}{1.327007in}}%
\pgfpathlineto{\pgfqpoint{1.715594in}{1.326791in}}%
\pgfpathlineto{\pgfqpoint{1.718067in}{1.326698in}}%
\pgfpathlineto{\pgfqpoint{1.719782in}{1.326384in}}%
\pgfpathlineto{\pgfqpoint{1.720649in}{0.504360in}}%
\pgfpathlineto{\pgfqpoint{1.726913in}{0.504452in}}%
\pgfpathlineto{\pgfqpoint{1.727877in}{0.504708in}}%
\pgfpathlineto{\pgfqpoint{1.727966in}{0.505366in}}%
\pgfpathlineto{\pgfqpoint{1.728056in}{0.536579in}}%
\pgfpathlineto{\pgfqpoint{1.728217in}{1.327035in}}%
\pgfpathlineto{\pgfqpoint{1.728946in}{1.327007in}}%
\pgfpathlineto{\pgfqpoint{1.732078in}{1.326791in}}%
\pgfpathlineto{\pgfqpoint{1.734551in}{1.326699in}}%
\pgfpathlineto{\pgfqpoint{1.736267in}{1.326385in}}%
\pgfpathlineto{\pgfqpoint{1.737134in}{0.504361in}}%
\pgfpathlineto{\pgfqpoint{1.743398in}{0.504453in}}%
\pgfpathlineto{\pgfqpoint{1.744362in}{0.504710in}}%
\pgfpathlineto{\pgfqpoint{1.744451in}{0.505371in}}%
\pgfpathlineto{\pgfqpoint{1.744541in}{0.537243in}}%
\pgfpathlineto{\pgfqpoint{1.744702in}{1.327036in}}%
\pgfpathlineto{\pgfqpoint{1.745432in}{1.327008in}}%
\pgfpathlineto{\pgfqpoint{1.748564in}{1.326792in}}%
\pgfpathlineto{\pgfqpoint{1.751037in}{1.326700in}}%
\pgfpathlineto{\pgfqpoint{1.752752in}{1.326386in}}%
\pgfpathlineto{\pgfqpoint{1.753618in}{0.504361in}}%
\pgfpathlineto{\pgfqpoint{1.759883in}{0.504453in}}%
\pgfpathlineto{\pgfqpoint{1.760847in}{0.504712in}}%
\pgfpathlineto{\pgfqpoint{1.760937in}{0.505382in}}%
\pgfpathlineto{\pgfqpoint{1.761109in}{0.856029in}}%
\pgfpathlineto{\pgfqpoint{1.761188in}{1.327037in}}%
\pgfpathlineto{\pgfqpoint{1.761919in}{1.327009in}}%
\pgfpathlineto{\pgfqpoint{1.765051in}{1.326793in}}%
\pgfpathlineto{\pgfqpoint{1.767524in}{1.326701in}}%
\pgfpathlineto{\pgfqpoint{1.769237in}{1.326388in}}%
\pgfpathlineto{\pgfqpoint{1.770103in}{0.504363in}}%
\pgfpathlineto{\pgfqpoint{1.776202in}{0.504457in}}%
\pgfpathlineto{\pgfqpoint{1.777332in}{0.504714in}}%
\pgfpathlineto{\pgfqpoint{1.777422in}{0.505395in}}%
\pgfpathlineto{\pgfqpoint{1.777595in}{0.858753in}}%
\pgfpathlineto{\pgfqpoint{1.777674in}{1.327038in}}%
\pgfpathlineto{\pgfqpoint{1.778406in}{1.327010in}}%
\pgfpathlineto{\pgfqpoint{1.781538in}{1.326795in}}%
\pgfpathlineto{\pgfqpoint{1.784011in}{1.326703in}}%
\pgfpathlineto{\pgfqpoint{1.785722in}{1.326391in}}%
\pgfpathlineto{\pgfqpoint{1.786587in}{0.504364in}}%
\pgfpathlineto{\pgfqpoint{1.792687in}{0.504458in}}%
\pgfpathlineto{\pgfqpoint{1.793818in}{0.504717in}}%
\pgfpathlineto{\pgfqpoint{1.793908in}{0.505410in}}%
\pgfpathlineto{\pgfqpoint{1.794081in}{0.861717in}}%
\pgfpathlineto{\pgfqpoint{1.794160in}{1.327039in}}%
\pgfpathlineto{\pgfqpoint{1.794894in}{1.327011in}}%
\pgfpathlineto{\pgfqpoint{1.798026in}{1.326796in}}%
\pgfpathlineto{\pgfqpoint{1.800499in}{1.326704in}}%
\pgfpathlineto{\pgfqpoint{1.802206in}{1.326393in}}%
\pgfpathlineto{\pgfqpoint{1.803072in}{0.504365in}}%
\pgfpathlineto{\pgfqpoint{1.809007in}{0.504445in}}%
\pgfpathlineto{\pgfqpoint{1.810303in}{0.504719in}}%
\pgfpathlineto{\pgfqpoint{1.810394in}{0.505424in}}%
\pgfpathlineto{\pgfqpoint{1.810567in}{0.864498in}}%
\pgfpathlineto{\pgfqpoint{1.810646in}{1.327040in}}%
\pgfpathlineto{\pgfqpoint{1.811381in}{1.327012in}}%
\pgfpathlineto{\pgfqpoint{1.814513in}{1.326797in}}%
\pgfpathlineto{\pgfqpoint{1.816986in}{1.326705in}}%
\pgfpathlineto{\pgfqpoint{1.818691in}{1.326394in}}%
\pgfpathlineto{\pgfqpoint{1.819557in}{0.504366in}}%
\pgfpathlineto{\pgfqpoint{1.825491in}{0.504446in}}%
\pgfpathlineto{\pgfqpoint{1.826788in}{0.504721in}}%
\pgfpathlineto{\pgfqpoint{1.826879in}{0.505435in}}%
\pgfpathlineto{\pgfqpoint{1.827052in}{0.866713in}}%
\pgfpathlineto{\pgfqpoint{1.827132in}{1.327041in}}%
\pgfpathlineto{\pgfqpoint{1.827867in}{1.327012in}}%
\pgfpathlineto{\pgfqpoint{1.830999in}{1.326798in}}%
\pgfpathlineto{\pgfqpoint{1.833472in}{1.326705in}}%
\pgfpathlineto{\pgfqpoint{1.835176in}{1.326395in}}%
\pgfpathlineto{\pgfqpoint{1.836041in}{0.504366in}}%
\pgfpathlineto{\pgfqpoint{1.841976in}{0.504446in}}%
\pgfpathlineto{\pgfqpoint{1.843273in}{0.504722in}}%
\pgfpathlineto{\pgfqpoint{1.843365in}{0.505442in}}%
\pgfpathlineto{\pgfqpoint{1.843538in}{0.868105in}}%
\pgfpathlineto{\pgfqpoint{1.843617in}{1.327041in}}%
\pgfpathlineto{\pgfqpoint{1.844353in}{1.327012in}}%
\pgfpathlineto{\pgfqpoint{1.847485in}{1.326799in}}%
\pgfpathlineto{\pgfqpoint{1.849958in}{1.326705in}}%
\pgfpathlineto{\pgfqpoint{1.851661in}{1.326395in}}%
\pgfpathlineto{\pgfqpoint{1.852526in}{0.504366in}}%
\pgfpathlineto{\pgfqpoint{1.858461in}{0.504446in}}%
\pgfpathlineto{\pgfqpoint{1.859758in}{0.504722in}}%
\pgfpathlineto{\pgfqpoint{1.859849in}{0.505444in}}%
\pgfpathlineto{\pgfqpoint{1.860023in}{0.868581in}}%
\pgfpathlineto{\pgfqpoint{1.860102in}{1.327041in}}%
\pgfpathlineto{\pgfqpoint{1.860838in}{1.327012in}}%
\pgfpathlineto{\pgfqpoint{1.863970in}{1.326799in}}%
\pgfpathlineto{\pgfqpoint{1.866443in}{1.326705in}}%
\pgfpathlineto{\pgfqpoint{1.868146in}{1.326395in}}%
\pgfpathlineto{\pgfqpoint{1.869011in}{0.504366in}}%
\pgfpathlineto{\pgfqpoint{1.874945in}{0.504446in}}%
\pgfpathlineto{\pgfqpoint{1.876243in}{0.504722in}}%
\pgfpathlineto{\pgfqpoint{1.876334in}{0.505442in}}%
\pgfpathlineto{\pgfqpoint{1.876508in}{0.868200in}}%
\pgfpathlineto{\pgfqpoint{1.876587in}{1.327041in}}%
\pgfpathlineto{\pgfqpoint{1.877323in}{1.327012in}}%
\pgfpathlineto{\pgfqpoint{1.880455in}{1.326798in}}%
\pgfpathlineto{\pgfqpoint{1.882928in}{1.326705in}}%
\pgfpathlineto{\pgfqpoint{1.884631in}{1.326395in}}%
\pgfpathlineto{\pgfqpoint{1.885496in}{0.504366in}}%
\pgfpathlineto{\pgfqpoint{1.891430in}{0.504446in}}%
\pgfpathlineto{\pgfqpoint{1.892727in}{0.504721in}}%
\pgfpathlineto{\pgfqpoint{1.892819in}{0.505437in}}%
\pgfpathlineto{\pgfqpoint{1.892992in}{0.867150in}}%
\pgfpathlineto{\pgfqpoint{1.893072in}{1.327040in}}%
\pgfpathlineto{\pgfqpoint{1.893807in}{1.327012in}}%
\pgfpathlineto{\pgfqpoint{1.896939in}{1.326798in}}%
\pgfpathlineto{\pgfqpoint{1.899412in}{1.326704in}}%
\pgfpathlineto{\pgfqpoint{1.901115in}{1.326394in}}%
\pgfpathlineto{\pgfqpoint{1.901981in}{0.504365in}}%
\pgfpathlineto{\pgfqpoint{1.907915in}{0.504445in}}%
\pgfpathlineto{\pgfqpoint{1.909212in}{0.504719in}}%
\pgfpathlineto{\pgfqpoint{1.909303in}{0.505429in}}%
\pgfpathlineto{\pgfqpoint{1.909476in}{0.865691in}}%
\pgfpathlineto{\pgfqpoint{1.909556in}{1.327040in}}%
\pgfpathlineto{\pgfqpoint{1.910291in}{1.327011in}}%
\pgfpathlineto{\pgfqpoint{1.913423in}{1.326797in}}%
\pgfpathlineto{\pgfqpoint{1.915896in}{1.326704in}}%
\pgfpathlineto{\pgfqpoint{1.917600in}{1.326392in}}%
\pgfpathlineto{\pgfqpoint{1.918466in}{0.504365in}}%
\pgfpathlineto{\pgfqpoint{1.924400in}{0.504445in}}%
\pgfpathlineto{\pgfqpoint{1.925697in}{0.504718in}}%
\pgfpathlineto{\pgfqpoint{1.925788in}{0.505421in}}%
\pgfpathlineto{\pgfqpoint{1.925961in}{0.864097in}}%
\pgfpathlineto{\pgfqpoint{1.926040in}{1.327039in}}%
\pgfpathlineto{\pgfqpoint{1.926774in}{1.327011in}}%
\pgfpathlineto{\pgfqpoint{1.929907in}{1.326796in}}%
\pgfpathlineto{\pgfqpoint{1.932379in}{1.326703in}}%
\pgfpathlineto{\pgfqpoint{1.934085in}{1.326391in}}%
\pgfpathlineto{\pgfqpoint{1.934951in}{0.504364in}}%
\pgfpathlineto{\pgfqpoint{1.941050in}{0.504458in}}%
\pgfpathlineto{\pgfqpoint{1.942181in}{0.504717in}}%
\pgfpathlineto{\pgfqpoint{1.942272in}{0.505413in}}%
\pgfpathlineto{\pgfqpoint{1.942445in}{0.862612in}}%
\pgfpathlineto{\pgfqpoint{1.942524in}{1.327039in}}%
\pgfpathlineto{\pgfqpoint{1.943258in}{1.327010in}}%
\pgfpathlineto{\pgfqpoint{1.946390in}{1.326795in}}%
\pgfpathlineto{\pgfqpoint{1.948863in}{1.326702in}}%
\pgfpathlineto{\pgfqpoint{1.950570in}{1.326391in}}%
\pgfpathlineto{\pgfqpoint{1.951436in}{0.504364in}}%
\pgfpathlineto{\pgfqpoint{1.957535in}{0.504457in}}%
\pgfpathlineto{\pgfqpoint{1.958666in}{0.504716in}}%
\pgfpathlineto{\pgfqpoint{1.958757in}{0.505408in}}%
\pgfpathlineto{\pgfqpoint{1.958929in}{0.861418in}}%
\pgfpathlineto{\pgfqpoint{1.959009in}{1.327038in}}%
\pgfpathlineto{\pgfqpoint{1.959742in}{1.327010in}}%
\pgfpathlineto{\pgfqpoint{1.962874in}{1.326795in}}%
\pgfpathlineto{\pgfqpoint{1.965347in}{1.326702in}}%
\pgfpathlineto{\pgfqpoint{1.967055in}{1.326390in}}%
\pgfpathlineto{\pgfqpoint{1.967921in}{0.504363in}}%
\pgfpathlineto{\pgfqpoint{1.974020in}{0.504457in}}%
\pgfpathlineto{\pgfqpoint{1.975151in}{0.504715in}}%
\pgfpathlineto{\pgfqpoint{1.975241in}{0.505404in}}%
\pgfpathlineto{\pgfqpoint{1.975414in}{0.860620in}}%
\pgfpathlineto{\pgfqpoint{1.975493in}{1.327038in}}%
\pgfpathlineto{\pgfqpoint{1.976226in}{1.327010in}}%
\pgfpathlineto{\pgfqpoint{1.979358in}{1.326795in}}%
\pgfpathlineto{\pgfqpoint{1.981831in}{1.326702in}}%
\pgfpathlineto{\pgfqpoint{1.983540in}{1.326390in}}%
\pgfpathlineto{\pgfqpoint{1.984406in}{0.504363in}}%
\pgfpathlineto{\pgfqpoint{1.990505in}{0.504457in}}%
\pgfpathlineto{\pgfqpoint{1.991635in}{0.504715in}}%
\pgfpathlineto{\pgfqpoint{1.991726in}{0.505402in}}%
\pgfpathlineto{\pgfqpoint{1.991899in}{0.860258in}}%
\pgfpathlineto{\pgfqpoint{1.991978in}{1.327038in}}%
\pgfpathlineto{\pgfqpoint{1.992710in}{1.327010in}}%
\pgfpathlineto{\pgfqpoint{1.995843in}{1.326795in}}%
\pgfpathlineto{\pgfqpoint{1.998315in}{1.326702in}}%
\pgfpathlineto{\pgfqpoint{2.000025in}{1.326390in}}%
\pgfpathlineto{\pgfqpoint{2.000891in}{0.504363in}}%
\pgfpathlineto{\pgfqpoint{2.006990in}{0.504457in}}%
\pgfpathlineto{\pgfqpoint{2.008120in}{0.504715in}}%
\pgfpathlineto{\pgfqpoint{2.008211in}{0.505402in}}%
\pgfpathlineto{\pgfqpoint{2.008384in}{0.860312in}}%
\pgfpathlineto{\pgfqpoint{2.008463in}{1.327038in}}%
\pgfpathlineto{\pgfqpoint{2.009195in}{1.327010in}}%
\pgfpathlineto{\pgfqpoint{2.012328in}{1.326795in}}%
\pgfpathlineto{\pgfqpoint{2.014800in}{1.326702in}}%
\pgfpathlineto{\pgfqpoint{2.016509in}{1.326390in}}%
\pgfpathlineto{\pgfqpoint{2.017375in}{0.504363in}}%
\pgfpathlineto{\pgfqpoint{2.023475in}{0.504457in}}%
\pgfpathlineto{\pgfqpoint{2.024605in}{0.504715in}}%
\pgfpathlineto{\pgfqpoint{2.024696in}{0.505404in}}%
\pgfpathlineto{\pgfqpoint{2.024869in}{0.860716in}}%
\pgfpathlineto{\pgfqpoint{2.024948in}{1.327038in}}%
\pgfpathlineto{\pgfqpoint{2.025681in}{1.327010in}}%
\pgfpathlineto{\pgfqpoint{2.028813in}{1.326795in}}%
\pgfpathlineto{\pgfqpoint{2.031286in}{1.326702in}}%
\pgfpathlineto{\pgfqpoint{2.032994in}{1.326390in}}%
\pgfpathlineto{\pgfqpoint{2.033860in}{0.504364in}}%
\pgfpathlineto{\pgfqpoint{2.039960in}{0.504458in}}%
\pgfpathlineto{\pgfqpoint{2.041090in}{0.504716in}}%
\pgfpathlineto{\pgfqpoint{2.041181in}{0.505408in}}%
\pgfpathlineto{\pgfqpoint{2.041354in}{0.861365in}}%
\pgfpathlineto{\pgfqpoint{2.041433in}{1.327039in}}%
\pgfpathlineto{\pgfqpoint{2.042166in}{1.327010in}}%
\pgfpathlineto{\pgfqpoint{2.045298in}{1.326795in}}%
\pgfpathlineto{\pgfqpoint{2.047771in}{1.326703in}}%
\pgfpathlineto{\pgfqpoint{2.049479in}{1.326391in}}%
\pgfpathlineto{\pgfqpoint{2.050345in}{0.504364in}}%
\pgfpathlineto{\pgfqpoint{2.056444in}{0.504458in}}%
\pgfpathlineto{\pgfqpoint{2.057575in}{0.504717in}}%
\pgfpathlineto{\pgfqpoint{2.057666in}{0.505412in}}%
\pgfpathlineto{\pgfqpoint{2.057839in}{0.862141in}}%
\pgfpathlineto{\pgfqpoint{2.057918in}{1.327039in}}%
\pgfpathlineto{\pgfqpoint{2.058652in}{1.327011in}}%
\pgfpathlineto{\pgfqpoint{2.061784in}{1.326796in}}%
\pgfpathlineto{\pgfqpoint{2.064257in}{1.326703in}}%
\pgfpathlineto{\pgfqpoint{2.065964in}{1.326392in}}%
\pgfpathlineto{\pgfqpoint{2.066830in}{0.504364in}}%
\pgfpathlineto{\pgfqpoint{2.072764in}{0.504445in}}%
\pgfpathlineto{\pgfqpoint{2.074060in}{0.504717in}}%
\pgfpathlineto{\pgfqpoint{2.074151in}{0.505416in}}%
\pgfpathlineto{\pgfqpoint{2.074324in}{0.862922in}}%
\pgfpathlineto{\pgfqpoint{2.074403in}{1.327039in}}%
\pgfpathlineto{\pgfqpoint{2.075137in}{1.327011in}}%
\pgfpathlineto{\pgfqpoint{2.078269in}{1.326796in}}%
\pgfpathlineto{\pgfqpoint{2.080742in}{1.326703in}}%
\pgfpathlineto{\pgfqpoint{2.082449in}{1.326392in}}%
\pgfpathlineto{\pgfqpoint{2.083315in}{0.504364in}}%
\pgfpathlineto{\pgfqpoint{2.089249in}{0.504445in}}%
\pgfpathlineto{\pgfqpoint{2.090545in}{0.504718in}}%
\pgfpathlineto{\pgfqpoint{2.090636in}{0.505419in}}%
\pgfpathlineto{\pgfqpoint{2.090809in}{0.863599in}}%
\pgfpathlineto{\pgfqpoint{2.090888in}{1.327039in}}%
\pgfpathlineto{\pgfqpoint{2.091623in}{1.327011in}}%
\pgfpathlineto{\pgfqpoint{2.094755in}{1.326796in}}%
\pgfpathlineto{\pgfqpoint{2.097227in}{1.326704in}}%
\pgfpathlineto{\pgfqpoint{2.098934in}{1.326392in}}%
\pgfpathlineto{\pgfqpoint{2.099799in}{0.504365in}}%
\pgfpathlineto{\pgfqpoint{2.105734in}{0.504445in}}%
\pgfpathlineto{\pgfqpoint{2.107030in}{0.504718in}}%
\pgfpathlineto{\pgfqpoint{2.107121in}{0.505421in}}%
\pgfpathlineto{\pgfqpoint{2.107294in}{0.864090in}}%
\pgfpathlineto{\pgfqpoint{2.107373in}{1.327039in}}%
\pgfpathlineto{\pgfqpoint{2.108108in}{1.327011in}}%
\pgfpathlineto{\pgfqpoint{2.111240in}{1.326797in}}%
\pgfpathlineto{\pgfqpoint{2.113713in}{1.326704in}}%
\pgfpathlineto{\pgfqpoint{2.115418in}{1.326393in}}%
\pgfpathlineto{\pgfqpoint{2.116284in}{0.504365in}}%
\pgfpathlineto{\pgfqpoint{2.122219in}{0.504445in}}%
\pgfpathlineto{\pgfqpoint{2.123515in}{0.504719in}}%
\pgfpathlineto{\pgfqpoint{2.123606in}{0.505423in}}%
\pgfpathlineto{\pgfqpoint{2.123779in}{0.864356in}}%
\pgfpathlineto{\pgfqpoint{2.123858in}{1.327040in}}%
\pgfpathlineto{\pgfqpoint{2.124593in}{1.327011in}}%
\pgfpathlineto{\pgfqpoint{2.127725in}{1.326797in}}%
\pgfpathlineto{\pgfqpoint{2.130198in}{1.326704in}}%
\pgfpathlineto{\pgfqpoint{2.131903in}{1.326393in}}%
\pgfpathlineto{\pgfqpoint{2.132769in}{0.504365in}}%
\pgfpathlineto{\pgfqpoint{2.138703in}{0.504445in}}%
\pgfpathlineto{\pgfqpoint{2.140000in}{0.504719in}}%
\pgfpathlineto{\pgfqpoint{2.140091in}{0.505423in}}%
\pgfpathlineto{\pgfqpoint{2.140264in}{0.864391in}}%
\pgfpathlineto{\pgfqpoint{2.140343in}{1.327039in}}%
\pgfpathlineto{\pgfqpoint{2.141078in}{1.327011in}}%
\pgfpathlineto{\pgfqpoint{2.144210in}{1.326797in}}%
\pgfpathlineto{\pgfqpoint{2.146683in}{1.326704in}}%
\pgfpathlineto{\pgfqpoint{2.148388in}{1.326393in}}%
\pgfpathlineto{\pgfqpoint{2.149254in}{0.504365in}}%
\pgfpathlineto{\pgfqpoint{2.155188in}{0.504445in}}%
\pgfpathlineto{\pgfqpoint{2.156485in}{0.504718in}}%
\pgfpathlineto{\pgfqpoint{2.156576in}{0.505422in}}%
\pgfpathlineto{\pgfqpoint{2.156749in}{0.864226in}}%
\pgfpathlineto{\pgfqpoint{2.156828in}{1.327039in}}%
\pgfpathlineto{\pgfqpoint{2.157562in}{1.327011in}}%
\pgfpathlineto{\pgfqpoint{2.160695in}{1.326796in}}%
\pgfpathlineto{\pgfqpoint{2.163167in}{1.326704in}}%
\pgfpathlineto{\pgfqpoint{2.164873in}{1.326392in}}%
\pgfpathlineto{\pgfqpoint{2.165739in}{0.504365in}}%
\pgfpathlineto{\pgfqpoint{2.171673in}{0.504445in}}%
\pgfpathlineto{\pgfqpoint{2.172969in}{0.504718in}}%
\pgfpathlineto{\pgfqpoint{2.173060in}{0.505420in}}%
\pgfpathlineto{\pgfqpoint{2.173233in}{0.863915in}}%
\pgfpathlineto{\pgfqpoint{2.173313in}{1.327039in}}%
\pgfpathlineto{\pgfqpoint{2.174047in}{1.327011in}}%
\pgfpathlineto{\pgfqpoint{2.177179in}{1.326796in}}%
\pgfpathlineto{\pgfqpoint{2.179652in}{1.326703in}}%
\pgfpathlineto{\pgfqpoint{2.181358in}{1.326392in}}%
\pgfpathlineto{\pgfqpoint{2.182224in}{0.504364in}}%
\pgfpathlineto{\pgfqpoint{2.188158in}{0.504445in}}%
\pgfpathlineto{\pgfqpoint{2.189454in}{0.504718in}}%
\pgfpathlineto{\pgfqpoint{2.189545in}{0.505418in}}%
\pgfpathlineto{\pgfqpoint{2.189718in}{0.863523in}}%
\pgfpathlineto{\pgfqpoint{2.189797in}{1.327039in}}%
\pgfpathlineto{\pgfqpoint{2.190532in}{1.327011in}}%
\pgfpathlineto{\pgfqpoint{2.193664in}{1.326796in}}%
\pgfpathlineto{\pgfqpoint{2.196136in}{1.326703in}}%
\pgfpathlineto{\pgfqpoint{2.197843in}{1.326392in}}%
\pgfpathlineto{\pgfqpoint{2.198708in}{0.504364in}}%
\pgfpathlineto{\pgfqpoint{2.204643in}{0.504445in}}%
\pgfpathlineto{\pgfqpoint{2.205939in}{0.504717in}}%
\pgfpathlineto{\pgfqpoint{2.206030in}{0.505416in}}%
\pgfpathlineto{\pgfqpoint{2.206203in}{0.863116in}}%
\pgfpathlineto{\pgfqpoint{2.206282in}{1.327039in}}%
\pgfpathlineto{\pgfqpoint{2.207016in}{1.327011in}}%
\pgfpathlineto{\pgfqpoint{2.210148in}{1.326796in}}%
\pgfpathlineto{\pgfqpoint{2.212621in}{1.326703in}}%
\pgfpathlineto{\pgfqpoint{2.214328in}{1.326391in}}%
\pgfpathlineto{\pgfqpoint{2.215193in}{0.504364in}}%
\pgfpathlineto{\pgfqpoint{2.221128in}{0.504444in}}%
\pgfpathlineto{\pgfqpoint{2.222424in}{0.504717in}}%
\pgfpathlineto{\pgfqpoint{2.222515in}{0.505414in}}%
\pgfpathlineto{\pgfqpoint{2.222687in}{0.862753in}}%
\pgfpathlineto{\pgfqpoint{2.222767in}{1.327039in}}%
\pgfpathlineto{\pgfqpoint{2.223501in}{1.327011in}}%
\pgfpathlineto{\pgfqpoint{2.226633in}{1.326796in}}%
\pgfpathlineto{\pgfqpoint{2.229106in}{1.326703in}}%
\pgfpathlineto{\pgfqpoint{2.230812in}{1.326391in}}%
\pgfpathlineto{\pgfqpoint{2.231678in}{0.504364in}}%
\pgfpathlineto{\pgfqpoint{2.237778in}{0.504458in}}%
\pgfpathlineto{\pgfqpoint{2.238909in}{0.504717in}}%
\pgfpathlineto{\pgfqpoint{2.238999in}{0.505413in}}%
\pgfpathlineto{\pgfqpoint{2.239172in}{0.862476in}}%
\pgfpathlineto{\pgfqpoint{2.239252in}{1.327039in}}%
\pgfpathlineto{\pgfqpoint{2.239985in}{1.327010in}}%
\pgfpathlineto{\pgfqpoint{2.243118in}{1.326796in}}%
\pgfpathlineto{\pgfqpoint{2.245590in}{1.326703in}}%
\pgfpathlineto{\pgfqpoint{2.247297in}{1.326391in}}%
\pgfpathlineto{\pgfqpoint{2.248163in}{0.504364in}}%
\pgfpathlineto{\pgfqpoint{2.254263in}{0.504458in}}%
\pgfpathlineto{\pgfqpoint{2.255393in}{0.504717in}}%
\pgfpathlineto{\pgfqpoint{2.255484in}{0.505412in}}%
\pgfpathlineto{\pgfqpoint{2.255657in}{0.862309in}}%
\pgfpathlineto{\pgfqpoint{2.255736in}{1.327039in}}%
\pgfpathlineto{\pgfqpoint{2.256470in}{1.327010in}}%
\pgfpathlineto{\pgfqpoint{2.259602in}{1.326796in}}%
\pgfpathlineto{\pgfqpoint{2.262075in}{1.326703in}}%
\pgfpathlineto{\pgfqpoint{2.263782in}{1.326391in}}%
\pgfpathlineto{\pgfqpoint{2.264648in}{0.504364in}}%
\pgfpathlineto{\pgfqpoint{2.270747in}{0.504458in}}%
\pgfpathlineto{\pgfqpoint{2.271878in}{0.504717in}}%
\pgfpathlineto{\pgfqpoint{2.271969in}{0.505412in}}%
\pgfpathlineto{\pgfqpoint{2.272142in}{0.862258in}}%
\pgfpathlineto{\pgfqpoint{2.272221in}{1.327039in}}%
\pgfpathlineto{\pgfqpoint{2.272955in}{1.327010in}}%
\pgfpathlineto{\pgfqpoint{2.276087in}{1.326796in}}%
\pgfpathlineto{\pgfqpoint{2.278560in}{1.326703in}}%
\pgfpathlineto{\pgfqpoint{2.280267in}{1.326391in}}%
\pgfpathlineto{\pgfqpoint{2.281133in}{0.504364in}}%
\pgfpathlineto{\pgfqpoint{2.287232in}{0.504458in}}%
\pgfpathlineto{\pgfqpoint{2.288363in}{0.504717in}}%
\pgfpathlineto{\pgfqpoint{2.288454in}{0.505412in}}%
\pgfpathlineto{\pgfqpoint{2.288627in}{0.862311in}}%
\pgfpathlineto{\pgfqpoint{2.288706in}{1.327039in}}%
\pgfpathlineto{\pgfqpoint{2.289440in}{1.327010in}}%
\pgfpathlineto{\pgfqpoint{2.292572in}{1.326796in}}%
\pgfpathlineto{\pgfqpoint{2.295045in}{1.326703in}}%
\pgfpathlineto{\pgfqpoint{2.296752in}{1.326391in}}%
\pgfpathlineto{\pgfqpoint{2.297618in}{0.504364in}}%
\pgfpathlineto{\pgfqpoint{2.303717in}{0.504458in}}%
\pgfpathlineto{\pgfqpoint{2.304848in}{0.504717in}}%
\pgfpathlineto{\pgfqpoint{2.304939in}{0.505413in}}%
\pgfpathlineto{\pgfqpoint{2.305112in}{0.862447in}}%
\pgfpathlineto{\pgfqpoint{2.305191in}{1.327039in}}%
\pgfpathlineto{\pgfqpoint{2.305925in}{1.327011in}}%
\pgfpathlineto{\pgfqpoint{2.309057in}{1.326796in}}%
\pgfpathlineto{\pgfqpoint{2.311530in}{1.326703in}}%
\pgfpathlineto{\pgfqpoint{2.313237in}{1.326391in}}%
\pgfpathlineto{\pgfqpoint{2.314102in}{0.504364in}}%
\pgfpathlineto{\pgfqpoint{2.320202in}{0.504458in}}%
\pgfpathlineto{\pgfqpoint{2.321333in}{0.504717in}}%
\pgfpathlineto{\pgfqpoint{2.321424in}{0.505414in}}%
\pgfpathlineto{\pgfqpoint{2.321597in}{0.862634in}}%
\pgfpathlineto{\pgfqpoint{2.321676in}{1.327039in}}%
\pgfpathlineto{\pgfqpoint{2.322410in}{1.327011in}}%
\pgfpathlineto{\pgfqpoint{2.325542in}{1.326796in}}%
\pgfpathlineto{\pgfqpoint{2.328015in}{1.326703in}}%
\pgfpathlineto{\pgfqpoint{2.329722in}{1.326391in}}%
\pgfpathlineto{\pgfqpoint{2.330587in}{0.504364in}}%
\pgfpathlineto{\pgfqpoint{2.336522in}{0.504444in}}%
\pgfpathlineto{\pgfqpoint{2.337818in}{0.504717in}}%
\pgfpathlineto{\pgfqpoint{2.337908in}{0.505415in}}%
\pgfpathlineto{\pgfqpoint{2.338081in}{0.862840in}}%
\pgfpathlineto{\pgfqpoint{2.338161in}{1.327039in}}%
\pgfpathlineto{\pgfqpoint{2.338895in}{1.327011in}}%
\pgfpathlineto{\pgfqpoint{2.342027in}{1.326796in}}%
\pgfpathlineto{\pgfqpoint{2.344500in}{1.326703in}}%
\pgfpathlineto{\pgfqpoint{2.346206in}{1.326392in}}%
\pgfpathlineto{\pgfqpoint{2.347072in}{0.504364in}}%
\pgfpathlineto{\pgfqpoint{2.353007in}{0.504445in}}%
\pgfpathlineto{\pgfqpoint{2.354303in}{0.504717in}}%
\pgfpathlineto{\pgfqpoint{2.354393in}{0.505416in}}%
\pgfpathlineto{\pgfqpoint{2.354566in}{0.863032in}}%
\pgfpathlineto{\pgfqpoint{2.354646in}{1.327039in}}%
\pgfpathlineto{\pgfqpoint{2.355380in}{1.327011in}}%
\pgfpathlineto{\pgfqpoint{2.358512in}{1.326796in}}%
\pgfpathlineto{\pgfqpoint{2.360985in}{1.326703in}}%
\pgfpathlineto{\pgfqpoint{2.362691in}{1.326392in}}%
\pgfpathlineto{\pgfqpoint{2.363557in}{0.504364in}}%
\pgfpathlineto{\pgfqpoint{2.369492in}{0.504445in}}%
\pgfpathlineto{\pgfqpoint{2.370787in}{0.504718in}}%
\pgfpathlineto{\pgfqpoint{2.370878in}{0.505417in}}%
\pgfpathlineto{\pgfqpoint{2.371051in}{0.863188in}}%
\pgfpathlineto{\pgfqpoint{2.371131in}{1.327039in}}%
\pgfpathlineto{\pgfqpoint{2.371865in}{1.327011in}}%
\pgfpathlineto{\pgfqpoint{2.374997in}{1.326796in}}%
\pgfpathlineto{\pgfqpoint{2.377470in}{1.326703in}}%
\pgfpathlineto{\pgfqpoint{2.379176in}{1.326392in}}%
\pgfpathlineto{\pgfqpoint{2.380042in}{0.504364in}}%
\pgfpathlineto{\pgfqpoint{2.385976in}{0.504445in}}%
\pgfpathlineto{\pgfqpoint{2.387272in}{0.504718in}}%
\pgfpathlineto{\pgfqpoint{2.387363in}{0.505417in}}%
\pgfpathlineto{\pgfqpoint{2.387536in}{0.863291in}}%
\pgfpathlineto{\pgfqpoint{2.387616in}{1.327039in}}%
\pgfpathlineto{\pgfqpoint{2.388350in}{1.327011in}}%
\pgfpathlineto{\pgfqpoint{2.391482in}{1.326796in}}%
\pgfpathlineto{\pgfqpoint{2.393954in}{1.326703in}}%
\pgfpathlineto{\pgfqpoint{2.395661in}{1.326392in}}%
\pgfpathlineto{\pgfqpoint{2.396527in}{0.504364in}}%
\pgfpathlineto{\pgfqpoint{2.402461in}{0.504445in}}%
\pgfpathlineto{\pgfqpoint{2.403757in}{0.504718in}}%
\pgfpathlineto{\pgfqpoint{2.403848in}{0.505417in}}%
\pgfpathlineto{\pgfqpoint{2.404021in}{0.863335in}}%
\pgfpathlineto{\pgfqpoint{2.404100in}{1.327039in}}%
\pgfpathlineto{\pgfqpoint{2.404835in}{1.327011in}}%
\pgfpathlineto{\pgfqpoint{2.407967in}{1.326796in}}%
\pgfpathlineto{\pgfqpoint{2.410439in}{1.326703in}}%
\pgfpathlineto{\pgfqpoint{2.412146in}{1.326392in}}%
\pgfpathlineto{\pgfqpoint{2.413011in}{0.504364in}}%
\pgfpathlineto{\pgfqpoint{2.418946in}{0.504445in}}%
\pgfpathlineto{\pgfqpoint{2.420242in}{0.504718in}}%
\pgfpathlineto{\pgfqpoint{2.420333in}{0.505417in}}%
\pgfpathlineto{\pgfqpoint{2.420506in}{0.863323in}}%
\pgfpathlineto{\pgfqpoint{2.420585in}{1.327039in}}%
\pgfpathlineto{\pgfqpoint{2.421319in}{1.327011in}}%
\pgfpathlineto{\pgfqpoint{2.424451in}{1.326796in}}%
\pgfpathlineto{\pgfqpoint{2.426924in}{1.326703in}}%
\pgfpathlineto{\pgfqpoint{2.428631in}{1.326392in}}%
\pgfpathlineto{\pgfqpoint{2.429496in}{0.504364in}}%
\pgfpathlineto{\pgfqpoint{2.435431in}{0.504445in}}%
\pgfpathlineto{\pgfqpoint{2.436727in}{0.504718in}}%
\pgfpathlineto{\pgfqpoint{2.436818in}{0.505417in}}%
\pgfpathlineto{\pgfqpoint{2.436991in}{0.863265in}}%
\pgfpathlineto{\pgfqpoint{2.437070in}{1.327039in}}%
\pgfpathlineto{\pgfqpoint{2.437804in}{1.327011in}}%
\pgfpathlineto{\pgfqpoint{2.440936in}{1.326796in}}%
\pgfpathlineto{\pgfqpoint{2.443409in}{1.326703in}}%
\pgfpathlineto{\pgfqpoint{2.445115in}{1.326392in}}%
\pgfpathlineto{\pgfqpoint{2.445981in}{0.504364in}}%
\pgfpathlineto{\pgfqpoint{2.451916in}{0.504445in}}%
\pgfpathlineto{\pgfqpoint{2.453212in}{0.504717in}}%
\pgfpathlineto{\pgfqpoint{2.453303in}{0.505417in}}%
\pgfpathlineto{\pgfqpoint{2.453475in}{0.863176in}}%
\pgfpathlineto{\pgfqpoint{2.453555in}{1.327039in}}%
\pgfpathlineto{\pgfqpoint{2.454289in}{1.327011in}}%
\pgfpathlineto{\pgfqpoint{2.457421in}{1.326796in}}%
\pgfpathlineto{\pgfqpoint{2.459894in}{1.326703in}}%
\pgfpathlineto{\pgfqpoint{2.461600in}{1.326392in}}%
\pgfpathlineto{\pgfqpoint{2.462466in}{0.504364in}}%
\pgfpathlineto{\pgfqpoint{2.468401in}{0.504445in}}%
\pgfpathlineto{\pgfqpoint{2.469696in}{0.504717in}}%
\pgfpathlineto{\pgfqpoint{2.469787in}{0.505416in}}%
\pgfpathlineto{\pgfqpoint{2.469960in}{0.863072in}}%
\pgfpathlineto{\pgfqpoint{2.470040in}{1.327039in}}%
\pgfpathlineto{\pgfqpoint{2.470774in}{1.327011in}}%
\pgfpathlineto{\pgfqpoint{2.473906in}{1.326796in}}%
\pgfpathlineto{\pgfqpoint{2.476379in}{1.326703in}}%
\pgfpathlineto{\pgfqpoint{2.478085in}{1.326392in}}%
\pgfpathlineto{\pgfqpoint{2.478951in}{0.504364in}}%
\pgfpathlineto{\pgfqpoint{2.484885in}{0.504444in}}%
\pgfpathlineto{\pgfqpoint{2.486181in}{0.504717in}}%
\pgfpathlineto{\pgfqpoint{2.486272in}{0.505416in}}%
\pgfpathlineto{\pgfqpoint{2.486445in}{0.862971in}}%
\pgfpathlineto{\pgfqpoint{2.486525in}{1.327039in}}%
\pgfpathlineto{\pgfqpoint{2.487259in}{1.327011in}}%
\pgfpathlineto{\pgfqpoint{2.490391in}{1.326796in}}%
\pgfpathlineto{\pgfqpoint{2.492863in}{1.326703in}}%
\pgfpathlineto{\pgfqpoint{2.494546in}{1.326559in}}%
\pgfpathlineto{\pgfqpoint{2.495284in}{0.504423in}}%
\pgfpathlineto{\pgfqpoint{2.497592in}{0.504335in}}%
\pgfpathlineto{\pgfqpoint{2.502663in}{0.504708in}}%
\pgfpathlineto{\pgfqpoint{2.502751in}{0.505332in}}%
\pgfpathlineto{\pgfqpoint{2.502841in}{0.531039in}}%
\pgfpathlineto{\pgfqpoint{2.503001in}{1.327040in}}%
\pgfpathlineto{\pgfqpoint{2.503727in}{1.327013in}}%
\pgfpathlineto{\pgfqpoint{2.506859in}{1.326797in}}%
\pgfpathlineto{\pgfqpoint{2.509331in}{1.326705in}}%
\pgfpathlineto{\pgfqpoint{2.511055in}{1.326392in}}%
\pgfpathlineto{\pgfqpoint{2.511921in}{0.504364in}}%
\pgfpathlineto{\pgfqpoint{2.517855in}{0.504445in}}%
\pgfpathlineto{\pgfqpoint{2.519151in}{0.504718in}}%
\pgfpathlineto{\pgfqpoint{2.519242in}{0.505419in}}%
\pgfpathlineto{\pgfqpoint{2.519415in}{0.863721in}}%
\pgfpathlineto{\pgfqpoint{2.519495in}{1.327039in}}%
\pgfpathlineto{\pgfqpoint{2.520229in}{1.327011in}}%
\pgfpathlineto{\pgfqpoint{2.523361in}{1.326796in}}%
\pgfpathlineto{\pgfqpoint{2.525834in}{1.326703in}}%
\pgfpathlineto{\pgfqpoint{2.527540in}{1.326392in}}%
\pgfpathlineto{\pgfqpoint{2.528405in}{0.504364in}}%
\pgfpathlineto{\pgfqpoint{2.534340in}{0.504445in}}%
\pgfpathlineto{\pgfqpoint{2.535636in}{0.504718in}}%
\pgfpathlineto{\pgfqpoint{2.535727in}{0.505418in}}%
\pgfpathlineto{\pgfqpoint{2.535900in}{0.863470in}}%
\pgfpathlineto{\pgfqpoint{2.535979in}{1.327039in}}%
\pgfpathlineto{\pgfqpoint{2.536713in}{1.327011in}}%
\pgfpathlineto{\pgfqpoint{2.539846in}{1.326796in}}%
\pgfpathlineto{\pgfqpoint{2.542318in}{1.326703in}}%
\pgfpathlineto{\pgfqpoint{2.544025in}{1.326392in}}%
\pgfpathlineto{\pgfqpoint{2.544890in}{0.504364in}}%
\pgfpathlineto{\pgfqpoint{2.550825in}{0.504445in}}%
\pgfpathlineto{\pgfqpoint{2.552121in}{0.504717in}}%
\pgfpathlineto{\pgfqpoint{2.552212in}{0.505417in}}%
\pgfpathlineto{\pgfqpoint{2.552385in}{0.863189in}}%
\pgfpathlineto{\pgfqpoint{2.552464in}{1.327039in}}%
\pgfpathlineto{\pgfqpoint{2.553198in}{1.327011in}}%
\pgfpathlineto{\pgfqpoint{2.556330in}{1.326796in}}%
\pgfpathlineto{\pgfqpoint{2.558803in}{1.326703in}}%
\pgfpathlineto{\pgfqpoint{2.560509in}{1.326392in}}%
\pgfpathlineto{\pgfqpoint{2.561375in}{0.504364in}}%
\pgfpathlineto{\pgfqpoint{2.567310in}{0.504444in}}%
\pgfpathlineto{\pgfqpoint{2.568606in}{0.504717in}}%
\pgfpathlineto{\pgfqpoint{2.568696in}{0.505415in}}%
\pgfpathlineto{\pgfqpoint{2.568869in}{0.862922in}}%
\pgfpathlineto{\pgfqpoint{2.568949in}{1.327039in}}%
\pgfpathlineto{\pgfqpoint{2.569683in}{1.327011in}}%
\pgfpathlineto{\pgfqpoint{2.572815in}{1.326796in}}%
\pgfpathlineto{\pgfqpoint{2.575288in}{1.326703in}}%
\pgfpathlineto{\pgfqpoint{2.576994in}{1.326391in}}%
\pgfpathlineto{\pgfqpoint{2.577860in}{0.504364in}}%
\pgfpathlineto{\pgfqpoint{2.583959in}{0.504458in}}%
\pgfpathlineto{\pgfqpoint{2.585090in}{0.504717in}}%
\pgfpathlineto{\pgfqpoint{2.585181in}{0.505414in}}%
\pgfpathlineto{\pgfqpoint{2.585354in}{0.862703in}}%
\pgfpathlineto{\pgfqpoint{2.585434in}{1.327039in}}%
\pgfpathlineto{\pgfqpoint{2.586167in}{1.327011in}}%
\pgfpathlineto{\pgfqpoint{2.589300in}{1.326796in}}%
\pgfpathlineto{\pgfqpoint{2.591772in}{1.326703in}}%
\pgfpathlineto{\pgfqpoint{2.593479in}{1.326391in}}%
\pgfpathlineto{\pgfqpoint{2.594345in}{0.504364in}}%
\pgfpathlineto{\pgfqpoint{2.600444in}{0.504458in}}%
\pgfpathlineto{\pgfqpoint{2.601575in}{0.504717in}}%
\pgfpathlineto{\pgfqpoint{2.601666in}{0.505414in}}%
\pgfpathlineto{\pgfqpoint{2.601839in}{0.862554in}}%
\pgfpathlineto{\pgfqpoint{2.601918in}{1.327039in}}%
\pgfpathlineto{\pgfqpoint{2.602652in}{1.327011in}}%
\pgfpathlineto{\pgfqpoint{2.605784in}{1.326796in}}%
\pgfpathlineto{\pgfqpoint{2.608257in}{1.326703in}}%
\pgfpathlineto{\pgfqpoint{2.609964in}{1.326391in}}%
\pgfpathlineto{\pgfqpoint{2.610830in}{0.504364in}}%
\pgfpathlineto{\pgfqpoint{2.616929in}{0.504458in}}%
\pgfpathlineto{\pgfqpoint{2.618060in}{0.504717in}}%
\pgfpathlineto{\pgfqpoint{2.618151in}{0.505413in}}%
\pgfpathlineto{\pgfqpoint{2.618324in}{0.862485in}}%
\pgfpathlineto{\pgfqpoint{2.618403in}{1.327039in}}%
\pgfpathlineto{\pgfqpoint{2.619137in}{1.327011in}}%
\pgfpathlineto{\pgfqpoint{2.622269in}{1.326796in}}%
\pgfpathlineto{\pgfqpoint{2.624742in}{1.326703in}}%
\pgfpathlineto{\pgfqpoint{2.626449in}{1.326391in}}%
\pgfpathlineto{\pgfqpoint{2.627315in}{0.504364in}}%
\pgfpathlineto{\pgfqpoint{2.633414in}{0.504458in}}%
\pgfpathlineto{\pgfqpoint{2.634545in}{0.504717in}}%
\pgfpathlineto{\pgfqpoint{2.634636in}{0.505413in}}%
\pgfpathlineto{\pgfqpoint{2.634809in}{0.862493in}}%
\pgfpathlineto{\pgfqpoint{2.634888in}{1.327039in}}%
\pgfpathlineto{\pgfqpoint{2.635622in}{1.327011in}}%
\pgfpathlineto{\pgfqpoint{2.638754in}{1.326796in}}%
\pgfpathlineto{\pgfqpoint{2.641227in}{1.326703in}}%
\pgfpathlineto{\pgfqpoint{2.642934in}{1.326391in}}%
\pgfpathlineto{\pgfqpoint{2.643799in}{0.504364in}}%
\pgfpathlineto{\pgfqpoint{2.649899in}{0.504458in}}%
\pgfpathlineto{\pgfqpoint{2.651030in}{0.504717in}}%
\pgfpathlineto{\pgfqpoint{2.651121in}{0.505414in}}%
\pgfpathlineto{\pgfqpoint{2.651293in}{0.862565in}}%
\pgfpathlineto{\pgfqpoint{2.651373in}{1.327039in}}%
\pgfpathlineto{\pgfqpoint{2.652107in}{1.327011in}}%
\pgfpathlineto{\pgfqpoint{2.655239in}{1.326796in}}%
\pgfpathlineto{\pgfqpoint{2.657712in}{1.326703in}}%
\pgfpathlineto{\pgfqpoint{2.659418in}{1.326391in}}%
\pgfpathlineto{\pgfqpoint{2.660284in}{0.504364in}}%
\pgfpathlineto{\pgfqpoint{2.666384in}{0.504458in}}%
\pgfpathlineto{\pgfqpoint{2.667515in}{0.504717in}}%
\pgfpathlineto{\pgfqpoint{2.667605in}{0.505414in}}%
\pgfpathlineto{\pgfqpoint{2.667778in}{0.862682in}}%
\pgfpathlineto{\pgfqpoint{2.667858in}{1.327039in}}%
\pgfpathlineto{\pgfqpoint{2.668592in}{1.327011in}}%
\pgfpathlineto{\pgfqpoint{2.671724in}{1.326796in}}%
\pgfpathlineto{\pgfqpoint{2.674197in}{1.326703in}}%
\pgfpathlineto{\pgfqpoint{2.675903in}{1.326391in}}%
\pgfpathlineto{\pgfqpoint{2.676769in}{0.504364in}}%
\pgfpathlineto{\pgfqpoint{2.682704in}{0.504444in}}%
\pgfpathlineto{\pgfqpoint{2.683999in}{0.504717in}}%
\pgfpathlineto{\pgfqpoint{2.684090in}{0.505415in}}%
\pgfpathlineto{\pgfqpoint{2.684263in}{0.862821in}}%
\pgfpathlineto{\pgfqpoint{2.684343in}{1.327039in}}%
\pgfpathlineto{\pgfqpoint{2.685077in}{1.327011in}}%
\pgfpathlineto{\pgfqpoint{2.688209in}{1.326796in}}%
\pgfpathlineto{\pgfqpoint{2.690681in}{1.326703in}}%
\pgfpathlineto{\pgfqpoint{2.692388in}{1.326392in}}%
\pgfpathlineto{\pgfqpoint{2.693254in}{0.504364in}}%
\pgfpathlineto{\pgfqpoint{2.699189in}{0.504445in}}%
\pgfpathlineto{\pgfqpoint{2.700484in}{0.504717in}}%
\pgfpathlineto{\pgfqpoint{2.700575in}{0.505416in}}%
\pgfpathlineto{\pgfqpoint{2.700748in}{0.862960in}}%
\pgfpathlineto{\pgfqpoint{2.700828in}{1.327039in}}%
\pgfpathlineto{\pgfqpoint{2.701562in}{1.327011in}}%
\pgfpathlineto{\pgfqpoint{2.704694in}{1.326796in}}%
\pgfpathlineto{\pgfqpoint{2.707166in}{1.326703in}}%
\pgfpathlineto{\pgfqpoint{2.708873in}{1.326392in}}%
\pgfpathlineto{\pgfqpoint{2.709739in}{0.504364in}}%
\pgfpathlineto{\pgfqpoint{2.715673in}{0.504445in}}%
\pgfpathlineto{\pgfqpoint{2.716969in}{0.504717in}}%
\pgfpathlineto{\pgfqpoint{2.717060in}{0.505416in}}%
\pgfpathlineto{\pgfqpoint{2.717233in}{0.863079in}}%
\pgfpathlineto{\pgfqpoint{2.717312in}{1.327039in}}%
\pgfpathlineto{\pgfqpoint{2.718046in}{1.327011in}}%
\pgfpathlineto{\pgfqpoint{2.721179in}{1.326796in}}%
\pgfpathlineto{\pgfqpoint{2.723651in}{1.326703in}}%
\pgfpathlineto{\pgfqpoint{2.725358in}{1.326392in}}%
\pgfpathlineto{\pgfqpoint{2.726224in}{0.504364in}}%
\pgfpathlineto{\pgfqpoint{2.732158in}{0.504445in}}%
\pgfpathlineto{\pgfqpoint{2.733454in}{0.504717in}}%
\pgfpathlineto{\pgfqpoint{2.733545in}{0.505417in}}%
\pgfpathlineto{\pgfqpoint{2.733718in}{0.863167in}}%
\pgfpathlineto{\pgfqpoint{2.733797in}{1.327039in}}%
\pgfpathlineto{\pgfqpoint{2.734531in}{1.327011in}}%
\pgfpathlineto{\pgfqpoint{2.737664in}{1.326796in}}%
\pgfpathlineto{\pgfqpoint{2.740136in}{1.326703in}}%
\pgfpathlineto{\pgfqpoint{2.741843in}{1.326392in}}%
\pgfpathlineto{\pgfqpoint{2.742708in}{0.504364in}}%
\pgfpathlineto{\pgfqpoint{2.748643in}{0.504445in}}%
\pgfpathlineto{\pgfqpoint{2.749939in}{0.504718in}}%
\pgfpathlineto{\pgfqpoint{2.750030in}{0.505417in}}%
\pgfpathlineto{\pgfqpoint{2.750203in}{0.863214in}}%
\pgfpathlineto{\pgfqpoint{2.750282in}{1.327039in}}%
\pgfpathlineto{\pgfqpoint{2.751016in}{1.327011in}}%
\pgfpathlineto{\pgfqpoint{2.754148in}{1.326796in}}%
\pgfpathlineto{\pgfqpoint{2.756621in}{1.326703in}}%
\pgfpathlineto{\pgfqpoint{2.758303in}{1.326559in}}%
\pgfpathlineto{\pgfqpoint{2.759042in}{0.504423in}}%
\pgfpathlineto{\pgfqpoint{2.761350in}{0.504335in}}%
\pgfpathlineto{\pgfqpoint{2.766421in}{0.504708in}}%
\pgfpathlineto{\pgfqpoint{2.766509in}{0.505332in}}%
\pgfpathlineto{\pgfqpoint{2.766599in}{0.531080in}}%
\pgfpathlineto{\pgfqpoint{2.766758in}{1.327040in}}%
\pgfpathlineto{\pgfqpoint{2.767484in}{1.327013in}}%
\pgfpathlineto{\pgfqpoint{2.770616in}{1.326797in}}%
\pgfpathlineto{\pgfqpoint{2.773089in}{1.326706in}}%
\pgfpathlineto{\pgfqpoint{2.774812in}{1.326392in}}%
\pgfpathlineto{\pgfqpoint{2.775678in}{0.504365in}}%
\pgfpathlineto{\pgfqpoint{2.781613in}{0.504445in}}%
\pgfpathlineto{\pgfqpoint{2.782909in}{0.504718in}}%
\pgfpathlineto{\pgfqpoint{2.783000in}{0.505421in}}%
\pgfpathlineto{\pgfqpoint{2.783173in}{0.864083in}}%
\pgfpathlineto{\pgfqpoint{2.783252in}{1.327039in}}%
\pgfpathlineto{\pgfqpoint{2.783987in}{1.327011in}}%
\pgfpathlineto{\pgfqpoint{2.787119in}{1.326796in}}%
\pgfpathlineto{\pgfqpoint{2.789591in}{1.326703in}}%
\pgfpathlineto{\pgfqpoint{2.791297in}{1.326392in}}%
\pgfpathlineto{\pgfqpoint{2.792163in}{0.504364in}}%
\pgfpathlineto{\pgfqpoint{2.798098in}{0.504445in}}%
\pgfpathlineto{\pgfqpoint{2.799394in}{0.504718in}}%
\pgfpathlineto{\pgfqpoint{2.799485in}{0.505420in}}%
\pgfpathlineto{\pgfqpoint{2.799658in}{0.863810in}}%
\pgfpathlineto{\pgfqpoint{2.799737in}{1.327039in}}%
\pgfpathlineto{\pgfqpoint{2.800471in}{1.327011in}}%
\pgfpathlineto{\pgfqpoint{2.803603in}{1.326796in}}%
\pgfpathlineto{\pgfqpoint{2.806076in}{1.326703in}}%
\pgfpathlineto{\pgfqpoint{2.807782in}{1.326392in}}%
\pgfpathlineto{\pgfqpoint{2.808648in}{0.504364in}}%
\pgfpathlineto{\pgfqpoint{2.814582in}{0.504445in}}%
\pgfpathlineto{\pgfqpoint{2.815878in}{0.504718in}}%
\pgfpathlineto{\pgfqpoint{2.815969in}{0.505418in}}%
\pgfpathlineto{\pgfqpoint{2.816142in}{0.863464in}}%
\pgfpathlineto{\pgfqpoint{2.816222in}{1.327039in}}%
\pgfpathlineto{\pgfqpoint{2.816956in}{1.327011in}}%
\pgfpathlineto{\pgfqpoint{2.820088in}{1.326796in}}%
\pgfpathlineto{\pgfqpoint{2.822561in}{1.326703in}}%
\pgfpathlineto{\pgfqpoint{2.824267in}{1.326392in}}%
\pgfpathlineto{\pgfqpoint{2.825133in}{0.504364in}}%
\pgfpathlineto{\pgfqpoint{2.831067in}{0.504445in}}%
\pgfpathlineto{\pgfqpoint{2.832363in}{0.504717in}}%
\pgfpathlineto{\pgfqpoint{2.832454in}{0.505416in}}%
\pgfpathlineto{\pgfqpoint{2.832627in}{0.863104in}}%
\pgfpathlineto{\pgfqpoint{2.832706in}{1.327039in}}%
\pgfpathlineto{\pgfqpoint{2.833440in}{1.327011in}}%
\pgfpathlineto{\pgfqpoint{2.836573in}{1.326796in}}%
\pgfpathlineto{\pgfqpoint{2.839045in}{1.326703in}}%
\pgfpathlineto{\pgfqpoint{2.840752in}{1.326391in}}%
\pgfpathlineto{\pgfqpoint{2.841618in}{0.504364in}}%
\pgfpathlineto{\pgfqpoint{2.847552in}{0.504444in}}%
\pgfpathlineto{\pgfqpoint{2.848848in}{0.504717in}}%
\pgfpathlineto{\pgfqpoint{2.848939in}{0.505415in}}%
\pgfpathlineto{\pgfqpoint{2.849112in}{0.862782in}}%
\pgfpathlineto{\pgfqpoint{2.849191in}{1.327039in}}%
\pgfpathlineto{\pgfqpoint{2.849925in}{1.327011in}}%
\pgfpathlineto{\pgfqpoint{2.853057in}{1.326796in}}%
\pgfpathlineto{\pgfqpoint{2.855530in}{1.326703in}}%
\pgfpathlineto{\pgfqpoint{2.857237in}{1.326391in}}%
\pgfpathlineto{\pgfqpoint{2.858102in}{0.504364in}}%
\pgfpathlineto{\pgfqpoint{2.864202in}{0.504458in}}%
\pgfpathlineto{\pgfqpoint{2.865333in}{0.504717in}}%
\pgfpathlineto{\pgfqpoint{2.865424in}{0.505413in}}%
\pgfpathlineto{\pgfqpoint{2.865596in}{0.862536in}}%
\pgfpathlineto{\pgfqpoint{2.865676in}{1.327039in}}%
\pgfpathlineto{\pgfqpoint{2.866410in}{1.327010in}}%
\pgfpathlineto{\pgfqpoint{2.869542in}{1.326796in}}%
\pgfpathlineto{\pgfqpoint{2.872015in}{1.326703in}}%
\pgfpathlineto{\pgfqpoint{2.873722in}{1.326391in}}%
\pgfpathlineto{\pgfqpoint{2.874587in}{0.504364in}}%
\pgfpathlineto{\pgfqpoint{2.880687in}{0.504458in}}%
\pgfpathlineto{\pgfqpoint{2.881818in}{0.504717in}}%
\pgfpathlineto{\pgfqpoint{2.881908in}{0.505413in}}%
\pgfpathlineto{\pgfqpoint{2.882081in}{0.862387in}}%
\pgfpathlineto{\pgfqpoint{2.882161in}{1.327039in}}%
\pgfpathlineto{\pgfqpoint{2.882894in}{1.327010in}}%
\pgfpathlineto{\pgfqpoint{2.886027in}{1.326796in}}%
\pgfpathlineto{\pgfqpoint{2.888499in}{1.326703in}}%
\pgfpathlineto{\pgfqpoint{2.890206in}{1.326391in}}%
\pgfpathlineto{\pgfqpoint{2.891072in}{0.504364in}}%
\pgfpathlineto{\pgfqpoint{2.897172in}{0.504458in}}%
\pgfpathlineto{\pgfqpoint{2.898302in}{0.504717in}}%
\pgfpathlineto{\pgfqpoint{2.898393in}{0.505412in}}%
\pgfpathlineto{\pgfqpoint{2.898566in}{0.862340in}}%
\pgfpathlineto{\pgfqpoint{2.898646in}{1.327039in}}%
\pgfpathlineto{\pgfqpoint{2.899379in}{1.327010in}}%
\pgfpathlineto{\pgfqpoint{2.902511in}{1.326796in}}%
\pgfpathlineto{\pgfqpoint{2.904984in}{1.326703in}}%
\pgfpathlineto{\pgfqpoint{2.906691in}{1.326391in}}%
\pgfpathlineto{\pgfqpoint{2.907557in}{0.504364in}}%
\pgfpathlineto{\pgfqpoint{2.913656in}{0.504458in}}%
\pgfpathlineto{\pgfqpoint{2.914787in}{0.504717in}}%
\pgfpathlineto{\pgfqpoint{2.914878in}{0.505413in}}%
\pgfpathlineto{\pgfqpoint{2.915051in}{0.862387in}}%
\pgfpathlineto{\pgfqpoint{2.915130in}{1.327039in}}%
\pgfpathlineto{\pgfqpoint{2.915864in}{1.327011in}}%
\pgfpathlineto{\pgfqpoint{2.918996in}{1.326796in}}%
\pgfpathlineto{\pgfqpoint{2.921469in}{1.326703in}}%
\pgfpathlineto{\pgfqpoint{2.923176in}{1.326391in}}%
\pgfpathlineto{\pgfqpoint{2.924042in}{0.504364in}}%
\pgfpathlineto{\pgfqpoint{2.930141in}{0.504458in}}%
\pgfpathlineto{\pgfqpoint{2.931272in}{0.504717in}}%
\pgfpathlineto{\pgfqpoint{2.931363in}{0.505413in}}%
\pgfpathlineto{\pgfqpoint{2.931536in}{0.862506in}}%
\pgfpathlineto{\pgfqpoint{2.931615in}{1.327039in}}%
\pgfpathlineto{\pgfqpoint{2.932349in}{1.327011in}}%
\pgfpathlineto{\pgfqpoint{2.935481in}{1.326796in}}%
\pgfpathlineto{\pgfqpoint{2.937954in}{1.326703in}}%
\pgfpathlineto{\pgfqpoint{2.939661in}{1.326391in}}%
\pgfpathlineto{\pgfqpoint{2.940527in}{0.504364in}}%
\pgfpathlineto{\pgfqpoint{2.946626in}{0.504458in}}%
\pgfpathlineto{\pgfqpoint{2.947757in}{0.504717in}}%
\pgfpathlineto{\pgfqpoint{2.947848in}{0.505414in}}%
\pgfpathlineto{\pgfqpoint{2.948021in}{0.862671in}}%
\pgfpathlineto{\pgfqpoint{2.948100in}{1.327039in}}%
\pgfpathlineto{\pgfqpoint{2.948834in}{1.327011in}}%
\pgfpathlineto{\pgfqpoint{2.951966in}{1.326796in}}%
\pgfpathlineto{\pgfqpoint{2.954439in}{1.326703in}}%
\pgfpathlineto{\pgfqpoint{2.956146in}{1.326392in}}%
\pgfpathlineto{\pgfqpoint{2.957012in}{0.504364in}}%
\pgfpathlineto{\pgfqpoint{2.962946in}{0.504444in}}%
\pgfpathlineto{\pgfqpoint{2.964242in}{0.504717in}}%
\pgfpathlineto{\pgfqpoint{2.964333in}{0.505415in}}%
\pgfpathlineto{\pgfqpoint{2.964506in}{0.862853in}}%
\pgfpathlineto{\pgfqpoint{2.964585in}{1.327039in}}%
\pgfpathlineto{\pgfqpoint{2.965319in}{1.327011in}}%
\pgfpathlineto{\pgfqpoint{2.968451in}{1.326796in}}%
\pgfpathlineto{\pgfqpoint{2.970924in}{1.326703in}}%
\pgfpathlineto{\pgfqpoint{2.972631in}{1.326392in}}%
\pgfpathlineto{\pgfqpoint{2.973496in}{0.504364in}}%
\pgfpathlineto{\pgfqpoint{2.979431in}{0.504445in}}%
\pgfpathlineto{\pgfqpoint{2.980727in}{0.504717in}}%
\pgfpathlineto{\pgfqpoint{2.980818in}{0.505416in}}%
\pgfpathlineto{\pgfqpoint{2.980991in}{0.863023in}}%
\pgfpathlineto{\pgfqpoint{2.981070in}{1.327039in}}%
\pgfpathlineto{\pgfqpoint{2.981804in}{1.327011in}}%
\pgfpathlineto{\pgfqpoint{2.984936in}{1.326796in}}%
\pgfpathlineto{\pgfqpoint{2.987409in}{1.326703in}}%
\pgfpathlineto{\pgfqpoint{2.989115in}{1.326392in}}%
\pgfpathlineto{\pgfqpoint{2.989981in}{0.504364in}}%
\pgfpathlineto{\pgfqpoint{2.995916in}{0.504445in}}%
\pgfpathlineto{\pgfqpoint{2.997212in}{0.504718in}}%
\pgfpathlineto{\pgfqpoint{2.997303in}{0.505417in}}%
\pgfpathlineto{\pgfqpoint{2.997475in}{0.863162in}}%
\pgfpathlineto{\pgfqpoint{2.997555in}{1.327039in}}%
\pgfpathlineto{\pgfqpoint{2.998289in}{1.327011in}}%
\pgfpathlineto{\pgfqpoint{3.001421in}{1.326796in}}%
\pgfpathlineto{\pgfqpoint{3.003894in}{1.326703in}}%
\pgfpathlineto{\pgfqpoint{3.005600in}{1.326392in}}%
\pgfpathlineto{\pgfqpoint{3.006466in}{0.504364in}}%
\pgfpathlineto{\pgfqpoint{3.012401in}{0.504445in}}%
\pgfpathlineto{\pgfqpoint{3.013697in}{0.504718in}}%
\pgfpathlineto{\pgfqpoint{3.013787in}{0.505417in}}%
\pgfpathlineto{\pgfqpoint{3.013960in}{0.863253in}}%
\pgfpathlineto{\pgfqpoint{3.014040in}{1.327039in}}%
\pgfpathlineto{\pgfqpoint{3.014774in}{1.327011in}}%
\pgfpathlineto{\pgfqpoint{3.017906in}{1.326796in}}%
\pgfpathlineto{\pgfqpoint{3.020379in}{1.326703in}}%
\pgfpathlineto{\pgfqpoint{3.022085in}{1.326392in}}%
\pgfpathlineto{\pgfqpoint{3.022951in}{0.504364in}}%
\pgfpathlineto{\pgfqpoint{3.028885in}{0.504445in}}%
\pgfpathlineto{\pgfqpoint{3.030181in}{0.504718in}}%
\pgfpathlineto{\pgfqpoint{3.030272in}{0.505417in}}%
\pgfpathlineto{\pgfqpoint{3.030445in}{0.863293in}}%
\pgfpathlineto{\pgfqpoint{3.030525in}{1.327039in}}%
\pgfpathlineto{\pgfqpoint{3.031259in}{1.327011in}}%
\pgfpathlineto{\pgfqpoint{3.034391in}{1.326796in}}%
\pgfpathlineto{\pgfqpoint{3.036864in}{1.326703in}}%
\pgfpathlineto{\pgfqpoint{3.038570in}{1.326392in}}%
\pgfpathlineto{\pgfqpoint{3.039436in}{0.504364in}}%
\pgfpathlineto{\pgfqpoint{3.045370in}{0.504445in}}%
\pgfpathlineto{\pgfqpoint{3.046666in}{0.504718in}}%
\pgfpathlineto{\pgfqpoint{3.046757in}{0.505417in}}%
\pgfpathlineto{\pgfqpoint{3.046930in}{0.863283in}}%
\pgfpathlineto{\pgfqpoint{3.047010in}{1.327039in}}%
\pgfpathlineto{\pgfqpoint{3.047744in}{1.327011in}}%
\pgfpathlineto{\pgfqpoint{3.050876in}{1.326796in}}%
\pgfpathlineto{\pgfqpoint{3.053348in}{1.326703in}}%
\pgfpathlineto{\pgfqpoint{3.055055in}{1.326392in}}%
\pgfpathlineto{\pgfqpoint{3.055921in}{0.504364in}}%
\pgfpathlineto{\pgfqpoint{3.061855in}{0.504445in}}%
\pgfpathlineto{\pgfqpoint{3.063151in}{0.504718in}}%
\pgfpathlineto{\pgfqpoint{3.063242in}{0.505417in}}%
\pgfpathlineto{\pgfqpoint{3.063415in}{0.863232in}}%
\pgfpathlineto{\pgfqpoint{3.063494in}{1.327039in}}%
\pgfpathlineto{\pgfqpoint{3.064228in}{1.327011in}}%
\pgfpathlineto{\pgfqpoint{3.067361in}{1.326796in}}%
\pgfpathlineto{\pgfqpoint{3.069833in}{1.326703in}}%
\pgfpathlineto{\pgfqpoint{3.071540in}{1.326392in}}%
\pgfpathlineto{\pgfqpoint{3.072405in}{0.504364in}}%
\pgfpathlineto{\pgfqpoint{3.078340in}{0.504445in}}%
\pgfpathlineto{\pgfqpoint{3.079636in}{0.504717in}}%
\pgfpathlineto{\pgfqpoint{3.079727in}{0.505417in}}%
\pgfpathlineto{\pgfqpoint{3.079900in}{0.863153in}}%
\pgfpathlineto{\pgfqpoint{3.079979in}{1.327039in}}%
\pgfpathlineto{\pgfqpoint{3.080713in}{1.327011in}}%
\pgfpathlineto{\pgfqpoint{3.083845in}{1.326796in}}%
\pgfpathlineto{\pgfqpoint{3.086318in}{1.326703in}}%
\pgfpathlineto{\pgfqpoint{3.088025in}{1.326392in}}%
\pgfpathlineto{\pgfqpoint{3.088890in}{0.504364in}}%
\pgfpathlineto{\pgfqpoint{3.094825in}{0.504445in}}%
\pgfpathlineto{\pgfqpoint{3.096121in}{0.504717in}}%
\pgfpathlineto{\pgfqpoint{3.096212in}{0.505416in}}%
\pgfpathlineto{\pgfqpoint{3.096385in}{0.863062in}}%
\pgfpathlineto{\pgfqpoint{3.096464in}{1.327039in}}%
\pgfpathlineto{\pgfqpoint{3.097198in}{1.327011in}}%
\pgfpathlineto{\pgfqpoint{3.100330in}{1.326796in}}%
\pgfpathlineto{\pgfqpoint{3.102803in}{1.326703in}}%
\pgfpathlineto{\pgfqpoint{3.104509in}{1.326392in}}%
\pgfpathlineto{\pgfqpoint{3.105375in}{0.504364in}}%
\pgfpathlineto{\pgfqpoint{3.111310in}{0.504444in}}%
\pgfpathlineto{\pgfqpoint{3.112606in}{0.504717in}}%
\pgfpathlineto{\pgfqpoint{3.112696in}{0.505416in}}%
\pgfpathlineto{\pgfqpoint{3.112869in}{0.862973in}}%
\pgfpathlineto{\pgfqpoint{3.112949in}{1.327039in}}%
\pgfpathlineto{\pgfqpoint{3.113683in}{1.327011in}}%
\pgfpathlineto{\pgfqpoint{3.116815in}{1.326796in}}%
\pgfpathlineto{\pgfqpoint{3.119288in}{1.326703in}}%
\pgfpathlineto{\pgfqpoint{3.120994in}{1.326392in}}%
\pgfpathlineto{\pgfqpoint{3.121860in}{0.504364in}}%
\pgfpathlineto{\pgfqpoint{3.127795in}{0.504444in}}%
\pgfpathlineto{\pgfqpoint{3.129090in}{0.504717in}}%
\pgfpathlineto{\pgfqpoint{3.129181in}{0.505415in}}%
\pgfpathlineto{\pgfqpoint{3.129354in}{0.862897in}}%
\pgfpathlineto{\pgfqpoint{3.129434in}{1.327039in}}%
\pgfpathlineto{\pgfqpoint{3.130168in}{1.327011in}}%
\pgfpathlineto{\pgfqpoint{3.133300in}{1.326796in}}%
\pgfpathlineto{\pgfqpoint{3.135772in}{1.326703in}}%
\pgfpathlineto{\pgfqpoint{3.137479in}{1.326392in}}%
\pgfpathlineto{\pgfqpoint{3.138345in}{0.504364in}}%
\pgfpathlineto{\pgfqpoint{3.144279in}{0.504444in}}%
\pgfpathlineto{\pgfqpoint{3.145575in}{0.504717in}}%
\pgfpathlineto{\pgfqpoint{3.145666in}{0.505415in}}%
\pgfpathlineto{\pgfqpoint{3.145839in}{0.862843in}}%
\pgfpathlineto{\pgfqpoint{3.145918in}{1.327039in}}%
\pgfpathlineto{\pgfqpoint{3.146652in}{1.327011in}}%
\pgfpathlineto{\pgfqpoint{3.149784in}{1.326796in}}%
\pgfpathlineto{\pgfqpoint{3.152257in}{1.326703in}}%
\pgfpathlineto{\pgfqpoint{3.153964in}{1.326391in}}%
\pgfpathlineto{\pgfqpoint{3.154830in}{0.504364in}}%
\pgfpathlineto{\pgfqpoint{3.160764in}{0.504444in}}%
\pgfpathlineto{\pgfqpoint{3.162060in}{0.504717in}}%
\pgfpathlineto{\pgfqpoint{3.162151in}{0.505415in}}%
\pgfpathlineto{\pgfqpoint{3.162324in}{0.862815in}}%
\pgfpathlineto{\pgfqpoint{3.162403in}{1.327039in}}%
\pgfpathlineto{\pgfqpoint{3.163137in}{1.327011in}}%
\pgfpathlineto{\pgfqpoint{3.166269in}{1.326796in}}%
\pgfpathlineto{\pgfqpoint{3.168742in}{1.326703in}}%
\pgfpathlineto{\pgfqpoint{3.170132in}{1.326543in}}%
\pgfpathlineto{\pgfqpoint{3.170132in}{1.326543in}}%
\pgfusepath{stroke}%
\end{pgfscope}%
\begin{pgfscope}%
\pgfsetrectcap%
\pgfsetmiterjoin%
\pgfsetlinewidth{0.803000pt}%
\definecolor{currentstroke}{rgb}{0.000000,0.000000,0.000000}%
\pgfsetstrokecolor{currentstroke}%
\pgfsetdash{}{0pt}%
\pgfpathmoveto{\pgfqpoint{0.573768in}{0.462654in}}%
\pgfpathlineto{\pgfqpoint{0.573768in}{1.368904in}}%
\pgfusepath{stroke}%
\end{pgfscope}%
\begin{pgfscope}%
\pgfsetrectcap%
\pgfsetmiterjoin%
\pgfsetlinewidth{0.803000pt}%
\definecolor{currentstroke}{rgb}{0.000000,0.000000,0.000000}%
\pgfsetstrokecolor{currentstroke}%
\pgfsetdash{}{0pt}%
\pgfpathmoveto{\pgfqpoint{3.293768in}{0.462654in}}%
\pgfpathlineto{\pgfqpoint{3.293768in}{1.368904in}}%
\pgfusepath{stroke}%
\end{pgfscope}%
\begin{pgfscope}%
\pgfsetrectcap%
\pgfsetmiterjoin%
\pgfsetlinewidth{0.803000pt}%
\definecolor{currentstroke}{rgb}{0.000000,0.000000,0.000000}%
\pgfsetstrokecolor{currentstroke}%
\pgfsetdash{}{0pt}%
\pgfpathmoveto{\pgfqpoint{0.573768in}{0.462654in}}%
\pgfpathlineto{\pgfqpoint{3.293768in}{0.462654in}}%
\pgfusepath{stroke}%
\end{pgfscope}%
\begin{pgfscope}%
\pgfsetrectcap%
\pgfsetmiterjoin%
\pgfsetlinewidth{0.803000pt}%
\definecolor{currentstroke}{rgb}{0.000000,0.000000,0.000000}%
\pgfsetstrokecolor{currentstroke}%
\pgfsetdash{}{0pt}%
\pgfpathmoveto{\pgfqpoint{0.573768in}{1.368904in}}%
\pgfpathlineto{\pgfqpoint{3.293768in}{1.368904in}}%
\pgfusepath{stroke}%
\end{pgfscope}%
\begin{pgfscope}%
\definecolor{textcolor}{rgb}{0.000000,0.000000,0.000000}%
\pgfsetstrokecolor{textcolor}%
\pgfsetfillcolor{textcolor}%
\pgftext[x=0.628168in,y=1.305467in,left,top]{\color{textcolor}{\rmfamily\fontsize{8.000000}{9.600000}\selectfont\catcode`\^=\active\def^{\ifmmode\sp\else\^{}\fi}\catcode`\%=\active\def%{\%}(c)}}%
\end{pgfscope}%
\begin{pgfscope}%
\pgfsetbuttcap%
\pgfsetmiterjoin%
\definecolor{currentfill}{rgb}{1.000000,1.000000,1.000000}%
\pgfsetfillcolor{currentfill}%
\pgfsetlinewidth{0.000000pt}%
\definecolor{currentstroke}{rgb}{0.000000,0.000000,0.000000}%
\pgfsetstrokecolor{currentstroke}%
\pgfsetstrokeopacity{0.000000}%
\pgfsetdash{}{0pt}%
\pgfpathmoveto{\pgfqpoint{2.259013in}{3.125849in}}%
\pgfpathlineto{\pgfqpoint{3.238213in}{3.125849in}}%
\pgfpathlineto{\pgfqpoint{3.238213in}{3.452099in}}%
\pgfpathlineto{\pgfqpoint{2.259013in}{3.452099in}}%
\pgfpathlineto{\pgfqpoint{2.259013in}{3.125849in}}%
\pgfpathclose%
\pgfusepath{fill}%
\end{pgfscope}%
\begin{pgfscope}%
\pgfpathrectangle{\pgfqpoint{2.259013in}{3.125849in}}{\pgfqpoint{0.979200in}{0.326250in}}%
\pgfusepath{clip}%
\pgfsetbuttcap%
\pgfsetroundjoin%
\pgfsetlinewidth{0.401500pt}%
\definecolor{currentstroke}{rgb}{0.690196,0.690196,0.690196}%
\pgfsetstrokecolor{currentstroke}%
\pgfsetstrokeopacity{0.250000}%
\pgfsetdash{{1.480000pt}{0.640000pt}}{0.000000pt}%
\pgfpathmoveto{\pgfqpoint{2.259013in}{3.125849in}}%
\pgfpathlineto{\pgfqpoint{2.259013in}{3.452099in}}%
\pgfusepath{stroke}%
\end{pgfscope}%
\begin{pgfscope}%
\pgfsetbuttcap%
\pgfsetroundjoin%
\definecolor{currentfill}{rgb}{0.000000,0.000000,0.000000}%
\pgfsetfillcolor{currentfill}%
\pgfsetlinewidth{0.803000pt}%
\definecolor{currentstroke}{rgb}{0.000000,0.000000,0.000000}%
\pgfsetstrokecolor{currentstroke}%
\pgfsetdash{}{0pt}%
\pgfsys@defobject{currentmarker}{\pgfqpoint{0.000000in}{0.000000in}}{\pgfqpoint{0.000000in}{0.027778in}}{%
\pgfpathmoveto{\pgfqpoint{0.000000in}{0.000000in}}%
\pgfpathlineto{\pgfqpoint{0.000000in}{0.027778in}}%
\pgfusepath{stroke,fill}%
}%
\begin{pgfscope}%
\pgfsys@transformshift{2.259013in}{3.125849in}%
\pgfsys@useobject{currentmarker}{}%
\end{pgfscope}%
\end{pgfscope}%
\begin{pgfscope}%
\definecolor{textcolor}{rgb}{0.000000,0.000000,0.000000}%
\pgfsetstrokecolor{textcolor}%
\pgfsetfillcolor{textcolor}%
\pgftext[x=2.259013in,y=3.111960in,,top]{\color{textcolor}{\rmfamily\fontsize{5.000000}{6.000000}\selectfont\catcode`\^=\active\def^{\ifmmode\sp\else\^{}\fi}\catcode`\%=\active\def%{\%}5.92}}%
\end{pgfscope}%
\begin{pgfscope}%
\pgfpathrectangle{\pgfqpoint{2.259013in}{3.125849in}}{\pgfqpoint{0.979200in}{0.326250in}}%
\pgfusepath{clip}%
\pgfsetbuttcap%
\pgfsetroundjoin%
\pgfsetlinewidth{0.401500pt}%
\definecolor{currentstroke}{rgb}{0.690196,0.690196,0.690196}%
\pgfsetstrokecolor{currentstroke}%
\pgfsetstrokeopacity{0.250000}%
\pgfsetdash{{1.480000pt}{0.640000pt}}{0.000000pt}%
\pgfpathmoveto{\pgfqpoint{2.503813in}{3.125849in}}%
\pgfpathlineto{\pgfqpoint{2.503813in}{3.452099in}}%
\pgfusepath{stroke}%
\end{pgfscope}%
\begin{pgfscope}%
\pgfsetbuttcap%
\pgfsetroundjoin%
\definecolor{currentfill}{rgb}{0.000000,0.000000,0.000000}%
\pgfsetfillcolor{currentfill}%
\pgfsetlinewidth{0.803000pt}%
\definecolor{currentstroke}{rgb}{0.000000,0.000000,0.000000}%
\pgfsetstrokecolor{currentstroke}%
\pgfsetdash{}{0pt}%
\pgfsys@defobject{currentmarker}{\pgfqpoint{0.000000in}{0.000000in}}{\pgfqpoint{0.000000in}{0.027778in}}{%
\pgfpathmoveto{\pgfqpoint{0.000000in}{0.000000in}}%
\pgfpathlineto{\pgfqpoint{0.000000in}{0.027778in}}%
\pgfusepath{stroke,fill}%
}%
\begin{pgfscope}%
\pgfsys@transformshift{2.503813in}{3.125849in}%
\pgfsys@useobject{currentmarker}{}%
\end{pgfscope}%
\end{pgfscope}%
\begin{pgfscope}%
\definecolor{textcolor}{rgb}{0.000000,0.000000,0.000000}%
\pgfsetstrokecolor{textcolor}%
\pgfsetfillcolor{textcolor}%
\pgftext[x=2.503813in,y=3.111960in,,top]{\color{textcolor}{\rmfamily\fontsize{5.000000}{6.000000}\selectfont\catcode`\^=\active\def^{\ifmmode\sp\else\^{}\fi}\catcode`\%=\active\def%{\%}5.94}}%
\end{pgfscope}%
\begin{pgfscope}%
\pgfpathrectangle{\pgfqpoint{2.259013in}{3.125849in}}{\pgfqpoint{0.979200in}{0.326250in}}%
\pgfusepath{clip}%
\pgfsetbuttcap%
\pgfsetroundjoin%
\pgfsetlinewidth{0.401500pt}%
\definecolor{currentstroke}{rgb}{0.690196,0.690196,0.690196}%
\pgfsetstrokecolor{currentstroke}%
\pgfsetstrokeopacity{0.250000}%
\pgfsetdash{{1.480000pt}{0.640000pt}}{0.000000pt}%
\pgfpathmoveto{\pgfqpoint{2.748613in}{3.125849in}}%
\pgfpathlineto{\pgfqpoint{2.748613in}{3.452099in}}%
\pgfusepath{stroke}%
\end{pgfscope}%
\begin{pgfscope}%
\pgfsetbuttcap%
\pgfsetroundjoin%
\definecolor{currentfill}{rgb}{0.000000,0.000000,0.000000}%
\pgfsetfillcolor{currentfill}%
\pgfsetlinewidth{0.803000pt}%
\definecolor{currentstroke}{rgb}{0.000000,0.000000,0.000000}%
\pgfsetstrokecolor{currentstroke}%
\pgfsetdash{}{0pt}%
\pgfsys@defobject{currentmarker}{\pgfqpoint{0.000000in}{0.000000in}}{\pgfqpoint{0.000000in}{0.027778in}}{%
\pgfpathmoveto{\pgfqpoint{0.000000in}{0.000000in}}%
\pgfpathlineto{\pgfqpoint{0.000000in}{0.027778in}}%
\pgfusepath{stroke,fill}%
}%
\begin{pgfscope}%
\pgfsys@transformshift{2.748613in}{3.125849in}%
\pgfsys@useobject{currentmarker}{}%
\end{pgfscope}%
\end{pgfscope}%
\begin{pgfscope}%
\definecolor{textcolor}{rgb}{0.000000,0.000000,0.000000}%
\pgfsetstrokecolor{textcolor}%
\pgfsetfillcolor{textcolor}%
\pgftext[x=2.748613in,y=3.111960in,,top]{\color{textcolor}{\rmfamily\fontsize{5.000000}{6.000000}\selectfont\catcode`\^=\active\def^{\ifmmode\sp\else\^{}\fi}\catcode`\%=\active\def%{\%}5.96}}%
\end{pgfscope}%
\begin{pgfscope}%
\pgfpathrectangle{\pgfqpoint{2.259013in}{3.125849in}}{\pgfqpoint{0.979200in}{0.326250in}}%
\pgfusepath{clip}%
\pgfsetbuttcap%
\pgfsetroundjoin%
\pgfsetlinewidth{0.401500pt}%
\definecolor{currentstroke}{rgb}{0.690196,0.690196,0.690196}%
\pgfsetstrokecolor{currentstroke}%
\pgfsetstrokeopacity{0.250000}%
\pgfsetdash{{1.480000pt}{0.640000pt}}{0.000000pt}%
\pgfpathmoveto{\pgfqpoint{2.993413in}{3.125849in}}%
\pgfpathlineto{\pgfqpoint{2.993413in}{3.452099in}}%
\pgfusepath{stroke}%
\end{pgfscope}%
\begin{pgfscope}%
\pgfsetbuttcap%
\pgfsetroundjoin%
\definecolor{currentfill}{rgb}{0.000000,0.000000,0.000000}%
\pgfsetfillcolor{currentfill}%
\pgfsetlinewidth{0.803000pt}%
\definecolor{currentstroke}{rgb}{0.000000,0.000000,0.000000}%
\pgfsetstrokecolor{currentstroke}%
\pgfsetdash{}{0pt}%
\pgfsys@defobject{currentmarker}{\pgfqpoint{0.000000in}{0.000000in}}{\pgfqpoint{0.000000in}{0.027778in}}{%
\pgfpathmoveto{\pgfqpoint{0.000000in}{0.000000in}}%
\pgfpathlineto{\pgfqpoint{0.000000in}{0.027778in}}%
\pgfusepath{stroke,fill}%
}%
\begin{pgfscope}%
\pgfsys@transformshift{2.993413in}{3.125849in}%
\pgfsys@useobject{currentmarker}{}%
\end{pgfscope}%
\end{pgfscope}%
\begin{pgfscope}%
\definecolor{textcolor}{rgb}{0.000000,0.000000,0.000000}%
\pgfsetstrokecolor{textcolor}%
\pgfsetfillcolor{textcolor}%
\pgftext[x=2.993413in,y=3.111960in,,top]{\color{textcolor}{\rmfamily\fontsize{5.000000}{6.000000}\selectfont\catcode`\^=\active\def^{\ifmmode\sp\else\^{}\fi}\catcode`\%=\active\def%{\%}5.98}}%
\end{pgfscope}%
\begin{pgfscope}%
\pgfpathrectangle{\pgfqpoint{2.259013in}{3.125849in}}{\pgfqpoint{0.979200in}{0.326250in}}%
\pgfusepath{clip}%
\pgfsetbuttcap%
\pgfsetroundjoin%
\pgfsetlinewidth{0.401500pt}%
\definecolor{currentstroke}{rgb}{0.690196,0.690196,0.690196}%
\pgfsetstrokecolor{currentstroke}%
\pgfsetstrokeopacity{0.250000}%
\pgfsetdash{{1.480000pt}{0.640000pt}}{0.000000pt}%
\pgfpathmoveto{\pgfqpoint{3.238213in}{3.125849in}}%
\pgfpathlineto{\pgfqpoint{3.238213in}{3.452099in}}%
\pgfusepath{stroke}%
\end{pgfscope}%
\begin{pgfscope}%
\pgfsetbuttcap%
\pgfsetroundjoin%
\definecolor{currentfill}{rgb}{0.000000,0.000000,0.000000}%
\pgfsetfillcolor{currentfill}%
\pgfsetlinewidth{0.803000pt}%
\definecolor{currentstroke}{rgb}{0.000000,0.000000,0.000000}%
\pgfsetstrokecolor{currentstroke}%
\pgfsetdash{}{0pt}%
\pgfsys@defobject{currentmarker}{\pgfqpoint{0.000000in}{0.000000in}}{\pgfqpoint{0.000000in}{0.027778in}}{%
\pgfpathmoveto{\pgfqpoint{0.000000in}{0.000000in}}%
\pgfpathlineto{\pgfqpoint{0.000000in}{0.027778in}}%
\pgfusepath{stroke,fill}%
}%
\begin{pgfscope}%
\pgfsys@transformshift{3.238213in}{3.125849in}%
\pgfsys@useobject{currentmarker}{}%
\end{pgfscope}%
\end{pgfscope}%
\begin{pgfscope}%
\definecolor{textcolor}{rgb}{0.000000,0.000000,0.000000}%
\pgfsetstrokecolor{textcolor}%
\pgfsetfillcolor{textcolor}%
\pgftext[x=3.238213in,y=3.111960in,,top]{\color{textcolor}{\rmfamily\fontsize{5.000000}{6.000000}\selectfont\catcode`\^=\active\def^{\ifmmode\sp\else\^{}\fi}\catcode`\%=\active\def%{\%}6.00}}%
\end{pgfscope}%
\begin{pgfscope}%
\pgfpathrectangle{\pgfqpoint{2.259013in}{3.125849in}}{\pgfqpoint{0.979200in}{0.326250in}}%
\pgfusepath{clip}%
\pgfsetbuttcap%
\pgfsetroundjoin%
\pgfsetlinewidth{0.401500pt}%
\definecolor{currentstroke}{rgb}{0.690196,0.690196,0.690196}%
\pgfsetstrokecolor{currentstroke}%
\pgfsetstrokeopacity{0.250000}%
\pgfsetdash{{1.480000pt}{0.640000pt}}{0.000000pt}%
\pgfpathmoveto{\pgfqpoint{2.259013in}{3.125849in}}%
\pgfpathlineto{\pgfqpoint{3.238213in}{3.125849in}}%
\pgfusepath{stroke}%
\end{pgfscope}%
\begin{pgfscope}%
\pgfsetbuttcap%
\pgfsetroundjoin%
\definecolor{currentfill}{rgb}{0.000000,0.000000,0.000000}%
\pgfsetfillcolor{currentfill}%
\pgfsetlinewidth{0.803000pt}%
\definecolor{currentstroke}{rgb}{0.000000,0.000000,0.000000}%
\pgfsetstrokecolor{currentstroke}%
\pgfsetdash{}{0pt}%
\pgfsys@defobject{currentmarker}{\pgfqpoint{0.000000in}{0.000000in}}{\pgfqpoint{0.027778in}{0.000000in}}{%
\pgfpathmoveto{\pgfqpoint{0.000000in}{0.000000in}}%
\pgfpathlineto{\pgfqpoint{0.027778in}{0.000000in}}%
\pgfusepath{stroke,fill}%
}%
\begin{pgfscope}%
\pgfsys@transformshift{2.259013in}{3.125849in}%
\pgfsys@useobject{currentmarker}{}%
\end{pgfscope}%
\end{pgfscope}%
\begin{pgfscope}%
\definecolor{textcolor}{rgb}{0.000000,0.000000,0.000000}%
\pgfsetstrokecolor{textcolor}%
\pgfsetfillcolor{textcolor}%
\pgftext[x=2.075368in, y=3.101736in, left, base]{\color{textcolor}{\rmfamily\fontsize{5.000000}{6.000000}\selectfont\catcode`\^=\active\def^{\ifmmode\sp\else\^{}\fi}\catcode`\%=\active\def%{\%}9.16}}%
\end{pgfscope}%
\begin{pgfscope}%
\pgfpathrectangle{\pgfqpoint{2.259013in}{3.125849in}}{\pgfqpoint{0.979200in}{0.326250in}}%
\pgfusepath{clip}%
\pgfsetbuttcap%
\pgfsetroundjoin%
\pgfsetlinewidth{0.401500pt}%
\definecolor{currentstroke}{rgb}{0.690196,0.690196,0.690196}%
\pgfsetstrokecolor{currentstroke}%
\pgfsetstrokeopacity{0.250000}%
\pgfsetdash{{1.480000pt}{0.640000pt}}{0.000000pt}%
\pgfpathmoveto{\pgfqpoint{2.259013in}{3.288974in}}%
\pgfpathlineto{\pgfqpoint{3.238213in}{3.288974in}}%
\pgfusepath{stroke}%
\end{pgfscope}%
\begin{pgfscope}%
\pgfsetbuttcap%
\pgfsetroundjoin%
\definecolor{currentfill}{rgb}{0.000000,0.000000,0.000000}%
\pgfsetfillcolor{currentfill}%
\pgfsetlinewidth{0.803000pt}%
\definecolor{currentstroke}{rgb}{0.000000,0.000000,0.000000}%
\pgfsetstrokecolor{currentstroke}%
\pgfsetdash{}{0pt}%
\pgfsys@defobject{currentmarker}{\pgfqpoint{0.000000in}{0.000000in}}{\pgfqpoint{0.027778in}{0.000000in}}{%
\pgfpathmoveto{\pgfqpoint{0.000000in}{0.000000in}}%
\pgfpathlineto{\pgfqpoint{0.027778in}{0.000000in}}%
\pgfusepath{stroke,fill}%
}%
\begin{pgfscope}%
\pgfsys@transformshift{2.259013in}{3.288974in}%
\pgfsys@useobject{currentmarker}{}%
\end{pgfscope}%
\end{pgfscope}%
\begin{pgfscope}%
\definecolor{textcolor}{rgb}{0.000000,0.000000,0.000000}%
\pgfsetstrokecolor{textcolor}%
\pgfsetfillcolor{textcolor}%
\pgftext[x=2.075368in, y=3.264861in, left, base]{\color{textcolor}{\rmfamily\fontsize{5.000000}{6.000000}\selectfont\catcode`\^=\active\def^{\ifmmode\sp\else\^{}\fi}\catcode`\%=\active\def%{\%}9.23}}%
\end{pgfscope}%
\begin{pgfscope}%
\pgfpathrectangle{\pgfqpoint{2.259013in}{3.125849in}}{\pgfqpoint{0.979200in}{0.326250in}}%
\pgfusepath{clip}%
\pgfsetbuttcap%
\pgfsetroundjoin%
\pgfsetlinewidth{0.401500pt}%
\definecolor{currentstroke}{rgb}{0.690196,0.690196,0.690196}%
\pgfsetstrokecolor{currentstroke}%
\pgfsetstrokeopacity{0.250000}%
\pgfsetdash{{1.480000pt}{0.640000pt}}{0.000000pt}%
\pgfpathmoveto{\pgfqpoint{2.259013in}{3.452099in}}%
\pgfpathlineto{\pgfqpoint{3.238213in}{3.452099in}}%
\pgfusepath{stroke}%
\end{pgfscope}%
\begin{pgfscope}%
\pgfsetbuttcap%
\pgfsetroundjoin%
\definecolor{currentfill}{rgb}{0.000000,0.000000,0.000000}%
\pgfsetfillcolor{currentfill}%
\pgfsetlinewidth{0.803000pt}%
\definecolor{currentstroke}{rgb}{0.000000,0.000000,0.000000}%
\pgfsetstrokecolor{currentstroke}%
\pgfsetdash{}{0pt}%
\pgfsys@defobject{currentmarker}{\pgfqpoint{0.000000in}{0.000000in}}{\pgfqpoint{0.027778in}{0.000000in}}{%
\pgfpathmoveto{\pgfqpoint{0.000000in}{0.000000in}}%
\pgfpathlineto{\pgfqpoint{0.027778in}{0.000000in}}%
\pgfusepath{stroke,fill}%
}%
\begin{pgfscope}%
\pgfsys@transformshift{2.259013in}{3.452099in}%
\pgfsys@useobject{currentmarker}{}%
\end{pgfscope}%
\end{pgfscope}%
\begin{pgfscope}%
\definecolor{textcolor}{rgb}{0.000000,0.000000,0.000000}%
\pgfsetstrokecolor{textcolor}%
\pgfsetfillcolor{textcolor}%
\pgftext[x=2.075368in, y=3.427986in, left, base]{\color{textcolor}{\rmfamily\fontsize{5.000000}{6.000000}\selectfont\catcode`\^=\active\def^{\ifmmode\sp\else\^{}\fi}\catcode`\%=\active\def%{\%}9.29}}%
\end{pgfscope}%
\begin{pgfscope}%
\pgfpathrectangle{\pgfqpoint{2.259013in}{3.125849in}}{\pgfqpoint{0.979200in}{0.326250in}}%
\pgfusepath{clip}%
\pgfsetrectcap%
\pgfsetroundjoin%
\pgfsetlinewidth{1.505625pt}%
\definecolor{currentstroke}{rgb}{0.000000,0.447059,0.698039}%
\pgfsetstrokecolor{currentstroke}%
\pgfsetdash{}{0pt}%
\pgfpathmoveto{\pgfqpoint{2.255507in}{3.316100in}}%
\pgfpathlineto{\pgfqpoint{2.271085in}{3.283388in}}%
\pgfpathlineto{\pgfqpoint{2.284340in}{3.255925in}}%
\pgfpathlineto{\pgfqpoint{2.299028in}{3.229323in}}%
\pgfpathlineto{\pgfqpoint{2.313716in}{3.206746in}}%
\pgfpathlineto{\pgfqpoint{2.323508in}{3.193920in}}%
\pgfpathlineto{\pgfqpoint{2.333300in}{3.182867in}}%
\pgfpathlineto{\pgfqpoint{2.343092in}{3.173581in}}%
\pgfpathlineto{\pgfqpoint{2.352884in}{3.166055in}}%
\pgfpathlineto{\pgfqpoint{2.362676in}{3.160284in}}%
\pgfpathlineto{\pgfqpoint{2.372468in}{3.156259in}}%
\pgfpathlineto{\pgfqpoint{2.382260in}{3.153975in}}%
\pgfpathlineto{\pgfqpoint{2.392052in}{3.153424in}}%
\pgfpathlineto{\pgfqpoint{2.401844in}{3.154600in}}%
\pgfpathlineto{\pgfqpoint{2.411636in}{3.157496in}}%
\pgfpathlineto{\pgfqpoint{2.421428in}{3.162104in}}%
\pgfpathlineto{\pgfqpoint{2.431220in}{3.168418in}}%
\pgfpathlineto{\pgfqpoint{2.441012in}{3.176429in}}%
\pgfpathlineto{\pgfqpoint{2.450804in}{3.186132in}}%
\pgfpathlineto{\pgfqpoint{2.460596in}{3.197517in}}%
\pgfpathlineto{\pgfqpoint{2.475284in}{3.217736in}}%
\pgfpathlineto{\pgfqpoint{2.489972in}{3.241699in}}%
\pgfpathlineto{\pgfqpoint{2.505833in}{3.271750in}}%
\pgfpathlineto{\pgfqpoint{2.523639in}{3.309683in}}%
\pgfpathlineto{\pgfqpoint{2.535969in}{3.333928in}}%
\pgfpathlineto{\pgfqpoint{2.550657in}{3.359039in}}%
\pgfpathlineto{\pgfqpoint{2.565345in}{3.380070in}}%
\pgfpathlineto{\pgfqpoint{2.575137in}{3.391835in}}%
\pgfpathlineto{\pgfqpoint{2.584929in}{3.401802in}}%
\pgfpathlineto{\pgfqpoint{2.594721in}{3.409978in}}%
\pgfpathlineto{\pgfqpoint{2.604513in}{3.416370in}}%
\pgfpathlineto{\pgfqpoint{2.614305in}{3.420985in}}%
\pgfpathlineto{\pgfqpoint{2.624097in}{3.423830in}}%
\pgfpathlineto{\pgfqpoint{2.633889in}{3.424911in}}%
\pgfpathlineto{\pgfqpoint{2.643681in}{3.424236in}}%
\pgfpathlineto{\pgfqpoint{2.653473in}{3.421812in}}%
\pgfpathlineto{\pgfqpoint{2.663265in}{3.417646in}}%
\pgfpathlineto{\pgfqpoint{2.673057in}{3.411745in}}%
\pgfpathlineto{\pgfqpoint{2.682849in}{3.404117in}}%
\pgfpathlineto{\pgfqpoint{2.692641in}{3.394769in}}%
\pgfpathlineto{\pgfqpoint{2.702433in}{3.383709in}}%
\pgfpathlineto{\pgfqpoint{2.712225in}{3.370944in}}%
\pgfpathlineto{\pgfqpoint{2.726913in}{3.348618in}}%
\pgfpathlineto{\pgfqpoint{2.741601in}{3.322499in}}%
\pgfpathlineto{\pgfqpoint{2.756534in}{3.292086in}}%
\pgfpathlineto{\pgfqpoint{2.773941in}{3.255533in}}%
\pgfpathlineto{\pgfqpoint{2.788629in}{3.228931in}}%
\pgfpathlineto{\pgfqpoint{2.803317in}{3.206354in}}%
\pgfpathlineto{\pgfqpoint{2.813109in}{3.193528in}}%
\pgfpathlineto{\pgfqpoint{2.822901in}{3.182476in}}%
\pgfpathlineto{\pgfqpoint{2.832693in}{3.173190in}}%
\pgfpathlineto{\pgfqpoint{2.842485in}{3.165665in}}%
\pgfpathlineto{\pgfqpoint{2.852277in}{3.159894in}}%
\pgfpathlineto{\pgfqpoint{2.862069in}{3.155870in}}%
\pgfpathlineto{\pgfqpoint{2.871861in}{3.153586in}}%
\pgfpathlineto{\pgfqpoint{2.881653in}{3.153036in}}%
\pgfpathlineto{\pgfqpoint{2.891445in}{3.154213in}}%
\pgfpathlineto{\pgfqpoint{2.901237in}{3.157109in}}%
\pgfpathlineto{\pgfqpoint{2.911029in}{3.161718in}}%
\pgfpathlineto{\pgfqpoint{2.920821in}{3.168032in}}%
\pgfpathlineto{\pgfqpoint{2.930613in}{3.176044in}}%
\pgfpathlineto{\pgfqpoint{2.940405in}{3.185747in}}%
\pgfpathlineto{\pgfqpoint{2.950197in}{3.197134in}}%
\pgfpathlineto{\pgfqpoint{2.964885in}{3.217353in}}%
\pgfpathlineto{\pgfqpoint{2.979573in}{3.241318in}}%
\pgfpathlineto{\pgfqpoint{2.995433in}{3.271370in}}%
\pgfpathlineto{\pgfqpoint{3.013238in}{3.309303in}}%
\pgfpathlineto{\pgfqpoint{3.025568in}{3.333550in}}%
\pgfpathlineto{\pgfqpoint{3.040256in}{3.358662in}}%
\pgfpathlineto{\pgfqpoint{3.054944in}{3.379694in}}%
\pgfpathlineto{\pgfqpoint{3.064736in}{3.391460in}}%
\pgfpathlineto{\pgfqpoint{3.074528in}{3.401428in}}%
\pgfpathlineto{\pgfqpoint{3.084320in}{3.409605in}}%
\pgfpathlineto{\pgfqpoint{3.094112in}{3.415998in}}%
\pgfpathlineto{\pgfqpoint{3.103904in}{3.420614in}}%
\pgfpathlineto{\pgfqpoint{3.113696in}{3.423460in}}%
\pgfpathlineto{\pgfqpoint{3.123488in}{3.424543in}}%
\pgfpathlineto{\pgfqpoint{3.133280in}{3.423869in}}%
\pgfpathlineto{\pgfqpoint{3.143072in}{3.421446in}}%
\pgfpathlineto{\pgfqpoint{3.152864in}{3.417281in}}%
\pgfpathlineto{\pgfqpoint{3.162656in}{3.411381in}}%
\pgfpathlineto{\pgfqpoint{3.172448in}{3.403754in}}%
\pgfpathlineto{\pgfqpoint{3.182240in}{3.394408in}}%
\pgfpathlineto{\pgfqpoint{3.192032in}{3.383349in}}%
\pgfpathlineto{\pgfqpoint{3.201824in}{3.370585in}}%
\pgfpathlineto{\pgfqpoint{3.216512in}{3.348260in}}%
\pgfpathlineto{\pgfqpoint{3.231200in}{3.322144in}}%
\pgfpathlineto{\pgfqpoint{3.238213in}{3.308346in}}%
\pgfpathlineto{\pgfqpoint{3.238213in}{3.308346in}}%
\pgfusepath{stroke}%
\end{pgfscope}%
\begin{pgfscope}%
\pgfsetrectcap%
\pgfsetmiterjoin%
\pgfsetlinewidth{0.803000pt}%
\definecolor{currentstroke}{rgb}{0.000000,0.000000,0.000000}%
\pgfsetstrokecolor{currentstroke}%
\pgfsetdash{}{0pt}%
\pgfpathmoveto{\pgfqpoint{2.259013in}{3.125849in}}%
\pgfpathlineto{\pgfqpoint{2.259013in}{3.452099in}}%
\pgfusepath{stroke}%
\end{pgfscope}%
\begin{pgfscope}%
\pgfsetrectcap%
\pgfsetmiterjoin%
\pgfsetlinewidth{0.803000pt}%
\definecolor{currentstroke}{rgb}{0.000000,0.000000,0.000000}%
\pgfsetstrokecolor{currentstroke}%
\pgfsetdash{}{0pt}%
\pgfpathmoveto{\pgfqpoint{3.238213in}{3.125849in}}%
\pgfpathlineto{\pgfqpoint{3.238213in}{3.452099in}}%
\pgfusepath{stroke}%
\end{pgfscope}%
\begin{pgfscope}%
\pgfsetrectcap%
\pgfsetmiterjoin%
\pgfsetlinewidth{0.803000pt}%
\definecolor{currentstroke}{rgb}{0.000000,0.000000,0.000000}%
\pgfsetstrokecolor{currentstroke}%
\pgfsetdash{}{0pt}%
\pgfpathmoveto{\pgfqpoint{2.259013in}{3.125849in}}%
\pgfpathlineto{\pgfqpoint{3.238213in}{3.125849in}}%
\pgfusepath{stroke}%
\end{pgfscope}%
\begin{pgfscope}%
\pgfsetrectcap%
\pgfsetmiterjoin%
\pgfsetlinewidth{0.803000pt}%
\definecolor{currentstroke}{rgb}{0.000000,0.000000,0.000000}%
\pgfsetstrokecolor{currentstroke}%
\pgfsetdash{}{0pt}%
\pgfpathmoveto{\pgfqpoint{2.259013in}{3.452099in}}%
\pgfpathlineto{\pgfqpoint{3.238213in}{3.452099in}}%
\pgfusepath{stroke}%
\end{pgfscope}%
\begin{pgfscope}%
\pgfsetbuttcap%
\pgfsetmiterjoin%
\definecolor{currentfill}{rgb}{1.000000,1.000000,1.000000}%
\pgfsetfillcolor{currentfill}%
\pgfsetlinewidth{0.000000pt}%
\definecolor{currentstroke}{rgb}{0.000000,0.000000,0.000000}%
\pgfsetstrokecolor{currentstroke}%
\pgfsetstrokeopacity{0.000000}%
\pgfsetdash{}{0pt}%
\pgfpathmoveto{\pgfqpoint{2.259013in}{2.056474in}}%
\pgfpathlineto{\pgfqpoint{3.238213in}{2.056474in}}%
\pgfpathlineto{\pgfqpoint{3.238213in}{2.382724in}}%
\pgfpathlineto{\pgfqpoint{2.259013in}{2.382724in}}%
\pgfpathlineto{\pgfqpoint{2.259013in}{2.056474in}}%
\pgfpathclose%
\pgfusepath{fill}%
\end{pgfscope}%
\begin{pgfscope}%
\pgfpathrectangle{\pgfqpoint{2.259013in}{2.056474in}}{\pgfqpoint{0.979200in}{0.326250in}}%
\pgfusepath{clip}%
\pgfsetbuttcap%
\pgfsetroundjoin%
\pgfsetlinewidth{0.401500pt}%
\definecolor{currentstroke}{rgb}{0.690196,0.690196,0.690196}%
\pgfsetstrokecolor{currentstroke}%
\pgfsetstrokeopacity{0.250000}%
\pgfsetdash{{1.480000pt}{0.640000pt}}{0.000000pt}%
\pgfpathmoveto{\pgfqpoint{2.259013in}{2.056474in}}%
\pgfpathlineto{\pgfqpoint{2.259013in}{2.382724in}}%
\pgfusepath{stroke}%
\end{pgfscope}%
\begin{pgfscope}%
\pgfsetbuttcap%
\pgfsetroundjoin%
\definecolor{currentfill}{rgb}{0.000000,0.000000,0.000000}%
\pgfsetfillcolor{currentfill}%
\pgfsetlinewidth{0.803000pt}%
\definecolor{currentstroke}{rgb}{0.000000,0.000000,0.000000}%
\pgfsetstrokecolor{currentstroke}%
\pgfsetdash{}{0pt}%
\pgfsys@defobject{currentmarker}{\pgfqpoint{0.000000in}{0.000000in}}{\pgfqpoint{0.000000in}{0.027778in}}{%
\pgfpathmoveto{\pgfqpoint{0.000000in}{0.000000in}}%
\pgfpathlineto{\pgfqpoint{0.000000in}{0.027778in}}%
\pgfusepath{stroke,fill}%
}%
\begin{pgfscope}%
\pgfsys@transformshift{2.259013in}{2.056474in}%
\pgfsys@useobject{currentmarker}{}%
\end{pgfscope}%
\end{pgfscope}%
\begin{pgfscope}%
\definecolor{textcolor}{rgb}{0.000000,0.000000,0.000000}%
\pgfsetstrokecolor{textcolor}%
\pgfsetfillcolor{textcolor}%
\pgftext[x=2.259013in,y=2.042585in,,top]{\color{textcolor}{\rmfamily\fontsize{5.000000}{6.000000}\selectfont\catcode`\^=\active\def^{\ifmmode\sp\else\^{}\fi}\catcode`\%=\active\def%{\%}5.92}}%
\end{pgfscope}%
\begin{pgfscope}%
\pgfpathrectangle{\pgfqpoint{2.259013in}{2.056474in}}{\pgfqpoint{0.979200in}{0.326250in}}%
\pgfusepath{clip}%
\pgfsetbuttcap%
\pgfsetroundjoin%
\pgfsetlinewidth{0.401500pt}%
\definecolor{currentstroke}{rgb}{0.690196,0.690196,0.690196}%
\pgfsetstrokecolor{currentstroke}%
\pgfsetstrokeopacity{0.250000}%
\pgfsetdash{{1.480000pt}{0.640000pt}}{0.000000pt}%
\pgfpathmoveto{\pgfqpoint{2.503813in}{2.056474in}}%
\pgfpathlineto{\pgfqpoint{2.503813in}{2.382724in}}%
\pgfusepath{stroke}%
\end{pgfscope}%
\begin{pgfscope}%
\pgfsetbuttcap%
\pgfsetroundjoin%
\definecolor{currentfill}{rgb}{0.000000,0.000000,0.000000}%
\pgfsetfillcolor{currentfill}%
\pgfsetlinewidth{0.803000pt}%
\definecolor{currentstroke}{rgb}{0.000000,0.000000,0.000000}%
\pgfsetstrokecolor{currentstroke}%
\pgfsetdash{}{0pt}%
\pgfsys@defobject{currentmarker}{\pgfqpoint{0.000000in}{0.000000in}}{\pgfqpoint{0.000000in}{0.027778in}}{%
\pgfpathmoveto{\pgfqpoint{0.000000in}{0.000000in}}%
\pgfpathlineto{\pgfqpoint{0.000000in}{0.027778in}}%
\pgfusepath{stroke,fill}%
}%
\begin{pgfscope}%
\pgfsys@transformshift{2.503813in}{2.056474in}%
\pgfsys@useobject{currentmarker}{}%
\end{pgfscope}%
\end{pgfscope}%
\begin{pgfscope}%
\definecolor{textcolor}{rgb}{0.000000,0.000000,0.000000}%
\pgfsetstrokecolor{textcolor}%
\pgfsetfillcolor{textcolor}%
\pgftext[x=2.503813in,y=2.042585in,,top]{\color{textcolor}{\rmfamily\fontsize{5.000000}{6.000000}\selectfont\catcode`\^=\active\def^{\ifmmode\sp\else\^{}\fi}\catcode`\%=\active\def%{\%}5.94}}%
\end{pgfscope}%
\begin{pgfscope}%
\pgfpathrectangle{\pgfqpoint{2.259013in}{2.056474in}}{\pgfqpoint{0.979200in}{0.326250in}}%
\pgfusepath{clip}%
\pgfsetbuttcap%
\pgfsetroundjoin%
\pgfsetlinewidth{0.401500pt}%
\definecolor{currentstroke}{rgb}{0.690196,0.690196,0.690196}%
\pgfsetstrokecolor{currentstroke}%
\pgfsetstrokeopacity{0.250000}%
\pgfsetdash{{1.480000pt}{0.640000pt}}{0.000000pt}%
\pgfpathmoveto{\pgfqpoint{2.748613in}{2.056474in}}%
\pgfpathlineto{\pgfqpoint{2.748613in}{2.382724in}}%
\pgfusepath{stroke}%
\end{pgfscope}%
\begin{pgfscope}%
\pgfsetbuttcap%
\pgfsetroundjoin%
\definecolor{currentfill}{rgb}{0.000000,0.000000,0.000000}%
\pgfsetfillcolor{currentfill}%
\pgfsetlinewidth{0.803000pt}%
\definecolor{currentstroke}{rgb}{0.000000,0.000000,0.000000}%
\pgfsetstrokecolor{currentstroke}%
\pgfsetdash{}{0pt}%
\pgfsys@defobject{currentmarker}{\pgfqpoint{0.000000in}{0.000000in}}{\pgfqpoint{0.000000in}{0.027778in}}{%
\pgfpathmoveto{\pgfqpoint{0.000000in}{0.000000in}}%
\pgfpathlineto{\pgfqpoint{0.000000in}{0.027778in}}%
\pgfusepath{stroke,fill}%
}%
\begin{pgfscope}%
\pgfsys@transformshift{2.748613in}{2.056474in}%
\pgfsys@useobject{currentmarker}{}%
\end{pgfscope}%
\end{pgfscope}%
\begin{pgfscope}%
\definecolor{textcolor}{rgb}{0.000000,0.000000,0.000000}%
\pgfsetstrokecolor{textcolor}%
\pgfsetfillcolor{textcolor}%
\pgftext[x=2.748613in,y=2.042585in,,top]{\color{textcolor}{\rmfamily\fontsize{5.000000}{6.000000}\selectfont\catcode`\^=\active\def^{\ifmmode\sp\else\^{}\fi}\catcode`\%=\active\def%{\%}5.96}}%
\end{pgfscope}%
\begin{pgfscope}%
\pgfpathrectangle{\pgfqpoint{2.259013in}{2.056474in}}{\pgfqpoint{0.979200in}{0.326250in}}%
\pgfusepath{clip}%
\pgfsetbuttcap%
\pgfsetroundjoin%
\pgfsetlinewidth{0.401500pt}%
\definecolor{currentstroke}{rgb}{0.690196,0.690196,0.690196}%
\pgfsetstrokecolor{currentstroke}%
\pgfsetstrokeopacity{0.250000}%
\pgfsetdash{{1.480000pt}{0.640000pt}}{0.000000pt}%
\pgfpathmoveto{\pgfqpoint{2.993413in}{2.056474in}}%
\pgfpathlineto{\pgfqpoint{2.993413in}{2.382724in}}%
\pgfusepath{stroke}%
\end{pgfscope}%
\begin{pgfscope}%
\pgfsetbuttcap%
\pgfsetroundjoin%
\definecolor{currentfill}{rgb}{0.000000,0.000000,0.000000}%
\pgfsetfillcolor{currentfill}%
\pgfsetlinewidth{0.803000pt}%
\definecolor{currentstroke}{rgb}{0.000000,0.000000,0.000000}%
\pgfsetstrokecolor{currentstroke}%
\pgfsetdash{}{0pt}%
\pgfsys@defobject{currentmarker}{\pgfqpoint{0.000000in}{0.000000in}}{\pgfqpoint{0.000000in}{0.027778in}}{%
\pgfpathmoveto{\pgfqpoint{0.000000in}{0.000000in}}%
\pgfpathlineto{\pgfqpoint{0.000000in}{0.027778in}}%
\pgfusepath{stroke,fill}%
}%
\begin{pgfscope}%
\pgfsys@transformshift{2.993413in}{2.056474in}%
\pgfsys@useobject{currentmarker}{}%
\end{pgfscope}%
\end{pgfscope}%
\begin{pgfscope}%
\definecolor{textcolor}{rgb}{0.000000,0.000000,0.000000}%
\pgfsetstrokecolor{textcolor}%
\pgfsetfillcolor{textcolor}%
\pgftext[x=2.993413in,y=2.042585in,,top]{\color{textcolor}{\rmfamily\fontsize{5.000000}{6.000000}\selectfont\catcode`\^=\active\def^{\ifmmode\sp\else\^{}\fi}\catcode`\%=\active\def%{\%}5.98}}%
\end{pgfscope}%
\begin{pgfscope}%
\pgfpathrectangle{\pgfqpoint{2.259013in}{2.056474in}}{\pgfqpoint{0.979200in}{0.326250in}}%
\pgfusepath{clip}%
\pgfsetbuttcap%
\pgfsetroundjoin%
\pgfsetlinewidth{0.401500pt}%
\definecolor{currentstroke}{rgb}{0.690196,0.690196,0.690196}%
\pgfsetstrokecolor{currentstroke}%
\pgfsetstrokeopacity{0.250000}%
\pgfsetdash{{1.480000pt}{0.640000pt}}{0.000000pt}%
\pgfpathmoveto{\pgfqpoint{3.238213in}{2.056474in}}%
\pgfpathlineto{\pgfqpoint{3.238213in}{2.382724in}}%
\pgfusepath{stroke}%
\end{pgfscope}%
\begin{pgfscope}%
\pgfsetbuttcap%
\pgfsetroundjoin%
\definecolor{currentfill}{rgb}{0.000000,0.000000,0.000000}%
\pgfsetfillcolor{currentfill}%
\pgfsetlinewidth{0.803000pt}%
\definecolor{currentstroke}{rgb}{0.000000,0.000000,0.000000}%
\pgfsetstrokecolor{currentstroke}%
\pgfsetdash{}{0pt}%
\pgfsys@defobject{currentmarker}{\pgfqpoint{0.000000in}{0.000000in}}{\pgfqpoint{0.000000in}{0.027778in}}{%
\pgfpathmoveto{\pgfqpoint{0.000000in}{0.000000in}}%
\pgfpathlineto{\pgfqpoint{0.000000in}{0.027778in}}%
\pgfusepath{stroke,fill}%
}%
\begin{pgfscope}%
\pgfsys@transformshift{3.238213in}{2.056474in}%
\pgfsys@useobject{currentmarker}{}%
\end{pgfscope}%
\end{pgfscope}%
\begin{pgfscope}%
\definecolor{textcolor}{rgb}{0.000000,0.000000,0.000000}%
\pgfsetstrokecolor{textcolor}%
\pgfsetfillcolor{textcolor}%
\pgftext[x=3.238213in,y=2.042585in,,top]{\color{textcolor}{\rmfamily\fontsize{5.000000}{6.000000}\selectfont\catcode`\^=\active\def^{\ifmmode\sp\else\^{}\fi}\catcode`\%=\active\def%{\%}6.00}}%
\end{pgfscope}%
\begin{pgfscope}%
\pgfpathrectangle{\pgfqpoint{2.259013in}{2.056474in}}{\pgfqpoint{0.979200in}{0.326250in}}%
\pgfusepath{clip}%
\pgfsetbuttcap%
\pgfsetroundjoin%
\pgfsetlinewidth{0.401500pt}%
\definecolor{currentstroke}{rgb}{0.690196,0.690196,0.690196}%
\pgfsetstrokecolor{currentstroke}%
\pgfsetstrokeopacity{0.250000}%
\pgfsetdash{{1.480000pt}{0.640000pt}}{0.000000pt}%
\pgfpathmoveto{\pgfqpoint{2.259013in}{2.056474in}}%
\pgfpathlineto{\pgfqpoint{3.238213in}{2.056474in}}%
\pgfusepath{stroke}%
\end{pgfscope}%
\begin{pgfscope}%
\pgfsetbuttcap%
\pgfsetroundjoin%
\definecolor{currentfill}{rgb}{0.000000,0.000000,0.000000}%
\pgfsetfillcolor{currentfill}%
\pgfsetlinewidth{0.803000pt}%
\definecolor{currentstroke}{rgb}{0.000000,0.000000,0.000000}%
\pgfsetstrokecolor{currentstroke}%
\pgfsetdash{}{0pt}%
\pgfsys@defobject{currentmarker}{\pgfqpoint{0.000000in}{0.000000in}}{\pgfqpoint{0.027778in}{0.000000in}}{%
\pgfpathmoveto{\pgfqpoint{0.000000in}{0.000000in}}%
\pgfpathlineto{\pgfqpoint{0.027778in}{0.000000in}}%
\pgfusepath{stroke,fill}%
}%
\begin{pgfscope}%
\pgfsys@transformshift{2.259013in}{2.056474in}%
\pgfsys@useobject{currentmarker}{}%
\end{pgfscope}%
\end{pgfscope}%
\begin{pgfscope}%
\definecolor{textcolor}{rgb}{0.000000,0.000000,0.000000}%
\pgfsetstrokecolor{textcolor}%
\pgfsetfillcolor{textcolor}%
\pgftext[x=2.075368in, y=2.032361in, left, base]{\color{textcolor}{\rmfamily\fontsize{5.000000}{6.000000}\selectfont\catcode`\^=\active\def^{\ifmmode\sp\else\^{}\fi}\catcode`\%=\active\def%{\%}3.32}}%
\end{pgfscope}%
\begin{pgfscope}%
\pgfpathrectangle{\pgfqpoint{2.259013in}{2.056474in}}{\pgfqpoint{0.979200in}{0.326250in}}%
\pgfusepath{clip}%
\pgfsetbuttcap%
\pgfsetroundjoin%
\pgfsetlinewidth{0.401500pt}%
\definecolor{currentstroke}{rgb}{0.690196,0.690196,0.690196}%
\pgfsetstrokecolor{currentstroke}%
\pgfsetstrokeopacity{0.250000}%
\pgfsetdash{{1.480000pt}{0.640000pt}}{0.000000pt}%
\pgfpathmoveto{\pgfqpoint{2.259013in}{2.219599in}}%
\pgfpathlineto{\pgfqpoint{3.238213in}{2.219599in}}%
\pgfusepath{stroke}%
\end{pgfscope}%
\begin{pgfscope}%
\pgfsetbuttcap%
\pgfsetroundjoin%
\definecolor{currentfill}{rgb}{0.000000,0.000000,0.000000}%
\pgfsetfillcolor{currentfill}%
\pgfsetlinewidth{0.803000pt}%
\definecolor{currentstroke}{rgb}{0.000000,0.000000,0.000000}%
\pgfsetstrokecolor{currentstroke}%
\pgfsetdash{}{0pt}%
\pgfsys@defobject{currentmarker}{\pgfqpoint{0.000000in}{0.000000in}}{\pgfqpoint{0.027778in}{0.000000in}}{%
\pgfpathmoveto{\pgfqpoint{0.000000in}{0.000000in}}%
\pgfpathlineto{\pgfqpoint{0.027778in}{0.000000in}}%
\pgfusepath{stroke,fill}%
}%
\begin{pgfscope}%
\pgfsys@transformshift{2.259013in}{2.219599in}%
\pgfsys@useobject{currentmarker}{}%
\end{pgfscope}%
\end{pgfscope}%
\begin{pgfscope}%
\definecolor{textcolor}{rgb}{0.000000,0.000000,0.000000}%
\pgfsetstrokecolor{textcolor}%
\pgfsetfillcolor{textcolor}%
\pgftext[x=2.075368in, y=2.195486in, left, base]{\color{textcolor}{\rmfamily\fontsize{5.000000}{6.000000}\selectfont\catcode`\^=\active\def^{\ifmmode\sp\else\^{}\fi}\catcode`\%=\active\def%{\%}4.61}}%
\end{pgfscope}%
\begin{pgfscope}%
\pgfpathrectangle{\pgfqpoint{2.259013in}{2.056474in}}{\pgfqpoint{0.979200in}{0.326250in}}%
\pgfusepath{clip}%
\pgfsetbuttcap%
\pgfsetroundjoin%
\pgfsetlinewidth{0.401500pt}%
\definecolor{currentstroke}{rgb}{0.690196,0.690196,0.690196}%
\pgfsetstrokecolor{currentstroke}%
\pgfsetstrokeopacity{0.250000}%
\pgfsetdash{{1.480000pt}{0.640000pt}}{0.000000pt}%
\pgfpathmoveto{\pgfqpoint{2.259013in}{2.382724in}}%
\pgfpathlineto{\pgfqpoint{3.238213in}{2.382724in}}%
\pgfusepath{stroke}%
\end{pgfscope}%
\begin{pgfscope}%
\pgfsetbuttcap%
\pgfsetroundjoin%
\definecolor{currentfill}{rgb}{0.000000,0.000000,0.000000}%
\pgfsetfillcolor{currentfill}%
\pgfsetlinewidth{0.803000pt}%
\definecolor{currentstroke}{rgb}{0.000000,0.000000,0.000000}%
\pgfsetstrokecolor{currentstroke}%
\pgfsetdash{}{0pt}%
\pgfsys@defobject{currentmarker}{\pgfqpoint{0.000000in}{0.000000in}}{\pgfqpoint{0.027778in}{0.000000in}}{%
\pgfpathmoveto{\pgfqpoint{0.000000in}{0.000000in}}%
\pgfpathlineto{\pgfqpoint{0.027778in}{0.000000in}}%
\pgfusepath{stroke,fill}%
}%
\begin{pgfscope}%
\pgfsys@transformshift{2.259013in}{2.382724in}%
\pgfsys@useobject{currentmarker}{}%
\end{pgfscope}%
\end{pgfscope}%
\begin{pgfscope}%
\definecolor{textcolor}{rgb}{0.000000,0.000000,0.000000}%
\pgfsetstrokecolor{textcolor}%
\pgfsetfillcolor{textcolor}%
\pgftext[x=2.075368in, y=2.358611in, left, base]{\color{textcolor}{\rmfamily\fontsize{5.000000}{6.000000}\selectfont\catcode`\^=\active\def^{\ifmmode\sp\else\^{}\fi}\catcode`\%=\active\def%{\%}5.90}}%
\end{pgfscope}%
\begin{pgfscope}%
\pgfpathrectangle{\pgfqpoint{2.259013in}{2.056474in}}{\pgfqpoint{0.979200in}{0.326250in}}%
\pgfusepath{clip}%
\pgfsetrectcap%
\pgfsetroundjoin%
\pgfsetlinewidth{1.505625pt}%
\definecolor{currentstroke}{rgb}{0.835294,0.368627,0.000000}%
\pgfsetstrokecolor{currentstroke}%
\pgfsetdash{}{0pt}%
\pgfpathmoveto{\pgfqpoint{2.255507in}{2.098818in}}%
\pgfpathlineto{\pgfqpoint{2.269873in}{2.083771in}}%
\pgfpathlineto{\pgfqpoint{2.271937in}{2.085150in}}%
\pgfpathlineto{\pgfqpoint{2.277908in}{2.091359in}}%
\pgfpathlineto{\pgfqpoint{2.516708in}{2.355536in}}%
\pgfpathlineto{\pgfqpoint{2.519067in}{2.354365in}}%
\pgfpathlineto{\pgfqpoint{2.759715in}{2.083872in}}%
\pgfpathlineto{\pgfqpoint{2.762570in}{2.086052in}}%
\pgfpathlineto{\pgfqpoint{2.773941in}{2.098483in}}%
\pgfpathlineto{\pgfqpoint{3.006308in}{2.355532in}}%
\pgfpathlineto{\pgfqpoint{3.008667in}{2.354360in}}%
\pgfpathlineto{\pgfqpoint{3.238213in}{2.094857in}}%
\pgfpathlineto{\pgfqpoint{3.238213in}{2.094857in}}%
\pgfusepath{stroke}%
\end{pgfscope}%
\begin{pgfscope}%
\pgfsetrectcap%
\pgfsetmiterjoin%
\pgfsetlinewidth{0.803000pt}%
\definecolor{currentstroke}{rgb}{0.000000,0.000000,0.000000}%
\pgfsetstrokecolor{currentstroke}%
\pgfsetdash{}{0pt}%
\pgfpathmoveto{\pgfqpoint{2.259013in}{2.056474in}}%
\pgfpathlineto{\pgfqpoint{2.259013in}{2.382724in}}%
\pgfusepath{stroke}%
\end{pgfscope}%
\begin{pgfscope}%
\pgfsetrectcap%
\pgfsetmiterjoin%
\pgfsetlinewidth{0.803000pt}%
\definecolor{currentstroke}{rgb}{0.000000,0.000000,0.000000}%
\pgfsetstrokecolor{currentstroke}%
\pgfsetdash{}{0pt}%
\pgfpathmoveto{\pgfqpoint{3.238213in}{2.056474in}}%
\pgfpathlineto{\pgfqpoint{3.238213in}{2.382724in}}%
\pgfusepath{stroke}%
\end{pgfscope}%
\begin{pgfscope}%
\pgfsetrectcap%
\pgfsetmiterjoin%
\pgfsetlinewidth{0.803000pt}%
\definecolor{currentstroke}{rgb}{0.000000,0.000000,0.000000}%
\pgfsetstrokecolor{currentstroke}%
\pgfsetdash{}{0pt}%
\pgfpathmoveto{\pgfqpoint{2.259013in}{2.056474in}}%
\pgfpathlineto{\pgfqpoint{3.238213in}{2.056474in}}%
\pgfusepath{stroke}%
\end{pgfscope}%
\begin{pgfscope}%
\pgfsetrectcap%
\pgfsetmiterjoin%
\pgfsetlinewidth{0.803000pt}%
\definecolor{currentstroke}{rgb}{0.000000,0.000000,0.000000}%
\pgfsetstrokecolor{currentstroke}%
\pgfsetdash{}{0pt}%
\pgfpathmoveto{\pgfqpoint{2.259013in}{2.382724in}}%
\pgfpathlineto{\pgfqpoint{3.238213in}{2.382724in}}%
\pgfusepath{stroke}%
\end{pgfscope}%
\begin{pgfscope}%
\pgfsetbuttcap%
\pgfsetmiterjoin%
\definecolor{currentfill}{rgb}{1.000000,1.000000,1.000000}%
\pgfsetfillcolor{currentfill}%
\pgfsetlinewidth{0.000000pt}%
\definecolor{currentstroke}{rgb}{0.000000,0.000000,0.000000}%
\pgfsetstrokecolor{currentstroke}%
\pgfsetstrokeopacity{0.000000}%
\pgfsetdash{}{0pt}%
\pgfpathmoveto{\pgfqpoint{2.259013in}{0.987099in}}%
\pgfpathlineto{\pgfqpoint{3.238213in}{0.987099in}}%
\pgfpathlineto{\pgfqpoint{3.238213in}{1.313349in}}%
\pgfpathlineto{\pgfqpoint{2.259013in}{1.313349in}}%
\pgfpathlineto{\pgfqpoint{2.259013in}{0.987099in}}%
\pgfpathclose%
\pgfusepath{fill}%
\end{pgfscope}%
\begin{pgfscope}%
\pgfpathrectangle{\pgfqpoint{2.259013in}{0.987099in}}{\pgfqpoint{0.979200in}{0.326250in}}%
\pgfusepath{clip}%
\pgfsetbuttcap%
\pgfsetroundjoin%
\pgfsetlinewidth{0.401500pt}%
\definecolor{currentstroke}{rgb}{0.690196,0.690196,0.690196}%
\pgfsetstrokecolor{currentstroke}%
\pgfsetstrokeopacity{0.250000}%
\pgfsetdash{{1.480000pt}{0.640000pt}}{0.000000pt}%
\pgfpathmoveto{\pgfqpoint{2.259013in}{0.987099in}}%
\pgfpathlineto{\pgfqpoint{2.259013in}{1.313349in}}%
\pgfusepath{stroke}%
\end{pgfscope}%
\begin{pgfscope}%
\pgfsetbuttcap%
\pgfsetroundjoin%
\definecolor{currentfill}{rgb}{0.000000,0.000000,0.000000}%
\pgfsetfillcolor{currentfill}%
\pgfsetlinewidth{0.803000pt}%
\definecolor{currentstroke}{rgb}{0.000000,0.000000,0.000000}%
\pgfsetstrokecolor{currentstroke}%
\pgfsetdash{}{0pt}%
\pgfsys@defobject{currentmarker}{\pgfqpoint{0.000000in}{0.000000in}}{\pgfqpoint{0.000000in}{0.027778in}}{%
\pgfpathmoveto{\pgfqpoint{0.000000in}{0.000000in}}%
\pgfpathlineto{\pgfqpoint{0.000000in}{0.027778in}}%
\pgfusepath{stroke,fill}%
}%
\begin{pgfscope}%
\pgfsys@transformshift{2.259013in}{0.987099in}%
\pgfsys@useobject{currentmarker}{}%
\end{pgfscope}%
\end{pgfscope}%
\begin{pgfscope}%
\definecolor{textcolor}{rgb}{0.000000,0.000000,0.000000}%
\pgfsetstrokecolor{textcolor}%
\pgfsetfillcolor{textcolor}%
\pgftext[x=2.259013in,y=0.973210in,,top]{\color{textcolor}{\rmfamily\fontsize{5.000000}{6.000000}\selectfont\catcode`\^=\active\def^{\ifmmode\sp\else\^{}\fi}\catcode`\%=\active\def%{\%}5.92}}%
\end{pgfscope}%
\begin{pgfscope}%
\pgfpathrectangle{\pgfqpoint{2.259013in}{0.987099in}}{\pgfqpoint{0.979200in}{0.326250in}}%
\pgfusepath{clip}%
\pgfsetbuttcap%
\pgfsetroundjoin%
\pgfsetlinewidth{0.401500pt}%
\definecolor{currentstroke}{rgb}{0.690196,0.690196,0.690196}%
\pgfsetstrokecolor{currentstroke}%
\pgfsetstrokeopacity{0.250000}%
\pgfsetdash{{1.480000pt}{0.640000pt}}{0.000000pt}%
\pgfpathmoveto{\pgfqpoint{2.503813in}{0.987099in}}%
\pgfpathlineto{\pgfqpoint{2.503813in}{1.313349in}}%
\pgfusepath{stroke}%
\end{pgfscope}%
\begin{pgfscope}%
\pgfsetbuttcap%
\pgfsetroundjoin%
\definecolor{currentfill}{rgb}{0.000000,0.000000,0.000000}%
\pgfsetfillcolor{currentfill}%
\pgfsetlinewidth{0.803000pt}%
\definecolor{currentstroke}{rgb}{0.000000,0.000000,0.000000}%
\pgfsetstrokecolor{currentstroke}%
\pgfsetdash{}{0pt}%
\pgfsys@defobject{currentmarker}{\pgfqpoint{0.000000in}{0.000000in}}{\pgfqpoint{0.000000in}{0.027778in}}{%
\pgfpathmoveto{\pgfqpoint{0.000000in}{0.000000in}}%
\pgfpathlineto{\pgfqpoint{0.000000in}{0.027778in}}%
\pgfusepath{stroke,fill}%
}%
\begin{pgfscope}%
\pgfsys@transformshift{2.503813in}{0.987099in}%
\pgfsys@useobject{currentmarker}{}%
\end{pgfscope}%
\end{pgfscope}%
\begin{pgfscope}%
\definecolor{textcolor}{rgb}{0.000000,0.000000,0.000000}%
\pgfsetstrokecolor{textcolor}%
\pgfsetfillcolor{textcolor}%
\pgftext[x=2.503813in,y=0.973210in,,top]{\color{textcolor}{\rmfamily\fontsize{5.000000}{6.000000}\selectfont\catcode`\^=\active\def^{\ifmmode\sp\else\^{}\fi}\catcode`\%=\active\def%{\%}5.94}}%
\end{pgfscope}%
\begin{pgfscope}%
\pgfpathrectangle{\pgfqpoint{2.259013in}{0.987099in}}{\pgfqpoint{0.979200in}{0.326250in}}%
\pgfusepath{clip}%
\pgfsetbuttcap%
\pgfsetroundjoin%
\pgfsetlinewidth{0.401500pt}%
\definecolor{currentstroke}{rgb}{0.690196,0.690196,0.690196}%
\pgfsetstrokecolor{currentstroke}%
\pgfsetstrokeopacity{0.250000}%
\pgfsetdash{{1.480000pt}{0.640000pt}}{0.000000pt}%
\pgfpathmoveto{\pgfqpoint{2.748613in}{0.987099in}}%
\pgfpathlineto{\pgfqpoint{2.748613in}{1.313349in}}%
\pgfusepath{stroke}%
\end{pgfscope}%
\begin{pgfscope}%
\pgfsetbuttcap%
\pgfsetroundjoin%
\definecolor{currentfill}{rgb}{0.000000,0.000000,0.000000}%
\pgfsetfillcolor{currentfill}%
\pgfsetlinewidth{0.803000pt}%
\definecolor{currentstroke}{rgb}{0.000000,0.000000,0.000000}%
\pgfsetstrokecolor{currentstroke}%
\pgfsetdash{}{0pt}%
\pgfsys@defobject{currentmarker}{\pgfqpoint{0.000000in}{0.000000in}}{\pgfqpoint{0.000000in}{0.027778in}}{%
\pgfpathmoveto{\pgfqpoint{0.000000in}{0.000000in}}%
\pgfpathlineto{\pgfqpoint{0.000000in}{0.027778in}}%
\pgfusepath{stroke,fill}%
}%
\begin{pgfscope}%
\pgfsys@transformshift{2.748613in}{0.987099in}%
\pgfsys@useobject{currentmarker}{}%
\end{pgfscope}%
\end{pgfscope}%
\begin{pgfscope}%
\definecolor{textcolor}{rgb}{0.000000,0.000000,0.000000}%
\pgfsetstrokecolor{textcolor}%
\pgfsetfillcolor{textcolor}%
\pgftext[x=2.748613in,y=0.973210in,,top]{\color{textcolor}{\rmfamily\fontsize{5.000000}{6.000000}\selectfont\catcode`\^=\active\def^{\ifmmode\sp\else\^{}\fi}\catcode`\%=\active\def%{\%}5.96}}%
\end{pgfscope}%
\begin{pgfscope}%
\pgfpathrectangle{\pgfqpoint{2.259013in}{0.987099in}}{\pgfqpoint{0.979200in}{0.326250in}}%
\pgfusepath{clip}%
\pgfsetbuttcap%
\pgfsetroundjoin%
\pgfsetlinewidth{0.401500pt}%
\definecolor{currentstroke}{rgb}{0.690196,0.690196,0.690196}%
\pgfsetstrokecolor{currentstroke}%
\pgfsetstrokeopacity{0.250000}%
\pgfsetdash{{1.480000pt}{0.640000pt}}{0.000000pt}%
\pgfpathmoveto{\pgfqpoint{2.993413in}{0.987099in}}%
\pgfpathlineto{\pgfqpoint{2.993413in}{1.313349in}}%
\pgfusepath{stroke}%
\end{pgfscope}%
\begin{pgfscope}%
\pgfsetbuttcap%
\pgfsetroundjoin%
\definecolor{currentfill}{rgb}{0.000000,0.000000,0.000000}%
\pgfsetfillcolor{currentfill}%
\pgfsetlinewidth{0.803000pt}%
\definecolor{currentstroke}{rgb}{0.000000,0.000000,0.000000}%
\pgfsetstrokecolor{currentstroke}%
\pgfsetdash{}{0pt}%
\pgfsys@defobject{currentmarker}{\pgfqpoint{0.000000in}{0.000000in}}{\pgfqpoint{0.000000in}{0.027778in}}{%
\pgfpathmoveto{\pgfqpoint{0.000000in}{0.000000in}}%
\pgfpathlineto{\pgfqpoint{0.000000in}{0.027778in}}%
\pgfusepath{stroke,fill}%
}%
\begin{pgfscope}%
\pgfsys@transformshift{2.993413in}{0.987099in}%
\pgfsys@useobject{currentmarker}{}%
\end{pgfscope}%
\end{pgfscope}%
\begin{pgfscope}%
\definecolor{textcolor}{rgb}{0.000000,0.000000,0.000000}%
\pgfsetstrokecolor{textcolor}%
\pgfsetfillcolor{textcolor}%
\pgftext[x=2.993413in,y=0.973210in,,top]{\color{textcolor}{\rmfamily\fontsize{5.000000}{6.000000}\selectfont\catcode`\^=\active\def^{\ifmmode\sp\else\^{}\fi}\catcode`\%=\active\def%{\%}5.98}}%
\end{pgfscope}%
\begin{pgfscope}%
\pgfpathrectangle{\pgfqpoint{2.259013in}{0.987099in}}{\pgfqpoint{0.979200in}{0.326250in}}%
\pgfusepath{clip}%
\pgfsetbuttcap%
\pgfsetroundjoin%
\pgfsetlinewidth{0.401500pt}%
\definecolor{currentstroke}{rgb}{0.690196,0.690196,0.690196}%
\pgfsetstrokecolor{currentstroke}%
\pgfsetstrokeopacity{0.250000}%
\pgfsetdash{{1.480000pt}{0.640000pt}}{0.000000pt}%
\pgfpathmoveto{\pgfqpoint{3.238213in}{0.987099in}}%
\pgfpathlineto{\pgfqpoint{3.238213in}{1.313349in}}%
\pgfusepath{stroke}%
\end{pgfscope}%
\begin{pgfscope}%
\pgfsetbuttcap%
\pgfsetroundjoin%
\definecolor{currentfill}{rgb}{0.000000,0.000000,0.000000}%
\pgfsetfillcolor{currentfill}%
\pgfsetlinewidth{0.803000pt}%
\definecolor{currentstroke}{rgb}{0.000000,0.000000,0.000000}%
\pgfsetstrokecolor{currentstroke}%
\pgfsetdash{}{0pt}%
\pgfsys@defobject{currentmarker}{\pgfqpoint{0.000000in}{0.000000in}}{\pgfqpoint{0.000000in}{0.027778in}}{%
\pgfpathmoveto{\pgfqpoint{0.000000in}{0.000000in}}%
\pgfpathlineto{\pgfqpoint{0.000000in}{0.027778in}}%
\pgfusepath{stroke,fill}%
}%
\begin{pgfscope}%
\pgfsys@transformshift{3.238213in}{0.987099in}%
\pgfsys@useobject{currentmarker}{}%
\end{pgfscope}%
\end{pgfscope}%
\begin{pgfscope}%
\definecolor{textcolor}{rgb}{0.000000,0.000000,0.000000}%
\pgfsetstrokecolor{textcolor}%
\pgfsetfillcolor{textcolor}%
\pgftext[x=3.238213in,y=0.973210in,,top]{\color{textcolor}{\rmfamily\fontsize{5.000000}{6.000000}\selectfont\catcode`\^=\active\def^{\ifmmode\sp\else\^{}\fi}\catcode`\%=\active\def%{\%}6.00}}%
\end{pgfscope}%
\begin{pgfscope}%
\pgfpathrectangle{\pgfqpoint{2.259013in}{0.987099in}}{\pgfqpoint{0.979200in}{0.326250in}}%
\pgfusepath{clip}%
\pgfsetbuttcap%
\pgfsetroundjoin%
\pgfsetlinewidth{0.401500pt}%
\definecolor{currentstroke}{rgb}{0.690196,0.690196,0.690196}%
\pgfsetstrokecolor{currentstroke}%
\pgfsetstrokeopacity{0.250000}%
\pgfsetdash{{1.480000pt}{0.640000pt}}{0.000000pt}%
\pgfpathmoveto{\pgfqpoint{2.259013in}{0.987099in}}%
\pgfpathlineto{\pgfqpoint{3.238213in}{0.987099in}}%
\pgfusepath{stroke}%
\end{pgfscope}%
\begin{pgfscope}%
\pgfsetbuttcap%
\pgfsetroundjoin%
\definecolor{currentfill}{rgb}{0.000000,0.000000,0.000000}%
\pgfsetfillcolor{currentfill}%
\pgfsetlinewidth{0.803000pt}%
\definecolor{currentstroke}{rgb}{0.000000,0.000000,0.000000}%
\pgfsetstrokecolor{currentstroke}%
\pgfsetdash{}{0pt}%
\pgfsys@defobject{currentmarker}{\pgfqpoint{0.000000in}{0.000000in}}{\pgfqpoint{0.027778in}{0.000000in}}{%
\pgfpathmoveto{\pgfqpoint{0.000000in}{0.000000in}}%
\pgfpathlineto{\pgfqpoint{0.027778in}{0.000000in}}%
\pgfusepath{stroke,fill}%
}%
\begin{pgfscope}%
\pgfsys@transformshift{2.259013in}{0.987099in}%
\pgfsys@useobject{currentmarker}{}%
\end{pgfscope}%
\end{pgfscope}%
\begin{pgfscope}%
\definecolor{textcolor}{rgb}{0.000000,0.000000,0.000000}%
\pgfsetstrokecolor{textcolor}%
\pgfsetfillcolor{textcolor}%
\pgftext[x=2.042575in, y=0.962986in, left, base]{\color{textcolor}{\rmfamily\fontsize{5.000000}{6.000000}\selectfont\catcode`\^=\active\def^{\ifmmode\sp\else\^{}\fi}\catcode`\%=\active\def%{\%}-2.00}}%
\end{pgfscope}%
\begin{pgfscope}%
\pgfpathrectangle{\pgfqpoint{2.259013in}{0.987099in}}{\pgfqpoint{0.979200in}{0.326250in}}%
\pgfusepath{clip}%
\pgfsetbuttcap%
\pgfsetroundjoin%
\pgfsetlinewidth{0.401500pt}%
\definecolor{currentstroke}{rgb}{0.690196,0.690196,0.690196}%
\pgfsetstrokecolor{currentstroke}%
\pgfsetstrokeopacity{0.250000}%
\pgfsetdash{{1.480000pt}{0.640000pt}}{0.000000pt}%
\pgfpathmoveto{\pgfqpoint{2.259013in}{1.150224in}}%
\pgfpathlineto{\pgfqpoint{3.238213in}{1.150224in}}%
\pgfusepath{stroke}%
\end{pgfscope}%
\begin{pgfscope}%
\pgfsetbuttcap%
\pgfsetroundjoin%
\definecolor{currentfill}{rgb}{0.000000,0.000000,0.000000}%
\pgfsetfillcolor{currentfill}%
\pgfsetlinewidth{0.803000pt}%
\definecolor{currentstroke}{rgb}{0.000000,0.000000,0.000000}%
\pgfsetstrokecolor{currentstroke}%
\pgfsetdash{}{0pt}%
\pgfsys@defobject{currentmarker}{\pgfqpoint{0.000000in}{0.000000in}}{\pgfqpoint{0.027778in}{0.000000in}}{%
\pgfpathmoveto{\pgfqpoint{0.000000in}{0.000000in}}%
\pgfpathlineto{\pgfqpoint{0.027778in}{0.000000in}}%
\pgfusepath{stroke,fill}%
}%
\begin{pgfscope}%
\pgfsys@transformshift{2.259013in}{1.150224in}%
\pgfsys@useobject{currentmarker}{}%
\end{pgfscope}%
\end{pgfscope}%
\begin{pgfscope}%
\definecolor{textcolor}{rgb}{0.000000,0.000000,0.000000}%
\pgfsetstrokecolor{textcolor}%
\pgfsetfillcolor{textcolor}%
\pgftext[x=2.075368in, y=1.126111in, left, base]{\color{textcolor}{\rmfamily\fontsize{5.000000}{6.000000}\selectfont\catcode`\^=\active\def^{\ifmmode\sp\else\^{}\fi}\catcode`\%=\active\def%{\%}0.00}}%
\end{pgfscope}%
\begin{pgfscope}%
\pgfpathrectangle{\pgfqpoint{2.259013in}{0.987099in}}{\pgfqpoint{0.979200in}{0.326250in}}%
\pgfusepath{clip}%
\pgfsetbuttcap%
\pgfsetroundjoin%
\pgfsetlinewidth{0.401500pt}%
\definecolor{currentstroke}{rgb}{0.690196,0.690196,0.690196}%
\pgfsetstrokecolor{currentstroke}%
\pgfsetstrokeopacity{0.250000}%
\pgfsetdash{{1.480000pt}{0.640000pt}}{0.000000pt}%
\pgfpathmoveto{\pgfqpoint{2.259013in}{1.313349in}}%
\pgfpathlineto{\pgfqpoint{3.238213in}{1.313349in}}%
\pgfusepath{stroke}%
\end{pgfscope}%
\begin{pgfscope}%
\pgfsetbuttcap%
\pgfsetroundjoin%
\definecolor{currentfill}{rgb}{0.000000,0.000000,0.000000}%
\pgfsetfillcolor{currentfill}%
\pgfsetlinewidth{0.803000pt}%
\definecolor{currentstroke}{rgb}{0.000000,0.000000,0.000000}%
\pgfsetstrokecolor{currentstroke}%
\pgfsetdash{}{0pt}%
\pgfsys@defobject{currentmarker}{\pgfqpoint{0.000000in}{0.000000in}}{\pgfqpoint{0.027778in}{0.000000in}}{%
\pgfpathmoveto{\pgfqpoint{0.000000in}{0.000000in}}%
\pgfpathlineto{\pgfqpoint{0.027778in}{0.000000in}}%
\pgfusepath{stroke,fill}%
}%
\begin{pgfscope}%
\pgfsys@transformshift{2.259013in}{1.313349in}%
\pgfsys@useobject{currentmarker}{}%
\end{pgfscope}%
\end{pgfscope}%
\begin{pgfscope}%
\definecolor{textcolor}{rgb}{0.000000,0.000000,0.000000}%
\pgfsetstrokecolor{textcolor}%
\pgfsetfillcolor{textcolor}%
\pgftext[x=2.075368in, y=1.289236in, left, base]{\color{textcolor}{\rmfamily\fontsize{5.000000}{6.000000}\selectfont\catcode`\^=\active\def^{\ifmmode\sp\else\^{}\fi}\catcode`\%=\active\def%{\%}2.00}}%
\end{pgfscope}%
\begin{pgfscope}%
\pgfpathrectangle{\pgfqpoint{2.259013in}{0.987099in}}{\pgfqpoint{0.979200in}{0.326250in}}%
\pgfusepath{clip}%
\pgfsetrectcap%
\pgfsetroundjoin%
\pgfsetlinewidth{1.505625pt}%
\definecolor{currentstroke}{rgb}{0.000000,0.619608,0.450980}%
\pgfsetstrokecolor{currentstroke}%
\pgfsetdash{}{0pt}%
\pgfpathmoveto{\pgfqpoint{2.255507in}{1.289040in}}%
\pgfpathlineto{\pgfqpoint{2.268418in}{1.288550in}}%
\pgfpathlineto{\pgfqpoint{2.268872in}{0.980154in}}%
\pgfpathmoveto{\pgfqpoint{2.518335in}{0.980154in}}%
\pgfpathlineto{\pgfqpoint{2.519067in}{1.289943in}}%
\pgfpathlineto{\pgfqpoint{2.540865in}{1.289882in}}%
\pgfpathlineto{\pgfqpoint{2.584929in}{1.289626in}}%
\pgfpathlineto{\pgfqpoint{2.628993in}{1.289553in}}%
\pgfpathlineto{\pgfqpoint{2.663265in}{1.289293in}}%
\pgfpathlineto{\pgfqpoint{2.697537in}{1.289265in}}%
\pgfpathlineto{\pgfqpoint{2.722017in}{1.289015in}}%
\pgfpathlineto{\pgfqpoint{2.751122in}{1.288946in}}%
\pgfpathlineto{\pgfqpoint{2.758018in}{1.288550in}}%
\pgfpathlineto{\pgfqpoint{2.758472in}{0.980154in}}%
\pgfpathmoveto{\pgfqpoint{3.007935in}{0.980154in}}%
\pgfpathlineto{\pgfqpoint{3.008667in}{1.289943in}}%
\pgfpathlineto{\pgfqpoint{3.030464in}{1.289882in}}%
\pgfpathlineto{\pgfqpoint{3.074528in}{1.289626in}}%
\pgfpathlineto{\pgfqpoint{3.118592in}{1.289553in}}%
\pgfpathlineto{\pgfqpoint{3.152864in}{1.289293in}}%
\pgfpathlineto{\pgfqpoint{3.187136in}{1.289265in}}%
\pgfpathlineto{\pgfqpoint{3.211616in}{1.289015in}}%
\pgfpathlineto{\pgfqpoint{3.238213in}{1.288876in}}%
\pgfpathlineto{\pgfqpoint{3.238213in}{1.288876in}}%
\pgfusepath{stroke}%
\end{pgfscope}%
\begin{pgfscope}%
\pgfsetrectcap%
\pgfsetmiterjoin%
\pgfsetlinewidth{0.803000pt}%
\definecolor{currentstroke}{rgb}{0.000000,0.000000,0.000000}%
\pgfsetstrokecolor{currentstroke}%
\pgfsetdash{}{0pt}%
\pgfpathmoveto{\pgfqpoint{2.259013in}{0.987099in}}%
\pgfpathlineto{\pgfqpoint{2.259013in}{1.313349in}}%
\pgfusepath{stroke}%
\end{pgfscope}%
\begin{pgfscope}%
\pgfsetrectcap%
\pgfsetmiterjoin%
\pgfsetlinewidth{0.803000pt}%
\definecolor{currentstroke}{rgb}{0.000000,0.000000,0.000000}%
\pgfsetstrokecolor{currentstroke}%
\pgfsetdash{}{0pt}%
\pgfpathmoveto{\pgfqpoint{3.238213in}{0.987099in}}%
\pgfpathlineto{\pgfqpoint{3.238213in}{1.313349in}}%
\pgfusepath{stroke}%
\end{pgfscope}%
\begin{pgfscope}%
\pgfsetrectcap%
\pgfsetmiterjoin%
\pgfsetlinewidth{0.803000pt}%
\definecolor{currentstroke}{rgb}{0.000000,0.000000,0.000000}%
\pgfsetstrokecolor{currentstroke}%
\pgfsetdash{}{0pt}%
\pgfpathmoveto{\pgfqpoint{2.259013in}{0.987099in}}%
\pgfpathlineto{\pgfqpoint{3.238213in}{0.987099in}}%
\pgfusepath{stroke}%
\end{pgfscope}%
\begin{pgfscope}%
\pgfsetrectcap%
\pgfsetmiterjoin%
\pgfsetlinewidth{0.803000pt}%
\definecolor{currentstroke}{rgb}{0.000000,0.000000,0.000000}%
\pgfsetstrokecolor{currentstroke}%
\pgfsetdash{}{0pt}%
\pgfpathmoveto{\pgfqpoint{2.259013in}{1.313349in}}%
\pgfpathlineto{\pgfqpoint{3.238213in}{1.313349in}}%
\pgfusepath{stroke}%
\end{pgfscope}%
\end{pgfpicture}%
\makeatother%
\endgroup%

	\caption{(a) tensão de saída, (b) corrente do indutor e (c) tensão no diodo do circuito Buck, com diodo não ideal e carga de $5$ A.}
	\label{fig:ex4_1}
\end{figure}

Comparando-se os valores da simulação com os esperados $V_o = 10$ V, $\Delta v_{c,pico} = 0{,}1$ V, $I_L = V_o/R_o = 10/2 = 5$ A e $\Delta i_L = 1$, percebe-se que há diferenças, que podem ser explicadas pela não idealidade do diodo. Enquanto este componente está conduzindo, há uma queda de tensão que não é levada em consideração nos cálculos das tensões e correntes do circuito Buck, como mostrado na Figura~\ref{fig:ex4_1} (c). Reduzindo a queda de tensão direta do diodo para próximo de zero, os valores de interesse do conversor em simulação convergem para os teorizados, como demonstrado na Figura~\ref{fig:ex4_2}.
\begin{figure}[htb!]
	\centering
	%% Creator: Matplotlib, PGF backend
%%
%% To include the figure in your LaTeX document, write
%%   \input{<filename>.pgf}
%%
%% Make sure the required packages are loaded in your preamble
%%   \usepackage{pgf}
%%
%% Also ensure that all the required font packages are loaded; for instance,
%% the lmodern package is sometimes necessary when using math font.
%%   \usepackage{lmodern}
%%
%% Figures using additional raster images can only be included by \input if
%% they are in the same directory as the main LaTeX file. For loading figures
%% from other directories you can use the `import` package
%%   \usepackage{import}
%%
%% and then include the figures with
%%   \import{<path to file>}{<filename>.pgf}
%%
%% Matplotlib used the following preamble
%%   \def\mathdefault#1{#1}
%%   \everymath=\expandafter{\the\everymath\displaystyle}
%%   \IfFileExists{scrextend.sty}{
%%     \usepackage[fontsize=8.000000pt]{scrextend}
%%   }{
%%     \renewcommand{\normalsize}{\fontsize{8.000000}{9.600000}\selectfont}
%%     \normalsize
%%   }
%%   
%%           \usepackage{amsmath}
%%           \usepackage{amssymb}
%%           \usepackage{mathtools}
%%           \usepackage{dsfont}
%%           \providecommand{\mathdefault}[1]{#1}
%%       
%%   \makeatletter\@ifpackageloaded{underscore}{}{\usepackage[strings]{underscore}}\makeatother
%%
\begingroup%
\makeatletter%
\begin{pgfpicture}%
\pgfpathrectangle{\pgfpointorigin}{\pgfqpoint{3.399460in}{3.646235in}}%
\pgfusepath{use as bounding box, clip}%
\begin{pgfscope}%
\pgfsetbuttcap%
\pgfsetmiterjoin%
\definecolor{currentfill}{rgb}{1.000000,1.000000,1.000000}%
\pgfsetfillcolor{currentfill}%
\pgfsetlinewidth{0.000000pt}%
\definecolor{currentstroke}{rgb}{1.000000,1.000000,1.000000}%
\pgfsetstrokecolor{currentstroke}%
\pgfsetdash{}{0pt}%
\pgfpathmoveto{\pgfqpoint{0.000000in}{0.000000in}}%
\pgfpathlineto{\pgfqpoint{3.399460in}{0.000000in}}%
\pgfpathlineto{\pgfqpoint{3.399460in}{3.646235in}}%
\pgfpathlineto{\pgfqpoint{0.000000in}{3.646235in}}%
\pgfpathlineto{\pgfqpoint{0.000000in}{0.000000in}}%
\pgfpathclose%
\pgfusepath{fill}%
\end{pgfscope}%
\begin{pgfscope}%
\pgfsetbuttcap%
\pgfsetmiterjoin%
\definecolor{currentfill}{rgb}{1.000000,1.000000,1.000000}%
\pgfsetfillcolor{currentfill}%
\pgfsetlinewidth{0.000000pt}%
\definecolor{currentstroke}{rgb}{0.000000,0.000000,0.000000}%
\pgfsetstrokecolor{currentstroke}%
\pgfsetstrokeopacity{0.000000}%
\pgfsetdash{}{0pt}%
\pgfpathmoveto{\pgfqpoint{0.573768in}{2.601404in}}%
\pgfpathlineto{\pgfqpoint{3.293768in}{2.601404in}}%
\pgfpathlineto{\pgfqpoint{3.293768in}{3.507654in}}%
\pgfpathlineto{\pgfqpoint{0.573768in}{3.507654in}}%
\pgfpathlineto{\pgfqpoint{0.573768in}{2.601404in}}%
\pgfpathclose%
\pgfusepath{fill}%
\end{pgfscope}%
\begin{pgfscope}%
\pgfpathrectangle{\pgfqpoint{0.573768in}{2.601404in}}{\pgfqpoint{2.720000in}{0.906250in}}%
\pgfusepath{clip}%
\pgfsetbuttcap%
\pgfsetroundjoin%
\pgfsetlinewidth{0.501875pt}%
\definecolor{currentstroke}{rgb}{0.690196,0.690196,0.690196}%
\pgfsetstrokecolor{currentstroke}%
\pgfsetstrokeopacity{0.250000}%
\pgfsetdash{{1.850000pt}{0.800000pt}}{0.000000pt}%
\pgfpathmoveto{\pgfqpoint{0.697405in}{2.601404in}}%
\pgfpathlineto{\pgfqpoint{0.697405in}{3.507654in}}%
\pgfusepath{stroke}%
\end{pgfscope}%
\begin{pgfscope}%
\pgfsetbuttcap%
\pgfsetroundjoin%
\definecolor{currentfill}{rgb}{0.000000,0.000000,0.000000}%
\pgfsetfillcolor{currentfill}%
\pgfsetlinewidth{0.803000pt}%
\definecolor{currentstroke}{rgb}{0.000000,0.000000,0.000000}%
\pgfsetstrokecolor{currentstroke}%
\pgfsetdash{}{0pt}%
\pgfsys@defobject{currentmarker}{\pgfqpoint{0.000000in}{-0.048611in}}{\pgfqpoint{0.000000in}{0.000000in}}{%
\pgfpathmoveto{\pgfqpoint{0.000000in}{0.000000in}}%
\pgfpathlineto{\pgfqpoint{0.000000in}{-0.048611in}}%
\pgfusepath{stroke,fill}%
}%
\begin{pgfscope}%
\pgfsys@transformshift{0.697405in}{2.601404in}%
\pgfsys@useobject{currentmarker}{}%
\end{pgfscope}%
\end{pgfscope}%
\begin{pgfscope}%
\pgfpathrectangle{\pgfqpoint{0.573768in}{2.601404in}}{\pgfqpoint{2.720000in}{0.906250in}}%
\pgfusepath{clip}%
\pgfsetbuttcap%
\pgfsetroundjoin%
\pgfsetlinewidth{0.501875pt}%
\definecolor{currentstroke}{rgb}{0.690196,0.690196,0.690196}%
\pgfsetstrokecolor{currentstroke}%
\pgfsetstrokeopacity{0.250000}%
\pgfsetdash{{1.850000pt}{0.800000pt}}{0.000000pt}%
\pgfpathmoveto{\pgfqpoint{1.109526in}{2.601404in}}%
\pgfpathlineto{\pgfqpoint{1.109526in}{3.507654in}}%
\pgfusepath{stroke}%
\end{pgfscope}%
\begin{pgfscope}%
\pgfsetbuttcap%
\pgfsetroundjoin%
\definecolor{currentfill}{rgb}{0.000000,0.000000,0.000000}%
\pgfsetfillcolor{currentfill}%
\pgfsetlinewidth{0.803000pt}%
\definecolor{currentstroke}{rgb}{0.000000,0.000000,0.000000}%
\pgfsetstrokecolor{currentstroke}%
\pgfsetdash{}{0pt}%
\pgfsys@defobject{currentmarker}{\pgfqpoint{0.000000in}{-0.048611in}}{\pgfqpoint{0.000000in}{0.000000in}}{%
\pgfpathmoveto{\pgfqpoint{0.000000in}{0.000000in}}%
\pgfpathlineto{\pgfqpoint{0.000000in}{-0.048611in}}%
\pgfusepath{stroke,fill}%
}%
\begin{pgfscope}%
\pgfsys@transformshift{1.109526in}{2.601404in}%
\pgfsys@useobject{currentmarker}{}%
\end{pgfscope}%
\end{pgfscope}%
\begin{pgfscope}%
\pgfpathrectangle{\pgfqpoint{0.573768in}{2.601404in}}{\pgfqpoint{2.720000in}{0.906250in}}%
\pgfusepath{clip}%
\pgfsetbuttcap%
\pgfsetroundjoin%
\pgfsetlinewidth{0.501875pt}%
\definecolor{currentstroke}{rgb}{0.690196,0.690196,0.690196}%
\pgfsetstrokecolor{currentstroke}%
\pgfsetstrokeopacity{0.250000}%
\pgfsetdash{{1.850000pt}{0.800000pt}}{0.000000pt}%
\pgfpathmoveto{\pgfqpoint{1.521647in}{2.601404in}}%
\pgfpathlineto{\pgfqpoint{1.521647in}{3.507654in}}%
\pgfusepath{stroke}%
\end{pgfscope}%
\begin{pgfscope}%
\pgfsetbuttcap%
\pgfsetroundjoin%
\definecolor{currentfill}{rgb}{0.000000,0.000000,0.000000}%
\pgfsetfillcolor{currentfill}%
\pgfsetlinewidth{0.803000pt}%
\definecolor{currentstroke}{rgb}{0.000000,0.000000,0.000000}%
\pgfsetstrokecolor{currentstroke}%
\pgfsetdash{}{0pt}%
\pgfsys@defobject{currentmarker}{\pgfqpoint{0.000000in}{-0.048611in}}{\pgfqpoint{0.000000in}{0.000000in}}{%
\pgfpathmoveto{\pgfqpoint{0.000000in}{0.000000in}}%
\pgfpathlineto{\pgfqpoint{0.000000in}{-0.048611in}}%
\pgfusepath{stroke,fill}%
}%
\begin{pgfscope}%
\pgfsys@transformshift{1.521647in}{2.601404in}%
\pgfsys@useobject{currentmarker}{}%
\end{pgfscope}%
\end{pgfscope}%
\begin{pgfscope}%
\pgfpathrectangle{\pgfqpoint{0.573768in}{2.601404in}}{\pgfqpoint{2.720000in}{0.906250in}}%
\pgfusepath{clip}%
\pgfsetbuttcap%
\pgfsetroundjoin%
\pgfsetlinewidth{0.501875pt}%
\definecolor{currentstroke}{rgb}{0.690196,0.690196,0.690196}%
\pgfsetstrokecolor{currentstroke}%
\pgfsetstrokeopacity{0.250000}%
\pgfsetdash{{1.850000pt}{0.800000pt}}{0.000000pt}%
\pgfpathmoveto{\pgfqpoint{1.933768in}{2.601404in}}%
\pgfpathlineto{\pgfqpoint{1.933768in}{3.507654in}}%
\pgfusepath{stroke}%
\end{pgfscope}%
\begin{pgfscope}%
\pgfsetbuttcap%
\pgfsetroundjoin%
\definecolor{currentfill}{rgb}{0.000000,0.000000,0.000000}%
\pgfsetfillcolor{currentfill}%
\pgfsetlinewidth{0.803000pt}%
\definecolor{currentstroke}{rgb}{0.000000,0.000000,0.000000}%
\pgfsetstrokecolor{currentstroke}%
\pgfsetdash{}{0pt}%
\pgfsys@defobject{currentmarker}{\pgfqpoint{0.000000in}{-0.048611in}}{\pgfqpoint{0.000000in}{0.000000in}}{%
\pgfpathmoveto{\pgfqpoint{0.000000in}{0.000000in}}%
\pgfpathlineto{\pgfqpoint{0.000000in}{-0.048611in}}%
\pgfusepath{stroke,fill}%
}%
\begin{pgfscope}%
\pgfsys@transformshift{1.933768in}{2.601404in}%
\pgfsys@useobject{currentmarker}{}%
\end{pgfscope}%
\end{pgfscope}%
\begin{pgfscope}%
\pgfpathrectangle{\pgfqpoint{0.573768in}{2.601404in}}{\pgfqpoint{2.720000in}{0.906250in}}%
\pgfusepath{clip}%
\pgfsetbuttcap%
\pgfsetroundjoin%
\pgfsetlinewidth{0.501875pt}%
\definecolor{currentstroke}{rgb}{0.690196,0.690196,0.690196}%
\pgfsetstrokecolor{currentstroke}%
\pgfsetstrokeopacity{0.250000}%
\pgfsetdash{{1.850000pt}{0.800000pt}}{0.000000pt}%
\pgfpathmoveto{\pgfqpoint{2.345890in}{2.601404in}}%
\pgfpathlineto{\pgfqpoint{2.345890in}{3.507654in}}%
\pgfusepath{stroke}%
\end{pgfscope}%
\begin{pgfscope}%
\pgfsetbuttcap%
\pgfsetroundjoin%
\definecolor{currentfill}{rgb}{0.000000,0.000000,0.000000}%
\pgfsetfillcolor{currentfill}%
\pgfsetlinewidth{0.803000pt}%
\definecolor{currentstroke}{rgb}{0.000000,0.000000,0.000000}%
\pgfsetstrokecolor{currentstroke}%
\pgfsetdash{}{0pt}%
\pgfsys@defobject{currentmarker}{\pgfqpoint{0.000000in}{-0.048611in}}{\pgfqpoint{0.000000in}{0.000000in}}{%
\pgfpathmoveto{\pgfqpoint{0.000000in}{0.000000in}}%
\pgfpathlineto{\pgfqpoint{0.000000in}{-0.048611in}}%
\pgfusepath{stroke,fill}%
}%
\begin{pgfscope}%
\pgfsys@transformshift{2.345890in}{2.601404in}%
\pgfsys@useobject{currentmarker}{}%
\end{pgfscope}%
\end{pgfscope}%
\begin{pgfscope}%
\pgfpathrectangle{\pgfqpoint{0.573768in}{2.601404in}}{\pgfqpoint{2.720000in}{0.906250in}}%
\pgfusepath{clip}%
\pgfsetbuttcap%
\pgfsetroundjoin%
\pgfsetlinewidth{0.501875pt}%
\definecolor{currentstroke}{rgb}{0.690196,0.690196,0.690196}%
\pgfsetstrokecolor{currentstroke}%
\pgfsetstrokeopacity{0.250000}%
\pgfsetdash{{1.850000pt}{0.800000pt}}{0.000000pt}%
\pgfpathmoveto{\pgfqpoint{2.758011in}{2.601404in}}%
\pgfpathlineto{\pgfqpoint{2.758011in}{3.507654in}}%
\pgfusepath{stroke}%
\end{pgfscope}%
\begin{pgfscope}%
\pgfsetbuttcap%
\pgfsetroundjoin%
\definecolor{currentfill}{rgb}{0.000000,0.000000,0.000000}%
\pgfsetfillcolor{currentfill}%
\pgfsetlinewidth{0.803000pt}%
\definecolor{currentstroke}{rgb}{0.000000,0.000000,0.000000}%
\pgfsetstrokecolor{currentstroke}%
\pgfsetdash{}{0pt}%
\pgfsys@defobject{currentmarker}{\pgfqpoint{0.000000in}{-0.048611in}}{\pgfqpoint{0.000000in}{0.000000in}}{%
\pgfpathmoveto{\pgfqpoint{0.000000in}{0.000000in}}%
\pgfpathlineto{\pgfqpoint{0.000000in}{-0.048611in}}%
\pgfusepath{stroke,fill}%
}%
\begin{pgfscope}%
\pgfsys@transformshift{2.758011in}{2.601404in}%
\pgfsys@useobject{currentmarker}{}%
\end{pgfscope}%
\end{pgfscope}%
\begin{pgfscope}%
\pgfpathrectangle{\pgfqpoint{0.573768in}{2.601404in}}{\pgfqpoint{2.720000in}{0.906250in}}%
\pgfusepath{clip}%
\pgfsetbuttcap%
\pgfsetroundjoin%
\pgfsetlinewidth{0.501875pt}%
\definecolor{currentstroke}{rgb}{0.690196,0.690196,0.690196}%
\pgfsetstrokecolor{currentstroke}%
\pgfsetstrokeopacity{0.250000}%
\pgfsetdash{{1.850000pt}{0.800000pt}}{0.000000pt}%
\pgfpathmoveto{\pgfqpoint{3.170132in}{2.601404in}}%
\pgfpathlineto{\pgfqpoint{3.170132in}{3.507654in}}%
\pgfusepath{stroke}%
\end{pgfscope}%
\begin{pgfscope}%
\pgfsetbuttcap%
\pgfsetroundjoin%
\definecolor{currentfill}{rgb}{0.000000,0.000000,0.000000}%
\pgfsetfillcolor{currentfill}%
\pgfsetlinewidth{0.803000pt}%
\definecolor{currentstroke}{rgb}{0.000000,0.000000,0.000000}%
\pgfsetstrokecolor{currentstroke}%
\pgfsetdash{}{0pt}%
\pgfsys@defobject{currentmarker}{\pgfqpoint{0.000000in}{-0.048611in}}{\pgfqpoint{0.000000in}{0.000000in}}{%
\pgfpathmoveto{\pgfqpoint{0.000000in}{0.000000in}}%
\pgfpathlineto{\pgfqpoint{0.000000in}{-0.048611in}}%
\pgfusepath{stroke,fill}%
}%
\begin{pgfscope}%
\pgfsys@transformshift{3.170132in}{2.601404in}%
\pgfsys@useobject{currentmarker}{}%
\end{pgfscope}%
\end{pgfscope}%
\begin{pgfscope}%
\pgfpathrectangle{\pgfqpoint{0.573768in}{2.601404in}}{\pgfqpoint{2.720000in}{0.906250in}}%
\pgfusepath{clip}%
\pgfsetbuttcap%
\pgfsetroundjoin%
\pgfsetlinewidth{0.501875pt}%
\definecolor{currentstroke}{rgb}{0.690196,0.690196,0.690196}%
\pgfsetstrokecolor{currentstroke}%
\pgfsetstrokeopacity{0.250000}%
\pgfsetdash{{1.850000pt}{0.800000pt}}{0.000000pt}%
\pgfpathmoveto{\pgfqpoint{0.573768in}{2.601404in}}%
\pgfpathlineto{\pgfqpoint{3.293768in}{2.601404in}}%
\pgfusepath{stroke}%
\end{pgfscope}%
\begin{pgfscope}%
\pgfsetbuttcap%
\pgfsetroundjoin%
\definecolor{currentfill}{rgb}{0.000000,0.000000,0.000000}%
\pgfsetfillcolor{currentfill}%
\pgfsetlinewidth{0.803000pt}%
\definecolor{currentstroke}{rgb}{0.000000,0.000000,0.000000}%
\pgfsetstrokecolor{currentstroke}%
\pgfsetdash{}{0pt}%
\pgfsys@defobject{currentmarker}{\pgfqpoint{-0.048611in}{0.000000in}}{\pgfqpoint{-0.000000in}{0.000000in}}{%
\pgfpathmoveto{\pgfqpoint{-0.000000in}{0.000000in}}%
\pgfpathlineto{\pgfqpoint{-0.048611in}{0.000000in}}%
\pgfusepath{stroke,fill}%
}%
\begin{pgfscope}%
\pgfsys@transformshift{0.573768in}{2.601404in}%
\pgfsys@useobject{currentmarker}{}%
\end{pgfscope}%
\end{pgfscope}%
\begin{pgfscope}%
\definecolor{textcolor}{rgb}{0.000000,0.000000,0.000000}%
\pgfsetstrokecolor{textcolor}%
\pgfsetfillcolor{textcolor}%
\pgftext[x=0.417518in, y=2.562824in, left, base]{\color{textcolor}{\rmfamily\fontsize{8.000000}{9.600000}\selectfont\catcode`\^=\active\def^{\ifmmode\sp\else\^{}\fi}\catcode`\%=\active\def%{\%}$\mathdefault{0}$}}%
\end{pgfscope}%
\begin{pgfscope}%
\pgfpathrectangle{\pgfqpoint{0.573768in}{2.601404in}}{\pgfqpoint{2.720000in}{0.906250in}}%
\pgfusepath{clip}%
\pgfsetbuttcap%
\pgfsetroundjoin%
\pgfsetlinewidth{0.501875pt}%
\definecolor{currentstroke}{rgb}{0.690196,0.690196,0.690196}%
\pgfsetstrokecolor{currentstroke}%
\pgfsetstrokeopacity{0.250000}%
\pgfsetdash{{1.850000pt}{0.800000pt}}{0.000000pt}%
\pgfpathmoveto{\pgfqpoint{0.573768in}{2.903488in}}%
\pgfpathlineto{\pgfqpoint{3.293768in}{2.903488in}}%
\pgfusepath{stroke}%
\end{pgfscope}%
\begin{pgfscope}%
\pgfsetbuttcap%
\pgfsetroundjoin%
\definecolor{currentfill}{rgb}{0.000000,0.000000,0.000000}%
\pgfsetfillcolor{currentfill}%
\pgfsetlinewidth{0.803000pt}%
\definecolor{currentstroke}{rgb}{0.000000,0.000000,0.000000}%
\pgfsetstrokecolor{currentstroke}%
\pgfsetdash{}{0pt}%
\pgfsys@defobject{currentmarker}{\pgfqpoint{-0.048611in}{0.000000in}}{\pgfqpoint{-0.000000in}{0.000000in}}{%
\pgfpathmoveto{\pgfqpoint{-0.000000in}{0.000000in}}%
\pgfpathlineto{\pgfqpoint{-0.048611in}{0.000000in}}%
\pgfusepath{stroke,fill}%
}%
\begin{pgfscope}%
\pgfsys@transformshift{0.573768in}{2.903488in}%
\pgfsys@useobject{currentmarker}{}%
\end{pgfscope}%
\end{pgfscope}%
\begin{pgfscope}%
\definecolor{textcolor}{rgb}{0.000000,0.000000,0.000000}%
\pgfsetstrokecolor{textcolor}%
\pgfsetfillcolor{textcolor}%
\pgftext[x=0.417518in, y=2.864908in, left, base]{\color{textcolor}{\rmfamily\fontsize{8.000000}{9.600000}\selectfont\catcode`\^=\active\def^{\ifmmode\sp\else\^{}\fi}\catcode`\%=\active\def%{\%}$\mathdefault{5}$}}%
\end{pgfscope}%
\begin{pgfscope}%
\pgfpathrectangle{\pgfqpoint{0.573768in}{2.601404in}}{\pgfqpoint{2.720000in}{0.906250in}}%
\pgfusepath{clip}%
\pgfsetbuttcap%
\pgfsetroundjoin%
\pgfsetlinewidth{0.501875pt}%
\definecolor{currentstroke}{rgb}{0.690196,0.690196,0.690196}%
\pgfsetstrokecolor{currentstroke}%
\pgfsetstrokeopacity{0.250000}%
\pgfsetdash{{1.850000pt}{0.800000pt}}{0.000000pt}%
\pgfpathmoveto{\pgfqpoint{0.573768in}{3.205571in}}%
\pgfpathlineto{\pgfqpoint{3.293768in}{3.205571in}}%
\pgfusepath{stroke}%
\end{pgfscope}%
\begin{pgfscope}%
\pgfsetbuttcap%
\pgfsetroundjoin%
\definecolor{currentfill}{rgb}{0.000000,0.000000,0.000000}%
\pgfsetfillcolor{currentfill}%
\pgfsetlinewidth{0.803000pt}%
\definecolor{currentstroke}{rgb}{0.000000,0.000000,0.000000}%
\pgfsetstrokecolor{currentstroke}%
\pgfsetdash{}{0pt}%
\pgfsys@defobject{currentmarker}{\pgfqpoint{-0.048611in}{0.000000in}}{\pgfqpoint{-0.000000in}{0.000000in}}{%
\pgfpathmoveto{\pgfqpoint{-0.000000in}{0.000000in}}%
\pgfpathlineto{\pgfqpoint{-0.048611in}{0.000000in}}%
\pgfusepath{stroke,fill}%
}%
\begin{pgfscope}%
\pgfsys@transformshift{0.573768in}{3.205571in}%
\pgfsys@useobject{currentmarker}{}%
\end{pgfscope}%
\end{pgfscope}%
\begin{pgfscope}%
\definecolor{textcolor}{rgb}{0.000000,0.000000,0.000000}%
\pgfsetstrokecolor{textcolor}%
\pgfsetfillcolor{textcolor}%
\pgftext[x=0.358489in, y=3.166991in, left, base]{\color{textcolor}{\rmfamily\fontsize{8.000000}{9.600000}\selectfont\catcode`\^=\active\def^{\ifmmode\sp\else\^{}\fi}\catcode`\%=\active\def%{\%}$\mathdefault{10}$}}%
\end{pgfscope}%
\begin{pgfscope}%
\pgfpathrectangle{\pgfqpoint{0.573768in}{2.601404in}}{\pgfqpoint{2.720000in}{0.906250in}}%
\pgfusepath{clip}%
\pgfsetbuttcap%
\pgfsetroundjoin%
\pgfsetlinewidth{0.501875pt}%
\definecolor{currentstroke}{rgb}{0.690196,0.690196,0.690196}%
\pgfsetstrokecolor{currentstroke}%
\pgfsetstrokeopacity{0.250000}%
\pgfsetdash{{1.850000pt}{0.800000pt}}{0.000000pt}%
\pgfpathmoveto{\pgfqpoint{0.573768in}{3.507654in}}%
\pgfpathlineto{\pgfqpoint{3.293768in}{3.507654in}}%
\pgfusepath{stroke}%
\end{pgfscope}%
\begin{pgfscope}%
\pgfsetbuttcap%
\pgfsetroundjoin%
\definecolor{currentfill}{rgb}{0.000000,0.000000,0.000000}%
\pgfsetfillcolor{currentfill}%
\pgfsetlinewidth{0.803000pt}%
\definecolor{currentstroke}{rgb}{0.000000,0.000000,0.000000}%
\pgfsetstrokecolor{currentstroke}%
\pgfsetdash{}{0pt}%
\pgfsys@defobject{currentmarker}{\pgfqpoint{-0.048611in}{0.000000in}}{\pgfqpoint{-0.000000in}{0.000000in}}{%
\pgfpathmoveto{\pgfqpoint{-0.000000in}{0.000000in}}%
\pgfpathlineto{\pgfqpoint{-0.048611in}{0.000000in}}%
\pgfusepath{stroke,fill}%
}%
\begin{pgfscope}%
\pgfsys@transformshift{0.573768in}{3.507654in}%
\pgfsys@useobject{currentmarker}{}%
\end{pgfscope}%
\end{pgfscope}%
\begin{pgfscope}%
\definecolor{textcolor}{rgb}{0.000000,0.000000,0.000000}%
\pgfsetstrokecolor{textcolor}%
\pgfsetfillcolor{textcolor}%
\pgftext[x=0.358489in, y=3.469074in, left, base]{\color{textcolor}{\rmfamily\fontsize{8.000000}{9.600000}\selectfont\catcode`\^=\active\def^{\ifmmode\sp\else\^{}\fi}\catcode`\%=\active\def%{\%}$\mathdefault{15}$}}%
\end{pgfscope}%
\begin{pgfscope}%
\definecolor{textcolor}{rgb}{0.000000,0.000000,0.000000}%
\pgfsetstrokecolor{textcolor}%
\pgfsetfillcolor{textcolor}%
\pgftext[x=0.302933in,y=3.054529in,,bottom,rotate=90.000000]{\color{textcolor}{\rmfamily\fontsize{8.000000}{9.600000}\selectfont\catcode`\^=\active\def^{\ifmmode\sp\else\^{}\fi}\catcode`\%=\active\def%{\%}$v_o$ (V)}}%
\end{pgfscope}%
\begin{pgfscope}%
\pgfpathrectangle{\pgfqpoint{0.573768in}{2.601404in}}{\pgfqpoint{2.720000in}{0.906250in}}%
\pgfusepath{clip}%
\pgfsetrectcap%
\pgfsetroundjoin%
\pgfsetlinewidth{1.505625pt}%
\definecolor{currentstroke}{rgb}{0.000000,0.447059,0.698039}%
\pgfsetstrokecolor{currentstroke}%
\pgfsetdash{}{0pt}%
\pgfpathmoveto{\pgfqpoint{0.697405in}{2.601404in}}%
\pgfpathlineto{\pgfqpoint{0.699185in}{2.602124in}}%
\pgfpathlineto{\pgfqpoint{0.700999in}{2.605087in}}%
\pgfpathlineto{\pgfqpoint{0.703472in}{2.612728in}}%
\pgfpathlineto{\pgfqpoint{0.706688in}{2.628518in}}%
\pgfpathlineto{\pgfqpoint{0.720436in}{2.708563in}}%
\pgfpathlineto{\pgfqpoint{0.735930in}{2.845398in}}%
\pgfpathlineto{\pgfqpoint{0.775591in}{3.245167in}}%
\pgfpathlineto{\pgfqpoint{0.780969in}{3.285231in}}%
\pgfpathlineto{\pgfqpoint{0.788003in}{3.337161in}}%
\pgfpathlineto{\pgfqpoint{0.793439in}{3.374904in}}%
\pgfpathlineto{\pgfqpoint{0.797614in}{3.393598in}}%
\pgfpathlineto{\pgfqpoint{0.803384in}{3.419640in}}%
\pgfpathlineto{\pgfqpoint{0.809266in}{3.445449in}}%
\pgfpathlineto{\pgfqpoint{0.812233in}{3.451783in}}%
\pgfpathlineto{\pgfqpoint{0.820009in}{3.463135in}}%
\pgfpathlineto{\pgfqpoint{0.824591in}{3.471288in}}%
\pgfpathlineto{\pgfqpoint{0.826405in}{3.471722in}}%
\pgfpathlineto{\pgfqpoint{0.828218in}{3.470454in}}%
\pgfpathlineto{\pgfqpoint{0.830950in}{3.466016in}}%
\pgfpathlineto{\pgfqpoint{0.834576in}{3.461667in}}%
\pgfpathlineto{\pgfqpoint{0.837698in}{3.460181in}}%
\pgfpathlineto{\pgfqpoint{0.839768in}{3.459139in}}%
\pgfpathlineto{\pgfqpoint{0.841911in}{3.455828in}}%
\pgfpathlineto{\pgfqpoint{0.844713in}{3.448143in}}%
\pgfpathlineto{\pgfqpoint{0.853935in}{3.419133in}}%
\pgfpathlineto{\pgfqpoint{0.857407in}{3.410475in}}%
\pgfpathlineto{\pgfqpoint{0.860869in}{3.396603in}}%
\pgfpathlineto{\pgfqpoint{0.870438in}{3.354450in}}%
\pgfpathlineto{\pgfqpoint{0.874648in}{3.339378in}}%
\pgfpathlineto{\pgfqpoint{0.879005in}{3.316625in}}%
\pgfpathlineto{\pgfqpoint{0.884258in}{3.289623in}}%
\pgfpathlineto{\pgfqpoint{0.890069in}{3.268710in}}%
\pgfpathlineto{\pgfqpoint{0.894355in}{3.248308in}}%
\pgfpathlineto{\pgfqpoint{0.901227in}{3.214672in}}%
\pgfpathlineto{\pgfqpoint{0.906686in}{3.198034in}}%
\pgfpathlineto{\pgfqpoint{0.910972in}{3.180354in}}%
\pgfpathlineto{\pgfqpoint{0.917058in}{3.154966in}}%
\pgfpathlineto{\pgfqpoint{0.920463in}{3.147357in}}%
\pgfpathlineto{\pgfqpoint{0.923501in}{3.141127in}}%
\pgfpathlineto{\pgfqpoint{0.927292in}{3.129555in}}%
\pgfpathlineto{\pgfqpoint{0.932993in}{3.111763in}}%
\pgfpathlineto{\pgfqpoint{0.935796in}{3.107952in}}%
\pgfpathlineto{\pgfqpoint{0.941306in}{3.103076in}}%
\pgfpathlineto{\pgfqpoint{0.944558in}{3.096365in}}%
\pgfpathlineto{\pgfqpoint{0.948226in}{3.088986in}}%
\pgfpathlineto{\pgfqpoint{0.950534in}{3.087329in}}%
\pgfpathlineto{\pgfqpoint{0.952512in}{3.087723in}}%
\pgfpathlineto{\pgfqpoint{0.957050in}{3.089829in}}%
\pgfpathlineto{\pgfqpoint{0.959357in}{3.088686in}}%
\pgfpathlineto{\pgfqpoint{0.965810in}{3.084081in}}%
\pgfpathlineto{\pgfqpoint{0.967788in}{3.085723in}}%
\pgfpathlineto{\pgfqpoint{0.970446in}{3.090313in}}%
\pgfpathlineto{\pgfqpoint{0.973578in}{3.094613in}}%
\pgfpathlineto{\pgfqpoint{0.976216in}{3.095881in}}%
\pgfpathlineto{\pgfqpoint{0.981773in}{3.096800in}}%
\pgfpathlineto{\pgfqpoint{0.984246in}{3.100768in}}%
\pgfpathlineto{\pgfqpoint{0.993078in}{3.117874in}}%
\pgfpathlineto{\pgfqpoint{0.998197in}{3.121967in}}%
\pgfpathlineto{\pgfqpoint{1.000999in}{3.128114in}}%
\pgfpathlineto{\pgfqpoint{1.008739in}{3.147331in}}%
\pgfpathlineto{\pgfqpoint{1.011631in}{3.149726in}}%
\pgfpathlineto{\pgfqpoint{1.013939in}{3.152450in}}%
\pgfpathlineto{\pgfqpoint{1.016742in}{3.158377in}}%
\pgfpathlineto{\pgfqpoint{1.026213in}{3.181308in}}%
\pgfpathlineto{\pgfqpoint{1.031501in}{3.186916in}}%
\pgfpathlineto{\pgfqpoint{1.034633in}{3.194468in}}%
\pgfpathlineto{\pgfqpoint{1.040225in}{3.208799in}}%
\pgfpathlineto{\pgfqpoint{1.042863in}{3.211571in}}%
\pgfpathlineto{\pgfqpoint{1.047662in}{3.214850in}}%
\pgfpathlineto{\pgfqpoint{1.050464in}{3.220084in}}%
\pgfpathlineto{\pgfqpoint{1.057205in}{3.234073in}}%
\pgfpathlineto{\pgfqpoint{1.059513in}{3.235116in}}%
\pgfpathlineto{\pgfqpoint{1.063777in}{3.235684in}}%
\pgfpathlineto{\pgfqpoint{1.066250in}{3.238584in}}%
\pgfpathlineto{\pgfqpoint{1.069733in}{3.245584in}}%
\pgfpathlineto{\pgfqpoint{1.072535in}{3.249409in}}%
\pgfpathlineto{\pgfqpoint{1.074679in}{3.250215in}}%
\pgfpathlineto{\pgfqpoint{1.076687in}{3.249346in}}%
\pgfpathlineto{\pgfqpoint{1.079521in}{3.247916in}}%
\pgfpathlineto{\pgfqpoint{1.081664in}{3.248586in}}%
\pgfpathlineto{\pgfqpoint{1.084137in}{3.251192in}}%
\pgfpathlineto{\pgfqpoint{1.088690in}{3.256663in}}%
\pgfpathlineto{\pgfqpoint{1.090669in}{3.256631in}}%
\pgfpathlineto{\pgfqpoint{1.092812in}{3.254880in}}%
\pgfpathlineto{\pgfqpoint{1.096781in}{3.251118in}}%
\pgfpathlineto{\pgfqpoint{1.098924in}{3.251205in}}%
\pgfpathlineto{\pgfqpoint{1.101196in}{3.252919in}}%
\pgfpathlineto{\pgfqpoint{1.104515in}{3.255532in}}%
\pgfpathlineto{\pgfqpoint{1.106494in}{3.255068in}}%
\pgfpathlineto{\pgfqpoint{1.108637in}{3.252866in}}%
\pgfpathlineto{\pgfqpoint{1.114459in}{3.245709in}}%
\pgfpathlineto{\pgfqpoint{1.116602in}{3.245764in}}%
\pgfpathlineto{\pgfqpoint{1.122186in}{3.247268in}}%
\pgfpathlineto{\pgfqpoint{1.124493in}{3.244769in}}%
\pgfpathlineto{\pgfqpoint{1.131949in}{3.234525in}}%
\pgfpathlineto{\pgfqpoint{1.134193in}{3.234835in}}%
\pgfpathlineto{\pgfqpoint{1.137155in}{3.235620in}}%
\pgfpathlineto{\pgfqpoint{1.139298in}{3.234178in}}%
\pgfpathlineto{\pgfqpoint{1.141771in}{3.230422in}}%
\pgfpathlineto{\pgfqpoint{1.147132in}{3.221406in}}%
\pgfpathlineto{\pgfqpoint{1.149440in}{3.220519in}}%
\pgfpathlineto{\pgfqpoint{1.151801in}{3.221372in}}%
\pgfpathlineto{\pgfqpoint{1.153944in}{3.221304in}}%
\pgfpathlineto{\pgfqpoint{1.156087in}{3.219547in}}%
\pgfpathlineto{\pgfqpoint{1.158725in}{3.215140in}}%
\pgfpathlineto{\pgfqpoint{1.162998in}{3.207687in}}%
\pgfpathlineto{\pgfqpoint{1.165471in}{3.206363in}}%
\pgfpathlineto{\pgfqpoint{1.167576in}{3.206952in}}%
\pgfpathlineto{\pgfqpoint{1.169795in}{3.207601in}}%
\pgfpathlineto{\pgfqpoint{1.171938in}{3.206532in}}%
\pgfpathlineto{\pgfqpoint{1.174411in}{3.203287in}}%
\pgfpathlineto{\pgfqpoint{1.180332in}{3.194328in}}%
\pgfpathlineto{\pgfqpoint{1.182475in}{3.194043in}}%
\pgfpathlineto{\pgfqpoint{1.184961in}{3.195531in}}%
\pgfpathlineto{\pgfqpoint{1.187104in}{3.195870in}}%
\pgfpathlineto{\pgfqpoint{1.189247in}{3.194579in}}%
\pgfpathlineto{\pgfqpoint{1.191835in}{3.190912in}}%
\pgfpathlineto{\pgfqpoint{1.195576in}{3.185461in}}%
\pgfpathlineto{\pgfqpoint{1.197884in}{3.184686in}}%
\pgfpathlineto{\pgfqpoint{1.200027in}{3.185714in}}%
\pgfpathlineto{\pgfqpoint{1.203589in}{3.187898in}}%
\pgfpathlineto{\pgfqpoint{1.205732in}{3.187119in}}%
\pgfpathlineto{\pgfqpoint{1.208205in}{3.184262in}}%
\pgfpathlineto{\pgfqpoint{1.211918in}{3.179644in}}%
\pgfpathlineto{\pgfqpoint{1.214061in}{3.179287in}}%
\pgfpathlineto{\pgfqpoint{1.216204in}{3.180618in}}%
\pgfpathlineto{\pgfqpoint{1.220550in}{3.184105in}}%
\pgfpathlineto{\pgfqpoint{1.222693in}{3.183501in}}%
\pgfpathlineto{\pgfqpoint{1.225051in}{3.181023in}}%
\pgfpathlineto{\pgfqpoint{1.228019in}{3.177937in}}%
\pgfpathlineto{\pgfqpoint{1.230162in}{3.177766in}}%
\pgfpathlineto{\pgfqpoint{1.232305in}{3.179283in}}%
\pgfpathlineto{\pgfqpoint{1.237713in}{3.184298in}}%
\pgfpathlineto{\pgfqpoint{1.239856in}{3.183643in}}%
\pgfpathlineto{\pgfqpoint{1.242419in}{3.180947in}}%
\pgfpathlineto{\pgfqpoint{1.244892in}{3.179443in}}%
\pgfpathlineto{\pgfqpoint{1.246870in}{3.179881in}}%
\pgfpathlineto{\pgfqpoint{1.249178in}{3.182186in}}%
\pgfpathlineto{\pgfqpoint{1.253703in}{3.187432in}}%
\pgfpathlineto{\pgfqpoint{1.255846in}{3.187478in}}%
\pgfpathlineto{\pgfqpoint{1.258119in}{3.185779in}}%
\pgfpathlineto{\pgfqpoint{1.260710in}{3.183851in}}%
\pgfpathlineto{\pgfqpoint{1.262853in}{3.184135in}}%
\pgfpathlineto{\pgfqpoint{1.264996in}{3.186079in}}%
\pgfpathlineto{\pgfqpoint{1.271177in}{3.193028in}}%
\pgfpathlineto{\pgfqpoint{1.273320in}{3.192507in}}%
\pgfpathlineto{\pgfqpoint{1.278809in}{3.189823in}}%
\pgfpathlineto{\pgfqpoint{1.280952in}{3.191469in}}%
\pgfpathlineto{\pgfqpoint{1.283871in}{3.196039in}}%
\pgfpathlineto{\pgfqpoint{1.286673in}{3.198935in}}%
\pgfpathlineto{\pgfqpoint{1.288816in}{3.199207in}}%
\pgfpathlineto{\pgfqpoint{1.290990in}{3.197809in}}%
\pgfpathlineto{\pgfqpoint{1.293799in}{3.195932in}}%
\pgfpathlineto{\pgfqpoint{1.295942in}{3.196413in}}%
\pgfpathlineto{\pgfqpoint{1.298250in}{3.198744in}}%
\pgfpathlineto{\pgfqpoint{1.303653in}{3.205215in}}%
\pgfpathlineto{\pgfqpoint{1.305796in}{3.205016in}}%
\pgfpathlineto{\pgfqpoint{1.308140in}{3.202986in}}%
\pgfpathlineto{\pgfqpoint{1.310613in}{3.201632in}}%
\pgfpathlineto{\pgfqpoint{1.312756in}{3.202234in}}%
\pgfpathlineto{\pgfqpoint{1.315063in}{3.204674in}}%
\pgfpathlineto{\pgfqpoint{1.319643in}{3.210131in}}%
\pgfpathlineto{\pgfqpoint{1.321786in}{3.210142in}}%
\pgfpathlineto{\pgfqpoint{1.323960in}{3.208458in}}%
\pgfpathlineto{\pgfqpoint{1.326935in}{3.206132in}}%
\pgfpathlineto{\pgfqpoint{1.329078in}{3.206409in}}%
\pgfpathlineto{\pgfqpoint{1.331385in}{3.208496in}}%
\pgfpathlineto{\pgfqpoint{1.336293in}{3.213750in}}%
\pgfpathlineto{\pgfqpoint{1.338436in}{3.213409in}}%
\pgfpathlineto{\pgfqpoint{1.340739in}{3.211194in}}%
\pgfpathlineto{\pgfqpoint{1.343264in}{3.208997in}}%
\pgfpathlineto{\pgfqpoint{1.345407in}{3.208926in}}%
\pgfpathlineto{\pgfqpoint{1.347550in}{3.210457in}}%
\pgfpathlineto{\pgfqpoint{1.353107in}{3.215566in}}%
\pgfpathlineto{\pgfqpoint{1.355085in}{3.214860in}}%
\pgfpathlineto{\pgfqpoint{1.357512in}{3.212116in}}%
\pgfpathlineto{\pgfqpoint{1.360111in}{3.209887in}}%
\pgfpathlineto{\pgfqpoint{1.362254in}{3.209868in}}%
\pgfpathlineto{\pgfqpoint{1.364398in}{3.211448in}}%
\pgfpathlineto{\pgfqpoint{1.369097in}{3.215633in}}%
\pgfpathlineto{\pgfqpoint{1.371075in}{3.215097in}}%
\pgfpathlineto{\pgfqpoint{1.373414in}{3.212642in}}%
\pgfpathlineto{\pgfqpoint{1.376436in}{3.209397in}}%
\pgfpathlineto{\pgfqpoint{1.378579in}{3.209075in}}%
\pgfpathlineto{\pgfqpoint{1.380722in}{3.210361in}}%
\pgfpathlineto{\pgfqpoint{1.385747in}{3.214344in}}%
\pgfpathlineto{\pgfqpoint{1.387890in}{3.213419in}}%
\pgfpathlineto{\pgfqpoint{1.390386in}{3.210295in}}%
\pgfpathlineto{\pgfqpoint{1.393090in}{3.207528in}}%
\pgfpathlineto{\pgfqpoint{1.395233in}{3.207211in}}%
\pgfpathlineto{\pgfqpoint{1.397376in}{3.208506in}}%
\pgfpathlineto{\pgfqpoint{1.401902in}{3.212099in}}%
\pgfpathlineto{\pgfqpoint{1.403880in}{3.211459in}}%
\pgfpathlineto{\pgfqpoint{1.406266in}{3.208825in}}%
\pgfpathlineto{\pgfqpoint{1.409740in}{3.204922in}}%
\pgfpathlineto{\pgfqpoint{1.411883in}{3.204673in}}%
\pgfpathlineto{\pgfqpoint{1.414026in}{3.206042in}}%
\pgfpathlineto{\pgfqpoint{1.418222in}{3.209385in}}%
\pgfpathlineto{\pgfqpoint{1.420365in}{3.208731in}}%
\pgfpathlineto{\pgfqpoint{1.422765in}{3.206056in}}%
\pgfpathlineto{\pgfqpoint{1.426062in}{3.202252in}}%
\pgfpathlineto{\pgfqpoint{1.428205in}{3.201881in}}%
\pgfpathlineto{\pgfqpoint{1.430348in}{3.203140in}}%
\pgfpathlineto{\pgfqpoint{1.435037in}{3.206697in}}%
\pgfpathlineto{\pgfqpoint{1.437180in}{3.205821in}}%
\pgfpathlineto{\pgfqpoint{1.439649in}{3.202800in}}%
\pgfpathlineto{\pgfqpoint{1.442383in}{3.199790in}}%
\pgfpathlineto{\pgfqpoint{1.444526in}{3.199345in}}%
\pgfpathlineto{\pgfqpoint{1.446669in}{3.200538in}}%
\pgfpathlineto{\pgfqpoint{1.451686in}{3.204395in}}%
\pgfpathlineto{\pgfqpoint{1.453829in}{3.203465in}}%
\pgfpathlineto{\pgfqpoint{1.456326in}{3.200369in}}%
\pgfpathlineto{\pgfqpoint{1.459032in}{3.197672in}}%
\pgfpathlineto{\pgfqpoint{1.461175in}{3.197440in}}%
\pgfpathlineto{\pgfqpoint{1.463318in}{3.198846in}}%
\pgfpathlineto{\pgfqpoint{1.468006in}{3.202752in}}%
\pgfpathlineto{\pgfqpoint{1.470149in}{3.202045in}}%
\pgfpathlineto{\pgfqpoint{1.472566in}{3.199308in}}%
\pgfpathlineto{\pgfqpoint{1.475350in}{3.196450in}}%
\pgfpathlineto{\pgfqpoint{1.477493in}{3.196190in}}%
\pgfpathlineto{\pgfqpoint{1.479636in}{3.197572in}}%
\pgfpathlineto{\pgfqpoint{1.484656in}{3.201887in}}%
\pgfpathlineto{\pgfqpoint{1.486799in}{3.201156in}}%
\pgfpathlineto{\pgfqpoint{1.489226in}{3.198390in}}%
\pgfpathlineto{\pgfqpoint{1.491836in}{3.195918in}}%
\pgfpathlineto{\pgfqpoint{1.493979in}{3.195757in}}%
\pgfpathlineto{\pgfqpoint{1.496122in}{3.197236in}}%
\pgfpathlineto{\pgfqpoint{1.501306in}{3.201781in}}%
\pgfpathlineto{\pgfqpoint{1.503449in}{3.201014in}}%
\pgfpathlineto{\pgfqpoint{1.505879in}{3.198236in}}%
\pgfpathlineto{\pgfqpoint{1.508483in}{3.196032in}}%
\pgfpathlineto{\pgfqpoint{1.510626in}{3.196080in}}%
\pgfpathlineto{\pgfqpoint{1.512769in}{3.197761in}}%
\pgfpathlineto{\pgfqpoint{1.517461in}{3.202262in}}%
\pgfpathlineto{\pgfqpoint{1.519604in}{3.201817in}}%
\pgfpathlineto{\pgfqpoint{1.521964in}{3.199448in}}%
\pgfpathlineto{\pgfqpoint{1.524637in}{3.196903in}}%
\pgfpathlineto{\pgfqpoint{1.526780in}{3.196751in}}%
\pgfpathlineto{\pgfqpoint{1.528923in}{3.198237in}}%
\pgfpathlineto{\pgfqpoint{1.534440in}{3.203270in}}%
\pgfpathlineto{\pgfqpoint{1.536583in}{3.202487in}}%
\pgfpathlineto{\pgfqpoint{1.539142in}{3.199583in}}%
\pgfpathlineto{\pgfqpoint{1.541615in}{3.197806in}}%
\pgfpathlineto{\pgfqpoint{1.543758in}{3.198064in}}%
\pgfpathlineto{\pgfqpoint{1.546065in}{3.200155in}}%
\pgfpathlineto{\pgfqpoint{1.550101in}{3.204307in}}%
\pgfpathlineto{\pgfqpoint{1.552244in}{3.204172in}}%
\pgfpathlineto{\pgfqpoint{1.554502in}{3.202257in}}%
\pgfpathlineto{\pgfqpoint{1.557934in}{3.199065in}}%
\pgfpathlineto{\pgfqpoint{1.560077in}{3.199196in}}%
\pgfpathlineto{\pgfqpoint{1.562220in}{3.200951in}}%
\pgfpathlineto{\pgfqpoint{1.566915in}{3.205588in}}%
\pgfpathlineto{\pgfqpoint{1.569058in}{3.205186in}}%
\pgfpathlineto{\pgfqpoint{1.571368in}{3.202922in}}%
\pgfpathlineto{\pgfqpoint{1.574088in}{3.200335in}}%
\pgfpathlineto{\pgfqpoint{1.576232in}{3.200191in}}%
\pgfpathlineto{\pgfqpoint{1.578375in}{3.201675in}}%
\pgfpathlineto{\pgfqpoint{1.583895in}{3.206670in}}%
\pgfpathlineto{\pgfqpoint{1.586038in}{3.205856in}}%
\pgfpathlineto{\pgfqpoint{1.588595in}{3.202905in}}%
\pgfpathlineto{\pgfqpoint{1.591068in}{3.201071in}}%
\pgfpathlineto{\pgfqpoint{1.593211in}{3.201270in}}%
\pgfpathlineto{\pgfqpoint{1.595519in}{3.203290in}}%
\pgfpathlineto{\pgfqpoint{1.599555in}{3.207303in}}%
\pgfpathlineto{\pgfqpoint{1.601698in}{3.207085in}}%
\pgfpathlineto{\pgfqpoint{1.603956in}{3.205076in}}%
\pgfpathlineto{\pgfqpoint{1.607389in}{3.201730in}}%
\pgfpathlineto{\pgfqpoint{1.609532in}{3.201760in}}%
\pgfpathlineto{\pgfqpoint{1.611675in}{3.203408in}}%
\pgfpathlineto{\pgfqpoint{1.616370in}{3.207799in}}%
\pgfpathlineto{\pgfqpoint{1.618513in}{3.207280in}}%
\pgfpathlineto{\pgfqpoint{1.620873in}{3.204815in}}%
\pgfpathlineto{\pgfqpoint{1.623545in}{3.202144in}}%
\pgfpathlineto{\pgfqpoint{1.625688in}{3.201877in}}%
\pgfpathlineto{\pgfqpoint{1.627831in}{3.203237in}}%
\pgfpathlineto{\pgfqpoint{1.633349in}{3.207904in}}%
\pgfpathlineto{\pgfqpoint{1.635492in}{3.206962in}}%
\pgfpathlineto{\pgfqpoint{1.638217in}{3.203656in}}%
\pgfpathlineto{\pgfqpoint{1.640690in}{3.201828in}}%
\pgfpathlineto{\pgfqpoint{1.642833in}{3.202030in}}%
\pgfpathlineto{\pgfqpoint{1.645141in}{3.204052in}}%
\pgfpathlineto{\pgfqpoint{1.648845in}{3.207586in}}%
\pgfpathlineto{\pgfqpoint{1.650988in}{3.207379in}}%
\pgfpathlineto{\pgfqpoint{1.653131in}{3.205525in}}%
\pgfpathlineto{\pgfqpoint{1.657175in}{3.201520in}}%
\pgfpathlineto{\pgfqpoint{1.659318in}{3.201697in}}%
\pgfpathlineto{\pgfqpoint{1.661461in}{3.203489in}}%
\pgfpathlineto{\pgfqpoint{1.665330in}{3.207188in}}%
\pgfpathlineto{\pgfqpoint{1.667473in}{3.206961in}}%
\pgfpathlineto{\pgfqpoint{1.669616in}{3.205088in}}%
\pgfpathlineto{\pgfqpoint{1.673661in}{3.201053in}}%
\pgfpathlineto{\pgfqpoint{1.675804in}{3.201217in}}%
\pgfpathlineto{\pgfqpoint{1.677947in}{3.202998in}}%
\pgfpathlineto{\pgfqpoint{1.681815in}{3.206678in}}%
\pgfpathlineto{\pgfqpoint{1.683958in}{3.206443in}}%
\pgfpathlineto{\pgfqpoint{1.686101in}{3.204565in}}%
\pgfpathlineto{\pgfqpoint{1.690146in}{3.200523in}}%
\pgfpathlineto{\pgfqpoint{1.692289in}{3.200686in}}%
\pgfpathlineto{\pgfqpoint{1.694432in}{3.202467in}}%
\pgfpathlineto{\pgfqpoint{1.698300in}{3.206149in}}%
\pgfpathlineto{\pgfqpoint{1.700443in}{3.205919in}}%
\pgfpathlineto{\pgfqpoint{1.702586in}{3.204046in}}%
\pgfpathlineto{\pgfqpoint{1.706631in}{3.200017in}}%
\pgfpathlineto{\pgfqpoint{1.708774in}{3.200188in}}%
\pgfpathlineto{\pgfqpoint{1.710917in}{3.201979in}}%
\pgfpathlineto{\pgfqpoint{1.714784in}{3.205682in}}%
\pgfpathlineto{\pgfqpoint{1.716927in}{3.205464in}}%
\pgfpathlineto{\pgfqpoint{1.719071in}{3.203605in}}%
\pgfpathlineto{\pgfqpoint{1.723116in}{3.199604in}}%
\pgfpathlineto{\pgfqpoint{1.725259in}{3.199791in}}%
\pgfpathlineto{\pgfqpoint{1.727402in}{3.201598in}}%
\pgfpathlineto{\pgfqpoint{1.731269in}{3.205332in}}%
\pgfpathlineto{\pgfqpoint{1.733412in}{3.205132in}}%
\pgfpathlineto{\pgfqpoint{1.735555in}{3.203290in}}%
\pgfpathlineto{\pgfqpoint{1.739601in}{3.199325in}}%
\pgfpathlineto{\pgfqpoint{1.741744in}{3.199531in}}%
\pgfpathlineto{\pgfqpoint{1.744052in}{3.201564in}}%
\pgfpathlineto{\pgfqpoint{1.747754in}{3.205127in}}%
\pgfpathlineto{\pgfqpoint{1.749897in}{3.204946in}}%
\pgfpathlineto{\pgfqpoint{1.752040in}{3.203125in}}%
\pgfpathlineto{\pgfqpoint{1.756085in}{3.199196in}}%
\pgfpathlineto{\pgfqpoint{1.758228in}{3.199421in}}%
\pgfpathlineto{\pgfqpoint{1.760536in}{3.201474in}}%
\pgfpathlineto{\pgfqpoint{1.764239in}{3.205070in}}%
\pgfpathlineto{\pgfqpoint{1.766382in}{3.204908in}}%
\pgfpathlineto{\pgfqpoint{1.768525in}{3.203104in}}%
\pgfpathlineto{\pgfqpoint{1.772570in}{3.199208in}}%
\pgfpathlineto{\pgfqpoint{1.774713in}{3.199450in}}%
\pgfpathlineto{\pgfqpoint{1.777020in}{3.201520in}}%
\pgfpathlineto{\pgfqpoint{1.780724in}{3.205143in}}%
\pgfpathlineto{\pgfqpoint{1.782867in}{3.204995in}}%
\pgfpathlineto{\pgfqpoint{1.785010in}{3.203205in}}%
\pgfpathlineto{\pgfqpoint{1.789054in}{3.199334in}}%
\pgfpathlineto{\pgfqpoint{1.791197in}{3.199587in}}%
\pgfpathlineto{\pgfqpoint{1.793505in}{3.201669in}}%
\pgfpathlineto{\pgfqpoint{1.797209in}{3.205310in}}%
\pgfpathlineto{\pgfqpoint{1.799352in}{3.205172in}}%
\pgfpathlineto{\pgfqpoint{1.801495in}{3.203391in}}%
\pgfpathlineto{\pgfqpoint{1.805539in}{3.199534in}}%
\pgfpathlineto{\pgfqpoint{1.807682in}{3.199793in}}%
\pgfpathlineto{\pgfqpoint{1.809989in}{3.201881in}}%
\pgfpathlineto{\pgfqpoint{1.813694in}{3.205531in}}%
\pgfpathlineto{\pgfqpoint{1.815837in}{3.205396in}}%
\pgfpathlineto{\pgfqpoint{1.817980in}{3.203618in}}%
\pgfpathlineto{\pgfqpoint{1.822023in}{3.199766in}}%
\pgfpathlineto{\pgfqpoint{1.824166in}{3.200026in}}%
\pgfpathlineto{\pgfqpoint{1.826474in}{3.202114in}}%
\pgfpathlineto{\pgfqpoint{1.830178in}{3.205764in}}%
\pgfpathlineto{\pgfqpoint{1.832321in}{3.205628in}}%
\pgfpathlineto{\pgfqpoint{1.834464in}{3.203848in}}%
\pgfpathlineto{\pgfqpoint{1.838508in}{3.199990in}}%
\pgfpathlineto{\pgfqpoint{1.840651in}{3.200247in}}%
\pgfpathlineto{\pgfqpoint{1.842959in}{3.202331in}}%
\pgfpathlineto{\pgfqpoint{1.846663in}{3.205973in}}%
\pgfpathlineto{\pgfqpoint{1.848806in}{3.205832in}}%
\pgfpathlineto{\pgfqpoint{1.850949in}{3.204047in}}%
\pgfpathlineto{\pgfqpoint{1.854993in}{3.200177in}}%
\pgfpathlineto{\pgfqpoint{1.857136in}{3.200428in}}%
\pgfpathlineto{\pgfqpoint{1.859444in}{3.202504in}}%
\pgfpathlineto{\pgfqpoint{1.863148in}{3.206133in}}%
\pgfpathlineto{\pgfqpoint{1.865291in}{3.205985in}}%
\pgfpathlineto{\pgfqpoint{1.867434in}{3.204191in}}%
\pgfpathlineto{\pgfqpoint{1.871478in}{3.200306in}}%
\pgfpathlineto{\pgfqpoint{1.873621in}{3.200548in}}%
\pgfpathlineto{\pgfqpoint{1.875929in}{3.202616in}}%
\pgfpathlineto{\pgfqpoint{1.879633in}{3.206230in}}%
\pgfpathlineto{\pgfqpoint{1.881776in}{3.206073in}}%
\pgfpathlineto{\pgfqpoint{1.883919in}{3.204271in}}%
\pgfpathlineto{\pgfqpoint{1.887963in}{3.200370in}}%
\pgfpathlineto{\pgfqpoint{1.890106in}{3.200603in}}%
\pgfpathlineto{\pgfqpoint{1.892414in}{3.202662in}}%
\pgfpathlineto{\pgfqpoint{1.896118in}{3.206261in}}%
\pgfpathlineto{\pgfqpoint{1.898261in}{3.206096in}}%
\pgfpathlineto{\pgfqpoint{1.900404in}{3.204286in}}%
\pgfpathlineto{\pgfqpoint{1.904448in}{3.200370in}}%
\pgfpathlineto{\pgfqpoint{1.906591in}{3.200597in}}%
\pgfpathlineto{\pgfqpoint{1.908899in}{3.202647in}}%
\pgfpathlineto{\pgfqpoint{1.912603in}{3.206235in}}%
\pgfpathlineto{\pgfqpoint{1.914746in}{3.206063in}}%
\pgfpathlineto{\pgfqpoint{1.916889in}{3.204246in}}%
\pgfpathlineto{\pgfqpoint{1.920933in}{3.200320in}}%
\pgfpathlineto{\pgfqpoint{1.923076in}{3.200541in}}%
\pgfpathlineto{\pgfqpoint{1.925384in}{3.202586in}}%
\pgfpathlineto{\pgfqpoint{1.929088in}{3.206165in}}%
\pgfpathlineto{\pgfqpoint{1.931231in}{3.205988in}}%
\pgfpathlineto{\pgfqpoint{1.933374in}{3.204168in}}%
\pgfpathlineto{\pgfqpoint{1.937418in}{3.200234in}}%
\pgfpathlineto{\pgfqpoint{1.939561in}{3.200452in}}%
\pgfpathlineto{\pgfqpoint{1.941869in}{3.202495in}}%
\pgfpathlineto{\pgfqpoint{1.945572in}{3.206069in}}%
\pgfpathlineto{\pgfqpoint{1.947715in}{3.205891in}}%
\pgfpathlineto{\pgfqpoint{1.949858in}{3.204069in}}%
\pgfpathlineto{\pgfqpoint{1.953903in}{3.200133in}}%
\pgfpathlineto{\pgfqpoint{1.956046in}{3.200350in}}%
\pgfpathlineto{\pgfqpoint{1.958354in}{3.202392in}}%
\pgfpathlineto{\pgfqpoint{1.962057in}{3.205967in}}%
\pgfpathlineto{\pgfqpoint{1.964200in}{3.205789in}}%
\pgfpathlineto{\pgfqpoint{1.966343in}{3.203967in}}%
\pgfpathlineto{\pgfqpoint{1.970388in}{3.200033in}}%
\pgfpathlineto{\pgfqpoint{1.972531in}{3.200252in}}%
\pgfpathlineto{\pgfqpoint{1.974839in}{3.202296in}}%
\pgfpathlineto{\pgfqpoint{1.978542in}{3.205873in}}%
\pgfpathlineto{\pgfqpoint{1.980685in}{3.205697in}}%
\pgfpathlineto{\pgfqpoint{1.982828in}{3.203878in}}%
\pgfpathlineto{\pgfqpoint{1.986872in}{3.199949in}}%
\pgfpathlineto{\pgfqpoint{1.989016in}{3.200170in}}%
\pgfpathlineto{\pgfqpoint{1.991323in}{3.202217in}}%
\pgfpathlineto{\pgfqpoint{1.995027in}{3.205800in}}%
\pgfpathlineto{\pgfqpoint{1.997170in}{3.205628in}}%
\pgfpathlineto{\pgfqpoint{1.999313in}{3.203812in}}%
\pgfpathlineto{\pgfqpoint{2.003357in}{3.199889in}}%
\pgfpathlineto{\pgfqpoint{2.005500in}{3.200114in}}%
\pgfpathlineto{\pgfqpoint{2.007808in}{3.202165in}}%
\pgfpathlineto{\pgfqpoint{2.011512in}{3.205755in}}%
\pgfpathlineto{\pgfqpoint{2.013655in}{3.205586in}}%
\pgfpathlineto{\pgfqpoint{2.015798in}{3.203774in}}%
\pgfpathlineto{\pgfqpoint{2.019842in}{3.199859in}}%
\pgfpathlineto{\pgfqpoint{2.021985in}{3.200087in}}%
\pgfpathlineto{\pgfqpoint{2.024293in}{3.202142in}}%
\pgfpathlineto{\pgfqpoint{2.027997in}{3.205738in}}%
\pgfpathlineto{\pgfqpoint{2.030140in}{3.205573in}}%
\pgfpathlineto{\pgfqpoint{2.032283in}{3.203765in}}%
\pgfpathlineto{\pgfqpoint{2.036327in}{3.199856in}}%
\pgfpathlineto{\pgfqpoint{2.038470in}{3.200088in}}%
\pgfpathlineto{\pgfqpoint{2.040778in}{3.202146in}}%
\pgfpathlineto{\pgfqpoint{2.044481in}{3.205747in}}%
\pgfpathlineto{\pgfqpoint{2.046624in}{3.205585in}}%
\pgfpathlineto{\pgfqpoint{2.048768in}{3.203780in}}%
\pgfpathlineto{\pgfqpoint{2.052812in}{3.199876in}}%
\pgfpathlineto{\pgfqpoint{2.054955in}{3.200110in}}%
\pgfpathlineto{\pgfqpoint{2.057263in}{3.202171in}}%
\pgfpathlineto{\pgfqpoint{2.060966in}{3.205777in}}%
\pgfpathlineto{\pgfqpoint{2.063109in}{3.205616in}}%
\pgfpathlineto{\pgfqpoint{2.065252in}{3.203813in}}%
\pgfpathlineto{\pgfqpoint{2.069296in}{3.199912in}}%
\pgfpathlineto{\pgfqpoint{2.071439in}{3.200148in}}%
\pgfpathlineto{\pgfqpoint{2.073747in}{3.202210in}}%
\pgfpathlineto{\pgfqpoint{2.077451in}{3.205818in}}%
\pgfpathlineto{\pgfqpoint{2.079594in}{3.205659in}}%
\pgfpathlineto{\pgfqpoint{2.081737in}{3.203856in}}%
\pgfpathlineto{\pgfqpoint{2.085781in}{3.199956in}}%
\pgfpathlineto{\pgfqpoint{2.087924in}{3.200193in}}%
\pgfpathlineto{\pgfqpoint{2.090232in}{3.202255in}}%
\pgfpathlineto{\pgfqpoint{2.093936in}{3.205863in}}%
\pgfpathlineto{\pgfqpoint{2.096079in}{3.205704in}}%
\pgfpathlineto{\pgfqpoint{2.098222in}{3.203901in}}%
\pgfpathlineto{\pgfqpoint{2.102266in}{3.200001in}}%
\pgfpathlineto{\pgfqpoint{2.104409in}{3.200236in}}%
\pgfpathlineto{\pgfqpoint{2.106717in}{3.202298in}}%
\pgfpathlineto{\pgfqpoint{2.110421in}{3.205905in}}%
\pgfpathlineto{\pgfqpoint{2.112564in}{3.205744in}}%
\pgfpathlineto{\pgfqpoint{2.114707in}{3.203941in}}%
\pgfpathlineto{\pgfqpoint{2.118751in}{3.200038in}}%
\pgfpathlineto{\pgfqpoint{2.120894in}{3.200273in}}%
\pgfpathlineto{\pgfqpoint{2.123202in}{3.202333in}}%
\pgfpathlineto{\pgfqpoint{2.126906in}{3.205938in}}%
\pgfpathlineto{\pgfqpoint{2.129049in}{3.205776in}}%
\pgfpathlineto{\pgfqpoint{2.131192in}{3.203971in}}%
\pgfpathlineto{\pgfqpoint{2.135236in}{3.200066in}}%
\pgfpathlineto{\pgfqpoint{2.137379in}{3.200299in}}%
\pgfpathlineto{\pgfqpoint{2.139687in}{3.202358in}}%
\pgfpathlineto{\pgfqpoint{2.143391in}{3.205959in}}%
\pgfpathlineto{\pgfqpoint{2.145534in}{3.205796in}}%
\pgfpathlineto{\pgfqpoint{2.147677in}{3.203989in}}%
\pgfpathlineto{\pgfqpoint{2.151721in}{3.200081in}}%
\pgfpathlineto{\pgfqpoint{2.153864in}{3.200312in}}%
\pgfpathlineto{\pgfqpoint{2.156172in}{3.202369in}}%
\pgfpathlineto{\pgfqpoint{2.159875in}{3.205968in}}%
\pgfpathlineto{\pgfqpoint{2.162018in}{3.205803in}}%
\pgfpathlineto{\pgfqpoint{2.164161in}{3.203994in}}%
\pgfpathlineto{\pgfqpoint{2.168206in}{3.200083in}}%
\pgfpathlineto{\pgfqpoint{2.170349in}{3.200313in}}%
\pgfpathlineto{\pgfqpoint{2.172656in}{3.202369in}}%
\pgfpathlineto{\pgfqpoint{2.176360in}{3.205965in}}%
\pgfpathlineto{\pgfqpoint{2.178503in}{3.205798in}}%
\pgfpathlineto{\pgfqpoint{2.180646in}{3.203989in}}%
\pgfpathlineto{\pgfqpoint{2.184690in}{3.200075in}}%
\pgfpathlineto{\pgfqpoint{2.186833in}{3.200304in}}%
\pgfpathlineto{\pgfqpoint{2.189141in}{3.202358in}}%
\pgfpathlineto{\pgfqpoint{2.192845in}{3.205953in}}%
\pgfpathlineto{\pgfqpoint{2.194988in}{3.205786in}}%
\pgfpathlineto{\pgfqpoint{2.197131in}{3.203975in}}%
\pgfpathlineto{\pgfqpoint{2.201175in}{3.200060in}}%
\pgfpathlineto{\pgfqpoint{2.203318in}{3.200288in}}%
\pgfpathlineto{\pgfqpoint{2.205626in}{3.202342in}}%
\pgfpathlineto{\pgfqpoint{2.209330in}{3.205935in}}%
\pgfpathlineto{\pgfqpoint{2.211473in}{3.205767in}}%
\pgfpathlineto{\pgfqpoint{2.213616in}{3.203956in}}%
\pgfpathlineto{\pgfqpoint{2.217660in}{3.200041in}}%
\pgfpathlineto{\pgfqpoint{2.219803in}{3.200269in}}%
\pgfpathlineto{\pgfqpoint{2.222111in}{3.202322in}}%
\pgfpathlineto{\pgfqpoint{2.225815in}{3.205915in}}%
\pgfpathlineto{\pgfqpoint{2.227958in}{3.205748in}}%
\pgfpathlineto{\pgfqpoint{2.230101in}{3.203936in}}%
\pgfpathlineto{\pgfqpoint{2.234145in}{3.200021in}}%
\pgfpathlineto{\pgfqpoint{2.236288in}{3.200249in}}%
\pgfpathlineto{\pgfqpoint{2.238596in}{3.202303in}}%
\pgfpathlineto{\pgfqpoint{2.242300in}{3.205896in}}%
\pgfpathlineto{\pgfqpoint{2.244443in}{3.205729in}}%
\pgfpathlineto{\pgfqpoint{2.246586in}{3.203918in}}%
\pgfpathlineto{\pgfqpoint{2.250630in}{3.200004in}}%
\pgfpathlineto{\pgfqpoint{2.252773in}{3.200233in}}%
\pgfpathlineto{\pgfqpoint{2.255081in}{3.202287in}}%
\pgfpathlineto{\pgfqpoint{2.258784in}{3.205881in}}%
\pgfpathlineto{\pgfqpoint{2.260928in}{3.205715in}}%
\pgfpathlineto{\pgfqpoint{2.263071in}{3.203905in}}%
\pgfpathlineto{\pgfqpoint{2.267115in}{3.199991in}}%
\pgfpathlineto{\pgfqpoint{2.269258in}{3.200221in}}%
\pgfpathlineto{\pgfqpoint{2.271566in}{3.202276in}}%
\pgfpathlineto{\pgfqpoint{2.275269in}{3.205871in}}%
\pgfpathlineto{\pgfqpoint{2.277412in}{3.205706in}}%
\pgfpathlineto{\pgfqpoint{2.279555in}{3.203896in}}%
\pgfpathlineto{\pgfqpoint{2.283600in}{3.199984in}}%
\pgfpathlineto{\pgfqpoint{2.285743in}{3.200214in}}%
\pgfpathlineto{\pgfqpoint{2.288050in}{3.202270in}}%
\pgfpathlineto{\pgfqpoint{2.291754in}{3.205867in}}%
\pgfpathlineto{\pgfqpoint{2.293897in}{3.205702in}}%
\pgfpathlineto{\pgfqpoint{2.296040in}{3.203893in}}%
\pgfpathlineto{\pgfqpoint{2.300084in}{3.199983in}}%
\pgfpathlineto{\pgfqpoint{2.302227in}{3.200213in}}%
\pgfpathlineto{\pgfqpoint{2.304535in}{3.202270in}}%
\pgfpathlineto{\pgfqpoint{2.308239in}{3.205868in}}%
\pgfpathlineto{\pgfqpoint{2.310382in}{3.205703in}}%
\pgfpathlineto{\pgfqpoint{2.312525in}{3.203895in}}%
\pgfpathlineto{\pgfqpoint{2.316569in}{3.199986in}}%
\pgfpathlineto{\pgfqpoint{2.318712in}{3.200217in}}%
\pgfpathlineto{\pgfqpoint{2.321020in}{3.202274in}}%
\pgfpathlineto{\pgfqpoint{2.324724in}{3.205873in}}%
\pgfpathlineto{\pgfqpoint{2.326867in}{3.205709in}}%
\pgfpathlineto{\pgfqpoint{2.329010in}{3.203901in}}%
\pgfpathlineto{\pgfqpoint{2.333054in}{3.199992in}}%
\pgfpathlineto{\pgfqpoint{2.335197in}{3.200224in}}%
\pgfpathlineto{\pgfqpoint{2.337505in}{3.202281in}}%
\pgfpathlineto{\pgfqpoint{2.341209in}{3.205881in}}%
\pgfpathlineto{\pgfqpoint{2.343352in}{3.205717in}}%
\pgfpathlineto{\pgfqpoint{2.345495in}{3.203909in}}%
\pgfpathlineto{\pgfqpoint{2.349539in}{3.200001in}}%
\pgfpathlineto{\pgfqpoint{2.351682in}{3.200232in}}%
\pgfpathlineto{\pgfqpoint{2.353990in}{3.202290in}}%
\pgfpathlineto{\pgfqpoint{2.357694in}{3.205889in}}%
\pgfpathlineto{\pgfqpoint{2.359837in}{3.205725in}}%
\pgfpathlineto{\pgfqpoint{2.361980in}{3.203918in}}%
\pgfpathlineto{\pgfqpoint{2.366024in}{3.200009in}}%
\pgfpathlineto{\pgfqpoint{2.368167in}{3.200241in}}%
\pgfpathlineto{\pgfqpoint{2.370475in}{3.202298in}}%
\pgfpathlineto{\pgfqpoint{2.374178in}{3.205897in}}%
\pgfpathlineto{\pgfqpoint{2.376321in}{3.205733in}}%
\pgfpathlineto{\pgfqpoint{2.378464in}{3.203926in}}%
\pgfpathlineto{\pgfqpoint{2.382509in}{3.200017in}}%
\pgfpathlineto{\pgfqpoint{2.384652in}{3.200248in}}%
\pgfpathlineto{\pgfqpoint{2.386959in}{3.202305in}}%
\pgfpathlineto{\pgfqpoint{2.390663in}{3.205904in}}%
\pgfpathlineto{\pgfqpoint{2.392806in}{3.205740in}}%
\pgfpathlineto{\pgfqpoint{2.394949in}{3.203932in}}%
\pgfpathlineto{\pgfqpoint{2.398993in}{3.200023in}}%
\pgfpathlineto{\pgfqpoint{2.401136in}{3.200254in}}%
\pgfpathlineto{\pgfqpoint{2.403444in}{3.202310in}}%
\pgfpathlineto{\pgfqpoint{2.407148in}{3.205909in}}%
\pgfpathlineto{\pgfqpoint{2.409291in}{3.205744in}}%
\pgfpathlineto{\pgfqpoint{2.411434in}{3.203936in}}%
\pgfpathlineto{\pgfqpoint{2.415478in}{3.200026in}}%
\pgfpathlineto{\pgfqpoint{2.417621in}{3.200257in}}%
\pgfpathlineto{\pgfqpoint{2.419929in}{3.202313in}}%
\pgfpathlineto{\pgfqpoint{2.423633in}{3.205911in}}%
\pgfpathlineto{\pgfqpoint{2.425776in}{3.205746in}}%
\pgfpathlineto{\pgfqpoint{2.427919in}{3.203938in}}%
\pgfpathlineto{\pgfqpoint{2.431963in}{3.200027in}}%
\pgfpathlineto{\pgfqpoint{2.434106in}{3.200257in}}%
\pgfpathlineto{\pgfqpoint{2.436414in}{3.202314in}}%
\pgfpathlineto{\pgfqpoint{2.440118in}{3.205911in}}%
\pgfpathlineto{\pgfqpoint{2.442261in}{3.205746in}}%
\pgfpathlineto{\pgfqpoint{2.444404in}{3.203937in}}%
\pgfpathlineto{\pgfqpoint{2.448448in}{3.200026in}}%
\pgfpathlineto{\pgfqpoint{2.450591in}{3.200256in}}%
\pgfpathlineto{\pgfqpoint{2.452899in}{3.202312in}}%
\pgfpathlineto{\pgfqpoint{2.456603in}{3.205909in}}%
\pgfpathlineto{\pgfqpoint{2.458746in}{3.205744in}}%
\pgfpathlineto{\pgfqpoint{2.460889in}{3.203935in}}%
\pgfpathlineto{\pgfqpoint{2.464933in}{3.200023in}}%
\pgfpathlineto{\pgfqpoint{2.467076in}{3.200253in}}%
\pgfpathlineto{\pgfqpoint{2.469384in}{3.202309in}}%
\pgfpathlineto{\pgfqpoint{2.473088in}{3.205906in}}%
\pgfpathlineto{\pgfqpoint{2.475231in}{3.205740in}}%
\pgfpathlineto{\pgfqpoint{2.477374in}{3.203931in}}%
\pgfpathlineto{\pgfqpoint{2.481418in}{3.200020in}}%
\pgfpathlineto{\pgfqpoint{2.483561in}{3.200250in}}%
\pgfpathlineto{\pgfqpoint{2.485869in}{3.202305in}}%
\pgfpathlineto{\pgfqpoint{2.489572in}{3.205902in}}%
\pgfpathlineto{\pgfqpoint{2.491715in}{3.205736in}}%
\pgfpathlineto{\pgfqpoint{2.493858in}{3.203927in}}%
\pgfpathlineto{\pgfqpoint{2.497904in}{3.200035in}}%
\pgfpathlineto{\pgfqpoint{2.500047in}{3.200279in}}%
\pgfpathlineto{\pgfqpoint{2.502355in}{3.202349in}}%
\pgfpathlineto{\pgfqpoint{2.506057in}{3.205963in}}%
\pgfpathlineto{\pgfqpoint{2.508200in}{3.205808in}}%
\pgfpathlineto{\pgfqpoint{2.510343in}{3.204009in}}%
\pgfpathlineto{\pgfqpoint{2.514387in}{3.200116in}}%
\pgfpathlineto{\pgfqpoint{2.516530in}{3.200355in}}%
\pgfpathlineto{\pgfqpoint{2.518838in}{3.202419in}}%
\pgfpathlineto{\pgfqpoint{2.522542in}{3.206029in}}%
\pgfpathlineto{\pgfqpoint{2.524685in}{3.205871in}}%
\pgfpathlineto{\pgfqpoint{2.526828in}{3.204068in}}%
\pgfpathlineto{\pgfqpoint{2.530872in}{3.200167in}}%
\pgfpathlineto{\pgfqpoint{2.533015in}{3.200401in}}%
\pgfpathlineto{\pgfqpoint{2.535323in}{3.202461in}}%
\pgfpathlineto{\pgfqpoint{2.539027in}{3.206064in}}%
\pgfpathlineto{\pgfqpoint{2.541170in}{3.205902in}}%
\pgfpathlineto{\pgfqpoint{2.543313in}{3.204095in}}%
\pgfpathlineto{\pgfqpoint{2.547357in}{3.200186in}}%
\pgfpathlineto{\pgfqpoint{2.549500in}{3.200417in}}%
\pgfpathlineto{\pgfqpoint{2.551808in}{3.202473in}}%
\pgfpathlineto{\pgfqpoint{2.555512in}{3.206069in}}%
\pgfpathlineto{\pgfqpoint{2.557655in}{3.205903in}}%
\pgfpathlineto{\pgfqpoint{2.559798in}{3.204093in}}%
\pgfpathlineto{\pgfqpoint{2.563842in}{3.200177in}}%
\pgfpathlineto{\pgfqpoint{2.565985in}{3.200405in}}%
\pgfpathlineto{\pgfqpoint{2.568293in}{3.202458in}}%
\pgfpathlineto{\pgfqpoint{2.571997in}{3.206049in}}%
\pgfpathlineto{\pgfqpoint{2.574140in}{3.205880in}}%
\pgfpathlineto{\pgfqpoint{2.576283in}{3.204066in}}%
\pgfpathlineto{\pgfqpoint{2.580327in}{3.200147in}}%
\pgfpathlineto{\pgfqpoint{2.582470in}{3.200372in}}%
\pgfpathlineto{\pgfqpoint{2.584778in}{3.202422in}}%
\pgfpathlineto{\pgfqpoint{2.588481in}{3.206010in}}%
\pgfpathlineto{\pgfqpoint{2.590624in}{3.205839in}}%
\pgfpathlineto{\pgfqpoint{2.592768in}{3.204025in}}%
\pgfpathlineto{\pgfqpoint{2.596812in}{3.200103in}}%
\pgfpathlineto{\pgfqpoint{2.598955in}{3.200327in}}%
\pgfpathlineto{\pgfqpoint{2.601263in}{3.202376in}}%
\pgfpathlineto{\pgfqpoint{2.604966in}{3.205963in}}%
\pgfpathlineto{\pgfqpoint{2.607109in}{3.205791in}}%
\pgfpathlineto{\pgfqpoint{2.609252in}{3.203976in}}%
\pgfpathlineto{\pgfqpoint{2.613297in}{3.200054in}}%
\pgfpathlineto{\pgfqpoint{2.615440in}{3.200278in}}%
\pgfpathlineto{\pgfqpoint{2.617747in}{3.202328in}}%
\pgfpathlineto{\pgfqpoint{2.621451in}{3.205915in}}%
\pgfpathlineto{\pgfqpoint{2.623594in}{3.205744in}}%
\pgfpathlineto{\pgfqpoint{2.625737in}{3.203929in}}%
\pgfpathlineto{\pgfqpoint{2.629781in}{3.200008in}}%
\pgfpathlineto{\pgfqpoint{2.631924in}{3.200233in}}%
\pgfpathlineto{\pgfqpoint{2.634232in}{3.202284in}}%
\pgfpathlineto{\pgfqpoint{2.637936in}{3.205873in}}%
\pgfpathlineto{\pgfqpoint{2.640079in}{3.205703in}}%
\pgfpathlineto{\pgfqpoint{2.642222in}{3.203890in}}%
\pgfpathlineto{\pgfqpoint{2.646266in}{3.199972in}}%
\pgfpathlineto{\pgfqpoint{2.648409in}{3.200198in}}%
\pgfpathlineto{\pgfqpoint{2.650717in}{3.202251in}}%
\pgfpathlineto{\pgfqpoint{2.654421in}{3.205842in}}%
\pgfpathlineto{\pgfqpoint{2.656564in}{3.205674in}}%
\pgfpathlineto{\pgfqpoint{2.658707in}{3.203863in}}%
\pgfpathlineto{\pgfqpoint{2.662751in}{3.199948in}}%
\pgfpathlineto{\pgfqpoint{2.664894in}{3.200176in}}%
\pgfpathlineto{\pgfqpoint{2.667202in}{3.202231in}}%
\pgfpathlineto{\pgfqpoint{2.670906in}{3.205825in}}%
\pgfpathlineto{\pgfqpoint{2.673049in}{3.205659in}}%
\pgfpathlineto{\pgfqpoint{2.675192in}{3.203850in}}%
\pgfpathlineto{\pgfqpoint{2.679236in}{3.199938in}}%
\pgfpathlineto{\pgfqpoint{2.681379in}{3.200168in}}%
\pgfpathlineto{\pgfqpoint{2.683687in}{3.202224in}}%
\pgfpathlineto{\pgfqpoint{2.687391in}{3.205822in}}%
\pgfpathlineto{\pgfqpoint{2.689534in}{3.205657in}}%
\pgfpathlineto{\pgfqpoint{2.691677in}{3.203849in}}%
\pgfpathlineto{\pgfqpoint{2.695721in}{3.199941in}}%
\pgfpathlineto{\pgfqpoint{2.697864in}{3.200172in}}%
\pgfpathlineto{\pgfqpoint{2.700172in}{3.202230in}}%
\pgfpathlineto{\pgfqpoint{2.703875in}{3.205830in}}%
\pgfpathlineto{\pgfqpoint{2.706018in}{3.205667in}}%
\pgfpathlineto{\pgfqpoint{2.708161in}{3.203860in}}%
\pgfpathlineto{\pgfqpoint{2.712206in}{3.199953in}}%
\pgfpathlineto{\pgfqpoint{2.714349in}{3.200186in}}%
\pgfpathlineto{\pgfqpoint{2.716656in}{3.202245in}}%
\pgfpathlineto{\pgfqpoint{2.720360in}{3.205846in}}%
\pgfpathlineto{\pgfqpoint{2.722503in}{3.205684in}}%
\pgfpathlineto{\pgfqpoint{2.724646in}{3.203878in}}%
\pgfpathlineto{\pgfqpoint{2.728690in}{3.199972in}}%
\pgfpathlineto{\pgfqpoint{2.730833in}{3.200205in}}%
\pgfpathlineto{\pgfqpoint{2.733141in}{3.202265in}}%
\pgfpathlineto{\pgfqpoint{2.736845in}{3.205867in}}%
\pgfpathlineto{\pgfqpoint{2.738988in}{3.205705in}}%
\pgfpathlineto{\pgfqpoint{2.741131in}{3.203899in}}%
\pgfpathlineto{\pgfqpoint{2.745175in}{3.199994in}}%
\pgfpathlineto{\pgfqpoint{2.747318in}{3.200227in}}%
\pgfpathlineto{\pgfqpoint{2.749626in}{3.202286in}}%
\pgfpathlineto{\pgfqpoint{2.753330in}{3.205888in}}%
\pgfpathlineto{\pgfqpoint{2.755473in}{3.205726in}}%
\pgfpathlineto{\pgfqpoint{2.757616in}{3.203920in}}%
\pgfpathlineto{\pgfqpoint{2.761662in}{3.200033in}}%
\pgfpathlineto{\pgfqpoint{2.763805in}{3.200280in}}%
\pgfpathlineto{\pgfqpoint{2.766113in}{3.202353in}}%
\pgfpathlineto{\pgfqpoint{2.769815in}{3.205972in}}%
\pgfpathlineto{\pgfqpoint{2.771958in}{3.205819in}}%
\pgfpathlineto{\pgfqpoint{2.774101in}{3.204023in}}%
\pgfpathlineto{\pgfqpoint{2.778145in}{3.200134in}}%
\pgfpathlineto{\pgfqpoint{2.780288in}{3.200375in}}%
\pgfpathlineto{\pgfqpoint{2.782596in}{3.202441in}}%
\pgfpathlineto{\pgfqpoint{2.786300in}{3.206055in}}%
\pgfpathlineto{\pgfqpoint{2.788443in}{3.205898in}}%
\pgfpathlineto{\pgfqpoint{2.790586in}{3.204097in}}%
\pgfpathlineto{\pgfqpoint{2.794630in}{3.200199in}}%
\pgfpathlineto{\pgfqpoint{2.796773in}{3.200435in}}%
\pgfpathlineto{\pgfqpoint{2.799080in}{3.202496in}}%
\pgfpathlineto{\pgfqpoint{2.802784in}{3.206101in}}%
\pgfpathlineto{\pgfqpoint{2.804928in}{3.205939in}}%
\pgfpathlineto{\pgfqpoint{2.807071in}{3.204132in}}%
\pgfpathlineto{\pgfqpoint{2.811114in}{3.200225in}}%
\pgfpathlineto{\pgfqpoint{2.813258in}{3.200456in}}%
\pgfpathlineto{\pgfqpoint{2.815565in}{3.202512in}}%
\pgfpathlineto{\pgfqpoint{2.819269in}{3.206109in}}%
\pgfpathlineto{\pgfqpoint{2.821412in}{3.205942in}}%
\pgfpathlineto{\pgfqpoint{2.823555in}{3.204132in}}%
\pgfpathlineto{\pgfqpoint{2.827599in}{3.200216in}}%
\pgfpathlineto{\pgfqpoint{2.829742in}{3.200443in}}%
\pgfpathlineto{\pgfqpoint{2.832050in}{3.202495in}}%
\pgfpathlineto{\pgfqpoint{2.835754in}{3.206085in}}%
\pgfpathlineto{\pgfqpoint{2.837897in}{3.205915in}}%
\pgfpathlineto{\pgfqpoint{2.840040in}{3.204101in}}%
\pgfpathlineto{\pgfqpoint{2.844084in}{3.200179in}}%
\pgfpathlineto{\pgfqpoint{2.846227in}{3.200404in}}%
\pgfpathlineto{\pgfqpoint{2.848535in}{3.202453in}}%
\pgfpathlineto{\pgfqpoint{2.852239in}{3.206039in}}%
\pgfpathlineto{\pgfqpoint{2.854382in}{3.205867in}}%
\pgfpathlineto{\pgfqpoint{2.856525in}{3.204051in}}%
\pgfpathlineto{\pgfqpoint{2.860569in}{3.200126in}}%
\pgfpathlineto{\pgfqpoint{2.862712in}{3.200349in}}%
\pgfpathlineto{\pgfqpoint{2.865020in}{3.202397in}}%
\pgfpathlineto{\pgfqpoint{2.868724in}{3.205981in}}%
\pgfpathlineto{\pgfqpoint{2.870867in}{3.205808in}}%
\pgfpathlineto{\pgfqpoint{2.873010in}{3.203992in}}%
\pgfpathlineto{\pgfqpoint{2.877054in}{3.200066in}}%
\pgfpathlineto{\pgfqpoint{2.879197in}{3.200289in}}%
\pgfpathlineto{\pgfqpoint{2.881505in}{3.202337in}}%
\pgfpathlineto{\pgfqpoint{2.885209in}{3.205922in}}%
\pgfpathlineto{\pgfqpoint{2.887352in}{3.205749in}}%
\pgfpathlineto{\pgfqpoint{2.889495in}{3.203934in}}%
\pgfpathlineto{\pgfqpoint{2.893539in}{3.200010in}}%
\pgfpathlineto{\pgfqpoint{2.895682in}{3.200234in}}%
\pgfpathlineto{\pgfqpoint{2.897990in}{3.202283in}}%
\pgfpathlineto{\pgfqpoint{2.901694in}{3.205870in}}%
\pgfpathlineto{\pgfqpoint{2.903837in}{3.205699in}}%
\pgfpathlineto{\pgfqpoint{2.905980in}{3.203885in}}%
\pgfpathlineto{\pgfqpoint{2.910024in}{3.199964in}}%
\pgfpathlineto{\pgfqpoint{2.912167in}{3.200190in}}%
\pgfpathlineto{\pgfqpoint{2.914475in}{3.202242in}}%
\pgfpathlineto{\pgfqpoint{2.918178in}{3.205832in}}%
\pgfpathlineto{\pgfqpoint{2.920321in}{3.205663in}}%
\pgfpathlineto{\pgfqpoint{2.922464in}{3.203851in}}%
\pgfpathlineto{\pgfqpoint{2.926509in}{3.199934in}}%
\pgfpathlineto{\pgfqpoint{2.928652in}{3.200162in}}%
\pgfpathlineto{\pgfqpoint{2.930960in}{3.202216in}}%
\pgfpathlineto{\pgfqpoint{2.934663in}{3.205810in}}%
\pgfpathlineto{\pgfqpoint{2.936806in}{3.205643in}}%
\pgfpathlineto{\pgfqpoint{2.938949in}{3.203833in}}%
\pgfpathlineto{\pgfqpoint{2.942994in}{3.199921in}}%
\pgfpathlineto{\pgfqpoint{2.945137in}{3.200151in}}%
\pgfpathlineto{\pgfqpoint{2.947444in}{3.202207in}}%
\pgfpathlineto{\pgfqpoint{2.951148in}{3.205805in}}%
\pgfpathlineto{\pgfqpoint{2.953291in}{3.205640in}}%
\pgfpathlineto{\pgfqpoint{2.955434in}{3.203832in}}%
\pgfpathlineto{\pgfqpoint{2.959478in}{3.199924in}}%
\pgfpathlineto{\pgfqpoint{2.961621in}{3.200155in}}%
\pgfpathlineto{\pgfqpoint{2.963929in}{3.202213in}}%
\pgfpathlineto{\pgfqpoint{2.967633in}{3.205814in}}%
\pgfpathlineto{\pgfqpoint{2.969776in}{3.205651in}}%
\pgfpathlineto{\pgfqpoint{2.971919in}{3.203845in}}%
\pgfpathlineto{\pgfqpoint{2.975963in}{3.199938in}}%
\pgfpathlineto{\pgfqpoint{2.978106in}{3.200172in}}%
\pgfpathlineto{\pgfqpoint{2.980414in}{3.202231in}}%
\pgfpathlineto{\pgfqpoint{2.984118in}{3.205833in}}%
\pgfpathlineto{\pgfqpoint{2.986261in}{3.205671in}}%
\pgfpathlineto{\pgfqpoint{2.988404in}{3.203866in}}%
\pgfpathlineto{\pgfqpoint{2.992448in}{3.199961in}}%
\pgfpathlineto{\pgfqpoint{2.994591in}{3.200195in}}%
\pgfpathlineto{\pgfqpoint{2.996899in}{3.202255in}}%
\pgfpathlineto{\pgfqpoint{3.000603in}{3.205858in}}%
\pgfpathlineto{\pgfqpoint{3.002746in}{3.205697in}}%
\pgfpathlineto{\pgfqpoint{3.004889in}{3.203892in}}%
\pgfpathlineto{\pgfqpoint{3.008933in}{3.199988in}}%
\pgfpathlineto{\pgfqpoint{3.011076in}{3.200221in}}%
\pgfpathlineto{\pgfqpoint{3.013384in}{3.202281in}}%
\pgfpathlineto{\pgfqpoint{3.017088in}{3.205885in}}%
\pgfpathlineto{\pgfqpoint{3.019231in}{3.205723in}}%
\pgfpathlineto{\pgfqpoint{3.021374in}{3.203917in}}%
\pgfpathlineto{\pgfqpoint{3.025418in}{3.200013in}}%
\pgfpathlineto{\pgfqpoint{3.027561in}{3.200246in}}%
\pgfpathlineto{\pgfqpoint{3.029869in}{3.202305in}}%
\pgfpathlineto{\pgfqpoint{3.033572in}{3.205908in}}%
\pgfpathlineto{\pgfqpoint{3.035715in}{3.205745in}}%
\pgfpathlineto{\pgfqpoint{3.037858in}{3.203939in}}%
\pgfpathlineto{\pgfqpoint{3.041902in}{3.200033in}}%
\pgfpathlineto{\pgfqpoint{3.044045in}{3.200266in}}%
\pgfpathlineto{\pgfqpoint{3.046353in}{3.202324in}}%
\pgfpathlineto{\pgfqpoint{3.050057in}{3.205925in}}%
\pgfpathlineto{\pgfqpoint{3.052200in}{3.205762in}}%
\pgfpathlineto{\pgfqpoint{3.054343in}{3.203955in}}%
\pgfpathlineto{\pgfqpoint{3.058387in}{3.200047in}}%
\pgfpathlineto{\pgfqpoint{3.060530in}{3.200279in}}%
\pgfpathlineto{\pgfqpoint{3.062838in}{3.202336in}}%
\pgfpathlineto{\pgfqpoint{3.066542in}{3.205936in}}%
\pgfpathlineto{\pgfqpoint{3.068685in}{3.205771in}}%
\pgfpathlineto{\pgfqpoint{3.070828in}{3.203963in}}%
\pgfpathlineto{\pgfqpoint{3.074872in}{3.200054in}}%
\pgfpathlineto{\pgfqpoint{3.077015in}{3.200285in}}%
\pgfpathlineto{\pgfqpoint{3.079323in}{3.202341in}}%
\pgfpathlineto{\pgfqpoint{3.083027in}{3.205939in}}%
\pgfpathlineto{\pgfqpoint{3.085170in}{3.205774in}}%
\pgfpathlineto{\pgfqpoint{3.087313in}{3.203965in}}%
\pgfpathlineto{\pgfqpoint{3.091357in}{3.200053in}}%
\pgfpathlineto{\pgfqpoint{3.093500in}{3.200283in}}%
\pgfpathlineto{\pgfqpoint{3.095808in}{3.202339in}}%
\pgfpathlineto{\pgfqpoint{3.099512in}{3.205935in}}%
\pgfpathlineto{\pgfqpoint{3.101655in}{3.205769in}}%
\pgfpathlineto{\pgfqpoint{3.103798in}{3.203960in}}%
\pgfpathlineto{\pgfqpoint{3.107842in}{3.200047in}}%
\pgfpathlineto{\pgfqpoint{3.109985in}{3.200277in}}%
\pgfpathlineto{\pgfqpoint{3.112293in}{3.202332in}}%
\pgfpathlineto{\pgfqpoint{3.115997in}{3.205927in}}%
\pgfpathlineto{\pgfqpoint{3.118140in}{3.205761in}}%
\pgfpathlineto{\pgfqpoint{3.120283in}{3.203951in}}%
\pgfpathlineto{\pgfqpoint{3.124327in}{3.200037in}}%
\pgfpathlineto{\pgfqpoint{3.126470in}{3.200267in}}%
\pgfpathlineto{\pgfqpoint{3.128778in}{3.202321in}}%
\pgfpathlineto{\pgfqpoint{3.132481in}{3.205916in}}%
\pgfpathlineto{\pgfqpoint{3.134624in}{3.205750in}}%
\pgfpathlineto{\pgfqpoint{3.136768in}{3.203940in}}%
\pgfpathlineto{\pgfqpoint{3.140812in}{3.200026in}}%
\pgfpathlineto{\pgfqpoint{3.142955in}{3.200255in}}%
\pgfpathlineto{\pgfqpoint{3.145263in}{3.202310in}}%
\pgfpathlineto{\pgfqpoint{3.148966in}{3.205905in}}%
\pgfpathlineto{\pgfqpoint{3.151109in}{3.205738in}}%
\pgfpathlineto{\pgfqpoint{3.153252in}{3.203928in}}%
\pgfpathlineto{\pgfqpoint{3.157296in}{3.200015in}}%
\pgfpathlineto{\pgfqpoint{3.159440in}{3.200244in}}%
\pgfpathlineto{\pgfqpoint{3.161747in}{3.202299in}}%
\pgfpathlineto{\pgfqpoint{3.165451in}{3.205894in}}%
\pgfpathlineto{\pgfqpoint{3.167594in}{3.205728in}}%
\pgfpathlineto{\pgfqpoint{3.169737in}{3.203918in}}%
\pgfpathlineto{\pgfqpoint{3.170132in}{3.203409in}}%
\pgfpathlineto{\pgfqpoint{3.170132in}{3.203409in}}%
\pgfusepath{stroke}%
\end{pgfscope}%
\begin{pgfscope}%
\pgfsetrectcap%
\pgfsetmiterjoin%
\pgfsetlinewidth{0.803000pt}%
\definecolor{currentstroke}{rgb}{0.000000,0.000000,0.000000}%
\pgfsetstrokecolor{currentstroke}%
\pgfsetdash{}{0pt}%
\pgfpathmoveto{\pgfqpoint{0.573768in}{2.601404in}}%
\pgfpathlineto{\pgfqpoint{0.573768in}{3.507654in}}%
\pgfusepath{stroke}%
\end{pgfscope}%
\begin{pgfscope}%
\pgfsetrectcap%
\pgfsetmiterjoin%
\pgfsetlinewidth{0.803000pt}%
\definecolor{currentstroke}{rgb}{0.000000,0.000000,0.000000}%
\pgfsetstrokecolor{currentstroke}%
\pgfsetdash{}{0pt}%
\pgfpathmoveto{\pgfqpoint{3.293768in}{2.601404in}}%
\pgfpathlineto{\pgfqpoint{3.293768in}{3.507654in}}%
\pgfusepath{stroke}%
\end{pgfscope}%
\begin{pgfscope}%
\pgfsetrectcap%
\pgfsetmiterjoin%
\pgfsetlinewidth{0.803000pt}%
\definecolor{currentstroke}{rgb}{0.000000,0.000000,0.000000}%
\pgfsetstrokecolor{currentstroke}%
\pgfsetdash{}{0pt}%
\pgfpathmoveto{\pgfqpoint{0.573768in}{2.601404in}}%
\pgfpathlineto{\pgfqpoint{3.293768in}{2.601404in}}%
\pgfusepath{stroke}%
\end{pgfscope}%
\begin{pgfscope}%
\pgfsetrectcap%
\pgfsetmiterjoin%
\pgfsetlinewidth{0.803000pt}%
\definecolor{currentstroke}{rgb}{0.000000,0.000000,0.000000}%
\pgfsetstrokecolor{currentstroke}%
\pgfsetdash{}{0pt}%
\pgfpathmoveto{\pgfqpoint{0.573768in}{3.507654in}}%
\pgfpathlineto{\pgfqpoint{3.293768in}{3.507654in}}%
\pgfusepath{stroke}%
\end{pgfscope}%
\begin{pgfscope}%
\definecolor{textcolor}{rgb}{0.000000,0.000000,0.000000}%
\pgfsetstrokecolor{textcolor}%
\pgfsetfillcolor{textcolor}%
\pgftext[x=0.628168in,y=3.444217in,left,top]{\color{textcolor}{\rmfamily\fontsize{8.000000}{9.600000}\selectfont\catcode`\^=\active\def^{\ifmmode\sp\else\^{}\fi}\catcode`\%=\active\def%{\%}(a)}}%
\end{pgfscope}%
\begin{pgfscope}%
\pgfsetbuttcap%
\pgfsetmiterjoin%
\definecolor{currentfill}{rgb}{1.000000,1.000000,1.000000}%
\pgfsetfillcolor{currentfill}%
\pgfsetlinewidth{0.000000pt}%
\definecolor{currentstroke}{rgb}{0.000000,0.000000,0.000000}%
\pgfsetstrokecolor{currentstroke}%
\pgfsetstrokeopacity{0.000000}%
\pgfsetdash{}{0pt}%
\pgfpathmoveto{\pgfqpoint{0.573768in}{1.532029in}}%
\pgfpathlineto{\pgfqpoint{3.293768in}{1.532029in}}%
\pgfpathlineto{\pgfqpoint{3.293768in}{2.438279in}}%
\pgfpathlineto{\pgfqpoint{0.573768in}{2.438279in}}%
\pgfpathlineto{\pgfqpoint{0.573768in}{1.532029in}}%
\pgfpathclose%
\pgfusepath{fill}%
\end{pgfscope}%
\begin{pgfscope}%
\pgfpathrectangle{\pgfqpoint{0.573768in}{1.532029in}}{\pgfqpoint{2.720000in}{0.906250in}}%
\pgfusepath{clip}%
\pgfsetbuttcap%
\pgfsetroundjoin%
\pgfsetlinewidth{0.501875pt}%
\definecolor{currentstroke}{rgb}{0.690196,0.690196,0.690196}%
\pgfsetstrokecolor{currentstroke}%
\pgfsetstrokeopacity{0.250000}%
\pgfsetdash{{1.850000pt}{0.800000pt}}{0.000000pt}%
\pgfpathmoveto{\pgfqpoint{0.697405in}{1.532029in}}%
\pgfpathlineto{\pgfqpoint{0.697405in}{2.438279in}}%
\pgfusepath{stroke}%
\end{pgfscope}%
\begin{pgfscope}%
\pgfsetbuttcap%
\pgfsetroundjoin%
\definecolor{currentfill}{rgb}{0.000000,0.000000,0.000000}%
\pgfsetfillcolor{currentfill}%
\pgfsetlinewidth{0.803000pt}%
\definecolor{currentstroke}{rgb}{0.000000,0.000000,0.000000}%
\pgfsetstrokecolor{currentstroke}%
\pgfsetdash{}{0pt}%
\pgfsys@defobject{currentmarker}{\pgfqpoint{0.000000in}{-0.048611in}}{\pgfqpoint{0.000000in}{0.000000in}}{%
\pgfpathmoveto{\pgfqpoint{0.000000in}{0.000000in}}%
\pgfpathlineto{\pgfqpoint{0.000000in}{-0.048611in}}%
\pgfusepath{stroke,fill}%
}%
\begin{pgfscope}%
\pgfsys@transformshift{0.697405in}{1.532029in}%
\pgfsys@useobject{currentmarker}{}%
\end{pgfscope}%
\end{pgfscope}%
\begin{pgfscope}%
\pgfpathrectangle{\pgfqpoint{0.573768in}{1.532029in}}{\pgfqpoint{2.720000in}{0.906250in}}%
\pgfusepath{clip}%
\pgfsetbuttcap%
\pgfsetroundjoin%
\pgfsetlinewidth{0.501875pt}%
\definecolor{currentstroke}{rgb}{0.690196,0.690196,0.690196}%
\pgfsetstrokecolor{currentstroke}%
\pgfsetstrokeopacity{0.250000}%
\pgfsetdash{{1.850000pt}{0.800000pt}}{0.000000pt}%
\pgfpathmoveto{\pgfqpoint{1.109526in}{1.532029in}}%
\pgfpathlineto{\pgfqpoint{1.109526in}{2.438279in}}%
\pgfusepath{stroke}%
\end{pgfscope}%
\begin{pgfscope}%
\pgfsetbuttcap%
\pgfsetroundjoin%
\definecolor{currentfill}{rgb}{0.000000,0.000000,0.000000}%
\pgfsetfillcolor{currentfill}%
\pgfsetlinewidth{0.803000pt}%
\definecolor{currentstroke}{rgb}{0.000000,0.000000,0.000000}%
\pgfsetstrokecolor{currentstroke}%
\pgfsetdash{}{0pt}%
\pgfsys@defobject{currentmarker}{\pgfqpoint{0.000000in}{-0.048611in}}{\pgfqpoint{0.000000in}{0.000000in}}{%
\pgfpathmoveto{\pgfqpoint{0.000000in}{0.000000in}}%
\pgfpathlineto{\pgfqpoint{0.000000in}{-0.048611in}}%
\pgfusepath{stroke,fill}%
}%
\begin{pgfscope}%
\pgfsys@transformshift{1.109526in}{1.532029in}%
\pgfsys@useobject{currentmarker}{}%
\end{pgfscope}%
\end{pgfscope}%
\begin{pgfscope}%
\pgfpathrectangle{\pgfqpoint{0.573768in}{1.532029in}}{\pgfqpoint{2.720000in}{0.906250in}}%
\pgfusepath{clip}%
\pgfsetbuttcap%
\pgfsetroundjoin%
\pgfsetlinewidth{0.501875pt}%
\definecolor{currentstroke}{rgb}{0.690196,0.690196,0.690196}%
\pgfsetstrokecolor{currentstroke}%
\pgfsetstrokeopacity{0.250000}%
\pgfsetdash{{1.850000pt}{0.800000pt}}{0.000000pt}%
\pgfpathmoveto{\pgfqpoint{1.521647in}{1.532029in}}%
\pgfpathlineto{\pgfqpoint{1.521647in}{2.438279in}}%
\pgfusepath{stroke}%
\end{pgfscope}%
\begin{pgfscope}%
\pgfsetbuttcap%
\pgfsetroundjoin%
\definecolor{currentfill}{rgb}{0.000000,0.000000,0.000000}%
\pgfsetfillcolor{currentfill}%
\pgfsetlinewidth{0.803000pt}%
\definecolor{currentstroke}{rgb}{0.000000,0.000000,0.000000}%
\pgfsetstrokecolor{currentstroke}%
\pgfsetdash{}{0pt}%
\pgfsys@defobject{currentmarker}{\pgfqpoint{0.000000in}{-0.048611in}}{\pgfqpoint{0.000000in}{0.000000in}}{%
\pgfpathmoveto{\pgfqpoint{0.000000in}{0.000000in}}%
\pgfpathlineto{\pgfqpoint{0.000000in}{-0.048611in}}%
\pgfusepath{stroke,fill}%
}%
\begin{pgfscope}%
\pgfsys@transformshift{1.521647in}{1.532029in}%
\pgfsys@useobject{currentmarker}{}%
\end{pgfscope}%
\end{pgfscope}%
\begin{pgfscope}%
\pgfpathrectangle{\pgfqpoint{0.573768in}{1.532029in}}{\pgfqpoint{2.720000in}{0.906250in}}%
\pgfusepath{clip}%
\pgfsetbuttcap%
\pgfsetroundjoin%
\pgfsetlinewidth{0.501875pt}%
\definecolor{currentstroke}{rgb}{0.690196,0.690196,0.690196}%
\pgfsetstrokecolor{currentstroke}%
\pgfsetstrokeopacity{0.250000}%
\pgfsetdash{{1.850000pt}{0.800000pt}}{0.000000pt}%
\pgfpathmoveto{\pgfqpoint{1.933768in}{1.532029in}}%
\pgfpathlineto{\pgfqpoint{1.933768in}{2.438279in}}%
\pgfusepath{stroke}%
\end{pgfscope}%
\begin{pgfscope}%
\pgfsetbuttcap%
\pgfsetroundjoin%
\definecolor{currentfill}{rgb}{0.000000,0.000000,0.000000}%
\pgfsetfillcolor{currentfill}%
\pgfsetlinewidth{0.803000pt}%
\definecolor{currentstroke}{rgb}{0.000000,0.000000,0.000000}%
\pgfsetstrokecolor{currentstroke}%
\pgfsetdash{}{0pt}%
\pgfsys@defobject{currentmarker}{\pgfqpoint{0.000000in}{-0.048611in}}{\pgfqpoint{0.000000in}{0.000000in}}{%
\pgfpathmoveto{\pgfqpoint{0.000000in}{0.000000in}}%
\pgfpathlineto{\pgfqpoint{0.000000in}{-0.048611in}}%
\pgfusepath{stroke,fill}%
}%
\begin{pgfscope}%
\pgfsys@transformshift{1.933768in}{1.532029in}%
\pgfsys@useobject{currentmarker}{}%
\end{pgfscope}%
\end{pgfscope}%
\begin{pgfscope}%
\pgfpathrectangle{\pgfqpoint{0.573768in}{1.532029in}}{\pgfqpoint{2.720000in}{0.906250in}}%
\pgfusepath{clip}%
\pgfsetbuttcap%
\pgfsetroundjoin%
\pgfsetlinewidth{0.501875pt}%
\definecolor{currentstroke}{rgb}{0.690196,0.690196,0.690196}%
\pgfsetstrokecolor{currentstroke}%
\pgfsetstrokeopacity{0.250000}%
\pgfsetdash{{1.850000pt}{0.800000pt}}{0.000000pt}%
\pgfpathmoveto{\pgfqpoint{2.345890in}{1.532029in}}%
\pgfpathlineto{\pgfqpoint{2.345890in}{2.438279in}}%
\pgfusepath{stroke}%
\end{pgfscope}%
\begin{pgfscope}%
\pgfsetbuttcap%
\pgfsetroundjoin%
\definecolor{currentfill}{rgb}{0.000000,0.000000,0.000000}%
\pgfsetfillcolor{currentfill}%
\pgfsetlinewidth{0.803000pt}%
\definecolor{currentstroke}{rgb}{0.000000,0.000000,0.000000}%
\pgfsetstrokecolor{currentstroke}%
\pgfsetdash{}{0pt}%
\pgfsys@defobject{currentmarker}{\pgfqpoint{0.000000in}{-0.048611in}}{\pgfqpoint{0.000000in}{0.000000in}}{%
\pgfpathmoveto{\pgfqpoint{0.000000in}{0.000000in}}%
\pgfpathlineto{\pgfqpoint{0.000000in}{-0.048611in}}%
\pgfusepath{stroke,fill}%
}%
\begin{pgfscope}%
\pgfsys@transformshift{2.345890in}{1.532029in}%
\pgfsys@useobject{currentmarker}{}%
\end{pgfscope}%
\end{pgfscope}%
\begin{pgfscope}%
\pgfpathrectangle{\pgfqpoint{0.573768in}{1.532029in}}{\pgfqpoint{2.720000in}{0.906250in}}%
\pgfusepath{clip}%
\pgfsetbuttcap%
\pgfsetroundjoin%
\pgfsetlinewidth{0.501875pt}%
\definecolor{currentstroke}{rgb}{0.690196,0.690196,0.690196}%
\pgfsetstrokecolor{currentstroke}%
\pgfsetstrokeopacity{0.250000}%
\pgfsetdash{{1.850000pt}{0.800000pt}}{0.000000pt}%
\pgfpathmoveto{\pgfqpoint{2.758011in}{1.532029in}}%
\pgfpathlineto{\pgfqpoint{2.758011in}{2.438279in}}%
\pgfusepath{stroke}%
\end{pgfscope}%
\begin{pgfscope}%
\pgfsetbuttcap%
\pgfsetroundjoin%
\definecolor{currentfill}{rgb}{0.000000,0.000000,0.000000}%
\pgfsetfillcolor{currentfill}%
\pgfsetlinewidth{0.803000pt}%
\definecolor{currentstroke}{rgb}{0.000000,0.000000,0.000000}%
\pgfsetstrokecolor{currentstroke}%
\pgfsetdash{}{0pt}%
\pgfsys@defobject{currentmarker}{\pgfqpoint{0.000000in}{-0.048611in}}{\pgfqpoint{0.000000in}{0.000000in}}{%
\pgfpathmoveto{\pgfqpoint{0.000000in}{0.000000in}}%
\pgfpathlineto{\pgfqpoint{0.000000in}{-0.048611in}}%
\pgfusepath{stroke,fill}%
}%
\begin{pgfscope}%
\pgfsys@transformshift{2.758011in}{1.532029in}%
\pgfsys@useobject{currentmarker}{}%
\end{pgfscope}%
\end{pgfscope}%
\begin{pgfscope}%
\pgfpathrectangle{\pgfqpoint{0.573768in}{1.532029in}}{\pgfqpoint{2.720000in}{0.906250in}}%
\pgfusepath{clip}%
\pgfsetbuttcap%
\pgfsetroundjoin%
\pgfsetlinewidth{0.501875pt}%
\definecolor{currentstroke}{rgb}{0.690196,0.690196,0.690196}%
\pgfsetstrokecolor{currentstroke}%
\pgfsetstrokeopacity{0.250000}%
\pgfsetdash{{1.850000pt}{0.800000pt}}{0.000000pt}%
\pgfpathmoveto{\pgfqpoint{3.170132in}{1.532029in}}%
\pgfpathlineto{\pgfqpoint{3.170132in}{2.438279in}}%
\pgfusepath{stroke}%
\end{pgfscope}%
\begin{pgfscope}%
\pgfsetbuttcap%
\pgfsetroundjoin%
\definecolor{currentfill}{rgb}{0.000000,0.000000,0.000000}%
\pgfsetfillcolor{currentfill}%
\pgfsetlinewidth{0.803000pt}%
\definecolor{currentstroke}{rgb}{0.000000,0.000000,0.000000}%
\pgfsetstrokecolor{currentstroke}%
\pgfsetdash{}{0pt}%
\pgfsys@defobject{currentmarker}{\pgfqpoint{0.000000in}{-0.048611in}}{\pgfqpoint{0.000000in}{0.000000in}}{%
\pgfpathmoveto{\pgfqpoint{0.000000in}{0.000000in}}%
\pgfpathlineto{\pgfqpoint{0.000000in}{-0.048611in}}%
\pgfusepath{stroke,fill}%
}%
\begin{pgfscope}%
\pgfsys@transformshift{3.170132in}{1.532029in}%
\pgfsys@useobject{currentmarker}{}%
\end{pgfscope}%
\end{pgfscope}%
\begin{pgfscope}%
\pgfpathrectangle{\pgfqpoint{0.573768in}{1.532029in}}{\pgfqpoint{2.720000in}{0.906250in}}%
\pgfusepath{clip}%
\pgfsetbuttcap%
\pgfsetroundjoin%
\pgfsetlinewidth{0.501875pt}%
\definecolor{currentstroke}{rgb}{0.690196,0.690196,0.690196}%
\pgfsetstrokecolor{currentstroke}%
\pgfsetstrokeopacity{0.250000}%
\pgfsetdash{{1.850000pt}{0.800000pt}}{0.000000pt}%
\pgfpathmoveto{\pgfqpoint{0.573768in}{1.532029in}}%
\pgfpathlineto{\pgfqpoint{3.293768in}{1.532029in}}%
\pgfusepath{stroke}%
\end{pgfscope}%
\begin{pgfscope}%
\pgfsetbuttcap%
\pgfsetroundjoin%
\definecolor{currentfill}{rgb}{0.000000,0.000000,0.000000}%
\pgfsetfillcolor{currentfill}%
\pgfsetlinewidth{0.803000pt}%
\definecolor{currentstroke}{rgb}{0.000000,0.000000,0.000000}%
\pgfsetstrokecolor{currentstroke}%
\pgfsetdash{}{0pt}%
\pgfsys@defobject{currentmarker}{\pgfqpoint{-0.048611in}{0.000000in}}{\pgfqpoint{-0.000000in}{0.000000in}}{%
\pgfpathmoveto{\pgfqpoint{-0.000000in}{0.000000in}}%
\pgfpathlineto{\pgfqpoint{-0.048611in}{0.000000in}}%
\pgfusepath{stroke,fill}%
}%
\begin{pgfscope}%
\pgfsys@transformshift{0.573768in}{1.532029in}%
\pgfsys@useobject{currentmarker}{}%
\end{pgfscope}%
\end{pgfscope}%
\begin{pgfscope}%
\definecolor{textcolor}{rgb}{0.000000,0.000000,0.000000}%
\pgfsetstrokecolor{textcolor}%
\pgfsetfillcolor{textcolor}%
\pgftext[x=0.417518in, y=1.493449in, left, base]{\color{textcolor}{\rmfamily\fontsize{8.000000}{9.600000}\selectfont\catcode`\^=\active\def^{\ifmmode\sp\else\^{}\fi}\catcode`\%=\active\def%{\%}$\mathdefault{0}$}}%
\end{pgfscope}%
\begin{pgfscope}%
\pgfpathrectangle{\pgfqpoint{0.573768in}{1.532029in}}{\pgfqpoint{2.720000in}{0.906250in}}%
\pgfusepath{clip}%
\pgfsetbuttcap%
\pgfsetroundjoin%
\pgfsetlinewidth{0.501875pt}%
\definecolor{currentstroke}{rgb}{0.690196,0.690196,0.690196}%
\pgfsetstrokecolor{currentstroke}%
\pgfsetstrokeopacity{0.250000}%
\pgfsetdash{{1.850000pt}{0.800000pt}}{0.000000pt}%
\pgfpathmoveto{\pgfqpoint{0.573768in}{1.834113in}}%
\pgfpathlineto{\pgfqpoint{3.293768in}{1.834113in}}%
\pgfusepath{stroke}%
\end{pgfscope}%
\begin{pgfscope}%
\pgfsetbuttcap%
\pgfsetroundjoin%
\definecolor{currentfill}{rgb}{0.000000,0.000000,0.000000}%
\pgfsetfillcolor{currentfill}%
\pgfsetlinewidth{0.803000pt}%
\definecolor{currentstroke}{rgb}{0.000000,0.000000,0.000000}%
\pgfsetstrokecolor{currentstroke}%
\pgfsetdash{}{0pt}%
\pgfsys@defobject{currentmarker}{\pgfqpoint{-0.048611in}{0.000000in}}{\pgfqpoint{-0.000000in}{0.000000in}}{%
\pgfpathmoveto{\pgfqpoint{-0.000000in}{0.000000in}}%
\pgfpathlineto{\pgfqpoint{-0.048611in}{0.000000in}}%
\pgfusepath{stroke,fill}%
}%
\begin{pgfscope}%
\pgfsys@transformshift{0.573768in}{1.834113in}%
\pgfsys@useobject{currentmarker}{}%
\end{pgfscope}%
\end{pgfscope}%
\begin{pgfscope}%
\definecolor{textcolor}{rgb}{0.000000,0.000000,0.000000}%
\pgfsetstrokecolor{textcolor}%
\pgfsetfillcolor{textcolor}%
\pgftext[x=0.417518in, y=1.795533in, left, base]{\color{textcolor}{\rmfamily\fontsize{8.000000}{9.600000}\selectfont\catcode`\^=\active\def^{\ifmmode\sp\else\^{}\fi}\catcode`\%=\active\def%{\%}$\mathdefault{5}$}}%
\end{pgfscope}%
\begin{pgfscope}%
\pgfpathrectangle{\pgfqpoint{0.573768in}{1.532029in}}{\pgfqpoint{2.720000in}{0.906250in}}%
\pgfusepath{clip}%
\pgfsetbuttcap%
\pgfsetroundjoin%
\pgfsetlinewidth{0.501875pt}%
\definecolor{currentstroke}{rgb}{0.690196,0.690196,0.690196}%
\pgfsetstrokecolor{currentstroke}%
\pgfsetstrokeopacity{0.250000}%
\pgfsetdash{{1.850000pt}{0.800000pt}}{0.000000pt}%
\pgfpathmoveto{\pgfqpoint{0.573768in}{2.136196in}}%
\pgfpathlineto{\pgfqpoint{3.293768in}{2.136196in}}%
\pgfusepath{stroke}%
\end{pgfscope}%
\begin{pgfscope}%
\pgfsetbuttcap%
\pgfsetroundjoin%
\definecolor{currentfill}{rgb}{0.000000,0.000000,0.000000}%
\pgfsetfillcolor{currentfill}%
\pgfsetlinewidth{0.803000pt}%
\definecolor{currentstroke}{rgb}{0.000000,0.000000,0.000000}%
\pgfsetstrokecolor{currentstroke}%
\pgfsetdash{}{0pt}%
\pgfsys@defobject{currentmarker}{\pgfqpoint{-0.048611in}{0.000000in}}{\pgfqpoint{-0.000000in}{0.000000in}}{%
\pgfpathmoveto{\pgfqpoint{-0.000000in}{0.000000in}}%
\pgfpathlineto{\pgfqpoint{-0.048611in}{0.000000in}}%
\pgfusepath{stroke,fill}%
}%
\begin{pgfscope}%
\pgfsys@transformshift{0.573768in}{2.136196in}%
\pgfsys@useobject{currentmarker}{}%
\end{pgfscope}%
\end{pgfscope}%
\begin{pgfscope}%
\definecolor{textcolor}{rgb}{0.000000,0.000000,0.000000}%
\pgfsetstrokecolor{textcolor}%
\pgfsetfillcolor{textcolor}%
\pgftext[x=0.358489in, y=2.097616in, left, base]{\color{textcolor}{\rmfamily\fontsize{8.000000}{9.600000}\selectfont\catcode`\^=\active\def^{\ifmmode\sp\else\^{}\fi}\catcode`\%=\active\def%{\%}$\mathdefault{10}$}}%
\end{pgfscope}%
\begin{pgfscope}%
\pgfpathrectangle{\pgfqpoint{0.573768in}{1.532029in}}{\pgfqpoint{2.720000in}{0.906250in}}%
\pgfusepath{clip}%
\pgfsetbuttcap%
\pgfsetroundjoin%
\pgfsetlinewidth{0.501875pt}%
\definecolor{currentstroke}{rgb}{0.690196,0.690196,0.690196}%
\pgfsetstrokecolor{currentstroke}%
\pgfsetstrokeopacity{0.250000}%
\pgfsetdash{{1.850000pt}{0.800000pt}}{0.000000pt}%
\pgfpathmoveto{\pgfqpoint{0.573768in}{2.438279in}}%
\pgfpathlineto{\pgfqpoint{3.293768in}{2.438279in}}%
\pgfusepath{stroke}%
\end{pgfscope}%
\begin{pgfscope}%
\pgfsetbuttcap%
\pgfsetroundjoin%
\definecolor{currentfill}{rgb}{0.000000,0.000000,0.000000}%
\pgfsetfillcolor{currentfill}%
\pgfsetlinewidth{0.803000pt}%
\definecolor{currentstroke}{rgb}{0.000000,0.000000,0.000000}%
\pgfsetstrokecolor{currentstroke}%
\pgfsetdash{}{0pt}%
\pgfsys@defobject{currentmarker}{\pgfqpoint{-0.048611in}{0.000000in}}{\pgfqpoint{-0.000000in}{0.000000in}}{%
\pgfpathmoveto{\pgfqpoint{-0.000000in}{0.000000in}}%
\pgfpathlineto{\pgfqpoint{-0.048611in}{0.000000in}}%
\pgfusepath{stroke,fill}%
}%
\begin{pgfscope}%
\pgfsys@transformshift{0.573768in}{2.438279in}%
\pgfsys@useobject{currentmarker}{}%
\end{pgfscope}%
\end{pgfscope}%
\begin{pgfscope}%
\definecolor{textcolor}{rgb}{0.000000,0.000000,0.000000}%
\pgfsetstrokecolor{textcolor}%
\pgfsetfillcolor{textcolor}%
\pgftext[x=0.358489in, y=2.399699in, left, base]{\color{textcolor}{\rmfamily\fontsize{8.000000}{9.600000}\selectfont\catcode`\^=\active\def^{\ifmmode\sp\else\^{}\fi}\catcode`\%=\active\def%{\%}$\mathdefault{15}$}}%
\end{pgfscope}%
\begin{pgfscope}%
\definecolor{textcolor}{rgb}{0.000000,0.000000,0.000000}%
\pgfsetstrokecolor{textcolor}%
\pgfsetfillcolor{textcolor}%
\pgftext[x=0.302933in,y=1.985154in,,bottom,rotate=90.000000]{\color{textcolor}{\rmfamily\fontsize{8.000000}{9.600000}\selectfont\catcode`\^=\active\def^{\ifmmode\sp\else\^{}\fi}\catcode`\%=\active\def%{\%}$i_L$ (A)}}%
\end{pgfscope}%
\begin{pgfscope}%
\pgfpathrectangle{\pgfqpoint{0.573768in}{1.532029in}}{\pgfqpoint{2.720000in}{0.906250in}}%
\pgfusepath{clip}%
\pgfsetrectcap%
\pgfsetroundjoin%
\pgfsetlinewidth{1.505625pt}%
\definecolor{currentstroke}{rgb}{0.835294,0.368627,0.000000}%
\pgfsetstrokecolor{currentstroke}%
\pgfsetdash{}{0pt}%
\pgfpathmoveto{\pgfqpoint{0.697405in}{1.532029in}}%
\pgfpathlineto{\pgfqpoint{0.697762in}{1.532615in}}%
\pgfpathlineto{\pgfqpoint{0.700010in}{1.597964in}}%
\pgfpathlineto{\pgfqpoint{0.706154in}{1.774190in}}%
\pgfpathlineto{\pgfqpoint{0.706688in}{1.773690in}}%
\pgfpathlineto{\pgfqpoint{0.710809in}{1.768570in}}%
\pgfpathlineto{\pgfqpoint{0.714214in}{1.762740in}}%
\pgfpathlineto{\pgfqpoint{0.714305in}{1.764161in}}%
\pgfpathlineto{\pgfqpoint{0.718952in}{1.891034in}}%
\pgfpathlineto{\pgfqpoint{0.722634in}{1.986798in}}%
\pgfpathlineto{\pgfqpoint{0.723147in}{1.985013in}}%
\pgfpathlineto{\pgfqpoint{0.728258in}{1.964089in}}%
\pgfpathlineto{\pgfqpoint{0.730704in}{1.952180in}}%
\pgfpathlineto{\pgfqpoint{0.730853in}{1.954516in}}%
\pgfpathlineto{\pgfqpoint{0.739119in}{2.146876in}}%
\pgfpathlineto{\pgfqpoint{0.739635in}{2.143167in}}%
\pgfpathlineto{\pgfqpoint{0.746229in}{2.089683in}}%
\pgfpathlineto{\pgfqpoint{0.747214in}{2.080978in}}%
\pgfpathlineto{\pgfqpoint{0.747452in}{2.084941in}}%
\pgfpathlineto{\pgfqpoint{0.755542in}{2.241464in}}%
\pgfpathlineto{\pgfqpoint{0.755735in}{2.239736in}}%
\pgfpathlineto{\pgfqpoint{0.763698in}{2.141851in}}%
\pgfpathlineto{\pgfqpoint{0.764324in}{2.151106in}}%
\pgfpathlineto{\pgfqpoint{0.772027in}{2.268838in}}%
\pgfpathlineto{\pgfqpoint{0.772156in}{2.267322in}}%
\pgfpathlineto{\pgfqpoint{0.780195in}{2.138544in}}%
\pgfpathlineto{\pgfqpoint{0.780804in}{2.145507in}}%
\pgfpathlineto{\pgfqpoint{0.788511in}{2.237365in}}%
\pgfpathlineto{\pgfqpoint{0.788654in}{2.235277in}}%
\pgfpathlineto{\pgfqpoint{0.796695in}{2.083322in}}%
\pgfpathlineto{\pgfqpoint{0.797285in}{2.088658in}}%
\pgfpathlineto{\pgfqpoint{0.804996in}{2.162375in}}%
\pgfpathlineto{\pgfqpoint{0.805140in}{2.159961in}}%
\pgfpathlineto{\pgfqpoint{0.813181in}{1.993552in}}%
\pgfpathlineto{\pgfqpoint{0.813745in}{1.997824in}}%
\pgfpathlineto{\pgfqpoint{0.821395in}{2.062542in}}%
\pgfpathlineto{\pgfqpoint{0.821623in}{2.060005in}}%
\pgfpathlineto{\pgfqpoint{0.829671in}{1.888486in}}%
\pgfpathlineto{\pgfqpoint{0.830125in}{1.891737in}}%
\pgfpathlineto{\pgfqpoint{0.837964in}{1.956716in}}%
\pgfpathlineto{\pgfqpoint{0.838201in}{1.952278in}}%
\pgfpathlineto{\pgfqpoint{0.846141in}{1.786040in}}%
\pgfpathlineto{\pgfqpoint{0.846564in}{1.789226in}}%
\pgfpathlineto{\pgfqpoint{0.854449in}{1.861266in}}%
\pgfpathlineto{\pgfqpoint{0.854687in}{1.857050in}}%
\pgfpathlineto{\pgfqpoint{0.862612in}{1.700430in}}%
\pgfpathlineto{\pgfqpoint{0.862948in}{1.703200in}}%
\pgfpathlineto{\pgfqpoint{0.870900in}{1.787897in}}%
\pgfpathlineto{\pgfqpoint{0.871200in}{1.783077in}}%
\pgfpathlineto{\pgfqpoint{0.879088in}{1.640327in}}%
\pgfpathlineto{\pgfqpoint{0.879401in}{1.643334in}}%
\pgfpathlineto{\pgfqpoint{0.887441in}{1.742823in}}%
\pgfpathlineto{\pgfqpoint{0.887761in}{1.737985in}}%
\pgfpathlineto{\pgfqpoint{0.895569in}{1.610966in}}%
\pgfpathlineto{\pgfqpoint{0.895865in}{1.614319in}}%
\pgfpathlineto{\pgfqpoint{0.903902in}{1.728358in}}%
\pgfpathlineto{\pgfqpoint{0.904378in}{1.721840in}}%
\pgfpathlineto{\pgfqpoint{0.912054in}{1.609905in}}%
\pgfpathlineto{\pgfqpoint{0.912353in}{1.613777in}}%
\pgfpathlineto{\pgfqpoint{0.920387in}{1.739994in}}%
\pgfpathlineto{\pgfqpoint{0.920863in}{1.734145in}}%
\pgfpathlineto{\pgfqpoint{0.928534in}{1.632296in}}%
\pgfpathlineto{\pgfqpoint{0.928829in}{1.636380in}}%
\pgfpathlineto{\pgfqpoint{0.936872in}{1.771487in}}%
\pgfpathlineto{\pgfqpoint{0.937349in}{1.766075in}}%
\pgfpathlineto{\pgfqpoint{0.945019in}{1.670610in}}%
\pgfpathlineto{\pgfqpoint{0.945334in}{1.675163in}}%
\pgfpathlineto{\pgfqpoint{0.953317in}{1.814248in}}%
\pgfpathlineto{\pgfqpoint{0.953753in}{1.809745in}}%
\pgfpathlineto{\pgfqpoint{0.961509in}{1.715748in}}%
\pgfpathlineto{\pgfqpoint{0.961789in}{1.719851in}}%
\pgfpathlineto{\pgfqpoint{0.969820in}{1.860365in}}%
\pgfpathlineto{\pgfqpoint{0.970281in}{1.855428in}}%
\pgfpathlineto{\pgfqpoint{0.978001in}{1.760918in}}%
\pgfpathlineto{\pgfqpoint{0.978355in}{1.766278in}}%
\pgfpathlineto{\pgfqpoint{0.986328in}{1.902894in}}%
\pgfpathlineto{\pgfqpoint{0.986649in}{1.899432in}}%
\pgfpathlineto{\pgfqpoint{0.994486in}{1.799605in}}%
\pgfpathlineto{\pgfqpoint{0.994809in}{1.804232in}}%
\pgfpathlineto{\pgfqpoint{1.002813in}{1.936388in}}%
\pgfpathlineto{\pgfqpoint{1.003134in}{1.932701in}}%
\pgfpathlineto{\pgfqpoint{1.010971in}{1.827242in}}%
\pgfpathlineto{\pgfqpoint{1.011329in}{1.832167in}}%
\pgfpathlineto{\pgfqpoint{1.019298in}{1.957561in}}%
\pgfpathlineto{\pgfqpoint{1.019619in}{1.953616in}}%
\pgfpathlineto{\pgfqpoint{1.027456in}{1.841898in}}%
\pgfpathlineto{\pgfqpoint{1.027874in}{1.847431in}}%
\pgfpathlineto{\pgfqpoint{1.035783in}{1.965630in}}%
\pgfpathlineto{\pgfqpoint{1.036104in}{1.961434in}}%
\pgfpathlineto{\pgfqpoint{1.043941in}{1.843843in}}%
\pgfpathlineto{\pgfqpoint{1.044365in}{1.849158in}}%
\pgfpathlineto{\pgfqpoint{1.052267in}{1.961852in}}%
\pgfpathlineto{\pgfqpoint{1.052589in}{1.957449in}}%
\pgfpathlineto{\pgfqpoint{1.060426in}{1.835141in}}%
\pgfpathlineto{\pgfqpoint{1.060811in}{1.839702in}}%
\pgfpathlineto{\pgfqpoint{1.068752in}{1.948953in}}%
\pgfpathlineto{\pgfqpoint{1.069074in}{1.944407in}}%
\pgfpathlineto{\pgfqpoint{1.076910in}{1.819002in}}%
\pgfpathlineto{\pgfqpoint{1.077261in}{1.822970in}}%
\pgfpathlineto{\pgfqpoint{1.085237in}{1.930446in}}%
\pgfpathlineto{\pgfqpoint{1.085558in}{1.925832in}}%
\pgfpathlineto{\pgfqpoint{1.093395in}{1.799098in}}%
\pgfpathlineto{\pgfqpoint{1.093814in}{1.803946in}}%
\pgfpathlineto{\pgfqpoint{1.101722in}{1.910026in}}%
\pgfpathlineto{\pgfqpoint{1.102043in}{1.905424in}}%
\pgfpathlineto{\pgfqpoint{1.109880in}{1.778990in}}%
\pgfpathlineto{\pgfqpoint{1.110226in}{1.782904in}}%
\pgfpathlineto{\pgfqpoint{1.118231in}{1.891012in}}%
\pgfpathlineto{\pgfqpoint{1.118559in}{1.886266in}}%
\pgfpathlineto{\pgfqpoint{1.126365in}{1.761928in}}%
\pgfpathlineto{\pgfqpoint{1.126718in}{1.766030in}}%
\pgfpathlineto{\pgfqpoint{1.134691in}{1.876083in}}%
\pgfpathlineto{\pgfqpoint{1.135012in}{1.871623in}}%
\pgfpathlineto{\pgfqpoint{1.142839in}{1.749370in}}%
\pgfpathlineto{\pgfqpoint{1.143196in}{1.753565in}}%
\pgfpathlineto{\pgfqpoint{1.151166in}{1.866352in}}%
\pgfpathlineto{\pgfqpoint{1.151472in}{1.862275in}}%
\pgfpathlineto{\pgfqpoint{1.159330in}{1.742160in}}%
\pgfpathlineto{\pgfqpoint{1.159710in}{1.746830in}}%
\pgfpathlineto{\pgfqpoint{1.167661in}{1.861981in}}%
\pgfpathlineto{\pgfqpoint{1.167982in}{1.857744in}}%
\pgfpathlineto{\pgfqpoint{1.175815in}{1.740620in}}%
\pgfpathlineto{\pgfqpoint{1.176210in}{1.745639in}}%
\pgfpathlineto{\pgfqpoint{1.184146in}{1.863005in}}%
\pgfpathlineto{\pgfqpoint{1.184466in}{1.858861in}}%
\pgfpathlineto{\pgfqpoint{1.192300in}{1.743875in}}%
\pgfpathlineto{\pgfqpoint{1.192648in}{1.748293in}}%
\pgfpathlineto{\pgfqpoint{1.200630in}{1.868200in}}%
\pgfpathlineto{\pgfqpoint{1.200951in}{1.864122in}}%
\pgfpathlineto{\pgfqpoint{1.208778in}{1.750771in}}%
\pgfpathlineto{\pgfqpoint{1.209136in}{1.755355in}}%
\pgfpathlineto{\pgfqpoint{1.217109in}{1.876282in}}%
\pgfpathlineto{\pgfqpoint{1.217418in}{1.872436in}}%
\pgfpathlineto{\pgfqpoint{1.225274in}{1.759302in}}%
\pgfpathlineto{\pgfqpoint{1.225607in}{1.763618in}}%
\pgfpathlineto{\pgfqpoint{1.233600in}{1.885074in}}%
\pgfpathlineto{\pgfqpoint{1.233921in}{1.881032in}}%
\pgfpathlineto{\pgfqpoint{1.241759in}{1.768112in}}%
\pgfpathlineto{\pgfqpoint{1.242122in}{1.772876in}}%
\pgfpathlineto{\pgfqpoint{1.250085in}{1.893501in}}%
\pgfpathlineto{\pgfqpoint{1.250406in}{1.889437in}}%
\pgfpathlineto{\pgfqpoint{1.258244in}{1.775879in}}%
\pgfpathlineto{\pgfqpoint{1.258604in}{1.780557in}}%
\pgfpathlineto{\pgfqpoint{1.266570in}{1.900361in}}%
\pgfpathlineto{\pgfqpoint{1.266891in}{1.896257in}}%
\pgfpathlineto{\pgfqpoint{1.274729in}{1.781678in}}%
\pgfpathlineto{\pgfqpoint{1.275073in}{1.786061in}}%
\pgfpathlineto{\pgfqpoint{1.283055in}{1.904953in}}%
\pgfpathlineto{\pgfqpoint{1.283376in}{1.900799in}}%
\pgfpathlineto{\pgfqpoint{1.291214in}{1.785032in}}%
\pgfpathlineto{\pgfqpoint{1.291551in}{1.789268in}}%
\pgfpathlineto{\pgfqpoint{1.299540in}{1.907038in}}%
\pgfpathlineto{\pgfqpoint{1.299861in}{1.902836in}}%
\pgfpathlineto{\pgfqpoint{1.307698in}{1.785916in}}%
\pgfpathlineto{\pgfqpoint{1.308036in}{1.790094in}}%
\pgfpathlineto{\pgfqpoint{1.316025in}{1.906779in}}%
\pgfpathlineto{\pgfqpoint{1.316346in}{1.902535in}}%
\pgfpathlineto{\pgfqpoint{1.324183in}{1.784655in}}%
\pgfpathlineto{\pgfqpoint{1.324521in}{1.788808in}}%
\pgfpathlineto{\pgfqpoint{1.332510in}{1.904643in}}%
\pgfpathlineto{\pgfqpoint{1.332831in}{1.900369in}}%
\pgfpathlineto{\pgfqpoint{1.340668in}{1.781826in}}%
\pgfpathlineto{\pgfqpoint{1.341012in}{1.786041in}}%
\pgfpathlineto{\pgfqpoint{1.348994in}{1.901283in}}%
\pgfpathlineto{\pgfqpoint{1.349315in}{1.896993in}}%
\pgfpathlineto{\pgfqpoint{1.357153in}{1.778122in}}%
\pgfpathlineto{\pgfqpoint{1.357512in}{1.782548in}}%
\pgfpathlineto{\pgfqpoint{1.365479in}{1.897410in}}%
\pgfpathlineto{\pgfqpoint{1.365800in}{1.893118in}}%
\pgfpathlineto{\pgfqpoint{1.373638in}{1.774241in}}%
\pgfpathlineto{\pgfqpoint{1.373998in}{1.778700in}}%
\pgfpathlineto{\pgfqpoint{1.381964in}{1.893670in}}%
\pgfpathlineto{\pgfqpoint{1.382285in}{1.889386in}}%
\pgfpathlineto{\pgfqpoint{1.390123in}{1.770766in}}%
\pgfpathlineto{\pgfqpoint{1.390485in}{1.775270in}}%
\pgfpathlineto{\pgfqpoint{1.398449in}{1.890575in}}%
\pgfpathlineto{\pgfqpoint{1.398770in}{1.886308in}}%
\pgfpathlineto{\pgfqpoint{1.406607in}{1.768123in}}%
\pgfpathlineto{\pgfqpoint{1.406970in}{1.772654in}}%
\pgfpathlineto{\pgfqpoint{1.414934in}{1.888448in}}%
\pgfpathlineto{\pgfqpoint{1.415255in}{1.884203in}}%
\pgfpathlineto{\pgfqpoint{1.423092in}{1.766535in}}%
\pgfpathlineto{\pgfqpoint{1.423456in}{1.771105in}}%
\pgfpathlineto{\pgfqpoint{1.431418in}{1.887419in}}%
\pgfpathlineto{\pgfqpoint{1.431740in}{1.883195in}}%
\pgfpathlineto{\pgfqpoint{1.439577in}{1.766039in}}%
\pgfpathlineto{\pgfqpoint{1.439941in}{1.770635in}}%
\pgfpathlineto{\pgfqpoint{1.447903in}{1.887433in}}%
\pgfpathlineto{\pgfqpoint{1.448224in}{1.883227in}}%
\pgfpathlineto{\pgfqpoint{1.456062in}{1.766503in}}%
\pgfpathlineto{\pgfqpoint{1.456426in}{1.771115in}}%
\pgfpathlineto{\pgfqpoint{1.464388in}{1.888295in}}%
\pgfpathlineto{\pgfqpoint{1.464709in}{1.884104in}}%
\pgfpathlineto{\pgfqpoint{1.472547in}{1.767686in}}%
\pgfpathlineto{\pgfqpoint{1.472910in}{1.772299in}}%
\pgfpathlineto{\pgfqpoint{1.480873in}{1.889727in}}%
\pgfpathlineto{\pgfqpoint{1.481194in}{1.885544in}}%
\pgfpathlineto{\pgfqpoint{1.489032in}{1.769285in}}%
\pgfpathlineto{\pgfqpoint{1.489395in}{1.773907in}}%
\pgfpathlineto{\pgfqpoint{1.497358in}{1.891420in}}%
\pgfpathlineto{\pgfqpoint{1.497679in}{1.887238in}}%
\pgfpathlineto{\pgfqpoint{1.505517in}{1.770995in}}%
\pgfpathlineto{\pgfqpoint{1.505879in}{1.775601in}}%
\pgfpathlineto{\pgfqpoint{1.513843in}{1.893081in}}%
\pgfpathlineto{\pgfqpoint{1.514164in}{1.888895in}}%
\pgfpathlineto{\pgfqpoint{1.522001in}{1.772550in}}%
\pgfpathlineto{\pgfqpoint{1.522363in}{1.777137in}}%
\pgfpathlineto{\pgfqpoint{1.530328in}{1.894478in}}%
\pgfpathlineto{\pgfqpoint{1.530649in}{1.890286in}}%
\pgfpathlineto{\pgfqpoint{1.538486in}{1.773756in}}%
\pgfpathlineto{\pgfqpoint{1.538847in}{1.778323in}}%
\pgfpathlineto{\pgfqpoint{1.546812in}{1.895460in}}%
\pgfpathlineto{\pgfqpoint{1.547134in}{1.891259in}}%
\pgfpathlineto{\pgfqpoint{1.554971in}{1.774503in}}%
\pgfpathlineto{\pgfqpoint{1.555332in}{1.779054in}}%
\pgfpathlineto{\pgfqpoint{1.563297in}{1.895962in}}%
\pgfpathlineto{\pgfqpoint{1.563618in}{1.891751in}}%
\pgfpathlineto{\pgfqpoint{1.571456in}{1.774769in}}%
\pgfpathlineto{\pgfqpoint{1.571816in}{1.779307in}}%
\pgfpathlineto{\pgfqpoint{1.579782in}{1.896001in}}%
\pgfpathlineto{\pgfqpoint{1.580103in}{1.891781in}}%
\pgfpathlineto{\pgfqpoint{1.587941in}{1.774604in}}%
\pgfpathlineto{\pgfqpoint{1.588301in}{1.779135in}}%
\pgfpathlineto{\pgfqpoint{1.596267in}{1.895655in}}%
\pgfpathlineto{\pgfqpoint{1.596588in}{1.891429in}}%
\pgfpathlineto{\pgfqpoint{1.604426in}{1.774111in}}%
\pgfpathlineto{\pgfqpoint{1.604786in}{1.778641in}}%
\pgfpathlineto{\pgfqpoint{1.612752in}{1.895045in}}%
\pgfpathlineto{\pgfqpoint{1.613073in}{1.890815in}}%
\pgfpathlineto{\pgfqpoint{1.620910in}{1.773421in}}%
\pgfpathlineto{\pgfqpoint{1.621272in}{1.777954in}}%
\pgfpathlineto{\pgfqpoint{1.629237in}{1.894307in}}%
\pgfpathlineto{\pgfqpoint{1.629558in}{1.890077in}}%
\pgfpathlineto{\pgfqpoint{1.637395in}{1.772669in}}%
\pgfpathlineto{\pgfqpoint{1.637757in}{1.777208in}}%
\pgfpathlineto{\pgfqpoint{1.645722in}{1.893571in}}%
\pgfpathlineto{\pgfqpoint{1.646043in}{1.889342in}}%
\pgfpathlineto{\pgfqpoint{1.653880in}{1.771974in}}%
\pgfpathlineto{\pgfqpoint{1.654242in}{1.776522in}}%
\pgfpathlineto{\pgfqpoint{1.662206in}{1.892942in}}%
\pgfpathlineto{\pgfqpoint{1.662528in}{1.888715in}}%
\pgfpathlineto{\pgfqpoint{1.670365in}{1.771426in}}%
\pgfpathlineto{\pgfqpoint{1.670727in}{1.775982in}}%
\pgfpathlineto{\pgfqpoint{1.678691in}{1.892490in}}%
\pgfpathlineto{\pgfqpoint{1.679012in}{1.888267in}}%
\pgfpathlineto{\pgfqpoint{1.686850in}{1.771076in}}%
\pgfpathlineto{\pgfqpoint{1.687212in}{1.775640in}}%
\pgfpathlineto{\pgfqpoint{1.695176in}{1.892247in}}%
\pgfpathlineto{\pgfqpoint{1.695497in}{1.888029in}}%
\pgfpathlineto{\pgfqpoint{1.703335in}{1.770938in}}%
\pgfpathlineto{\pgfqpoint{1.703697in}{1.775507in}}%
\pgfpathlineto{\pgfqpoint{1.711661in}{1.892210in}}%
\pgfpathlineto{\pgfqpoint{1.711982in}{1.887996in}}%
\pgfpathlineto{\pgfqpoint{1.719820in}{1.770993in}}%
\pgfpathlineto{\pgfqpoint{1.720182in}{1.775565in}}%
\pgfpathlineto{\pgfqpoint{1.728146in}{1.892347in}}%
\pgfpathlineto{\pgfqpoint{1.728467in}{1.888136in}}%
\pgfpathlineto{\pgfqpoint{1.736304in}{1.771197in}}%
\pgfpathlineto{\pgfqpoint{1.736667in}{1.775771in}}%
\pgfpathlineto{\pgfqpoint{1.744631in}{1.892606in}}%
\pgfpathlineto{\pgfqpoint{1.744952in}{1.888396in}}%
\pgfpathlineto{\pgfqpoint{1.752789in}{1.771494in}}%
\pgfpathlineto{\pgfqpoint{1.753151in}{1.776066in}}%
\pgfpathlineto{\pgfqpoint{1.761115in}{1.892927in}}%
\pgfpathlineto{\pgfqpoint{1.761437in}{1.888717in}}%
\pgfpathlineto{\pgfqpoint{1.769274in}{1.771824in}}%
\pgfpathlineto{\pgfqpoint{1.769636in}{1.776393in}}%
\pgfpathlineto{\pgfqpoint{1.777600in}{1.893252in}}%
\pgfpathlineto{\pgfqpoint{1.777921in}{1.889043in}}%
\pgfpathlineto{\pgfqpoint{1.785759in}{1.772134in}}%
\pgfpathlineto{\pgfqpoint{1.786121in}{1.776700in}}%
\pgfpathlineto{\pgfqpoint{1.794085in}{1.893535in}}%
\pgfpathlineto{\pgfqpoint{1.794406in}{1.889324in}}%
\pgfpathlineto{\pgfqpoint{1.802244in}{1.772382in}}%
\pgfpathlineto{\pgfqpoint{1.802605in}{1.776944in}}%
\pgfpathlineto{\pgfqpoint{1.810570in}{1.893743in}}%
\pgfpathlineto{\pgfqpoint{1.810891in}{1.889530in}}%
\pgfpathlineto{\pgfqpoint{1.818729in}{1.772546in}}%
\pgfpathlineto{\pgfqpoint{1.819090in}{1.777104in}}%
\pgfpathlineto{\pgfqpoint{1.827055in}{1.893859in}}%
\pgfpathlineto{\pgfqpoint{1.827376in}{1.889645in}}%
\pgfpathlineto{\pgfqpoint{1.835214in}{1.772616in}}%
\pgfpathlineto{\pgfqpoint{1.835575in}{1.777172in}}%
\pgfpathlineto{\pgfqpoint{1.843540in}{1.893884in}}%
\pgfpathlineto{\pgfqpoint{1.843861in}{1.889668in}}%
\pgfpathlineto{\pgfqpoint{1.851698in}{1.772600in}}%
\pgfpathlineto{\pgfqpoint{1.852060in}{1.777154in}}%
\pgfpathlineto{\pgfqpoint{1.860025in}{1.893831in}}%
\pgfpathlineto{\pgfqpoint{1.860346in}{1.889613in}}%
\pgfpathlineto{\pgfqpoint{1.868183in}{1.772516in}}%
\pgfpathlineto{\pgfqpoint{1.868545in}{1.777070in}}%
\pgfpathlineto{\pgfqpoint{1.876509in}{1.893722in}}%
\pgfpathlineto{\pgfqpoint{1.876831in}{1.889503in}}%
\pgfpathlineto{\pgfqpoint{1.884668in}{1.772388in}}%
\pgfpathlineto{\pgfqpoint{1.885030in}{1.776943in}}%
\pgfpathlineto{\pgfqpoint{1.892994in}{1.893582in}}%
\pgfpathlineto{\pgfqpoint{1.893315in}{1.889364in}}%
\pgfpathlineto{\pgfqpoint{1.901153in}{1.772244in}}%
\pgfpathlineto{\pgfqpoint{1.901515in}{1.776799in}}%
\pgfpathlineto{\pgfqpoint{1.909479in}{1.893439in}}%
\pgfpathlineto{\pgfqpoint{1.909800in}{1.889220in}}%
\pgfpathlineto{\pgfqpoint{1.917638in}{1.772106in}}%
\pgfpathlineto{\pgfqpoint{1.918000in}{1.776663in}}%
\pgfpathlineto{\pgfqpoint{1.925964in}{1.893312in}}%
\pgfpathlineto{\pgfqpoint{1.926285in}{1.889093in}}%
\pgfpathlineto{\pgfqpoint{1.934123in}{1.771993in}}%
\pgfpathlineto{\pgfqpoint{1.934484in}{1.776552in}}%
\pgfpathlineto{\pgfqpoint{1.942449in}{1.893216in}}%
\pgfpathlineto{\pgfqpoint{1.942770in}{1.888999in}}%
\pgfpathlineto{\pgfqpoint{1.950607in}{1.771917in}}%
\pgfpathlineto{\pgfqpoint{1.950969in}{1.776478in}}%
\pgfpathlineto{\pgfqpoint{1.958934in}{1.893161in}}%
\pgfpathlineto{\pgfqpoint{1.959255in}{1.888945in}}%
\pgfpathlineto{\pgfqpoint{1.967092in}{1.771882in}}%
\pgfpathlineto{\pgfqpoint{1.967454in}{1.776444in}}%
\pgfpathlineto{\pgfqpoint{1.975419in}{1.893146in}}%
\pgfpathlineto{\pgfqpoint{1.975740in}{1.888931in}}%
\pgfpathlineto{\pgfqpoint{1.983577in}{1.771886in}}%
\pgfpathlineto{\pgfqpoint{1.983939in}{1.776448in}}%
\pgfpathlineto{\pgfqpoint{1.991903in}{1.893166in}}%
\pgfpathlineto{\pgfqpoint{1.992224in}{1.888951in}}%
\pgfpathlineto{\pgfqpoint{2.000062in}{1.771920in}}%
\pgfpathlineto{\pgfqpoint{2.000424in}{1.776483in}}%
\pgfpathlineto{\pgfqpoint{2.008388in}{1.893212in}}%
\pgfpathlineto{\pgfqpoint{2.008709in}{1.888998in}}%
\pgfpathlineto{\pgfqpoint{2.016547in}{1.771975in}}%
\pgfpathlineto{\pgfqpoint{2.016909in}{1.776537in}}%
\pgfpathlineto{\pgfqpoint{2.024873in}{1.893273in}}%
\pgfpathlineto{\pgfqpoint{2.025194in}{1.889058in}}%
\pgfpathlineto{\pgfqpoint{2.033032in}{1.772038in}}%
\pgfpathlineto{\pgfqpoint{2.033394in}{1.776600in}}%
\pgfpathlineto{\pgfqpoint{2.041358in}{1.893336in}}%
\pgfpathlineto{\pgfqpoint{2.041679in}{1.889122in}}%
\pgfpathlineto{\pgfqpoint{2.049517in}{1.772099in}}%
\pgfpathlineto{\pgfqpoint{2.049878in}{1.776660in}}%
\pgfpathlineto{\pgfqpoint{2.057843in}{1.893393in}}%
\pgfpathlineto{\pgfqpoint{2.058164in}{1.889178in}}%
\pgfpathlineto{\pgfqpoint{2.066001in}{1.772150in}}%
\pgfpathlineto{\pgfqpoint{2.066363in}{1.776711in}}%
\pgfpathlineto{\pgfqpoint{2.074328in}{1.893437in}}%
\pgfpathlineto{\pgfqpoint{2.074649in}{1.889221in}}%
\pgfpathlineto{\pgfqpoint{2.082486in}{1.772185in}}%
\pgfpathlineto{\pgfqpoint{2.082848in}{1.776745in}}%
\pgfpathlineto{\pgfqpoint{2.090812in}{1.893463in}}%
\pgfpathlineto{\pgfqpoint{2.091134in}{1.889247in}}%
\pgfpathlineto{\pgfqpoint{2.098971in}{1.772202in}}%
\pgfpathlineto{\pgfqpoint{2.099333in}{1.776762in}}%
\pgfpathlineto{\pgfqpoint{2.107297in}{1.893471in}}%
\pgfpathlineto{\pgfqpoint{2.107618in}{1.889255in}}%
\pgfpathlineto{\pgfqpoint{2.115456in}{1.772202in}}%
\pgfpathlineto{\pgfqpoint{2.115818in}{1.776761in}}%
\pgfpathlineto{\pgfqpoint{2.123782in}{1.893464in}}%
\pgfpathlineto{\pgfqpoint{2.124103in}{1.889247in}}%
\pgfpathlineto{\pgfqpoint{2.131941in}{1.772189in}}%
\pgfpathlineto{\pgfqpoint{2.132303in}{1.776747in}}%
\pgfpathlineto{\pgfqpoint{2.140267in}{1.893444in}}%
\pgfpathlineto{\pgfqpoint{2.140588in}{1.889228in}}%
\pgfpathlineto{\pgfqpoint{2.148426in}{1.772165in}}%
\pgfpathlineto{\pgfqpoint{2.148787in}{1.776724in}}%
\pgfpathlineto{\pgfqpoint{2.156752in}{1.893418in}}%
\pgfpathlineto{\pgfqpoint{2.157073in}{1.889202in}}%
\pgfpathlineto{\pgfqpoint{2.164910in}{1.772138in}}%
\pgfpathlineto{\pgfqpoint{2.165272in}{1.776697in}}%
\pgfpathlineto{\pgfqpoint{2.173237in}{1.893390in}}%
\pgfpathlineto{\pgfqpoint{2.173558in}{1.889174in}}%
\pgfpathlineto{\pgfqpoint{2.181395in}{1.772111in}}%
\pgfpathlineto{\pgfqpoint{2.181757in}{1.776670in}}%
\pgfpathlineto{\pgfqpoint{2.189722in}{1.893365in}}%
\pgfpathlineto{\pgfqpoint{2.190043in}{1.889149in}}%
\pgfpathlineto{\pgfqpoint{2.197880in}{1.772088in}}%
\pgfpathlineto{\pgfqpoint{2.198242in}{1.776647in}}%
\pgfpathlineto{\pgfqpoint{2.206206in}{1.893345in}}%
\pgfpathlineto{\pgfqpoint{2.206528in}{1.889129in}}%
\pgfpathlineto{\pgfqpoint{2.214365in}{1.772071in}}%
\pgfpathlineto{\pgfqpoint{2.214727in}{1.776631in}}%
\pgfpathlineto{\pgfqpoint{2.222691in}{1.893333in}}%
\pgfpathlineto{\pgfqpoint{2.223012in}{1.889117in}}%
\pgfpathlineto{\pgfqpoint{2.230850in}{1.772063in}}%
\pgfpathlineto{\pgfqpoint{2.231212in}{1.776623in}}%
\pgfpathlineto{\pgfqpoint{2.239176in}{1.893328in}}%
\pgfpathlineto{\pgfqpoint{2.239497in}{1.889112in}}%
\pgfpathlineto{\pgfqpoint{2.247335in}{1.772062in}}%
\pgfpathlineto{\pgfqpoint{2.247697in}{1.776623in}}%
\pgfpathlineto{\pgfqpoint{2.255661in}{1.893331in}}%
\pgfpathlineto{\pgfqpoint{2.255982in}{1.889115in}}%
\pgfpathlineto{\pgfqpoint{2.263820in}{1.772068in}}%
\pgfpathlineto{\pgfqpoint{2.264181in}{1.776628in}}%
\pgfpathlineto{\pgfqpoint{2.272146in}{1.893339in}}%
\pgfpathlineto{\pgfqpoint{2.272467in}{1.889123in}}%
\pgfpathlineto{\pgfqpoint{2.280304in}{1.772078in}}%
\pgfpathlineto{\pgfqpoint{2.280666in}{1.776638in}}%
\pgfpathlineto{\pgfqpoint{2.288631in}{1.893350in}}%
\pgfpathlineto{\pgfqpoint{2.288952in}{1.889135in}}%
\pgfpathlineto{\pgfqpoint{2.296789in}{1.772090in}}%
\pgfpathlineto{\pgfqpoint{2.297151in}{1.776650in}}%
\pgfpathlineto{\pgfqpoint{2.305115in}{1.893363in}}%
\pgfpathlineto{\pgfqpoint{2.305437in}{1.889147in}}%
\pgfpathlineto{\pgfqpoint{2.313274in}{1.772102in}}%
\pgfpathlineto{\pgfqpoint{2.313636in}{1.776662in}}%
\pgfpathlineto{\pgfqpoint{2.321600in}{1.893374in}}%
\pgfpathlineto{\pgfqpoint{2.321921in}{1.889158in}}%
\pgfpathlineto{\pgfqpoint{2.329759in}{1.772112in}}%
\pgfpathlineto{\pgfqpoint{2.330121in}{1.776672in}}%
\pgfpathlineto{\pgfqpoint{2.338085in}{1.893383in}}%
\pgfpathlineto{\pgfqpoint{2.338406in}{1.889167in}}%
\pgfpathlineto{\pgfqpoint{2.346244in}{1.772120in}}%
\pgfpathlineto{\pgfqpoint{2.346606in}{1.776679in}}%
\pgfpathlineto{\pgfqpoint{2.354570in}{1.893389in}}%
\pgfpathlineto{\pgfqpoint{2.354891in}{1.889173in}}%
\pgfpathlineto{\pgfqpoint{2.362729in}{1.772124in}}%
\pgfpathlineto{\pgfqpoint{2.363090in}{1.776683in}}%
\pgfpathlineto{\pgfqpoint{2.371055in}{1.893391in}}%
\pgfpathlineto{\pgfqpoint{2.371376in}{1.889175in}}%
\pgfpathlineto{\pgfqpoint{2.379214in}{1.772124in}}%
\pgfpathlineto{\pgfqpoint{2.379575in}{1.776684in}}%
\pgfpathlineto{\pgfqpoint{2.387540in}{1.893390in}}%
\pgfpathlineto{\pgfqpoint{2.387861in}{1.889174in}}%
\pgfpathlineto{\pgfqpoint{2.395698in}{1.772122in}}%
\pgfpathlineto{\pgfqpoint{2.396060in}{1.776682in}}%
\pgfpathlineto{\pgfqpoint{2.404025in}{1.893387in}}%
\pgfpathlineto{\pgfqpoint{2.404346in}{1.889171in}}%
\pgfpathlineto{\pgfqpoint{2.412183in}{1.772118in}}%
\pgfpathlineto{\pgfqpoint{2.412545in}{1.776678in}}%
\pgfpathlineto{\pgfqpoint{2.420509in}{1.893382in}}%
\pgfpathlineto{\pgfqpoint{2.420831in}{1.889166in}}%
\pgfpathlineto{\pgfqpoint{2.428668in}{1.772113in}}%
\pgfpathlineto{\pgfqpoint{2.429030in}{1.776672in}}%
\pgfpathlineto{\pgfqpoint{2.436994in}{1.893377in}}%
\pgfpathlineto{\pgfqpoint{2.437315in}{1.889161in}}%
\pgfpathlineto{\pgfqpoint{2.445153in}{1.772107in}}%
\pgfpathlineto{\pgfqpoint{2.445515in}{1.776667in}}%
\pgfpathlineto{\pgfqpoint{2.453479in}{1.893371in}}%
\pgfpathlineto{\pgfqpoint{2.453800in}{1.889155in}}%
\pgfpathlineto{\pgfqpoint{2.461638in}{1.772103in}}%
\pgfpathlineto{\pgfqpoint{2.462000in}{1.776663in}}%
\pgfpathlineto{\pgfqpoint{2.469964in}{1.893367in}}%
\pgfpathlineto{\pgfqpoint{2.470285in}{1.889151in}}%
\pgfpathlineto{\pgfqpoint{2.478123in}{1.772099in}}%
\pgfpathlineto{\pgfqpoint{2.478484in}{1.776659in}}%
\pgfpathlineto{\pgfqpoint{2.486449in}{1.893365in}}%
\pgfpathlineto{\pgfqpoint{2.486770in}{1.889149in}}%
\pgfpathlineto{\pgfqpoint{2.494597in}{1.772286in}}%
\pgfpathlineto{\pgfqpoint{2.494957in}{1.776724in}}%
\pgfpathlineto{\pgfqpoint{2.502934in}{1.893611in}}%
\pgfpathlineto{\pgfqpoint{2.503255in}{1.889395in}}%
\pgfpathlineto{\pgfqpoint{2.511092in}{1.772330in}}%
\pgfpathlineto{\pgfqpoint{2.511454in}{1.776887in}}%
\pgfpathlineto{\pgfqpoint{2.519419in}{1.893575in}}%
\pgfpathlineto{\pgfqpoint{2.519740in}{1.889358in}}%
\pgfpathlineto{\pgfqpoint{2.527577in}{1.772282in}}%
\pgfpathlineto{\pgfqpoint{2.527939in}{1.776839in}}%
\pgfpathlineto{\pgfqpoint{2.535903in}{1.893517in}}%
\pgfpathlineto{\pgfqpoint{2.536224in}{1.889300in}}%
\pgfpathlineto{\pgfqpoint{2.544062in}{1.772217in}}%
\pgfpathlineto{\pgfqpoint{2.544424in}{1.776774in}}%
\pgfpathlineto{\pgfqpoint{2.552388in}{1.893448in}}%
\pgfpathlineto{\pgfqpoint{2.552709in}{1.889231in}}%
\pgfpathlineto{\pgfqpoint{2.560547in}{1.772148in}}%
\pgfpathlineto{\pgfqpoint{2.560909in}{1.776706in}}%
\pgfpathlineto{\pgfqpoint{2.568873in}{1.893382in}}%
\pgfpathlineto{\pgfqpoint{2.569194in}{1.889164in}}%
\pgfpathlineto{\pgfqpoint{2.577032in}{1.772085in}}%
\pgfpathlineto{\pgfqpoint{2.577394in}{1.776644in}}%
\pgfpathlineto{\pgfqpoint{2.585358in}{1.893326in}}%
\pgfpathlineto{\pgfqpoint{2.585679in}{1.889109in}}%
\pgfpathlineto{\pgfqpoint{2.593517in}{1.772037in}}%
\pgfpathlineto{\pgfqpoint{2.593878in}{1.776597in}}%
\pgfpathlineto{\pgfqpoint{2.601843in}{1.893287in}}%
\pgfpathlineto{\pgfqpoint{2.602164in}{1.889070in}}%
\pgfpathlineto{\pgfqpoint{2.610001in}{1.772008in}}%
\pgfpathlineto{\pgfqpoint{2.610363in}{1.776568in}}%
\pgfpathlineto{\pgfqpoint{2.618328in}{1.893267in}}%
\pgfpathlineto{\pgfqpoint{2.618649in}{1.889051in}}%
\pgfpathlineto{\pgfqpoint{2.626486in}{1.771998in}}%
\pgfpathlineto{\pgfqpoint{2.626848in}{1.776559in}}%
\pgfpathlineto{\pgfqpoint{2.634812in}{1.893266in}}%
\pgfpathlineto{\pgfqpoint{2.635134in}{1.889050in}}%
\pgfpathlineto{\pgfqpoint{2.642971in}{1.772005in}}%
\pgfpathlineto{\pgfqpoint{2.643333in}{1.776566in}}%
\pgfpathlineto{\pgfqpoint{2.651297in}{1.893280in}}%
\pgfpathlineto{\pgfqpoint{2.651618in}{1.889065in}}%
\pgfpathlineto{\pgfqpoint{2.659456in}{1.772025in}}%
\pgfpathlineto{\pgfqpoint{2.659818in}{1.776586in}}%
\pgfpathlineto{\pgfqpoint{2.667782in}{1.893305in}}%
\pgfpathlineto{\pgfqpoint{2.668103in}{1.889090in}}%
\pgfpathlineto{\pgfqpoint{2.675941in}{1.772053in}}%
\pgfpathlineto{\pgfqpoint{2.676303in}{1.776614in}}%
\pgfpathlineto{\pgfqpoint{2.684267in}{1.893335in}}%
\pgfpathlineto{\pgfqpoint{2.684588in}{1.889120in}}%
\pgfpathlineto{\pgfqpoint{2.692426in}{1.772083in}}%
\pgfpathlineto{\pgfqpoint{2.692787in}{1.776644in}}%
\pgfpathlineto{\pgfqpoint{2.700752in}{1.893365in}}%
\pgfpathlineto{\pgfqpoint{2.701073in}{1.889149in}}%
\pgfpathlineto{\pgfqpoint{2.708910in}{1.772111in}}%
\pgfpathlineto{\pgfqpoint{2.709272in}{1.776672in}}%
\pgfpathlineto{\pgfqpoint{2.717237in}{1.893390in}}%
\pgfpathlineto{\pgfqpoint{2.717558in}{1.889174in}}%
\pgfpathlineto{\pgfqpoint{2.725395in}{1.772133in}}%
\pgfpathlineto{\pgfqpoint{2.725757in}{1.776693in}}%
\pgfpathlineto{\pgfqpoint{2.733722in}{1.893408in}}%
\pgfpathlineto{\pgfqpoint{2.734043in}{1.889192in}}%
\pgfpathlineto{\pgfqpoint{2.741880in}{1.772147in}}%
\pgfpathlineto{\pgfqpoint{2.742242in}{1.776707in}}%
\pgfpathlineto{\pgfqpoint{2.750206in}{1.893417in}}%
\pgfpathlineto{\pgfqpoint{2.750528in}{1.889202in}}%
\pgfpathlineto{\pgfqpoint{2.758354in}{1.772341in}}%
\pgfpathlineto{\pgfqpoint{2.758714in}{1.776779in}}%
\pgfpathlineto{\pgfqpoint{2.766691in}{1.893666in}}%
\pgfpathlineto{\pgfqpoint{2.767012in}{1.889450in}}%
\pgfpathlineto{\pgfqpoint{2.774850in}{1.772383in}}%
\pgfpathlineto{\pgfqpoint{2.775212in}{1.776940in}}%
\pgfpathlineto{\pgfqpoint{2.783176in}{1.893624in}}%
\pgfpathlineto{\pgfqpoint{2.783497in}{1.889407in}}%
\pgfpathlineto{\pgfqpoint{2.791335in}{1.772326in}}%
\pgfpathlineto{\pgfqpoint{2.791696in}{1.776882in}}%
\pgfpathlineto{\pgfqpoint{2.799661in}{1.893554in}}%
\pgfpathlineto{\pgfqpoint{2.799982in}{1.889337in}}%
\pgfpathlineto{\pgfqpoint{2.807820in}{1.772247in}}%
\pgfpathlineto{\pgfqpoint{2.808181in}{1.776804in}}%
\pgfpathlineto{\pgfqpoint{2.816146in}{1.893471in}}%
\pgfpathlineto{\pgfqpoint{2.816467in}{1.889253in}}%
\pgfpathlineto{\pgfqpoint{2.824304in}{1.772162in}}%
\pgfpathlineto{\pgfqpoint{2.824666in}{1.776720in}}%
\pgfpathlineto{\pgfqpoint{2.832631in}{1.893388in}}%
\pgfpathlineto{\pgfqpoint{2.832952in}{1.889171in}}%
\pgfpathlineto{\pgfqpoint{2.840789in}{1.772085in}}%
\pgfpathlineto{\pgfqpoint{2.841151in}{1.776644in}}%
\pgfpathlineto{\pgfqpoint{2.849115in}{1.893319in}}%
\pgfpathlineto{\pgfqpoint{2.849437in}{1.889102in}}%
\pgfpathlineto{\pgfqpoint{2.857274in}{1.772025in}}%
\pgfpathlineto{\pgfqpoint{2.857636in}{1.776585in}}%
\pgfpathlineto{\pgfqpoint{2.865600in}{1.893270in}}%
\pgfpathlineto{\pgfqpoint{2.865921in}{1.889053in}}%
\pgfpathlineto{\pgfqpoint{2.873759in}{1.771987in}}%
\pgfpathlineto{\pgfqpoint{2.874121in}{1.776548in}}%
\pgfpathlineto{\pgfqpoint{2.882085in}{1.893244in}}%
\pgfpathlineto{\pgfqpoint{2.882406in}{1.889028in}}%
\pgfpathlineto{\pgfqpoint{2.890244in}{1.771974in}}%
\pgfpathlineto{\pgfqpoint{2.890606in}{1.776535in}}%
\pgfpathlineto{\pgfqpoint{2.898570in}{1.893242in}}%
\pgfpathlineto{\pgfqpoint{2.898891in}{1.889026in}}%
\pgfpathlineto{\pgfqpoint{2.906729in}{1.771981in}}%
\pgfpathlineto{\pgfqpoint{2.907091in}{1.776543in}}%
\pgfpathlineto{\pgfqpoint{2.915055in}{1.893258in}}%
\pgfpathlineto{\pgfqpoint{2.915376in}{1.889043in}}%
\pgfpathlineto{\pgfqpoint{2.923214in}{1.772005in}}%
\pgfpathlineto{\pgfqpoint{2.923575in}{1.776567in}}%
\pgfpathlineto{\pgfqpoint{2.931540in}{1.893288in}}%
\pgfpathlineto{\pgfqpoint{2.931861in}{1.889073in}}%
\pgfpathlineto{\pgfqpoint{2.939698in}{1.772039in}}%
\pgfpathlineto{\pgfqpoint{2.940060in}{1.776601in}}%
\pgfpathlineto{\pgfqpoint{2.948025in}{1.893325in}}%
\pgfpathlineto{\pgfqpoint{2.948346in}{1.889109in}}%
\pgfpathlineto{\pgfqpoint{2.956183in}{1.772076in}}%
\pgfpathlineto{\pgfqpoint{2.956545in}{1.776637in}}%
\pgfpathlineto{\pgfqpoint{2.964509in}{1.893361in}}%
\pgfpathlineto{\pgfqpoint{2.964831in}{1.889146in}}%
\pgfpathlineto{\pgfqpoint{2.972668in}{1.772111in}}%
\pgfpathlineto{\pgfqpoint{2.973030in}{1.776671in}}%
\pgfpathlineto{\pgfqpoint{2.980994in}{1.893392in}}%
\pgfpathlineto{\pgfqpoint{2.981315in}{1.889177in}}%
\pgfpathlineto{\pgfqpoint{2.989153in}{1.772138in}}%
\pgfpathlineto{\pgfqpoint{2.989515in}{1.776698in}}%
\pgfpathlineto{\pgfqpoint{2.997479in}{1.893415in}}%
\pgfpathlineto{\pgfqpoint{2.997800in}{1.889199in}}%
\pgfpathlineto{\pgfqpoint{3.005638in}{1.772156in}}%
\pgfpathlineto{\pgfqpoint{3.006000in}{1.776716in}}%
\pgfpathlineto{\pgfqpoint{3.013964in}{1.893427in}}%
\pgfpathlineto{\pgfqpoint{3.014285in}{1.889211in}}%
\pgfpathlineto{\pgfqpoint{3.022123in}{1.772163in}}%
\pgfpathlineto{\pgfqpoint{3.022484in}{1.776722in}}%
\pgfpathlineto{\pgfqpoint{3.030449in}{1.893429in}}%
\pgfpathlineto{\pgfqpoint{3.030770in}{1.889213in}}%
\pgfpathlineto{\pgfqpoint{3.038607in}{1.772160in}}%
\pgfpathlineto{\pgfqpoint{3.038969in}{1.776720in}}%
\pgfpathlineto{\pgfqpoint{3.046934in}{1.893423in}}%
\pgfpathlineto{\pgfqpoint{3.047255in}{1.889207in}}%
\pgfpathlineto{\pgfqpoint{3.055092in}{1.772150in}}%
\pgfpathlineto{\pgfqpoint{3.055454in}{1.776710in}}%
\pgfpathlineto{\pgfqpoint{3.063419in}{1.893410in}}%
\pgfpathlineto{\pgfqpoint{3.063740in}{1.889194in}}%
\pgfpathlineto{\pgfqpoint{3.071577in}{1.772136in}}%
\pgfpathlineto{\pgfqpoint{3.071939in}{1.776695in}}%
\pgfpathlineto{\pgfqpoint{3.079903in}{1.893394in}}%
\pgfpathlineto{\pgfqpoint{3.080224in}{1.889178in}}%
\pgfpathlineto{\pgfqpoint{3.088062in}{1.772120in}}%
\pgfpathlineto{\pgfqpoint{3.088424in}{1.776679in}}%
\pgfpathlineto{\pgfqpoint{3.096388in}{1.893378in}}%
\pgfpathlineto{\pgfqpoint{3.096709in}{1.889162in}}%
\pgfpathlineto{\pgfqpoint{3.104547in}{1.772104in}}%
\pgfpathlineto{\pgfqpoint{3.104909in}{1.776664in}}%
\pgfpathlineto{\pgfqpoint{3.112873in}{1.893364in}}%
\pgfpathlineto{\pgfqpoint{3.113194in}{1.889148in}}%
\pgfpathlineto{\pgfqpoint{3.121032in}{1.772092in}}%
\pgfpathlineto{\pgfqpoint{3.121394in}{1.776652in}}%
\pgfpathlineto{\pgfqpoint{3.129358in}{1.893354in}}%
\pgfpathlineto{\pgfqpoint{3.129679in}{1.889138in}}%
\pgfpathlineto{\pgfqpoint{3.137517in}{1.772084in}}%
\pgfpathlineto{\pgfqpoint{3.137878in}{1.776644in}}%
\pgfpathlineto{\pgfqpoint{3.145843in}{1.893348in}}%
\pgfpathlineto{\pgfqpoint{3.146164in}{1.889132in}}%
\pgfpathlineto{\pgfqpoint{3.154001in}{1.772080in}}%
\pgfpathlineto{\pgfqpoint{3.154363in}{1.776640in}}%
\pgfpathlineto{\pgfqpoint{3.162328in}{1.893347in}}%
\pgfpathlineto{\pgfqpoint{3.162649in}{1.889131in}}%
\pgfpathlineto{\pgfqpoint{3.170132in}{1.777082in}}%
\pgfpathlineto{\pgfqpoint{3.170132in}{1.777082in}}%
\pgfusepath{stroke}%
\end{pgfscope}%
\begin{pgfscope}%
\pgfsetrectcap%
\pgfsetmiterjoin%
\pgfsetlinewidth{0.803000pt}%
\definecolor{currentstroke}{rgb}{0.000000,0.000000,0.000000}%
\pgfsetstrokecolor{currentstroke}%
\pgfsetdash{}{0pt}%
\pgfpathmoveto{\pgfqpoint{0.573768in}{1.532029in}}%
\pgfpathlineto{\pgfqpoint{0.573768in}{2.438279in}}%
\pgfusepath{stroke}%
\end{pgfscope}%
\begin{pgfscope}%
\pgfsetrectcap%
\pgfsetmiterjoin%
\pgfsetlinewidth{0.803000pt}%
\definecolor{currentstroke}{rgb}{0.000000,0.000000,0.000000}%
\pgfsetstrokecolor{currentstroke}%
\pgfsetdash{}{0pt}%
\pgfpathmoveto{\pgfqpoint{3.293768in}{1.532029in}}%
\pgfpathlineto{\pgfqpoint{3.293768in}{2.438279in}}%
\pgfusepath{stroke}%
\end{pgfscope}%
\begin{pgfscope}%
\pgfsetrectcap%
\pgfsetmiterjoin%
\pgfsetlinewidth{0.803000pt}%
\definecolor{currentstroke}{rgb}{0.000000,0.000000,0.000000}%
\pgfsetstrokecolor{currentstroke}%
\pgfsetdash{}{0pt}%
\pgfpathmoveto{\pgfqpoint{0.573768in}{1.532029in}}%
\pgfpathlineto{\pgfqpoint{3.293768in}{1.532029in}}%
\pgfusepath{stroke}%
\end{pgfscope}%
\begin{pgfscope}%
\pgfsetrectcap%
\pgfsetmiterjoin%
\pgfsetlinewidth{0.803000pt}%
\definecolor{currentstroke}{rgb}{0.000000,0.000000,0.000000}%
\pgfsetstrokecolor{currentstroke}%
\pgfsetdash{}{0pt}%
\pgfpathmoveto{\pgfqpoint{0.573768in}{2.438279in}}%
\pgfpathlineto{\pgfqpoint{3.293768in}{2.438279in}}%
\pgfusepath{stroke}%
\end{pgfscope}%
\begin{pgfscope}%
\definecolor{textcolor}{rgb}{0.000000,0.000000,0.000000}%
\pgfsetstrokecolor{textcolor}%
\pgfsetfillcolor{textcolor}%
\pgftext[x=0.628168in,y=2.374842in,left,top]{\color{textcolor}{\rmfamily\fontsize{8.000000}{9.600000}\selectfont\catcode`\^=\active\def^{\ifmmode\sp\else\^{}\fi}\catcode`\%=\active\def%{\%}(b)}}%
\end{pgfscope}%
\begin{pgfscope}%
\pgfsetbuttcap%
\pgfsetmiterjoin%
\definecolor{currentfill}{rgb}{1.000000,1.000000,1.000000}%
\pgfsetfillcolor{currentfill}%
\pgfsetlinewidth{0.000000pt}%
\definecolor{currentstroke}{rgb}{0.000000,0.000000,0.000000}%
\pgfsetstrokecolor{currentstroke}%
\pgfsetstrokeopacity{0.000000}%
\pgfsetdash{}{0pt}%
\pgfpathmoveto{\pgfqpoint{0.573768in}{0.462654in}}%
\pgfpathlineto{\pgfqpoint{3.293768in}{0.462654in}}%
\pgfpathlineto{\pgfqpoint{3.293768in}{1.368904in}}%
\pgfpathlineto{\pgfqpoint{0.573768in}{1.368904in}}%
\pgfpathlineto{\pgfqpoint{0.573768in}{0.462654in}}%
\pgfpathclose%
\pgfusepath{fill}%
\end{pgfscope}%
\begin{pgfscope}%
\pgfpathrectangle{\pgfqpoint{0.573768in}{0.462654in}}{\pgfqpoint{2.720000in}{0.906250in}}%
\pgfusepath{clip}%
\pgfsetbuttcap%
\pgfsetroundjoin%
\pgfsetlinewidth{0.501875pt}%
\definecolor{currentstroke}{rgb}{0.690196,0.690196,0.690196}%
\pgfsetstrokecolor{currentstroke}%
\pgfsetstrokeopacity{0.250000}%
\pgfsetdash{{1.850000pt}{0.800000pt}}{0.000000pt}%
\pgfpathmoveto{\pgfqpoint{0.697405in}{0.462654in}}%
\pgfpathlineto{\pgfqpoint{0.697405in}{1.368904in}}%
\pgfusepath{stroke}%
\end{pgfscope}%
\begin{pgfscope}%
\pgfsetbuttcap%
\pgfsetroundjoin%
\definecolor{currentfill}{rgb}{0.000000,0.000000,0.000000}%
\pgfsetfillcolor{currentfill}%
\pgfsetlinewidth{0.803000pt}%
\definecolor{currentstroke}{rgb}{0.000000,0.000000,0.000000}%
\pgfsetstrokecolor{currentstroke}%
\pgfsetdash{}{0pt}%
\pgfsys@defobject{currentmarker}{\pgfqpoint{0.000000in}{-0.048611in}}{\pgfqpoint{0.000000in}{0.000000in}}{%
\pgfpathmoveto{\pgfqpoint{0.000000in}{0.000000in}}%
\pgfpathlineto{\pgfqpoint{0.000000in}{-0.048611in}}%
\pgfusepath{stroke,fill}%
}%
\begin{pgfscope}%
\pgfsys@transformshift{0.697405in}{0.462654in}%
\pgfsys@useobject{currentmarker}{}%
\end{pgfscope}%
\end{pgfscope}%
\begin{pgfscope}%
\definecolor{textcolor}{rgb}{0.000000,0.000000,0.000000}%
\pgfsetstrokecolor{textcolor}%
\pgfsetfillcolor{textcolor}%
\pgftext[x=0.697405in,y=0.365432in,,top]{\color{textcolor}{\rmfamily\fontsize{8.000000}{9.600000}\selectfont\catcode`\^=\active\def^{\ifmmode\sp\else\^{}\fi}\catcode`\%=\active\def%{\%}$\mathdefault{0}$}}%
\end{pgfscope}%
\begin{pgfscope}%
\pgfpathrectangle{\pgfqpoint{0.573768in}{0.462654in}}{\pgfqpoint{2.720000in}{0.906250in}}%
\pgfusepath{clip}%
\pgfsetbuttcap%
\pgfsetroundjoin%
\pgfsetlinewidth{0.501875pt}%
\definecolor{currentstroke}{rgb}{0.690196,0.690196,0.690196}%
\pgfsetstrokecolor{currentstroke}%
\pgfsetstrokeopacity{0.250000}%
\pgfsetdash{{1.850000pt}{0.800000pt}}{0.000000pt}%
\pgfpathmoveto{\pgfqpoint{1.109526in}{0.462654in}}%
\pgfpathlineto{\pgfqpoint{1.109526in}{1.368904in}}%
\pgfusepath{stroke}%
\end{pgfscope}%
\begin{pgfscope}%
\pgfsetbuttcap%
\pgfsetroundjoin%
\definecolor{currentfill}{rgb}{0.000000,0.000000,0.000000}%
\pgfsetfillcolor{currentfill}%
\pgfsetlinewidth{0.803000pt}%
\definecolor{currentstroke}{rgb}{0.000000,0.000000,0.000000}%
\pgfsetstrokecolor{currentstroke}%
\pgfsetdash{}{0pt}%
\pgfsys@defobject{currentmarker}{\pgfqpoint{0.000000in}{-0.048611in}}{\pgfqpoint{0.000000in}{0.000000in}}{%
\pgfpathmoveto{\pgfqpoint{0.000000in}{0.000000in}}%
\pgfpathlineto{\pgfqpoint{0.000000in}{-0.048611in}}%
\pgfusepath{stroke,fill}%
}%
\begin{pgfscope}%
\pgfsys@transformshift{1.109526in}{0.462654in}%
\pgfsys@useobject{currentmarker}{}%
\end{pgfscope}%
\end{pgfscope}%
\begin{pgfscope}%
\definecolor{textcolor}{rgb}{0.000000,0.000000,0.000000}%
\pgfsetstrokecolor{textcolor}%
\pgfsetfillcolor{textcolor}%
\pgftext[x=1.109526in,y=0.365432in,,top]{\color{textcolor}{\rmfamily\fontsize{8.000000}{9.600000}\selectfont\catcode`\^=\active\def^{\ifmmode\sp\else\^{}\fi}\catcode`\%=\active\def%{\%}$\mathdefault{1}$}}%
\end{pgfscope}%
\begin{pgfscope}%
\pgfpathrectangle{\pgfqpoint{0.573768in}{0.462654in}}{\pgfqpoint{2.720000in}{0.906250in}}%
\pgfusepath{clip}%
\pgfsetbuttcap%
\pgfsetroundjoin%
\pgfsetlinewidth{0.501875pt}%
\definecolor{currentstroke}{rgb}{0.690196,0.690196,0.690196}%
\pgfsetstrokecolor{currentstroke}%
\pgfsetstrokeopacity{0.250000}%
\pgfsetdash{{1.850000pt}{0.800000pt}}{0.000000pt}%
\pgfpathmoveto{\pgfqpoint{1.521647in}{0.462654in}}%
\pgfpathlineto{\pgfqpoint{1.521647in}{1.368904in}}%
\pgfusepath{stroke}%
\end{pgfscope}%
\begin{pgfscope}%
\pgfsetbuttcap%
\pgfsetroundjoin%
\definecolor{currentfill}{rgb}{0.000000,0.000000,0.000000}%
\pgfsetfillcolor{currentfill}%
\pgfsetlinewidth{0.803000pt}%
\definecolor{currentstroke}{rgb}{0.000000,0.000000,0.000000}%
\pgfsetstrokecolor{currentstroke}%
\pgfsetdash{}{0pt}%
\pgfsys@defobject{currentmarker}{\pgfqpoint{0.000000in}{-0.048611in}}{\pgfqpoint{0.000000in}{0.000000in}}{%
\pgfpathmoveto{\pgfqpoint{0.000000in}{0.000000in}}%
\pgfpathlineto{\pgfqpoint{0.000000in}{-0.048611in}}%
\pgfusepath{stroke,fill}%
}%
\begin{pgfscope}%
\pgfsys@transformshift{1.521647in}{0.462654in}%
\pgfsys@useobject{currentmarker}{}%
\end{pgfscope}%
\end{pgfscope}%
\begin{pgfscope}%
\definecolor{textcolor}{rgb}{0.000000,0.000000,0.000000}%
\pgfsetstrokecolor{textcolor}%
\pgfsetfillcolor{textcolor}%
\pgftext[x=1.521647in,y=0.365432in,,top]{\color{textcolor}{\rmfamily\fontsize{8.000000}{9.600000}\selectfont\catcode`\^=\active\def^{\ifmmode\sp\else\^{}\fi}\catcode`\%=\active\def%{\%}$\mathdefault{2}$}}%
\end{pgfscope}%
\begin{pgfscope}%
\pgfpathrectangle{\pgfqpoint{0.573768in}{0.462654in}}{\pgfqpoint{2.720000in}{0.906250in}}%
\pgfusepath{clip}%
\pgfsetbuttcap%
\pgfsetroundjoin%
\pgfsetlinewidth{0.501875pt}%
\definecolor{currentstroke}{rgb}{0.690196,0.690196,0.690196}%
\pgfsetstrokecolor{currentstroke}%
\pgfsetstrokeopacity{0.250000}%
\pgfsetdash{{1.850000pt}{0.800000pt}}{0.000000pt}%
\pgfpathmoveto{\pgfqpoint{1.933768in}{0.462654in}}%
\pgfpathlineto{\pgfqpoint{1.933768in}{1.368904in}}%
\pgfusepath{stroke}%
\end{pgfscope}%
\begin{pgfscope}%
\pgfsetbuttcap%
\pgfsetroundjoin%
\definecolor{currentfill}{rgb}{0.000000,0.000000,0.000000}%
\pgfsetfillcolor{currentfill}%
\pgfsetlinewidth{0.803000pt}%
\definecolor{currentstroke}{rgb}{0.000000,0.000000,0.000000}%
\pgfsetstrokecolor{currentstroke}%
\pgfsetdash{}{0pt}%
\pgfsys@defobject{currentmarker}{\pgfqpoint{0.000000in}{-0.048611in}}{\pgfqpoint{0.000000in}{0.000000in}}{%
\pgfpathmoveto{\pgfqpoint{0.000000in}{0.000000in}}%
\pgfpathlineto{\pgfqpoint{0.000000in}{-0.048611in}}%
\pgfusepath{stroke,fill}%
}%
\begin{pgfscope}%
\pgfsys@transformshift{1.933768in}{0.462654in}%
\pgfsys@useobject{currentmarker}{}%
\end{pgfscope}%
\end{pgfscope}%
\begin{pgfscope}%
\definecolor{textcolor}{rgb}{0.000000,0.000000,0.000000}%
\pgfsetstrokecolor{textcolor}%
\pgfsetfillcolor{textcolor}%
\pgftext[x=1.933768in,y=0.365432in,,top]{\color{textcolor}{\rmfamily\fontsize{8.000000}{9.600000}\selectfont\catcode`\^=\active\def^{\ifmmode\sp\else\^{}\fi}\catcode`\%=\active\def%{\%}$\mathdefault{3}$}}%
\end{pgfscope}%
\begin{pgfscope}%
\pgfpathrectangle{\pgfqpoint{0.573768in}{0.462654in}}{\pgfqpoint{2.720000in}{0.906250in}}%
\pgfusepath{clip}%
\pgfsetbuttcap%
\pgfsetroundjoin%
\pgfsetlinewidth{0.501875pt}%
\definecolor{currentstroke}{rgb}{0.690196,0.690196,0.690196}%
\pgfsetstrokecolor{currentstroke}%
\pgfsetstrokeopacity{0.250000}%
\pgfsetdash{{1.850000pt}{0.800000pt}}{0.000000pt}%
\pgfpathmoveto{\pgfqpoint{2.345890in}{0.462654in}}%
\pgfpathlineto{\pgfqpoint{2.345890in}{1.368904in}}%
\pgfusepath{stroke}%
\end{pgfscope}%
\begin{pgfscope}%
\pgfsetbuttcap%
\pgfsetroundjoin%
\definecolor{currentfill}{rgb}{0.000000,0.000000,0.000000}%
\pgfsetfillcolor{currentfill}%
\pgfsetlinewidth{0.803000pt}%
\definecolor{currentstroke}{rgb}{0.000000,0.000000,0.000000}%
\pgfsetstrokecolor{currentstroke}%
\pgfsetdash{}{0pt}%
\pgfsys@defobject{currentmarker}{\pgfqpoint{0.000000in}{-0.048611in}}{\pgfqpoint{0.000000in}{0.000000in}}{%
\pgfpathmoveto{\pgfqpoint{0.000000in}{0.000000in}}%
\pgfpathlineto{\pgfqpoint{0.000000in}{-0.048611in}}%
\pgfusepath{stroke,fill}%
}%
\begin{pgfscope}%
\pgfsys@transformshift{2.345890in}{0.462654in}%
\pgfsys@useobject{currentmarker}{}%
\end{pgfscope}%
\end{pgfscope}%
\begin{pgfscope}%
\definecolor{textcolor}{rgb}{0.000000,0.000000,0.000000}%
\pgfsetstrokecolor{textcolor}%
\pgfsetfillcolor{textcolor}%
\pgftext[x=2.345890in,y=0.365432in,,top]{\color{textcolor}{\rmfamily\fontsize{8.000000}{9.600000}\selectfont\catcode`\^=\active\def^{\ifmmode\sp\else\^{}\fi}\catcode`\%=\active\def%{\%}$\mathdefault{4}$}}%
\end{pgfscope}%
\begin{pgfscope}%
\pgfpathrectangle{\pgfqpoint{0.573768in}{0.462654in}}{\pgfqpoint{2.720000in}{0.906250in}}%
\pgfusepath{clip}%
\pgfsetbuttcap%
\pgfsetroundjoin%
\pgfsetlinewidth{0.501875pt}%
\definecolor{currentstroke}{rgb}{0.690196,0.690196,0.690196}%
\pgfsetstrokecolor{currentstroke}%
\pgfsetstrokeopacity{0.250000}%
\pgfsetdash{{1.850000pt}{0.800000pt}}{0.000000pt}%
\pgfpathmoveto{\pgfqpoint{2.758011in}{0.462654in}}%
\pgfpathlineto{\pgfqpoint{2.758011in}{1.368904in}}%
\pgfusepath{stroke}%
\end{pgfscope}%
\begin{pgfscope}%
\pgfsetbuttcap%
\pgfsetroundjoin%
\definecolor{currentfill}{rgb}{0.000000,0.000000,0.000000}%
\pgfsetfillcolor{currentfill}%
\pgfsetlinewidth{0.803000pt}%
\definecolor{currentstroke}{rgb}{0.000000,0.000000,0.000000}%
\pgfsetstrokecolor{currentstroke}%
\pgfsetdash{}{0pt}%
\pgfsys@defobject{currentmarker}{\pgfqpoint{0.000000in}{-0.048611in}}{\pgfqpoint{0.000000in}{0.000000in}}{%
\pgfpathmoveto{\pgfqpoint{0.000000in}{0.000000in}}%
\pgfpathlineto{\pgfqpoint{0.000000in}{-0.048611in}}%
\pgfusepath{stroke,fill}%
}%
\begin{pgfscope}%
\pgfsys@transformshift{2.758011in}{0.462654in}%
\pgfsys@useobject{currentmarker}{}%
\end{pgfscope}%
\end{pgfscope}%
\begin{pgfscope}%
\definecolor{textcolor}{rgb}{0.000000,0.000000,0.000000}%
\pgfsetstrokecolor{textcolor}%
\pgfsetfillcolor{textcolor}%
\pgftext[x=2.758011in,y=0.365432in,,top]{\color{textcolor}{\rmfamily\fontsize{8.000000}{9.600000}\selectfont\catcode`\^=\active\def^{\ifmmode\sp\else\^{}\fi}\catcode`\%=\active\def%{\%}$\mathdefault{5}$}}%
\end{pgfscope}%
\begin{pgfscope}%
\pgfpathrectangle{\pgfqpoint{0.573768in}{0.462654in}}{\pgfqpoint{2.720000in}{0.906250in}}%
\pgfusepath{clip}%
\pgfsetbuttcap%
\pgfsetroundjoin%
\pgfsetlinewidth{0.501875pt}%
\definecolor{currentstroke}{rgb}{0.690196,0.690196,0.690196}%
\pgfsetstrokecolor{currentstroke}%
\pgfsetstrokeopacity{0.250000}%
\pgfsetdash{{1.850000pt}{0.800000pt}}{0.000000pt}%
\pgfpathmoveto{\pgfqpoint{3.170132in}{0.462654in}}%
\pgfpathlineto{\pgfqpoint{3.170132in}{1.368904in}}%
\pgfusepath{stroke}%
\end{pgfscope}%
\begin{pgfscope}%
\pgfsetbuttcap%
\pgfsetroundjoin%
\definecolor{currentfill}{rgb}{0.000000,0.000000,0.000000}%
\pgfsetfillcolor{currentfill}%
\pgfsetlinewidth{0.803000pt}%
\definecolor{currentstroke}{rgb}{0.000000,0.000000,0.000000}%
\pgfsetstrokecolor{currentstroke}%
\pgfsetdash{}{0pt}%
\pgfsys@defobject{currentmarker}{\pgfqpoint{0.000000in}{-0.048611in}}{\pgfqpoint{0.000000in}{0.000000in}}{%
\pgfpathmoveto{\pgfqpoint{0.000000in}{0.000000in}}%
\pgfpathlineto{\pgfqpoint{0.000000in}{-0.048611in}}%
\pgfusepath{stroke,fill}%
}%
\begin{pgfscope}%
\pgfsys@transformshift{3.170132in}{0.462654in}%
\pgfsys@useobject{currentmarker}{}%
\end{pgfscope}%
\end{pgfscope}%
\begin{pgfscope}%
\definecolor{textcolor}{rgb}{0.000000,0.000000,0.000000}%
\pgfsetstrokecolor{textcolor}%
\pgfsetfillcolor{textcolor}%
\pgftext[x=3.170132in,y=0.365432in,,top]{\color{textcolor}{\rmfamily\fontsize{8.000000}{9.600000}\selectfont\catcode`\^=\active\def^{\ifmmode\sp\else\^{}\fi}\catcode`\%=\active\def%{\%}$\mathdefault{6}$}}%
\end{pgfscope}%
\begin{pgfscope}%
\definecolor{textcolor}{rgb}{0.000000,0.000000,0.000000}%
\pgfsetstrokecolor{textcolor}%
\pgfsetfillcolor{textcolor}%
\pgftext[x=1.933768in,y=0.211111in,,top]{\color{textcolor}{\rmfamily\fontsize{8.000000}{9.600000}\selectfont\catcode`\^=\active\def^{\ifmmode\sp\else\^{}\fi}\catcode`\%=\active\def%{\%}Time (ms)}}%
\end{pgfscope}%
\begin{pgfscope}%
\pgfpathrectangle{\pgfqpoint{0.573768in}{0.462654in}}{\pgfqpoint{2.720000in}{0.906250in}}%
\pgfusepath{clip}%
\pgfsetbuttcap%
\pgfsetroundjoin%
\pgfsetlinewidth{0.501875pt}%
\definecolor{currentstroke}{rgb}{0.690196,0.690196,0.690196}%
\pgfsetstrokecolor{currentstroke}%
\pgfsetstrokeopacity{0.250000}%
\pgfsetdash{{1.850000pt}{0.800000pt}}{0.000000pt}%
\pgfpathmoveto{\pgfqpoint{0.573768in}{0.503789in}}%
\pgfpathlineto{\pgfqpoint{3.293768in}{0.503789in}}%
\pgfusepath{stroke}%
\end{pgfscope}%
\begin{pgfscope}%
\pgfsetbuttcap%
\pgfsetroundjoin%
\definecolor{currentfill}{rgb}{0.000000,0.000000,0.000000}%
\pgfsetfillcolor{currentfill}%
\pgfsetlinewidth{0.803000pt}%
\definecolor{currentstroke}{rgb}{0.000000,0.000000,0.000000}%
\pgfsetstrokecolor{currentstroke}%
\pgfsetdash{}{0pt}%
\pgfsys@defobject{currentmarker}{\pgfqpoint{-0.048611in}{0.000000in}}{\pgfqpoint{-0.000000in}{0.000000in}}{%
\pgfpathmoveto{\pgfqpoint{-0.000000in}{0.000000in}}%
\pgfpathlineto{\pgfqpoint{-0.048611in}{0.000000in}}%
\pgfusepath{stroke,fill}%
}%
\begin{pgfscope}%
\pgfsys@transformshift{0.573768in}{0.503789in}%
\pgfsys@useobject{currentmarker}{}%
\end{pgfscope}%
\end{pgfscope}%
\begin{pgfscope}%
\definecolor{textcolor}{rgb}{0.000000,0.000000,0.000000}%
\pgfsetstrokecolor{textcolor}%
\pgfsetfillcolor{textcolor}%
\pgftext[x=0.266667in, y=0.465209in, left, base]{\color{textcolor}{\rmfamily\fontsize{8.000000}{9.600000}\selectfont\catcode`\^=\active\def^{\ifmmode\sp\else\^{}\fi}\catcode`\%=\active\def%{\%}$\mathdefault{\ensuremath{-}20}$}}%
\end{pgfscope}%
\begin{pgfscope}%
\pgfpathrectangle{\pgfqpoint{0.573768in}{0.462654in}}{\pgfqpoint{2.720000in}{0.906250in}}%
\pgfusepath{clip}%
\pgfsetbuttcap%
\pgfsetroundjoin%
\pgfsetlinewidth{0.501875pt}%
\definecolor{currentstroke}{rgb}{0.690196,0.690196,0.690196}%
\pgfsetstrokecolor{currentstroke}%
\pgfsetstrokeopacity{0.250000}%
\pgfsetdash{{1.850000pt}{0.800000pt}}{0.000000pt}%
\pgfpathmoveto{\pgfqpoint{0.573768in}{0.910790in}}%
\pgfpathlineto{\pgfqpoint{3.293768in}{0.910790in}}%
\pgfusepath{stroke}%
\end{pgfscope}%
\begin{pgfscope}%
\pgfsetbuttcap%
\pgfsetroundjoin%
\definecolor{currentfill}{rgb}{0.000000,0.000000,0.000000}%
\pgfsetfillcolor{currentfill}%
\pgfsetlinewidth{0.803000pt}%
\definecolor{currentstroke}{rgb}{0.000000,0.000000,0.000000}%
\pgfsetstrokecolor{currentstroke}%
\pgfsetdash{}{0pt}%
\pgfsys@defobject{currentmarker}{\pgfqpoint{-0.048611in}{0.000000in}}{\pgfqpoint{-0.000000in}{0.000000in}}{%
\pgfpathmoveto{\pgfqpoint{-0.000000in}{0.000000in}}%
\pgfpathlineto{\pgfqpoint{-0.048611in}{0.000000in}}%
\pgfusepath{stroke,fill}%
}%
\begin{pgfscope}%
\pgfsys@transformshift{0.573768in}{0.910790in}%
\pgfsys@useobject{currentmarker}{}%
\end{pgfscope}%
\end{pgfscope}%
\begin{pgfscope}%
\definecolor{textcolor}{rgb}{0.000000,0.000000,0.000000}%
\pgfsetstrokecolor{textcolor}%
\pgfsetfillcolor{textcolor}%
\pgftext[x=0.266667in, y=0.872210in, left, base]{\color{textcolor}{\rmfamily\fontsize{8.000000}{9.600000}\selectfont\catcode`\^=\active\def^{\ifmmode\sp\else\^{}\fi}\catcode`\%=\active\def%{\%}$\mathdefault{\ensuremath{-}10}$}}%
\end{pgfscope}%
\begin{pgfscope}%
\pgfpathrectangle{\pgfqpoint{0.573768in}{0.462654in}}{\pgfqpoint{2.720000in}{0.906250in}}%
\pgfusepath{clip}%
\pgfsetbuttcap%
\pgfsetroundjoin%
\pgfsetlinewidth{0.501875pt}%
\definecolor{currentstroke}{rgb}{0.690196,0.690196,0.690196}%
\pgfsetstrokecolor{currentstroke}%
\pgfsetstrokeopacity{0.250000}%
\pgfsetdash{{1.850000pt}{0.800000pt}}{0.000000pt}%
\pgfpathmoveto{\pgfqpoint{0.573768in}{1.317791in}}%
\pgfpathlineto{\pgfqpoint{3.293768in}{1.317791in}}%
\pgfusepath{stroke}%
\end{pgfscope}%
\begin{pgfscope}%
\pgfsetbuttcap%
\pgfsetroundjoin%
\definecolor{currentfill}{rgb}{0.000000,0.000000,0.000000}%
\pgfsetfillcolor{currentfill}%
\pgfsetlinewidth{0.803000pt}%
\definecolor{currentstroke}{rgb}{0.000000,0.000000,0.000000}%
\pgfsetstrokecolor{currentstroke}%
\pgfsetdash{}{0pt}%
\pgfsys@defobject{currentmarker}{\pgfqpoint{-0.048611in}{0.000000in}}{\pgfqpoint{-0.000000in}{0.000000in}}{%
\pgfpathmoveto{\pgfqpoint{-0.000000in}{0.000000in}}%
\pgfpathlineto{\pgfqpoint{-0.048611in}{0.000000in}}%
\pgfusepath{stroke,fill}%
}%
\begin{pgfscope}%
\pgfsys@transformshift{0.573768in}{1.317791in}%
\pgfsys@useobject{currentmarker}{}%
\end{pgfscope}%
\end{pgfscope}%
\begin{pgfscope}%
\definecolor{textcolor}{rgb}{0.000000,0.000000,0.000000}%
\pgfsetstrokecolor{textcolor}%
\pgfsetfillcolor{textcolor}%
\pgftext[x=0.417518in, y=1.279211in, left, base]{\color{textcolor}{\rmfamily\fontsize{8.000000}{9.600000}\selectfont\catcode`\^=\active\def^{\ifmmode\sp\else\^{}\fi}\catcode`\%=\active\def%{\%}$\mathdefault{0}$}}%
\end{pgfscope}%
\begin{pgfscope}%
\definecolor{textcolor}{rgb}{0.000000,0.000000,0.000000}%
\pgfsetstrokecolor{textcolor}%
\pgfsetfillcolor{textcolor}%
\pgftext[x=0.211111in,y=0.915779in,,bottom,rotate=90.000000]{\color{textcolor}{\rmfamily\fontsize{8.000000}{9.600000}\selectfont\catcode`\^=\active\def^{\ifmmode\sp\else\^{}\fi}\catcode`\%=\active\def%{\%}$v_D$ (V)}}%
\end{pgfscope}%
\begin{pgfscope}%
\pgfpathrectangle{\pgfqpoint{0.573768in}{0.462654in}}{\pgfqpoint{2.720000in}{0.906250in}}%
\pgfusepath{clip}%
\pgfsetrectcap%
\pgfsetroundjoin%
\pgfsetlinewidth{1.505625pt}%
\definecolor{currentstroke}{rgb}{0.000000,0.619608,0.450980}%
\pgfsetstrokecolor{currentstroke}%
\pgfsetdash{}{0pt}%
\pgfpathmoveto{\pgfqpoint{0.697405in}{1.317791in}}%
\pgfpathlineto{\pgfqpoint{0.697448in}{1.320818in}}%
\pgfpathlineto{\pgfqpoint{0.697702in}{1.320101in}}%
\pgfpathlineto{\pgfqpoint{0.698196in}{0.503848in}}%
\pgfpathlineto{\pgfqpoint{0.698526in}{0.503855in}}%
\pgfpathlineto{\pgfqpoint{0.705819in}{0.504464in}}%
\pgfpathlineto{\pgfqpoint{0.705911in}{0.505016in}}%
\pgfpathlineto{\pgfqpoint{0.706084in}{0.889206in}}%
\pgfpathlineto{\pgfqpoint{0.706853in}{1.326562in}}%
\pgfpathlineto{\pgfqpoint{0.708666in}{1.326480in}}%
\pgfpathlineto{\pgfqpoint{0.710809in}{1.326533in}}%
\pgfpathlineto{\pgfqpoint{0.712952in}{1.326451in}}%
\pgfpathlineto{\pgfqpoint{0.714156in}{1.326477in}}%
\pgfpathlineto{\pgfqpoint{0.714206in}{1.326337in}}%
\pgfpathlineto{\pgfqpoint{0.715490in}{0.504424in}}%
\pgfpathlineto{\pgfqpoint{0.722200in}{0.504761in}}%
\pgfpathlineto{\pgfqpoint{0.722396in}{0.506157in}}%
\pgfpathlineto{\pgfqpoint{0.722571in}{0.910333in}}%
\pgfpathlineto{\pgfqpoint{0.722634in}{1.327201in}}%
\pgfpathlineto{\pgfqpoint{0.723312in}{1.327198in}}%
\pgfpathlineto{\pgfqpoint{0.728093in}{1.327122in}}%
\pgfpathlineto{\pgfqpoint{0.730691in}{1.327029in}}%
\pgfpathlineto{\pgfqpoint{0.731644in}{0.504933in}}%
\pgfpathlineto{\pgfqpoint{0.735271in}{0.504934in}}%
\pgfpathlineto{\pgfqpoint{0.738789in}{0.505584in}}%
\pgfpathlineto{\pgfqpoint{0.738881in}{0.507022in}}%
\pgfpathlineto{\pgfqpoint{0.739056in}{0.923415in}}%
\pgfpathlineto{\pgfqpoint{0.739119in}{1.327515in}}%
\pgfpathlineto{\pgfqpoint{0.739800in}{1.327508in}}%
\pgfpathlineto{\pgfqpoint{0.746229in}{1.327392in}}%
\pgfpathlineto{\pgfqpoint{0.747176in}{1.327328in}}%
\pgfpathlineto{\pgfqpoint{0.747189in}{1.324602in}}%
\pgfpathlineto{\pgfqpoint{0.747996in}{0.505356in}}%
\pgfpathlineto{\pgfqpoint{0.749644in}{0.505110in}}%
\pgfpathlineto{\pgfqpoint{0.755228in}{0.505560in}}%
\pgfpathlineto{\pgfqpoint{0.755366in}{0.507535in}}%
\pgfpathlineto{\pgfqpoint{0.755542in}{0.930359in}}%
\pgfpathlineto{\pgfqpoint{0.755604in}{1.327672in}}%
\pgfpathlineto{\pgfqpoint{0.756286in}{1.327654in}}%
\pgfpathlineto{\pgfqpoint{0.763661in}{1.327445in}}%
\pgfpathlineto{\pgfqpoint{0.763673in}{1.325563in}}%
\pgfpathlineto{\pgfqpoint{0.764489in}{0.505508in}}%
\pgfpathlineto{\pgfqpoint{0.766138in}{0.505222in}}%
\pgfpathlineto{\pgfqpoint{0.771712in}{0.505633in}}%
\pgfpathlineto{\pgfqpoint{0.771851in}{0.507686in}}%
\pgfpathlineto{\pgfqpoint{0.772027in}{0.932337in}}%
\pgfpathlineto{\pgfqpoint{0.772090in}{1.327711in}}%
\pgfpathlineto{\pgfqpoint{0.772788in}{1.327689in}}%
\pgfpathlineto{\pgfqpoint{0.780146in}{1.327439in}}%
\pgfpathlineto{\pgfqpoint{0.780158in}{1.325529in}}%
\pgfpathlineto{\pgfqpoint{0.780969in}{0.505496in}}%
\pgfpathlineto{\pgfqpoint{0.782617in}{0.505196in}}%
\pgfpathlineto{\pgfqpoint{0.788197in}{0.505555in}}%
\pgfpathlineto{\pgfqpoint{0.788336in}{0.507520in}}%
\pgfpathlineto{\pgfqpoint{0.788511in}{0.930168in}}%
\pgfpathlineto{\pgfqpoint{0.788585in}{1.327659in}}%
\pgfpathlineto{\pgfqpoint{0.789318in}{1.327641in}}%
\pgfpathlineto{\pgfqpoint{0.795088in}{1.327434in}}%
\pgfpathlineto{\pgfqpoint{0.796631in}{1.327334in}}%
\pgfpathlineto{\pgfqpoint{0.796643in}{1.324667in}}%
\pgfpathlineto{\pgfqpoint{0.797450in}{0.505341in}}%
\pgfpathlineto{\pgfqpoint{0.799098in}{0.505060in}}%
\pgfpathlineto{\pgfqpoint{0.804682in}{0.505366in}}%
\pgfpathlineto{\pgfqpoint{0.804821in}{0.507118in}}%
\pgfpathlineto{\pgfqpoint{0.804996in}{0.924774in}}%
\pgfpathlineto{\pgfqpoint{0.805068in}{1.327542in}}%
\pgfpathlineto{\pgfqpoint{0.805804in}{1.327519in}}%
\pgfpathlineto{\pgfqpoint{0.811244in}{1.327276in}}%
\pgfpathlineto{\pgfqpoint{0.813115in}{1.327136in}}%
\pgfpathlineto{\pgfqpoint{0.814404in}{0.504909in}}%
\pgfpathlineto{\pgfqpoint{0.821167in}{0.505112in}}%
\pgfpathlineto{\pgfqpoint{0.821305in}{0.506582in}}%
\pgfpathlineto{\pgfqpoint{0.821480in}{0.917159in}}%
\pgfpathlineto{\pgfqpoint{0.821552in}{1.327362in}}%
\pgfpathlineto{\pgfqpoint{0.822283in}{1.327333in}}%
\pgfpathlineto{\pgfqpoint{0.826734in}{1.327082in}}%
\pgfpathlineto{\pgfqpoint{0.829600in}{1.326844in}}%
\pgfpathlineto{\pgfqpoint{0.830455in}{0.504768in}}%
\pgfpathlineto{\pgfqpoint{0.833422in}{0.504629in}}%
\pgfpathlineto{\pgfqpoint{0.837698in}{0.505027in}}%
\pgfpathlineto{\pgfqpoint{0.837790in}{0.506014in}}%
\pgfpathlineto{\pgfqpoint{0.837964in}{0.908116in}}%
\pgfpathlineto{\pgfqpoint{0.838040in}{1.327132in}}%
\pgfpathlineto{\pgfqpoint{0.838779in}{1.327097in}}%
\pgfpathlineto{\pgfqpoint{0.842241in}{1.326842in}}%
\pgfpathlineto{\pgfqpoint{0.844713in}{1.326718in}}%
\pgfpathlineto{\pgfqpoint{0.846085in}{1.326453in}}%
\pgfpathlineto{\pgfqpoint{0.846894in}{0.504516in}}%
\pgfpathlineto{\pgfqpoint{0.849696in}{0.504409in}}%
\pgfpathlineto{\pgfqpoint{0.854183in}{0.504740in}}%
\pgfpathlineto{\pgfqpoint{0.854275in}{0.505502in}}%
\pgfpathlineto{\pgfqpoint{0.854449in}{0.899235in}}%
\pgfpathlineto{\pgfqpoint{0.854524in}{1.326870in}}%
\pgfpathlineto{\pgfqpoint{0.855264in}{1.326829in}}%
\pgfpathlineto{\pgfqpoint{0.857737in}{1.326583in}}%
\pgfpathlineto{\pgfqpoint{0.859880in}{1.326468in}}%
\pgfpathlineto{\pgfqpoint{0.861693in}{1.326200in}}%
\pgfpathlineto{\pgfqpoint{0.862520in}{1.326155in}}%
\pgfpathlineto{\pgfqpoint{0.862570in}{1.325942in}}%
\pgfpathlineto{\pgfqpoint{0.863416in}{0.504263in}}%
\pgfpathlineto{\pgfqpoint{0.870614in}{0.504395in}}%
\pgfpathlineto{\pgfqpoint{0.870727in}{0.504793in}}%
\pgfpathlineto{\pgfqpoint{0.870817in}{0.507598in}}%
\pgfpathlineto{\pgfqpoint{0.870975in}{1.326597in}}%
\pgfpathlineto{\pgfqpoint{0.871681in}{1.326545in}}%
\pgfpathlineto{\pgfqpoint{0.875472in}{1.326143in}}%
\pgfpathlineto{\pgfqpoint{0.877945in}{1.325890in}}%
\pgfpathlineto{\pgfqpoint{0.879055in}{1.325353in}}%
\pgfpathlineto{\pgfqpoint{0.879807in}{0.504134in}}%
\pgfpathlineto{\pgfqpoint{0.882939in}{0.504119in}}%
\pgfpathlineto{\pgfqpoint{0.887268in}{0.505225in}}%
\pgfpathlineto{\pgfqpoint{0.887517in}{1.326400in}}%
\pgfpathlineto{\pgfqpoint{0.888750in}{1.326242in}}%
\pgfpathlineto{\pgfqpoint{0.890893in}{1.326087in}}%
\pgfpathlineto{\pgfqpoint{0.892707in}{1.325776in}}%
\pgfpathlineto{\pgfqpoint{0.894190in}{1.325659in}}%
\pgfpathlineto{\pgfqpoint{0.895540in}{1.324861in}}%
\pgfpathlineto{\pgfqpoint{0.896282in}{0.504043in}}%
\pgfpathlineto{\pgfqpoint{0.903637in}{0.504339in}}%
\pgfpathlineto{\pgfqpoint{0.903730in}{0.504787in}}%
\pgfpathlineto{\pgfqpoint{0.903902in}{0.884116in}}%
\pgfpathlineto{\pgfqpoint{0.903978in}{1.326343in}}%
\pgfpathlineto{\pgfqpoint{0.904707in}{1.326287in}}%
\pgfpathlineto{\pgfqpoint{0.906191in}{1.326053in}}%
\pgfpathlineto{\pgfqpoint{0.907675in}{1.326015in}}%
\pgfpathlineto{\pgfqpoint{0.909158in}{1.325705in}}%
\pgfpathlineto{\pgfqpoint{0.910312in}{1.325702in}}%
\pgfpathlineto{\pgfqpoint{0.911466in}{1.325331in}}%
\pgfpathlineto{\pgfqpoint{0.911975in}{1.325351in}}%
\pgfpathlineto{\pgfqpoint{0.912025in}{1.324837in}}%
\pgfpathlineto{\pgfqpoint{0.912772in}{0.504043in}}%
\pgfpathlineto{\pgfqpoint{0.920122in}{0.504373in}}%
\pgfpathlineto{\pgfqpoint{0.920214in}{0.504847in}}%
\pgfpathlineto{\pgfqpoint{0.920387in}{0.885563in}}%
\pgfpathlineto{\pgfqpoint{0.920463in}{1.326402in}}%
\pgfpathlineto{\pgfqpoint{0.921193in}{1.326354in}}%
\pgfpathlineto{\pgfqpoint{0.923006in}{1.326122in}}%
\pgfpathlineto{\pgfqpoint{0.924490in}{1.326099in}}%
\pgfpathlineto{\pgfqpoint{0.925973in}{1.325833in}}%
\pgfpathlineto{\pgfqpoint{0.927127in}{1.325843in}}%
\pgfpathlineto{\pgfqpoint{0.928434in}{1.325615in}}%
\pgfpathlineto{\pgfqpoint{0.928509in}{1.325233in}}%
\pgfpathlineto{\pgfqpoint{0.929367in}{0.504095in}}%
\pgfpathlineto{\pgfqpoint{0.936607in}{0.504466in}}%
\pgfpathlineto{\pgfqpoint{0.936699in}{0.505015in}}%
\pgfpathlineto{\pgfqpoint{0.936872in}{0.889373in}}%
\pgfpathlineto{\pgfqpoint{0.936948in}{1.326546in}}%
\pgfpathlineto{\pgfqpoint{0.937679in}{1.326507in}}%
\pgfpathlineto{\pgfqpoint{0.939492in}{1.326318in}}%
\pgfpathlineto{\pgfqpoint{0.941306in}{1.326284in}}%
\pgfpathlineto{\pgfqpoint{0.942789in}{1.326075in}}%
\pgfpathlineto{\pgfqpoint{0.944273in}{1.326061in}}%
\pgfpathlineto{\pgfqpoint{0.944994in}{1.325687in}}%
\pgfpathlineto{\pgfqpoint{0.945754in}{0.504230in}}%
\pgfpathlineto{\pgfqpoint{0.948721in}{0.504235in}}%
\pgfpathlineto{\pgfqpoint{0.953144in}{0.504843in}}%
\pgfpathlineto{\pgfqpoint{0.953235in}{0.507170in}}%
\pgfpathlineto{\pgfqpoint{0.953393in}{1.326686in}}%
\pgfpathlineto{\pgfqpoint{0.954247in}{1.326643in}}%
\pgfpathlineto{\pgfqpoint{0.961479in}{1.326050in}}%
\pgfpathlineto{\pgfqpoint{0.962348in}{0.504315in}}%
\pgfpathlineto{\pgfqpoint{0.969526in}{0.504577in}}%
\pgfpathlineto{\pgfqpoint{0.969648in}{0.505203in}}%
\pgfpathlineto{\pgfqpoint{0.969738in}{0.512496in}}%
\pgfpathlineto{\pgfqpoint{0.969896in}{1.326870in}}%
\pgfpathlineto{\pgfqpoint{0.970611in}{1.326850in}}%
\pgfpathlineto{\pgfqpoint{0.973084in}{1.326680in}}%
\pgfpathlineto{\pgfqpoint{0.975227in}{1.326644in}}%
\pgfpathlineto{\pgfqpoint{0.977370in}{1.326461in}}%
\pgfpathlineto{\pgfqpoint{0.977914in}{1.326472in}}%
\pgfpathlineto{\pgfqpoint{0.977964in}{1.326329in}}%
\pgfpathlineto{\pgfqpoint{0.978806in}{0.504454in}}%
\pgfpathlineto{\pgfqpoint{0.981938in}{0.504431in}}%
\pgfpathlineto{\pgfqpoint{0.986062in}{0.504859in}}%
\pgfpathlineto{\pgfqpoint{0.986154in}{0.505718in}}%
\pgfpathlineto{\pgfqpoint{0.986328in}{0.903049in}}%
\pgfpathlineto{\pgfqpoint{0.986403in}{1.326994in}}%
\pgfpathlineto{\pgfqpoint{0.987143in}{1.326973in}}%
\pgfpathlineto{\pgfqpoint{0.989946in}{1.326811in}}%
\pgfpathlineto{\pgfqpoint{0.992419in}{1.326757in}}%
\pgfpathlineto{\pgfqpoint{0.994449in}{1.326512in}}%
\pgfpathlineto{\pgfqpoint{0.995229in}{0.504581in}}%
\pgfpathlineto{\pgfqpoint{0.997702in}{0.504487in}}%
\pgfpathlineto{\pgfqpoint{1.002546in}{0.504959in}}%
\pgfpathlineto{\pgfqpoint{1.002639in}{0.505898in}}%
\pgfpathlineto{\pgfqpoint{1.002813in}{0.906140in}}%
\pgfpathlineto{\pgfqpoint{1.002889in}{1.327083in}}%
\pgfpathlineto{\pgfqpoint{1.003629in}{1.327062in}}%
\pgfpathlineto{\pgfqpoint{1.006761in}{1.326895in}}%
\pgfpathlineto{\pgfqpoint{1.009563in}{1.326825in}}%
\pgfpathlineto{\pgfqpoint{1.010934in}{1.326626in}}%
\pgfpathlineto{\pgfqpoint{1.011796in}{0.504622in}}%
\pgfpathlineto{\pgfqpoint{1.014763in}{0.504563in}}%
\pgfpathlineto{\pgfqpoint{1.019031in}{0.505023in}}%
\pgfpathlineto{\pgfqpoint{1.019124in}{0.506012in}}%
\pgfpathlineto{\pgfqpoint{1.019298in}{0.908040in}}%
\pgfpathlineto{\pgfqpoint{1.019374in}{1.327136in}}%
\pgfpathlineto{\pgfqpoint{1.020114in}{1.327114in}}%
\pgfpathlineto{\pgfqpoint{1.023576in}{1.326935in}}%
\pgfpathlineto{\pgfqpoint{1.026378in}{1.326861in}}%
\pgfpathlineto{\pgfqpoint{1.027418in}{1.326682in}}%
\pgfpathlineto{\pgfqpoint{1.028204in}{0.504691in}}%
\pgfpathlineto{\pgfqpoint{1.030841in}{0.504573in}}%
\pgfpathlineto{\pgfqpoint{1.035516in}{0.505048in}}%
\pgfpathlineto{\pgfqpoint{1.035608in}{0.506056in}}%
\pgfpathlineto{\pgfqpoint{1.035783in}{0.908762in}}%
\pgfpathlineto{\pgfqpoint{1.035858in}{1.327155in}}%
\pgfpathlineto{\pgfqpoint{1.036599in}{1.327132in}}%
\pgfpathlineto{\pgfqpoint{1.040061in}{1.326949in}}%
\pgfpathlineto{\pgfqpoint{1.042863in}{1.326869in}}%
\pgfpathlineto{\pgfqpoint{1.043903in}{1.326689in}}%
\pgfpathlineto{\pgfqpoint{1.044694in}{0.504691in}}%
\pgfpathlineto{\pgfqpoint{1.047332in}{0.504572in}}%
\pgfpathlineto{\pgfqpoint{1.052001in}{0.505037in}}%
\pgfpathlineto{\pgfqpoint{1.052093in}{0.506036in}}%
\pgfpathlineto{\pgfqpoint{1.052267in}{0.908446in}}%
\pgfpathlineto{\pgfqpoint{1.052343in}{1.327146in}}%
\pgfpathlineto{\pgfqpoint{1.053084in}{1.327122in}}%
\pgfpathlineto{\pgfqpoint{1.056545in}{1.326930in}}%
\pgfpathlineto{\pgfqpoint{1.059348in}{1.326843in}}%
\pgfpathlineto{\pgfqpoint{1.060388in}{1.326657in}}%
\pgfpathlineto{\pgfqpoint{1.061304in}{0.504621in}}%
\pgfpathlineto{\pgfqpoint{1.064601in}{0.504574in}}%
\pgfpathlineto{\pgfqpoint{1.068486in}{0.504999in}}%
\pgfpathlineto{\pgfqpoint{1.068578in}{0.505968in}}%
\pgfpathlineto{\pgfqpoint{1.068752in}{0.907320in}}%
\pgfpathlineto{\pgfqpoint{1.068828in}{1.327114in}}%
\pgfpathlineto{\pgfqpoint{1.069568in}{1.327089in}}%
\pgfpathlineto{\pgfqpoint{1.072700in}{1.326902in}}%
\pgfpathlineto{\pgfqpoint{1.075503in}{1.326809in}}%
\pgfpathlineto{\pgfqpoint{1.076873in}{1.326594in}}%
\pgfpathlineto{\pgfqpoint{1.077708in}{0.504604in}}%
\pgfpathlineto{\pgfqpoint{1.080511in}{0.504518in}}%
\pgfpathlineto{\pgfqpoint{1.084971in}{0.504944in}}%
\pgfpathlineto{\pgfqpoint{1.085063in}{0.505869in}}%
\pgfpathlineto{\pgfqpoint{1.085237in}{0.905669in}}%
\pgfpathlineto{\pgfqpoint{1.085313in}{1.327067in}}%
\pgfpathlineto{\pgfqpoint{1.086053in}{1.327040in}}%
\pgfpathlineto{\pgfqpoint{1.089185in}{1.326842in}}%
\pgfpathlineto{\pgfqpoint{1.091658in}{1.326760in}}%
\pgfpathlineto{\pgfqpoint{1.093358in}{1.326510in}}%
\pgfpathlineto{\pgfqpoint{1.094143in}{0.504571in}}%
\pgfpathlineto{\pgfqpoint{1.096616in}{0.504463in}}%
\pgfpathlineto{\pgfqpoint{1.101457in}{0.504889in}}%
\pgfpathlineto{\pgfqpoint{1.101549in}{0.505784in}}%
\pgfpathlineto{\pgfqpoint{1.101722in}{0.904534in}}%
\pgfpathlineto{\pgfqpoint{1.101798in}{1.327012in}}%
\pgfpathlineto{\pgfqpoint{1.102537in}{1.326984in}}%
\pgfpathlineto{\pgfqpoint{1.105669in}{1.326774in}}%
\pgfpathlineto{\pgfqpoint{1.108142in}{1.326687in}}%
\pgfpathlineto{\pgfqpoint{1.109843in}{1.326419in}}%
\pgfpathlineto{\pgfqpoint{1.110667in}{0.504501in}}%
\pgfpathlineto{\pgfqpoint{1.113470in}{0.504429in}}%
\pgfpathlineto{\pgfqpoint{1.117972in}{0.504993in}}%
\pgfpathlineto{\pgfqpoint{1.118060in}{0.506425in}}%
\pgfpathlineto{\pgfqpoint{1.118310in}{1.326953in}}%
\pgfpathlineto{\pgfqpoint{1.119383in}{1.326904in}}%
\pgfpathlineto{\pgfqpoint{1.123175in}{1.326658in}}%
\pgfpathlineto{\pgfqpoint{1.126278in}{1.326478in}}%
\pgfpathlineto{\pgfqpoint{1.126328in}{1.326335in}}%
\pgfpathlineto{\pgfqpoint{1.127168in}{0.504450in}}%
\pgfpathlineto{\pgfqpoint{1.130300in}{0.504404in}}%
\pgfpathlineto{\pgfqpoint{1.134425in}{0.504781in}}%
\pgfpathlineto{\pgfqpoint{1.134517in}{0.505577in}}%
\pgfpathlineto{\pgfqpoint{1.134691in}{0.900569in}}%
\pgfpathlineto{\pgfqpoint{1.134767in}{1.326916in}}%
\pgfpathlineto{\pgfqpoint{1.135507in}{1.326887in}}%
\pgfpathlineto{\pgfqpoint{1.138309in}{1.326677in}}%
\pgfpathlineto{\pgfqpoint{1.140452in}{1.326609in}}%
\pgfpathlineto{\pgfqpoint{1.142639in}{1.326428in}}%
\pgfpathlineto{\pgfqpoint{1.142788in}{1.326416in}}%
\pgfpathlineto{\pgfqpoint{1.143671in}{0.504400in}}%
\pgfpathlineto{\pgfqpoint{1.149275in}{0.504454in}}%
\pgfpathlineto{\pgfqpoint{1.150906in}{0.504737in}}%
\pgfpathlineto{\pgfqpoint{1.150994in}{0.505398in}}%
\pgfpathlineto{\pgfqpoint{1.151084in}{0.533355in}}%
\pgfpathlineto{\pgfqpoint{1.151241in}{1.326888in}}%
\pgfpathlineto{\pgfqpoint{1.151966in}{1.326861in}}%
\pgfpathlineto{\pgfqpoint{1.154439in}{1.326663in}}%
\pgfpathlineto{\pgfqpoint{1.156582in}{1.326602in}}%
\pgfpathlineto{\pgfqpoint{1.158395in}{1.326403in}}%
\pgfpathlineto{\pgfqpoint{1.159247in}{1.326384in}}%
\pgfpathlineto{\pgfqpoint{1.159297in}{1.326217in}}%
\pgfpathlineto{\pgfqpoint{1.160031in}{0.504427in}}%
\pgfpathlineto{\pgfqpoint{1.162998in}{0.504360in}}%
\pgfpathlineto{\pgfqpoint{1.167395in}{0.504738in}}%
\pgfpathlineto{\pgfqpoint{1.167487in}{0.505501in}}%
\pgfpathlineto{\pgfqpoint{1.167661in}{0.899166in}}%
\pgfpathlineto{\pgfqpoint{1.167737in}{1.326873in}}%
\pgfpathlineto{\pgfqpoint{1.168476in}{1.326844in}}%
\pgfpathlineto{\pgfqpoint{1.170949in}{1.326652in}}%
\pgfpathlineto{\pgfqpoint{1.173092in}{1.326588in}}%
\pgfpathlineto{\pgfqpoint{1.175235in}{1.326370in}}%
\pgfpathlineto{\pgfqpoint{1.175732in}{1.326376in}}%
\pgfpathlineto{\pgfqpoint{1.175782in}{1.326208in}}%
\pgfpathlineto{\pgfqpoint{1.176540in}{0.504414in}}%
\pgfpathlineto{\pgfqpoint{1.179837in}{0.504370in}}%
\pgfpathlineto{\pgfqpoint{1.183880in}{0.504741in}}%
\pgfpathlineto{\pgfqpoint{1.183972in}{0.505506in}}%
\pgfpathlineto{\pgfqpoint{1.184146in}{0.899260in}}%
\pgfpathlineto{\pgfqpoint{1.184221in}{1.326876in}}%
\pgfpathlineto{\pgfqpoint{1.184961in}{1.326848in}}%
\pgfpathlineto{\pgfqpoint{1.187434in}{1.326659in}}%
\pgfpathlineto{\pgfqpoint{1.189577in}{1.326598in}}%
\pgfpathlineto{\pgfqpoint{1.191720in}{1.326386in}}%
\pgfpathlineto{\pgfqpoint{1.192217in}{1.326392in}}%
\pgfpathlineto{\pgfqpoint{1.192267in}{1.326227in}}%
\pgfpathlineto{\pgfqpoint{1.193104in}{0.504401in}}%
\pgfpathlineto{\pgfqpoint{1.196730in}{0.504394in}}%
\pgfpathlineto{\pgfqpoint{1.200365in}{0.504757in}}%
\pgfpathlineto{\pgfqpoint{1.200457in}{0.505534in}}%
\pgfpathlineto{\pgfqpoint{1.200630in}{0.899768in}}%
\pgfpathlineto{\pgfqpoint{1.200706in}{1.326892in}}%
\pgfpathlineto{\pgfqpoint{1.201446in}{1.326865in}}%
\pgfpathlineto{\pgfqpoint{1.204248in}{1.326663in}}%
\pgfpathlineto{\pgfqpoint{1.206391in}{1.326604in}}%
\pgfpathlineto{\pgfqpoint{1.208579in}{1.326434in}}%
\pgfpathlineto{\pgfqpoint{1.208727in}{1.326423in}}%
\pgfpathlineto{\pgfqpoint{1.209610in}{0.504406in}}%
\pgfpathlineto{\pgfqpoint{1.215215in}{0.504472in}}%
\pgfpathlineto{\pgfqpoint{1.216847in}{0.504770in}}%
\pgfpathlineto{\pgfqpoint{1.216936in}{0.505487in}}%
\pgfpathlineto{\pgfqpoint{1.217026in}{0.540211in}}%
\pgfpathlineto{\pgfqpoint{1.217184in}{1.326918in}}%
\pgfpathlineto{\pgfqpoint{1.217912in}{1.326893in}}%
\pgfpathlineto{\pgfqpoint{1.220385in}{1.326712in}}%
\pgfpathlineto{\pgfqpoint{1.222528in}{1.326659in}}%
\pgfpathlineto{\pgfqpoint{1.224671in}{1.326458in}}%
\pgfpathlineto{\pgfqpoint{1.225187in}{1.326465in}}%
\pgfpathlineto{\pgfqpoint{1.225237in}{1.326312in}}%
\pgfpathlineto{\pgfqpoint{1.226041in}{0.504457in}}%
\pgfpathlineto{\pgfqpoint{1.228843in}{0.504399in}}%
\pgfpathlineto{\pgfqpoint{1.233334in}{0.504807in}}%
\pgfpathlineto{\pgfqpoint{1.233427in}{0.505624in}}%
\pgfpathlineto{\pgfqpoint{1.233600in}{0.901403in}}%
\pgfpathlineto{\pgfqpoint{1.233676in}{1.326943in}}%
\pgfpathlineto{\pgfqpoint{1.234416in}{1.326917in}}%
\pgfpathlineto{\pgfqpoint{1.237218in}{1.326727in}}%
\pgfpathlineto{\pgfqpoint{1.239691in}{1.326654in}}%
\pgfpathlineto{\pgfqpoint{1.241721in}{1.326366in}}%
\pgfpathlineto{\pgfqpoint{1.242584in}{0.504464in}}%
\pgfpathlineto{\pgfqpoint{1.245716in}{0.504430in}}%
\pgfpathlineto{\pgfqpoint{1.249819in}{0.504832in}}%
\pgfpathlineto{\pgfqpoint{1.249911in}{0.505669in}}%
\pgfpathlineto{\pgfqpoint{1.250085in}{0.902209in}}%
\pgfpathlineto{\pgfqpoint{1.250161in}{1.326967in}}%
\pgfpathlineto{\pgfqpoint{1.250901in}{1.326942in}}%
\pgfpathlineto{\pgfqpoint{1.253703in}{1.326756in}}%
\pgfpathlineto{\pgfqpoint{1.256176in}{1.326684in}}%
\pgfpathlineto{\pgfqpoint{1.258206in}{1.326404in}}%
\pgfpathlineto{\pgfqpoint{1.259062in}{0.504486in}}%
\pgfpathlineto{\pgfqpoint{1.262029in}{0.504446in}}%
\pgfpathlineto{\pgfqpoint{1.266304in}{0.504853in}}%
\pgfpathlineto{\pgfqpoint{1.266396in}{0.505706in}}%
\pgfpathlineto{\pgfqpoint{1.266570in}{0.902859in}}%
\pgfpathlineto{\pgfqpoint{1.266646in}{1.326986in}}%
\pgfpathlineto{\pgfqpoint{1.267386in}{1.326961in}}%
\pgfpathlineto{\pgfqpoint{1.270188in}{1.326778in}}%
\pgfpathlineto{\pgfqpoint{1.272661in}{1.326706in}}%
\pgfpathlineto{\pgfqpoint{1.274691in}{1.326431in}}%
\pgfpathlineto{\pgfqpoint{1.275512in}{0.504513in}}%
\pgfpathlineto{\pgfqpoint{1.278314in}{0.504446in}}%
\pgfpathlineto{\pgfqpoint{1.282789in}{0.504867in}}%
\pgfpathlineto{\pgfqpoint{1.282881in}{0.505731in}}%
\pgfpathlineto{\pgfqpoint{1.283055in}{0.903293in}}%
\pgfpathlineto{\pgfqpoint{1.283131in}{1.326999in}}%
\pgfpathlineto{\pgfqpoint{1.283871in}{1.326974in}}%
\pgfpathlineto{\pgfqpoint{1.286673in}{1.326792in}}%
\pgfpathlineto{\pgfqpoint{1.289146in}{1.326720in}}%
\pgfpathlineto{\pgfqpoint{1.291176in}{1.326447in}}%
\pgfpathlineto{\pgfqpoint{1.291986in}{0.504525in}}%
\pgfpathlineto{\pgfqpoint{1.294788in}{0.504452in}}%
\pgfpathlineto{\pgfqpoint{1.299274in}{0.504873in}}%
\pgfpathlineto{\pgfqpoint{1.299366in}{0.505742in}}%
\pgfpathlineto{\pgfqpoint{1.299540in}{0.903491in}}%
\pgfpathlineto{\pgfqpoint{1.299616in}{1.327005in}}%
\pgfpathlineto{\pgfqpoint{1.300356in}{1.326979in}}%
\pgfpathlineto{\pgfqpoint{1.303158in}{1.326797in}}%
\pgfpathlineto{\pgfqpoint{1.305631in}{1.326724in}}%
\pgfpathlineto{\pgfqpoint{1.307661in}{1.326451in}}%
\pgfpathlineto{\pgfqpoint{1.308470in}{0.504528in}}%
\pgfpathlineto{\pgfqpoint{1.311272in}{0.504453in}}%
\pgfpathlineto{\pgfqpoint{1.315759in}{0.504872in}}%
\pgfpathlineto{\pgfqpoint{1.315851in}{0.505741in}}%
\pgfpathlineto{\pgfqpoint{1.316025in}{0.903470in}}%
\pgfpathlineto{\pgfqpoint{1.316100in}{1.327004in}}%
\pgfpathlineto{\pgfqpoint{1.316840in}{1.326978in}}%
\pgfpathlineto{\pgfqpoint{1.319643in}{1.326795in}}%
\pgfpathlineto{\pgfqpoint{1.322116in}{1.326720in}}%
\pgfpathlineto{\pgfqpoint{1.324146in}{1.326445in}}%
\pgfpathlineto{\pgfqpoint{1.324956in}{0.504523in}}%
\pgfpathlineto{\pgfqpoint{1.327759in}{0.504449in}}%
\pgfpathlineto{\pgfqpoint{1.332243in}{0.504866in}}%
\pgfpathlineto{\pgfqpoint{1.332336in}{0.505729in}}%
\pgfpathlineto{\pgfqpoint{1.332510in}{0.903273in}}%
\pgfpathlineto{\pgfqpoint{1.332585in}{1.326998in}}%
\pgfpathlineto{\pgfqpoint{1.333325in}{1.326972in}}%
\pgfpathlineto{\pgfqpoint{1.336128in}{1.326786in}}%
\pgfpathlineto{\pgfqpoint{1.338600in}{1.326710in}}%
\pgfpathlineto{\pgfqpoint{1.340631in}{1.326432in}}%
\pgfpathlineto{\pgfqpoint{1.341451in}{0.504512in}}%
\pgfpathlineto{\pgfqpoint{1.344253in}{0.504443in}}%
\pgfpathlineto{\pgfqpoint{1.348728in}{0.504856in}}%
\pgfpathlineto{\pgfqpoint{1.348821in}{0.505712in}}%
\pgfpathlineto{\pgfqpoint{1.348994in}{0.902960in}}%
\pgfpathlineto{\pgfqpoint{1.349070in}{1.326989in}}%
\pgfpathlineto{\pgfqpoint{1.349810in}{1.326962in}}%
\pgfpathlineto{\pgfqpoint{1.352612in}{1.326774in}}%
\pgfpathlineto{\pgfqpoint{1.355085in}{1.326697in}}%
\pgfpathlineto{\pgfqpoint{1.357115in}{1.326414in}}%
\pgfpathlineto{\pgfqpoint{1.357968in}{0.504492in}}%
\pgfpathlineto{\pgfqpoint{1.360936in}{0.504445in}}%
\pgfpathlineto{\pgfqpoint{1.365213in}{0.504844in}}%
\pgfpathlineto{\pgfqpoint{1.365305in}{0.505691in}}%
\pgfpathlineto{\pgfqpoint{1.365479in}{0.902596in}}%
\pgfpathlineto{\pgfqpoint{1.365555in}{1.326978in}}%
\pgfpathlineto{\pgfqpoint{1.366295in}{1.326951in}}%
\pgfpathlineto{\pgfqpoint{1.369097in}{1.326761in}}%
\pgfpathlineto{\pgfqpoint{1.371570in}{1.326683in}}%
\pgfpathlineto{\pgfqpoint{1.373600in}{1.326396in}}%
\pgfpathlineto{\pgfqpoint{1.374458in}{0.504480in}}%
\pgfpathlineto{\pgfqpoint{1.377425in}{0.504437in}}%
\pgfpathlineto{\pgfqpoint{1.381698in}{0.504833in}}%
\pgfpathlineto{\pgfqpoint{1.381790in}{0.505671in}}%
\pgfpathlineto{\pgfqpoint{1.381964in}{0.902242in}}%
\pgfpathlineto{\pgfqpoint{1.382040in}{1.326967in}}%
\pgfpathlineto{\pgfqpoint{1.382780in}{1.326940in}}%
\pgfpathlineto{\pgfqpoint{1.385582in}{1.326748in}}%
\pgfpathlineto{\pgfqpoint{1.388055in}{1.326669in}}%
\pgfpathlineto{\pgfqpoint{1.390085in}{1.326379in}}%
\pgfpathlineto{\pgfqpoint{1.390947in}{0.504469in}}%
\pgfpathlineto{\pgfqpoint{1.394079in}{0.504430in}}%
\pgfpathlineto{\pgfqpoint{1.398183in}{0.504824in}}%
\pgfpathlineto{\pgfqpoint{1.398275in}{0.505654in}}%
\pgfpathlineto{\pgfqpoint{1.398449in}{0.901946in}}%
\pgfpathlineto{\pgfqpoint{1.398525in}{1.326959in}}%
\pgfpathlineto{\pgfqpoint{1.399264in}{1.326931in}}%
\pgfpathlineto{\pgfqpoint{1.402067in}{1.326738in}}%
\pgfpathlineto{\pgfqpoint{1.404540in}{1.326659in}}%
\pgfpathlineto{\pgfqpoint{1.406570in}{1.326366in}}%
\pgfpathlineto{\pgfqpoint{1.407432in}{0.504462in}}%
\pgfpathlineto{\pgfqpoint{1.410564in}{0.504424in}}%
\pgfpathlineto{\pgfqpoint{1.414668in}{0.504818in}}%
\pgfpathlineto{\pgfqpoint{1.414760in}{0.505643in}}%
\pgfpathlineto{\pgfqpoint{1.414934in}{0.901742in}}%
\pgfpathlineto{\pgfqpoint{1.415009in}{1.326952in}}%
\pgfpathlineto{\pgfqpoint{1.415749in}{1.326925in}}%
\pgfpathlineto{\pgfqpoint{1.418552in}{1.326731in}}%
\pgfpathlineto{\pgfqpoint{1.421024in}{1.326652in}}%
\pgfpathlineto{\pgfqpoint{1.423055in}{1.326358in}}%
\pgfpathlineto{\pgfqpoint{1.423919in}{0.504458in}}%
\pgfpathlineto{\pgfqpoint{1.427051in}{0.504421in}}%
\pgfpathlineto{\pgfqpoint{1.431152in}{0.504814in}}%
\pgfpathlineto{\pgfqpoint{1.431245in}{0.505637in}}%
\pgfpathlineto{\pgfqpoint{1.431418in}{0.901642in}}%
\pgfpathlineto{\pgfqpoint{1.431494in}{1.326950in}}%
\pgfpathlineto{\pgfqpoint{1.432234in}{1.326923in}}%
\pgfpathlineto{\pgfqpoint{1.435037in}{1.326728in}}%
\pgfpathlineto{\pgfqpoint{1.437509in}{1.326650in}}%
\pgfpathlineto{\pgfqpoint{1.439540in}{1.326355in}}%
\pgfpathlineto{\pgfqpoint{1.440405in}{0.504456in}}%
\pgfpathlineto{\pgfqpoint{1.443537in}{0.504421in}}%
\pgfpathlineto{\pgfqpoint{1.447637in}{0.504814in}}%
\pgfpathlineto{\pgfqpoint{1.447730in}{0.505637in}}%
\pgfpathlineto{\pgfqpoint{1.447903in}{0.901642in}}%
\pgfpathlineto{\pgfqpoint{1.447979in}{1.326950in}}%
\pgfpathlineto{\pgfqpoint{1.448719in}{1.326923in}}%
\pgfpathlineto{\pgfqpoint{1.451521in}{1.326729in}}%
\pgfpathlineto{\pgfqpoint{1.453994in}{1.326651in}}%
\pgfpathlineto{\pgfqpoint{1.456025in}{1.326358in}}%
\pgfpathlineto{\pgfqpoint{1.456889in}{0.504458in}}%
\pgfpathlineto{\pgfqpoint{1.460021in}{0.504422in}}%
\pgfpathlineto{\pgfqpoint{1.464122in}{0.504817in}}%
\pgfpathlineto{\pgfqpoint{1.464214in}{0.505642in}}%
\pgfpathlineto{\pgfqpoint{1.464388in}{0.901723in}}%
\pgfpathlineto{\pgfqpoint{1.464464in}{1.326952in}}%
\pgfpathlineto{\pgfqpoint{1.465204in}{1.326925in}}%
\pgfpathlineto{\pgfqpoint{1.468006in}{1.326733in}}%
\pgfpathlineto{\pgfqpoint{1.470479in}{1.326655in}}%
\pgfpathlineto{\pgfqpoint{1.472509in}{1.326363in}}%
\pgfpathlineto{\pgfqpoint{1.473372in}{0.504461in}}%
\pgfpathlineto{\pgfqpoint{1.476504in}{0.504425in}}%
\pgfpathlineto{\pgfqpoint{1.480607in}{0.504821in}}%
\pgfpathlineto{\pgfqpoint{1.480699in}{0.505649in}}%
\pgfpathlineto{\pgfqpoint{1.480873in}{0.901859in}}%
\pgfpathlineto{\pgfqpoint{1.480949in}{1.326956in}}%
\pgfpathlineto{\pgfqpoint{1.481689in}{1.326930in}}%
\pgfpathlineto{\pgfqpoint{1.484491in}{1.326738in}}%
\pgfpathlineto{\pgfqpoint{1.486964in}{1.326661in}}%
\pgfpathlineto{\pgfqpoint{1.488994in}{1.326371in}}%
\pgfpathlineto{\pgfqpoint{1.489858in}{0.504466in}}%
\pgfpathlineto{\pgfqpoint{1.492990in}{0.504429in}}%
\pgfpathlineto{\pgfqpoint{1.497092in}{0.504826in}}%
\pgfpathlineto{\pgfqpoint{1.497184in}{0.505658in}}%
\pgfpathlineto{\pgfqpoint{1.497358in}{0.902020in}}%
\pgfpathlineto{\pgfqpoint{1.497434in}{1.326961in}}%
\pgfpathlineto{\pgfqpoint{1.498174in}{1.326935in}}%
\pgfpathlineto{\pgfqpoint{1.500976in}{1.326744in}}%
\pgfpathlineto{\pgfqpoint{1.503449in}{1.326668in}}%
\pgfpathlineto{\pgfqpoint{1.505479in}{1.326380in}}%
\pgfpathlineto{\pgfqpoint{1.506340in}{0.504471in}}%
\pgfpathlineto{\pgfqpoint{1.509472in}{0.504432in}}%
\pgfpathlineto{\pgfqpoint{1.513577in}{0.504831in}}%
\pgfpathlineto{\pgfqpoint{1.513669in}{0.505667in}}%
\pgfpathlineto{\pgfqpoint{1.513843in}{0.902178in}}%
\pgfpathlineto{\pgfqpoint{1.513919in}{1.326966in}}%
\pgfpathlineto{\pgfqpoint{1.514658in}{1.326939in}}%
\pgfpathlineto{\pgfqpoint{1.517461in}{1.326750in}}%
\pgfpathlineto{\pgfqpoint{1.519934in}{1.326674in}}%
\pgfpathlineto{\pgfqpoint{1.521964in}{1.326387in}}%
\pgfpathlineto{\pgfqpoint{1.522823in}{0.504475in}}%
\pgfpathlineto{\pgfqpoint{1.525955in}{0.504436in}}%
\pgfpathlineto{\pgfqpoint{1.530062in}{0.504835in}}%
\pgfpathlineto{\pgfqpoint{1.530154in}{0.505675in}}%
\pgfpathlineto{\pgfqpoint{1.530328in}{0.902311in}}%
\pgfpathlineto{\pgfqpoint{1.530403in}{1.326970in}}%
\pgfpathlineto{\pgfqpoint{1.531143in}{1.326943in}}%
\pgfpathlineto{\pgfqpoint{1.533946in}{1.326755in}}%
\pgfpathlineto{\pgfqpoint{1.536418in}{1.326679in}}%
\pgfpathlineto{\pgfqpoint{1.538449in}{1.326393in}}%
\pgfpathlineto{\pgfqpoint{1.539307in}{0.504479in}}%
\pgfpathlineto{\pgfqpoint{1.542439in}{0.504438in}}%
\pgfpathlineto{\pgfqpoint{1.546546in}{0.504838in}}%
\pgfpathlineto{\pgfqpoint{1.546639in}{0.505680in}}%
\pgfpathlineto{\pgfqpoint{1.546812in}{0.902405in}}%
\pgfpathlineto{\pgfqpoint{1.546888in}{1.326973in}}%
\pgfpathlineto{\pgfqpoint{1.547628in}{1.326946in}}%
\pgfpathlineto{\pgfqpoint{1.550431in}{1.326758in}}%
\pgfpathlineto{\pgfqpoint{1.552903in}{1.326682in}}%
\pgfpathlineto{\pgfqpoint{1.554934in}{1.326397in}}%
\pgfpathlineto{\pgfqpoint{1.555791in}{0.504481in}}%
\pgfpathlineto{\pgfqpoint{1.558758in}{0.504440in}}%
\pgfpathlineto{\pgfqpoint{1.563031in}{0.504840in}}%
\pgfpathlineto{\pgfqpoint{1.563124in}{0.505683in}}%
\pgfpathlineto{\pgfqpoint{1.563297in}{0.902453in}}%
\pgfpathlineto{\pgfqpoint{1.563373in}{1.326974in}}%
\pgfpathlineto{\pgfqpoint{1.564113in}{1.326948in}}%
\pgfpathlineto{\pgfqpoint{1.566915in}{1.326759in}}%
\pgfpathlineto{\pgfqpoint{1.569388in}{1.326683in}}%
\pgfpathlineto{\pgfqpoint{1.571418in}{1.326398in}}%
\pgfpathlineto{\pgfqpoint{1.572275in}{0.504482in}}%
\pgfpathlineto{\pgfqpoint{1.575242in}{0.504440in}}%
\pgfpathlineto{\pgfqpoint{1.579516in}{0.504840in}}%
\pgfpathlineto{\pgfqpoint{1.579608in}{0.505683in}}%
\pgfpathlineto{\pgfqpoint{1.579782in}{0.902457in}}%
\pgfpathlineto{\pgfqpoint{1.579858in}{1.326974in}}%
\pgfpathlineto{\pgfqpoint{1.580598in}{1.326948in}}%
\pgfpathlineto{\pgfqpoint{1.583400in}{1.326759in}}%
\pgfpathlineto{\pgfqpoint{1.585873in}{1.326682in}}%
\pgfpathlineto{\pgfqpoint{1.587903in}{1.326397in}}%
\pgfpathlineto{\pgfqpoint{1.588760in}{0.504481in}}%
\pgfpathlineto{\pgfqpoint{1.591727in}{0.504440in}}%
\pgfpathlineto{\pgfqpoint{1.596001in}{0.504839in}}%
\pgfpathlineto{\pgfqpoint{1.596093in}{0.505681in}}%
\pgfpathlineto{\pgfqpoint{1.596267in}{0.902425in}}%
\pgfpathlineto{\pgfqpoint{1.596343in}{1.326973in}}%
\pgfpathlineto{\pgfqpoint{1.597083in}{1.326947in}}%
\pgfpathlineto{\pgfqpoint{1.599885in}{1.326757in}}%
\pgfpathlineto{\pgfqpoint{1.602358in}{1.326681in}}%
\pgfpathlineto{\pgfqpoint{1.604388in}{1.326395in}}%
\pgfpathlineto{\pgfqpoint{1.605246in}{0.504480in}}%
\pgfpathlineto{\pgfqpoint{1.608213in}{0.504439in}}%
\pgfpathlineto{\pgfqpoint{1.612486in}{0.504837in}}%
\pgfpathlineto{\pgfqpoint{1.612578in}{0.505678in}}%
\pgfpathlineto{\pgfqpoint{1.612752in}{0.902368in}}%
\pgfpathlineto{\pgfqpoint{1.612828in}{1.326971in}}%
\pgfpathlineto{\pgfqpoint{1.613568in}{1.326945in}}%
\pgfpathlineto{\pgfqpoint{1.616370in}{1.326755in}}%
\pgfpathlineto{\pgfqpoint{1.618843in}{1.326678in}}%
\pgfpathlineto{\pgfqpoint{1.620873in}{1.326392in}}%
\pgfpathlineto{\pgfqpoint{1.621731in}{0.504478in}}%
\pgfpathlineto{\pgfqpoint{1.624863in}{0.504436in}}%
\pgfpathlineto{\pgfqpoint{1.628971in}{0.504835in}}%
\pgfpathlineto{\pgfqpoint{1.629063in}{0.505674in}}%
\pgfpathlineto{\pgfqpoint{1.629237in}{0.902298in}}%
\pgfpathlineto{\pgfqpoint{1.629313in}{1.326969in}}%
\pgfpathlineto{\pgfqpoint{1.630052in}{1.326943in}}%
\pgfpathlineto{\pgfqpoint{1.632855in}{1.326752in}}%
\pgfpathlineto{\pgfqpoint{1.635328in}{1.326675in}}%
\pgfpathlineto{\pgfqpoint{1.637358in}{1.326388in}}%
\pgfpathlineto{\pgfqpoint{1.638217in}{0.504476in}}%
\pgfpathlineto{\pgfqpoint{1.641349in}{0.504435in}}%
\pgfpathlineto{\pgfqpoint{1.645455in}{0.504833in}}%
\pgfpathlineto{\pgfqpoint{1.645548in}{0.505670in}}%
\pgfpathlineto{\pgfqpoint{1.645722in}{0.902228in}}%
\pgfpathlineto{\pgfqpoint{1.645797in}{1.326967in}}%
\pgfpathlineto{\pgfqpoint{1.646537in}{1.326941in}}%
\pgfpathlineto{\pgfqpoint{1.649340in}{1.326750in}}%
\pgfpathlineto{\pgfqpoint{1.651812in}{1.326673in}}%
\pgfpathlineto{\pgfqpoint{1.653843in}{1.326385in}}%
\pgfpathlineto{\pgfqpoint{1.654703in}{0.504473in}}%
\pgfpathlineto{\pgfqpoint{1.657835in}{0.504433in}}%
\pgfpathlineto{\pgfqpoint{1.661940in}{0.504831in}}%
\pgfpathlineto{\pgfqpoint{1.662033in}{0.505667in}}%
\pgfpathlineto{\pgfqpoint{1.662206in}{0.902168in}}%
\pgfpathlineto{\pgfqpoint{1.662282in}{1.326965in}}%
\pgfpathlineto{\pgfqpoint{1.663022in}{1.326939in}}%
\pgfpathlineto{\pgfqpoint{1.665824in}{1.326748in}}%
\pgfpathlineto{\pgfqpoint{1.668297in}{1.326671in}}%
\pgfpathlineto{\pgfqpoint{1.670328in}{1.326382in}}%
\pgfpathlineto{\pgfqpoint{1.671188in}{0.504472in}}%
\pgfpathlineto{\pgfqpoint{1.674320in}{0.504432in}}%
\pgfpathlineto{\pgfqpoint{1.678425in}{0.504830in}}%
\pgfpathlineto{\pgfqpoint{1.678517in}{0.505664in}}%
\pgfpathlineto{\pgfqpoint{1.678691in}{0.902125in}}%
\pgfpathlineto{\pgfqpoint{1.678767in}{1.326964in}}%
\pgfpathlineto{\pgfqpoint{1.679507in}{1.326937in}}%
\pgfpathlineto{\pgfqpoint{1.682309in}{1.326746in}}%
\pgfpathlineto{\pgfqpoint{1.684782in}{1.326669in}}%
\pgfpathlineto{\pgfqpoint{1.686812in}{1.326380in}}%
\pgfpathlineto{\pgfqpoint{1.687673in}{0.504471in}}%
\pgfpathlineto{\pgfqpoint{1.690806in}{0.504432in}}%
\pgfpathlineto{\pgfqpoint{1.694910in}{0.504829in}}%
\pgfpathlineto{\pgfqpoint{1.695002in}{0.505663in}}%
\pgfpathlineto{\pgfqpoint{1.695176in}{0.902102in}}%
\pgfpathlineto{\pgfqpoint{1.695252in}{1.326963in}}%
\pgfpathlineto{\pgfqpoint{1.695992in}{1.326937in}}%
\pgfpathlineto{\pgfqpoint{1.698794in}{1.326746in}}%
\pgfpathlineto{\pgfqpoint{1.701267in}{1.326669in}}%
\pgfpathlineto{\pgfqpoint{1.703297in}{1.326380in}}%
\pgfpathlineto{\pgfqpoint{1.704159in}{0.504470in}}%
\pgfpathlineto{\pgfqpoint{1.707291in}{0.504431in}}%
\pgfpathlineto{\pgfqpoint{1.711395in}{0.504829in}}%
\pgfpathlineto{\pgfqpoint{1.711487in}{0.505663in}}%
\pgfpathlineto{\pgfqpoint{1.711661in}{0.902098in}}%
\pgfpathlineto{\pgfqpoint{1.711737in}{1.326963in}}%
\pgfpathlineto{\pgfqpoint{1.712477in}{1.326937in}}%
\pgfpathlineto{\pgfqpoint{1.715279in}{1.326746in}}%
\pgfpathlineto{\pgfqpoint{1.717752in}{1.326669in}}%
\pgfpathlineto{\pgfqpoint{1.719782in}{1.326380in}}%
\pgfpathlineto{\pgfqpoint{1.720643in}{0.504471in}}%
\pgfpathlineto{\pgfqpoint{1.723775in}{0.504432in}}%
\pgfpathlineto{\pgfqpoint{1.727880in}{0.504829in}}%
\pgfpathlineto{\pgfqpoint{1.727972in}{0.505663in}}%
\pgfpathlineto{\pgfqpoint{1.728146in}{0.902111in}}%
\pgfpathlineto{\pgfqpoint{1.728222in}{1.326964in}}%
\pgfpathlineto{\pgfqpoint{1.728961in}{1.326937in}}%
\pgfpathlineto{\pgfqpoint{1.731764in}{1.326746in}}%
\pgfpathlineto{\pgfqpoint{1.734237in}{1.326669in}}%
\pgfpathlineto{\pgfqpoint{1.736267in}{1.326381in}}%
\pgfpathlineto{\pgfqpoint{1.737128in}{0.504471in}}%
\pgfpathlineto{\pgfqpoint{1.740260in}{0.504432in}}%
\pgfpathlineto{\pgfqpoint{1.744365in}{0.504830in}}%
\pgfpathlineto{\pgfqpoint{1.744457in}{0.505665in}}%
\pgfpathlineto{\pgfqpoint{1.744631in}{0.902135in}}%
\pgfpathlineto{\pgfqpoint{1.744706in}{1.326964in}}%
\pgfpathlineto{\pgfqpoint{1.745446in}{1.326938in}}%
\pgfpathlineto{\pgfqpoint{1.748249in}{1.326747in}}%
\pgfpathlineto{\pgfqpoint{1.750721in}{1.326670in}}%
\pgfpathlineto{\pgfqpoint{1.752752in}{1.326382in}}%
\pgfpathlineto{\pgfqpoint{1.753612in}{0.504472in}}%
\pgfpathlineto{\pgfqpoint{1.756744in}{0.504433in}}%
\pgfpathlineto{\pgfqpoint{1.760849in}{0.504831in}}%
\pgfpathlineto{\pgfqpoint{1.760942in}{0.505667in}}%
\pgfpathlineto{\pgfqpoint{1.761115in}{0.902165in}}%
\pgfpathlineto{\pgfqpoint{1.761191in}{1.326965in}}%
\pgfpathlineto{\pgfqpoint{1.761931in}{1.326939in}}%
\pgfpathlineto{\pgfqpoint{1.764734in}{1.326748in}}%
\pgfpathlineto{\pgfqpoint{1.767206in}{1.326672in}}%
\pgfpathlineto{\pgfqpoint{1.769237in}{1.326384in}}%
\pgfpathlineto{\pgfqpoint{1.770097in}{0.504473in}}%
\pgfpathlineto{\pgfqpoint{1.773229in}{0.504434in}}%
\pgfpathlineto{\pgfqpoint{1.777334in}{0.504832in}}%
\pgfpathlineto{\pgfqpoint{1.777427in}{0.505668in}}%
\pgfpathlineto{\pgfqpoint{1.777600in}{0.902196in}}%
\pgfpathlineto{\pgfqpoint{1.777676in}{1.326966in}}%
\pgfpathlineto{\pgfqpoint{1.778416in}{1.326940in}}%
\pgfpathlineto{\pgfqpoint{1.781218in}{1.326750in}}%
\pgfpathlineto{\pgfqpoint{1.783691in}{1.326673in}}%
\pgfpathlineto{\pgfqpoint{1.785721in}{1.326385in}}%
\pgfpathlineto{\pgfqpoint{1.786581in}{0.504474in}}%
\pgfpathlineto{\pgfqpoint{1.789713in}{0.504434in}}%
\pgfpathlineto{\pgfqpoint{1.793819in}{0.504833in}}%
\pgfpathlineto{\pgfqpoint{1.793911in}{0.505670in}}%
\pgfpathlineto{\pgfqpoint{1.794085in}{0.902223in}}%
\pgfpathlineto{\pgfqpoint{1.794161in}{1.326967in}}%
\pgfpathlineto{\pgfqpoint{1.794901in}{1.326941in}}%
\pgfpathlineto{\pgfqpoint{1.797703in}{1.326750in}}%
\pgfpathlineto{\pgfqpoint{1.800176in}{1.326674in}}%
\pgfpathlineto{\pgfqpoint{1.802206in}{1.326387in}}%
\pgfpathlineto{\pgfqpoint{1.803066in}{0.504475in}}%
\pgfpathlineto{\pgfqpoint{1.806198in}{0.504435in}}%
\pgfpathlineto{\pgfqpoint{1.810304in}{0.504833in}}%
\pgfpathlineto{\pgfqpoint{1.810396in}{0.505671in}}%
\pgfpathlineto{\pgfqpoint{1.810570in}{0.902243in}}%
\pgfpathlineto{\pgfqpoint{1.810646in}{1.326968in}}%
\pgfpathlineto{\pgfqpoint{1.811386in}{1.326941in}}%
\pgfpathlineto{\pgfqpoint{1.814188in}{1.326751in}}%
\pgfpathlineto{\pgfqpoint{1.816661in}{1.326675in}}%
\pgfpathlineto{\pgfqpoint{1.818691in}{1.326387in}}%
\pgfpathlineto{\pgfqpoint{1.819550in}{0.504475in}}%
\pgfpathlineto{\pgfqpoint{1.822683in}{0.504435in}}%
\pgfpathlineto{\pgfqpoint{1.826789in}{0.504834in}}%
\pgfpathlineto{\pgfqpoint{1.826881in}{0.505672in}}%
\pgfpathlineto{\pgfqpoint{1.827055in}{0.902254in}}%
\pgfpathlineto{\pgfqpoint{1.827131in}{1.326968in}}%
\pgfpathlineto{\pgfqpoint{1.827871in}{1.326941in}}%
\pgfpathlineto{\pgfqpoint{1.830673in}{1.326751in}}%
\pgfpathlineto{\pgfqpoint{1.833146in}{1.326675in}}%
\pgfpathlineto{\pgfqpoint{1.835176in}{1.326388in}}%
\pgfpathlineto{\pgfqpoint{1.836035in}{0.504475in}}%
\pgfpathlineto{\pgfqpoint{1.839167in}{0.504435in}}%
\pgfpathlineto{\pgfqpoint{1.843274in}{0.504834in}}%
\pgfpathlineto{\pgfqpoint{1.843366in}{0.505672in}}%
\pgfpathlineto{\pgfqpoint{1.843540in}{0.902257in}}%
\pgfpathlineto{\pgfqpoint{1.843616in}{1.326968in}}%
\pgfpathlineto{\pgfqpoint{1.844355in}{1.326942in}}%
\pgfpathlineto{\pgfqpoint{1.847158in}{1.326752in}}%
\pgfpathlineto{\pgfqpoint{1.849631in}{1.326675in}}%
\pgfpathlineto{\pgfqpoint{1.851661in}{1.326388in}}%
\pgfpathlineto{\pgfqpoint{1.852520in}{0.504475in}}%
\pgfpathlineto{\pgfqpoint{1.855652in}{0.504435in}}%
\pgfpathlineto{\pgfqpoint{1.859759in}{0.504834in}}%
\pgfpathlineto{\pgfqpoint{1.859851in}{0.505671in}}%
\pgfpathlineto{\pgfqpoint{1.860025in}{0.902252in}}%
\pgfpathlineto{\pgfqpoint{1.860100in}{1.326968in}}%
\pgfpathlineto{\pgfqpoint{1.860840in}{1.326941in}}%
\pgfpathlineto{\pgfqpoint{1.863643in}{1.326751in}}%
\pgfpathlineto{\pgfqpoint{1.866115in}{1.326675in}}%
\pgfpathlineto{\pgfqpoint{1.868146in}{1.326387in}}%
\pgfpathlineto{\pgfqpoint{1.869005in}{0.504475in}}%
\pgfpathlineto{\pgfqpoint{1.872137in}{0.504435in}}%
\pgfpathlineto{\pgfqpoint{1.876243in}{0.504833in}}%
\pgfpathlineto{\pgfqpoint{1.876336in}{0.505671in}}%
\pgfpathlineto{\pgfqpoint{1.876509in}{0.902241in}}%
\pgfpathlineto{\pgfqpoint{1.876585in}{1.326968in}}%
\pgfpathlineto{\pgfqpoint{1.877325in}{1.326941in}}%
\pgfpathlineto{\pgfqpoint{1.880128in}{1.326751in}}%
\pgfpathlineto{\pgfqpoint{1.882600in}{1.326674in}}%
\pgfpathlineto{\pgfqpoint{1.884631in}{1.326387in}}%
\pgfpathlineto{\pgfqpoint{1.885490in}{0.504475in}}%
\pgfpathlineto{\pgfqpoint{1.888622in}{0.504435in}}%
\pgfpathlineto{\pgfqpoint{1.892728in}{0.504833in}}%
\pgfpathlineto{\pgfqpoint{1.892821in}{0.505670in}}%
\pgfpathlineto{\pgfqpoint{1.892994in}{0.902228in}}%
\pgfpathlineto{\pgfqpoint{1.893070in}{1.326967in}}%
\pgfpathlineto{\pgfqpoint{1.893810in}{1.326941in}}%
\pgfpathlineto{\pgfqpoint{1.896612in}{1.326750in}}%
\pgfpathlineto{\pgfqpoint{1.899085in}{1.326673in}}%
\pgfpathlineto{\pgfqpoint{1.901115in}{1.326386in}}%
\pgfpathlineto{\pgfqpoint{1.901975in}{0.504474in}}%
\pgfpathlineto{\pgfqpoint{1.905107in}{0.504434in}}%
\pgfpathlineto{\pgfqpoint{1.909213in}{0.504832in}}%
\pgfpathlineto{\pgfqpoint{1.909305in}{0.505669in}}%
\pgfpathlineto{\pgfqpoint{1.909479in}{0.902215in}}%
\pgfpathlineto{\pgfqpoint{1.909555in}{1.326967in}}%
\pgfpathlineto{\pgfqpoint{1.910295in}{1.326940in}}%
\pgfpathlineto{\pgfqpoint{1.913097in}{1.326750in}}%
\pgfpathlineto{\pgfqpoint{1.915570in}{1.326673in}}%
\pgfpathlineto{\pgfqpoint{1.917600in}{1.326385in}}%
\pgfpathlineto{\pgfqpoint{1.918460in}{0.504474in}}%
\pgfpathlineto{\pgfqpoint{1.921592in}{0.504434in}}%
\pgfpathlineto{\pgfqpoint{1.925698in}{0.504832in}}%
\pgfpathlineto{\pgfqpoint{1.925790in}{0.505669in}}%
\pgfpathlineto{\pgfqpoint{1.925964in}{0.902202in}}%
\pgfpathlineto{\pgfqpoint{1.926040in}{1.326966in}}%
\pgfpathlineto{\pgfqpoint{1.926780in}{1.326940in}}%
\pgfpathlineto{\pgfqpoint{1.929582in}{1.326749in}}%
\pgfpathlineto{\pgfqpoint{1.932055in}{1.326673in}}%
\pgfpathlineto{\pgfqpoint{1.934085in}{1.326385in}}%
\pgfpathlineto{\pgfqpoint{1.934945in}{0.504474in}}%
\pgfpathlineto{\pgfqpoint{1.938077in}{0.504434in}}%
\pgfpathlineto{\pgfqpoint{1.942183in}{0.504832in}}%
\pgfpathlineto{\pgfqpoint{1.942275in}{0.505668in}}%
\pgfpathlineto{\pgfqpoint{1.942449in}{0.902193in}}%
\pgfpathlineto{\pgfqpoint{1.942525in}{1.326966in}}%
\pgfpathlineto{\pgfqpoint{1.943264in}{1.326940in}}%
\pgfpathlineto{\pgfqpoint{1.946067in}{1.326749in}}%
\pgfpathlineto{\pgfqpoint{1.948540in}{1.326672in}}%
\pgfpathlineto{\pgfqpoint{1.950570in}{1.326384in}}%
\pgfpathlineto{\pgfqpoint{1.951430in}{0.504473in}}%
\pgfpathlineto{\pgfqpoint{1.954562in}{0.504434in}}%
\pgfpathlineto{\pgfqpoint{1.958668in}{0.504832in}}%
\pgfpathlineto{\pgfqpoint{1.958760in}{0.505668in}}%
\pgfpathlineto{\pgfqpoint{1.958934in}{0.902188in}}%
\pgfpathlineto{\pgfqpoint{1.959009in}{1.326966in}}%
\pgfpathlineto{\pgfqpoint{1.959749in}{1.326939in}}%
\pgfpathlineto{\pgfqpoint{1.962552in}{1.326749in}}%
\pgfpathlineto{\pgfqpoint{1.965024in}{1.326672in}}%
\pgfpathlineto{\pgfqpoint{1.967055in}{1.326384in}}%
\pgfpathlineto{\pgfqpoint{1.967915in}{0.504473in}}%
\pgfpathlineto{\pgfqpoint{1.971047in}{0.504433in}}%
\pgfpathlineto{\pgfqpoint{1.975152in}{0.504832in}}%
\pgfpathlineto{\pgfqpoint{1.975245in}{0.505668in}}%
\pgfpathlineto{\pgfqpoint{1.975419in}{0.902187in}}%
\pgfpathlineto{\pgfqpoint{1.975494in}{1.326966in}}%
\pgfpathlineto{\pgfqpoint{1.976234in}{1.326939in}}%
\pgfpathlineto{\pgfqpoint{1.979037in}{1.326749in}}%
\pgfpathlineto{\pgfqpoint{1.981509in}{1.326672in}}%
\pgfpathlineto{\pgfqpoint{1.983540in}{1.326384in}}%
\pgfpathlineto{\pgfqpoint{1.984400in}{0.504473in}}%
\pgfpathlineto{\pgfqpoint{1.987532in}{0.504434in}}%
\pgfpathlineto{\pgfqpoint{1.991637in}{0.504832in}}%
\pgfpathlineto{\pgfqpoint{1.991730in}{0.505668in}}%
\pgfpathlineto{\pgfqpoint{1.991903in}{0.902189in}}%
\pgfpathlineto{\pgfqpoint{1.991979in}{1.326966in}}%
\pgfpathlineto{\pgfqpoint{1.992719in}{1.326939in}}%
\pgfpathlineto{\pgfqpoint{1.995521in}{1.326749in}}%
\pgfpathlineto{\pgfqpoint{1.997994in}{1.326672in}}%
\pgfpathlineto{\pgfqpoint{2.000025in}{1.326384in}}%
\pgfpathlineto{\pgfqpoint{2.000885in}{0.504473in}}%
\pgfpathlineto{\pgfqpoint{2.004017in}{0.504434in}}%
\pgfpathlineto{\pgfqpoint{2.008122in}{0.504832in}}%
\pgfpathlineto{\pgfqpoint{2.008214in}{0.505668in}}%
\pgfpathlineto{\pgfqpoint{2.008388in}{0.902193in}}%
\pgfpathlineto{\pgfqpoint{2.008464in}{1.326966in}}%
\pgfpathlineto{\pgfqpoint{2.009204in}{1.326940in}}%
\pgfpathlineto{\pgfqpoint{2.012006in}{1.326749in}}%
\pgfpathlineto{\pgfqpoint{2.014479in}{1.326672in}}%
\pgfpathlineto{\pgfqpoint{2.016509in}{1.326385in}}%
\pgfpathlineto{\pgfqpoint{2.017369in}{0.504474in}}%
\pgfpathlineto{\pgfqpoint{2.020501in}{0.504434in}}%
\pgfpathlineto{\pgfqpoint{2.024607in}{0.504832in}}%
\pgfpathlineto{\pgfqpoint{2.024699in}{0.505668in}}%
\pgfpathlineto{\pgfqpoint{2.024873in}{0.902199in}}%
\pgfpathlineto{\pgfqpoint{2.024949in}{1.326966in}}%
\pgfpathlineto{\pgfqpoint{2.025689in}{1.326940in}}%
\pgfpathlineto{\pgfqpoint{2.028491in}{1.326749in}}%
\pgfpathlineto{\pgfqpoint{2.030964in}{1.326673in}}%
\pgfpathlineto{\pgfqpoint{2.032994in}{1.326385in}}%
\pgfpathlineto{\pgfqpoint{2.033854in}{0.504474in}}%
\pgfpathlineto{\pgfqpoint{2.036986in}{0.504434in}}%
\pgfpathlineto{\pgfqpoint{2.041092in}{0.504832in}}%
\pgfpathlineto{\pgfqpoint{2.041184in}{0.505669in}}%
\pgfpathlineto{\pgfqpoint{2.041358in}{0.902205in}}%
\pgfpathlineto{\pgfqpoint{2.041434in}{1.326967in}}%
\pgfpathlineto{\pgfqpoint{2.042174in}{1.326940in}}%
\pgfpathlineto{\pgfqpoint{2.044976in}{1.326750in}}%
\pgfpathlineto{\pgfqpoint{2.047449in}{1.326673in}}%
\pgfpathlineto{\pgfqpoint{2.049479in}{1.326385in}}%
\pgfpathlineto{\pgfqpoint{2.050339in}{0.504474in}}%
\pgfpathlineto{\pgfqpoint{2.053471in}{0.504434in}}%
\pgfpathlineto{\pgfqpoint{2.057577in}{0.504832in}}%
\pgfpathlineto{\pgfqpoint{2.057669in}{0.505669in}}%
\pgfpathlineto{\pgfqpoint{2.057843in}{0.902210in}}%
\pgfpathlineto{\pgfqpoint{2.057919in}{1.326967in}}%
\pgfpathlineto{\pgfqpoint{2.058658in}{1.326940in}}%
\pgfpathlineto{\pgfqpoint{2.061461in}{1.326750in}}%
\pgfpathlineto{\pgfqpoint{2.063934in}{1.326673in}}%
\pgfpathlineto{\pgfqpoint{2.065964in}{1.326386in}}%
\pgfpathlineto{\pgfqpoint{2.066824in}{0.504474in}}%
\pgfpathlineto{\pgfqpoint{2.069956in}{0.504434in}}%
\pgfpathlineto{\pgfqpoint{2.074062in}{0.504832in}}%
\pgfpathlineto{\pgfqpoint{2.074154in}{0.505669in}}%
\pgfpathlineto{\pgfqpoint{2.074328in}{0.902214in}}%
\pgfpathlineto{\pgfqpoint{2.074403in}{1.326967in}}%
\pgfpathlineto{\pgfqpoint{2.075143in}{1.326940in}}%
\pgfpathlineto{\pgfqpoint{2.077946in}{1.326750in}}%
\pgfpathlineto{\pgfqpoint{2.080418in}{1.326673in}}%
\pgfpathlineto{\pgfqpoint{2.082449in}{1.326386in}}%
\pgfpathlineto{\pgfqpoint{2.083308in}{0.504474in}}%
\pgfpathlineto{\pgfqpoint{2.086441in}{0.504434in}}%
\pgfpathlineto{\pgfqpoint{2.090546in}{0.504833in}}%
\pgfpathlineto{\pgfqpoint{2.090639in}{0.505669in}}%
\pgfpathlineto{\pgfqpoint{2.090812in}{0.902217in}}%
\pgfpathlineto{\pgfqpoint{2.090888in}{1.326967in}}%
\pgfpathlineto{\pgfqpoint{2.091628in}{1.326940in}}%
\pgfpathlineto{\pgfqpoint{2.094431in}{1.326750in}}%
\pgfpathlineto{\pgfqpoint{2.096903in}{1.326673in}}%
\pgfpathlineto{\pgfqpoint{2.098934in}{1.326386in}}%
\pgfpathlineto{\pgfqpoint{2.099793in}{0.504474in}}%
\pgfpathlineto{\pgfqpoint{2.102925in}{0.504434in}}%
\pgfpathlineto{\pgfqpoint{2.107031in}{0.504833in}}%
\pgfpathlineto{\pgfqpoint{2.107124in}{0.505670in}}%
\pgfpathlineto{\pgfqpoint{2.107297in}{0.902217in}}%
\pgfpathlineto{\pgfqpoint{2.107373in}{1.326967in}}%
\pgfpathlineto{\pgfqpoint{2.108113in}{1.326940in}}%
\pgfpathlineto{\pgfqpoint{2.110915in}{1.326750in}}%
\pgfpathlineto{\pgfqpoint{2.113388in}{1.326673in}}%
\pgfpathlineto{\pgfqpoint{2.115418in}{1.326386in}}%
\pgfpathlineto{\pgfqpoint{2.116278in}{0.504474in}}%
\pgfpathlineto{\pgfqpoint{2.119410in}{0.504434in}}%
\pgfpathlineto{\pgfqpoint{2.123516in}{0.504833in}}%
\pgfpathlineto{\pgfqpoint{2.123608in}{0.505669in}}%
\pgfpathlineto{\pgfqpoint{2.123782in}{0.902217in}}%
\pgfpathlineto{\pgfqpoint{2.123858in}{1.326967in}}%
\pgfpathlineto{\pgfqpoint{2.124598in}{1.326940in}}%
\pgfpathlineto{\pgfqpoint{2.127400in}{1.326750in}}%
\pgfpathlineto{\pgfqpoint{2.129873in}{1.326673in}}%
\pgfpathlineto{\pgfqpoint{2.131903in}{1.326386in}}%
\pgfpathlineto{\pgfqpoint{2.132763in}{0.504474in}}%
\pgfpathlineto{\pgfqpoint{2.135895in}{0.504434in}}%
\pgfpathlineto{\pgfqpoint{2.140001in}{0.504832in}}%
\pgfpathlineto{\pgfqpoint{2.140093in}{0.505669in}}%
\pgfpathlineto{\pgfqpoint{2.140267in}{0.902215in}}%
\pgfpathlineto{\pgfqpoint{2.140343in}{1.326967in}}%
\pgfpathlineto{\pgfqpoint{2.141083in}{1.326940in}}%
\pgfpathlineto{\pgfqpoint{2.143885in}{1.326750in}}%
\pgfpathlineto{\pgfqpoint{2.146358in}{1.326673in}}%
\pgfpathlineto{\pgfqpoint{2.148388in}{1.326386in}}%
\pgfpathlineto{\pgfqpoint{2.149248in}{0.504474in}}%
\pgfpathlineto{\pgfqpoint{2.152380in}{0.504434in}}%
\pgfpathlineto{\pgfqpoint{2.156486in}{0.504832in}}%
\pgfpathlineto{\pgfqpoint{2.156578in}{0.505669in}}%
\pgfpathlineto{\pgfqpoint{2.156752in}{0.902212in}}%
\pgfpathlineto{\pgfqpoint{2.156828in}{1.326967in}}%
\pgfpathlineto{\pgfqpoint{2.157568in}{1.326940in}}%
\pgfpathlineto{\pgfqpoint{2.160370in}{1.326750in}}%
\pgfpathlineto{\pgfqpoint{2.162843in}{1.326673in}}%
\pgfpathlineto{\pgfqpoint{2.164873in}{1.326385in}}%
\pgfpathlineto{\pgfqpoint{2.165733in}{0.504474in}}%
\pgfpathlineto{\pgfqpoint{2.168865in}{0.504434in}}%
\pgfpathlineto{\pgfqpoint{2.172971in}{0.504832in}}%
\pgfpathlineto{\pgfqpoint{2.173063in}{0.505669in}}%
\pgfpathlineto{\pgfqpoint{2.173237in}{0.902210in}}%
\pgfpathlineto{\pgfqpoint{2.173313in}{1.326967in}}%
\pgfpathlineto{\pgfqpoint{2.174052in}{1.326940in}}%
\pgfpathlineto{\pgfqpoint{2.176855in}{1.326750in}}%
\pgfpathlineto{\pgfqpoint{2.179328in}{1.326673in}}%
\pgfpathlineto{\pgfqpoint{2.181358in}{1.326385in}}%
\pgfpathlineto{\pgfqpoint{2.182218in}{0.504474in}}%
\pgfpathlineto{\pgfqpoint{2.185350in}{0.504434in}}%
\pgfpathlineto{\pgfqpoint{2.189455in}{0.504832in}}%
\pgfpathlineto{\pgfqpoint{2.189548in}{0.505669in}}%
\pgfpathlineto{\pgfqpoint{2.189722in}{0.902207in}}%
\pgfpathlineto{\pgfqpoint{2.189797in}{1.326967in}}%
\pgfpathlineto{\pgfqpoint{2.190537in}{1.326940in}}%
\pgfpathlineto{\pgfqpoint{2.193340in}{1.326750in}}%
\pgfpathlineto{\pgfqpoint{2.195812in}{1.326673in}}%
\pgfpathlineto{\pgfqpoint{2.197843in}{1.326385in}}%
\pgfpathlineto{\pgfqpoint{2.198703in}{0.504474in}}%
\pgfpathlineto{\pgfqpoint{2.201835in}{0.504434in}}%
\pgfpathlineto{\pgfqpoint{2.205940in}{0.504832in}}%
\pgfpathlineto{\pgfqpoint{2.206033in}{0.505669in}}%
\pgfpathlineto{\pgfqpoint{2.206206in}{0.902206in}}%
\pgfpathlineto{\pgfqpoint{2.206282in}{1.326967in}}%
\pgfpathlineto{\pgfqpoint{2.207022in}{1.326940in}}%
\pgfpathlineto{\pgfqpoint{2.209824in}{1.326750in}}%
\pgfpathlineto{\pgfqpoint{2.212297in}{1.326673in}}%
\pgfpathlineto{\pgfqpoint{2.214328in}{1.326385in}}%
\pgfpathlineto{\pgfqpoint{2.215187in}{0.504474in}}%
\pgfpathlineto{\pgfqpoint{2.218320in}{0.504434in}}%
\pgfpathlineto{\pgfqpoint{2.222425in}{0.504832in}}%
\pgfpathlineto{\pgfqpoint{2.222517in}{0.505669in}}%
\pgfpathlineto{\pgfqpoint{2.222691in}{0.902204in}}%
\pgfpathlineto{\pgfqpoint{2.222767in}{1.326967in}}%
\pgfpathlineto{\pgfqpoint{2.223507in}{1.326940in}}%
\pgfpathlineto{\pgfqpoint{2.226309in}{1.326750in}}%
\pgfpathlineto{\pgfqpoint{2.228782in}{1.326673in}}%
\pgfpathlineto{\pgfqpoint{2.230812in}{1.326385in}}%
\pgfpathlineto{\pgfqpoint{2.231672in}{0.504474in}}%
\pgfpathlineto{\pgfqpoint{2.234804in}{0.504434in}}%
\pgfpathlineto{\pgfqpoint{2.238910in}{0.504832in}}%
\pgfpathlineto{\pgfqpoint{2.239002in}{0.505669in}}%
\pgfpathlineto{\pgfqpoint{2.239176in}{0.902204in}}%
\pgfpathlineto{\pgfqpoint{2.239252in}{1.326967in}}%
\pgfpathlineto{\pgfqpoint{2.239992in}{1.326940in}}%
\pgfpathlineto{\pgfqpoint{2.242794in}{1.326750in}}%
\pgfpathlineto{\pgfqpoint{2.245267in}{1.326673in}}%
\pgfpathlineto{\pgfqpoint{2.247297in}{1.326385in}}%
\pgfpathlineto{\pgfqpoint{2.248157in}{0.504474in}}%
\pgfpathlineto{\pgfqpoint{2.251289in}{0.504434in}}%
\pgfpathlineto{\pgfqpoint{2.255395in}{0.504832in}}%
\pgfpathlineto{\pgfqpoint{2.255487in}{0.505669in}}%
\pgfpathlineto{\pgfqpoint{2.255661in}{0.902204in}}%
\pgfpathlineto{\pgfqpoint{2.255737in}{1.326967in}}%
\pgfpathlineto{\pgfqpoint{2.256477in}{1.326940in}}%
\pgfpathlineto{\pgfqpoint{2.259279in}{1.326750in}}%
\pgfpathlineto{\pgfqpoint{2.261752in}{1.326673in}}%
\pgfpathlineto{\pgfqpoint{2.263782in}{1.326385in}}%
\pgfpathlineto{\pgfqpoint{2.264642in}{0.504474in}}%
\pgfpathlineto{\pgfqpoint{2.267774in}{0.504434in}}%
\pgfpathlineto{\pgfqpoint{2.271880in}{0.504832in}}%
\pgfpathlineto{\pgfqpoint{2.271972in}{0.505669in}}%
\pgfpathlineto{\pgfqpoint{2.272146in}{0.902205in}}%
\pgfpathlineto{\pgfqpoint{2.272222in}{1.326967in}}%
\pgfpathlineto{\pgfqpoint{2.272961in}{1.326940in}}%
\pgfpathlineto{\pgfqpoint{2.275764in}{1.326750in}}%
\pgfpathlineto{\pgfqpoint{2.278237in}{1.326673in}}%
\pgfpathlineto{\pgfqpoint{2.280267in}{1.326385in}}%
\pgfpathlineto{\pgfqpoint{2.281127in}{0.504474in}}%
\pgfpathlineto{\pgfqpoint{2.284259in}{0.504434in}}%
\pgfpathlineto{\pgfqpoint{2.288365in}{0.504832in}}%
\pgfpathlineto{\pgfqpoint{2.288457in}{0.505669in}}%
\pgfpathlineto{\pgfqpoint{2.288631in}{0.902206in}}%
\pgfpathlineto{\pgfqpoint{2.288706in}{1.326967in}}%
\pgfpathlineto{\pgfqpoint{2.289446in}{1.326940in}}%
\pgfpathlineto{\pgfqpoint{2.292249in}{1.326750in}}%
\pgfpathlineto{\pgfqpoint{2.294721in}{1.326673in}}%
\pgfpathlineto{\pgfqpoint{2.296752in}{1.326385in}}%
\pgfpathlineto{\pgfqpoint{2.297612in}{0.504474in}}%
\pgfpathlineto{\pgfqpoint{2.300744in}{0.504434in}}%
\pgfpathlineto{\pgfqpoint{2.304849in}{0.504832in}}%
\pgfpathlineto{\pgfqpoint{2.304942in}{0.505669in}}%
\pgfpathlineto{\pgfqpoint{2.305115in}{0.902207in}}%
\pgfpathlineto{\pgfqpoint{2.305191in}{1.326967in}}%
\pgfpathlineto{\pgfqpoint{2.305931in}{1.326940in}}%
\pgfpathlineto{\pgfqpoint{2.308734in}{1.326750in}}%
\pgfpathlineto{\pgfqpoint{2.311206in}{1.326673in}}%
\pgfpathlineto{\pgfqpoint{2.313237in}{1.326385in}}%
\pgfpathlineto{\pgfqpoint{2.314096in}{0.504474in}}%
\pgfpathlineto{\pgfqpoint{2.317229in}{0.504434in}}%
\pgfpathlineto{\pgfqpoint{2.321334in}{0.504832in}}%
\pgfpathlineto{\pgfqpoint{2.321427in}{0.505669in}}%
\pgfpathlineto{\pgfqpoint{2.321600in}{0.902208in}}%
\pgfpathlineto{\pgfqpoint{2.321676in}{1.326967in}}%
\pgfpathlineto{\pgfqpoint{2.322416in}{1.326940in}}%
\pgfpathlineto{\pgfqpoint{2.325218in}{1.326750in}}%
\pgfpathlineto{\pgfqpoint{2.327691in}{1.326673in}}%
\pgfpathlineto{\pgfqpoint{2.329721in}{1.326385in}}%
\pgfpathlineto{\pgfqpoint{2.330581in}{0.504474in}}%
\pgfpathlineto{\pgfqpoint{2.333713in}{0.504434in}}%
\pgfpathlineto{\pgfqpoint{2.337819in}{0.504832in}}%
\pgfpathlineto{\pgfqpoint{2.337911in}{0.505669in}}%
\pgfpathlineto{\pgfqpoint{2.338085in}{0.902209in}}%
\pgfpathlineto{\pgfqpoint{2.338161in}{1.326967in}}%
\pgfpathlineto{\pgfqpoint{2.338901in}{1.326940in}}%
\pgfpathlineto{\pgfqpoint{2.341703in}{1.326750in}}%
\pgfpathlineto{\pgfqpoint{2.344176in}{1.326673in}}%
\pgfpathlineto{\pgfqpoint{2.346206in}{1.326385in}}%
\pgfpathlineto{\pgfqpoint{2.347066in}{0.504474in}}%
\pgfpathlineto{\pgfqpoint{2.350198in}{0.504434in}}%
\pgfpathlineto{\pgfqpoint{2.354304in}{0.504832in}}%
\pgfpathlineto{\pgfqpoint{2.354396in}{0.505669in}}%
\pgfpathlineto{\pgfqpoint{2.354570in}{0.902210in}}%
\pgfpathlineto{\pgfqpoint{2.354646in}{1.326967in}}%
\pgfpathlineto{\pgfqpoint{2.355386in}{1.326940in}}%
\pgfpathlineto{\pgfqpoint{2.358188in}{1.326750in}}%
\pgfpathlineto{\pgfqpoint{2.360661in}{1.326673in}}%
\pgfpathlineto{\pgfqpoint{2.362691in}{1.326385in}}%
\pgfpathlineto{\pgfqpoint{2.363551in}{0.504474in}}%
\pgfpathlineto{\pgfqpoint{2.366683in}{0.504434in}}%
\pgfpathlineto{\pgfqpoint{2.370789in}{0.504832in}}%
\pgfpathlineto{\pgfqpoint{2.370881in}{0.505669in}}%
\pgfpathlineto{\pgfqpoint{2.371055in}{0.902210in}}%
\pgfpathlineto{\pgfqpoint{2.371131in}{1.326967in}}%
\pgfpathlineto{\pgfqpoint{2.371871in}{1.326940in}}%
\pgfpathlineto{\pgfqpoint{2.374673in}{1.326750in}}%
\pgfpathlineto{\pgfqpoint{2.377146in}{1.326673in}}%
\pgfpathlineto{\pgfqpoint{2.379176in}{1.326385in}}%
\pgfpathlineto{\pgfqpoint{2.380036in}{0.504474in}}%
\pgfpathlineto{\pgfqpoint{2.383168in}{0.504434in}}%
\pgfpathlineto{\pgfqpoint{2.387274in}{0.504832in}}%
\pgfpathlineto{\pgfqpoint{2.387366in}{0.505669in}}%
\pgfpathlineto{\pgfqpoint{2.387540in}{0.902210in}}%
\pgfpathlineto{\pgfqpoint{2.387616in}{1.326967in}}%
\pgfpathlineto{\pgfqpoint{2.388355in}{1.326940in}}%
\pgfpathlineto{\pgfqpoint{2.391158in}{1.326750in}}%
\pgfpathlineto{\pgfqpoint{2.393631in}{1.326673in}}%
\pgfpathlineto{\pgfqpoint{2.395661in}{1.326385in}}%
\pgfpathlineto{\pgfqpoint{2.396521in}{0.504474in}}%
\pgfpathlineto{\pgfqpoint{2.399653in}{0.504434in}}%
\pgfpathlineto{\pgfqpoint{2.403759in}{0.504832in}}%
\pgfpathlineto{\pgfqpoint{2.403851in}{0.505669in}}%
\pgfpathlineto{\pgfqpoint{2.404025in}{0.902209in}}%
\pgfpathlineto{\pgfqpoint{2.404100in}{1.326967in}}%
\pgfpathlineto{\pgfqpoint{2.404840in}{1.326940in}}%
\pgfpathlineto{\pgfqpoint{2.407643in}{1.326750in}}%
\pgfpathlineto{\pgfqpoint{2.410115in}{1.326673in}}%
\pgfpathlineto{\pgfqpoint{2.412146in}{1.326385in}}%
\pgfpathlineto{\pgfqpoint{2.413006in}{0.504474in}}%
\pgfpathlineto{\pgfqpoint{2.416138in}{0.504434in}}%
\pgfpathlineto{\pgfqpoint{2.420243in}{0.504832in}}%
\pgfpathlineto{\pgfqpoint{2.420336in}{0.505669in}}%
\pgfpathlineto{\pgfqpoint{2.420509in}{0.902209in}}%
\pgfpathlineto{\pgfqpoint{2.420585in}{1.326967in}}%
\pgfpathlineto{\pgfqpoint{2.421325in}{1.326940in}}%
\pgfpathlineto{\pgfqpoint{2.424128in}{1.326750in}}%
\pgfpathlineto{\pgfqpoint{2.426600in}{1.326673in}}%
\pgfpathlineto{\pgfqpoint{2.428631in}{1.326385in}}%
\pgfpathlineto{\pgfqpoint{2.429490in}{0.504474in}}%
\pgfpathlineto{\pgfqpoint{2.432623in}{0.504434in}}%
\pgfpathlineto{\pgfqpoint{2.436728in}{0.504832in}}%
\pgfpathlineto{\pgfqpoint{2.436821in}{0.505669in}}%
\pgfpathlineto{\pgfqpoint{2.436994in}{0.902208in}}%
\pgfpathlineto{\pgfqpoint{2.437070in}{1.326967in}}%
\pgfpathlineto{\pgfqpoint{2.437810in}{1.326940in}}%
\pgfpathlineto{\pgfqpoint{2.440612in}{1.326750in}}%
\pgfpathlineto{\pgfqpoint{2.443085in}{1.326673in}}%
\pgfpathlineto{\pgfqpoint{2.445115in}{1.326385in}}%
\pgfpathlineto{\pgfqpoint{2.445975in}{0.504474in}}%
\pgfpathlineto{\pgfqpoint{2.449107in}{0.504434in}}%
\pgfpathlineto{\pgfqpoint{2.453213in}{0.504832in}}%
\pgfpathlineto{\pgfqpoint{2.453305in}{0.505669in}}%
\pgfpathlineto{\pgfqpoint{2.453479in}{0.902208in}}%
\pgfpathlineto{\pgfqpoint{2.453555in}{1.326967in}}%
\pgfpathlineto{\pgfqpoint{2.454295in}{1.326940in}}%
\pgfpathlineto{\pgfqpoint{2.457097in}{1.326750in}}%
\pgfpathlineto{\pgfqpoint{2.459570in}{1.326673in}}%
\pgfpathlineto{\pgfqpoint{2.461600in}{1.326385in}}%
\pgfpathlineto{\pgfqpoint{2.462460in}{0.504474in}}%
\pgfpathlineto{\pgfqpoint{2.465592in}{0.504434in}}%
\pgfpathlineto{\pgfqpoint{2.469698in}{0.504832in}}%
\pgfpathlineto{\pgfqpoint{2.469790in}{0.505669in}}%
\pgfpathlineto{\pgfqpoint{2.469964in}{0.902208in}}%
\pgfpathlineto{\pgfqpoint{2.470040in}{1.326967in}}%
\pgfpathlineto{\pgfqpoint{2.470780in}{1.326940in}}%
\pgfpathlineto{\pgfqpoint{2.473582in}{1.326750in}}%
\pgfpathlineto{\pgfqpoint{2.476055in}{1.326673in}}%
\pgfpathlineto{\pgfqpoint{2.478085in}{1.326385in}}%
\pgfpathlineto{\pgfqpoint{2.478945in}{0.504474in}}%
\pgfpathlineto{\pgfqpoint{2.482077in}{0.504434in}}%
\pgfpathlineto{\pgfqpoint{2.486183in}{0.504832in}}%
\pgfpathlineto{\pgfqpoint{2.486275in}{0.505669in}}%
\pgfpathlineto{\pgfqpoint{2.486449in}{0.902207in}}%
\pgfpathlineto{\pgfqpoint{2.486525in}{1.326967in}}%
\pgfpathlineto{\pgfqpoint{2.487264in}{1.326940in}}%
\pgfpathlineto{\pgfqpoint{2.490067in}{1.326750in}}%
\pgfpathlineto{\pgfqpoint{2.492540in}{1.326673in}}%
\pgfpathlineto{\pgfqpoint{2.494546in}{1.326521in}}%
\pgfpathlineto{\pgfqpoint{2.495431in}{0.504463in}}%
\pgfpathlineto{\pgfqpoint{2.500377in}{0.504491in}}%
\pgfpathlineto{\pgfqpoint{2.502668in}{0.504833in}}%
\pgfpathlineto{\pgfqpoint{2.502760in}{0.505670in}}%
\pgfpathlineto{\pgfqpoint{2.502934in}{0.902228in}}%
\pgfpathlineto{\pgfqpoint{2.503009in}{1.326967in}}%
\pgfpathlineto{\pgfqpoint{2.503749in}{1.326941in}}%
\pgfpathlineto{\pgfqpoint{2.506552in}{1.326751in}}%
\pgfpathlineto{\pgfqpoint{2.509024in}{1.326674in}}%
\pgfpathlineto{\pgfqpoint{2.511055in}{1.326386in}}%
\pgfpathlineto{\pgfqpoint{2.511914in}{0.504475in}}%
\pgfpathlineto{\pgfqpoint{2.515047in}{0.504434in}}%
\pgfpathlineto{\pgfqpoint{2.519152in}{0.504833in}}%
\pgfpathlineto{\pgfqpoint{2.519245in}{0.505670in}}%
\pgfpathlineto{\pgfqpoint{2.519419in}{0.902227in}}%
\pgfpathlineto{\pgfqpoint{2.519494in}{1.326967in}}%
\pgfpathlineto{\pgfqpoint{2.520234in}{1.326941in}}%
\pgfpathlineto{\pgfqpoint{2.523037in}{1.326750in}}%
\pgfpathlineto{\pgfqpoint{2.525509in}{1.326674in}}%
\pgfpathlineto{\pgfqpoint{2.527540in}{1.326386in}}%
\pgfpathlineto{\pgfqpoint{2.528399in}{0.504474in}}%
\pgfpathlineto{\pgfqpoint{2.531531in}{0.504434in}}%
\pgfpathlineto{\pgfqpoint{2.535637in}{0.504833in}}%
\pgfpathlineto{\pgfqpoint{2.535730in}{0.505670in}}%
\pgfpathlineto{\pgfqpoint{2.535903in}{0.902222in}}%
\pgfpathlineto{\pgfqpoint{2.535979in}{1.326967in}}%
\pgfpathlineto{\pgfqpoint{2.536719in}{1.326940in}}%
\pgfpathlineto{\pgfqpoint{2.539521in}{1.326750in}}%
\pgfpathlineto{\pgfqpoint{2.541994in}{1.326673in}}%
\pgfpathlineto{\pgfqpoint{2.544025in}{1.326386in}}%
\pgfpathlineto{\pgfqpoint{2.544884in}{0.504474in}}%
\pgfpathlineto{\pgfqpoint{2.548016in}{0.504434in}}%
\pgfpathlineto{\pgfqpoint{2.552122in}{0.504832in}}%
\pgfpathlineto{\pgfqpoint{2.552214in}{0.505669in}}%
\pgfpathlineto{\pgfqpoint{2.552388in}{0.902215in}}%
\pgfpathlineto{\pgfqpoint{2.552464in}{1.326967in}}%
\pgfpathlineto{\pgfqpoint{2.553204in}{1.326940in}}%
\pgfpathlineto{\pgfqpoint{2.556006in}{1.326750in}}%
\pgfpathlineto{\pgfqpoint{2.558479in}{1.326673in}}%
\pgfpathlineto{\pgfqpoint{2.560509in}{1.326386in}}%
\pgfpathlineto{\pgfqpoint{2.561369in}{0.504474in}}%
\pgfpathlineto{\pgfqpoint{2.564501in}{0.504434in}}%
\pgfpathlineto{\pgfqpoint{2.568607in}{0.504832in}}%
\pgfpathlineto{\pgfqpoint{2.568699in}{0.505669in}}%
\pgfpathlineto{\pgfqpoint{2.568873in}{0.902209in}}%
\pgfpathlineto{\pgfqpoint{2.568949in}{1.326967in}}%
\pgfpathlineto{\pgfqpoint{2.569689in}{1.326940in}}%
\pgfpathlineto{\pgfqpoint{2.572491in}{1.326750in}}%
\pgfpathlineto{\pgfqpoint{2.574964in}{1.326673in}}%
\pgfpathlineto{\pgfqpoint{2.576994in}{1.326385in}}%
\pgfpathlineto{\pgfqpoint{2.577854in}{0.504474in}}%
\pgfpathlineto{\pgfqpoint{2.580986in}{0.504434in}}%
\pgfpathlineto{\pgfqpoint{2.585092in}{0.504832in}}%
\pgfpathlineto{\pgfqpoint{2.585184in}{0.505669in}}%
\pgfpathlineto{\pgfqpoint{2.585358in}{0.902204in}}%
\pgfpathlineto{\pgfqpoint{2.585434in}{1.326967in}}%
\pgfpathlineto{\pgfqpoint{2.586174in}{1.326940in}}%
\pgfpathlineto{\pgfqpoint{2.588976in}{1.326750in}}%
\pgfpathlineto{\pgfqpoint{2.591449in}{1.326673in}}%
\pgfpathlineto{\pgfqpoint{2.593479in}{1.326385in}}%
\pgfpathlineto{\pgfqpoint{2.594339in}{0.504474in}}%
\pgfpathlineto{\pgfqpoint{2.597471in}{0.504434in}}%
\pgfpathlineto{\pgfqpoint{2.601577in}{0.504832in}}%
\pgfpathlineto{\pgfqpoint{2.601669in}{0.505669in}}%
\pgfpathlineto{\pgfqpoint{2.601843in}{0.902200in}}%
\pgfpathlineto{\pgfqpoint{2.601919in}{1.326966in}}%
\pgfpathlineto{\pgfqpoint{2.602658in}{1.326940in}}%
\pgfpathlineto{\pgfqpoint{2.605461in}{1.326749in}}%
\pgfpathlineto{\pgfqpoint{2.607934in}{1.326673in}}%
\pgfpathlineto{\pgfqpoint{2.609964in}{1.326385in}}%
\pgfpathlineto{\pgfqpoint{2.610824in}{0.504474in}}%
\pgfpathlineto{\pgfqpoint{2.613956in}{0.504434in}}%
\pgfpathlineto{\pgfqpoint{2.618062in}{0.504832in}}%
\pgfpathlineto{\pgfqpoint{2.618154in}{0.505668in}}%
\pgfpathlineto{\pgfqpoint{2.618328in}{0.902198in}}%
\pgfpathlineto{\pgfqpoint{2.618403in}{1.326966in}}%
\pgfpathlineto{\pgfqpoint{2.619143in}{1.326940in}}%
\pgfpathlineto{\pgfqpoint{2.621946in}{1.326749in}}%
\pgfpathlineto{\pgfqpoint{2.624418in}{1.326673in}}%
\pgfpathlineto{\pgfqpoint{2.626449in}{1.326385in}}%
\pgfpathlineto{\pgfqpoint{2.627309in}{0.504474in}}%
\pgfpathlineto{\pgfqpoint{2.630441in}{0.504434in}}%
\pgfpathlineto{\pgfqpoint{2.634546in}{0.504832in}}%
\pgfpathlineto{\pgfqpoint{2.634639in}{0.505668in}}%
\pgfpathlineto{\pgfqpoint{2.634812in}{0.902198in}}%
\pgfpathlineto{\pgfqpoint{2.634888in}{1.326966in}}%
\pgfpathlineto{\pgfqpoint{2.635628in}{1.326940in}}%
\pgfpathlineto{\pgfqpoint{2.638431in}{1.326749in}}%
\pgfpathlineto{\pgfqpoint{2.640903in}{1.326673in}}%
\pgfpathlineto{\pgfqpoint{2.642934in}{1.326385in}}%
\pgfpathlineto{\pgfqpoint{2.643794in}{0.504474in}}%
\pgfpathlineto{\pgfqpoint{2.646926in}{0.504434in}}%
\pgfpathlineto{\pgfqpoint{2.651031in}{0.504832in}}%
\pgfpathlineto{\pgfqpoint{2.651124in}{0.505668in}}%
\pgfpathlineto{\pgfqpoint{2.651297in}{0.902199in}}%
\pgfpathlineto{\pgfqpoint{2.651373in}{1.326966in}}%
\pgfpathlineto{\pgfqpoint{2.652113in}{1.326940in}}%
\pgfpathlineto{\pgfqpoint{2.654915in}{1.326749in}}%
\pgfpathlineto{\pgfqpoint{2.657388in}{1.326673in}}%
\pgfpathlineto{\pgfqpoint{2.659418in}{1.326385in}}%
\pgfpathlineto{\pgfqpoint{2.660278in}{0.504474in}}%
\pgfpathlineto{\pgfqpoint{2.663411in}{0.504434in}}%
\pgfpathlineto{\pgfqpoint{2.667516in}{0.504832in}}%
\pgfpathlineto{\pgfqpoint{2.667608in}{0.505669in}}%
\pgfpathlineto{\pgfqpoint{2.667782in}{0.902202in}}%
\pgfpathlineto{\pgfqpoint{2.667858in}{1.326966in}}%
\pgfpathlineto{\pgfqpoint{2.668598in}{1.326940in}}%
\pgfpathlineto{\pgfqpoint{2.671400in}{1.326750in}}%
\pgfpathlineto{\pgfqpoint{2.673873in}{1.326673in}}%
\pgfpathlineto{\pgfqpoint{2.675903in}{1.326385in}}%
\pgfpathlineto{\pgfqpoint{2.676763in}{0.504474in}}%
\pgfpathlineto{\pgfqpoint{2.679895in}{0.504434in}}%
\pgfpathlineto{\pgfqpoint{2.684001in}{0.504832in}}%
\pgfpathlineto{\pgfqpoint{2.684093in}{0.505669in}}%
\pgfpathlineto{\pgfqpoint{2.684267in}{0.902205in}}%
\pgfpathlineto{\pgfqpoint{2.684343in}{1.326967in}}%
\pgfpathlineto{\pgfqpoint{2.685083in}{1.326940in}}%
\pgfpathlineto{\pgfqpoint{2.687885in}{1.326750in}}%
\pgfpathlineto{\pgfqpoint{2.690358in}{1.326673in}}%
\pgfpathlineto{\pgfqpoint{2.692388in}{1.326385in}}%
\pgfpathlineto{\pgfqpoint{2.693248in}{0.504474in}}%
\pgfpathlineto{\pgfqpoint{2.696380in}{0.504434in}}%
\pgfpathlineto{\pgfqpoint{2.700486in}{0.504832in}}%
\pgfpathlineto{\pgfqpoint{2.700578in}{0.505669in}}%
\pgfpathlineto{\pgfqpoint{2.700752in}{0.902207in}}%
\pgfpathlineto{\pgfqpoint{2.700828in}{1.326967in}}%
\pgfpathlineto{\pgfqpoint{2.701568in}{1.326940in}}%
\pgfpathlineto{\pgfqpoint{2.704370in}{1.326750in}}%
\pgfpathlineto{\pgfqpoint{2.706843in}{1.326673in}}%
\pgfpathlineto{\pgfqpoint{2.708873in}{1.326385in}}%
\pgfpathlineto{\pgfqpoint{2.709733in}{0.504474in}}%
\pgfpathlineto{\pgfqpoint{2.712865in}{0.504434in}}%
\pgfpathlineto{\pgfqpoint{2.716971in}{0.504832in}}%
\pgfpathlineto{\pgfqpoint{2.717063in}{0.505669in}}%
\pgfpathlineto{\pgfqpoint{2.717237in}{0.902210in}}%
\pgfpathlineto{\pgfqpoint{2.717313in}{1.326967in}}%
\pgfpathlineto{\pgfqpoint{2.718052in}{1.326940in}}%
\pgfpathlineto{\pgfqpoint{2.720855in}{1.326750in}}%
\pgfpathlineto{\pgfqpoint{2.723328in}{1.326673in}}%
\pgfpathlineto{\pgfqpoint{2.725358in}{1.326385in}}%
\pgfpathlineto{\pgfqpoint{2.726218in}{0.504474in}}%
\pgfpathlineto{\pgfqpoint{2.729350in}{0.504434in}}%
\pgfpathlineto{\pgfqpoint{2.733455in}{0.504832in}}%
\pgfpathlineto{\pgfqpoint{2.733548in}{0.505669in}}%
\pgfpathlineto{\pgfqpoint{2.733722in}{0.902211in}}%
\pgfpathlineto{\pgfqpoint{2.733797in}{1.326967in}}%
\pgfpathlineto{\pgfqpoint{2.734537in}{1.326940in}}%
\pgfpathlineto{\pgfqpoint{2.737340in}{1.326750in}}%
\pgfpathlineto{\pgfqpoint{2.739812in}{1.326673in}}%
\pgfpathlineto{\pgfqpoint{2.741843in}{1.326386in}}%
\pgfpathlineto{\pgfqpoint{2.742702in}{0.504474in}}%
\pgfpathlineto{\pgfqpoint{2.745835in}{0.504434in}}%
\pgfpathlineto{\pgfqpoint{2.749940in}{0.504832in}}%
\pgfpathlineto{\pgfqpoint{2.750033in}{0.505669in}}%
\pgfpathlineto{\pgfqpoint{2.750206in}{0.902212in}}%
\pgfpathlineto{\pgfqpoint{2.750282in}{1.326967in}}%
\pgfpathlineto{\pgfqpoint{2.751022in}{1.326940in}}%
\pgfpathlineto{\pgfqpoint{2.753824in}{1.326750in}}%
\pgfpathlineto{\pgfqpoint{2.756297in}{1.326673in}}%
\pgfpathlineto{\pgfqpoint{2.758303in}{1.326521in}}%
\pgfpathlineto{\pgfqpoint{2.759189in}{0.504463in}}%
\pgfpathlineto{\pgfqpoint{2.764135in}{0.504491in}}%
\pgfpathlineto{\pgfqpoint{2.766425in}{0.504833in}}%
\pgfpathlineto{\pgfqpoint{2.766517in}{0.505671in}}%
\pgfpathlineto{\pgfqpoint{2.766691in}{0.902233in}}%
\pgfpathlineto{\pgfqpoint{2.766767in}{1.326967in}}%
\pgfpathlineto{\pgfqpoint{2.767507in}{1.326941in}}%
\pgfpathlineto{\pgfqpoint{2.770309in}{1.326751in}}%
\pgfpathlineto{\pgfqpoint{2.772782in}{1.326674in}}%
\pgfpathlineto{\pgfqpoint{2.774812in}{1.326387in}}%
\pgfpathlineto{\pgfqpoint{2.775672in}{0.504475in}}%
\pgfpathlineto{\pgfqpoint{2.778804in}{0.504435in}}%
\pgfpathlineto{\pgfqpoint{2.782910in}{0.504833in}}%
\pgfpathlineto{\pgfqpoint{2.783002in}{0.505670in}}%
\pgfpathlineto{\pgfqpoint{2.783176in}{0.902232in}}%
\pgfpathlineto{\pgfqpoint{2.783252in}{1.326967in}}%
\pgfpathlineto{\pgfqpoint{2.783992in}{1.326941in}}%
\pgfpathlineto{\pgfqpoint{2.786794in}{1.326751in}}%
\pgfpathlineto{\pgfqpoint{2.789267in}{1.326674in}}%
\pgfpathlineto{\pgfqpoint{2.791297in}{1.326386in}}%
\pgfpathlineto{\pgfqpoint{2.792157in}{0.504475in}}%
\pgfpathlineto{\pgfqpoint{2.795289in}{0.504434in}}%
\pgfpathlineto{\pgfqpoint{2.799395in}{0.504833in}}%
\pgfpathlineto{\pgfqpoint{2.799487in}{0.505670in}}%
\pgfpathlineto{\pgfqpoint{2.799661in}{0.902225in}}%
\pgfpathlineto{\pgfqpoint{2.799737in}{1.326967in}}%
\pgfpathlineto{\pgfqpoint{2.800477in}{1.326941in}}%
\pgfpathlineto{\pgfqpoint{2.803279in}{1.326750in}}%
\pgfpathlineto{\pgfqpoint{2.805752in}{1.326673in}}%
\pgfpathlineto{\pgfqpoint{2.807782in}{1.326386in}}%
\pgfpathlineto{\pgfqpoint{2.808642in}{0.504474in}}%
\pgfpathlineto{\pgfqpoint{2.811774in}{0.504434in}}%
\pgfpathlineto{\pgfqpoint{2.815880in}{0.504833in}}%
\pgfpathlineto{\pgfqpoint{2.815972in}{0.505670in}}%
\pgfpathlineto{\pgfqpoint{2.816146in}{0.902218in}}%
\pgfpathlineto{\pgfqpoint{2.816222in}{1.326967in}}%
\pgfpathlineto{\pgfqpoint{2.816961in}{1.326940in}}%
\pgfpathlineto{\pgfqpoint{2.819764in}{1.326750in}}%
\pgfpathlineto{\pgfqpoint{2.822237in}{1.326673in}}%
\pgfpathlineto{\pgfqpoint{2.824267in}{1.326386in}}%
\pgfpathlineto{\pgfqpoint{2.825127in}{0.504474in}}%
\pgfpathlineto{\pgfqpoint{2.828259in}{0.504434in}}%
\pgfpathlineto{\pgfqpoint{2.832365in}{0.504832in}}%
\pgfpathlineto{\pgfqpoint{2.832457in}{0.505669in}}%
\pgfpathlineto{\pgfqpoint{2.832631in}{0.902210in}}%
\pgfpathlineto{\pgfqpoint{2.832706in}{1.326967in}}%
\pgfpathlineto{\pgfqpoint{2.833446in}{1.326940in}}%
\pgfpathlineto{\pgfqpoint{2.836249in}{1.326750in}}%
\pgfpathlineto{\pgfqpoint{2.838721in}{1.326673in}}%
\pgfpathlineto{\pgfqpoint{2.840752in}{1.326385in}}%
\pgfpathlineto{\pgfqpoint{2.841612in}{0.504474in}}%
\pgfpathlineto{\pgfqpoint{2.844744in}{0.504434in}}%
\pgfpathlineto{\pgfqpoint{2.848849in}{0.504832in}}%
\pgfpathlineto{\pgfqpoint{2.848942in}{0.505669in}}%
\pgfpathlineto{\pgfqpoint{2.849115in}{0.902203in}}%
\pgfpathlineto{\pgfqpoint{2.849191in}{1.326966in}}%
\pgfpathlineto{\pgfqpoint{2.849931in}{1.326940in}}%
\pgfpathlineto{\pgfqpoint{2.852734in}{1.326750in}}%
\pgfpathlineto{\pgfqpoint{2.855206in}{1.326673in}}%
\pgfpathlineto{\pgfqpoint{2.857237in}{1.326385in}}%
\pgfpathlineto{\pgfqpoint{2.858097in}{0.504474in}}%
\pgfpathlineto{\pgfqpoint{2.861229in}{0.504434in}}%
\pgfpathlineto{\pgfqpoint{2.865334in}{0.504832in}}%
\pgfpathlineto{\pgfqpoint{2.865427in}{0.505668in}}%
\pgfpathlineto{\pgfqpoint{2.865600in}{0.902198in}}%
\pgfpathlineto{\pgfqpoint{2.865676in}{1.326966in}}%
\pgfpathlineto{\pgfqpoint{2.866416in}{1.326940in}}%
\pgfpathlineto{\pgfqpoint{2.869218in}{1.326749in}}%
\pgfpathlineto{\pgfqpoint{2.871691in}{1.326672in}}%
\pgfpathlineto{\pgfqpoint{2.873721in}{1.326385in}}%
\pgfpathlineto{\pgfqpoint{2.874581in}{0.504474in}}%
\pgfpathlineto{\pgfqpoint{2.877714in}{0.504434in}}%
\pgfpathlineto{\pgfqpoint{2.881819in}{0.504832in}}%
\pgfpathlineto{\pgfqpoint{2.881911in}{0.505668in}}%
\pgfpathlineto{\pgfqpoint{2.882085in}{0.902196in}}%
\pgfpathlineto{\pgfqpoint{2.882161in}{1.326966in}}%
\pgfpathlineto{\pgfqpoint{2.882901in}{1.326940in}}%
\pgfpathlineto{\pgfqpoint{2.885703in}{1.326749in}}%
\pgfpathlineto{\pgfqpoint{2.888176in}{1.326672in}}%
\pgfpathlineto{\pgfqpoint{2.890206in}{1.326385in}}%
\pgfpathlineto{\pgfqpoint{2.891066in}{0.504474in}}%
\pgfpathlineto{\pgfqpoint{2.894198in}{0.504434in}}%
\pgfpathlineto{\pgfqpoint{2.898304in}{0.504832in}}%
\pgfpathlineto{\pgfqpoint{2.898396in}{0.505668in}}%
\pgfpathlineto{\pgfqpoint{2.898570in}{0.902196in}}%
\pgfpathlineto{\pgfqpoint{2.898646in}{1.326966in}}%
\pgfpathlineto{\pgfqpoint{2.899386in}{1.326940in}}%
\pgfpathlineto{\pgfqpoint{2.902188in}{1.326749in}}%
\pgfpathlineto{\pgfqpoint{2.904661in}{1.326672in}}%
\pgfpathlineto{\pgfqpoint{2.906691in}{1.326385in}}%
\pgfpathlineto{\pgfqpoint{2.907551in}{0.504474in}}%
\pgfpathlineto{\pgfqpoint{2.910683in}{0.504434in}}%
\pgfpathlineto{\pgfqpoint{2.914789in}{0.504832in}}%
\pgfpathlineto{\pgfqpoint{2.914881in}{0.505668in}}%
\pgfpathlineto{\pgfqpoint{2.915055in}{0.902197in}}%
\pgfpathlineto{\pgfqpoint{2.915131in}{1.326966in}}%
\pgfpathlineto{\pgfqpoint{2.915871in}{1.326940in}}%
\pgfpathlineto{\pgfqpoint{2.918673in}{1.326749in}}%
\pgfpathlineto{\pgfqpoint{2.921146in}{1.326673in}}%
\pgfpathlineto{\pgfqpoint{2.923176in}{1.326385in}}%
\pgfpathlineto{\pgfqpoint{2.924036in}{0.504474in}}%
\pgfpathlineto{\pgfqpoint{2.927168in}{0.504434in}}%
\pgfpathlineto{\pgfqpoint{2.931274in}{0.504832in}}%
\pgfpathlineto{\pgfqpoint{2.931366in}{0.505669in}}%
\pgfpathlineto{\pgfqpoint{2.931540in}{0.902200in}}%
\pgfpathlineto{\pgfqpoint{2.931616in}{1.326966in}}%
\pgfpathlineto{\pgfqpoint{2.932355in}{1.326940in}}%
\pgfpathlineto{\pgfqpoint{2.935158in}{1.326749in}}%
\pgfpathlineto{\pgfqpoint{2.937631in}{1.326673in}}%
\pgfpathlineto{\pgfqpoint{2.939661in}{1.326385in}}%
\pgfpathlineto{\pgfqpoint{2.940521in}{0.504474in}}%
\pgfpathlineto{\pgfqpoint{2.943653in}{0.504434in}}%
\pgfpathlineto{\pgfqpoint{2.947759in}{0.504832in}}%
\pgfpathlineto{\pgfqpoint{2.947851in}{0.505669in}}%
\pgfpathlineto{\pgfqpoint{2.948025in}{0.902204in}}%
\pgfpathlineto{\pgfqpoint{2.948100in}{1.326967in}}%
\pgfpathlineto{\pgfqpoint{2.948840in}{1.326940in}}%
\pgfpathlineto{\pgfqpoint{2.951643in}{1.326750in}}%
\pgfpathlineto{\pgfqpoint{2.954115in}{1.326673in}}%
\pgfpathlineto{\pgfqpoint{2.956146in}{1.326385in}}%
\pgfpathlineto{\pgfqpoint{2.957006in}{0.504474in}}%
\pgfpathlineto{\pgfqpoint{2.960138in}{0.504434in}}%
\pgfpathlineto{\pgfqpoint{2.964243in}{0.504832in}}%
\pgfpathlineto{\pgfqpoint{2.964336in}{0.505669in}}%
\pgfpathlineto{\pgfqpoint{2.964509in}{0.902207in}}%
\pgfpathlineto{\pgfqpoint{2.964585in}{1.326967in}}%
\pgfpathlineto{\pgfqpoint{2.965325in}{1.326940in}}%
\pgfpathlineto{\pgfqpoint{2.968128in}{1.326750in}}%
\pgfpathlineto{\pgfqpoint{2.970600in}{1.326673in}}%
\pgfpathlineto{\pgfqpoint{2.972631in}{1.326385in}}%
\pgfpathlineto{\pgfqpoint{2.973490in}{0.504474in}}%
\pgfpathlineto{\pgfqpoint{2.976623in}{0.504434in}}%
\pgfpathlineto{\pgfqpoint{2.980728in}{0.504832in}}%
\pgfpathlineto{\pgfqpoint{2.980821in}{0.505669in}}%
\pgfpathlineto{\pgfqpoint{2.980994in}{0.902210in}}%
\pgfpathlineto{\pgfqpoint{2.981070in}{1.326967in}}%
\pgfpathlineto{\pgfqpoint{2.981810in}{1.326940in}}%
\pgfpathlineto{\pgfqpoint{2.984612in}{1.326750in}}%
\pgfpathlineto{\pgfqpoint{2.987085in}{1.326673in}}%
\pgfpathlineto{\pgfqpoint{2.989115in}{1.326385in}}%
\pgfpathlineto{\pgfqpoint{2.989975in}{0.504474in}}%
\pgfpathlineto{\pgfqpoint{2.993107in}{0.504434in}}%
\pgfpathlineto{\pgfqpoint{2.997213in}{0.504832in}}%
\pgfpathlineto{\pgfqpoint{2.997305in}{0.505669in}}%
\pgfpathlineto{\pgfqpoint{2.997479in}{0.902212in}}%
\pgfpathlineto{\pgfqpoint{2.997555in}{1.326967in}}%
\pgfpathlineto{\pgfqpoint{2.998295in}{1.326940in}}%
\pgfpathlineto{\pgfqpoint{3.001097in}{1.326750in}}%
\pgfpathlineto{\pgfqpoint{3.003570in}{1.326673in}}%
\pgfpathlineto{\pgfqpoint{3.005600in}{1.326386in}}%
\pgfpathlineto{\pgfqpoint{3.006460in}{0.504474in}}%
\pgfpathlineto{\pgfqpoint{3.009592in}{0.504434in}}%
\pgfpathlineto{\pgfqpoint{3.013698in}{0.504832in}}%
\pgfpathlineto{\pgfqpoint{3.013790in}{0.505669in}}%
\pgfpathlineto{\pgfqpoint{3.013964in}{0.902213in}}%
\pgfpathlineto{\pgfqpoint{3.014040in}{1.326967in}}%
\pgfpathlineto{\pgfqpoint{3.014780in}{1.326940in}}%
\pgfpathlineto{\pgfqpoint{3.017582in}{1.326750in}}%
\pgfpathlineto{\pgfqpoint{3.020055in}{1.326673in}}%
\pgfpathlineto{\pgfqpoint{3.022085in}{1.326386in}}%
\pgfpathlineto{\pgfqpoint{3.022945in}{0.504474in}}%
\pgfpathlineto{\pgfqpoint{3.026077in}{0.504434in}}%
\pgfpathlineto{\pgfqpoint{3.030183in}{0.504832in}}%
\pgfpathlineto{\pgfqpoint{3.030275in}{0.505669in}}%
\pgfpathlineto{\pgfqpoint{3.030449in}{0.902213in}}%
\pgfpathlineto{\pgfqpoint{3.030525in}{1.326967in}}%
\pgfpathlineto{\pgfqpoint{3.031264in}{1.326940in}}%
\pgfpathlineto{\pgfqpoint{3.034067in}{1.326750in}}%
\pgfpathlineto{\pgfqpoint{3.036540in}{1.326673in}}%
\pgfpathlineto{\pgfqpoint{3.038570in}{1.326386in}}%
\pgfpathlineto{\pgfqpoint{3.039430in}{0.504474in}}%
\pgfpathlineto{\pgfqpoint{3.042562in}{0.504434in}}%
\pgfpathlineto{\pgfqpoint{3.046668in}{0.504832in}}%
\pgfpathlineto{\pgfqpoint{3.046760in}{0.505669in}}%
\pgfpathlineto{\pgfqpoint{3.046934in}{0.902213in}}%
\pgfpathlineto{\pgfqpoint{3.047009in}{1.326967in}}%
\pgfpathlineto{\pgfqpoint{3.047749in}{1.326940in}}%
\pgfpathlineto{\pgfqpoint{3.050552in}{1.326750in}}%
\pgfpathlineto{\pgfqpoint{3.053024in}{1.326673in}}%
\pgfpathlineto{\pgfqpoint{3.055055in}{1.326386in}}%
\pgfpathlineto{\pgfqpoint{3.055915in}{0.504474in}}%
\pgfpathlineto{\pgfqpoint{3.059047in}{0.504434in}}%
\pgfpathlineto{\pgfqpoint{3.063152in}{0.504832in}}%
\pgfpathlineto{\pgfqpoint{3.063245in}{0.505669in}}%
\pgfpathlineto{\pgfqpoint{3.063419in}{0.902212in}}%
\pgfpathlineto{\pgfqpoint{3.063494in}{1.326967in}}%
\pgfpathlineto{\pgfqpoint{3.064234in}{1.326940in}}%
\pgfpathlineto{\pgfqpoint{3.067037in}{1.326750in}}%
\pgfpathlineto{\pgfqpoint{3.069509in}{1.326673in}}%
\pgfpathlineto{\pgfqpoint{3.071540in}{1.326385in}}%
\pgfpathlineto{\pgfqpoint{3.072399in}{0.504474in}}%
\pgfpathlineto{\pgfqpoint{3.075532in}{0.504434in}}%
\pgfpathlineto{\pgfqpoint{3.079637in}{0.504832in}}%
\pgfpathlineto{\pgfqpoint{3.079730in}{0.505669in}}%
\pgfpathlineto{\pgfqpoint{3.079903in}{0.902210in}}%
\pgfpathlineto{\pgfqpoint{3.079979in}{1.326967in}}%
\pgfpathlineto{\pgfqpoint{3.080719in}{1.326940in}}%
\pgfpathlineto{\pgfqpoint{3.083521in}{1.326750in}}%
\pgfpathlineto{\pgfqpoint{3.085994in}{1.326673in}}%
\pgfpathlineto{\pgfqpoint{3.088025in}{1.326385in}}%
\pgfpathlineto{\pgfqpoint{3.088884in}{0.504474in}}%
\pgfpathlineto{\pgfqpoint{3.092016in}{0.504434in}}%
\pgfpathlineto{\pgfqpoint{3.096122in}{0.504832in}}%
\pgfpathlineto{\pgfqpoint{3.096214in}{0.505669in}}%
\pgfpathlineto{\pgfqpoint{3.096388in}{0.902209in}}%
\pgfpathlineto{\pgfqpoint{3.096464in}{1.326967in}}%
\pgfpathlineto{\pgfqpoint{3.097204in}{1.326940in}}%
\pgfpathlineto{\pgfqpoint{3.100006in}{1.326750in}}%
\pgfpathlineto{\pgfqpoint{3.102479in}{1.326673in}}%
\pgfpathlineto{\pgfqpoint{3.104509in}{1.326385in}}%
\pgfpathlineto{\pgfqpoint{3.105369in}{0.504474in}}%
\pgfpathlineto{\pgfqpoint{3.108501in}{0.504434in}}%
\pgfpathlineto{\pgfqpoint{3.112607in}{0.504832in}}%
\pgfpathlineto{\pgfqpoint{3.112699in}{0.505669in}}%
\pgfpathlineto{\pgfqpoint{3.112873in}{0.902207in}}%
\pgfpathlineto{\pgfqpoint{3.112949in}{1.326967in}}%
\pgfpathlineto{\pgfqpoint{3.113689in}{1.326940in}}%
\pgfpathlineto{\pgfqpoint{3.116491in}{1.326750in}}%
\pgfpathlineto{\pgfqpoint{3.118964in}{1.326673in}}%
\pgfpathlineto{\pgfqpoint{3.120994in}{1.326385in}}%
\pgfpathlineto{\pgfqpoint{3.121854in}{0.504474in}}%
\pgfpathlineto{\pgfqpoint{3.124986in}{0.504434in}}%
\pgfpathlineto{\pgfqpoint{3.129092in}{0.504832in}}%
\pgfpathlineto{\pgfqpoint{3.129184in}{0.505669in}}%
\pgfpathlineto{\pgfqpoint{3.129358in}{0.902206in}}%
\pgfpathlineto{\pgfqpoint{3.129434in}{1.326967in}}%
\pgfpathlineto{\pgfqpoint{3.130174in}{1.326940in}}%
\pgfpathlineto{\pgfqpoint{3.132976in}{1.326750in}}%
\pgfpathlineto{\pgfqpoint{3.135449in}{1.326673in}}%
\pgfpathlineto{\pgfqpoint{3.137479in}{1.326385in}}%
\pgfpathlineto{\pgfqpoint{3.138339in}{0.504474in}}%
\pgfpathlineto{\pgfqpoint{3.141471in}{0.504434in}}%
\pgfpathlineto{\pgfqpoint{3.145577in}{0.504832in}}%
\pgfpathlineto{\pgfqpoint{3.145669in}{0.505669in}}%
\pgfpathlineto{\pgfqpoint{3.145843in}{0.902206in}}%
\pgfpathlineto{\pgfqpoint{3.145919in}{1.326967in}}%
\pgfpathlineto{\pgfqpoint{3.146658in}{1.326940in}}%
\pgfpathlineto{\pgfqpoint{3.149461in}{1.326750in}}%
\pgfpathlineto{\pgfqpoint{3.151934in}{1.326673in}}%
\pgfpathlineto{\pgfqpoint{3.153964in}{1.326385in}}%
\pgfpathlineto{\pgfqpoint{3.154824in}{0.504474in}}%
\pgfpathlineto{\pgfqpoint{3.157956in}{0.504434in}}%
\pgfpathlineto{\pgfqpoint{3.162062in}{0.504832in}}%
\pgfpathlineto{\pgfqpoint{3.162154in}{0.505669in}}%
\pgfpathlineto{\pgfqpoint{3.162328in}{0.902206in}}%
\pgfpathlineto{\pgfqpoint{3.162403in}{1.326967in}}%
\pgfpathlineto{\pgfqpoint{3.163143in}{1.326940in}}%
\pgfpathlineto{\pgfqpoint{3.165946in}{1.326750in}}%
\pgfpathlineto{\pgfqpoint{3.168418in}{1.326673in}}%
\pgfpathlineto{\pgfqpoint{3.170132in}{1.326503in}}%
\pgfpathlineto{\pgfqpoint{3.170132in}{1.326503in}}%
\pgfusepath{stroke}%
\end{pgfscope}%
\begin{pgfscope}%
\pgfsetrectcap%
\pgfsetmiterjoin%
\pgfsetlinewidth{0.803000pt}%
\definecolor{currentstroke}{rgb}{0.000000,0.000000,0.000000}%
\pgfsetstrokecolor{currentstroke}%
\pgfsetdash{}{0pt}%
\pgfpathmoveto{\pgfqpoint{0.573768in}{0.462654in}}%
\pgfpathlineto{\pgfqpoint{0.573768in}{1.368904in}}%
\pgfusepath{stroke}%
\end{pgfscope}%
\begin{pgfscope}%
\pgfsetrectcap%
\pgfsetmiterjoin%
\pgfsetlinewidth{0.803000pt}%
\definecolor{currentstroke}{rgb}{0.000000,0.000000,0.000000}%
\pgfsetstrokecolor{currentstroke}%
\pgfsetdash{}{0pt}%
\pgfpathmoveto{\pgfqpoint{3.293768in}{0.462654in}}%
\pgfpathlineto{\pgfqpoint{3.293768in}{1.368904in}}%
\pgfusepath{stroke}%
\end{pgfscope}%
\begin{pgfscope}%
\pgfsetrectcap%
\pgfsetmiterjoin%
\pgfsetlinewidth{0.803000pt}%
\definecolor{currentstroke}{rgb}{0.000000,0.000000,0.000000}%
\pgfsetstrokecolor{currentstroke}%
\pgfsetdash{}{0pt}%
\pgfpathmoveto{\pgfqpoint{0.573768in}{0.462654in}}%
\pgfpathlineto{\pgfqpoint{3.293768in}{0.462654in}}%
\pgfusepath{stroke}%
\end{pgfscope}%
\begin{pgfscope}%
\pgfsetrectcap%
\pgfsetmiterjoin%
\pgfsetlinewidth{0.803000pt}%
\definecolor{currentstroke}{rgb}{0.000000,0.000000,0.000000}%
\pgfsetstrokecolor{currentstroke}%
\pgfsetdash{}{0pt}%
\pgfpathmoveto{\pgfqpoint{0.573768in}{1.368904in}}%
\pgfpathlineto{\pgfqpoint{3.293768in}{1.368904in}}%
\pgfusepath{stroke}%
\end{pgfscope}%
\begin{pgfscope}%
\definecolor{textcolor}{rgb}{0.000000,0.000000,0.000000}%
\pgfsetstrokecolor{textcolor}%
\pgfsetfillcolor{textcolor}%
\pgftext[x=0.628168in,y=1.305467in,left,top]{\color{textcolor}{\rmfamily\fontsize{8.000000}{9.600000}\selectfont\catcode`\^=\active\def^{\ifmmode\sp\else\^{}\fi}\catcode`\%=\active\def%{\%}(c)}}%
\end{pgfscope}%
\begin{pgfscope}%
\pgfsetbuttcap%
\pgfsetmiterjoin%
\definecolor{currentfill}{rgb}{1.000000,1.000000,1.000000}%
\pgfsetfillcolor{currentfill}%
\pgfsetlinewidth{0.000000pt}%
\definecolor{currentstroke}{rgb}{0.000000,0.000000,0.000000}%
\pgfsetstrokecolor{currentstroke}%
\pgfsetstrokeopacity{0.000000}%
\pgfsetdash{}{0pt}%
\pgfpathmoveto{\pgfqpoint{2.259013in}{3.125849in}}%
\pgfpathlineto{\pgfqpoint{3.238213in}{3.125849in}}%
\pgfpathlineto{\pgfqpoint{3.238213in}{3.452099in}}%
\pgfpathlineto{\pgfqpoint{2.259013in}{3.452099in}}%
\pgfpathlineto{\pgfqpoint{2.259013in}{3.125849in}}%
\pgfpathclose%
\pgfusepath{fill}%
\end{pgfscope}%
\begin{pgfscope}%
\pgfpathrectangle{\pgfqpoint{2.259013in}{3.125849in}}{\pgfqpoint{0.979200in}{0.326250in}}%
\pgfusepath{clip}%
\pgfsetbuttcap%
\pgfsetroundjoin%
\pgfsetlinewidth{0.401500pt}%
\definecolor{currentstroke}{rgb}{0.690196,0.690196,0.690196}%
\pgfsetstrokecolor{currentstroke}%
\pgfsetstrokeopacity{0.250000}%
\pgfsetdash{{1.480000pt}{0.640000pt}}{0.000000pt}%
\pgfpathmoveto{\pgfqpoint{2.259013in}{3.125849in}}%
\pgfpathlineto{\pgfqpoint{2.259013in}{3.452099in}}%
\pgfusepath{stroke}%
\end{pgfscope}%
\begin{pgfscope}%
\pgfsetbuttcap%
\pgfsetroundjoin%
\definecolor{currentfill}{rgb}{0.000000,0.000000,0.000000}%
\pgfsetfillcolor{currentfill}%
\pgfsetlinewidth{0.803000pt}%
\definecolor{currentstroke}{rgb}{0.000000,0.000000,0.000000}%
\pgfsetstrokecolor{currentstroke}%
\pgfsetdash{}{0pt}%
\pgfsys@defobject{currentmarker}{\pgfqpoint{0.000000in}{0.000000in}}{\pgfqpoint{0.000000in}{0.027778in}}{%
\pgfpathmoveto{\pgfqpoint{0.000000in}{0.000000in}}%
\pgfpathlineto{\pgfqpoint{0.000000in}{0.027778in}}%
\pgfusepath{stroke,fill}%
}%
\begin{pgfscope}%
\pgfsys@transformshift{2.259013in}{3.125849in}%
\pgfsys@useobject{currentmarker}{}%
\end{pgfscope}%
\end{pgfscope}%
\begin{pgfscope}%
\definecolor{textcolor}{rgb}{0.000000,0.000000,0.000000}%
\pgfsetstrokecolor{textcolor}%
\pgfsetfillcolor{textcolor}%
\pgftext[x=2.259013in,y=3.111960in,,top]{\color{textcolor}{\rmfamily\fontsize{5.000000}{6.000000}\selectfont\catcode`\^=\active\def^{\ifmmode\sp\else\^{}\fi}\catcode`\%=\active\def%{\%}5.9}}%
\end{pgfscope}%
\begin{pgfscope}%
\pgfpathrectangle{\pgfqpoint{2.259013in}{3.125849in}}{\pgfqpoint{0.979200in}{0.326250in}}%
\pgfusepath{clip}%
\pgfsetbuttcap%
\pgfsetroundjoin%
\pgfsetlinewidth{0.401500pt}%
\definecolor{currentstroke}{rgb}{0.690196,0.690196,0.690196}%
\pgfsetstrokecolor{currentstroke}%
\pgfsetstrokeopacity{0.250000}%
\pgfsetdash{{1.480000pt}{0.640000pt}}{0.000000pt}%
\pgfpathmoveto{\pgfqpoint{2.503813in}{3.125849in}}%
\pgfpathlineto{\pgfqpoint{2.503813in}{3.452099in}}%
\pgfusepath{stroke}%
\end{pgfscope}%
\begin{pgfscope}%
\pgfsetbuttcap%
\pgfsetroundjoin%
\definecolor{currentfill}{rgb}{0.000000,0.000000,0.000000}%
\pgfsetfillcolor{currentfill}%
\pgfsetlinewidth{0.803000pt}%
\definecolor{currentstroke}{rgb}{0.000000,0.000000,0.000000}%
\pgfsetstrokecolor{currentstroke}%
\pgfsetdash{}{0pt}%
\pgfsys@defobject{currentmarker}{\pgfqpoint{0.000000in}{0.000000in}}{\pgfqpoint{0.000000in}{0.027778in}}{%
\pgfpathmoveto{\pgfqpoint{0.000000in}{0.000000in}}%
\pgfpathlineto{\pgfqpoint{0.000000in}{0.027778in}}%
\pgfusepath{stroke,fill}%
}%
\begin{pgfscope}%
\pgfsys@transformshift{2.503813in}{3.125849in}%
\pgfsys@useobject{currentmarker}{}%
\end{pgfscope}%
\end{pgfscope}%
\begin{pgfscope}%
\definecolor{textcolor}{rgb}{0.000000,0.000000,0.000000}%
\pgfsetstrokecolor{textcolor}%
\pgfsetfillcolor{textcolor}%
\pgftext[x=2.503813in,y=3.111960in,,top]{\color{textcolor}{\rmfamily\fontsize{5.000000}{6.000000}\selectfont\catcode`\^=\active\def^{\ifmmode\sp\else\^{}\fi}\catcode`\%=\active\def%{\%}5.9}}%
\end{pgfscope}%
\begin{pgfscope}%
\pgfpathrectangle{\pgfqpoint{2.259013in}{3.125849in}}{\pgfqpoint{0.979200in}{0.326250in}}%
\pgfusepath{clip}%
\pgfsetbuttcap%
\pgfsetroundjoin%
\pgfsetlinewidth{0.401500pt}%
\definecolor{currentstroke}{rgb}{0.690196,0.690196,0.690196}%
\pgfsetstrokecolor{currentstroke}%
\pgfsetstrokeopacity{0.250000}%
\pgfsetdash{{1.480000pt}{0.640000pt}}{0.000000pt}%
\pgfpathmoveto{\pgfqpoint{2.748613in}{3.125849in}}%
\pgfpathlineto{\pgfqpoint{2.748613in}{3.452099in}}%
\pgfusepath{stroke}%
\end{pgfscope}%
\begin{pgfscope}%
\pgfsetbuttcap%
\pgfsetroundjoin%
\definecolor{currentfill}{rgb}{0.000000,0.000000,0.000000}%
\pgfsetfillcolor{currentfill}%
\pgfsetlinewidth{0.803000pt}%
\definecolor{currentstroke}{rgb}{0.000000,0.000000,0.000000}%
\pgfsetstrokecolor{currentstroke}%
\pgfsetdash{}{0pt}%
\pgfsys@defobject{currentmarker}{\pgfqpoint{0.000000in}{0.000000in}}{\pgfqpoint{0.000000in}{0.027778in}}{%
\pgfpathmoveto{\pgfqpoint{0.000000in}{0.000000in}}%
\pgfpathlineto{\pgfqpoint{0.000000in}{0.027778in}}%
\pgfusepath{stroke,fill}%
}%
\begin{pgfscope}%
\pgfsys@transformshift{2.748613in}{3.125849in}%
\pgfsys@useobject{currentmarker}{}%
\end{pgfscope}%
\end{pgfscope}%
\begin{pgfscope}%
\definecolor{textcolor}{rgb}{0.000000,0.000000,0.000000}%
\pgfsetstrokecolor{textcolor}%
\pgfsetfillcolor{textcolor}%
\pgftext[x=2.748613in,y=3.111960in,,top]{\color{textcolor}{\rmfamily\fontsize{5.000000}{6.000000}\selectfont\catcode`\^=\active\def^{\ifmmode\sp\else\^{}\fi}\catcode`\%=\active\def%{\%}6.0}}%
\end{pgfscope}%
\begin{pgfscope}%
\pgfpathrectangle{\pgfqpoint{2.259013in}{3.125849in}}{\pgfqpoint{0.979200in}{0.326250in}}%
\pgfusepath{clip}%
\pgfsetbuttcap%
\pgfsetroundjoin%
\pgfsetlinewidth{0.401500pt}%
\definecolor{currentstroke}{rgb}{0.690196,0.690196,0.690196}%
\pgfsetstrokecolor{currentstroke}%
\pgfsetstrokeopacity{0.250000}%
\pgfsetdash{{1.480000pt}{0.640000pt}}{0.000000pt}%
\pgfpathmoveto{\pgfqpoint{2.993413in}{3.125849in}}%
\pgfpathlineto{\pgfqpoint{2.993413in}{3.452099in}}%
\pgfusepath{stroke}%
\end{pgfscope}%
\begin{pgfscope}%
\pgfsetbuttcap%
\pgfsetroundjoin%
\definecolor{currentfill}{rgb}{0.000000,0.000000,0.000000}%
\pgfsetfillcolor{currentfill}%
\pgfsetlinewidth{0.803000pt}%
\definecolor{currentstroke}{rgb}{0.000000,0.000000,0.000000}%
\pgfsetstrokecolor{currentstroke}%
\pgfsetdash{}{0pt}%
\pgfsys@defobject{currentmarker}{\pgfqpoint{0.000000in}{0.000000in}}{\pgfqpoint{0.000000in}{0.027778in}}{%
\pgfpathmoveto{\pgfqpoint{0.000000in}{0.000000in}}%
\pgfpathlineto{\pgfqpoint{0.000000in}{0.027778in}}%
\pgfusepath{stroke,fill}%
}%
\begin{pgfscope}%
\pgfsys@transformshift{2.993413in}{3.125849in}%
\pgfsys@useobject{currentmarker}{}%
\end{pgfscope}%
\end{pgfscope}%
\begin{pgfscope}%
\definecolor{textcolor}{rgb}{0.000000,0.000000,0.000000}%
\pgfsetstrokecolor{textcolor}%
\pgfsetfillcolor{textcolor}%
\pgftext[x=2.993413in,y=3.111960in,,top]{\color{textcolor}{\rmfamily\fontsize{5.000000}{6.000000}\selectfont\catcode`\^=\active\def^{\ifmmode\sp\else\^{}\fi}\catcode`\%=\active\def%{\%}6.0}}%
\end{pgfscope}%
\begin{pgfscope}%
\pgfpathrectangle{\pgfqpoint{2.259013in}{3.125849in}}{\pgfqpoint{0.979200in}{0.326250in}}%
\pgfusepath{clip}%
\pgfsetbuttcap%
\pgfsetroundjoin%
\pgfsetlinewidth{0.401500pt}%
\definecolor{currentstroke}{rgb}{0.690196,0.690196,0.690196}%
\pgfsetstrokecolor{currentstroke}%
\pgfsetstrokeopacity{0.250000}%
\pgfsetdash{{1.480000pt}{0.640000pt}}{0.000000pt}%
\pgfpathmoveto{\pgfqpoint{3.238213in}{3.125849in}}%
\pgfpathlineto{\pgfqpoint{3.238213in}{3.452099in}}%
\pgfusepath{stroke}%
\end{pgfscope}%
\begin{pgfscope}%
\pgfsetbuttcap%
\pgfsetroundjoin%
\definecolor{currentfill}{rgb}{0.000000,0.000000,0.000000}%
\pgfsetfillcolor{currentfill}%
\pgfsetlinewidth{0.803000pt}%
\definecolor{currentstroke}{rgb}{0.000000,0.000000,0.000000}%
\pgfsetstrokecolor{currentstroke}%
\pgfsetdash{}{0pt}%
\pgfsys@defobject{currentmarker}{\pgfqpoint{0.000000in}{0.000000in}}{\pgfqpoint{0.000000in}{0.027778in}}{%
\pgfpathmoveto{\pgfqpoint{0.000000in}{0.000000in}}%
\pgfpathlineto{\pgfqpoint{0.000000in}{0.027778in}}%
\pgfusepath{stroke,fill}%
}%
\begin{pgfscope}%
\pgfsys@transformshift{3.238213in}{3.125849in}%
\pgfsys@useobject{currentmarker}{}%
\end{pgfscope}%
\end{pgfscope}%
\begin{pgfscope}%
\definecolor{textcolor}{rgb}{0.000000,0.000000,0.000000}%
\pgfsetstrokecolor{textcolor}%
\pgfsetfillcolor{textcolor}%
\pgftext[x=3.238213in,y=3.111960in,,top]{\color{textcolor}{\rmfamily\fontsize{5.000000}{6.000000}\selectfont\catcode`\^=\active\def^{\ifmmode\sp\else\^{}\fi}\catcode`\%=\active\def%{\%}6.0}}%
\end{pgfscope}%
\begin{pgfscope}%
\pgfpathrectangle{\pgfqpoint{2.259013in}{3.125849in}}{\pgfqpoint{0.979200in}{0.326250in}}%
\pgfusepath{clip}%
\pgfsetbuttcap%
\pgfsetroundjoin%
\pgfsetlinewidth{0.401500pt}%
\definecolor{currentstroke}{rgb}{0.690196,0.690196,0.690196}%
\pgfsetstrokecolor{currentstroke}%
\pgfsetstrokeopacity{0.250000}%
\pgfsetdash{{1.480000pt}{0.640000pt}}{0.000000pt}%
\pgfpathmoveto{\pgfqpoint{2.259013in}{3.125849in}}%
\pgfpathlineto{\pgfqpoint{3.238213in}{3.125849in}}%
\pgfusepath{stroke}%
\end{pgfscope}%
\begin{pgfscope}%
\pgfsetbuttcap%
\pgfsetroundjoin%
\definecolor{currentfill}{rgb}{0.000000,0.000000,0.000000}%
\pgfsetfillcolor{currentfill}%
\pgfsetlinewidth{0.803000pt}%
\definecolor{currentstroke}{rgb}{0.000000,0.000000,0.000000}%
\pgfsetstrokecolor{currentstroke}%
\pgfsetdash{}{0pt}%
\pgfsys@defobject{currentmarker}{\pgfqpoint{0.000000in}{0.000000in}}{\pgfqpoint{0.027778in}{0.000000in}}{%
\pgfpathmoveto{\pgfqpoint{0.000000in}{0.000000in}}%
\pgfpathlineto{\pgfqpoint{0.027778in}{0.000000in}}%
\pgfusepath{stroke,fill}%
}%
\begin{pgfscope}%
\pgfsys@transformshift{2.259013in}{3.125849in}%
\pgfsys@useobject{currentmarker}{}%
\end{pgfscope}%
\end{pgfscope}%
\begin{pgfscope}%
\definecolor{textcolor}{rgb}{0.000000,0.000000,0.000000}%
\pgfsetstrokecolor{textcolor}%
\pgfsetfillcolor{textcolor}%
\pgftext[x=2.075368in, y=3.101736in, left, base]{\color{textcolor}{\rmfamily\fontsize{5.000000}{6.000000}\selectfont\catcode`\^=\active\def^{\ifmmode\sp\else\^{}\fi}\catcode`\%=\active\def%{\%}9.88}}%
\end{pgfscope}%
\begin{pgfscope}%
\pgfpathrectangle{\pgfqpoint{2.259013in}{3.125849in}}{\pgfqpoint{0.979200in}{0.326250in}}%
\pgfusepath{clip}%
\pgfsetbuttcap%
\pgfsetroundjoin%
\pgfsetlinewidth{0.401500pt}%
\definecolor{currentstroke}{rgb}{0.690196,0.690196,0.690196}%
\pgfsetstrokecolor{currentstroke}%
\pgfsetstrokeopacity{0.250000}%
\pgfsetdash{{1.480000pt}{0.640000pt}}{0.000000pt}%
\pgfpathmoveto{\pgfqpoint{2.259013in}{3.288974in}}%
\pgfpathlineto{\pgfqpoint{3.238213in}{3.288974in}}%
\pgfusepath{stroke}%
\end{pgfscope}%
\begin{pgfscope}%
\pgfsetbuttcap%
\pgfsetroundjoin%
\definecolor{currentfill}{rgb}{0.000000,0.000000,0.000000}%
\pgfsetfillcolor{currentfill}%
\pgfsetlinewidth{0.803000pt}%
\definecolor{currentstroke}{rgb}{0.000000,0.000000,0.000000}%
\pgfsetstrokecolor{currentstroke}%
\pgfsetdash{}{0pt}%
\pgfsys@defobject{currentmarker}{\pgfqpoint{0.000000in}{0.000000in}}{\pgfqpoint{0.027778in}{0.000000in}}{%
\pgfpathmoveto{\pgfqpoint{0.000000in}{0.000000in}}%
\pgfpathlineto{\pgfqpoint{0.027778in}{0.000000in}}%
\pgfusepath{stroke,fill}%
}%
\begin{pgfscope}%
\pgfsys@transformshift{2.259013in}{3.288974in}%
\pgfsys@useobject{currentmarker}{}%
\end{pgfscope}%
\end{pgfscope}%
\begin{pgfscope}%
\definecolor{textcolor}{rgb}{0.000000,0.000000,0.000000}%
\pgfsetstrokecolor{textcolor}%
\pgfsetfillcolor{textcolor}%
\pgftext[x=2.075368in, y=3.264861in, left, base]{\color{textcolor}{\rmfamily\fontsize{5.000000}{6.000000}\selectfont\catcode`\^=\active\def^{\ifmmode\sp\else\^{}\fi}\catcode`\%=\active\def%{\%}9.95}}%
\end{pgfscope}%
\begin{pgfscope}%
\pgfpathrectangle{\pgfqpoint{2.259013in}{3.125849in}}{\pgfqpoint{0.979200in}{0.326250in}}%
\pgfusepath{clip}%
\pgfsetbuttcap%
\pgfsetroundjoin%
\pgfsetlinewidth{0.401500pt}%
\definecolor{currentstroke}{rgb}{0.690196,0.690196,0.690196}%
\pgfsetstrokecolor{currentstroke}%
\pgfsetstrokeopacity{0.250000}%
\pgfsetdash{{1.480000pt}{0.640000pt}}{0.000000pt}%
\pgfpathmoveto{\pgfqpoint{2.259013in}{3.452099in}}%
\pgfpathlineto{\pgfqpoint{3.238213in}{3.452099in}}%
\pgfusepath{stroke}%
\end{pgfscope}%
\begin{pgfscope}%
\pgfsetbuttcap%
\pgfsetroundjoin%
\definecolor{currentfill}{rgb}{0.000000,0.000000,0.000000}%
\pgfsetfillcolor{currentfill}%
\pgfsetlinewidth{0.803000pt}%
\definecolor{currentstroke}{rgb}{0.000000,0.000000,0.000000}%
\pgfsetstrokecolor{currentstroke}%
\pgfsetdash{}{0pt}%
\pgfsys@defobject{currentmarker}{\pgfqpoint{0.000000in}{0.000000in}}{\pgfqpoint{0.027778in}{0.000000in}}{%
\pgfpathmoveto{\pgfqpoint{0.000000in}{0.000000in}}%
\pgfpathlineto{\pgfqpoint{0.027778in}{0.000000in}}%
\pgfusepath{stroke,fill}%
}%
\begin{pgfscope}%
\pgfsys@transformshift{2.259013in}{3.452099in}%
\pgfsys@useobject{currentmarker}{}%
\end{pgfscope}%
\end{pgfscope}%
\begin{pgfscope}%
\definecolor{textcolor}{rgb}{0.000000,0.000000,0.000000}%
\pgfsetstrokecolor{textcolor}%
\pgfsetfillcolor{textcolor}%
\pgftext[x=2.028107in, y=3.427986in, left, base]{\color{textcolor}{\rmfamily\fontsize{5.000000}{6.000000}\selectfont\catcode`\^=\active\def^{\ifmmode\sp\else\^{}\fi}\catcode`\%=\active\def%{\%}10.02}}%
\end{pgfscope}%
\begin{pgfscope}%
\pgfpathrectangle{\pgfqpoint{2.259013in}{3.125849in}}{\pgfqpoint{0.979200in}{0.326250in}}%
\pgfusepath{clip}%
\pgfsetrectcap%
\pgfsetroundjoin%
\pgfsetlinewidth{1.505625pt}%
\definecolor{currentstroke}{rgb}{0.000000,0.447059,0.698039}%
\pgfsetstrokecolor{currentstroke}%
\pgfsetdash{}{0pt}%
\pgfpathmoveto{\pgfqpoint{2.255597in}{3.328854in}}%
\pgfpathlineto{\pgfqpoint{2.271652in}{3.299497in}}%
\pgfpathlineto{\pgfqpoint{2.284163in}{3.276987in}}%
\pgfpathlineto{\pgfqpoint{2.298851in}{3.253804in}}%
\pgfpathlineto{\pgfqpoint{2.313539in}{3.234122in}}%
\pgfpathlineto{\pgfqpoint{2.328227in}{3.217922in}}%
\pgfpathlineto{\pgfqpoint{2.338019in}{3.209049in}}%
\pgfpathlineto{\pgfqpoint{2.347811in}{3.201709in}}%
\pgfpathlineto{\pgfqpoint{2.357603in}{3.195898in}}%
\pgfpathlineto{\pgfqpoint{2.367395in}{3.191610in}}%
\pgfpathlineto{\pgfqpoint{2.377187in}{3.188839in}}%
\pgfpathlineto{\pgfqpoint{2.386979in}{3.187579in}}%
\pgfpathlineto{\pgfqpoint{2.396771in}{3.187823in}}%
\pgfpathlineto{\pgfqpoint{2.406563in}{3.189567in}}%
\pgfpathlineto{\pgfqpoint{2.416355in}{3.192804in}}%
\pgfpathlineto{\pgfqpoint{2.426147in}{3.197527in}}%
\pgfpathlineto{\pgfqpoint{2.435939in}{3.203730in}}%
\pgfpathlineto{\pgfqpoint{2.445731in}{3.211408in}}%
\pgfpathlineto{\pgfqpoint{2.455523in}{3.220553in}}%
\pgfpathlineto{\pgfqpoint{2.470211in}{3.237008in}}%
\pgfpathlineto{\pgfqpoint{2.484899in}{3.256728in}}%
\pgfpathlineto{\pgfqpoint{2.499587in}{3.279692in}}%
\pgfpathlineto{\pgfqpoint{2.514314in}{3.305944in}}%
\pgfpathlineto{\pgfqpoint{2.531252in}{3.336713in}}%
\pgfpathlineto{\pgfqpoint{2.545940in}{3.359670in}}%
\pgfpathlineto{\pgfqpoint{2.560628in}{3.379076in}}%
\pgfpathlineto{\pgfqpoint{2.575316in}{3.394949in}}%
\pgfpathlineto{\pgfqpoint{2.585108in}{3.403578in}}%
\pgfpathlineto{\pgfqpoint{2.594900in}{3.410650in}}%
\pgfpathlineto{\pgfqpoint{2.604692in}{3.416173in}}%
\pgfpathlineto{\pgfqpoint{2.614484in}{3.420152in}}%
\pgfpathlineto{\pgfqpoint{2.624276in}{3.422592in}}%
\pgfpathlineto{\pgfqpoint{2.634068in}{3.423501in}}%
\pgfpathlineto{\pgfqpoint{2.643860in}{3.422883in}}%
\pgfpathlineto{\pgfqpoint{2.653652in}{3.420746in}}%
\pgfpathlineto{\pgfqpoint{2.663444in}{3.417095in}}%
\pgfpathlineto{\pgfqpoint{2.673236in}{3.411936in}}%
\pgfpathlineto{\pgfqpoint{2.683028in}{3.405277in}}%
\pgfpathlineto{\pgfqpoint{2.692820in}{3.397124in}}%
\pgfpathlineto{\pgfqpoint{2.702612in}{3.387483in}}%
\pgfpathlineto{\pgfqpoint{2.717300in}{3.370246in}}%
\pgfpathlineto{\pgfqpoint{2.731988in}{3.349700in}}%
\pgfpathlineto{\pgfqpoint{2.745197in}{3.328416in}}%
\pgfpathlineto{\pgfqpoint{2.761252in}{3.299060in}}%
\pgfpathlineto{\pgfqpoint{2.773763in}{3.276551in}}%
\pgfpathlineto{\pgfqpoint{2.788451in}{3.253369in}}%
\pgfpathlineto{\pgfqpoint{2.803139in}{3.233687in}}%
\pgfpathlineto{\pgfqpoint{2.817827in}{3.217489in}}%
\pgfpathlineto{\pgfqpoint{2.827619in}{3.208616in}}%
\pgfpathlineto{\pgfqpoint{2.837411in}{3.201278in}}%
\pgfpathlineto{\pgfqpoint{2.847203in}{3.195468in}}%
\pgfpathlineto{\pgfqpoint{2.856995in}{3.191180in}}%
\pgfpathlineto{\pgfqpoint{2.866787in}{3.188410in}}%
\pgfpathlineto{\pgfqpoint{2.876579in}{3.187151in}}%
\pgfpathlineto{\pgfqpoint{2.886371in}{3.187396in}}%
\pgfpathlineto{\pgfqpoint{2.896163in}{3.189141in}}%
\pgfpathlineto{\pgfqpoint{2.905955in}{3.192379in}}%
\pgfpathlineto{\pgfqpoint{2.915747in}{3.197103in}}%
\pgfpathlineto{\pgfqpoint{2.925539in}{3.203307in}}%
\pgfpathlineto{\pgfqpoint{2.935331in}{3.210986in}}%
\pgfpathlineto{\pgfqpoint{2.945123in}{3.220132in}}%
\pgfpathlineto{\pgfqpoint{2.959811in}{3.236589in}}%
\pgfpathlineto{\pgfqpoint{2.974499in}{3.256311in}}%
\pgfpathlineto{\pgfqpoint{2.989187in}{3.279276in}}%
\pgfpathlineto{\pgfqpoint{3.003914in}{3.305530in}}%
\pgfpathlineto{\pgfqpoint{3.020852in}{3.336302in}}%
\pgfpathlineto{\pgfqpoint{3.035540in}{3.359261in}}%
\pgfpathlineto{\pgfqpoint{3.050228in}{3.378668in}}%
\pgfpathlineto{\pgfqpoint{3.064916in}{3.394543in}}%
\pgfpathlineto{\pgfqpoint{3.074708in}{3.403173in}}%
\pgfpathlineto{\pgfqpoint{3.084500in}{3.410248in}}%
\pgfpathlineto{\pgfqpoint{3.094292in}{3.415772in}}%
\pgfpathlineto{\pgfqpoint{3.104084in}{3.419752in}}%
\pgfpathlineto{\pgfqpoint{3.113876in}{3.422194in}}%
\pgfpathlineto{\pgfqpoint{3.123668in}{3.423104in}}%
\pgfpathlineto{\pgfqpoint{3.133460in}{3.422488in}}%
\pgfpathlineto{\pgfqpoint{3.143252in}{3.420352in}}%
\pgfpathlineto{\pgfqpoint{3.153044in}{3.416702in}}%
\pgfpathlineto{\pgfqpoint{3.162836in}{3.411546in}}%
\pgfpathlineto{\pgfqpoint{3.172628in}{3.404888in}}%
\pgfpathlineto{\pgfqpoint{3.182420in}{3.396736in}}%
\pgfpathlineto{\pgfqpoint{3.192212in}{3.387097in}}%
\pgfpathlineto{\pgfqpoint{3.206900in}{3.369863in}}%
\pgfpathlineto{\pgfqpoint{3.221588in}{3.349318in}}%
\pgfpathlineto{\pgfqpoint{3.234797in}{3.328037in}}%
\pgfpathlineto{\pgfqpoint{3.238213in}{3.322102in}}%
\pgfpathlineto{\pgfqpoint{3.238213in}{3.322102in}}%
\pgfusepath{stroke}%
\end{pgfscope}%
\begin{pgfscope}%
\pgfsetrectcap%
\pgfsetmiterjoin%
\pgfsetlinewidth{0.803000pt}%
\definecolor{currentstroke}{rgb}{0.000000,0.000000,0.000000}%
\pgfsetstrokecolor{currentstroke}%
\pgfsetdash{}{0pt}%
\pgfpathmoveto{\pgfqpoint{2.259013in}{3.125849in}}%
\pgfpathlineto{\pgfqpoint{2.259013in}{3.452099in}}%
\pgfusepath{stroke}%
\end{pgfscope}%
\begin{pgfscope}%
\pgfsetrectcap%
\pgfsetmiterjoin%
\pgfsetlinewidth{0.803000pt}%
\definecolor{currentstroke}{rgb}{0.000000,0.000000,0.000000}%
\pgfsetstrokecolor{currentstroke}%
\pgfsetdash{}{0pt}%
\pgfpathmoveto{\pgfqpoint{3.238213in}{3.125849in}}%
\pgfpathlineto{\pgfqpoint{3.238213in}{3.452099in}}%
\pgfusepath{stroke}%
\end{pgfscope}%
\begin{pgfscope}%
\pgfsetrectcap%
\pgfsetmiterjoin%
\pgfsetlinewidth{0.803000pt}%
\definecolor{currentstroke}{rgb}{0.000000,0.000000,0.000000}%
\pgfsetstrokecolor{currentstroke}%
\pgfsetdash{}{0pt}%
\pgfpathmoveto{\pgfqpoint{2.259013in}{3.125849in}}%
\pgfpathlineto{\pgfqpoint{3.238213in}{3.125849in}}%
\pgfusepath{stroke}%
\end{pgfscope}%
\begin{pgfscope}%
\pgfsetrectcap%
\pgfsetmiterjoin%
\pgfsetlinewidth{0.803000pt}%
\definecolor{currentstroke}{rgb}{0.000000,0.000000,0.000000}%
\pgfsetstrokecolor{currentstroke}%
\pgfsetdash{}{0pt}%
\pgfpathmoveto{\pgfqpoint{2.259013in}{3.452099in}}%
\pgfpathlineto{\pgfqpoint{3.238213in}{3.452099in}}%
\pgfusepath{stroke}%
\end{pgfscope}%
\begin{pgfscope}%
\pgfsetbuttcap%
\pgfsetmiterjoin%
\definecolor{currentfill}{rgb}{1.000000,1.000000,1.000000}%
\pgfsetfillcolor{currentfill}%
\pgfsetlinewidth{0.000000pt}%
\definecolor{currentstroke}{rgb}{0.000000,0.000000,0.000000}%
\pgfsetstrokecolor{currentstroke}%
\pgfsetstrokeopacity{0.000000}%
\pgfsetdash{}{0pt}%
\pgfpathmoveto{\pgfqpoint{2.259013in}{2.056474in}}%
\pgfpathlineto{\pgfqpoint{3.238213in}{2.056474in}}%
\pgfpathlineto{\pgfqpoint{3.238213in}{2.382724in}}%
\pgfpathlineto{\pgfqpoint{2.259013in}{2.382724in}}%
\pgfpathlineto{\pgfqpoint{2.259013in}{2.056474in}}%
\pgfpathclose%
\pgfusepath{fill}%
\end{pgfscope}%
\begin{pgfscope}%
\pgfpathrectangle{\pgfqpoint{2.259013in}{2.056474in}}{\pgfqpoint{0.979200in}{0.326250in}}%
\pgfusepath{clip}%
\pgfsetbuttcap%
\pgfsetroundjoin%
\pgfsetlinewidth{0.401500pt}%
\definecolor{currentstroke}{rgb}{0.690196,0.690196,0.690196}%
\pgfsetstrokecolor{currentstroke}%
\pgfsetstrokeopacity{0.250000}%
\pgfsetdash{{1.480000pt}{0.640000pt}}{0.000000pt}%
\pgfpathmoveto{\pgfqpoint{2.259013in}{2.056474in}}%
\pgfpathlineto{\pgfqpoint{2.259013in}{2.382724in}}%
\pgfusepath{stroke}%
\end{pgfscope}%
\begin{pgfscope}%
\pgfsetbuttcap%
\pgfsetroundjoin%
\definecolor{currentfill}{rgb}{0.000000,0.000000,0.000000}%
\pgfsetfillcolor{currentfill}%
\pgfsetlinewidth{0.803000pt}%
\definecolor{currentstroke}{rgb}{0.000000,0.000000,0.000000}%
\pgfsetstrokecolor{currentstroke}%
\pgfsetdash{}{0pt}%
\pgfsys@defobject{currentmarker}{\pgfqpoint{0.000000in}{0.000000in}}{\pgfqpoint{0.000000in}{0.027778in}}{%
\pgfpathmoveto{\pgfqpoint{0.000000in}{0.000000in}}%
\pgfpathlineto{\pgfqpoint{0.000000in}{0.027778in}}%
\pgfusepath{stroke,fill}%
}%
\begin{pgfscope}%
\pgfsys@transformshift{2.259013in}{2.056474in}%
\pgfsys@useobject{currentmarker}{}%
\end{pgfscope}%
\end{pgfscope}%
\begin{pgfscope}%
\definecolor{textcolor}{rgb}{0.000000,0.000000,0.000000}%
\pgfsetstrokecolor{textcolor}%
\pgfsetfillcolor{textcolor}%
\pgftext[x=2.259013in,y=2.042585in,,top]{\color{textcolor}{\rmfamily\fontsize{5.000000}{6.000000}\selectfont\catcode`\^=\active\def^{\ifmmode\sp\else\^{}\fi}\catcode`\%=\active\def%{\%}5.9}}%
\end{pgfscope}%
\begin{pgfscope}%
\pgfpathrectangle{\pgfqpoint{2.259013in}{2.056474in}}{\pgfqpoint{0.979200in}{0.326250in}}%
\pgfusepath{clip}%
\pgfsetbuttcap%
\pgfsetroundjoin%
\pgfsetlinewidth{0.401500pt}%
\definecolor{currentstroke}{rgb}{0.690196,0.690196,0.690196}%
\pgfsetstrokecolor{currentstroke}%
\pgfsetstrokeopacity{0.250000}%
\pgfsetdash{{1.480000pt}{0.640000pt}}{0.000000pt}%
\pgfpathmoveto{\pgfqpoint{2.503813in}{2.056474in}}%
\pgfpathlineto{\pgfqpoint{2.503813in}{2.382724in}}%
\pgfusepath{stroke}%
\end{pgfscope}%
\begin{pgfscope}%
\pgfsetbuttcap%
\pgfsetroundjoin%
\definecolor{currentfill}{rgb}{0.000000,0.000000,0.000000}%
\pgfsetfillcolor{currentfill}%
\pgfsetlinewidth{0.803000pt}%
\definecolor{currentstroke}{rgb}{0.000000,0.000000,0.000000}%
\pgfsetstrokecolor{currentstroke}%
\pgfsetdash{}{0pt}%
\pgfsys@defobject{currentmarker}{\pgfqpoint{0.000000in}{0.000000in}}{\pgfqpoint{0.000000in}{0.027778in}}{%
\pgfpathmoveto{\pgfqpoint{0.000000in}{0.000000in}}%
\pgfpathlineto{\pgfqpoint{0.000000in}{0.027778in}}%
\pgfusepath{stroke,fill}%
}%
\begin{pgfscope}%
\pgfsys@transformshift{2.503813in}{2.056474in}%
\pgfsys@useobject{currentmarker}{}%
\end{pgfscope}%
\end{pgfscope}%
\begin{pgfscope}%
\definecolor{textcolor}{rgb}{0.000000,0.000000,0.000000}%
\pgfsetstrokecolor{textcolor}%
\pgfsetfillcolor{textcolor}%
\pgftext[x=2.503813in,y=2.042585in,,top]{\color{textcolor}{\rmfamily\fontsize{5.000000}{6.000000}\selectfont\catcode`\^=\active\def^{\ifmmode\sp\else\^{}\fi}\catcode`\%=\active\def%{\%}5.9}}%
\end{pgfscope}%
\begin{pgfscope}%
\pgfpathrectangle{\pgfqpoint{2.259013in}{2.056474in}}{\pgfqpoint{0.979200in}{0.326250in}}%
\pgfusepath{clip}%
\pgfsetbuttcap%
\pgfsetroundjoin%
\pgfsetlinewidth{0.401500pt}%
\definecolor{currentstroke}{rgb}{0.690196,0.690196,0.690196}%
\pgfsetstrokecolor{currentstroke}%
\pgfsetstrokeopacity{0.250000}%
\pgfsetdash{{1.480000pt}{0.640000pt}}{0.000000pt}%
\pgfpathmoveto{\pgfqpoint{2.748613in}{2.056474in}}%
\pgfpathlineto{\pgfqpoint{2.748613in}{2.382724in}}%
\pgfusepath{stroke}%
\end{pgfscope}%
\begin{pgfscope}%
\pgfsetbuttcap%
\pgfsetroundjoin%
\definecolor{currentfill}{rgb}{0.000000,0.000000,0.000000}%
\pgfsetfillcolor{currentfill}%
\pgfsetlinewidth{0.803000pt}%
\definecolor{currentstroke}{rgb}{0.000000,0.000000,0.000000}%
\pgfsetstrokecolor{currentstroke}%
\pgfsetdash{}{0pt}%
\pgfsys@defobject{currentmarker}{\pgfqpoint{0.000000in}{0.000000in}}{\pgfqpoint{0.000000in}{0.027778in}}{%
\pgfpathmoveto{\pgfqpoint{0.000000in}{0.000000in}}%
\pgfpathlineto{\pgfqpoint{0.000000in}{0.027778in}}%
\pgfusepath{stroke,fill}%
}%
\begin{pgfscope}%
\pgfsys@transformshift{2.748613in}{2.056474in}%
\pgfsys@useobject{currentmarker}{}%
\end{pgfscope}%
\end{pgfscope}%
\begin{pgfscope}%
\definecolor{textcolor}{rgb}{0.000000,0.000000,0.000000}%
\pgfsetstrokecolor{textcolor}%
\pgfsetfillcolor{textcolor}%
\pgftext[x=2.748613in,y=2.042585in,,top]{\color{textcolor}{\rmfamily\fontsize{5.000000}{6.000000}\selectfont\catcode`\^=\active\def^{\ifmmode\sp\else\^{}\fi}\catcode`\%=\active\def%{\%}6.0}}%
\end{pgfscope}%
\begin{pgfscope}%
\pgfpathrectangle{\pgfqpoint{2.259013in}{2.056474in}}{\pgfqpoint{0.979200in}{0.326250in}}%
\pgfusepath{clip}%
\pgfsetbuttcap%
\pgfsetroundjoin%
\pgfsetlinewidth{0.401500pt}%
\definecolor{currentstroke}{rgb}{0.690196,0.690196,0.690196}%
\pgfsetstrokecolor{currentstroke}%
\pgfsetstrokeopacity{0.250000}%
\pgfsetdash{{1.480000pt}{0.640000pt}}{0.000000pt}%
\pgfpathmoveto{\pgfqpoint{2.993413in}{2.056474in}}%
\pgfpathlineto{\pgfqpoint{2.993413in}{2.382724in}}%
\pgfusepath{stroke}%
\end{pgfscope}%
\begin{pgfscope}%
\pgfsetbuttcap%
\pgfsetroundjoin%
\definecolor{currentfill}{rgb}{0.000000,0.000000,0.000000}%
\pgfsetfillcolor{currentfill}%
\pgfsetlinewidth{0.803000pt}%
\definecolor{currentstroke}{rgb}{0.000000,0.000000,0.000000}%
\pgfsetstrokecolor{currentstroke}%
\pgfsetdash{}{0pt}%
\pgfsys@defobject{currentmarker}{\pgfqpoint{0.000000in}{0.000000in}}{\pgfqpoint{0.000000in}{0.027778in}}{%
\pgfpathmoveto{\pgfqpoint{0.000000in}{0.000000in}}%
\pgfpathlineto{\pgfqpoint{0.000000in}{0.027778in}}%
\pgfusepath{stroke,fill}%
}%
\begin{pgfscope}%
\pgfsys@transformshift{2.993413in}{2.056474in}%
\pgfsys@useobject{currentmarker}{}%
\end{pgfscope}%
\end{pgfscope}%
\begin{pgfscope}%
\definecolor{textcolor}{rgb}{0.000000,0.000000,0.000000}%
\pgfsetstrokecolor{textcolor}%
\pgfsetfillcolor{textcolor}%
\pgftext[x=2.993413in,y=2.042585in,,top]{\color{textcolor}{\rmfamily\fontsize{5.000000}{6.000000}\selectfont\catcode`\^=\active\def^{\ifmmode\sp\else\^{}\fi}\catcode`\%=\active\def%{\%}6.0}}%
\end{pgfscope}%
\begin{pgfscope}%
\pgfpathrectangle{\pgfqpoint{2.259013in}{2.056474in}}{\pgfqpoint{0.979200in}{0.326250in}}%
\pgfusepath{clip}%
\pgfsetbuttcap%
\pgfsetroundjoin%
\pgfsetlinewidth{0.401500pt}%
\definecolor{currentstroke}{rgb}{0.690196,0.690196,0.690196}%
\pgfsetstrokecolor{currentstroke}%
\pgfsetstrokeopacity{0.250000}%
\pgfsetdash{{1.480000pt}{0.640000pt}}{0.000000pt}%
\pgfpathmoveto{\pgfqpoint{3.238213in}{2.056474in}}%
\pgfpathlineto{\pgfqpoint{3.238213in}{2.382724in}}%
\pgfusepath{stroke}%
\end{pgfscope}%
\begin{pgfscope}%
\pgfsetbuttcap%
\pgfsetroundjoin%
\definecolor{currentfill}{rgb}{0.000000,0.000000,0.000000}%
\pgfsetfillcolor{currentfill}%
\pgfsetlinewidth{0.803000pt}%
\definecolor{currentstroke}{rgb}{0.000000,0.000000,0.000000}%
\pgfsetstrokecolor{currentstroke}%
\pgfsetdash{}{0pt}%
\pgfsys@defobject{currentmarker}{\pgfqpoint{0.000000in}{0.000000in}}{\pgfqpoint{0.000000in}{0.027778in}}{%
\pgfpathmoveto{\pgfqpoint{0.000000in}{0.000000in}}%
\pgfpathlineto{\pgfqpoint{0.000000in}{0.027778in}}%
\pgfusepath{stroke,fill}%
}%
\begin{pgfscope}%
\pgfsys@transformshift{3.238213in}{2.056474in}%
\pgfsys@useobject{currentmarker}{}%
\end{pgfscope}%
\end{pgfscope}%
\begin{pgfscope}%
\definecolor{textcolor}{rgb}{0.000000,0.000000,0.000000}%
\pgfsetstrokecolor{textcolor}%
\pgfsetfillcolor{textcolor}%
\pgftext[x=3.238213in,y=2.042585in,,top]{\color{textcolor}{\rmfamily\fontsize{5.000000}{6.000000}\selectfont\catcode`\^=\active\def^{\ifmmode\sp\else\^{}\fi}\catcode`\%=\active\def%{\%}6.0}}%
\end{pgfscope}%
\begin{pgfscope}%
\pgfpathrectangle{\pgfqpoint{2.259013in}{2.056474in}}{\pgfqpoint{0.979200in}{0.326250in}}%
\pgfusepath{clip}%
\pgfsetbuttcap%
\pgfsetroundjoin%
\pgfsetlinewidth{0.401500pt}%
\definecolor{currentstroke}{rgb}{0.690196,0.690196,0.690196}%
\pgfsetstrokecolor{currentstroke}%
\pgfsetstrokeopacity{0.250000}%
\pgfsetdash{{1.480000pt}{0.640000pt}}{0.000000pt}%
\pgfpathmoveto{\pgfqpoint{2.259013in}{2.056474in}}%
\pgfpathlineto{\pgfqpoint{3.238213in}{2.056474in}}%
\pgfusepath{stroke}%
\end{pgfscope}%
\begin{pgfscope}%
\pgfsetbuttcap%
\pgfsetroundjoin%
\definecolor{currentfill}{rgb}{0.000000,0.000000,0.000000}%
\pgfsetfillcolor{currentfill}%
\pgfsetlinewidth{0.803000pt}%
\definecolor{currentstroke}{rgb}{0.000000,0.000000,0.000000}%
\pgfsetstrokecolor{currentstroke}%
\pgfsetdash{}{0pt}%
\pgfsys@defobject{currentmarker}{\pgfqpoint{0.000000in}{0.000000in}}{\pgfqpoint{0.027778in}{0.000000in}}{%
\pgfpathmoveto{\pgfqpoint{0.000000in}{0.000000in}}%
\pgfpathlineto{\pgfqpoint{0.027778in}{0.000000in}}%
\pgfusepath{stroke,fill}%
}%
\begin{pgfscope}%
\pgfsys@transformshift{2.259013in}{2.056474in}%
\pgfsys@useobject{currentmarker}{}%
\end{pgfscope}%
\end{pgfscope}%
\begin{pgfscope}%
\definecolor{textcolor}{rgb}{0.000000,0.000000,0.000000}%
\pgfsetstrokecolor{textcolor}%
\pgfsetfillcolor{textcolor}%
\pgftext[x=2.075368in, y=2.032361in, left, base]{\color{textcolor}{\rmfamily\fontsize{5.000000}{6.000000}\selectfont\catcode`\^=\active\def^{\ifmmode\sp\else\^{}\fi}\catcode`\%=\active\def%{\%}3.95}}%
\end{pgfscope}%
\begin{pgfscope}%
\pgfpathrectangle{\pgfqpoint{2.259013in}{2.056474in}}{\pgfqpoint{0.979200in}{0.326250in}}%
\pgfusepath{clip}%
\pgfsetbuttcap%
\pgfsetroundjoin%
\pgfsetlinewidth{0.401500pt}%
\definecolor{currentstroke}{rgb}{0.690196,0.690196,0.690196}%
\pgfsetstrokecolor{currentstroke}%
\pgfsetstrokeopacity{0.250000}%
\pgfsetdash{{1.480000pt}{0.640000pt}}{0.000000pt}%
\pgfpathmoveto{\pgfqpoint{2.259013in}{2.219599in}}%
\pgfpathlineto{\pgfqpoint{3.238213in}{2.219599in}}%
\pgfusepath{stroke}%
\end{pgfscope}%
\begin{pgfscope}%
\pgfsetbuttcap%
\pgfsetroundjoin%
\definecolor{currentfill}{rgb}{0.000000,0.000000,0.000000}%
\pgfsetfillcolor{currentfill}%
\pgfsetlinewidth{0.803000pt}%
\definecolor{currentstroke}{rgb}{0.000000,0.000000,0.000000}%
\pgfsetstrokecolor{currentstroke}%
\pgfsetdash{}{0pt}%
\pgfsys@defobject{currentmarker}{\pgfqpoint{0.000000in}{0.000000in}}{\pgfqpoint{0.027778in}{0.000000in}}{%
\pgfpathmoveto{\pgfqpoint{0.000000in}{0.000000in}}%
\pgfpathlineto{\pgfqpoint{0.027778in}{0.000000in}}%
\pgfusepath{stroke,fill}%
}%
\begin{pgfscope}%
\pgfsys@transformshift{2.259013in}{2.219599in}%
\pgfsys@useobject{currentmarker}{}%
\end{pgfscope}%
\end{pgfscope}%
\begin{pgfscope}%
\definecolor{textcolor}{rgb}{0.000000,0.000000,0.000000}%
\pgfsetstrokecolor{textcolor}%
\pgfsetfillcolor{textcolor}%
\pgftext[x=2.075368in, y=2.195486in, left, base]{\color{textcolor}{\rmfamily\fontsize{5.000000}{6.000000}\selectfont\catcode`\^=\active\def^{\ifmmode\sp\else\^{}\fi}\catcode`\%=\active\def%{\%}5.00}}%
\end{pgfscope}%
\begin{pgfscope}%
\pgfpathrectangle{\pgfqpoint{2.259013in}{2.056474in}}{\pgfqpoint{0.979200in}{0.326250in}}%
\pgfusepath{clip}%
\pgfsetbuttcap%
\pgfsetroundjoin%
\pgfsetlinewidth{0.401500pt}%
\definecolor{currentstroke}{rgb}{0.690196,0.690196,0.690196}%
\pgfsetstrokecolor{currentstroke}%
\pgfsetstrokeopacity{0.250000}%
\pgfsetdash{{1.480000pt}{0.640000pt}}{0.000000pt}%
\pgfpathmoveto{\pgfqpoint{2.259013in}{2.382724in}}%
\pgfpathlineto{\pgfqpoint{3.238213in}{2.382724in}}%
\pgfusepath{stroke}%
\end{pgfscope}%
\begin{pgfscope}%
\pgfsetbuttcap%
\pgfsetroundjoin%
\definecolor{currentfill}{rgb}{0.000000,0.000000,0.000000}%
\pgfsetfillcolor{currentfill}%
\pgfsetlinewidth{0.803000pt}%
\definecolor{currentstroke}{rgb}{0.000000,0.000000,0.000000}%
\pgfsetstrokecolor{currentstroke}%
\pgfsetdash{}{0pt}%
\pgfsys@defobject{currentmarker}{\pgfqpoint{0.000000in}{0.000000in}}{\pgfqpoint{0.027778in}{0.000000in}}{%
\pgfpathmoveto{\pgfqpoint{0.000000in}{0.000000in}}%
\pgfpathlineto{\pgfqpoint{0.027778in}{0.000000in}}%
\pgfusepath{stroke,fill}%
}%
\begin{pgfscope}%
\pgfsys@transformshift{2.259013in}{2.382724in}%
\pgfsys@useobject{currentmarker}{}%
\end{pgfscope}%
\end{pgfscope}%
\begin{pgfscope}%
\definecolor{textcolor}{rgb}{0.000000,0.000000,0.000000}%
\pgfsetstrokecolor{textcolor}%
\pgfsetfillcolor{textcolor}%
\pgftext[x=2.075368in, y=2.358611in, left, base]{\color{textcolor}{\rmfamily\fontsize{5.000000}{6.000000}\selectfont\catcode`\^=\active\def^{\ifmmode\sp\else\^{}\fi}\catcode`\%=\active\def%{\%}6.05}}%
\end{pgfscope}%
\begin{pgfscope}%
\pgfpathrectangle{\pgfqpoint{2.259013in}{2.056474in}}{\pgfqpoint{0.979200in}{0.326250in}}%
\pgfusepath{clip}%
\pgfsetrectcap%
\pgfsetroundjoin%
\pgfsetlinewidth{1.505625pt}%
\definecolor{currentstroke}{rgb}{0.835294,0.368627,0.000000}%
\pgfsetstrokecolor{currentstroke}%
\pgfsetdash{}{0pt}%
\pgfpathmoveto{\pgfqpoint{2.255597in}{2.077370in}}%
\pgfpathlineto{\pgfqpoint{2.269716in}{2.060146in}}%
\pgfpathlineto{\pgfqpoint{2.271652in}{2.061438in}}%
\pgfpathlineto{\pgfqpoint{2.276347in}{2.066827in}}%
\pgfpathlineto{\pgfqpoint{2.421251in}{2.251945in}}%
\pgfpathlineto{\pgfqpoint{2.514314in}{2.370477in}}%
\pgfpathlineto{\pgfqpoint{2.516820in}{2.371918in}}%
\pgfpathlineto{\pgfqpoint{2.519071in}{2.370498in}}%
\pgfpathlineto{\pgfqpoint{2.759687in}{2.060285in}}%
\pgfpathlineto{\pgfqpoint{2.762132in}{2.062260in}}%
\pgfpathlineto{\pgfqpoint{2.769877in}{2.071812in}}%
\pgfpathlineto{\pgfqpoint{3.006420in}{2.371914in}}%
\pgfpathlineto{\pgfqpoint{3.008671in}{2.370495in}}%
\pgfpathlineto{\pgfqpoint{3.238213in}{2.072948in}}%
\pgfpathlineto{\pgfqpoint{3.238213in}{2.072948in}}%
\pgfusepath{stroke}%
\end{pgfscope}%
\begin{pgfscope}%
\pgfsetrectcap%
\pgfsetmiterjoin%
\pgfsetlinewidth{0.803000pt}%
\definecolor{currentstroke}{rgb}{0.000000,0.000000,0.000000}%
\pgfsetstrokecolor{currentstroke}%
\pgfsetdash{}{0pt}%
\pgfpathmoveto{\pgfqpoint{2.259013in}{2.056474in}}%
\pgfpathlineto{\pgfqpoint{2.259013in}{2.382724in}}%
\pgfusepath{stroke}%
\end{pgfscope}%
\begin{pgfscope}%
\pgfsetrectcap%
\pgfsetmiterjoin%
\pgfsetlinewidth{0.803000pt}%
\definecolor{currentstroke}{rgb}{0.000000,0.000000,0.000000}%
\pgfsetstrokecolor{currentstroke}%
\pgfsetdash{}{0pt}%
\pgfpathmoveto{\pgfqpoint{3.238213in}{2.056474in}}%
\pgfpathlineto{\pgfqpoint{3.238213in}{2.382724in}}%
\pgfusepath{stroke}%
\end{pgfscope}%
\begin{pgfscope}%
\pgfsetrectcap%
\pgfsetmiterjoin%
\pgfsetlinewidth{0.803000pt}%
\definecolor{currentstroke}{rgb}{0.000000,0.000000,0.000000}%
\pgfsetstrokecolor{currentstroke}%
\pgfsetdash{}{0pt}%
\pgfpathmoveto{\pgfqpoint{2.259013in}{2.056474in}}%
\pgfpathlineto{\pgfqpoint{3.238213in}{2.056474in}}%
\pgfusepath{stroke}%
\end{pgfscope}%
\begin{pgfscope}%
\pgfsetrectcap%
\pgfsetmiterjoin%
\pgfsetlinewidth{0.803000pt}%
\definecolor{currentstroke}{rgb}{0.000000,0.000000,0.000000}%
\pgfsetstrokecolor{currentstroke}%
\pgfsetdash{}{0pt}%
\pgfpathmoveto{\pgfqpoint{2.259013in}{2.382724in}}%
\pgfpathlineto{\pgfqpoint{3.238213in}{2.382724in}}%
\pgfusepath{stroke}%
\end{pgfscope}%
\begin{pgfscope}%
\pgfsetbuttcap%
\pgfsetmiterjoin%
\definecolor{currentfill}{rgb}{1.000000,1.000000,1.000000}%
\pgfsetfillcolor{currentfill}%
\pgfsetlinewidth{0.000000pt}%
\definecolor{currentstroke}{rgb}{0.000000,0.000000,0.000000}%
\pgfsetstrokecolor{currentstroke}%
\pgfsetstrokeopacity{0.000000}%
\pgfsetdash{}{0pt}%
\pgfpathmoveto{\pgfqpoint{2.259013in}{0.987099in}}%
\pgfpathlineto{\pgfqpoint{3.238213in}{0.987099in}}%
\pgfpathlineto{\pgfqpoint{3.238213in}{1.313349in}}%
\pgfpathlineto{\pgfqpoint{2.259013in}{1.313349in}}%
\pgfpathlineto{\pgfqpoint{2.259013in}{0.987099in}}%
\pgfpathclose%
\pgfusepath{fill}%
\end{pgfscope}%
\begin{pgfscope}%
\pgfpathrectangle{\pgfqpoint{2.259013in}{0.987099in}}{\pgfqpoint{0.979200in}{0.326250in}}%
\pgfusepath{clip}%
\pgfsetbuttcap%
\pgfsetroundjoin%
\pgfsetlinewidth{0.401500pt}%
\definecolor{currentstroke}{rgb}{0.690196,0.690196,0.690196}%
\pgfsetstrokecolor{currentstroke}%
\pgfsetstrokeopacity{0.250000}%
\pgfsetdash{{1.480000pt}{0.640000pt}}{0.000000pt}%
\pgfpathmoveto{\pgfqpoint{2.259013in}{0.987099in}}%
\pgfpathlineto{\pgfqpoint{2.259013in}{1.313349in}}%
\pgfusepath{stroke}%
\end{pgfscope}%
\begin{pgfscope}%
\pgfsetbuttcap%
\pgfsetroundjoin%
\definecolor{currentfill}{rgb}{0.000000,0.000000,0.000000}%
\pgfsetfillcolor{currentfill}%
\pgfsetlinewidth{0.803000pt}%
\definecolor{currentstroke}{rgb}{0.000000,0.000000,0.000000}%
\pgfsetstrokecolor{currentstroke}%
\pgfsetdash{}{0pt}%
\pgfsys@defobject{currentmarker}{\pgfqpoint{0.000000in}{0.000000in}}{\pgfqpoint{0.000000in}{0.027778in}}{%
\pgfpathmoveto{\pgfqpoint{0.000000in}{0.000000in}}%
\pgfpathlineto{\pgfqpoint{0.000000in}{0.027778in}}%
\pgfusepath{stroke,fill}%
}%
\begin{pgfscope}%
\pgfsys@transformshift{2.259013in}{0.987099in}%
\pgfsys@useobject{currentmarker}{}%
\end{pgfscope}%
\end{pgfscope}%
\begin{pgfscope}%
\definecolor{textcolor}{rgb}{0.000000,0.000000,0.000000}%
\pgfsetstrokecolor{textcolor}%
\pgfsetfillcolor{textcolor}%
\pgftext[x=2.259013in,y=0.973210in,,top]{\color{textcolor}{\rmfamily\fontsize{5.000000}{6.000000}\selectfont\catcode`\^=\active\def^{\ifmmode\sp\else\^{}\fi}\catcode`\%=\active\def%{\%}5.9}}%
\end{pgfscope}%
\begin{pgfscope}%
\pgfpathrectangle{\pgfqpoint{2.259013in}{0.987099in}}{\pgfqpoint{0.979200in}{0.326250in}}%
\pgfusepath{clip}%
\pgfsetbuttcap%
\pgfsetroundjoin%
\pgfsetlinewidth{0.401500pt}%
\definecolor{currentstroke}{rgb}{0.690196,0.690196,0.690196}%
\pgfsetstrokecolor{currentstroke}%
\pgfsetstrokeopacity{0.250000}%
\pgfsetdash{{1.480000pt}{0.640000pt}}{0.000000pt}%
\pgfpathmoveto{\pgfqpoint{2.503813in}{0.987099in}}%
\pgfpathlineto{\pgfqpoint{2.503813in}{1.313349in}}%
\pgfusepath{stroke}%
\end{pgfscope}%
\begin{pgfscope}%
\pgfsetbuttcap%
\pgfsetroundjoin%
\definecolor{currentfill}{rgb}{0.000000,0.000000,0.000000}%
\pgfsetfillcolor{currentfill}%
\pgfsetlinewidth{0.803000pt}%
\definecolor{currentstroke}{rgb}{0.000000,0.000000,0.000000}%
\pgfsetstrokecolor{currentstroke}%
\pgfsetdash{}{0pt}%
\pgfsys@defobject{currentmarker}{\pgfqpoint{0.000000in}{0.000000in}}{\pgfqpoint{0.000000in}{0.027778in}}{%
\pgfpathmoveto{\pgfqpoint{0.000000in}{0.000000in}}%
\pgfpathlineto{\pgfqpoint{0.000000in}{0.027778in}}%
\pgfusepath{stroke,fill}%
}%
\begin{pgfscope}%
\pgfsys@transformshift{2.503813in}{0.987099in}%
\pgfsys@useobject{currentmarker}{}%
\end{pgfscope}%
\end{pgfscope}%
\begin{pgfscope}%
\definecolor{textcolor}{rgb}{0.000000,0.000000,0.000000}%
\pgfsetstrokecolor{textcolor}%
\pgfsetfillcolor{textcolor}%
\pgftext[x=2.503813in,y=0.973210in,,top]{\color{textcolor}{\rmfamily\fontsize{5.000000}{6.000000}\selectfont\catcode`\^=\active\def^{\ifmmode\sp\else\^{}\fi}\catcode`\%=\active\def%{\%}5.9}}%
\end{pgfscope}%
\begin{pgfscope}%
\pgfpathrectangle{\pgfqpoint{2.259013in}{0.987099in}}{\pgfqpoint{0.979200in}{0.326250in}}%
\pgfusepath{clip}%
\pgfsetbuttcap%
\pgfsetroundjoin%
\pgfsetlinewidth{0.401500pt}%
\definecolor{currentstroke}{rgb}{0.690196,0.690196,0.690196}%
\pgfsetstrokecolor{currentstroke}%
\pgfsetstrokeopacity{0.250000}%
\pgfsetdash{{1.480000pt}{0.640000pt}}{0.000000pt}%
\pgfpathmoveto{\pgfqpoint{2.748613in}{0.987099in}}%
\pgfpathlineto{\pgfqpoint{2.748613in}{1.313349in}}%
\pgfusepath{stroke}%
\end{pgfscope}%
\begin{pgfscope}%
\pgfsetbuttcap%
\pgfsetroundjoin%
\definecolor{currentfill}{rgb}{0.000000,0.000000,0.000000}%
\pgfsetfillcolor{currentfill}%
\pgfsetlinewidth{0.803000pt}%
\definecolor{currentstroke}{rgb}{0.000000,0.000000,0.000000}%
\pgfsetstrokecolor{currentstroke}%
\pgfsetdash{}{0pt}%
\pgfsys@defobject{currentmarker}{\pgfqpoint{0.000000in}{0.000000in}}{\pgfqpoint{0.000000in}{0.027778in}}{%
\pgfpathmoveto{\pgfqpoint{0.000000in}{0.000000in}}%
\pgfpathlineto{\pgfqpoint{0.000000in}{0.027778in}}%
\pgfusepath{stroke,fill}%
}%
\begin{pgfscope}%
\pgfsys@transformshift{2.748613in}{0.987099in}%
\pgfsys@useobject{currentmarker}{}%
\end{pgfscope}%
\end{pgfscope}%
\begin{pgfscope}%
\definecolor{textcolor}{rgb}{0.000000,0.000000,0.000000}%
\pgfsetstrokecolor{textcolor}%
\pgfsetfillcolor{textcolor}%
\pgftext[x=2.748613in,y=0.973210in,,top]{\color{textcolor}{\rmfamily\fontsize{5.000000}{6.000000}\selectfont\catcode`\^=\active\def^{\ifmmode\sp\else\^{}\fi}\catcode`\%=\active\def%{\%}6.0}}%
\end{pgfscope}%
\begin{pgfscope}%
\pgfpathrectangle{\pgfqpoint{2.259013in}{0.987099in}}{\pgfqpoint{0.979200in}{0.326250in}}%
\pgfusepath{clip}%
\pgfsetbuttcap%
\pgfsetroundjoin%
\pgfsetlinewidth{0.401500pt}%
\definecolor{currentstroke}{rgb}{0.690196,0.690196,0.690196}%
\pgfsetstrokecolor{currentstroke}%
\pgfsetstrokeopacity{0.250000}%
\pgfsetdash{{1.480000pt}{0.640000pt}}{0.000000pt}%
\pgfpathmoveto{\pgfqpoint{2.993413in}{0.987099in}}%
\pgfpathlineto{\pgfqpoint{2.993413in}{1.313349in}}%
\pgfusepath{stroke}%
\end{pgfscope}%
\begin{pgfscope}%
\pgfsetbuttcap%
\pgfsetroundjoin%
\definecolor{currentfill}{rgb}{0.000000,0.000000,0.000000}%
\pgfsetfillcolor{currentfill}%
\pgfsetlinewidth{0.803000pt}%
\definecolor{currentstroke}{rgb}{0.000000,0.000000,0.000000}%
\pgfsetstrokecolor{currentstroke}%
\pgfsetdash{}{0pt}%
\pgfsys@defobject{currentmarker}{\pgfqpoint{0.000000in}{0.000000in}}{\pgfqpoint{0.000000in}{0.027778in}}{%
\pgfpathmoveto{\pgfqpoint{0.000000in}{0.000000in}}%
\pgfpathlineto{\pgfqpoint{0.000000in}{0.027778in}}%
\pgfusepath{stroke,fill}%
}%
\begin{pgfscope}%
\pgfsys@transformshift{2.993413in}{0.987099in}%
\pgfsys@useobject{currentmarker}{}%
\end{pgfscope}%
\end{pgfscope}%
\begin{pgfscope}%
\definecolor{textcolor}{rgb}{0.000000,0.000000,0.000000}%
\pgfsetstrokecolor{textcolor}%
\pgfsetfillcolor{textcolor}%
\pgftext[x=2.993413in,y=0.973210in,,top]{\color{textcolor}{\rmfamily\fontsize{5.000000}{6.000000}\selectfont\catcode`\^=\active\def^{\ifmmode\sp\else\^{}\fi}\catcode`\%=\active\def%{\%}6.0}}%
\end{pgfscope}%
\begin{pgfscope}%
\pgfpathrectangle{\pgfqpoint{2.259013in}{0.987099in}}{\pgfqpoint{0.979200in}{0.326250in}}%
\pgfusepath{clip}%
\pgfsetbuttcap%
\pgfsetroundjoin%
\pgfsetlinewidth{0.401500pt}%
\definecolor{currentstroke}{rgb}{0.690196,0.690196,0.690196}%
\pgfsetstrokecolor{currentstroke}%
\pgfsetstrokeopacity{0.250000}%
\pgfsetdash{{1.480000pt}{0.640000pt}}{0.000000pt}%
\pgfpathmoveto{\pgfqpoint{3.238213in}{0.987099in}}%
\pgfpathlineto{\pgfqpoint{3.238213in}{1.313349in}}%
\pgfusepath{stroke}%
\end{pgfscope}%
\begin{pgfscope}%
\pgfsetbuttcap%
\pgfsetroundjoin%
\definecolor{currentfill}{rgb}{0.000000,0.000000,0.000000}%
\pgfsetfillcolor{currentfill}%
\pgfsetlinewidth{0.803000pt}%
\definecolor{currentstroke}{rgb}{0.000000,0.000000,0.000000}%
\pgfsetstrokecolor{currentstroke}%
\pgfsetdash{}{0pt}%
\pgfsys@defobject{currentmarker}{\pgfqpoint{0.000000in}{0.000000in}}{\pgfqpoint{0.000000in}{0.027778in}}{%
\pgfpathmoveto{\pgfqpoint{0.000000in}{0.000000in}}%
\pgfpathlineto{\pgfqpoint{0.000000in}{0.027778in}}%
\pgfusepath{stroke,fill}%
}%
\begin{pgfscope}%
\pgfsys@transformshift{3.238213in}{0.987099in}%
\pgfsys@useobject{currentmarker}{}%
\end{pgfscope}%
\end{pgfscope}%
\begin{pgfscope}%
\definecolor{textcolor}{rgb}{0.000000,0.000000,0.000000}%
\pgfsetstrokecolor{textcolor}%
\pgfsetfillcolor{textcolor}%
\pgftext[x=3.238213in,y=0.973210in,,top]{\color{textcolor}{\rmfamily\fontsize{5.000000}{6.000000}\selectfont\catcode`\^=\active\def^{\ifmmode\sp\else\^{}\fi}\catcode`\%=\active\def%{\%}6.0}}%
\end{pgfscope}%
\begin{pgfscope}%
\pgfpathrectangle{\pgfqpoint{2.259013in}{0.987099in}}{\pgfqpoint{0.979200in}{0.326250in}}%
\pgfusepath{clip}%
\pgfsetbuttcap%
\pgfsetroundjoin%
\pgfsetlinewidth{0.401500pt}%
\definecolor{currentstroke}{rgb}{0.690196,0.690196,0.690196}%
\pgfsetstrokecolor{currentstroke}%
\pgfsetstrokeopacity{0.250000}%
\pgfsetdash{{1.480000pt}{0.640000pt}}{0.000000pt}%
\pgfpathmoveto{\pgfqpoint{2.259013in}{0.987099in}}%
\pgfpathlineto{\pgfqpoint{3.238213in}{0.987099in}}%
\pgfusepath{stroke}%
\end{pgfscope}%
\begin{pgfscope}%
\pgfsetbuttcap%
\pgfsetroundjoin%
\definecolor{currentfill}{rgb}{0.000000,0.000000,0.000000}%
\pgfsetfillcolor{currentfill}%
\pgfsetlinewidth{0.803000pt}%
\definecolor{currentstroke}{rgb}{0.000000,0.000000,0.000000}%
\pgfsetstrokecolor{currentstroke}%
\pgfsetdash{}{0pt}%
\pgfsys@defobject{currentmarker}{\pgfqpoint{0.000000in}{0.000000in}}{\pgfqpoint{0.027778in}{0.000000in}}{%
\pgfpathmoveto{\pgfqpoint{0.000000in}{0.000000in}}%
\pgfpathlineto{\pgfqpoint{0.027778in}{0.000000in}}%
\pgfusepath{stroke,fill}%
}%
\begin{pgfscope}%
\pgfsys@transformshift{2.259013in}{0.987099in}%
\pgfsys@useobject{currentmarker}{}%
\end{pgfscope}%
\end{pgfscope}%
\begin{pgfscope}%
\definecolor{textcolor}{rgb}{0.000000,0.000000,0.000000}%
\pgfsetstrokecolor{textcolor}%
\pgfsetfillcolor{textcolor}%
\pgftext[x=2.042575in, y=0.962986in, left, base]{\color{textcolor}{\rmfamily\fontsize{5.000000}{6.000000}\selectfont\catcode`\^=\active\def^{\ifmmode\sp\else\^{}\fi}\catcode`\%=\active\def%{\%}-2.00}}%
\end{pgfscope}%
\begin{pgfscope}%
\pgfpathrectangle{\pgfqpoint{2.259013in}{0.987099in}}{\pgfqpoint{0.979200in}{0.326250in}}%
\pgfusepath{clip}%
\pgfsetbuttcap%
\pgfsetroundjoin%
\pgfsetlinewidth{0.401500pt}%
\definecolor{currentstroke}{rgb}{0.690196,0.690196,0.690196}%
\pgfsetstrokecolor{currentstroke}%
\pgfsetstrokeopacity{0.250000}%
\pgfsetdash{{1.480000pt}{0.640000pt}}{0.000000pt}%
\pgfpathmoveto{\pgfqpoint{2.259013in}{1.150224in}}%
\pgfpathlineto{\pgfqpoint{3.238213in}{1.150224in}}%
\pgfusepath{stroke}%
\end{pgfscope}%
\begin{pgfscope}%
\pgfsetbuttcap%
\pgfsetroundjoin%
\definecolor{currentfill}{rgb}{0.000000,0.000000,0.000000}%
\pgfsetfillcolor{currentfill}%
\pgfsetlinewidth{0.803000pt}%
\definecolor{currentstroke}{rgb}{0.000000,0.000000,0.000000}%
\pgfsetstrokecolor{currentstroke}%
\pgfsetdash{}{0pt}%
\pgfsys@defobject{currentmarker}{\pgfqpoint{0.000000in}{0.000000in}}{\pgfqpoint{0.027778in}{0.000000in}}{%
\pgfpathmoveto{\pgfqpoint{0.000000in}{0.000000in}}%
\pgfpathlineto{\pgfqpoint{0.027778in}{0.000000in}}%
\pgfusepath{stroke,fill}%
}%
\begin{pgfscope}%
\pgfsys@transformshift{2.259013in}{1.150224in}%
\pgfsys@useobject{currentmarker}{}%
\end{pgfscope}%
\end{pgfscope}%
\begin{pgfscope}%
\definecolor{textcolor}{rgb}{0.000000,0.000000,0.000000}%
\pgfsetstrokecolor{textcolor}%
\pgfsetfillcolor{textcolor}%
\pgftext[x=2.075368in, y=1.126111in, left, base]{\color{textcolor}{\rmfamily\fontsize{5.000000}{6.000000}\selectfont\catcode`\^=\active\def^{\ifmmode\sp\else\^{}\fi}\catcode`\%=\active\def%{\%}0.00}}%
\end{pgfscope}%
\begin{pgfscope}%
\pgfpathrectangle{\pgfqpoint{2.259013in}{0.987099in}}{\pgfqpoint{0.979200in}{0.326250in}}%
\pgfusepath{clip}%
\pgfsetbuttcap%
\pgfsetroundjoin%
\pgfsetlinewidth{0.401500pt}%
\definecolor{currentstroke}{rgb}{0.690196,0.690196,0.690196}%
\pgfsetstrokecolor{currentstroke}%
\pgfsetstrokeopacity{0.250000}%
\pgfsetdash{{1.480000pt}{0.640000pt}}{0.000000pt}%
\pgfpathmoveto{\pgfqpoint{2.259013in}{1.313349in}}%
\pgfpathlineto{\pgfqpoint{3.238213in}{1.313349in}}%
\pgfusepath{stroke}%
\end{pgfscope}%
\begin{pgfscope}%
\pgfsetbuttcap%
\pgfsetroundjoin%
\definecolor{currentfill}{rgb}{0.000000,0.000000,0.000000}%
\pgfsetfillcolor{currentfill}%
\pgfsetlinewidth{0.803000pt}%
\definecolor{currentstroke}{rgb}{0.000000,0.000000,0.000000}%
\pgfsetstrokecolor{currentstroke}%
\pgfsetdash{}{0pt}%
\pgfsys@defobject{currentmarker}{\pgfqpoint{0.000000in}{0.000000in}}{\pgfqpoint{0.027778in}{0.000000in}}{%
\pgfpathmoveto{\pgfqpoint{0.000000in}{0.000000in}}%
\pgfpathlineto{\pgfqpoint{0.027778in}{0.000000in}}%
\pgfusepath{stroke,fill}%
}%
\begin{pgfscope}%
\pgfsys@transformshift{2.259013in}{1.313349in}%
\pgfsys@useobject{currentmarker}{}%
\end{pgfscope}%
\end{pgfscope}%
\begin{pgfscope}%
\definecolor{textcolor}{rgb}{0.000000,0.000000,0.000000}%
\pgfsetstrokecolor{textcolor}%
\pgfsetfillcolor{textcolor}%
\pgftext[x=2.075368in, y=1.289236in, left, base]{\color{textcolor}{\rmfamily\fontsize{5.000000}{6.000000}\selectfont\catcode`\^=\active\def^{\ifmmode\sp\else\^{}\fi}\catcode`\%=\active\def%{\%}2.00}}%
\end{pgfscope}%
\begin{pgfscope}%
\pgfpathrectangle{\pgfqpoint{2.259013in}{0.987099in}}{\pgfqpoint{0.979200in}{0.326250in}}%
\pgfusepath{clip}%
\pgfsetrectcap%
\pgfsetroundjoin%
\pgfsetlinewidth{1.505625pt}%
\definecolor{currentstroke}{rgb}{0.000000,0.619608,0.450980}%
\pgfsetstrokecolor{currentstroke}%
\pgfsetdash{}{0pt}%
\pgfpathmoveto{\pgfqpoint{2.255597in}{1.167833in}}%
\pgfpathlineto{\pgfqpoint{2.268418in}{1.167446in}}%
\pgfpathlineto{\pgfqpoint{2.268789in}{0.980154in}}%
\pgfpathmoveto{\pgfqpoint{2.518573in}{0.980154in}}%
\pgfpathlineto{\pgfqpoint{2.519071in}{1.168611in}}%
\pgfpathlineto{\pgfqpoint{2.536148in}{1.168467in}}%
\pgfpathlineto{\pgfqpoint{2.570420in}{1.168468in}}%
\pgfpathlineto{\pgfqpoint{2.604692in}{1.168244in}}%
\pgfpathlineto{\pgfqpoint{2.638964in}{1.168235in}}%
\pgfpathlineto{\pgfqpoint{2.673236in}{1.167994in}}%
\pgfpathlineto{\pgfqpoint{2.707508in}{1.167984in}}%
\pgfpathlineto{\pgfqpoint{2.741780in}{1.167712in}}%
\pgfpathlineto{\pgfqpoint{2.756533in}{1.167721in}}%
\pgfpathlineto{\pgfqpoint{2.758018in}{1.167446in}}%
\pgfpathlineto{\pgfqpoint{2.758389in}{0.980154in}}%
\pgfpathmoveto{\pgfqpoint{3.008173in}{0.980154in}}%
\pgfpathlineto{\pgfqpoint{3.008671in}{1.168611in}}%
\pgfpathlineto{\pgfqpoint{3.025748in}{1.168467in}}%
\pgfpathlineto{\pgfqpoint{3.060020in}{1.168468in}}%
\pgfpathlineto{\pgfqpoint{3.094292in}{1.168244in}}%
\pgfpathlineto{\pgfqpoint{3.128564in}{1.168235in}}%
\pgfpathlineto{\pgfqpoint{3.162836in}{1.167994in}}%
\pgfpathlineto{\pgfqpoint{3.197108in}{1.167984in}}%
\pgfpathlineto{\pgfqpoint{3.231380in}{1.167712in}}%
\pgfpathlineto{\pgfqpoint{3.238213in}{1.167681in}}%
\pgfpathlineto{\pgfqpoint{3.238213in}{1.167681in}}%
\pgfusepath{stroke}%
\end{pgfscope}%
\begin{pgfscope}%
\pgfsetrectcap%
\pgfsetmiterjoin%
\pgfsetlinewidth{0.803000pt}%
\definecolor{currentstroke}{rgb}{0.000000,0.000000,0.000000}%
\pgfsetstrokecolor{currentstroke}%
\pgfsetdash{}{0pt}%
\pgfpathmoveto{\pgfqpoint{2.259013in}{0.987099in}}%
\pgfpathlineto{\pgfqpoint{2.259013in}{1.313349in}}%
\pgfusepath{stroke}%
\end{pgfscope}%
\begin{pgfscope}%
\pgfsetrectcap%
\pgfsetmiterjoin%
\pgfsetlinewidth{0.803000pt}%
\definecolor{currentstroke}{rgb}{0.000000,0.000000,0.000000}%
\pgfsetstrokecolor{currentstroke}%
\pgfsetdash{}{0pt}%
\pgfpathmoveto{\pgfqpoint{3.238213in}{0.987099in}}%
\pgfpathlineto{\pgfqpoint{3.238213in}{1.313349in}}%
\pgfusepath{stroke}%
\end{pgfscope}%
\begin{pgfscope}%
\pgfsetrectcap%
\pgfsetmiterjoin%
\pgfsetlinewidth{0.803000pt}%
\definecolor{currentstroke}{rgb}{0.000000,0.000000,0.000000}%
\pgfsetstrokecolor{currentstroke}%
\pgfsetdash{}{0pt}%
\pgfpathmoveto{\pgfqpoint{2.259013in}{0.987099in}}%
\pgfpathlineto{\pgfqpoint{3.238213in}{0.987099in}}%
\pgfusepath{stroke}%
\end{pgfscope}%
\begin{pgfscope}%
\pgfsetrectcap%
\pgfsetmiterjoin%
\pgfsetlinewidth{0.803000pt}%
\definecolor{currentstroke}{rgb}{0.000000,0.000000,0.000000}%
\pgfsetstrokecolor{currentstroke}%
\pgfsetdash{}{0pt}%
\pgfpathmoveto{\pgfqpoint{2.259013in}{1.313349in}}%
\pgfpathlineto{\pgfqpoint{3.238213in}{1.313349in}}%
\pgfusepath{stroke}%
\end{pgfscope}%
\end{pgfpicture}%
\makeatother%
\endgroup%

	\caption{(a) tensão de saída, (b) corrente do indutor e (c) tensão no diodo do circuito Buck, com diodo ideal e carga de $5$ A.}
	\label{fig:ex4_2}
\end{figure}

A Figura~\ref{fig:ex4_3} mostra os resultados de simulação obtidos para a carga com corrente de $1$ A.
\begin{figure}[htb!]
	\centering
	%% Creator: Matplotlib, PGF backend
%%
%% To include the figure in your LaTeX document, write
%%   \input{<filename>.pgf}
%%
%% Make sure the required packages are loaded in your preamble
%%   \usepackage{pgf}
%%
%% Also ensure that all the required font packages are loaded; for instance,
%% the lmodern package is sometimes necessary when using math font.
%%   \usepackage{lmodern}
%%
%% Figures using additional raster images can only be included by \input if
%% they are in the same directory as the main LaTeX file. For loading figures
%% from other directories you can use the `import` package
%%   \usepackage{import}
%%
%% and then include the figures with
%%   \import{<path to file>}{<filename>.pgf}
%%
%% Matplotlib used the following preamble
%%   \def\mathdefault#1{#1}
%%   \everymath=\expandafter{\the\everymath\displaystyle}
%%   \IfFileExists{scrextend.sty}{
%%     \usepackage[fontsize=8.000000pt]{scrextend}
%%   }{
%%     \renewcommand{\normalsize}{\fontsize{8.000000}{9.600000}\selectfont}
%%     \normalsize
%%   }
%%   
%%           \usepackage{amsmath}
%%           \usepackage{amssymb}
%%           \usepackage{mathtools}
%%           \usepackage{dsfont}
%%           \providecommand{\mathdefault}[1]{#1}
%%       
%%   \makeatletter\@ifpackageloaded{underscore}{}{\usepackage[strings]{underscore}}\makeatother
%%
\begingroup%
\makeatletter%
\begin{pgfpicture}%
\pgfpathrectangle{\pgfpointorigin}{\pgfqpoint{3.399460in}{3.646235in}}%
\pgfusepath{use as bounding box, clip}%
\begin{pgfscope}%
\pgfsetbuttcap%
\pgfsetmiterjoin%
\definecolor{currentfill}{rgb}{1.000000,1.000000,1.000000}%
\pgfsetfillcolor{currentfill}%
\pgfsetlinewidth{0.000000pt}%
\definecolor{currentstroke}{rgb}{1.000000,1.000000,1.000000}%
\pgfsetstrokecolor{currentstroke}%
\pgfsetdash{}{0pt}%
\pgfpathmoveto{\pgfqpoint{0.000000in}{0.000000in}}%
\pgfpathlineto{\pgfqpoint{3.399460in}{0.000000in}}%
\pgfpathlineto{\pgfqpoint{3.399460in}{3.646235in}}%
\pgfpathlineto{\pgfqpoint{0.000000in}{3.646235in}}%
\pgfpathlineto{\pgfqpoint{0.000000in}{0.000000in}}%
\pgfpathclose%
\pgfusepath{fill}%
\end{pgfscope}%
\begin{pgfscope}%
\pgfsetbuttcap%
\pgfsetmiterjoin%
\definecolor{currentfill}{rgb}{1.000000,1.000000,1.000000}%
\pgfsetfillcolor{currentfill}%
\pgfsetlinewidth{0.000000pt}%
\definecolor{currentstroke}{rgb}{0.000000,0.000000,0.000000}%
\pgfsetstrokecolor{currentstroke}%
\pgfsetstrokeopacity{0.000000}%
\pgfsetdash{}{0pt}%
\pgfpathmoveto{\pgfqpoint{0.573768in}{2.601404in}}%
\pgfpathlineto{\pgfqpoint{3.293768in}{2.601404in}}%
\pgfpathlineto{\pgfqpoint{3.293768in}{3.507654in}}%
\pgfpathlineto{\pgfqpoint{0.573768in}{3.507654in}}%
\pgfpathlineto{\pgfqpoint{0.573768in}{2.601404in}}%
\pgfpathclose%
\pgfusepath{fill}%
\end{pgfscope}%
\begin{pgfscope}%
\pgfpathrectangle{\pgfqpoint{0.573768in}{2.601404in}}{\pgfqpoint{2.720000in}{0.906250in}}%
\pgfusepath{clip}%
\pgfsetbuttcap%
\pgfsetroundjoin%
\pgfsetlinewidth{0.501875pt}%
\definecolor{currentstroke}{rgb}{0.690196,0.690196,0.690196}%
\pgfsetstrokecolor{currentstroke}%
\pgfsetstrokeopacity{0.250000}%
\pgfsetdash{{1.850000pt}{0.800000pt}}{0.000000pt}%
\pgfpathmoveto{\pgfqpoint{0.697405in}{2.601404in}}%
\pgfpathlineto{\pgfqpoint{0.697405in}{3.507654in}}%
\pgfusepath{stroke}%
\end{pgfscope}%
\begin{pgfscope}%
\pgfsetbuttcap%
\pgfsetroundjoin%
\definecolor{currentfill}{rgb}{0.000000,0.000000,0.000000}%
\pgfsetfillcolor{currentfill}%
\pgfsetlinewidth{0.803000pt}%
\definecolor{currentstroke}{rgb}{0.000000,0.000000,0.000000}%
\pgfsetstrokecolor{currentstroke}%
\pgfsetdash{}{0pt}%
\pgfsys@defobject{currentmarker}{\pgfqpoint{0.000000in}{-0.048611in}}{\pgfqpoint{0.000000in}{0.000000in}}{%
\pgfpathmoveto{\pgfqpoint{0.000000in}{0.000000in}}%
\pgfpathlineto{\pgfqpoint{0.000000in}{-0.048611in}}%
\pgfusepath{stroke,fill}%
}%
\begin{pgfscope}%
\pgfsys@transformshift{0.697405in}{2.601404in}%
\pgfsys@useobject{currentmarker}{}%
\end{pgfscope}%
\end{pgfscope}%
\begin{pgfscope}%
\pgfpathrectangle{\pgfqpoint{0.573768in}{2.601404in}}{\pgfqpoint{2.720000in}{0.906250in}}%
\pgfusepath{clip}%
\pgfsetbuttcap%
\pgfsetroundjoin%
\pgfsetlinewidth{0.501875pt}%
\definecolor{currentstroke}{rgb}{0.690196,0.690196,0.690196}%
\pgfsetstrokecolor{currentstroke}%
\pgfsetstrokeopacity{0.250000}%
\pgfsetdash{{1.850000pt}{0.800000pt}}{0.000000pt}%
\pgfpathmoveto{\pgfqpoint{1.109526in}{2.601404in}}%
\pgfpathlineto{\pgfqpoint{1.109526in}{3.507654in}}%
\pgfusepath{stroke}%
\end{pgfscope}%
\begin{pgfscope}%
\pgfsetbuttcap%
\pgfsetroundjoin%
\definecolor{currentfill}{rgb}{0.000000,0.000000,0.000000}%
\pgfsetfillcolor{currentfill}%
\pgfsetlinewidth{0.803000pt}%
\definecolor{currentstroke}{rgb}{0.000000,0.000000,0.000000}%
\pgfsetstrokecolor{currentstroke}%
\pgfsetdash{}{0pt}%
\pgfsys@defobject{currentmarker}{\pgfqpoint{0.000000in}{-0.048611in}}{\pgfqpoint{0.000000in}{0.000000in}}{%
\pgfpathmoveto{\pgfqpoint{0.000000in}{0.000000in}}%
\pgfpathlineto{\pgfqpoint{0.000000in}{-0.048611in}}%
\pgfusepath{stroke,fill}%
}%
\begin{pgfscope}%
\pgfsys@transformshift{1.109526in}{2.601404in}%
\pgfsys@useobject{currentmarker}{}%
\end{pgfscope}%
\end{pgfscope}%
\begin{pgfscope}%
\pgfpathrectangle{\pgfqpoint{0.573768in}{2.601404in}}{\pgfqpoint{2.720000in}{0.906250in}}%
\pgfusepath{clip}%
\pgfsetbuttcap%
\pgfsetroundjoin%
\pgfsetlinewidth{0.501875pt}%
\definecolor{currentstroke}{rgb}{0.690196,0.690196,0.690196}%
\pgfsetstrokecolor{currentstroke}%
\pgfsetstrokeopacity{0.250000}%
\pgfsetdash{{1.850000pt}{0.800000pt}}{0.000000pt}%
\pgfpathmoveto{\pgfqpoint{1.521647in}{2.601404in}}%
\pgfpathlineto{\pgfqpoint{1.521647in}{3.507654in}}%
\pgfusepath{stroke}%
\end{pgfscope}%
\begin{pgfscope}%
\pgfsetbuttcap%
\pgfsetroundjoin%
\definecolor{currentfill}{rgb}{0.000000,0.000000,0.000000}%
\pgfsetfillcolor{currentfill}%
\pgfsetlinewidth{0.803000pt}%
\definecolor{currentstroke}{rgb}{0.000000,0.000000,0.000000}%
\pgfsetstrokecolor{currentstroke}%
\pgfsetdash{}{0pt}%
\pgfsys@defobject{currentmarker}{\pgfqpoint{0.000000in}{-0.048611in}}{\pgfqpoint{0.000000in}{0.000000in}}{%
\pgfpathmoveto{\pgfqpoint{0.000000in}{0.000000in}}%
\pgfpathlineto{\pgfqpoint{0.000000in}{-0.048611in}}%
\pgfusepath{stroke,fill}%
}%
\begin{pgfscope}%
\pgfsys@transformshift{1.521647in}{2.601404in}%
\pgfsys@useobject{currentmarker}{}%
\end{pgfscope}%
\end{pgfscope}%
\begin{pgfscope}%
\pgfpathrectangle{\pgfqpoint{0.573768in}{2.601404in}}{\pgfqpoint{2.720000in}{0.906250in}}%
\pgfusepath{clip}%
\pgfsetbuttcap%
\pgfsetroundjoin%
\pgfsetlinewidth{0.501875pt}%
\definecolor{currentstroke}{rgb}{0.690196,0.690196,0.690196}%
\pgfsetstrokecolor{currentstroke}%
\pgfsetstrokeopacity{0.250000}%
\pgfsetdash{{1.850000pt}{0.800000pt}}{0.000000pt}%
\pgfpathmoveto{\pgfqpoint{1.933768in}{2.601404in}}%
\pgfpathlineto{\pgfqpoint{1.933768in}{3.507654in}}%
\pgfusepath{stroke}%
\end{pgfscope}%
\begin{pgfscope}%
\pgfsetbuttcap%
\pgfsetroundjoin%
\definecolor{currentfill}{rgb}{0.000000,0.000000,0.000000}%
\pgfsetfillcolor{currentfill}%
\pgfsetlinewidth{0.803000pt}%
\definecolor{currentstroke}{rgb}{0.000000,0.000000,0.000000}%
\pgfsetstrokecolor{currentstroke}%
\pgfsetdash{}{0pt}%
\pgfsys@defobject{currentmarker}{\pgfqpoint{0.000000in}{-0.048611in}}{\pgfqpoint{0.000000in}{0.000000in}}{%
\pgfpathmoveto{\pgfqpoint{0.000000in}{0.000000in}}%
\pgfpathlineto{\pgfqpoint{0.000000in}{-0.048611in}}%
\pgfusepath{stroke,fill}%
}%
\begin{pgfscope}%
\pgfsys@transformshift{1.933768in}{2.601404in}%
\pgfsys@useobject{currentmarker}{}%
\end{pgfscope}%
\end{pgfscope}%
\begin{pgfscope}%
\pgfpathrectangle{\pgfqpoint{0.573768in}{2.601404in}}{\pgfqpoint{2.720000in}{0.906250in}}%
\pgfusepath{clip}%
\pgfsetbuttcap%
\pgfsetroundjoin%
\pgfsetlinewidth{0.501875pt}%
\definecolor{currentstroke}{rgb}{0.690196,0.690196,0.690196}%
\pgfsetstrokecolor{currentstroke}%
\pgfsetstrokeopacity{0.250000}%
\pgfsetdash{{1.850000pt}{0.800000pt}}{0.000000pt}%
\pgfpathmoveto{\pgfqpoint{2.345890in}{2.601404in}}%
\pgfpathlineto{\pgfqpoint{2.345890in}{3.507654in}}%
\pgfusepath{stroke}%
\end{pgfscope}%
\begin{pgfscope}%
\pgfsetbuttcap%
\pgfsetroundjoin%
\definecolor{currentfill}{rgb}{0.000000,0.000000,0.000000}%
\pgfsetfillcolor{currentfill}%
\pgfsetlinewidth{0.803000pt}%
\definecolor{currentstroke}{rgb}{0.000000,0.000000,0.000000}%
\pgfsetstrokecolor{currentstroke}%
\pgfsetdash{}{0pt}%
\pgfsys@defobject{currentmarker}{\pgfqpoint{0.000000in}{-0.048611in}}{\pgfqpoint{0.000000in}{0.000000in}}{%
\pgfpathmoveto{\pgfqpoint{0.000000in}{0.000000in}}%
\pgfpathlineto{\pgfqpoint{0.000000in}{-0.048611in}}%
\pgfusepath{stroke,fill}%
}%
\begin{pgfscope}%
\pgfsys@transformshift{2.345890in}{2.601404in}%
\pgfsys@useobject{currentmarker}{}%
\end{pgfscope}%
\end{pgfscope}%
\begin{pgfscope}%
\pgfpathrectangle{\pgfqpoint{0.573768in}{2.601404in}}{\pgfqpoint{2.720000in}{0.906250in}}%
\pgfusepath{clip}%
\pgfsetbuttcap%
\pgfsetroundjoin%
\pgfsetlinewidth{0.501875pt}%
\definecolor{currentstroke}{rgb}{0.690196,0.690196,0.690196}%
\pgfsetstrokecolor{currentstroke}%
\pgfsetstrokeopacity{0.250000}%
\pgfsetdash{{1.850000pt}{0.800000pt}}{0.000000pt}%
\pgfpathmoveto{\pgfqpoint{2.758011in}{2.601404in}}%
\pgfpathlineto{\pgfqpoint{2.758011in}{3.507654in}}%
\pgfusepath{stroke}%
\end{pgfscope}%
\begin{pgfscope}%
\pgfsetbuttcap%
\pgfsetroundjoin%
\definecolor{currentfill}{rgb}{0.000000,0.000000,0.000000}%
\pgfsetfillcolor{currentfill}%
\pgfsetlinewidth{0.803000pt}%
\definecolor{currentstroke}{rgb}{0.000000,0.000000,0.000000}%
\pgfsetstrokecolor{currentstroke}%
\pgfsetdash{}{0pt}%
\pgfsys@defobject{currentmarker}{\pgfqpoint{0.000000in}{-0.048611in}}{\pgfqpoint{0.000000in}{0.000000in}}{%
\pgfpathmoveto{\pgfqpoint{0.000000in}{0.000000in}}%
\pgfpathlineto{\pgfqpoint{0.000000in}{-0.048611in}}%
\pgfusepath{stroke,fill}%
}%
\begin{pgfscope}%
\pgfsys@transformshift{2.758011in}{2.601404in}%
\pgfsys@useobject{currentmarker}{}%
\end{pgfscope}%
\end{pgfscope}%
\begin{pgfscope}%
\pgfpathrectangle{\pgfqpoint{0.573768in}{2.601404in}}{\pgfqpoint{2.720000in}{0.906250in}}%
\pgfusepath{clip}%
\pgfsetbuttcap%
\pgfsetroundjoin%
\pgfsetlinewidth{0.501875pt}%
\definecolor{currentstroke}{rgb}{0.690196,0.690196,0.690196}%
\pgfsetstrokecolor{currentstroke}%
\pgfsetstrokeopacity{0.250000}%
\pgfsetdash{{1.850000pt}{0.800000pt}}{0.000000pt}%
\pgfpathmoveto{\pgfqpoint{3.170132in}{2.601404in}}%
\pgfpathlineto{\pgfqpoint{3.170132in}{3.507654in}}%
\pgfusepath{stroke}%
\end{pgfscope}%
\begin{pgfscope}%
\pgfsetbuttcap%
\pgfsetroundjoin%
\definecolor{currentfill}{rgb}{0.000000,0.000000,0.000000}%
\pgfsetfillcolor{currentfill}%
\pgfsetlinewidth{0.803000pt}%
\definecolor{currentstroke}{rgb}{0.000000,0.000000,0.000000}%
\pgfsetstrokecolor{currentstroke}%
\pgfsetdash{}{0pt}%
\pgfsys@defobject{currentmarker}{\pgfqpoint{0.000000in}{-0.048611in}}{\pgfqpoint{0.000000in}{0.000000in}}{%
\pgfpathmoveto{\pgfqpoint{0.000000in}{0.000000in}}%
\pgfpathlineto{\pgfqpoint{0.000000in}{-0.048611in}}%
\pgfusepath{stroke,fill}%
}%
\begin{pgfscope}%
\pgfsys@transformshift{3.170132in}{2.601404in}%
\pgfsys@useobject{currentmarker}{}%
\end{pgfscope}%
\end{pgfscope}%
\begin{pgfscope}%
\pgfpathrectangle{\pgfqpoint{0.573768in}{2.601404in}}{\pgfqpoint{2.720000in}{0.906250in}}%
\pgfusepath{clip}%
\pgfsetbuttcap%
\pgfsetroundjoin%
\pgfsetlinewidth{0.501875pt}%
\definecolor{currentstroke}{rgb}{0.690196,0.690196,0.690196}%
\pgfsetstrokecolor{currentstroke}%
\pgfsetstrokeopacity{0.250000}%
\pgfsetdash{{1.850000pt}{0.800000pt}}{0.000000pt}%
\pgfpathmoveto{\pgfqpoint{0.573768in}{2.601404in}}%
\pgfpathlineto{\pgfqpoint{3.293768in}{2.601404in}}%
\pgfusepath{stroke}%
\end{pgfscope}%
\begin{pgfscope}%
\pgfsetbuttcap%
\pgfsetroundjoin%
\definecolor{currentfill}{rgb}{0.000000,0.000000,0.000000}%
\pgfsetfillcolor{currentfill}%
\pgfsetlinewidth{0.803000pt}%
\definecolor{currentstroke}{rgb}{0.000000,0.000000,0.000000}%
\pgfsetstrokecolor{currentstroke}%
\pgfsetdash{}{0pt}%
\pgfsys@defobject{currentmarker}{\pgfqpoint{-0.048611in}{0.000000in}}{\pgfqpoint{-0.000000in}{0.000000in}}{%
\pgfpathmoveto{\pgfqpoint{-0.000000in}{0.000000in}}%
\pgfpathlineto{\pgfqpoint{-0.048611in}{0.000000in}}%
\pgfusepath{stroke,fill}%
}%
\begin{pgfscope}%
\pgfsys@transformshift{0.573768in}{2.601404in}%
\pgfsys@useobject{currentmarker}{}%
\end{pgfscope}%
\end{pgfscope}%
\begin{pgfscope}%
\definecolor{textcolor}{rgb}{0.000000,0.000000,0.000000}%
\pgfsetstrokecolor{textcolor}%
\pgfsetfillcolor{textcolor}%
\pgftext[x=0.417518in, y=2.562824in, left, base]{\color{textcolor}{\rmfamily\fontsize{8.000000}{9.600000}\selectfont\catcode`\^=\active\def^{\ifmmode\sp\else\^{}\fi}\catcode`\%=\active\def%{\%}$\mathdefault{0}$}}%
\end{pgfscope}%
\begin{pgfscope}%
\pgfpathrectangle{\pgfqpoint{0.573768in}{2.601404in}}{\pgfqpoint{2.720000in}{0.906250in}}%
\pgfusepath{clip}%
\pgfsetbuttcap%
\pgfsetroundjoin%
\pgfsetlinewidth{0.501875pt}%
\definecolor{currentstroke}{rgb}{0.690196,0.690196,0.690196}%
\pgfsetstrokecolor{currentstroke}%
\pgfsetstrokeopacity{0.250000}%
\pgfsetdash{{1.850000pt}{0.800000pt}}{0.000000pt}%
\pgfpathmoveto{\pgfqpoint{0.573768in}{2.827967in}}%
\pgfpathlineto{\pgfqpoint{3.293768in}{2.827967in}}%
\pgfusepath{stroke}%
\end{pgfscope}%
\begin{pgfscope}%
\pgfsetbuttcap%
\pgfsetroundjoin%
\definecolor{currentfill}{rgb}{0.000000,0.000000,0.000000}%
\pgfsetfillcolor{currentfill}%
\pgfsetlinewidth{0.803000pt}%
\definecolor{currentstroke}{rgb}{0.000000,0.000000,0.000000}%
\pgfsetstrokecolor{currentstroke}%
\pgfsetdash{}{0pt}%
\pgfsys@defobject{currentmarker}{\pgfqpoint{-0.048611in}{0.000000in}}{\pgfqpoint{-0.000000in}{0.000000in}}{%
\pgfpathmoveto{\pgfqpoint{-0.000000in}{0.000000in}}%
\pgfpathlineto{\pgfqpoint{-0.048611in}{0.000000in}}%
\pgfusepath{stroke,fill}%
}%
\begin{pgfscope}%
\pgfsys@transformshift{0.573768in}{2.827967in}%
\pgfsys@useobject{currentmarker}{}%
\end{pgfscope}%
\end{pgfscope}%
\begin{pgfscope}%
\definecolor{textcolor}{rgb}{0.000000,0.000000,0.000000}%
\pgfsetstrokecolor{textcolor}%
\pgfsetfillcolor{textcolor}%
\pgftext[x=0.417518in, y=2.789387in, left, base]{\color{textcolor}{\rmfamily\fontsize{8.000000}{9.600000}\selectfont\catcode`\^=\active\def^{\ifmmode\sp\else\^{}\fi}\catcode`\%=\active\def%{\%}$\mathdefault{5}$}}%
\end{pgfscope}%
\begin{pgfscope}%
\pgfpathrectangle{\pgfqpoint{0.573768in}{2.601404in}}{\pgfqpoint{2.720000in}{0.906250in}}%
\pgfusepath{clip}%
\pgfsetbuttcap%
\pgfsetroundjoin%
\pgfsetlinewidth{0.501875pt}%
\definecolor{currentstroke}{rgb}{0.690196,0.690196,0.690196}%
\pgfsetstrokecolor{currentstroke}%
\pgfsetstrokeopacity{0.250000}%
\pgfsetdash{{1.850000pt}{0.800000pt}}{0.000000pt}%
\pgfpathmoveto{\pgfqpoint{0.573768in}{3.054529in}}%
\pgfpathlineto{\pgfqpoint{3.293768in}{3.054529in}}%
\pgfusepath{stroke}%
\end{pgfscope}%
\begin{pgfscope}%
\pgfsetbuttcap%
\pgfsetroundjoin%
\definecolor{currentfill}{rgb}{0.000000,0.000000,0.000000}%
\pgfsetfillcolor{currentfill}%
\pgfsetlinewidth{0.803000pt}%
\definecolor{currentstroke}{rgb}{0.000000,0.000000,0.000000}%
\pgfsetstrokecolor{currentstroke}%
\pgfsetdash{}{0pt}%
\pgfsys@defobject{currentmarker}{\pgfqpoint{-0.048611in}{0.000000in}}{\pgfqpoint{-0.000000in}{0.000000in}}{%
\pgfpathmoveto{\pgfqpoint{-0.000000in}{0.000000in}}%
\pgfpathlineto{\pgfqpoint{-0.048611in}{0.000000in}}%
\pgfusepath{stroke,fill}%
}%
\begin{pgfscope}%
\pgfsys@transformshift{0.573768in}{3.054529in}%
\pgfsys@useobject{currentmarker}{}%
\end{pgfscope}%
\end{pgfscope}%
\begin{pgfscope}%
\definecolor{textcolor}{rgb}{0.000000,0.000000,0.000000}%
\pgfsetstrokecolor{textcolor}%
\pgfsetfillcolor{textcolor}%
\pgftext[x=0.358489in, y=3.015949in, left, base]{\color{textcolor}{\rmfamily\fontsize{8.000000}{9.600000}\selectfont\catcode`\^=\active\def^{\ifmmode\sp\else\^{}\fi}\catcode`\%=\active\def%{\%}$\mathdefault{10}$}}%
\end{pgfscope}%
\begin{pgfscope}%
\pgfpathrectangle{\pgfqpoint{0.573768in}{2.601404in}}{\pgfqpoint{2.720000in}{0.906250in}}%
\pgfusepath{clip}%
\pgfsetbuttcap%
\pgfsetroundjoin%
\pgfsetlinewidth{0.501875pt}%
\definecolor{currentstroke}{rgb}{0.690196,0.690196,0.690196}%
\pgfsetstrokecolor{currentstroke}%
\pgfsetstrokeopacity{0.250000}%
\pgfsetdash{{1.850000pt}{0.800000pt}}{0.000000pt}%
\pgfpathmoveto{\pgfqpoint{0.573768in}{3.281092in}}%
\pgfpathlineto{\pgfqpoint{3.293768in}{3.281092in}}%
\pgfusepath{stroke}%
\end{pgfscope}%
\begin{pgfscope}%
\pgfsetbuttcap%
\pgfsetroundjoin%
\definecolor{currentfill}{rgb}{0.000000,0.000000,0.000000}%
\pgfsetfillcolor{currentfill}%
\pgfsetlinewidth{0.803000pt}%
\definecolor{currentstroke}{rgb}{0.000000,0.000000,0.000000}%
\pgfsetstrokecolor{currentstroke}%
\pgfsetdash{}{0pt}%
\pgfsys@defobject{currentmarker}{\pgfqpoint{-0.048611in}{0.000000in}}{\pgfqpoint{-0.000000in}{0.000000in}}{%
\pgfpathmoveto{\pgfqpoint{-0.000000in}{0.000000in}}%
\pgfpathlineto{\pgfqpoint{-0.048611in}{0.000000in}}%
\pgfusepath{stroke,fill}%
}%
\begin{pgfscope}%
\pgfsys@transformshift{0.573768in}{3.281092in}%
\pgfsys@useobject{currentmarker}{}%
\end{pgfscope}%
\end{pgfscope}%
\begin{pgfscope}%
\definecolor{textcolor}{rgb}{0.000000,0.000000,0.000000}%
\pgfsetstrokecolor{textcolor}%
\pgfsetfillcolor{textcolor}%
\pgftext[x=0.358489in, y=3.242512in, left, base]{\color{textcolor}{\rmfamily\fontsize{8.000000}{9.600000}\selectfont\catcode`\^=\active\def^{\ifmmode\sp\else\^{}\fi}\catcode`\%=\active\def%{\%}$\mathdefault{15}$}}%
\end{pgfscope}%
\begin{pgfscope}%
\pgfpathrectangle{\pgfqpoint{0.573768in}{2.601404in}}{\pgfqpoint{2.720000in}{0.906250in}}%
\pgfusepath{clip}%
\pgfsetbuttcap%
\pgfsetroundjoin%
\pgfsetlinewidth{0.501875pt}%
\definecolor{currentstroke}{rgb}{0.690196,0.690196,0.690196}%
\pgfsetstrokecolor{currentstroke}%
\pgfsetstrokeopacity{0.250000}%
\pgfsetdash{{1.850000pt}{0.800000pt}}{0.000000pt}%
\pgfpathmoveto{\pgfqpoint{0.573768in}{3.507654in}}%
\pgfpathlineto{\pgfqpoint{3.293768in}{3.507654in}}%
\pgfusepath{stroke}%
\end{pgfscope}%
\begin{pgfscope}%
\pgfsetbuttcap%
\pgfsetroundjoin%
\definecolor{currentfill}{rgb}{0.000000,0.000000,0.000000}%
\pgfsetfillcolor{currentfill}%
\pgfsetlinewidth{0.803000pt}%
\definecolor{currentstroke}{rgb}{0.000000,0.000000,0.000000}%
\pgfsetstrokecolor{currentstroke}%
\pgfsetdash{}{0pt}%
\pgfsys@defobject{currentmarker}{\pgfqpoint{-0.048611in}{0.000000in}}{\pgfqpoint{-0.000000in}{0.000000in}}{%
\pgfpathmoveto{\pgfqpoint{-0.000000in}{0.000000in}}%
\pgfpathlineto{\pgfqpoint{-0.048611in}{0.000000in}}%
\pgfusepath{stroke,fill}%
}%
\begin{pgfscope}%
\pgfsys@transformshift{0.573768in}{3.507654in}%
\pgfsys@useobject{currentmarker}{}%
\end{pgfscope}%
\end{pgfscope}%
\begin{pgfscope}%
\definecolor{textcolor}{rgb}{0.000000,0.000000,0.000000}%
\pgfsetstrokecolor{textcolor}%
\pgfsetfillcolor{textcolor}%
\pgftext[x=0.358489in, y=3.469074in, left, base]{\color{textcolor}{\rmfamily\fontsize{8.000000}{9.600000}\selectfont\catcode`\^=\active\def^{\ifmmode\sp\else\^{}\fi}\catcode`\%=\active\def%{\%}$\mathdefault{20}$}}%
\end{pgfscope}%
\begin{pgfscope}%
\definecolor{textcolor}{rgb}{0.000000,0.000000,0.000000}%
\pgfsetstrokecolor{textcolor}%
\pgfsetfillcolor{textcolor}%
\pgftext[x=0.302933in,y=3.054529in,,bottom,rotate=90.000000]{\color{textcolor}{\rmfamily\fontsize{8.000000}{9.600000}\selectfont\catcode`\^=\active\def^{\ifmmode\sp\else\^{}\fi}\catcode`\%=\active\def%{\%}$v_o$ (V)}}%
\end{pgfscope}%
\begin{pgfscope}%
\pgfpathrectangle{\pgfqpoint{0.573768in}{2.601404in}}{\pgfqpoint{2.720000in}{0.906250in}}%
\pgfusepath{clip}%
\pgfsetrectcap%
\pgfsetroundjoin%
\pgfsetlinewidth{1.505625pt}%
\definecolor{currentstroke}{rgb}{0.000000,0.447059,0.698039}%
\pgfsetstrokecolor{currentstroke}%
\pgfsetdash{}{0pt}%
\pgfpathmoveto{\pgfqpoint{0.697405in}{2.601404in}}%
\pgfpathlineto{\pgfqpoint{0.699438in}{2.602153in}}%
\pgfpathlineto{\pgfqpoint{0.701416in}{2.604954in}}%
\pgfpathlineto{\pgfqpoint{0.704054in}{2.611890in}}%
\pgfpathlineto{\pgfqpoint{0.707841in}{2.627309in}}%
\pgfpathlineto{\pgfqpoint{0.718791in}{2.677610in}}%
\pgfpathlineto{\pgfqpoint{0.723477in}{2.711765in}}%
\pgfpathlineto{\pgfqpoint{0.739470in}{2.844160in}}%
\pgfpathlineto{\pgfqpoint{0.749312in}{2.937746in}}%
\pgfpathlineto{\pgfqpoint{0.756615in}{3.014541in}}%
\pgfpathlineto{\pgfqpoint{0.764316in}{3.090671in}}%
\pgfpathlineto{\pgfqpoint{0.772618in}{3.172585in}}%
\pgfpathlineto{\pgfqpoint{0.779047in}{3.231215in}}%
\pgfpathlineto{\pgfqpoint{0.795243in}{3.352252in}}%
\pgfpathlineto{\pgfqpoint{0.806448in}{3.407679in}}%
\pgfpathlineto{\pgfqpoint{0.810240in}{3.421122in}}%
\pgfpathlineto{\pgfqpoint{0.813233in}{3.426717in}}%
\pgfpathlineto{\pgfqpoint{0.820702in}{3.436596in}}%
\pgfpathlineto{\pgfqpoint{0.823766in}{3.439843in}}%
\pgfpathlineto{\pgfqpoint{0.825580in}{3.439643in}}%
\pgfpathlineto{\pgfqpoint{0.827558in}{3.437561in}}%
\pgfpathlineto{\pgfqpoint{0.831331in}{3.430188in}}%
\pgfpathlineto{\pgfqpoint{0.850351in}{3.394823in}}%
\pgfpathlineto{\pgfqpoint{0.887585in}{3.334248in}}%
\pgfpathlineto{\pgfqpoint{0.889921in}{3.330708in}}%
\pgfpathlineto{\pgfqpoint{0.901043in}{3.312716in}}%
\pgfpathlineto{\pgfqpoint{0.908373in}{3.303213in}}%
\pgfpathlineto{\pgfqpoint{0.917050in}{3.290009in}}%
\pgfpathlineto{\pgfqpoint{0.925618in}{3.279829in}}%
\pgfpathlineto{\pgfqpoint{0.933018in}{3.269033in}}%
\pgfpathlineto{\pgfqpoint{0.937019in}{3.266089in}}%
\pgfpathlineto{\pgfqpoint{0.939154in}{3.264068in}}%
\pgfpathlineto{\pgfqpoint{0.952211in}{3.247082in}}%
\pgfpathlineto{\pgfqpoint{0.956001in}{3.244379in}}%
\pgfpathlineto{\pgfqpoint{0.968071in}{3.229072in}}%
\pgfpathlineto{\pgfqpoint{0.973110in}{3.225786in}}%
\pgfpathlineto{\pgfqpoint{0.982841in}{3.213110in}}%
\pgfpathlineto{\pgfqpoint{0.986155in}{3.211829in}}%
\pgfpathlineto{\pgfqpoint{0.988442in}{3.210857in}}%
\pgfpathlineto{\pgfqpoint{0.990988in}{3.207671in}}%
\pgfpathlineto{\pgfqpoint{0.998367in}{3.198265in}}%
\pgfpathlineto{\pgfqpoint{1.001664in}{3.196805in}}%
\pgfpathlineto{\pgfqpoint{1.005623in}{3.195572in}}%
\pgfpathlineto{\pgfqpoint{1.008492in}{3.191837in}}%
\pgfpathlineto{\pgfqpoint{1.014624in}{3.184190in}}%
\pgfpathlineto{\pgfqpoint{1.017756in}{3.182722in}}%
\pgfpathlineto{\pgfqpoint{1.022566in}{3.181249in}}%
\pgfpathlineto{\pgfqpoint{1.026181in}{3.176421in}}%
\pgfpathlineto{\pgfqpoint{1.031165in}{3.170617in}}%
\pgfpathlineto{\pgfqpoint{1.034297in}{3.169368in}}%
\pgfpathlineto{\pgfqpoint{1.039087in}{3.168312in}}%
\pgfpathlineto{\pgfqpoint{1.042151in}{3.164431in}}%
\pgfpathlineto{\pgfqpoint{1.047376in}{3.158437in}}%
\pgfpathlineto{\pgfqpoint{1.050343in}{3.157328in}}%
\pgfpathlineto{\pgfqpoint{1.056204in}{3.156303in}}%
\pgfpathlineto{\pgfqpoint{1.059953in}{3.151484in}}%
\pgfpathlineto{\pgfqpoint{1.063587in}{3.147749in}}%
\pgfpathlineto{\pgfqpoint{1.066554in}{3.146778in}}%
\pgfpathlineto{\pgfqpoint{1.069748in}{3.147558in}}%
\pgfpathlineto{\pgfqpoint{1.071892in}{3.147094in}}%
\pgfpathlineto{\pgfqpoint{1.074199in}{3.144932in}}%
\pgfpathlineto{\pgfqpoint{1.080348in}{3.138060in}}%
\pgfpathlineto{\pgfqpoint{1.083150in}{3.137284in}}%
\pgfpathlineto{\pgfqpoint{1.089539in}{3.136951in}}%
\pgfpathlineto{\pgfqpoint{1.092942in}{3.132812in}}%
\pgfpathlineto{\pgfqpoint{1.096826in}{3.128968in}}%
\pgfpathlineto{\pgfqpoint{1.099628in}{3.128270in}}%
\pgfpathlineto{\pgfqpoint{1.106183in}{3.128143in}}%
\pgfpathlineto{\pgfqpoint{1.109611in}{3.124046in}}%
\pgfpathlineto{\pgfqpoint{1.113224in}{3.120578in}}%
\pgfpathlineto{\pgfqpoint{1.116026in}{3.119938in}}%
\pgfpathlineto{\pgfqpoint{1.118876in}{3.121004in}}%
\pgfpathlineto{\pgfqpoint{1.121184in}{3.121057in}}%
\pgfpathlineto{\pgfqpoint{1.123492in}{3.119461in}}%
\pgfpathlineto{\pgfqpoint{1.130831in}{3.112693in}}%
\pgfpathlineto{\pgfqpoint{1.133468in}{3.112963in}}%
\pgfpathlineto{\pgfqpoint{1.138164in}{3.114188in}}%
\pgfpathlineto{\pgfqpoint{1.140472in}{3.112517in}}%
\pgfpathlineto{\pgfqpoint{1.146948in}{3.106493in}}%
\pgfpathlineto{\pgfqpoint{1.149585in}{3.106590in}}%
\pgfpathlineto{\pgfqpoint{1.155142in}{3.107920in}}%
\pgfpathlineto{\pgfqpoint{1.157615in}{3.105823in}}%
\pgfpathlineto{\pgfqpoint{1.162677in}{3.100893in}}%
\pgfpathlineto{\pgfqpoint{1.165314in}{3.100566in}}%
\pgfpathlineto{\pgfqpoint{1.167966in}{3.101886in}}%
\pgfpathlineto{\pgfqpoint{1.170438in}{3.102591in}}%
\pgfpathlineto{\pgfqpoint{1.172746in}{3.101597in}}%
\pgfpathlineto{\pgfqpoint{1.175895in}{3.098106in}}%
\pgfpathlineto{\pgfqpoint{1.179078in}{3.095401in}}%
\pgfpathlineto{\pgfqpoint{1.181716in}{3.095031in}}%
\pgfpathlineto{\pgfqpoint{1.184312in}{3.096305in}}%
\pgfpathlineto{\pgfqpoint{1.186793in}{3.097176in}}%
\pgfpathlineto{\pgfqpoint{1.189101in}{3.096347in}}%
\pgfpathlineto{\pgfqpoint{1.192012in}{3.093294in}}%
\pgfpathlineto{\pgfqpoint{1.195524in}{3.090182in}}%
\pgfpathlineto{\pgfqpoint{1.198162in}{3.089833in}}%
\pgfpathlineto{\pgfqpoint{1.200775in}{3.091164in}}%
\pgfpathlineto{\pgfqpoint{1.203242in}{3.092120in}}%
\pgfpathlineto{\pgfqpoint{1.205550in}{3.091396in}}%
\pgfpathlineto{\pgfqpoint{1.208310in}{3.088623in}}%
\pgfpathlineto{\pgfqpoint{1.212120in}{3.085333in}}%
\pgfpathlineto{\pgfqpoint{1.214758in}{3.085140in}}%
\pgfpathlineto{\pgfqpoint{1.217419in}{3.086667in}}%
\pgfpathlineto{\pgfqpoint{1.219892in}{3.087620in}}%
\pgfpathlineto{\pgfqpoint{1.222200in}{3.086903in}}%
\pgfpathlineto{\pgfqpoint{1.225004in}{3.084108in}}%
\pgfpathlineto{\pgfqpoint{1.228447in}{3.081253in}}%
\pgfpathlineto{\pgfqpoint{1.231085in}{3.081073in}}%
\pgfpathlineto{\pgfqpoint{1.233680in}{3.082592in}}%
\pgfpathlineto{\pgfqpoint{1.236219in}{3.083858in}}%
\pgfpathlineto{\pgfqpoint{1.238527in}{3.083390in}}%
\pgfpathlineto{\pgfqpoint{1.240985in}{3.081231in}}%
\pgfpathlineto{\pgfqpoint{1.244650in}{3.078027in}}%
\pgfpathlineto{\pgfqpoint{1.247288in}{3.077781in}}%
\pgfpathlineto{\pgfqpoint{1.249832in}{3.079208in}}%
\pgfpathlineto{\pgfqpoint{1.252866in}{3.080943in}}%
\pgfpathlineto{\pgfqpoint{1.255174in}{3.080503in}}%
\pgfpathlineto{\pgfqpoint{1.257722in}{3.078277in}}%
\pgfpathlineto{\pgfqpoint{1.261031in}{3.075437in}}%
\pgfpathlineto{\pgfqpoint{1.263669in}{3.075158in}}%
\pgfpathlineto{\pgfqpoint{1.266241in}{3.076586in}}%
\pgfpathlineto{\pgfqpoint{1.269547in}{3.078493in}}%
\pgfpathlineto{\pgfqpoint{1.271855in}{3.077992in}}%
\pgfpathlineto{\pgfqpoint{1.274459in}{3.075646in}}%
\pgfpathlineto{\pgfqpoint{1.277623in}{3.073023in}}%
\pgfpathlineto{\pgfqpoint{1.280096in}{3.072756in}}%
\pgfpathlineto{\pgfqpoint{1.282671in}{3.074137in}}%
\pgfpathlineto{\pgfqpoint{1.285986in}{3.076066in}}%
\pgfpathlineto{\pgfqpoint{1.288294in}{3.075588in}}%
\pgfpathlineto{\pgfqpoint{1.290921in}{3.073252in}}%
\pgfpathlineto{\pgfqpoint{1.294023in}{3.070672in}}%
\pgfpathlineto{\pgfqpoint{1.296496in}{3.070389in}}%
\pgfpathlineto{\pgfqpoint{1.299056in}{3.071750in}}%
\pgfpathlineto{\pgfqpoint{1.302646in}{3.073868in}}%
\pgfpathlineto{\pgfqpoint{1.304953in}{3.073342in}}%
\pgfpathlineto{\pgfqpoint{1.307608in}{3.070931in}}%
\pgfpathlineto{\pgfqpoint{1.310669in}{3.068470in}}%
\pgfpathlineto{\pgfqpoint{1.313142in}{3.068255in}}%
\pgfpathlineto{\pgfqpoint{1.315699in}{3.069692in}}%
\pgfpathlineto{\pgfqpoint{1.319001in}{3.071652in}}%
\pgfpathlineto{\pgfqpoint{1.321309in}{3.071217in}}%
\pgfpathlineto{\pgfqpoint{1.323914in}{3.068961in}}%
\pgfpathlineto{\pgfqpoint{1.327261in}{3.066281in}}%
\pgfpathlineto{\pgfqpoint{1.329733in}{3.066167in}}%
\pgfpathlineto{\pgfqpoint{1.332307in}{3.067730in}}%
\pgfpathlineto{\pgfqpoint{1.335232in}{3.069549in}}%
\pgfpathlineto{\pgfqpoint{1.337540in}{3.069294in}}%
\pgfpathlineto{\pgfqpoint{1.340013in}{3.067374in}}%
\pgfpathlineto{\pgfqpoint{1.343819in}{3.064265in}}%
\pgfpathlineto{\pgfqpoint{1.346292in}{3.064237in}}%
\pgfpathlineto{\pgfqpoint{1.348821in}{3.065852in}}%
\pgfpathlineto{\pgfqpoint{1.351771in}{3.067741in}}%
\pgfpathlineto{\pgfqpoint{1.354079in}{3.067525in}}%
\pgfpathlineto{\pgfqpoint{1.356552in}{3.065660in}}%
\pgfpathlineto{\pgfqpoint{1.360468in}{3.062501in}}%
\pgfpathlineto{\pgfqpoint{1.362941in}{3.062531in}}%
\pgfpathlineto{\pgfqpoint{1.365477in}{3.064221in}}%
\pgfpathlineto{\pgfqpoint{1.368256in}{3.065904in}}%
\pgfpathlineto{\pgfqpoint{1.370564in}{3.065666in}}%
\pgfpathlineto{\pgfqpoint{1.373037in}{3.063783in}}%
\pgfpathlineto{\pgfqpoint{1.376875in}{3.060676in}}%
\pgfpathlineto{\pgfqpoint{1.379348in}{3.060685in}}%
\pgfpathlineto{\pgfqpoint{1.381862in}{3.062337in}}%
\pgfpathlineto{\pgfqpoint{1.384705in}{3.064144in}}%
\pgfpathlineto{\pgfqpoint{1.387013in}{3.063961in}}%
\pgfpathlineto{\pgfqpoint{1.389485in}{3.062138in}}%
\pgfpathlineto{\pgfqpoint{1.393531in}{3.058962in}}%
\pgfpathlineto{\pgfqpoint{1.396004in}{3.059099in}}%
\pgfpathlineto{\pgfqpoint{1.398593in}{3.060959in}}%
\pgfpathlineto{\pgfqpoint{1.401391in}{3.062613in}}%
\pgfpathlineto{\pgfqpoint{1.403699in}{3.062346in}}%
\pgfpathlineto{\pgfqpoint{1.406280in}{3.060328in}}%
\pgfpathlineto{\pgfqpoint{1.409814in}{3.057579in}}%
\pgfpathlineto{\pgfqpoint{1.412286in}{3.057611in}}%
\pgfpathlineto{\pgfqpoint{1.414863in}{3.059361in}}%
\pgfpathlineto{\pgfqpoint{1.417756in}{3.061260in}}%
\pgfpathlineto{\pgfqpoint{1.420229in}{3.061052in}}%
\pgfpathlineto{\pgfqpoint{1.422800in}{3.059071in}}%
\pgfpathlineto{\pgfqpoint{1.426278in}{3.056402in}}%
\pgfpathlineto{\pgfqpoint{1.428751in}{3.056449in}}%
\pgfpathlineto{\pgfqpoint{1.431253in}{3.058143in}}%
\pgfpathlineto{\pgfqpoint{1.434195in}{3.060158in}}%
\pgfpathlineto{\pgfqpoint{1.436667in}{3.060014in}}%
\pgfpathlineto{\pgfqpoint{1.439235in}{3.058101in}}%
\pgfpathlineto{\pgfqpoint{1.442749in}{3.055397in}}%
\pgfpathlineto{\pgfqpoint{1.445222in}{3.055450in}}%
\pgfpathlineto{\pgfqpoint{1.447764in}{3.057194in}}%
\pgfpathlineto{\pgfqpoint{1.450718in}{3.059239in}}%
\pgfpathlineto{\pgfqpoint{1.453190in}{3.059110in}}%
\pgfpathlineto{\pgfqpoint{1.455752in}{3.057218in}}%
\pgfpathlineto{\pgfqpoint{1.459199in}{3.054592in}}%
\pgfpathlineto{\pgfqpoint{1.461671in}{3.054648in}}%
\pgfpathlineto{\pgfqpoint{1.464196in}{3.056381in}}%
\pgfpathlineto{\pgfqpoint{1.467293in}{3.058506in}}%
\pgfpathlineto{\pgfqpoint{1.469765in}{3.058320in}}%
\pgfpathlineto{\pgfqpoint{1.472311in}{3.056396in}}%
\pgfpathlineto{\pgfqpoint{1.475572in}{3.053925in}}%
\pgfpathlineto{\pgfqpoint{1.478044in}{3.053939in}}%
\pgfpathlineto{\pgfqpoint{1.480607in}{3.055671in}}%
\pgfpathlineto{\pgfqpoint{1.483815in}{3.057944in}}%
\pgfpathlineto{\pgfqpoint{1.486288in}{3.057784in}}%
\pgfpathlineto{\pgfqpoint{1.488808in}{3.055915in}}%
\pgfpathlineto{\pgfqpoint{1.492064in}{3.053506in}}%
\pgfpathlineto{\pgfqpoint{1.494537in}{3.053565in}}%
\pgfpathlineto{\pgfqpoint{1.497099in}{3.055343in}}%
\pgfpathlineto{\pgfqpoint{1.500343in}{3.057691in}}%
\pgfpathlineto{\pgfqpoint{1.502816in}{3.057557in}}%
\pgfpathlineto{\pgfqpoint{1.505385in}{3.055658in}}%
\pgfpathlineto{\pgfqpoint{1.508436in}{3.053415in}}%
\pgfpathlineto{\pgfqpoint{1.510909in}{3.053423in}}%
\pgfpathlineto{\pgfqpoint{1.513473in}{3.055152in}}%
\pgfpathlineto{\pgfqpoint{1.516950in}{3.057664in}}%
\pgfpathlineto{\pgfqpoint{1.519423in}{3.057459in}}%
\pgfpathlineto{\pgfqpoint{1.522011in}{3.055472in}}%
\pgfpathlineto{\pgfqpoint{1.524908in}{3.053419in}}%
\pgfpathlineto{\pgfqpoint{1.527380in}{3.053447in}}%
\pgfpathlineto{\pgfqpoint{1.529943in}{3.055195in}}%
\pgfpathlineto{\pgfqpoint{1.533436in}{3.057761in}}%
\pgfpathlineto{\pgfqpoint{1.535909in}{3.057583in}}%
\pgfpathlineto{\pgfqpoint{1.538481in}{3.055638in}}%
\pgfpathlineto{\pgfqpoint{1.541396in}{3.053588in}}%
\pgfpathlineto{\pgfqpoint{1.543868in}{3.053633in}}%
\pgfpathlineto{\pgfqpoint{1.546428in}{3.055395in}}%
\pgfpathlineto{\pgfqpoint{1.549921in}{3.057981in}}%
\pgfpathlineto{\pgfqpoint{1.552393in}{3.057815in}}%
\pgfpathlineto{\pgfqpoint{1.554965in}{3.055884in}}%
\pgfpathlineto{\pgfqpoint{1.557903in}{3.053834in}}%
\pgfpathlineto{\pgfqpoint{1.560376in}{3.053900in}}%
\pgfpathlineto{\pgfqpoint{1.562953in}{3.055701in}}%
\pgfpathlineto{\pgfqpoint{1.566411in}{3.058256in}}%
\pgfpathlineto{\pgfqpoint{1.568884in}{3.058092in}}%
\pgfpathlineto{\pgfqpoint{1.571455in}{3.056160in}}%
\pgfpathlineto{\pgfqpoint{1.574361in}{3.054124in}}%
\pgfpathlineto{\pgfqpoint{1.576834in}{3.054170in}}%
\pgfpathlineto{\pgfqpoint{1.579398in}{3.055936in}}%
\pgfpathlineto{\pgfqpoint{1.582875in}{3.058511in}}%
\pgfpathlineto{\pgfqpoint{1.585348in}{3.058348in}}%
\pgfpathlineto{\pgfqpoint{1.587925in}{3.056413in}}%
\pgfpathlineto{\pgfqpoint{1.590849in}{3.054362in}}%
\pgfpathlineto{\pgfqpoint{1.593322in}{3.054413in}}%
\pgfpathlineto{\pgfqpoint{1.595897in}{3.056194in}}%
\pgfpathlineto{\pgfqpoint{1.599375in}{3.058764in}}%
\pgfpathlineto{\pgfqpoint{1.601848in}{3.058595in}}%
\pgfpathlineto{\pgfqpoint{1.604420in}{3.056657in}}%
\pgfpathlineto{\pgfqpoint{1.607356in}{3.054598in}}%
\pgfpathlineto{\pgfqpoint{1.609829in}{3.054651in}}%
\pgfpathlineto{\pgfqpoint{1.612401in}{3.056431in}}%
\pgfpathlineto{\pgfqpoint{1.615866in}{3.058969in}}%
\pgfpathlineto{\pgfqpoint{1.618339in}{3.058785in}}%
\pgfpathlineto{\pgfqpoint{1.620917in}{3.056822in}}%
\pgfpathlineto{\pgfqpoint{1.623845in}{3.054754in}}%
\pgfpathlineto{\pgfqpoint{1.626318in}{3.054792in}}%
\pgfpathlineto{\pgfqpoint{1.628893in}{3.056558in}}%
\pgfpathlineto{\pgfqpoint{1.632348in}{3.059060in}}%
\pgfpathlineto{\pgfqpoint{1.634821in}{3.058858in}}%
\pgfpathlineto{\pgfqpoint{1.637404in}{3.056872in}}%
\pgfpathlineto{\pgfqpoint{1.640254in}{3.054813in}}%
\pgfpathlineto{\pgfqpoint{1.642727in}{3.054786in}}%
\pgfpathlineto{\pgfqpoint{1.645180in}{3.056364in}}%
\pgfpathlineto{\pgfqpoint{1.649026in}{3.059108in}}%
\pgfpathlineto{\pgfqpoint{1.651499in}{3.058777in}}%
\pgfpathlineto{\pgfqpoint{1.654119in}{3.056626in}}%
\pgfpathlineto{\pgfqpoint{1.656873in}{3.054746in}}%
\pgfpathlineto{\pgfqpoint{1.659346in}{3.054786in}}%
\pgfpathlineto{\pgfqpoint{1.661887in}{3.056518in}}%
\pgfpathlineto{\pgfqpoint{1.665260in}{3.058900in}}%
\pgfpathlineto{\pgfqpoint{1.667732in}{3.058677in}}%
\pgfpathlineto{\pgfqpoint{1.670314in}{3.056667in}}%
\pgfpathlineto{\pgfqpoint{1.673242in}{3.054444in}}%
\pgfpathlineto{\pgfqpoint{1.675715in}{3.054351in}}%
\pgfpathlineto{\pgfqpoint{1.678159in}{3.055858in}}%
\pgfpathlineto{\pgfqpoint{1.681963in}{3.058449in}}%
\pgfpathlineto{\pgfqpoint{1.684271in}{3.058130in}}%
\pgfpathlineto{\pgfqpoint{1.686865in}{3.056056in}}%
\pgfpathlineto{\pgfqpoint{1.689715in}{3.053946in}}%
\pgfpathlineto{\pgfqpoint{1.692188in}{3.053879in}}%
\pgfpathlineto{\pgfqpoint{1.694637in}{3.055418in}}%
\pgfpathlineto{\pgfqpoint{1.698492in}{3.058096in}}%
\pgfpathlineto{\pgfqpoint{1.700965in}{3.057719in}}%
\pgfpathlineto{\pgfqpoint{1.703618in}{3.055499in}}%
\pgfpathlineto{\pgfqpoint{1.706455in}{3.053621in}}%
\pgfpathlineto{\pgfqpoint{1.708928in}{3.053744in}}%
\pgfpathlineto{\pgfqpoint{1.711444in}{3.055539in}}%
\pgfpathlineto{\pgfqpoint{1.714653in}{3.057859in}}%
\pgfpathlineto{\pgfqpoint{1.717125in}{3.057732in}}%
\pgfpathlineto{\pgfqpoint{1.719688in}{3.055845in}}%
\pgfpathlineto{\pgfqpoint{1.722732in}{3.053590in}}%
\pgfpathlineto{\pgfqpoint{1.725204in}{3.053582in}}%
\pgfpathlineto{\pgfqpoint{1.727761in}{3.055286in}}%
\pgfpathlineto{\pgfqpoint{1.731253in}{3.057794in}}%
\pgfpathlineto{\pgfqpoint{1.733726in}{3.057576in}}%
\pgfpathlineto{\pgfqpoint{1.736311in}{3.055578in}}%
\pgfpathlineto{\pgfqpoint{1.739198in}{3.053514in}}%
\pgfpathlineto{\pgfqpoint{1.741670in}{3.053521in}}%
\pgfpathlineto{\pgfqpoint{1.744234in}{3.055248in}}%
\pgfpathlineto{\pgfqpoint{1.747774in}{3.057834in}}%
\pgfpathlineto{\pgfqpoint{1.750247in}{3.057633in}}%
\pgfpathlineto{\pgfqpoint{1.752817in}{3.055664in}}%
\pgfpathlineto{\pgfqpoint{1.755740in}{3.053611in}}%
\pgfpathlineto{\pgfqpoint{1.758213in}{3.053663in}}%
\pgfpathlineto{\pgfqpoint{1.760776in}{3.055436in}}%
\pgfpathlineto{\pgfqpoint{1.764156in}{3.057895in}}%
\pgfpathlineto{\pgfqpoint{1.766629in}{3.057730in}}%
\pgfpathlineto{\pgfqpoint{1.769187in}{3.055809in}}%
\pgfpathlineto{\pgfqpoint{1.772164in}{3.053602in}}%
\pgfpathlineto{\pgfqpoint{1.774637in}{3.053574in}}%
\pgfpathlineto{\pgfqpoint{1.777074in}{3.055140in}}%
\pgfpathlineto{\pgfqpoint{1.780874in}{3.057821in}}%
\pgfpathlineto{\pgfqpoint{1.783347in}{3.057489in}}%
\pgfpathlineto{\pgfqpoint{1.785959in}{3.055352in}}%
\pgfpathlineto{\pgfqpoint{1.788793in}{3.053435in}}%
\pgfpathlineto{\pgfqpoint{1.791265in}{3.053520in}}%
\pgfpathlineto{\pgfqpoint{1.793819in}{3.055318in}}%
\pgfpathlineto{\pgfqpoint{1.797027in}{3.057679in}}%
\pgfpathlineto{\pgfqpoint{1.799500in}{3.057586in}}%
\pgfpathlineto{\pgfqpoint{1.802020in}{3.055784in}}%
\pgfpathlineto{\pgfqpoint{1.805149in}{3.053483in}}%
\pgfpathlineto{\pgfqpoint{1.807622in}{3.053510in}}%
\pgfpathlineto{\pgfqpoint{1.810186in}{3.055258in}}%
\pgfpathlineto{\pgfqpoint{1.813678in}{3.057823in}}%
\pgfpathlineto{\pgfqpoint{1.816151in}{3.057645in}}%
\pgfpathlineto{\pgfqpoint{1.818724in}{3.055699in}}%
\pgfpathlineto{\pgfqpoint{1.821637in}{3.053648in}}%
\pgfpathlineto{\pgfqpoint{1.824110in}{3.053692in}}%
\pgfpathlineto{\pgfqpoint{1.826671in}{3.055452in}}%
\pgfpathlineto{\pgfqpoint{1.830163in}{3.058036in}}%
\pgfpathlineto{\pgfqpoint{1.832636in}{3.057868in}}%
\pgfpathlineto{\pgfqpoint{1.835207in}{3.055934in}}%
\pgfpathlineto{\pgfqpoint{1.838144in}{3.053883in}}%
\pgfpathlineto{\pgfqpoint{1.840617in}{3.053947in}}%
\pgfpathlineto{\pgfqpoint{1.843189in}{3.055739in}}%
\pgfpathlineto{\pgfqpoint{1.846654in}{3.058298in}}%
\pgfpathlineto{\pgfqpoint{1.849127in}{3.058132in}}%
\pgfpathlineto{\pgfqpoint{1.851702in}{3.056193in}}%
\pgfpathlineto{\pgfqpoint{1.854559in}{3.054175in}}%
\pgfpathlineto{\pgfqpoint{1.857032in}{3.054193in}}%
\pgfpathlineto{\pgfqpoint{1.859484in}{3.055816in}}%
\pgfpathlineto{\pgfqpoint{1.863339in}{3.058623in}}%
\pgfpathlineto{\pgfqpoint{1.865812in}{3.058327in}}%
\pgfpathlineto{\pgfqpoint{1.868420in}{3.056229in}}%
\pgfpathlineto{\pgfqpoint{1.871243in}{3.054386in}}%
\pgfpathlineto{\pgfqpoint{1.873716in}{3.054522in}}%
\pgfpathlineto{\pgfqpoint{1.876245in}{3.056341in}}%
\pgfpathlineto{\pgfqpoint{1.879453in}{3.058734in}}%
\pgfpathlineto{\pgfqpoint{1.881926in}{3.058663in}}%
\pgfpathlineto{\pgfqpoint{1.884445in}{3.056877in}}%
\pgfpathlineto{\pgfqpoint{1.887740in}{3.054532in}}%
\pgfpathlineto{\pgfqpoint{1.890213in}{3.054669in}}%
\pgfpathlineto{\pgfqpoint{1.892730in}{3.056476in}}%
\pgfpathlineto{\pgfqpoint{1.895938in}{3.058858in}}%
\pgfpathlineto{\pgfqpoint{1.898410in}{3.058778in}}%
\pgfpathlineto{\pgfqpoint{1.900929in}{3.056984in}}%
\pgfpathlineto{\pgfqpoint{1.904219in}{3.054631in}}%
\pgfpathlineto{\pgfqpoint{1.906692in}{3.054756in}}%
\pgfpathlineto{\pgfqpoint{1.909214in}{3.056558in}}%
\pgfpathlineto{\pgfqpoint{1.912423in}{3.058930in}}%
\pgfpathlineto{\pgfqpoint{1.914895in}{3.058841in}}%
\pgfpathlineto{\pgfqpoint{1.917414in}{3.057039in}}%
\pgfpathlineto{\pgfqpoint{1.920704in}{3.054678in}}%
\pgfpathlineto{\pgfqpoint{1.923177in}{3.054797in}}%
\pgfpathlineto{\pgfqpoint{1.925699in}{3.056592in}}%
\pgfpathlineto{\pgfqpoint{1.928907in}{3.058956in}}%
\pgfpathlineto{\pgfqpoint{1.931380in}{3.058860in}}%
\pgfpathlineto{\pgfqpoint{1.933899in}{3.057051in}}%
\pgfpathlineto{\pgfqpoint{1.937171in}{3.054688in}}%
\pgfpathlineto{\pgfqpoint{1.939644in}{3.054791in}}%
\pgfpathlineto{\pgfqpoint{1.942186in}{3.056589in}}%
\pgfpathlineto{\pgfqpoint{1.945400in}{3.058950in}}%
\pgfpathlineto{\pgfqpoint{1.947872in}{3.058847in}}%
\pgfpathlineto{\pgfqpoint{1.950384in}{3.057037in}}%
\pgfpathlineto{\pgfqpoint{1.953650in}{3.054670in}}%
\pgfpathlineto{\pgfqpoint{1.956123in}{3.054764in}}%
\pgfpathlineto{\pgfqpoint{1.958696in}{3.056586in}}%
\pgfpathlineto{\pgfqpoint{1.961907in}{3.058935in}}%
\pgfpathlineto{\pgfqpoint{1.964379in}{3.058824in}}%
\pgfpathlineto{\pgfqpoint{1.966869in}{3.057029in}}%
\pgfpathlineto{\pgfqpoint{1.970086in}{3.054651in}}%
\pgfpathlineto{\pgfqpoint{1.972559in}{3.054695in}}%
\pgfpathlineto{\pgfqpoint{1.975106in}{3.056438in}}%
\pgfpathlineto{\pgfqpoint{1.978526in}{3.058857in}}%
\pgfpathlineto{\pgfqpoint{1.980998in}{3.058624in}}%
\pgfpathlineto{\pgfqpoint{1.983582in}{3.056609in}}%
\pgfpathlineto{\pgfqpoint{1.986492in}{3.054513in}}%
\pgfpathlineto{\pgfqpoint{1.988964in}{3.054511in}}%
\pgfpathlineto{\pgfqpoint{1.991519in}{3.056216in}}%
\pgfpathlineto{\pgfqpoint{1.994999in}{3.058713in}}%
\pgfpathlineto{\pgfqpoint{1.997472in}{3.058493in}}%
\pgfpathlineto{\pgfqpoint{2.000062in}{3.056486in}}%
\pgfpathlineto{\pgfqpoint{2.002864in}{3.054436in}}%
\pgfpathlineto{\pgfqpoint{2.005336in}{3.054371in}}%
\pgfpathlineto{\pgfqpoint{2.007809in}{3.055933in}}%
\pgfpathlineto{\pgfqpoint{2.011694in}{3.058713in}}%
\pgfpathlineto{\pgfqpoint{2.014166in}{3.058383in}}%
\pgfpathlineto{\pgfqpoint{2.016773in}{3.056248in}}%
\pgfpathlineto{\pgfqpoint{2.019569in}{3.054351in}}%
\pgfpathlineto{\pgfqpoint{2.022042in}{3.054417in}}%
\pgfpathlineto{\pgfqpoint{2.024618in}{3.056213in}}%
\pgfpathlineto{\pgfqpoint{2.027823in}{3.058543in}}%
\pgfpathlineto{\pgfqpoint{2.030295in}{3.058427in}}%
\pgfpathlineto{\pgfqpoint{2.032808in}{3.056604in}}%
\pgfpathlineto{\pgfqpoint{2.036080in}{3.054218in}}%
\pgfpathlineto{\pgfqpoint{2.038553in}{3.054305in}}%
\pgfpathlineto{\pgfqpoint{2.041093in}{3.056087in}}%
\pgfpathlineto{\pgfqpoint{2.044301in}{3.058429in}}%
\pgfpathlineto{\pgfqpoint{2.046774in}{3.058319in}}%
\pgfpathlineto{\pgfqpoint{2.049293in}{3.056498in}}%
\pgfpathlineto{\pgfqpoint{2.052577in}{3.054124in}}%
\pgfpathlineto{\pgfqpoint{2.055050in}{3.054231in}}%
\pgfpathlineto{\pgfqpoint{2.057578in}{3.056022in}}%
\pgfpathlineto{\pgfqpoint{2.060786in}{3.058380in}}%
\pgfpathlineto{\pgfqpoint{2.063259in}{3.058283in}}%
\pgfpathlineto{\pgfqpoint{2.065778in}{3.056474in}}%
\pgfpathlineto{\pgfqpoint{2.068977in}{3.054137in}}%
\pgfpathlineto{\pgfqpoint{2.071450in}{3.054197in}}%
\pgfpathlineto{\pgfqpoint{2.074003in}{3.055965in}}%
\pgfpathlineto{\pgfqpoint{2.077459in}{3.058471in}}%
\pgfpathlineto{\pgfqpoint{2.079932in}{3.058273in}}%
\pgfpathlineto{\pgfqpoint{2.082497in}{3.056310in}}%
\pgfpathlineto{\pgfqpoint{2.085441in}{3.054219in}}%
\pgfpathlineto{\pgfqpoint{2.087913in}{3.054254in}}%
\pgfpathlineto{\pgfqpoint{2.090486in}{3.056017in}}%
\pgfpathlineto{\pgfqpoint{2.093952in}{3.058509in}}%
\pgfpathlineto{\pgfqpoint{2.096425in}{3.058292in}}%
\pgfpathlineto{\pgfqpoint{2.098993in}{3.056306in}}%
\pgfpathlineto{\pgfqpoint{2.101935in}{3.054208in}}%
\pgfpathlineto{\pgfqpoint{2.104408in}{3.054237in}}%
\pgfpathlineto{\pgfqpoint{2.106961in}{3.055972in}}%
\pgfpathlineto{\pgfqpoint{2.110392in}{3.058420in}}%
\pgfpathlineto{\pgfqpoint{2.112865in}{3.058200in}}%
\pgfpathlineto{\pgfqpoint{2.115458in}{3.056194in}}%
\pgfpathlineto{\pgfqpoint{2.118321in}{3.054130in}}%
\pgfpathlineto{\pgfqpoint{2.120793in}{3.054114in}}%
\pgfpathlineto{\pgfqpoint{2.123243in}{3.055702in}}%
\pgfpathlineto{\pgfqpoint{2.127091in}{3.058460in}}%
\pgfpathlineto{\pgfqpoint{2.129563in}{3.058139in}}%
\pgfpathlineto{\pgfqpoint{2.132177in}{3.056009in}}%
\pgfpathlineto{\pgfqpoint{2.134960in}{3.054137in}}%
\pgfpathlineto{\pgfqpoint{2.137433in}{3.054214in}}%
\pgfpathlineto{\pgfqpoint{2.140001in}{3.056014in}}%
\pgfpathlineto{\pgfqpoint{2.143209in}{3.058368in}}%
\pgfpathlineto{\pgfqpoint{2.145682in}{3.058268in}}%
\pgfpathlineto{\pgfqpoint{2.148202in}{3.056455in}}%
\pgfpathlineto{\pgfqpoint{2.151348in}{3.054134in}}%
\pgfpathlineto{\pgfqpoint{2.153820in}{3.054157in}}%
\pgfpathlineto{\pgfqpoint{2.156382in}{3.055897in}}%
\pgfpathlineto{\pgfqpoint{2.159833in}{3.058403in}}%
\pgfpathlineto{\pgfqpoint{2.162306in}{3.058214in}}%
\pgfpathlineto{\pgfqpoint{2.164895in}{3.056241in}}%
\pgfpathlineto{\pgfqpoint{2.167819in}{3.054168in}}%
\pgfpathlineto{\pgfqpoint{2.170291in}{3.054201in}}%
\pgfpathlineto{\pgfqpoint{2.172852in}{3.055948in}}%
\pgfpathlineto{\pgfqpoint{2.176345in}{3.058511in}}%
\pgfpathlineto{\pgfqpoint{2.178818in}{3.058327in}}%
\pgfpathlineto{\pgfqpoint{2.181391in}{3.056374in}}%
\pgfpathlineto{\pgfqpoint{2.184302in}{3.054314in}}%
\pgfpathlineto{\pgfqpoint{2.186775in}{3.054345in}}%
\pgfpathlineto{\pgfqpoint{2.189337in}{3.056092in}}%
\pgfpathlineto{\pgfqpoint{2.192830in}{3.058652in}}%
\pgfpathlineto{\pgfqpoint{2.195303in}{3.058466in}}%
\pgfpathlineto{\pgfqpoint{2.197876in}{3.056509in}}%
\pgfpathlineto{\pgfqpoint{2.200784in}{3.054447in}}%
\pgfpathlineto{\pgfqpoint{2.203257in}{3.054473in}}%
\pgfpathlineto{\pgfqpoint{2.205822in}{3.056217in}}%
\pgfpathlineto{\pgfqpoint{2.209315in}{3.058771in}}%
\pgfpathlineto{\pgfqpoint{2.211787in}{3.058581in}}%
\pgfpathlineto{\pgfqpoint{2.214361in}{3.056618in}}%
\pgfpathlineto{\pgfqpoint{2.217266in}{3.054552in}}%
\pgfpathlineto{\pgfqpoint{2.219739in}{3.054572in}}%
\pgfpathlineto{\pgfqpoint{2.222307in}{3.056312in}}%
\pgfpathlineto{\pgfqpoint{2.225800in}{3.058859in}}%
\pgfpathlineto{\pgfqpoint{2.228272in}{3.058662in}}%
\pgfpathlineto{\pgfqpoint{2.230846in}{3.056694in}}%
\pgfpathlineto{\pgfqpoint{2.233751in}{3.054623in}}%
\pgfpathlineto{\pgfqpoint{2.236224in}{3.054637in}}%
\pgfpathlineto{\pgfqpoint{2.238792in}{3.056372in}}%
\pgfpathlineto{\pgfqpoint{2.242284in}{3.058912in}}%
\pgfpathlineto{\pgfqpoint{2.244757in}{3.058710in}}%
\pgfpathlineto{\pgfqpoint{2.247330in}{3.056736in}}%
\pgfpathlineto{\pgfqpoint{2.250236in}{3.054660in}}%
\pgfpathlineto{\pgfqpoint{2.252709in}{3.054671in}}%
\pgfpathlineto{\pgfqpoint{2.255277in}{3.056401in}}%
\pgfpathlineto{\pgfqpoint{2.258769in}{3.058934in}}%
\pgfpathlineto{\pgfqpoint{2.261242in}{3.058728in}}%
\pgfpathlineto{\pgfqpoint{2.263826in}{3.056738in}}%
\pgfpathlineto{\pgfqpoint{2.266752in}{3.054659in}}%
\pgfpathlineto{\pgfqpoint{2.269225in}{3.054687in}}%
\pgfpathlineto{\pgfqpoint{2.271795in}{3.056437in}}%
\pgfpathlineto{\pgfqpoint{2.275256in}{3.058935in}}%
\pgfpathlineto{\pgfqpoint{2.277729in}{3.058726in}}%
\pgfpathlineto{\pgfqpoint{2.280308in}{3.056737in}}%
\pgfpathlineto{\pgfqpoint{2.283198in}{3.054655in}}%
\pgfpathlineto{\pgfqpoint{2.285670in}{3.054646in}}%
\pgfpathlineto{\pgfqpoint{2.288234in}{3.056352in}}%
\pgfpathlineto{\pgfqpoint{2.291782in}{3.058902in}}%
\pgfpathlineto{\pgfqpoint{2.294255in}{3.058663in}}%
\pgfpathlineto{\pgfqpoint{2.296843in}{3.056636in}}%
\pgfpathlineto{\pgfqpoint{2.299701in}{3.054599in}}%
\pgfpathlineto{\pgfqpoint{2.302174in}{3.054599in}}%
\pgfpathlineto{\pgfqpoint{2.304731in}{3.056309in}}%
\pgfpathlineto{\pgfqpoint{2.308264in}{3.058855in}}%
\pgfpathlineto{\pgfqpoint{2.310737in}{3.058626in}}%
\pgfpathlineto{\pgfqpoint{2.313317in}{3.056612in}}%
\pgfpathlineto{\pgfqpoint{2.316214in}{3.054534in}}%
\pgfpathlineto{\pgfqpoint{2.318686in}{3.054536in}}%
\pgfpathlineto{\pgfqpoint{2.321253in}{3.056257in}}%
\pgfpathlineto{\pgfqpoint{2.324707in}{3.058706in}}%
\pgfpathlineto{\pgfqpoint{2.327180in}{3.058470in}}%
\pgfpathlineto{\pgfqpoint{2.329771in}{3.056445in}}%
\pgfpathlineto{\pgfqpoint{2.332592in}{3.054382in}}%
\pgfpathlineto{\pgfqpoint{2.335065in}{3.054326in}}%
\pgfpathlineto{\pgfqpoint{2.337537in}{3.055895in}}%
\pgfpathlineto{\pgfqpoint{2.341400in}{3.058625in}}%
\pgfpathlineto{\pgfqpoint{2.343873in}{3.058274in}}%
\pgfpathlineto{\pgfqpoint{2.346477in}{3.056127in}}%
\pgfpathlineto{\pgfqpoint{2.349259in}{3.054246in}}%
\pgfpathlineto{\pgfqpoint{2.351731in}{3.054312in}}%
\pgfpathlineto{\pgfqpoint{2.354273in}{3.056075in}}%
\pgfpathlineto{\pgfqpoint{2.357544in}{3.058478in}}%
\pgfpathlineto{\pgfqpoint{2.360017in}{3.058363in}}%
\pgfpathlineto{\pgfqpoint{2.362505in}{3.056566in}}%
\pgfpathlineto{\pgfqpoint{2.365703in}{3.054193in}}%
\pgfpathlineto{\pgfqpoint{2.368175in}{3.054224in}}%
\pgfpathlineto{\pgfqpoint{2.370732in}{3.055966in}}%
\pgfpathlineto{\pgfqpoint{2.374208in}{3.058439in}}%
\pgfpathlineto{\pgfqpoint{2.376680in}{3.058202in}}%
\pgfpathlineto{\pgfqpoint{2.379264in}{3.056177in}}%
\pgfpathlineto{\pgfqpoint{2.382120in}{3.054118in}}%
\pgfpathlineto{\pgfqpoint{2.384593in}{3.054099in}}%
\pgfpathlineto{\pgfqpoint{2.387155in}{3.055796in}}%
\pgfpathlineto{\pgfqpoint{2.390647in}{3.058288in}}%
\pgfpathlineto{\pgfqpoint{2.393119in}{3.058057in}}%
\pgfpathlineto{\pgfqpoint{2.395717in}{3.056033in}}%
\pgfpathlineto{\pgfqpoint{2.398582in}{3.053968in}}%
\pgfpathlineto{\pgfqpoint{2.401054in}{3.053954in}}%
\pgfpathlineto{\pgfqpoint{2.403495in}{3.055536in}}%
\pgfpathlineto{\pgfqpoint{2.407333in}{3.058265in}}%
\pgfpathlineto{\pgfqpoint{2.409806in}{3.057933in}}%
\pgfpathlineto{\pgfqpoint{2.412411in}{3.055801in}}%
\pgfpathlineto{\pgfqpoint{2.415196in}{3.053919in}}%
\pgfpathlineto{\pgfqpoint{2.417669in}{3.053988in}}%
\pgfpathlineto{\pgfqpoint{2.420207in}{3.055752in}}%
\pgfpathlineto{\pgfqpoint{2.423624in}{3.058195in}}%
\pgfpathlineto{\pgfqpoint{2.426097in}{3.057984in}}%
\pgfpathlineto{\pgfqpoint{2.428677in}{3.055993in}}%
\pgfpathlineto{\pgfqpoint{2.431530in}{3.053927in}}%
\pgfpathlineto{\pgfqpoint{2.434002in}{3.053899in}}%
\pgfpathlineto{\pgfqpoint{2.436454in}{3.055478in}}%
\pgfpathlineto{\pgfqpoint{2.440306in}{3.058228in}}%
\pgfpathlineto{\pgfqpoint{2.442779in}{3.057901in}}%
\pgfpathlineto{\pgfqpoint{2.445403in}{3.055753in}}%
\pgfpathlineto{\pgfqpoint{2.448251in}{3.053870in}}%
\pgfpathlineto{\pgfqpoint{2.450724in}{3.053997in}}%
\pgfpathlineto{\pgfqpoint{2.453305in}{3.055864in}}%
\pgfpathlineto{\pgfqpoint{2.456421in}{3.058122in}}%
\pgfpathlineto{\pgfqpoint{2.458894in}{3.058023in}}%
\pgfpathlineto{\pgfqpoint{2.461414in}{3.056213in}}%
\pgfpathlineto{\pgfqpoint{2.464595in}{3.053883in}}%
\pgfpathlineto{\pgfqpoint{2.467068in}{3.053935in}}%
\pgfpathlineto{\pgfqpoint{2.469642in}{3.055715in}}%
\pgfpathlineto{\pgfqpoint{2.473124in}{3.058234in}}%
\pgfpathlineto{\pgfqpoint{2.475596in}{3.058024in}}%
\pgfpathlineto{\pgfqpoint{2.478161in}{3.056050in}}%
\pgfpathlineto{\pgfqpoint{2.481013in}{3.054018in}}%
\pgfpathlineto{\pgfqpoint{2.483486in}{3.054020in}}%
\pgfpathlineto{\pgfqpoint{2.486052in}{3.055742in}}%
\pgfpathlineto{\pgfqpoint{2.489559in}{3.058297in}}%
\pgfpathlineto{\pgfqpoint{2.492032in}{3.058099in}}%
\pgfpathlineto{\pgfqpoint{2.494609in}{3.056124in}}%
\pgfpathlineto{\pgfqpoint{2.497527in}{3.054039in}}%
\pgfpathlineto{\pgfqpoint{2.500000in}{3.054056in}}%
\pgfpathlineto{\pgfqpoint{2.502564in}{3.055791in}}%
\pgfpathlineto{\pgfqpoint{2.505994in}{3.058272in}}%
\pgfpathlineto{\pgfqpoint{2.508466in}{3.058082in}}%
\pgfpathlineto{\pgfqpoint{2.511005in}{3.056158in}}%
\pgfpathlineto{\pgfqpoint{2.513981in}{3.053953in}}%
\pgfpathlineto{\pgfqpoint{2.516454in}{3.053923in}}%
\pgfpathlineto{\pgfqpoint{2.518892in}{3.055487in}}%
\pgfpathlineto{\pgfqpoint{2.522613in}{3.058110in}}%
\pgfpathlineto{\pgfqpoint{2.525086in}{3.057797in}}%
\pgfpathlineto{\pgfqpoint{2.527670in}{3.055697in}}%
\pgfpathlineto{\pgfqpoint{2.530467in}{3.053644in}}%
\pgfpathlineto{\pgfqpoint{2.532940in}{3.053575in}}%
\pgfpathlineto{\pgfqpoint{2.535377in}{3.055101in}}%
\pgfpathlineto{\pgfqpoint{2.539176in}{3.057718in}}%
\pgfpathlineto{\pgfqpoint{2.541649in}{3.057344in}}%
\pgfpathlineto{\pgfqpoint{2.544292in}{3.055140in}}%
\pgfpathlineto{\pgfqpoint{2.547096in}{3.053234in}}%
\pgfpathlineto{\pgfqpoint{2.549568in}{3.053303in}}%
\pgfpathlineto{\pgfqpoint{2.552129in}{3.055091in}}%
\pgfpathlineto{\pgfqpoint{2.555373in}{3.057453in}}%
\pgfpathlineto{\pgfqpoint{2.557845in}{3.057331in}}%
\pgfpathlineto{\pgfqpoint{2.560323in}{3.055543in}}%
\pgfpathlineto{\pgfqpoint{2.563567in}{3.053184in}}%
\pgfpathlineto{\pgfqpoint{2.566040in}{3.053270in}}%
\pgfpathlineto{\pgfqpoint{2.568616in}{3.055093in}}%
\pgfpathlineto{\pgfqpoint{2.571861in}{3.057490in}}%
\pgfpathlineto{\pgfqpoint{2.574333in}{3.057394in}}%
\pgfpathlineto{\pgfqpoint{2.576884in}{3.055554in}}%
\pgfpathlineto{\pgfqpoint{2.579920in}{3.053339in}}%
\pgfpathlineto{\pgfqpoint{2.582393in}{3.053356in}}%
\pgfpathlineto{\pgfqpoint{2.584961in}{3.055098in}}%
\pgfpathlineto{\pgfqpoint{2.588474in}{3.057685in}}%
\pgfpathlineto{\pgfqpoint{2.590947in}{3.057506in}}%
\pgfpathlineto{\pgfqpoint{2.593524in}{3.055554in}}%
\pgfpathlineto{\pgfqpoint{2.596446in}{3.053513in}}%
\pgfpathlineto{\pgfqpoint{2.598919in}{3.053572in}}%
\pgfpathlineto{\pgfqpoint{2.601492in}{3.055363in}}%
\pgfpathlineto{\pgfqpoint{2.604957in}{3.057921in}}%
\pgfpathlineto{\pgfqpoint{2.607430in}{3.057756in}}%
\pgfpathlineto{\pgfqpoint{2.610007in}{3.055818in}}%
\pgfpathlineto{\pgfqpoint{2.612786in}{3.053836in}}%
\pgfpathlineto{\pgfqpoint{2.615259in}{3.053812in}}%
\pgfpathlineto{\pgfqpoint{2.617731in}{3.055415in}}%
\pgfpathlineto{\pgfqpoint{2.621602in}{3.058274in}}%
\pgfpathlineto{\pgfqpoint{2.624075in}{3.058009in}}%
\pgfpathlineto{\pgfqpoint{2.626660in}{3.055965in}}%
\pgfpathlineto{\pgfqpoint{2.629488in}{3.054080in}}%
\pgfpathlineto{\pgfqpoint{2.631961in}{3.054187in}}%
\pgfpathlineto{\pgfqpoint{2.634517in}{3.056005in}}%
\pgfpathlineto{\pgfqpoint{2.637953in}{3.058521in}}%
\pgfpathlineto{\pgfqpoint{2.640426in}{3.058343in}}%
\pgfpathlineto{\pgfqpoint{2.642976in}{3.056415in}}%
\pgfpathlineto{\pgfqpoint{2.645846in}{3.054373in}}%
\pgfpathlineto{\pgfqpoint{2.648319in}{3.054380in}}%
\pgfpathlineto{\pgfqpoint{2.650764in}{3.055985in}}%
\pgfpathlineto{\pgfqpoint{2.654573in}{3.058733in}}%
\pgfpathlineto{\pgfqpoint{2.657046in}{3.058433in}}%
\pgfpathlineto{\pgfqpoint{2.659665in}{3.056317in}}%
\pgfpathlineto{\pgfqpoint{2.662550in}{3.054415in}}%
\pgfpathlineto{\pgfqpoint{2.665023in}{3.054556in}}%
\pgfpathlineto{\pgfqpoint{2.667527in}{3.056354in}}%
\pgfpathlineto{\pgfqpoint{2.670732in}{3.058720in}}%
\pgfpathlineto{\pgfqpoint{2.673204in}{3.058630in}}%
\pgfpathlineto{\pgfqpoint{2.675717in}{3.056832in}}%
\pgfpathlineto{\pgfqpoint{2.678862in}{3.054507in}}%
\pgfpathlineto{\pgfqpoint{2.681335in}{3.054525in}}%
\pgfpathlineto{\pgfqpoint{2.683897in}{3.056259in}}%
\pgfpathlineto{\pgfqpoint{2.687349in}{3.058756in}}%
\pgfpathlineto{\pgfqpoint{2.689822in}{3.058557in}}%
\pgfpathlineto{\pgfqpoint{2.692410in}{3.056574in}}%
\pgfpathlineto{\pgfqpoint{2.695327in}{3.054493in}}%
\pgfpathlineto{\pgfqpoint{2.697800in}{3.054510in}}%
\pgfpathlineto{\pgfqpoint{2.700368in}{3.056248in}}%
\pgfpathlineto{\pgfqpoint{2.703860in}{3.058792in}}%
\pgfpathlineto{\pgfqpoint{2.706333in}{3.058594in}}%
\pgfpathlineto{\pgfqpoint{2.708906in}{3.056625in}}%
\pgfpathlineto{\pgfqpoint{2.711812in}{3.054553in}}%
\pgfpathlineto{\pgfqpoint{2.714284in}{3.054567in}}%
\pgfpathlineto{\pgfqpoint{2.716852in}{3.056302in}}%
\pgfpathlineto{\pgfqpoint{2.720345in}{3.058842in}}%
\pgfpathlineto{\pgfqpoint{2.722818in}{3.058641in}}%
\pgfpathlineto{\pgfqpoint{2.725391in}{3.056668in}}%
\pgfpathlineto{\pgfqpoint{2.728297in}{3.054593in}}%
\pgfpathlineto{\pgfqpoint{2.730770in}{3.054605in}}%
\pgfpathlineto{\pgfqpoint{2.733337in}{3.056337in}}%
\pgfpathlineto{\pgfqpoint{2.736830in}{3.058873in}}%
\pgfpathlineto{\pgfqpoint{2.739303in}{3.058668in}}%
\pgfpathlineto{\pgfqpoint{2.741889in}{3.056679in}}%
\pgfpathlineto{\pgfqpoint{2.744805in}{3.054607in}}%
\pgfpathlineto{\pgfqpoint{2.747277in}{3.054631in}}%
\pgfpathlineto{\pgfqpoint{2.749844in}{3.056374in}}%
\pgfpathlineto{\pgfqpoint{2.753315in}{3.058884in}}%
\pgfpathlineto{\pgfqpoint{2.755787in}{3.058677in}}%
\pgfpathlineto{\pgfqpoint{2.758369in}{3.056689in}}%
\pgfpathlineto{\pgfqpoint{2.761253in}{3.054620in}}%
\pgfpathlineto{\pgfqpoint{2.763726in}{3.054617in}}%
\pgfpathlineto{\pgfqpoint{2.766164in}{3.056206in}}%
\pgfpathlineto{\pgfqpoint{2.770008in}{3.058935in}}%
\pgfpathlineto{\pgfqpoint{2.772481in}{3.058594in}}%
\pgfpathlineto{\pgfqpoint{2.775118in}{3.056423in}}%
\pgfpathlineto{\pgfqpoint{2.778003in}{3.054550in}}%
\pgfpathlineto{\pgfqpoint{2.780476in}{3.054716in}}%
\pgfpathlineto{\pgfqpoint{2.783023in}{3.056584in}}%
\pgfpathlineto{\pgfqpoint{2.786150in}{3.058842in}}%
\pgfpathlineto{\pgfqpoint{2.788623in}{3.058730in}}%
\pgfpathlineto{\pgfqpoint{2.791111in}{3.056935in}}%
\pgfpathlineto{\pgfqpoint{2.794341in}{3.054555in}}%
\pgfpathlineto{\pgfqpoint{2.796814in}{3.054609in}}%
\pgfpathlineto{\pgfqpoint{2.799349in}{3.056349in}}%
\pgfpathlineto{\pgfqpoint{2.802768in}{3.058772in}}%
\pgfpathlineto{\pgfqpoint{2.805241in}{3.058542in}}%
\pgfpathlineto{\pgfqpoint{2.807828in}{3.056526in}}%
\pgfpathlineto{\pgfqpoint{2.810682in}{3.054454in}}%
\pgfpathlineto{\pgfqpoint{2.813155in}{3.054421in}}%
\pgfpathlineto{\pgfqpoint{2.815606in}{3.055992in}}%
\pgfpathlineto{\pgfqpoint{2.819456in}{3.058725in}}%
\pgfpathlineto{\pgfqpoint{2.821929in}{3.058385in}}%
\pgfpathlineto{\pgfqpoint{2.824557in}{3.056219in}}%
\pgfpathlineto{\pgfqpoint{2.827334in}{3.054339in}}%
\pgfpathlineto{\pgfqpoint{2.829807in}{3.054402in}}%
\pgfpathlineto{\pgfqpoint{2.832350in}{3.056162in}}%
\pgfpathlineto{\pgfqpoint{2.835518in}{3.058462in}}%
\pgfpathlineto{\pgfqpoint{2.837991in}{3.058355in}}%
\pgfpathlineto{\pgfqpoint{2.840520in}{3.056523in}}%
\pgfpathlineto{\pgfqpoint{2.843842in}{3.054009in}}%
\pgfpathlineto{\pgfqpoint{2.846314in}{3.054042in}}%
\pgfpathlineto{\pgfqpoint{2.848851in}{3.055766in}}%
\pgfpathlineto{\pgfqpoint{2.852059in}{3.058038in}}%
\pgfpathlineto{\pgfqpoint{2.854531in}{3.057877in}}%
\pgfpathlineto{\pgfqpoint{2.857051in}{3.056008in}}%
\pgfpathlineto{\pgfqpoint{2.860307in}{3.053596in}}%
\pgfpathlineto{\pgfqpoint{2.862780in}{3.053652in}}%
\pgfpathlineto{\pgfqpoint{2.865341in}{3.055426in}}%
\pgfpathlineto{\pgfqpoint{2.868585in}{3.057770in}}%
\pgfpathlineto{\pgfqpoint{2.871057in}{3.057632in}}%
\pgfpathlineto{\pgfqpoint{2.873628in}{3.055729in}}%
\pgfpathlineto{\pgfqpoint{2.876763in}{3.053454in}}%
\pgfpathlineto{\pgfqpoint{2.879236in}{3.053515in}}%
\pgfpathlineto{\pgfqpoint{2.881786in}{3.055284in}}%
\pgfpathlineto{\pgfqpoint{2.885060in}{3.057698in}}%
\pgfpathlineto{\pgfqpoint{2.887533in}{3.057592in}}%
\pgfpathlineto{\pgfqpoint{2.890020in}{3.055809in}}%
\pgfpathlineto{\pgfqpoint{2.893291in}{3.053437in}}%
\pgfpathlineto{\pgfqpoint{2.895764in}{3.053537in}}%
\pgfpathlineto{\pgfqpoint{2.898315in}{3.055347in}}%
\pgfpathlineto{\pgfqpoint{2.901518in}{3.057709in}}%
\pgfpathlineto{\pgfqpoint{2.903991in}{3.057621in}}%
\pgfpathlineto{\pgfqpoint{2.906505in}{3.055828in}}%
\pgfpathlineto{\pgfqpoint{2.909685in}{3.053504in}}%
\pgfpathlineto{\pgfqpoint{2.912157in}{3.053558in}}%
\pgfpathlineto{\pgfqpoint{2.914715in}{3.055328in}}%
\pgfpathlineto{\pgfqpoint{2.918170in}{3.057804in}}%
\pgfpathlineto{\pgfqpoint{2.920643in}{3.057585in}}%
\pgfpathlineto{\pgfqpoint{2.923214in}{3.055598in}}%
\pgfpathlineto{\pgfqpoint{2.926050in}{3.053522in}}%
\pgfpathlineto{\pgfqpoint{2.928522in}{3.053471in}}%
\pgfpathlineto{\pgfqpoint{2.930987in}{3.055041in}}%
\pgfpathlineto{\pgfqpoint{2.934844in}{3.057790in}}%
\pgfpathlineto{\pgfqpoint{2.937317in}{3.057461in}}%
\pgfpathlineto{\pgfqpoint{2.939932in}{3.055324in}}%
\pgfpathlineto{\pgfqpoint{2.942752in}{3.053442in}}%
\pgfpathlineto{\pgfqpoint{2.945225in}{3.053546in}}%
\pgfpathlineto{\pgfqpoint{2.947760in}{3.055342in}}%
\pgfpathlineto{\pgfqpoint{2.950968in}{3.057709in}}%
\pgfpathlineto{\pgfqpoint{2.953441in}{3.057621in}}%
\pgfpathlineto{\pgfqpoint{2.955960in}{3.055823in}}%
\pgfpathlineto{\pgfqpoint{2.959250in}{3.053476in}}%
\pgfpathlineto{\pgfqpoint{2.961722in}{3.053611in}}%
\pgfpathlineto{\pgfqpoint{2.964245in}{3.055425in}}%
\pgfpathlineto{\pgfqpoint{2.967453in}{3.057820in}}%
\pgfpathlineto{\pgfqpoint{2.969926in}{3.057754in}}%
\pgfpathlineto{\pgfqpoint{2.972445in}{3.055977in}}%
\pgfpathlineto{\pgfqpoint{2.975577in}{3.053693in}}%
\pgfpathlineto{\pgfqpoint{2.978050in}{3.053738in}}%
\pgfpathlineto{\pgfqpoint{2.980610in}{3.055498in}}%
\pgfpathlineto{\pgfqpoint{2.984103in}{3.058080in}}%
\pgfpathlineto{\pgfqpoint{2.986575in}{3.057912in}}%
\pgfpathlineto{\pgfqpoint{2.989147in}{3.055978in}}%
\pgfpathlineto{\pgfqpoint{2.992083in}{3.053925in}}%
\pgfpathlineto{\pgfqpoint{2.994556in}{3.053988in}}%
\pgfpathlineto{\pgfqpoint{2.997128in}{3.055779in}}%
\pgfpathlineto{\pgfqpoint{3.000587in}{3.058331in}}%
\pgfpathlineto{\pgfqpoint{3.003060in}{3.058164in}}%
\pgfpathlineto{\pgfqpoint{3.005632in}{3.056230in}}%
\pgfpathlineto{\pgfqpoint{3.008569in}{3.054177in}}%
\pgfpathlineto{\pgfqpoint{3.011042in}{3.054238in}}%
\pgfpathlineto{\pgfqpoint{3.013613in}{3.056026in}}%
\pgfpathlineto{\pgfqpoint{3.017078in}{3.058577in}}%
\pgfpathlineto{\pgfqpoint{3.019551in}{3.058405in}}%
\pgfpathlineto{\pgfqpoint{3.022119in}{3.056466in}}%
\pgfpathlineto{\pgfqpoint{3.025014in}{3.054408in}}%
\pgfpathlineto{\pgfqpoint{3.027486in}{3.054426in}}%
\pgfpathlineto{\pgfqpoint{3.030056in}{3.056166in}}%
\pgfpathlineto{\pgfqpoint{3.033557in}{3.058729in}}%
\pgfpathlineto{\pgfqpoint{3.036030in}{3.058539in}}%
\pgfpathlineto{\pgfqpoint{3.038603in}{3.056577in}}%
\pgfpathlineto{\pgfqpoint{3.041509in}{3.054512in}}%
\pgfpathlineto{\pgfqpoint{3.043982in}{3.054532in}}%
\pgfpathlineto{\pgfqpoint{3.046549in}{3.056273in}}%
\pgfpathlineto{\pgfqpoint{3.050042in}{3.058822in}}%
\pgfpathlineto{\pgfqpoint{3.052515in}{3.058626in}}%
\pgfpathlineto{\pgfqpoint{3.055088in}{3.056660in}}%
\pgfpathlineto{\pgfqpoint{3.057994in}{3.054589in}}%
\pgfpathlineto{\pgfqpoint{3.060466in}{3.054606in}}%
\pgfpathlineto{\pgfqpoint{3.063034in}{3.056342in}}%
\pgfpathlineto{\pgfqpoint{3.066527in}{3.058883in}}%
\pgfpathlineto{\pgfqpoint{3.069000in}{3.058683in}}%
\pgfpathlineto{\pgfqpoint{3.071573in}{3.056711in}}%
\pgfpathlineto{\pgfqpoint{3.074479in}{3.054636in}}%
\pgfpathlineto{\pgfqpoint{3.076951in}{3.054648in}}%
\pgfpathlineto{\pgfqpoint{3.079519in}{3.056380in}}%
\pgfpathlineto{\pgfqpoint{3.083012in}{3.058916in}}%
\pgfpathlineto{\pgfqpoint{3.085484in}{3.058711in}}%
\pgfpathlineto{\pgfqpoint{3.088070in}{3.056721in}}%
\pgfpathlineto{\pgfqpoint{3.090987in}{3.054648in}}%
\pgfpathlineto{\pgfqpoint{3.093460in}{3.054672in}}%
\pgfpathlineto{\pgfqpoint{3.096033in}{3.056421in}}%
\pgfpathlineto{\pgfqpoint{3.099496in}{3.058922in}}%
\pgfpathlineto{\pgfqpoint{3.101969in}{3.058714in}}%
\pgfpathlineto{\pgfqpoint{3.104550in}{3.056725in}}%
\pgfpathlineto{\pgfqpoint{3.107397in}{3.054668in}}%
\pgfpathlineto{\pgfqpoint{3.109870in}{3.054640in}}%
\pgfpathlineto{\pgfqpoint{3.112327in}{3.056222in}}%
\pgfpathlineto{\pgfqpoint{3.116176in}{3.058970in}}%
\pgfpathlineto{\pgfqpoint{3.118649in}{3.058641in}}%
\pgfpathlineto{\pgfqpoint{3.121254in}{3.056509in}}%
\pgfpathlineto{\pgfqpoint{3.124027in}{3.054620in}}%
\pgfpathlineto{\pgfqpoint{3.126499in}{3.054668in}}%
\pgfpathlineto{\pgfqpoint{3.129046in}{3.056416in}}%
\pgfpathlineto{\pgfqpoint{3.132465in}{3.058842in}}%
\pgfpathlineto{\pgfqpoint{3.134938in}{3.058614in}}%
\pgfpathlineto{\pgfqpoint{3.137522in}{3.056603in}}%
\pgfpathlineto{\pgfqpoint{3.140339in}{3.054542in}}%
\pgfpathlineto{\pgfqpoint{3.142811in}{3.054482in}}%
\pgfpathlineto{\pgfqpoint{3.145284in}{3.056048in}}%
\pgfpathlineto{\pgfqpoint{3.149121in}{3.058773in}}%
\pgfpathlineto{\pgfqpoint{3.151594in}{3.058440in}}%
\pgfpathlineto{\pgfqpoint{3.154215in}{3.056289in}}%
\pgfpathlineto{\pgfqpoint{3.157119in}{3.054351in}}%
\pgfpathlineto{\pgfqpoint{3.159592in}{3.054479in}}%
\pgfpathlineto{\pgfqpoint{3.162154in}{3.056325in}}%
\pgfpathlineto{\pgfqpoint{3.165270in}{3.058563in}}%
\pgfpathlineto{\pgfqpoint{3.167742in}{3.058448in}}%
\pgfpathlineto{\pgfqpoint{3.170132in}{3.056755in}}%
\pgfpathlineto{\pgfqpoint{3.170132in}{3.056755in}}%
\pgfusepath{stroke}%
\end{pgfscope}%
\begin{pgfscope}%
\pgfsetrectcap%
\pgfsetmiterjoin%
\pgfsetlinewidth{0.803000pt}%
\definecolor{currentstroke}{rgb}{0.000000,0.000000,0.000000}%
\pgfsetstrokecolor{currentstroke}%
\pgfsetdash{}{0pt}%
\pgfpathmoveto{\pgfqpoint{0.573768in}{2.601404in}}%
\pgfpathlineto{\pgfqpoint{0.573768in}{3.507654in}}%
\pgfusepath{stroke}%
\end{pgfscope}%
\begin{pgfscope}%
\pgfsetrectcap%
\pgfsetmiterjoin%
\pgfsetlinewidth{0.803000pt}%
\definecolor{currentstroke}{rgb}{0.000000,0.000000,0.000000}%
\pgfsetstrokecolor{currentstroke}%
\pgfsetdash{}{0pt}%
\pgfpathmoveto{\pgfqpoint{3.293768in}{2.601404in}}%
\pgfpathlineto{\pgfqpoint{3.293768in}{3.507654in}}%
\pgfusepath{stroke}%
\end{pgfscope}%
\begin{pgfscope}%
\pgfsetrectcap%
\pgfsetmiterjoin%
\pgfsetlinewidth{0.803000pt}%
\definecolor{currentstroke}{rgb}{0.000000,0.000000,0.000000}%
\pgfsetstrokecolor{currentstroke}%
\pgfsetdash{}{0pt}%
\pgfpathmoveto{\pgfqpoint{0.573768in}{2.601404in}}%
\pgfpathlineto{\pgfqpoint{3.293768in}{2.601404in}}%
\pgfusepath{stroke}%
\end{pgfscope}%
\begin{pgfscope}%
\pgfsetrectcap%
\pgfsetmiterjoin%
\pgfsetlinewidth{0.803000pt}%
\definecolor{currentstroke}{rgb}{0.000000,0.000000,0.000000}%
\pgfsetstrokecolor{currentstroke}%
\pgfsetdash{}{0pt}%
\pgfpathmoveto{\pgfqpoint{0.573768in}{3.507654in}}%
\pgfpathlineto{\pgfqpoint{3.293768in}{3.507654in}}%
\pgfusepath{stroke}%
\end{pgfscope}%
\begin{pgfscope}%
\definecolor{textcolor}{rgb}{0.000000,0.000000,0.000000}%
\pgfsetstrokecolor{textcolor}%
\pgfsetfillcolor{textcolor}%
\pgftext[x=0.628168in,y=3.444217in,left,top]{\color{textcolor}{\rmfamily\fontsize{8.000000}{9.600000}\selectfont\catcode`\^=\active\def^{\ifmmode\sp\else\^{}\fi}\catcode`\%=\active\def%{\%}(a)}}%
\end{pgfscope}%
\begin{pgfscope}%
\pgfsetbuttcap%
\pgfsetmiterjoin%
\definecolor{currentfill}{rgb}{1.000000,1.000000,1.000000}%
\pgfsetfillcolor{currentfill}%
\pgfsetlinewidth{0.000000pt}%
\definecolor{currentstroke}{rgb}{0.000000,0.000000,0.000000}%
\pgfsetstrokecolor{currentstroke}%
\pgfsetstrokeopacity{0.000000}%
\pgfsetdash{}{0pt}%
\pgfpathmoveto{\pgfqpoint{0.573768in}{1.532029in}}%
\pgfpathlineto{\pgfqpoint{3.293768in}{1.532029in}}%
\pgfpathlineto{\pgfqpoint{3.293768in}{2.438279in}}%
\pgfpathlineto{\pgfqpoint{0.573768in}{2.438279in}}%
\pgfpathlineto{\pgfqpoint{0.573768in}{1.532029in}}%
\pgfpathclose%
\pgfusepath{fill}%
\end{pgfscope}%
\begin{pgfscope}%
\pgfpathrectangle{\pgfqpoint{0.573768in}{1.532029in}}{\pgfqpoint{2.720000in}{0.906250in}}%
\pgfusepath{clip}%
\pgfsetbuttcap%
\pgfsetroundjoin%
\pgfsetlinewidth{0.501875pt}%
\definecolor{currentstroke}{rgb}{0.690196,0.690196,0.690196}%
\pgfsetstrokecolor{currentstroke}%
\pgfsetstrokeopacity{0.250000}%
\pgfsetdash{{1.850000pt}{0.800000pt}}{0.000000pt}%
\pgfpathmoveto{\pgfqpoint{0.697405in}{1.532029in}}%
\pgfpathlineto{\pgfqpoint{0.697405in}{2.438279in}}%
\pgfusepath{stroke}%
\end{pgfscope}%
\begin{pgfscope}%
\pgfsetbuttcap%
\pgfsetroundjoin%
\definecolor{currentfill}{rgb}{0.000000,0.000000,0.000000}%
\pgfsetfillcolor{currentfill}%
\pgfsetlinewidth{0.803000pt}%
\definecolor{currentstroke}{rgb}{0.000000,0.000000,0.000000}%
\pgfsetstrokecolor{currentstroke}%
\pgfsetdash{}{0pt}%
\pgfsys@defobject{currentmarker}{\pgfqpoint{0.000000in}{-0.048611in}}{\pgfqpoint{0.000000in}{0.000000in}}{%
\pgfpathmoveto{\pgfqpoint{0.000000in}{0.000000in}}%
\pgfpathlineto{\pgfqpoint{0.000000in}{-0.048611in}}%
\pgfusepath{stroke,fill}%
}%
\begin{pgfscope}%
\pgfsys@transformshift{0.697405in}{1.532029in}%
\pgfsys@useobject{currentmarker}{}%
\end{pgfscope}%
\end{pgfscope}%
\begin{pgfscope}%
\pgfpathrectangle{\pgfqpoint{0.573768in}{1.532029in}}{\pgfqpoint{2.720000in}{0.906250in}}%
\pgfusepath{clip}%
\pgfsetbuttcap%
\pgfsetroundjoin%
\pgfsetlinewidth{0.501875pt}%
\definecolor{currentstroke}{rgb}{0.690196,0.690196,0.690196}%
\pgfsetstrokecolor{currentstroke}%
\pgfsetstrokeopacity{0.250000}%
\pgfsetdash{{1.850000pt}{0.800000pt}}{0.000000pt}%
\pgfpathmoveto{\pgfqpoint{1.109526in}{1.532029in}}%
\pgfpathlineto{\pgfqpoint{1.109526in}{2.438279in}}%
\pgfusepath{stroke}%
\end{pgfscope}%
\begin{pgfscope}%
\pgfsetbuttcap%
\pgfsetroundjoin%
\definecolor{currentfill}{rgb}{0.000000,0.000000,0.000000}%
\pgfsetfillcolor{currentfill}%
\pgfsetlinewidth{0.803000pt}%
\definecolor{currentstroke}{rgb}{0.000000,0.000000,0.000000}%
\pgfsetstrokecolor{currentstroke}%
\pgfsetdash{}{0pt}%
\pgfsys@defobject{currentmarker}{\pgfqpoint{0.000000in}{-0.048611in}}{\pgfqpoint{0.000000in}{0.000000in}}{%
\pgfpathmoveto{\pgfqpoint{0.000000in}{0.000000in}}%
\pgfpathlineto{\pgfqpoint{0.000000in}{-0.048611in}}%
\pgfusepath{stroke,fill}%
}%
\begin{pgfscope}%
\pgfsys@transformshift{1.109526in}{1.532029in}%
\pgfsys@useobject{currentmarker}{}%
\end{pgfscope}%
\end{pgfscope}%
\begin{pgfscope}%
\pgfpathrectangle{\pgfqpoint{0.573768in}{1.532029in}}{\pgfqpoint{2.720000in}{0.906250in}}%
\pgfusepath{clip}%
\pgfsetbuttcap%
\pgfsetroundjoin%
\pgfsetlinewidth{0.501875pt}%
\definecolor{currentstroke}{rgb}{0.690196,0.690196,0.690196}%
\pgfsetstrokecolor{currentstroke}%
\pgfsetstrokeopacity{0.250000}%
\pgfsetdash{{1.850000pt}{0.800000pt}}{0.000000pt}%
\pgfpathmoveto{\pgfqpoint{1.521647in}{1.532029in}}%
\pgfpathlineto{\pgfqpoint{1.521647in}{2.438279in}}%
\pgfusepath{stroke}%
\end{pgfscope}%
\begin{pgfscope}%
\pgfsetbuttcap%
\pgfsetroundjoin%
\definecolor{currentfill}{rgb}{0.000000,0.000000,0.000000}%
\pgfsetfillcolor{currentfill}%
\pgfsetlinewidth{0.803000pt}%
\definecolor{currentstroke}{rgb}{0.000000,0.000000,0.000000}%
\pgfsetstrokecolor{currentstroke}%
\pgfsetdash{}{0pt}%
\pgfsys@defobject{currentmarker}{\pgfqpoint{0.000000in}{-0.048611in}}{\pgfqpoint{0.000000in}{0.000000in}}{%
\pgfpathmoveto{\pgfqpoint{0.000000in}{0.000000in}}%
\pgfpathlineto{\pgfqpoint{0.000000in}{-0.048611in}}%
\pgfusepath{stroke,fill}%
}%
\begin{pgfscope}%
\pgfsys@transformshift{1.521647in}{1.532029in}%
\pgfsys@useobject{currentmarker}{}%
\end{pgfscope}%
\end{pgfscope}%
\begin{pgfscope}%
\pgfpathrectangle{\pgfqpoint{0.573768in}{1.532029in}}{\pgfqpoint{2.720000in}{0.906250in}}%
\pgfusepath{clip}%
\pgfsetbuttcap%
\pgfsetroundjoin%
\pgfsetlinewidth{0.501875pt}%
\definecolor{currentstroke}{rgb}{0.690196,0.690196,0.690196}%
\pgfsetstrokecolor{currentstroke}%
\pgfsetstrokeopacity{0.250000}%
\pgfsetdash{{1.850000pt}{0.800000pt}}{0.000000pt}%
\pgfpathmoveto{\pgfqpoint{1.933768in}{1.532029in}}%
\pgfpathlineto{\pgfqpoint{1.933768in}{2.438279in}}%
\pgfusepath{stroke}%
\end{pgfscope}%
\begin{pgfscope}%
\pgfsetbuttcap%
\pgfsetroundjoin%
\definecolor{currentfill}{rgb}{0.000000,0.000000,0.000000}%
\pgfsetfillcolor{currentfill}%
\pgfsetlinewidth{0.803000pt}%
\definecolor{currentstroke}{rgb}{0.000000,0.000000,0.000000}%
\pgfsetstrokecolor{currentstroke}%
\pgfsetdash{}{0pt}%
\pgfsys@defobject{currentmarker}{\pgfqpoint{0.000000in}{-0.048611in}}{\pgfqpoint{0.000000in}{0.000000in}}{%
\pgfpathmoveto{\pgfqpoint{0.000000in}{0.000000in}}%
\pgfpathlineto{\pgfqpoint{0.000000in}{-0.048611in}}%
\pgfusepath{stroke,fill}%
}%
\begin{pgfscope}%
\pgfsys@transformshift{1.933768in}{1.532029in}%
\pgfsys@useobject{currentmarker}{}%
\end{pgfscope}%
\end{pgfscope}%
\begin{pgfscope}%
\pgfpathrectangle{\pgfqpoint{0.573768in}{1.532029in}}{\pgfqpoint{2.720000in}{0.906250in}}%
\pgfusepath{clip}%
\pgfsetbuttcap%
\pgfsetroundjoin%
\pgfsetlinewidth{0.501875pt}%
\definecolor{currentstroke}{rgb}{0.690196,0.690196,0.690196}%
\pgfsetstrokecolor{currentstroke}%
\pgfsetstrokeopacity{0.250000}%
\pgfsetdash{{1.850000pt}{0.800000pt}}{0.000000pt}%
\pgfpathmoveto{\pgfqpoint{2.345890in}{1.532029in}}%
\pgfpathlineto{\pgfqpoint{2.345890in}{2.438279in}}%
\pgfusepath{stroke}%
\end{pgfscope}%
\begin{pgfscope}%
\pgfsetbuttcap%
\pgfsetroundjoin%
\definecolor{currentfill}{rgb}{0.000000,0.000000,0.000000}%
\pgfsetfillcolor{currentfill}%
\pgfsetlinewidth{0.803000pt}%
\definecolor{currentstroke}{rgb}{0.000000,0.000000,0.000000}%
\pgfsetstrokecolor{currentstroke}%
\pgfsetdash{}{0pt}%
\pgfsys@defobject{currentmarker}{\pgfqpoint{0.000000in}{-0.048611in}}{\pgfqpoint{0.000000in}{0.000000in}}{%
\pgfpathmoveto{\pgfqpoint{0.000000in}{0.000000in}}%
\pgfpathlineto{\pgfqpoint{0.000000in}{-0.048611in}}%
\pgfusepath{stroke,fill}%
}%
\begin{pgfscope}%
\pgfsys@transformshift{2.345890in}{1.532029in}%
\pgfsys@useobject{currentmarker}{}%
\end{pgfscope}%
\end{pgfscope}%
\begin{pgfscope}%
\pgfpathrectangle{\pgfqpoint{0.573768in}{1.532029in}}{\pgfqpoint{2.720000in}{0.906250in}}%
\pgfusepath{clip}%
\pgfsetbuttcap%
\pgfsetroundjoin%
\pgfsetlinewidth{0.501875pt}%
\definecolor{currentstroke}{rgb}{0.690196,0.690196,0.690196}%
\pgfsetstrokecolor{currentstroke}%
\pgfsetstrokeopacity{0.250000}%
\pgfsetdash{{1.850000pt}{0.800000pt}}{0.000000pt}%
\pgfpathmoveto{\pgfqpoint{2.758011in}{1.532029in}}%
\pgfpathlineto{\pgfqpoint{2.758011in}{2.438279in}}%
\pgfusepath{stroke}%
\end{pgfscope}%
\begin{pgfscope}%
\pgfsetbuttcap%
\pgfsetroundjoin%
\definecolor{currentfill}{rgb}{0.000000,0.000000,0.000000}%
\pgfsetfillcolor{currentfill}%
\pgfsetlinewidth{0.803000pt}%
\definecolor{currentstroke}{rgb}{0.000000,0.000000,0.000000}%
\pgfsetstrokecolor{currentstroke}%
\pgfsetdash{}{0pt}%
\pgfsys@defobject{currentmarker}{\pgfqpoint{0.000000in}{-0.048611in}}{\pgfqpoint{0.000000in}{0.000000in}}{%
\pgfpathmoveto{\pgfqpoint{0.000000in}{0.000000in}}%
\pgfpathlineto{\pgfqpoint{0.000000in}{-0.048611in}}%
\pgfusepath{stroke,fill}%
}%
\begin{pgfscope}%
\pgfsys@transformshift{2.758011in}{1.532029in}%
\pgfsys@useobject{currentmarker}{}%
\end{pgfscope}%
\end{pgfscope}%
\begin{pgfscope}%
\pgfpathrectangle{\pgfqpoint{0.573768in}{1.532029in}}{\pgfqpoint{2.720000in}{0.906250in}}%
\pgfusepath{clip}%
\pgfsetbuttcap%
\pgfsetroundjoin%
\pgfsetlinewidth{0.501875pt}%
\definecolor{currentstroke}{rgb}{0.690196,0.690196,0.690196}%
\pgfsetstrokecolor{currentstroke}%
\pgfsetstrokeopacity{0.250000}%
\pgfsetdash{{1.850000pt}{0.800000pt}}{0.000000pt}%
\pgfpathmoveto{\pgfqpoint{3.170132in}{1.532029in}}%
\pgfpathlineto{\pgfqpoint{3.170132in}{2.438279in}}%
\pgfusepath{stroke}%
\end{pgfscope}%
\begin{pgfscope}%
\pgfsetbuttcap%
\pgfsetroundjoin%
\definecolor{currentfill}{rgb}{0.000000,0.000000,0.000000}%
\pgfsetfillcolor{currentfill}%
\pgfsetlinewidth{0.803000pt}%
\definecolor{currentstroke}{rgb}{0.000000,0.000000,0.000000}%
\pgfsetstrokecolor{currentstroke}%
\pgfsetdash{}{0pt}%
\pgfsys@defobject{currentmarker}{\pgfqpoint{0.000000in}{-0.048611in}}{\pgfqpoint{0.000000in}{0.000000in}}{%
\pgfpathmoveto{\pgfqpoint{0.000000in}{0.000000in}}%
\pgfpathlineto{\pgfqpoint{0.000000in}{-0.048611in}}%
\pgfusepath{stroke,fill}%
}%
\begin{pgfscope}%
\pgfsys@transformshift{3.170132in}{1.532029in}%
\pgfsys@useobject{currentmarker}{}%
\end{pgfscope}%
\end{pgfscope}%
\begin{pgfscope}%
\pgfpathrectangle{\pgfqpoint{0.573768in}{1.532029in}}{\pgfqpoint{2.720000in}{0.906250in}}%
\pgfusepath{clip}%
\pgfsetbuttcap%
\pgfsetroundjoin%
\pgfsetlinewidth{0.501875pt}%
\definecolor{currentstroke}{rgb}{0.690196,0.690196,0.690196}%
\pgfsetstrokecolor{currentstroke}%
\pgfsetstrokeopacity{0.250000}%
\pgfsetdash{{1.850000pt}{0.800000pt}}{0.000000pt}%
\pgfpathmoveto{\pgfqpoint{0.573768in}{1.532029in}}%
\pgfpathlineto{\pgfqpoint{3.293768in}{1.532029in}}%
\pgfusepath{stroke}%
\end{pgfscope}%
\begin{pgfscope}%
\pgfsetbuttcap%
\pgfsetroundjoin%
\definecolor{currentfill}{rgb}{0.000000,0.000000,0.000000}%
\pgfsetfillcolor{currentfill}%
\pgfsetlinewidth{0.803000pt}%
\definecolor{currentstroke}{rgb}{0.000000,0.000000,0.000000}%
\pgfsetstrokecolor{currentstroke}%
\pgfsetdash{}{0pt}%
\pgfsys@defobject{currentmarker}{\pgfqpoint{-0.048611in}{0.000000in}}{\pgfqpoint{-0.000000in}{0.000000in}}{%
\pgfpathmoveto{\pgfqpoint{-0.000000in}{0.000000in}}%
\pgfpathlineto{\pgfqpoint{-0.048611in}{0.000000in}}%
\pgfusepath{stroke,fill}%
}%
\begin{pgfscope}%
\pgfsys@transformshift{0.573768in}{1.532029in}%
\pgfsys@useobject{currentmarker}{}%
\end{pgfscope}%
\end{pgfscope}%
\begin{pgfscope}%
\definecolor{textcolor}{rgb}{0.000000,0.000000,0.000000}%
\pgfsetstrokecolor{textcolor}%
\pgfsetfillcolor{textcolor}%
\pgftext[x=0.417518in, y=1.493449in, left, base]{\color{textcolor}{\rmfamily\fontsize{8.000000}{9.600000}\selectfont\catcode`\^=\active\def^{\ifmmode\sp\else\^{}\fi}\catcode`\%=\active\def%{\%}$\mathdefault{0}$}}%
\end{pgfscope}%
\begin{pgfscope}%
\pgfpathrectangle{\pgfqpoint{0.573768in}{1.532029in}}{\pgfqpoint{2.720000in}{0.906250in}}%
\pgfusepath{clip}%
\pgfsetbuttcap%
\pgfsetroundjoin%
\pgfsetlinewidth{0.501875pt}%
\definecolor{currentstroke}{rgb}{0.690196,0.690196,0.690196}%
\pgfsetstrokecolor{currentstroke}%
\pgfsetstrokeopacity{0.250000}%
\pgfsetdash{{1.850000pt}{0.800000pt}}{0.000000pt}%
\pgfpathmoveto{\pgfqpoint{0.573768in}{1.834113in}}%
\pgfpathlineto{\pgfqpoint{3.293768in}{1.834113in}}%
\pgfusepath{stroke}%
\end{pgfscope}%
\begin{pgfscope}%
\pgfsetbuttcap%
\pgfsetroundjoin%
\definecolor{currentfill}{rgb}{0.000000,0.000000,0.000000}%
\pgfsetfillcolor{currentfill}%
\pgfsetlinewidth{0.803000pt}%
\definecolor{currentstroke}{rgb}{0.000000,0.000000,0.000000}%
\pgfsetstrokecolor{currentstroke}%
\pgfsetdash{}{0pt}%
\pgfsys@defobject{currentmarker}{\pgfqpoint{-0.048611in}{0.000000in}}{\pgfqpoint{-0.000000in}{0.000000in}}{%
\pgfpathmoveto{\pgfqpoint{-0.000000in}{0.000000in}}%
\pgfpathlineto{\pgfqpoint{-0.048611in}{0.000000in}}%
\pgfusepath{stroke,fill}%
}%
\begin{pgfscope}%
\pgfsys@transformshift{0.573768in}{1.834113in}%
\pgfsys@useobject{currentmarker}{}%
\end{pgfscope}%
\end{pgfscope}%
\begin{pgfscope}%
\definecolor{textcolor}{rgb}{0.000000,0.000000,0.000000}%
\pgfsetstrokecolor{textcolor}%
\pgfsetfillcolor{textcolor}%
\pgftext[x=0.417518in, y=1.795533in, left, base]{\color{textcolor}{\rmfamily\fontsize{8.000000}{9.600000}\selectfont\catcode`\^=\active\def^{\ifmmode\sp\else\^{}\fi}\catcode`\%=\active\def%{\%}$\mathdefault{5}$}}%
\end{pgfscope}%
\begin{pgfscope}%
\pgfpathrectangle{\pgfqpoint{0.573768in}{1.532029in}}{\pgfqpoint{2.720000in}{0.906250in}}%
\pgfusepath{clip}%
\pgfsetbuttcap%
\pgfsetroundjoin%
\pgfsetlinewidth{0.501875pt}%
\definecolor{currentstroke}{rgb}{0.690196,0.690196,0.690196}%
\pgfsetstrokecolor{currentstroke}%
\pgfsetstrokeopacity{0.250000}%
\pgfsetdash{{1.850000pt}{0.800000pt}}{0.000000pt}%
\pgfpathmoveto{\pgfqpoint{0.573768in}{2.136196in}}%
\pgfpathlineto{\pgfqpoint{3.293768in}{2.136196in}}%
\pgfusepath{stroke}%
\end{pgfscope}%
\begin{pgfscope}%
\pgfsetbuttcap%
\pgfsetroundjoin%
\definecolor{currentfill}{rgb}{0.000000,0.000000,0.000000}%
\pgfsetfillcolor{currentfill}%
\pgfsetlinewidth{0.803000pt}%
\definecolor{currentstroke}{rgb}{0.000000,0.000000,0.000000}%
\pgfsetstrokecolor{currentstroke}%
\pgfsetdash{}{0pt}%
\pgfsys@defobject{currentmarker}{\pgfqpoint{-0.048611in}{0.000000in}}{\pgfqpoint{-0.000000in}{0.000000in}}{%
\pgfpathmoveto{\pgfqpoint{-0.000000in}{0.000000in}}%
\pgfpathlineto{\pgfqpoint{-0.048611in}{0.000000in}}%
\pgfusepath{stroke,fill}%
}%
\begin{pgfscope}%
\pgfsys@transformshift{0.573768in}{2.136196in}%
\pgfsys@useobject{currentmarker}{}%
\end{pgfscope}%
\end{pgfscope}%
\begin{pgfscope}%
\definecolor{textcolor}{rgb}{0.000000,0.000000,0.000000}%
\pgfsetstrokecolor{textcolor}%
\pgfsetfillcolor{textcolor}%
\pgftext[x=0.358489in, y=2.097616in, left, base]{\color{textcolor}{\rmfamily\fontsize{8.000000}{9.600000}\selectfont\catcode`\^=\active\def^{\ifmmode\sp\else\^{}\fi}\catcode`\%=\active\def%{\%}$\mathdefault{10}$}}%
\end{pgfscope}%
\begin{pgfscope}%
\pgfpathrectangle{\pgfqpoint{0.573768in}{1.532029in}}{\pgfqpoint{2.720000in}{0.906250in}}%
\pgfusepath{clip}%
\pgfsetbuttcap%
\pgfsetroundjoin%
\pgfsetlinewidth{0.501875pt}%
\definecolor{currentstroke}{rgb}{0.690196,0.690196,0.690196}%
\pgfsetstrokecolor{currentstroke}%
\pgfsetstrokeopacity{0.250000}%
\pgfsetdash{{1.850000pt}{0.800000pt}}{0.000000pt}%
\pgfpathmoveto{\pgfqpoint{0.573768in}{2.438279in}}%
\pgfpathlineto{\pgfqpoint{3.293768in}{2.438279in}}%
\pgfusepath{stroke}%
\end{pgfscope}%
\begin{pgfscope}%
\pgfsetbuttcap%
\pgfsetroundjoin%
\definecolor{currentfill}{rgb}{0.000000,0.000000,0.000000}%
\pgfsetfillcolor{currentfill}%
\pgfsetlinewidth{0.803000pt}%
\definecolor{currentstroke}{rgb}{0.000000,0.000000,0.000000}%
\pgfsetstrokecolor{currentstroke}%
\pgfsetdash{}{0pt}%
\pgfsys@defobject{currentmarker}{\pgfqpoint{-0.048611in}{0.000000in}}{\pgfqpoint{-0.000000in}{0.000000in}}{%
\pgfpathmoveto{\pgfqpoint{-0.000000in}{0.000000in}}%
\pgfpathlineto{\pgfqpoint{-0.048611in}{0.000000in}}%
\pgfusepath{stroke,fill}%
}%
\begin{pgfscope}%
\pgfsys@transformshift{0.573768in}{2.438279in}%
\pgfsys@useobject{currentmarker}{}%
\end{pgfscope}%
\end{pgfscope}%
\begin{pgfscope}%
\definecolor{textcolor}{rgb}{0.000000,0.000000,0.000000}%
\pgfsetstrokecolor{textcolor}%
\pgfsetfillcolor{textcolor}%
\pgftext[x=0.358489in, y=2.399699in, left, base]{\color{textcolor}{\rmfamily\fontsize{8.000000}{9.600000}\selectfont\catcode`\^=\active\def^{\ifmmode\sp\else\^{}\fi}\catcode`\%=\active\def%{\%}$\mathdefault{15}$}}%
\end{pgfscope}%
\begin{pgfscope}%
\definecolor{textcolor}{rgb}{0.000000,0.000000,0.000000}%
\pgfsetstrokecolor{textcolor}%
\pgfsetfillcolor{textcolor}%
\pgftext[x=0.302933in,y=1.985154in,,bottom,rotate=90.000000]{\color{textcolor}{\rmfamily\fontsize{8.000000}{9.600000}\selectfont\catcode`\^=\active\def^{\ifmmode\sp\else\^{}\fi}\catcode`\%=\active\def%{\%}$i_L$ (A)}}%
\end{pgfscope}%
\begin{pgfscope}%
\pgfpathrectangle{\pgfqpoint{0.573768in}{1.532029in}}{\pgfqpoint{2.720000in}{0.906250in}}%
\pgfusepath{clip}%
\pgfsetrectcap%
\pgfsetroundjoin%
\pgfsetlinewidth{1.505625pt}%
\definecolor{currentstroke}{rgb}{0.835294,0.368627,0.000000}%
\pgfsetstrokecolor{currentstroke}%
\pgfsetdash{}{0pt}%
\pgfpathmoveto{\pgfqpoint{0.697405in}{1.532029in}}%
\pgfpathlineto{\pgfqpoint{0.697771in}{1.532717in}}%
\pgfpathlineto{\pgfqpoint{0.700757in}{1.619744in}}%
\pgfpathlineto{\pgfqpoint{0.706154in}{1.774080in}}%
\pgfpathlineto{\pgfqpoint{0.706687in}{1.773572in}}%
\pgfpathlineto{\pgfqpoint{0.710643in}{1.768563in}}%
\pgfpathlineto{\pgfqpoint{0.714214in}{1.762215in}}%
\pgfpathlineto{\pgfqpoint{0.714304in}{1.763621in}}%
\pgfpathlineto{\pgfqpoint{0.718956in}{1.889950in}}%
\pgfpathlineto{\pgfqpoint{0.722634in}{1.984773in}}%
\pgfpathlineto{\pgfqpoint{0.723147in}{1.982843in}}%
\pgfpathlineto{\pgfqpoint{0.728093in}{1.960855in}}%
\pgfpathlineto{\pgfqpoint{0.730704in}{1.946861in}}%
\pgfpathlineto{\pgfqpoint{0.730909in}{1.950380in}}%
\pgfpathlineto{\pgfqpoint{0.739056in}{2.135271in}}%
\pgfpathlineto{\pgfqpoint{0.739329in}{2.133518in}}%
\pgfpathlineto{\pgfqpoint{0.745570in}{2.076344in}}%
\pgfpathlineto{\pgfqpoint{0.747214in}{2.059472in}}%
\pgfpathlineto{\pgfqpoint{0.747499in}{2.063968in}}%
\pgfpathlineto{\pgfqpoint{0.755541in}{2.205138in}}%
\pgfpathlineto{\pgfqpoint{0.755665in}{2.203937in}}%
\pgfpathlineto{\pgfqpoint{0.763475in}{2.089112in}}%
\pgfpathlineto{\pgfqpoint{0.763702in}{2.085709in}}%
\pgfpathlineto{\pgfqpoint{0.764041in}{2.089433in}}%
\pgfpathlineto{\pgfqpoint{0.772026in}{2.186711in}}%
\pgfpathlineto{\pgfqpoint{0.773278in}{2.163591in}}%
\pgfpathlineto{\pgfqpoint{0.780210in}{2.025073in}}%
\pgfpathlineto{\pgfqpoint{0.780951in}{2.030779in}}%
\pgfpathlineto{\pgfqpoint{0.788197in}{2.085257in}}%
\pgfpathlineto{\pgfqpoint{0.788426in}{2.086647in}}%
\pgfpathlineto{\pgfqpoint{0.788651in}{2.083753in}}%
\pgfpathlineto{\pgfqpoint{0.796755in}{1.890595in}}%
\pgfpathlineto{\pgfqpoint{0.797498in}{1.893867in}}%
\pgfpathlineto{\pgfqpoint{0.804256in}{1.921616in}}%
\pgfpathlineto{\pgfqpoint{0.804910in}{1.923816in}}%
\pgfpathlineto{\pgfqpoint{0.805066in}{1.922098in}}%
\pgfpathlineto{\pgfqpoint{0.813269in}{1.706203in}}%
\pgfpathlineto{\pgfqpoint{0.813943in}{1.707871in}}%
\pgfpathlineto{\pgfqpoint{0.821353in}{1.725806in}}%
\pgfpathlineto{\pgfqpoint{0.821435in}{1.725666in}}%
\pgfpathlineto{\pgfqpoint{0.822612in}{1.694252in}}%
\pgfpathlineto{\pgfqpoint{0.828771in}{1.529553in}}%
\pgfpathlineto{\pgfqpoint{0.829333in}{1.531139in}}%
\pgfpathlineto{\pgfqpoint{0.832814in}{1.539054in}}%
\pgfpathlineto{\pgfqpoint{0.837880in}{1.552926in}}%
\pgfpathlineto{\pgfqpoint{0.838032in}{1.551413in}}%
\pgfpathlineto{\pgfqpoint{0.839056in}{1.530738in}}%
\pgfpathlineto{\pgfqpoint{0.839446in}{1.532043in}}%
\pgfpathlineto{\pgfqpoint{0.839772in}{1.532640in}}%
\pgfpathlineto{\pgfqpoint{0.840102in}{1.531875in}}%
\pgfpathlineto{\pgfqpoint{0.840417in}{1.531323in}}%
\pgfpathlineto{\pgfqpoint{0.840747in}{1.532042in}}%
\pgfpathlineto{\pgfqpoint{0.841076in}{1.532623in}}%
\pgfpathlineto{\pgfqpoint{0.841406in}{1.531853in}}%
\pgfpathlineto{\pgfqpoint{0.841722in}{1.531343in}}%
\pgfpathlineto{\pgfqpoint{0.842051in}{1.532071in}}%
\pgfpathlineto{\pgfqpoint{0.842381in}{1.532605in}}%
\pgfpathlineto{\pgfqpoint{0.842711in}{1.531811in}}%
\pgfpathlineto{\pgfqpoint{0.843027in}{1.531369in}}%
\pgfpathlineto{\pgfqpoint{0.843357in}{1.532118in}}%
\pgfpathlineto{\pgfqpoint{0.843687in}{1.532580in}}%
\pgfpathlineto{\pgfqpoint{0.844017in}{1.531752in}}%
\pgfpathlineto{\pgfqpoint{0.844335in}{1.531406in}}%
\pgfpathlineto{\pgfqpoint{0.844665in}{1.532182in}}%
\pgfpathlineto{\pgfqpoint{0.844994in}{1.532543in}}%
\pgfpathlineto{\pgfqpoint{0.845159in}{1.532185in}}%
\pgfpathlineto{\pgfqpoint{0.846139in}{1.531257in}}%
\pgfpathlineto{\pgfqpoint{0.846222in}{1.531343in}}%
\pgfpathlineto{\pgfqpoint{0.854362in}{1.560947in}}%
\pgfpathlineto{\pgfqpoint{0.854585in}{1.557759in}}%
\pgfpathlineto{\pgfqpoint{0.855961in}{1.530884in}}%
\pgfpathlineto{\pgfqpoint{0.856241in}{1.531742in}}%
\pgfpathlineto{\pgfqpoint{0.856570in}{1.532970in}}%
\pgfpathlineto{\pgfqpoint{0.856900in}{1.532137in}}%
\pgfpathlineto{\pgfqpoint{0.857226in}{1.530932in}}%
\pgfpathlineto{\pgfqpoint{0.857543in}{1.531885in}}%
\pgfpathlineto{\pgfqpoint{0.857873in}{1.532986in}}%
\pgfpathlineto{\pgfqpoint{0.858202in}{1.531941in}}%
\pgfpathlineto{\pgfqpoint{0.858517in}{1.530970in}}%
\pgfpathlineto{\pgfqpoint{0.858846in}{1.532055in}}%
\pgfpathlineto{\pgfqpoint{0.859176in}{1.532972in}}%
\pgfpathlineto{\pgfqpoint{0.859506in}{1.531717in}}%
\pgfpathlineto{\pgfqpoint{0.859796in}{1.531013in}}%
\pgfpathlineto{\pgfqpoint{0.860126in}{1.532156in}}%
\pgfpathlineto{\pgfqpoint{0.860455in}{1.532932in}}%
\pgfpathlineto{\pgfqpoint{0.860620in}{1.532440in}}%
\pgfpathlineto{\pgfqpoint{0.861079in}{1.531082in}}%
\pgfpathlineto{\pgfqpoint{0.861409in}{1.532289in}}%
\pgfpathlineto{\pgfqpoint{0.861739in}{1.532860in}}%
\pgfpathlineto{\pgfqpoint{0.861903in}{1.532260in}}%
\pgfpathlineto{\pgfqpoint{0.862621in}{1.530864in}}%
\pgfpathlineto{\pgfqpoint{0.862806in}{1.531488in}}%
\pgfpathlineto{\pgfqpoint{0.870852in}{1.568202in}}%
\pgfpathlineto{\pgfqpoint{0.871071in}{1.565208in}}%
\pgfpathlineto{\pgfqpoint{0.872751in}{1.531194in}}%
\pgfpathlineto{\pgfqpoint{0.873020in}{1.531847in}}%
\pgfpathlineto{\pgfqpoint{0.873348in}{1.532723in}}%
\pgfpathlineto{\pgfqpoint{0.873677in}{1.531895in}}%
\pgfpathlineto{\pgfqpoint{0.873970in}{1.531231in}}%
\pgfpathlineto{\pgfqpoint{0.874297in}{1.532100in}}%
\pgfpathlineto{\pgfqpoint{0.874619in}{1.532690in}}%
\pgfpathlineto{\pgfqpoint{0.874949in}{1.531643in}}%
\pgfpathlineto{\pgfqpoint{0.875247in}{1.531330in}}%
\pgfpathlineto{\pgfqpoint{0.875412in}{1.531783in}}%
\pgfpathlineto{\pgfqpoint{0.875738in}{1.532676in}}%
\pgfpathlineto{\pgfqpoint{0.876067in}{1.531977in}}%
\pgfpathlineto{\pgfqpoint{0.876383in}{1.531266in}}%
\pgfpathlineto{\pgfqpoint{0.876710in}{1.532110in}}%
\pgfpathlineto{\pgfqpoint{0.877030in}{1.532651in}}%
\pgfpathlineto{\pgfqpoint{0.877360in}{1.531640in}}%
\pgfpathlineto{\pgfqpoint{0.877506in}{1.531301in}}%
\pgfpathlineto{\pgfqpoint{0.877824in}{1.531822in}}%
\pgfpathlineto{\pgfqpoint{0.878149in}{1.532650in}}%
\pgfpathlineto{\pgfqpoint{0.878477in}{1.531936in}}%
\pgfpathlineto{\pgfqpoint{0.879106in}{1.530794in}}%
\pgfpathlineto{\pgfqpoint{0.879286in}{1.531498in}}%
\pgfpathlineto{\pgfqpoint{0.887361in}{1.575260in}}%
\pgfpathlineto{\pgfqpoint{0.887585in}{1.572096in}}%
\pgfpathlineto{\pgfqpoint{0.889398in}{1.530713in}}%
\pgfpathlineto{\pgfqpoint{0.889921in}{1.532856in}}%
\pgfpathlineto{\pgfqpoint{0.890086in}{1.533216in}}%
\pgfpathlineto{\pgfqpoint{0.890251in}{1.532676in}}%
\pgfpathlineto{\pgfqpoint{0.890695in}{1.530811in}}%
\pgfpathlineto{\pgfqpoint{0.891023in}{1.532370in}}%
\pgfpathlineto{\pgfqpoint{0.891346in}{1.533087in}}%
\pgfpathlineto{\pgfqpoint{0.891511in}{1.532252in}}%
\pgfpathlineto{\pgfqpoint{0.891817in}{1.530774in}}%
\pgfpathlineto{\pgfqpoint{0.892131in}{1.531884in}}%
\pgfpathlineto{\pgfqpoint{0.892457in}{1.533156in}}%
\pgfpathlineto{\pgfqpoint{0.892787in}{1.531689in}}%
\pgfpathlineto{\pgfqpoint{0.893067in}{1.530842in}}%
\pgfpathlineto{\pgfqpoint{0.893389in}{1.532264in}}%
\pgfpathlineto{\pgfqpoint{0.893714in}{1.533055in}}%
\pgfpathlineto{\pgfqpoint{0.893879in}{1.532310in}}%
\pgfpathlineto{\pgfqpoint{0.894183in}{1.530845in}}%
\pgfpathlineto{\pgfqpoint{0.894497in}{1.531836in}}%
\pgfpathlineto{\pgfqpoint{0.894821in}{1.533091in}}%
\pgfpathlineto{\pgfqpoint{0.895105in}{1.532014in}}%
\pgfpathlineto{\pgfqpoint{0.895592in}{1.530073in}}%
\pgfpathlineto{\pgfqpoint{0.895840in}{1.531334in}}%
\pgfpathlineto{\pgfqpoint{0.903819in}{1.581139in}}%
\pgfpathlineto{\pgfqpoint{0.904048in}{1.578350in}}%
\pgfpathlineto{\pgfqpoint{0.906390in}{1.529522in}}%
\pgfpathlineto{\pgfqpoint{0.906825in}{1.531780in}}%
\pgfpathlineto{\pgfqpoint{0.907121in}{1.533104in}}%
\pgfpathlineto{\pgfqpoint{0.907451in}{1.531919in}}%
\pgfpathlineto{\pgfqpoint{0.907741in}{1.530835in}}%
\pgfpathlineto{\pgfqpoint{0.908053in}{1.532146in}}%
\pgfpathlineto{\pgfqpoint{0.908373in}{1.533071in}}%
\pgfpathlineto{\pgfqpoint{0.908538in}{1.532400in}}%
\pgfpathlineto{\pgfqpoint{0.908839in}{1.530878in}}%
\pgfpathlineto{\pgfqpoint{0.909151in}{1.531787in}}%
\pgfpathlineto{\pgfqpoint{0.909472in}{1.533071in}}%
\pgfpathlineto{\pgfqpoint{0.909802in}{1.531752in}}%
\pgfpathlineto{\pgfqpoint{0.910081in}{1.530913in}}%
\pgfpathlineto{\pgfqpoint{0.910399in}{1.532272in}}%
\pgfpathlineto{\pgfqpoint{0.910714in}{1.532970in}}%
\pgfpathlineto{\pgfqpoint{0.910879in}{1.532240in}}%
\pgfpathlineto{\pgfqpoint{0.911184in}{1.530913in}}%
\pgfpathlineto{\pgfqpoint{0.911602in}{1.532506in}}%
\pgfpathlineto{\pgfqpoint{0.911758in}{1.532987in}}%
\pgfpathlineto{\pgfqpoint{0.912002in}{1.531156in}}%
\pgfpathlineto{\pgfqpoint{0.912055in}{1.530871in}}%
\pgfpathlineto{\pgfqpoint{0.912434in}{1.533085in}}%
\pgfpathlineto{\pgfqpoint{0.920323in}{1.588221in}}%
\pgfpathlineto{\pgfqpoint{0.920553in}{1.585352in}}%
\pgfpathlineto{\pgfqpoint{0.923298in}{1.531419in}}%
\pgfpathlineto{\pgfqpoint{0.923622in}{1.531936in}}%
\pgfpathlineto{\pgfqpoint{0.923938in}{1.532519in}}%
\pgfpathlineto{\pgfqpoint{0.924253in}{1.531800in}}%
\pgfpathlineto{\pgfqpoint{0.924541in}{1.531466in}}%
\pgfpathlineto{\pgfqpoint{0.924867in}{1.532261in}}%
\pgfpathlineto{\pgfqpoint{0.925022in}{1.532503in}}%
\pgfpathlineto{\pgfqpoint{0.925331in}{1.531953in}}%
\pgfpathlineto{\pgfqpoint{0.925618in}{1.531442in}}%
\pgfpathlineto{\pgfqpoint{0.925936in}{1.532108in}}%
\pgfpathlineto{\pgfqpoint{0.926221in}{1.532476in}}%
\pgfpathlineto{\pgfqpoint{0.926539in}{1.531712in}}%
\pgfpathlineto{\pgfqpoint{0.926696in}{1.531457in}}%
\pgfpathlineto{\pgfqpoint{0.927006in}{1.531955in}}%
\pgfpathlineto{\pgfqpoint{0.927319in}{1.532481in}}%
\pgfpathlineto{\pgfqpoint{0.927634in}{1.531800in}}%
\pgfpathlineto{\pgfqpoint{0.927923in}{1.531508in}}%
\pgfpathlineto{\pgfqpoint{0.928254in}{1.532254in}}%
\pgfpathlineto{\pgfqpoint{0.928457in}{1.531521in}}%
\pgfpathlineto{\pgfqpoint{0.928530in}{1.531145in}}%
\pgfpathlineto{\pgfqpoint{0.928897in}{1.533485in}}%
\pgfpathlineto{\pgfqpoint{0.936874in}{1.594282in}}%
\pgfpathlineto{\pgfqpoint{0.937096in}{1.590068in}}%
\pgfpathlineto{\pgfqpoint{0.940301in}{1.530573in}}%
\pgfpathlineto{\pgfqpoint{0.940459in}{1.531254in}}%
\pgfpathlineto{\pgfqpoint{0.940938in}{1.533259in}}%
\pgfpathlineto{\pgfqpoint{0.941096in}{1.532278in}}%
\pgfpathlineto{\pgfqpoint{0.941399in}{1.530582in}}%
\pgfpathlineto{\pgfqpoint{0.941707in}{1.532032in}}%
\pgfpathlineto{\pgfqpoint{0.942025in}{1.533325in}}%
\pgfpathlineto{\pgfqpoint{0.942355in}{1.531292in}}%
\pgfpathlineto{\pgfqpoint{0.942485in}{1.530667in}}%
\pgfpathlineto{\pgfqpoint{0.942793in}{1.531722in}}%
\pgfpathlineto{\pgfqpoint{0.943110in}{1.533298in}}%
\pgfpathlineto{\pgfqpoint{0.943440in}{1.531637in}}%
\pgfpathlineto{\pgfqpoint{0.943712in}{1.530716in}}%
\pgfpathlineto{\pgfqpoint{0.944026in}{1.532422in}}%
\pgfpathlineto{\pgfqpoint{0.944176in}{1.533157in}}%
\pgfpathlineto{\pgfqpoint{0.944584in}{1.531449in}}%
\pgfpathlineto{\pgfqpoint{0.944941in}{1.530355in}}%
\pgfpathlineto{\pgfqpoint{0.945260in}{1.531987in}}%
\pgfpathlineto{\pgfqpoint{0.953293in}{1.598891in}}%
\pgfpathlineto{\pgfqpoint{0.953608in}{1.594518in}}%
\pgfpathlineto{\pgfqpoint{0.957010in}{1.531325in}}%
\pgfpathlineto{\pgfqpoint{0.957175in}{1.531528in}}%
\pgfpathlineto{\pgfqpoint{0.957627in}{1.532591in}}%
\pgfpathlineto{\pgfqpoint{0.957940in}{1.531692in}}%
\pgfpathlineto{\pgfqpoint{0.958223in}{1.531433in}}%
\pgfpathlineto{\pgfqpoint{0.958388in}{1.531891in}}%
\pgfpathlineto{\pgfqpoint{0.958701in}{1.532589in}}%
\pgfpathlineto{\pgfqpoint{0.959012in}{1.531791in}}%
\pgfpathlineto{\pgfqpoint{0.959291in}{1.531399in}}%
\pgfpathlineto{\pgfqpoint{0.959610in}{1.532290in}}%
\pgfpathlineto{\pgfqpoint{0.959760in}{1.532566in}}%
\pgfpathlineto{\pgfqpoint{0.960067in}{1.531943in}}%
\pgfpathlineto{\pgfqpoint{0.960348in}{1.531379in}}%
\pgfpathlineto{\pgfqpoint{0.960661in}{1.532149in}}%
\pgfpathlineto{\pgfqpoint{0.960938in}{1.532530in}}%
\pgfpathlineto{\pgfqpoint{0.961224in}{1.531644in}}%
\pgfpathlineto{\pgfqpoint{0.961502in}{1.530191in}}%
\pgfpathlineto{\pgfqpoint{0.961798in}{1.532332in}}%
\pgfpathlineto{\pgfqpoint{0.969841in}{1.603664in}}%
\pgfpathlineto{\pgfqpoint{0.970064in}{1.599710in}}%
\pgfpathlineto{\pgfqpoint{0.973857in}{1.530173in}}%
\pgfpathlineto{\pgfqpoint{0.974005in}{1.531115in}}%
\pgfpathlineto{\pgfqpoint{0.974460in}{1.533726in}}%
\pgfpathlineto{\pgfqpoint{0.974617in}{1.532490in}}%
\pgfpathlineto{\pgfqpoint{0.974911in}{1.530139in}}%
\pgfpathlineto{\pgfqpoint{0.975216in}{1.531926in}}%
\pgfpathlineto{\pgfqpoint{0.975527in}{1.533778in}}%
\pgfpathlineto{\pgfqpoint{0.975856in}{1.531198in}}%
\pgfpathlineto{\pgfqpoint{0.975981in}{1.530297in}}%
\pgfpathlineto{\pgfqpoint{0.976424in}{1.532716in}}%
\pgfpathlineto{\pgfqpoint{0.976572in}{1.533625in}}%
\pgfpathlineto{\pgfqpoint{0.976894in}{1.531990in}}%
\pgfpathlineto{\pgfqpoint{0.977177in}{1.530261in}}%
\pgfpathlineto{\pgfqpoint{0.977562in}{1.532925in}}%
\pgfpathlineto{\pgfqpoint{0.977979in}{1.530995in}}%
\pgfpathlineto{\pgfqpoint{0.981193in}{1.560750in}}%
\pgfpathlineto{\pgfqpoint{0.986326in}{1.609076in}}%
\pgfpathlineto{\pgfqpoint{0.986549in}{1.605238in}}%
\pgfpathlineto{\pgfqpoint{0.990683in}{1.530021in}}%
\pgfpathlineto{\pgfqpoint{0.990823in}{1.530625in}}%
\pgfpathlineto{\pgfqpoint{0.991292in}{1.533875in}}%
\pgfpathlineto{\pgfqpoint{0.991622in}{1.531000in}}%
\pgfpathlineto{\pgfqpoint{0.991745in}{1.530149in}}%
\pgfpathlineto{\pgfqpoint{0.992049in}{1.531624in}}%
\pgfpathlineto{\pgfqpoint{0.992364in}{1.533821in}}%
\pgfpathlineto{\pgfqpoint{0.992694in}{1.531412in}}%
\pgfpathlineto{\pgfqpoint{0.992957in}{1.530253in}}%
\pgfpathlineto{\pgfqpoint{0.993244in}{1.532456in}}%
\pgfpathlineto{\pgfqpoint{0.993547in}{1.533646in}}%
\pgfpathlineto{\pgfqpoint{0.993701in}{1.532428in}}%
\pgfpathlineto{\pgfqpoint{0.993999in}{1.530199in}}%
\pgfpathlineto{\pgfqpoint{0.994456in}{1.532261in}}%
\pgfpathlineto{\pgfqpoint{0.994905in}{1.536169in}}%
\pgfpathlineto{\pgfqpoint{1.002833in}{1.614572in}}%
\pgfpathlineto{\pgfqpoint{1.003063in}{1.610611in}}%
\pgfpathlineto{\pgfqpoint{1.007578in}{1.530491in}}%
\pgfpathlineto{\pgfqpoint{1.007713in}{1.530848in}}%
\pgfpathlineto{\pgfqpoint{1.008173in}{1.533422in}}%
\pgfpathlineto{\pgfqpoint{1.008492in}{1.531339in}}%
\pgfpathlineto{\pgfqpoint{1.008616in}{1.530615in}}%
\pgfpathlineto{\pgfqpoint{1.008911in}{1.531626in}}%
\pgfpathlineto{\pgfqpoint{1.009203in}{1.533359in}}%
\pgfpathlineto{\pgfqpoint{1.009519in}{1.531793in}}%
\pgfpathlineto{\pgfqpoint{1.009789in}{1.530606in}}%
\pgfpathlineto{\pgfqpoint{1.010076in}{1.532283in}}%
\pgfpathlineto{\pgfqpoint{1.010367in}{1.533281in}}%
\pgfpathlineto{\pgfqpoint{1.010661in}{1.531169in}}%
\pgfpathlineto{\pgfqpoint{1.010922in}{1.529728in}}%
\pgfpathlineto{\pgfqpoint{1.011243in}{1.532254in}}%
\pgfpathlineto{\pgfqpoint{1.019295in}{1.615852in}}%
\pgfpathlineto{\pgfqpoint{1.019518in}{1.612212in}}%
\pgfpathlineto{\pgfqpoint{1.024074in}{1.529402in}}%
\pgfpathlineto{\pgfqpoint{1.024397in}{1.531625in}}%
\pgfpathlineto{\pgfqpoint{1.024693in}{1.534003in}}%
\pgfpathlineto{\pgfqpoint{1.025022in}{1.531375in}}%
\pgfpathlineto{\pgfqpoint{1.025284in}{1.530075in}}%
\pgfpathlineto{\pgfqpoint{1.025564in}{1.532472in}}%
\pgfpathlineto{\pgfqpoint{1.025864in}{1.533832in}}%
\pgfpathlineto{\pgfqpoint{1.026017in}{1.532510in}}%
\pgfpathlineto{\pgfqpoint{1.026310in}{1.530020in}}%
\pgfpathlineto{\pgfqpoint{1.026612in}{1.532054in}}%
\pgfpathlineto{\pgfqpoint{1.026915in}{1.533877in}}%
\pgfpathlineto{\pgfqpoint{1.027186in}{1.531201in}}%
\pgfpathlineto{\pgfqpoint{1.027433in}{1.528897in}}%
\pgfpathlineto{\pgfqpoint{1.027754in}{1.531793in}}%
\pgfpathlineto{\pgfqpoint{1.035800in}{1.618595in}}%
\pgfpathlineto{\pgfqpoint{1.036033in}{1.614801in}}%
\pgfpathlineto{\pgfqpoint{1.040840in}{1.529368in}}%
\pgfpathlineto{\pgfqpoint{1.041135in}{1.531344in}}%
\pgfpathlineto{\pgfqpoint{1.041429in}{1.534119in}}%
\pgfpathlineto{\pgfqpoint{1.041750in}{1.531790in}}%
\pgfpathlineto{\pgfqpoint{1.042019in}{1.529807in}}%
\pgfpathlineto{\pgfqpoint{1.042290in}{1.532134in}}%
\pgfpathlineto{\pgfqpoint{1.042583in}{1.534136in}}%
\pgfpathlineto{\pgfqpoint{1.042911in}{1.530843in}}%
\pgfpathlineto{\pgfqpoint{1.043031in}{1.529898in}}%
\pgfpathlineto{\pgfqpoint{1.043333in}{1.531634in}}%
\pgfpathlineto{\pgfqpoint{1.052270in}{1.623623in}}%
\pgfpathlineto{\pgfqpoint{1.043923in}{1.530626in}}%
\pgfpathlineto{\pgfqpoint{1.052416in}{1.621560in}}%
\pgfpathlineto{\pgfqpoint{1.057726in}{1.529509in}}%
\pgfpathlineto{\pgfqpoint{1.058034in}{1.532284in}}%
\pgfpathlineto{\pgfqpoint{1.058333in}{1.534326in}}%
\pgfpathlineto{\pgfqpoint{1.058497in}{1.532855in}}%
\pgfpathlineto{\pgfqpoint{1.058781in}{1.529641in}}%
\pgfpathlineto{\pgfqpoint{1.059083in}{1.531781in}}%
\pgfpathlineto{\pgfqpoint{1.059383in}{1.534289in}}%
\pgfpathlineto{\pgfqpoint{1.059712in}{1.531050in}}%
\pgfpathlineto{\pgfqpoint{1.059953in}{1.529844in}}%
\pgfpathlineto{\pgfqpoint{1.060302in}{1.532779in}}%
\pgfpathlineto{\pgfqpoint{1.060784in}{1.536986in}}%
\pgfpathlineto{\pgfqpoint{1.068776in}{1.629089in}}%
\pgfpathlineto{\pgfqpoint{1.069000in}{1.625586in}}%
\pgfpathlineto{\pgfqpoint{1.074485in}{1.531288in}}%
\pgfpathlineto{\pgfqpoint{1.074848in}{1.532189in}}%
\pgfpathlineto{\pgfqpoint{1.075111in}{1.532620in}}%
\pgfpathlineto{\pgfqpoint{1.075403in}{1.531599in}}%
\pgfpathlineto{\pgfqpoint{1.075549in}{1.531292in}}%
\pgfpathlineto{\pgfqpoint{1.075845in}{1.532017in}}%
\pgfpathlineto{\pgfqpoint{1.076122in}{1.532635in}}%
\pgfpathlineto{\pgfqpoint{1.076407in}{1.531745in}}%
\pgfpathlineto{\pgfqpoint{1.076879in}{1.530824in}}%
\pgfpathlineto{\pgfqpoint{1.077018in}{1.531866in}}%
\pgfpathlineto{\pgfqpoint{1.085274in}{1.629366in}}%
\pgfpathlineto{\pgfqpoint{1.085583in}{1.624335in}}%
\pgfpathlineto{\pgfqpoint{1.091270in}{1.531028in}}%
\pgfpathlineto{\pgfqpoint{1.091578in}{1.532214in}}%
\pgfpathlineto{\pgfqpoint{1.091848in}{1.532849in}}%
\pgfpathlineto{\pgfqpoint{1.092133in}{1.531507in}}%
\pgfpathlineto{\pgfqpoint{1.092260in}{1.531082in}}%
\pgfpathlineto{\pgfqpoint{1.092559in}{1.531922in}}%
\pgfpathlineto{\pgfqpoint{1.092842in}{1.532871in}}%
\pgfpathlineto{\pgfqpoint{1.093172in}{1.531146in}}%
\pgfpathlineto{\pgfqpoint{1.093378in}{1.529865in}}%
\pgfpathlineto{\pgfqpoint{1.093643in}{1.532575in}}%
\pgfpathlineto{\pgfqpoint{1.101757in}{1.630813in}}%
\pgfpathlineto{\pgfqpoint{1.102062in}{1.625952in}}%
\pgfpathlineto{\pgfqpoint{1.107799in}{1.530839in}}%
\pgfpathlineto{\pgfqpoint{1.108078in}{1.532489in}}%
\pgfpathlineto{\pgfqpoint{1.108216in}{1.533117in}}%
\pgfpathlineto{\pgfqpoint{1.108511in}{1.532039in}}%
\pgfpathlineto{\pgfqpoint{1.108783in}{1.530786in}}%
\pgfpathlineto{\pgfqpoint{1.109079in}{1.532242in}}%
\pgfpathlineto{\pgfqpoint{1.109352in}{1.533092in}}%
\pgfpathlineto{\pgfqpoint{1.109611in}{1.531359in}}%
\pgfpathlineto{\pgfqpoint{1.109841in}{1.529656in}}%
\pgfpathlineto{\pgfqpoint{1.110156in}{1.532796in}}%
\pgfpathlineto{\pgfqpoint{1.118226in}{1.632880in}}%
\pgfpathlineto{\pgfqpoint{1.118546in}{1.627946in}}%
\pgfpathlineto{\pgfqpoint{1.124603in}{1.530062in}}%
\pgfpathlineto{\pgfqpoint{1.124893in}{1.531701in}}%
\pgfpathlineto{\pgfqpoint{1.125182in}{1.533887in}}%
\pgfpathlineto{\pgfqpoint{1.125498in}{1.531353in}}%
\pgfpathlineto{\pgfqpoint{1.125754in}{1.530215in}}%
\pgfpathlineto{\pgfqpoint{1.126072in}{1.533013in}}%
\pgfpathlineto{\pgfqpoint{1.126332in}{1.532340in}}%
\pgfpathlineto{\pgfqpoint{1.127699in}{1.549023in}}%
\pgfpathlineto{\pgfqpoint{1.134711in}{1.637599in}}%
\pgfpathlineto{\pgfqpoint{1.135032in}{1.632730in}}%
\pgfpathlineto{\pgfqpoint{1.141399in}{1.529796in}}%
\pgfpathlineto{\pgfqpoint{1.141683in}{1.531377in}}%
\pgfpathlineto{\pgfqpoint{1.141966in}{1.534135in}}%
\pgfpathlineto{\pgfqpoint{1.142278in}{1.531758in}}%
\pgfpathlineto{\pgfqpoint{1.142496in}{1.529767in}}%
\pgfpathlineto{\pgfqpoint{1.142965in}{1.532734in}}%
\pgfpathlineto{\pgfqpoint{1.151195in}{1.638451in}}%
\pgfpathlineto{\pgfqpoint{1.151515in}{1.633651in}}%
\pgfpathlineto{\pgfqpoint{1.158026in}{1.530486in}}%
\pgfpathlineto{\pgfqpoint{1.158356in}{1.532495in}}%
\pgfpathlineto{\pgfqpoint{1.158493in}{1.533320in}}%
\pgfpathlineto{\pgfqpoint{1.158880in}{1.531284in}}%
\pgfpathlineto{\pgfqpoint{1.159111in}{1.530411in}}%
\pgfpathlineto{\pgfqpoint{1.159447in}{1.532026in}}%
\pgfpathlineto{\pgfqpoint{1.167660in}{1.639312in}}%
\pgfpathlineto{\pgfqpoint{1.167966in}{1.634884in}}%
\pgfpathlineto{\pgfqpoint{1.174645in}{1.530004in}}%
\pgfpathlineto{\pgfqpoint{1.174927in}{1.532251in}}%
\pgfpathlineto{\pgfqpoint{1.175212in}{1.533863in}}%
\pgfpathlineto{\pgfqpoint{1.175515in}{1.530740in}}%
\pgfpathlineto{\pgfqpoint{1.175727in}{1.529317in}}%
\pgfpathlineto{\pgfqpoint{1.176056in}{1.532349in}}%
\pgfpathlineto{\pgfqpoint{1.184164in}{1.639480in}}%
\pgfpathlineto{\pgfqpoint{1.184485in}{1.634795in}}%
\pgfpathlineto{\pgfqpoint{1.191183in}{1.529696in}}%
\pgfpathlineto{\pgfqpoint{1.191470in}{1.531474in}}%
\pgfpathlineto{\pgfqpoint{1.191754in}{1.534255in}}%
\pgfpathlineto{\pgfqpoint{1.192081in}{1.531021in}}%
\pgfpathlineto{\pgfqpoint{1.192279in}{1.529056in}}%
\pgfpathlineto{\pgfqpoint{1.192596in}{1.532623in}}%
\pgfpathlineto{\pgfqpoint{1.200629in}{1.640363in}}%
\pgfpathlineto{\pgfqpoint{1.200934in}{1.636046in}}%
\pgfpathlineto{\pgfqpoint{1.207718in}{1.530076in}}%
\pgfpathlineto{\pgfqpoint{1.208042in}{1.531663in}}%
\pgfpathlineto{\pgfqpoint{1.208310in}{1.534015in}}%
\pgfpathlineto{\pgfqpoint{1.208677in}{1.530230in}}%
\pgfpathlineto{\pgfqpoint{1.208764in}{1.529415in}}%
\pgfpathlineto{\pgfqpoint{1.209058in}{1.532740in}}%
\pgfpathlineto{\pgfqpoint{1.217114in}{1.642028in}}%
\pgfpathlineto{\pgfqpoint{1.217419in}{1.637754in}}%
\pgfpathlineto{\pgfqpoint{1.224437in}{1.530271in}}%
\pgfpathlineto{\pgfqpoint{1.224735in}{1.532004in}}%
\pgfpathlineto{\pgfqpoint{1.225253in}{1.530703in}}%
\pgfpathlineto{\pgfqpoint{1.228447in}{1.574023in}}%
\pgfpathlineto{\pgfqpoint{1.233604in}{1.644420in}}%
\pgfpathlineto{\pgfqpoint{1.233911in}{1.640160in}}%
\pgfpathlineto{\pgfqpoint{1.241149in}{1.530501in}}%
\pgfpathlineto{\pgfqpoint{1.241466in}{1.532272in}}%
\pgfpathlineto{\pgfqpoint{1.241728in}{1.532202in}}%
\pgfpathlineto{\pgfqpoint{1.242672in}{1.544632in}}%
\pgfpathlineto{\pgfqpoint{1.250086in}{1.646929in}}%
\pgfpathlineto{\pgfqpoint{1.250393in}{1.642709in}}%
\pgfpathlineto{\pgfqpoint{1.257806in}{1.529996in}}%
\pgfpathlineto{\pgfqpoint{1.258159in}{1.531963in}}%
\pgfpathlineto{\pgfqpoint{1.258561in}{1.536462in}}%
\pgfpathlineto{\pgfqpoint{1.266588in}{1.647700in}}%
\pgfpathlineto{\pgfqpoint{1.266910in}{1.643201in}}%
\pgfpathlineto{\pgfqpoint{1.274401in}{1.530329in}}%
\pgfpathlineto{\pgfqpoint{1.274756in}{1.531548in}}%
\pgfpathlineto{\pgfqpoint{1.276964in}{1.562296in}}%
\pgfpathlineto{\pgfqpoint{1.283047in}{1.647201in}}%
\pgfpathlineto{\pgfqpoint{1.283351in}{1.643097in}}%
\pgfpathlineto{\pgfqpoint{1.290872in}{1.529912in}}%
\pgfpathlineto{\pgfqpoint{1.291237in}{1.531742in}}%
\pgfpathlineto{\pgfqpoint{1.293529in}{1.563804in}}%
\pgfpathlineto{\pgfqpoint{1.299540in}{1.648072in}}%
\pgfpathlineto{\pgfqpoint{1.299844in}{1.643961in}}%
\pgfpathlineto{\pgfqpoint{1.307567in}{1.529689in}}%
\pgfpathlineto{\pgfqpoint{1.307791in}{1.531189in}}%
\pgfpathlineto{\pgfqpoint{1.316049in}{1.647378in}}%
\pgfpathlineto{\pgfqpoint{1.316528in}{1.640484in}}%
\pgfpathlineto{\pgfqpoint{1.324148in}{1.530233in}}%
\pgfpathlineto{\pgfqpoint{1.324300in}{1.531621in}}%
\pgfpathlineto{\pgfqpoint{1.332478in}{1.647829in}}%
\pgfpathlineto{\pgfqpoint{1.332924in}{1.641714in}}%
\pgfpathlineto{\pgfqpoint{1.340445in}{1.530017in}}%
\pgfpathlineto{\pgfqpoint{1.340751in}{1.531497in}}%
\pgfpathlineto{\pgfqpoint{1.348992in}{1.648825in}}%
\pgfpathlineto{\pgfqpoint{1.349628in}{1.639692in}}%
\pgfpathlineto{\pgfqpoint{1.357101in}{1.529082in}}%
\pgfpathlineto{\pgfqpoint{1.357385in}{1.532064in}}%
\pgfpathlineto{\pgfqpoint{1.365477in}{1.647787in}}%
\pgfpathlineto{\pgfqpoint{1.365783in}{1.643727in}}%
\pgfpathlineto{\pgfqpoint{1.373587in}{1.529101in}}%
\pgfpathlineto{\pgfqpoint{1.373845in}{1.531841in}}%
\pgfpathlineto{\pgfqpoint{1.381943in}{1.648313in}}%
\pgfpathlineto{\pgfqpoint{1.382247in}{1.644384in}}%
\pgfpathlineto{\pgfqpoint{1.390080in}{1.529126in}}%
\pgfpathlineto{\pgfqpoint{1.390371in}{1.532350in}}%
\pgfpathlineto{\pgfqpoint{1.398447in}{1.648750in}}%
\pgfpathlineto{\pgfqpoint{1.398753in}{1.644722in}}%
\pgfpathlineto{\pgfqpoint{1.406586in}{1.529350in}}%
\pgfpathlineto{\pgfqpoint{1.406836in}{1.532280in}}%
\pgfpathlineto{\pgfqpoint{1.414951in}{1.649361in}}%
\pgfpathlineto{\pgfqpoint{1.415283in}{1.644888in}}%
\pgfpathlineto{\pgfqpoint{1.423078in}{1.529703in}}%
\pgfpathlineto{\pgfqpoint{1.423373in}{1.533346in}}%
\pgfpathlineto{\pgfqpoint{1.431418in}{1.649968in}}%
\pgfpathlineto{\pgfqpoint{1.431723in}{1.645993in}}%
\pgfpathlineto{\pgfqpoint{1.439564in}{1.529769in}}%
\pgfpathlineto{\pgfqpoint{1.439855in}{1.533358in}}%
\pgfpathlineto{\pgfqpoint{1.447929in}{1.650364in}}%
\pgfpathlineto{\pgfqpoint{1.448245in}{1.646127in}}%
\pgfpathlineto{\pgfqpoint{1.456048in}{1.530116in}}%
\pgfpathlineto{\pgfqpoint{1.456405in}{1.534688in}}%
\pgfpathlineto{\pgfqpoint{1.464367in}{1.650761in}}%
\pgfpathlineto{\pgfqpoint{1.464670in}{1.646899in}}%
\pgfpathlineto{\pgfqpoint{1.472531in}{1.530655in}}%
\pgfpathlineto{\pgfqpoint{1.472827in}{1.534292in}}%
\pgfpathlineto{\pgfqpoint{1.480871in}{1.651558in}}%
\pgfpathlineto{\pgfqpoint{1.481178in}{1.647576in}}%
\pgfpathlineto{\pgfqpoint{1.489018in}{1.531428in}}%
\pgfpathlineto{\pgfqpoint{1.489316in}{1.535135in}}%
\pgfpathlineto{\pgfqpoint{1.497379in}{1.652472in}}%
\pgfpathlineto{\pgfqpoint{1.497706in}{1.648095in}}%
\pgfpathlineto{\pgfqpoint{1.505498in}{1.531951in}}%
\pgfpathlineto{\pgfqpoint{1.505837in}{1.536166in}}%
\pgfpathlineto{\pgfqpoint{1.513841in}{1.652968in}}%
\pgfpathlineto{\pgfqpoint{1.514147in}{1.648987in}}%
\pgfpathlineto{\pgfqpoint{1.521986in}{1.532569in}}%
\pgfpathlineto{\pgfqpoint{1.522305in}{1.536551in}}%
\pgfpathlineto{\pgfqpoint{1.530326in}{1.653593in}}%
\pgfpathlineto{\pgfqpoint{1.530633in}{1.649603in}}%
\pgfpathlineto{\pgfqpoint{1.538471in}{1.532966in}}%
\pgfpathlineto{\pgfqpoint{1.538792in}{1.536965in}}%
\pgfpathlineto{\pgfqpoint{1.546811in}{1.653935in}}%
\pgfpathlineto{\pgfqpoint{1.547118in}{1.649943in}}%
\pgfpathlineto{\pgfqpoint{1.554954in}{1.533174in}}%
\pgfpathlineto{\pgfqpoint{1.555287in}{1.537327in}}%
\pgfpathlineto{\pgfqpoint{1.563300in}{1.654089in}}%
\pgfpathlineto{\pgfqpoint{1.563609in}{1.650059in}}%
\pgfpathlineto{\pgfqpoint{1.571434in}{1.533220in}}%
\pgfpathlineto{\pgfqpoint{1.571763in}{1.537223in}}%
\pgfpathlineto{\pgfqpoint{1.579773in}{1.654003in}}%
\pgfpathlineto{\pgfqpoint{1.580076in}{1.650096in}}%
\pgfpathlineto{\pgfqpoint{1.587925in}{1.533311in}}%
\pgfpathlineto{\pgfqpoint{1.588246in}{1.537301in}}%
\pgfpathlineto{\pgfqpoint{1.596266in}{1.654072in}}%
\pgfpathlineto{\pgfqpoint{1.596573in}{1.650072in}}%
\pgfpathlineto{\pgfqpoint{1.604409in}{1.533121in}}%
\pgfpathlineto{\pgfqpoint{1.604742in}{1.537259in}}%
\pgfpathlineto{\pgfqpoint{1.612755in}{1.653833in}}%
\pgfpathlineto{\pgfqpoint{1.613064in}{1.649788in}}%
\pgfpathlineto{\pgfqpoint{1.620893in}{1.532754in}}%
\pgfpathlineto{\pgfqpoint{1.621228in}{1.536908in}}%
\pgfpathlineto{\pgfqpoint{1.629238in}{1.653424in}}%
\pgfpathlineto{\pgfqpoint{1.629546in}{1.649413in}}%
\pgfpathlineto{\pgfqpoint{1.637378in}{1.532416in}}%
\pgfpathlineto{\pgfqpoint{1.637683in}{1.536157in}}%
\pgfpathlineto{\pgfqpoint{1.645755in}{1.653166in}}%
\pgfpathlineto{\pgfqpoint{1.646059in}{1.649031in}}%
\pgfpathlineto{\pgfqpoint{1.653862in}{1.531936in}}%
\pgfpathlineto{\pgfqpoint{1.654240in}{1.536702in}}%
\pgfpathlineto{\pgfqpoint{1.662175in}{1.652399in}}%
\pgfpathlineto{\pgfqpoint{1.662472in}{1.648688in}}%
\pgfpathlineto{\pgfqpoint{1.670355in}{1.530626in}}%
\pgfpathlineto{\pgfqpoint{1.670667in}{1.534545in}}%
\pgfpathlineto{\pgfqpoint{1.678689in}{1.651363in}}%
\pgfpathlineto{\pgfqpoint{1.678996in}{1.647376in}}%
\pgfpathlineto{\pgfqpoint{1.686836in}{1.531187in}}%
\pgfpathlineto{\pgfqpoint{1.687133in}{1.534875in}}%
\pgfpathlineto{\pgfqpoint{1.695198in}{1.652134in}}%
\pgfpathlineto{\pgfqpoint{1.695525in}{1.647747in}}%
\pgfpathlineto{\pgfqpoint{1.703319in}{1.531620in}}%
\pgfpathlineto{\pgfqpoint{1.703684in}{1.536251in}}%
\pgfpathlineto{\pgfqpoint{1.711701in}{1.652619in}}%
\pgfpathlineto{\pgfqpoint{1.712015in}{1.648286in}}%
\pgfpathlineto{\pgfqpoint{1.719806in}{1.531745in}}%
\pgfpathlineto{\pgfqpoint{1.720130in}{1.535800in}}%
\pgfpathlineto{\pgfqpoint{1.728144in}{1.652709in}}%
\pgfpathlineto{\pgfqpoint{1.728450in}{1.648728in}}%
\pgfpathlineto{\pgfqpoint{1.736283in}{1.532354in}}%
\pgfpathlineto{\pgfqpoint{1.736609in}{1.536323in}}%
\pgfpathlineto{\pgfqpoint{1.744650in}{1.653380in}}%
\pgfpathlineto{\pgfqpoint{1.744971in}{1.649099in}}%
\pgfpathlineto{\pgfqpoint{1.752773in}{1.532534in}}%
\pgfpathlineto{\pgfqpoint{1.753138in}{1.537146in}}%
\pgfpathlineto{\pgfqpoint{1.761080in}{1.653249in}}%
\pgfpathlineto{\pgfqpoint{1.761372in}{1.649650in}}%
\pgfpathlineto{\pgfqpoint{1.769252in}{1.531646in}}%
\pgfpathlineto{\pgfqpoint{1.769575in}{1.535544in}}%
\pgfpathlineto{\pgfqpoint{1.777602in}{1.652606in}}%
\pgfpathlineto{\pgfqpoint{1.777907in}{1.648636in}}%
\pgfpathlineto{\pgfqpoint{1.785741in}{1.531999in}}%
\pgfpathlineto{\pgfqpoint{1.786043in}{1.535684in}}%
\pgfpathlineto{\pgfqpoint{1.794083in}{1.653006in}}%
\pgfpathlineto{\pgfqpoint{1.794390in}{1.649025in}}%
\pgfpathlineto{\pgfqpoint{1.802228in}{1.532585in}}%
\pgfpathlineto{\pgfqpoint{1.802546in}{1.536550in}}%
\pgfpathlineto{\pgfqpoint{1.810569in}{1.653596in}}%
\pgfpathlineto{\pgfqpoint{1.810876in}{1.649605in}}%
\pgfpathlineto{\pgfqpoint{1.818713in}{1.532956in}}%
\pgfpathlineto{\pgfqpoint{1.819034in}{1.536951in}}%
\pgfpathlineto{\pgfqpoint{1.827054in}{1.653909in}}%
\pgfpathlineto{\pgfqpoint{1.827361in}{1.649916in}}%
\pgfpathlineto{\pgfqpoint{1.835197in}{1.533145in}}%
\pgfpathlineto{\pgfqpoint{1.835530in}{1.537292in}}%
\pgfpathlineto{\pgfqpoint{1.843543in}{1.654048in}}%
\pgfpathlineto{\pgfqpoint{1.843852in}{1.650010in}}%
\pgfpathlineto{\pgfqpoint{1.851678in}{1.533222in}}%
\pgfpathlineto{\pgfqpoint{1.851977in}{1.536809in}}%
\pgfpathlineto{\pgfqpoint{1.860044in}{1.654058in}}%
\pgfpathlineto{\pgfqpoint{1.860372in}{1.649664in}}%
\pgfpathlineto{\pgfqpoint{1.868169in}{1.533039in}}%
\pgfpathlineto{\pgfqpoint{1.868482in}{1.536900in}}%
\pgfpathlineto{\pgfqpoint{1.876508in}{1.653750in}}%
\pgfpathlineto{\pgfqpoint{1.876815in}{1.649750in}}%
\pgfpathlineto{\pgfqpoint{1.884653in}{1.532881in}}%
\pgfpathlineto{\pgfqpoint{1.884972in}{1.536846in}}%
\pgfpathlineto{\pgfqpoint{1.892993in}{1.653591in}}%
\pgfpathlineto{\pgfqpoint{1.893300in}{1.649589in}}%
\pgfpathlineto{\pgfqpoint{1.901138in}{1.532741in}}%
\pgfpathlineto{\pgfqpoint{1.901454in}{1.536664in}}%
\pgfpathlineto{\pgfqpoint{1.909478in}{1.653428in}}%
\pgfpathlineto{\pgfqpoint{1.909785in}{1.649426in}}%
\pgfpathlineto{\pgfqpoint{1.917623in}{1.532613in}}%
\pgfpathlineto{\pgfqpoint{1.917938in}{1.536536in}}%
\pgfpathlineto{\pgfqpoint{1.925963in}{1.653289in}}%
\pgfpathlineto{\pgfqpoint{1.926270in}{1.649287in}}%
\pgfpathlineto{\pgfqpoint{1.934105in}{1.532508in}}%
\pgfpathlineto{\pgfqpoint{1.934425in}{1.536448in}}%
\pgfpathlineto{\pgfqpoint{1.942451in}{1.653186in}}%
\pgfpathlineto{\pgfqpoint{1.942762in}{1.649115in}}%
\pgfpathlineto{\pgfqpoint{1.950590in}{1.532388in}}%
\pgfpathlineto{\pgfqpoint{1.950905in}{1.536266in}}%
\pgfpathlineto{\pgfqpoint{1.958955in}{1.653141in}}%
\pgfpathlineto{\pgfqpoint{1.959269in}{1.648941in}}%
\pgfpathlineto{\pgfqpoint{1.967074in}{1.531979in}}%
\pgfpathlineto{\pgfqpoint{1.967449in}{1.536697in}}%
\pgfpathlineto{\pgfqpoint{1.975417in}{1.652652in}}%
\pgfpathlineto{\pgfqpoint{1.975723in}{1.648660in}}%
\pgfpathlineto{\pgfqpoint{1.983558in}{1.532080in}}%
\pgfpathlineto{\pgfqpoint{1.983890in}{1.536179in}}%
\pgfpathlineto{\pgfqpoint{1.991895in}{1.652797in}}%
\pgfpathlineto{\pgfqpoint{1.992199in}{1.648870in}}%
\pgfpathlineto{\pgfqpoint{2.000041in}{1.532351in}}%
\pgfpathlineto{\pgfqpoint{2.000399in}{1.536772in}}%
\pgfpathlineto{\pgfqpoint{2.008412in}{1.653151in}}%
\pgfpathlineto{\pgfqpoint{2.008726in}{1.648941in}}%
\pgfpathlineto{\pgfqpoint{2.016525in}{1.531962in}}%
\pgfpathlineto{\pgfqpoint{2.016932in}{1.537109in}}%
\pgfpathlineto{\pgfqpoint{2.024876in}{1.652744in}}%
\pgfpathlineto{\pgfqpoint{2.025185in}{1.648705in}}%
\pgfpathlineto{\pgfqpoint{2.033016in}{1.532019in}}%
\pgfpathlineto{\pgfqpoint{2.033320in}{1.535769in}}%
\pgfpathlineto{\pgfqpoint{2.041357in}{1.652830in}}%
\pgfpathlineto{\pgfqpoint{2.041664in}{1.648833in}}%
\pgfpathlineto{\pgfqpoint{2.049511in}{1.532279in}}%
\pgfpathlineto{\pgfqpoint{2.049812in}{1.536116in}}%
\pgfpathlineto{\pgfqpoint{2.057842in}{1.653095in}}%
\pgfpathlineto{\pgfqpoint{2.058149in}{1.649099in}}%
\pgfpathlineto{\pgfqpoint{2.065986in}{1.532480in}}%
\pgfpathlineto{\pgfqpoint{2.066361in}{1.537255in}}%
\pgfpathlineto{\pgfqpoint{2.074357in}{1.653389in}}%
\pgfpathlineto{\pgfqpoint{2.074657in}{1.649346in}}%
\pgfpathlineto{\pgfqpoint{2.082479in}{1.532273in}}%
\pgfpathlineto{\pgfqpoint{2.082830in}{1.536816in}}%
\pgfpathlineto{\pgfqpoint{2.090836in}{1.653139in}}%
\pgfpathlineto{\pgfqpoint{2.091150in}{1.648934in}}%
\pgfpathlineto{\pgfqpoint{2.098960in}{1.531974in}}%
\pgfpathlineto{\pgfqpoint{2.099312in}{1.536467in}}%
\pgfpathlineto{\pgfqpoint{2.107288in}{1.652739in}}%
\pgfpathlineto{\pgfqpoint{2.107593in}{1.648819in}}%
\pgfpathlineto{\pgfqpoint{2.115438in}{1.532374in}}%
\pgfpathlineto{\pgfqpoint{2.115746in}{1.536162in}}%
\pgfpathlineto{\pgfqpoint{2.123808in}{1.653311in}}%
\pgfpathlineto{\pgfqpoint{2.124123in}{1.649084in}}%
\pgfpathlineto{\pgfqpoint{2.131930in}{1.532208in}}%
\pgfpathlineto{\pgfqpoint{2.132322in}{1.537302in}}%
\pgfpathlineto{\pgfqpoint{2.140265in}{1.653014in}}%
\pgfpathlineto{\pgfqpoint{2.140572in}{1.649026in}}%
\pgfpathlineto{\pgfqpoint{2.148409in}{1.532455in}}%
\pgfpathlineto{\pgfqpoint{2.148762in}{1.536883in}}%
\pgfpathlineto{\pgfqpoint{2.156739in}{1.653237in}}%
\pgfpathlineto{\pgfqpoint{2.157038in}{1.649404in}}%
\pgfpathlineto{\pgfqpoint{2.164895in}{1.532845in}}%
\pgfpathlineto{\pgfqpoint{2.165214in}{1.536826in}}%
\pgfpathlineto{\pgfqpoint{2.173236in}{1.653673in}}%
\pgfpathlineto{\pgfqpoint{2.173543in}{1.649675in}}%
\pgfpathlineto{\pgfqpoint{2.181380in}{1.532877in}}%
\pgfpathlineto{\pgfqpoint{2.181699in}{1.536842in}}%
\pgfpathlineto{\pgfqpoint{2.189720in}{1.653659in}}%
\pgfpathlineto{\pgfqpoint{2.190027in}{1.649660in}}%
\pgfpathlineto{\pgfqpoint{2.197865in}{1.532839in}}%
\pgfpathlineto{\pgfqpoint{2.198182in}{1.536786in}}%
\pgfpathlineto{\pgfqpoint{2.206205in}{1.653587in}}%
\pgfpathlineto{\pgfqpoint{2.206512in}{1.649586in}}%
\pgfpathlineto{\pgfqpoint{2.214350in}{1.532763in}}%
\pgfpathlineto{\pgfqpoint{2.214666in}{1.536688in}}%
\pgfpathlineto{\pgfqpoint{2.222690in}{1.653484in}}%
\pgfpathlineto{\pgfqpoint{2.222997in}{1.649483in}}%
\pgfpathlineto{\pgfqpoint{2.230835in}{1.532675in}}%
\pgfpathlineto{\pgfqpoint{2.231150in}{1.536598in}}%
\pgfpathlineto{\pgfqpoint{2.239175in}{1.653378in}}%
\pgfpathlineto{\pgfqpoint{2.239482in}{1.649377in}}%
\pgfpathlineto{\pgfqpoint{2.247319in}{1.532589in}}%
\pgfpathlineto{\pgfqpoint{2.247635in}{1.536515in}}%
\pgfpathlineto{\pgfqpoint{2.255660in}{1.653285in}}%
\pgfpathlineto{\pgfqpoint{2.255967in}{1.649283in}}%
\pgfpathlineto{\pgfqpoint{2.263802in}{1.532518in}}%
\pgfpathlineto{\pgfqpoint{2.264135in}{1.536649in}}%
\pgfpathlineto{\pgfqpoint{2.272147in}{1.653223in}}%
\pgfpathlineto{\pgfqpoint{2.272454in}{1.649222in}}%
\pgfpathlineto{\pgfqpoint{2.280282in}{1.532297in}}%
\pgfpathlineto{\pgfqpoint{2.280603in}{1.536157in}}%
\pgfpathlineto{\pgfqpoint{2.288648in}{1.652992in}}%
\pgfpathlineto{\pgfqpoint{2.288979in}{1.648560in}}%
\pgfpathlineto{\pgfqpoint{2.296780in}{1.532205in}}%
\pgfpathlineto{\pgfqpoint{2.297112in}{1.536436in}}%
\pgfpathlineto{\pgfqpoint{2.305152in}{1.652989in}}%
\pgfpathlineto{\pgfqpoint{2.305462in}{1.648757in}}%
\pgfpathlineto{\pgfqpoint{2.313262in}{1.531845in}}%
\pgfpathlineto{\pgfqpoint{2.313601in}{1.536107in}}%
\pgfpathlineto{\pgfqpoint{2.321598in}{1.652532in}}%
\pgfpathlineto{\pgfqpoint{2.321905in}{1.648541in}}%
\pgfpathlineto{\pgfqpoint{2.329743in}{1.532015in}}%
\pgfpathlineto{\pgfqpoint{2.330030in}{1.535524in}}%
\pgfpathlineto{\pgfqpoint{2.338106in}{1.652841in}}%
\pgfpathlineto{\pgfqpoint{2.338433in}{1.648454in}}%
\pgfpathlineto{\pgfqpoint{2.346234in}{1.532062in}}%
\pgfpathlineto{\pgfqpoint{2.346525in}{1.535700in}}%
\pgfpathlineto{\pgfqpoint{2.354593in}{1.652928in}}%
\pgfpathlineto{\pgfqpoint{2.354907in}{1.648728in}}%
\pgfpathlineto{\pgfqpoint{2.362713in}{1.531884in}}%
\pgfpathlineto{\pgfqpoint{2.363071in}{1.536387in}}%
\pgfpathlineto{\pgfqpoint{2.371093in}{1.652739in}}%
\pgfpathlineto{\pgfqpoint{2.371405in}{1.648456in}}%
\pgfpathlineto{\pgfqpoint{2.379202in}{1.531722in}}%
\pgfpathlineto{\pgfqpoint{2.379528in}{1.535844in}}%
\pgfpathlineto{\pgfqpoint{2.387538in}{1.652550in}}%
\pgfpathlineto{\pgfqpoint{2.387844in}{1.648563in}}%
\pgfpathlineto{\pgfqpoint{2.395683in}{1.532115in}}%
\pgfpathlineto{\pgfqpoint{2.395998in}{1.536028in}}%
\pgfpathlineto{\pgfqpoint{2.404044in}{1.653053in}}%
\pgfpathlineto{\pgfqpoint{2.404366in}{1.648763in}}%
\pgfpathlineto{\pgfqpoint{2.412170in}{1.532101in}}%
\pgfpathlineto{\pgfqpoint{2.412558in}{1.537088in}}%
\pgfpathlineto{\pgfqpoint{2.420514in}{1.652985in}}%
\pgfpathlineto{\pgfqpoint{2.420827in}{1.648887in}}%
\pgfpathlineto{\pgfqpoint{2.428653in}{1.532130in}}%
\pgfpathlineto{\pgfqpoint{2.428958in}{1.535890in}}%
\pgfpathlineto{\pgfqpoint{2.437021in}{1.653115in}}%
\pgfpathlineto{\pgfqpoint{2.437339in}{1.648865in}}%
\pgfpathlineto{\pgfqpoint{2.445135in}{1.532098in}}%
\pgfpathlineto{\pgfqpoint{2.445473in}{1.536305in}}%
\pgfpathlineto{\pgfqpoint{2.453477in}{1.652978in}}%
\pgfpathlineto{\pgfqpoint{2.453784in}{1.648992in}}%
\pgfpathlineto{\pgfqpoint{2.461622in}{1.532478in}}%
\pgfpathlineto{\pgfqpoint{2.461989in}{1.537132in}}%
\pgfpathlineto{\pgfqpoint{2.470005in}{1.653393in}}%
\pgfpathlineto{\pgfqpoint{2.470321in}{1.649008in}}%
\pgfpathlineto{\pgfqpoint{2.478101in}{1.532428in}}%
\pgfpathlineto{\pgfqpoint{2.478424in}{1.536323in}}%
\pgfpathlineto{\pgfqpoint{2.486451in}{1.653270in}}%
\pgfpathlineto{\pgfqpoint{2.486757in}{1.649285in}}%
\pgfpathlineto{\pgfqpoint{2.494589in}{1.532288in}}%
\pgfpathlineto{\pgfqpoint{2.494938in}{1.536673in}}%
\pgfpathlineto{\pgfqpoint{2.502907in}{1.653009in}}%
\pgfpathlineto{\pgfqpoint{2.503206in}{1.649239in}}%
\pgfpathlineto{\pgfqpoint{2.511082in}{1.531698in}}%
\pgfpathlineto{\pgfqpoint{2.511403in}{1.535769in}}%
\pgfpathlineto{\pgfqpoint{2.519381in}{1.652345in}}%
\pgfpathlineto{\pgfqpoint{2.519669in}{1.648811in}}%
\pgfpathlineto{\pgfqpoint{2.527562in}{1.530708in}}%
\pgfpathlineto{\pgfqpoint{2.527871in}{1.534549in}}%
\pgfpathlineto{\pgfqpoint{2.535902in}{1.651690in}}%
\pgfpathlineto{\pgfqpoint{2.536209in}{1.647702in}}%
\pgfpathlineto{\pgfqpoint{2.544049in}{1.531564in}}%
\pgfpathlineto{\pgfqpoint{2.544347in}{1.535282in}}%
\pgfpathlineto{\pgfqpoint{2.552408in}{1.652677in}}%
\pgfpathlineto{\pgfqpoint{2.552735in}{1.648306in}}%
\pgfpathlineto{\pgfqpoint{2.560535in}{1.532174in}}%
\pgfpathlineto{\pgfqpoint{2.560838in}{1.535947in}}%
\pgfpathlineto{\pgfqpoint{2.568895in}{1.653271in}}%
\pgfpathlineto{\pgfqpoint{2.569223in}{1.648882in}}%
\pgfpathlineto{\pgfqpoint{2.577009in}{1.532620in}}%
\pgfpathlineto{\pgfqpoint{2.577327in}{1.536432in}}%
\pgfpathlineto{\pgfqpoint{2.585360in}{1.653629in}}%
\pgfpathlineto{\pgfqpoint{2.585672in}{1.649563in}}%
\pgfpathlineto{\pgfqpoint{2.593500in}{1.532968in}}%
\pgfpathlineto{\pgfqpoint{2.593832in}{1.537114in}}%
\pgfpathlineto{\pgfqpoint{2.601846in}{1.653971in}}%
\pgfpathlineto{\pgfqpoint{2.602155in}{1.649936in}}%
\pgfpathlineto{\pgfqpoint{2.609985in}{1.533299in}}%
\pgfpathlineto{\pgfqpoint{2.610332in}{1.537678in}}%
\pgfpathlineto{\pgfqpoint{2.618329in}{1.654239in}}%
\pgfpathlineto{\pgfqpoint{2.618635in}{1.650258in}}%
\pgfpathlineto{\pgfqpoint{2.626471in}{1.533127in}}%
\pgfpathlineto{\pgfqpoint{2.626860in}{1.538096in}}%
\pgfpathlineto{\pgfqpoint{2.634836in}{1.654011in}}%
\pgfpathlineto{\pgfqpoint{2.635151in}{1.649803in}}%
\pgfpathlineto{\pgfqpoint{2.642957in}{1.532847in}}%
\pgfpathlineto{\pgfqpoint{2.643279in}{1.536849in}}%
\pgfpathlineto{\pgfqpoint{2.651299in}{1.653608in}}%
\pgfpathlineto{\pgfqpoint{2.651606in}{1.649601in}}%
\pgfpathlineto{\pgfqpoint{2.659439in}{1.532624in}}%
\pgfpathlineto{\pgfqpoint{2.659771in}{1.536755in}}%
\pgfpathlineto{\pgfqpoint{2.667785in}{1.653376in}}%
\pgfpathlineto{\pgfqpoint{2.668094in}{1.649336in}}%
\pgfpathlineto{\pgfqpoint{2.675924in}{1.532456in}}%
\pgfpathlineto{\pgfqpoint{2.676277in}{1.536871in}}%
\pgfpathlineto{\pgfqpoint{2.684255in}{1.653136in}}%
\pgfpathlineto{\pgfqpoint{2.684553in}{1.649296in}}%
\pgfpathlineto{\pgfqpoint{2.692410in}{1.532690in}}%
\pgfpathlineto{\pgfqpoint{2.692726in}{1.536619in}}%
\pgfpathlineto{\pgfqpoint{2.700751in}{1.653433in}}%
\pgfpathlineto{\pgfqpoint{2.701058in}{1.649433in}}%
\pgfpathlineto{\pgfqpoint{2.708895in}{1.532651in}}%
\pgfpathlineto{\pgfqpoint{2.709211in}{1.536576in}}%
\pgfpathlineto{\pgfqpoint{2.717236in}{1.653374in}}%
\pgfpathlineto{\pgfqpoint{2.717543in}{1.649373in}}%
\pgfpathlineto{\pgfqpoint{2.725380in}{1.532600in}}%
\pgfpathlineto{\pgfqpoint{2.725696in}{1.536525in}}%
\pgfpathlineto{\pgfqpoint{2.733720in}{1.653312in}}%
\pgfpathlineto{\pgfqpoint{2.734027in}{1.649311in}}%
\pgfpathlineto{\pgfqpoint{2.741864in}{1.532549in}}%
\pgfpathlineto{\pgfqpoint{2.742194in}{1.536661in}}%
\pgfpathlineto{\pgfqpoint{2.750205in}{1.653254in}}%
\pgfpathlineto{\pgfqpoint{2.750512in}{1.649253in}}%
\pgfpathlineto{\pgfqpoint{2.758361in}{1.532488in}}%
\pgfpathlineto{\pgfqpoint{2.758659in}{1.536304in}}%
\pgfpathlineto{\pgfqpoint{2.766708in}{1.653201in}}%
\pgfpathlineto{\pgfqpoint{2.767041in}{1.648756in}}%
\pgfpathlineto{\pgfqpoint{2.774829in}{1.532391in}}%
\pgfpathlineto{\pgfqpoint{2.775201in}{1.537002in}}%
\pgfpathlineto{\pgfqpoint{2.783197in}{1.653107in}}%
\pgfpathlineto{\pgfqpoint{2.783513in}{1.648885in}}%
\pgfpathlineto{\pgfqpoint{2.791316in}{1.531993in}}%
\pgfpathlineto{\pgfqpoint{2.791704in}{1.536896in}}%
\pgfpathlineto{\pgfqpoint{2.799659in}{1.652690in}}%
\pgfpathlineto{\pgfqpoint{2.799966in}{1.648698in}}%
\pgfpathlineto{\pgfqpoint{2.807805in}{1.532124in}}%
\pgfpathlineto{\pgfqpoint{2.808109in}{1.535889in}}%
\pgfpathlineto{\pgfqpoint{2.816172in}{1.652973in}}%
\pgfpathlineto{\pgfqpoint{2.816489in}{1.648732in}}%
\pgfpathlineto{\pgfqpoint{2.824296in}{1.531831in}}%
\pgfpathlineto{\pgfqpoint{2.824603in}{1.535696in}}%
\pgfpathlineto{\pgfqpoint{2.832599in}{1.652380in}}%
\pgfpathlineto{\pgfqpoint{2.832896in}{1.648678in}}%
\pgfpathlineto{\pgfqpoint{2.840775in}{1.530703in}}%
\pgfpathlineto{\pgfqpoint{2.841083in}{1.534526in}}%
\pgfpathlineto{\pgfqpoint{2.849114in}{1.651572in}}%
\pgfpathlineto{\pgfqpoint{2.849421in}{1.647580in}}%
\pgfpathlineto{\pgfqpoint{2.857261in}{1.531392in}}%
\pgfpathlineto{\pgfqpoint{2.857559in}{1.535103in}}%
\pgfpathlineto{\pgfqpoint{2.865620in}{1.652409in}}%
\pgfpathlineto{\pgfqpoint{2.865947in}{1.648035in}}%
\pgfpathlineto{\pgfqpoint{2.873743in}{1.531899in}}%
\pgfpathlineto{\pgfqpoint{2.874126in}{1.536782in}}%
\pgfpathlineto{\pgfqpoint{2.882106in}{1.652952in}}%
\pgfpathlineto{\pgfqpoint{2.882423in}{1.648728in}}%
\pgfpathlineto{\pgfqpoint{2.890227in}{1.532131in}}%
\pgfpathlineto{\pgfqpoint{2.890539in}{1.535982in}}%
\pgfpathlineto{\pgfqpoint{2.898572in}{1.653150in}}%
\pgfpathlineto{\pgfqpoint{2.898881in}{1.649135in}}%
\pgfpathlineto{\pgfqpoint{2.906713in}{1.532481in}}%
\pgfpathlineto{\pgfqpoint{2.907079in}{1.537126in}}%
\pgfpathlineto{\pgfqpoint{2.915012in}{1.653201in}}%
\pgfpathlineto{\pgfqpoint{2.915368in}{1.648710in}}%
\pgfpathlineto{\pgfqpoint{2.923198in}{1.531892in}}%
\pgfpathlineto{\pgfqpoint{2.923577in}{1.536721in}}%
\pgfpathlineto{\pgfqpoint{2.931561in}{1.652938in}}%
\pgfpathlineto{\pgfqpoint{2.931877in}{1.648724in}}%
\pgfpathlineto{\pgfqpoint{2.939680in}{1.532109in}}%
\pgfpathlineto{\pgfqpoint{2.939993in}{1.535955in}}%
\pgfpathlineto{\pgfqpoint{2.948023in}{1.653116in}}%
\pgfpathlineto{\pgfqpoint{2.948330in}{1.649126in}}%
\pgfpathlineto{\pgfqpoint{2.956168in}{1.532623in}}%
\pgfpathlineto{\pgfqpoint{2.956484in}{1.536561in}}%
\pgfpathlineto{\pgfqpoint{2.964508in}{1.653622in}}%
\pgfpathlineto{\pgfqpoint{2.964815in}{1.649630in}}%
\pgfpathlineto{\pgfqpoint{2.972652in}{1.532966in}}%
\pgfpathlineto{\pgfqpoint{2.972974in}{1.536964in}}%
\pgfpathlineto{\pgfqpoint{2.980993in}{1.653907in}}%
\pgfpathlineto{\pgfqpoint{2.981300in}{1.649913in}}%
\pgfpathlineto{\pgfqpoint{2.989136in}{1.533134in}}%
\pgfpathlineto{\pgfqpoint{2.989469in}{1.537280in}}%
\pgfpathlineto{\pgfqpoint{2.997478in}{1.654009in}}%
\pgfpathlineto{\pgfqpoint{2.997785in}{1.650013in}}%
\pgfpathlineto{\pgfqpoint{3.005621in}{1.533160in}}%
\pgfpathlineto{\pgfqpoint{3.005954in}{1.537306in}}%
\pgfpathlineto{\pgfqpoint{3.013967in}{1.653983in}}%
\pgfpathlineto{\pgfqpoint{3.014276in}{1.649947in}}%
\pgfpathlineto{\pgfqpoint{3.022107in}{1.532850in}}%
\pgfpathlineto{\pgfqpoint{3.022416in}{1.536646in}}%
\pgfpathlineto{\pgfqpoint{3.030448in}{1.653585in}}%
\pgfpathlineto{\pgfqpoint{3.030755in}{1.649585in}}%
\pgfpathlineto{\pgfqpoint{3.038592in}{1.532770in}}%
\pgfpathlineto{\pgfqpoint{3.038908in}{1.536696in}}%
\pgfpathlineto{\pgfqpoint{3.046933in}{1.653502in}}%
\pgfpathlineto{\pgfqpoint{3.047240in}{1.649501in}}%
\pgfpathlineto{\pgfqpoint{3.055077in}{1.532694in}}%
\pgfpathlineto{\pgfqpoint{3.055393in}{1.536618in}}%
\pgfpathlineto{\pgfqpoint{3.063417in}{1.653406in}}%
\pgfpathlineto{\pgfqpoint{3.063724in}{1.649405in}}%
\pgfpathlineto{\pgfqpoint{3.071562in}{1.532615in}}%
\pgfpathlineto{\pgfqpoint{3.071878in}{1.536539in}}%
\pgfpathlineto{\pgfqpoint{3.079902in}{1.653316in}}%
\pgfpathlineto{\pgfqpoint{3.080209in}{1.649314in}}%
\pgfpathlineto{\pgfqpoint{3.088046in}{1.532543in}}%
\pgfpathlineto{\pgfqpoint{3.088375in}{1.536650in}}%
\pgfpathlineto{\pgfqpoint{3.096387in}{1.653237in}}%
\pgfpathlineto{\pgfqpoint{3.096694in}{1.649235in}}%
\pgfpathlineto{\pgfqpoint{3.104532in}{1.532486in}}%
\pgfpathlineto{\pgfqpoint{3.104847in}{1.536384in}}%
\pgfpathlineto{\pgfqpoint{3.112895in}{1.653235in}}%
\pgfpathlineto{\pgfqpoint{3.113209in}{1.649035in}}%
\pgfpathlineto{\pgfqpoint{3.121022in}{1.532069in}}%
\pgfpathlineto{\pgfqpoint{3.121389in}{1.536812in}}%
\pgfpathlineto{\pgfqpoint{3.129356in}{1.652757in}}%
\pgfpathlineto{\pgfqpoint{3.129662in}{1.648765in}}%
\pgfpathlineto{\pgfqpoint{3.137505in}{1.532163in}}%
\pgfpathlineto{\pgfqpoint{3.137800in}{1.535814in}}%
\pgfpathlineto{\pgfqpoint{3.145847in}{1.652906in}}%
\pgfpathlineto{\pgfqpoint{3.146154in}{1.648895in}}%
\pgfpathlineto{\pgfqpoint{3.153983in}{1.531972in}}%
\pgfpathlineto{\pgfqpoint{3.154329in}{1.536284in}}%
\pgfpathlineto{\pgfqpoint{3.162326in}{1.652722in}}%
\pgfpathlineto{\pgfqpoint{3.162632in}{1.648733in}}%
\pgfpathlineto{\pgfqpoint{3.170132in}{1.536774in}}%
\pgfpathlineto{\pgfqpoint{3.170132in}{1.536774in}}%
\pgfusepath{stroke}%
\end{pgfscope}%
\begin{pgfscope}%
\pgfsetrectcap%
\pgfsetmiterjoin%
\pgfsetlinewidth{0.803000pt}%
\definecolor{currentstroke}{rgb}{0.000000,0.000000,0.000000}%
\pgfsetstrokecolor{currentstroke}%
\pgfsetdash{}{0pt}%
\pgfpathmoveto{\pgfqpoint{0.573768in}{1.532029in}}%
\pgfpathlineto{\pgfqpoint{0.573768in}{2.438279in}}%
\pgfusepath{stroke}%
\end{pgfscope}%
\begin{pgfscope}%
\pgfsetrectcap%
\pgfsetmiterjoin%
\pgfsetlinewidth{0.803000pt}%
\definecolor{currentstroke}{rgb}{0.000000,0.000000,0.000000}%
\pgfsetstrokecolor{currentstroke}%
\pgfsetdash{}{0pt}%
\pgfpathmoveto{\pgfqpoint{3.293768in}{1.532029in}}%
\pgfpathlineto{\pgfqpoint{3.293768in}{2.438279in}}%
\pgfusepath{stroke}%
\end{pgfscope}%
\begin{pgfscope}%
\pgfsetrectcap%
\pgfsetmiterjoin%
\pgfsetlinewidth{0.803000pt}%
\definecolor{currentstroke}{rgb}{0.000000,0.000000,0.000000}%
\pgfsetstrokecolor{currentstroke}%
\pgfsetdash{}{0pt}%
\pgfpathmoveto{\pgfqpoint{0.573768in}{1.532029in}}%
\pgfpathlineto{\pgfqpoint{3.293768in}{1.532029in}}%
\pgfusepath{stroke}%
\end{pgfscope}%
\begin{pgfscope}%
\pgfsetrectcap%
\pgfsetmiterjoin%
\pgfsetlinewidth{0.803000pt}%
\definecolor{currentstroke}{rgb}{0.000000,0.000000,0.000000}%
\pgfsetstrokecolor{currentstroke}%
\pgfsetdash{}{0pt}%
\pgfpathmoveto{\pgfqpoint{0.573768in}{2.438279in}}%
\pgfpathlineto{\pgfqpoint{3.293768in}{2.438279in}}%
\pgfusepath{stroke}%
\end{pgfscope}%
\begin{pgfscope}%
\definecolor{textcolor}{rgb}{0.000000,0.000000,0.000000}%
\pgfsetstrokecolor{textcolor}%
\pgfsetfillcolor{textcolor}%
\pgftext[x=0.628168in,y=2.374842in,left,top]{\color{textcolor}{\rmfamily\fontsize{8.000000}{9.600000}\selectfont\catcode`\^=\active\def^{\ifmmode\sp\else\^{}\fi}\catcode`\%=\active\def%{\%}(b)}}%
\end{pgfscope}%
\begin{pgfscope}%
\pgfsetbuttcap%
\pgfsetmiterjoin%
\definecolor{currentfill}{rgb}{1.000000,1.000000,1.000000}%
\pgfsetfillcolor{currentfill}%
\pgfsetlinewidth{0.000000pt}%
\definecolor{currentstroke}{rgb}{0.000000,0.000000,0.000000}%
\pgfsetstrokecolor{currentstroke}%
\pgfsetstrokeopacity{0.000000}%
\pgfsetdash{}{0pt}%
\pgfpathmoveto{\pgfqpoint{0.573768in}{0.462654in}}%
\pgfpathlineto{\pgfqpoint{3.293768in}{0.462654in}}%
\pgfpathlineto{\pgfqpoint{3.293768in}{1.368904in}}%
\pgfpathlineto{\pgfqpoint{0.573768in}{1.368904in}}%
\pgfpathlineto{\pgfqpoint{0.573768in}{0.462654in}}%
\pgfpathclose%
\pgfusepath{fill}%
\end{pgfscope}%
\begin{pgfscope}%
\pgfpathrectangle{\pgfqpoint{0.573768in}{0.462654in}}{\pgfqpoint{2.720000in}{0.906250in}}%
\pgfusepath{clip}%
\pgfsetbuttcap%
\pgfsetroundjoin%
\pgfsetlinewidth{0.501875pt}%
\definecolor{currentstroke}{rgb}{0.690196,0.690196,0.690196}%
\pgfsetstrokecolor{currentstroke}%
\pgfsetstrokeopacity{0.250000}%
\pgfsetdash{{1.850000pt}{0.800000pt}}{0.000000pt}%
\pgfpathmoveto{\pgfqpoint{0.697405in}{0.462654in}}%
\pgfpathlineto{\pgfqpoint{0.697405in}{1.368904in}}%
\pgfusepath{stroke}%
\end{pgfscope}%
\begin{pgfscope}%
\pgfsetbuttcap%
\pgfsetroundjoin%
\definecolor{currentfill}{rgb}{0.000000,0.000000,0.000000}%
\pgfsetfillcolor{currentfill}%
\pgfsetlinewidth{0.803000pt}%
\definecolor{currentstroke}{rgb}{0.000000,0.000000,0.000000}%
\pgfsetstrokecolor{currentstroke}%
\pgfsetdash{}{0pt}%
\pgfsys@defobject{currentmarker}{\pgfqpoint{0.000000in}{-0.048611in}}{\pgfqpoint{0.000000in}{0.000000in}}{%
\pgfpathmoveto{\pgfqpoint{0.000000in}{0.000000in}}%
\pgfpathlineto{\pgfqpoint{0.000000in}{-0.048611in}}%
\pgfusepath{stroke,fill}%
}%
\begin{pgfscope}%
\pgfsys@transformshift{0.697405in}{0.462654in}%
\pgfsys@useobject{currentmarker}{}%
\end{pgfscope}%
\end{pgfscope}%
\begin{pgfscope}%
\definecolor{textcolor}{rgb}{0.000000,0.000000,0.000000}%
\pgfsetstrokecolor{textcolor}%
\pgfsetfillcolor{textcolor}%
\pgftext[x=0.697405in,y=0.365432in,,top]{\color{textcolor}{\rmfamily\fontsize{8.000000}{9.600000}\selectfont\catcode`\^=\active\def^{\ifmmode\sp\else\^{}\fi}\catcode`\%=\active\def%{\%}$\mathdefault{0}$}}%
\end{pgfscope}%
\begin{pgfscope}%
\pgfpathrectangle{\pgfqpoint{0.573768in}{0.462654in}}{\pgfqpoint{2.720000in}{0.906250in}}%
\pgfusepath{clip}%
\pgfsetbuttcap%
\pgfsetroundjoin%
\pgfsetlinewidth{0.501875pt}%
\definecolor{currentstroke}{rgb}{0.690196,0.690196,0.690196}%
\pgfsetstrokecolor{currentstroke}%
\pgfsetstrokeopacity{0.250000}%
\pgfsetdash{{1.850000pt}{0.800000pt}}{0.000000pt}%
\pgfpathmoveto{\pgfqpoint{1.109526in}{0.462654in}}%
\pgfpathlineto{\pgfqpoint{1.109526in}{1.368904in}}%
\pgfusepath{stroke}%
\end{pgfscope}%
\begin{pgfscope}%
\pgfsetbuttcap%
\pgfsetroundjoin%
\definecolor{currentfill}{rgb}{0.000000,0.000000,0.000000}%
\pgfsetfillcolor{currentfill}%
\pgfsetlinewidth{0.803000pt}%
\definecolor{currentstroke}{rgb}{0.000000,0.000000,0.000000}%
\pgfsetstrokecolor{currentstroke}%
\pgfsetdash{}{0pt}%
\pgfsys@defobject{currentmarker}{\pgfqpoint{0.000000in}{-0.048611in}}{\pgfqpoint{0.000000in}{0.000000in}}{%
\pgfpathmoveto{\pgfqpoint{0.000000in}{0.000000in}}%
\pgfpathlineto{\pgfqpoint{0.000000in}{-0.048611in}}%
\pgfusepath{stroke,fill}%
}%
\begin{pgfscope}%
\pgfsys@transformshift{1.109526in}{0.462654in}%
\pgfsys@useobject{currentmarker}{}%
\end{pgfscope}%
\end{pgfscope}%
\begin{pgfscope}%
\definecolor{textcolor}{rgb}{0.000000,0.000000,0.000000}%
\pgfsetstrokecolor{textcolor}%
\pgfsetfillcolor{textcolor}%
\pgftext[x=1.109526in,y=0.365432in,,top]{\color{textcolor}{\rmfamily\fontsize{8.000000}{9.600000}\selectfont\catcode`\^=\active\def^{\ifmmode\sp\else\^{}\fi}\catcode`\%=\active\def%{\%}$\mathdefault{1}$}}%
\end{pgfscope}%
\begin{pgfscope}%
\pgfpathrectangle{\pgfqpoint{0.573768in}{0.462654in}}{\pgfqpoint{2.720000in}{0.906250in}}%
\pgfusepath{clip}%
\pgfsetbuttcap%
\pgfsetroundjoin%
\pgfsetlinewidth{0.501875pt}%
\definecolor{currentstroke}{rgb}{0.690196,0.690196,0.690196}%
\pgfsetstrokecolor{currentstroke}%
\pgfsetstrokeopacity{0.250000}%
\pgfsetdash{{1.850000pt}{0.800000pt}}{0.000000pt}%
\pgfpathmoveto{\pgfqpoint{1.521647in}{0.462654in}}%
\pgfpathlineto{\pgfqpoint{1.521647in}{1.368904in}}%
\pgfusepath{stroke}%
\end{pgfscope}%
\begin{pgfscope}%
\pgfsetbuttcap%
\pgfsetroundjoin%
\definecolor{currentfill}{rgb}{0.000000,0.000000,0.000000}%
\pgfsetfillcolor{currentfill}%
\pgfsetlinewidth{0.803000pt}%
\definecolor{currentstroke}{rgb}{0.000000,0.000000,0.000000}%
\pgfsetstrokecolor{currentstroke}%
\pgfsetdash{}{0pt}%
\pgfsys@defobject{currentmarker}{\pgfqpoint{0.000000in}{-0.048611in}}{\pgfqpoint{0.000000in}{0.000000in}}{%
\pgfpathmoveto{\pgfqpoint{0.000000in}{0.000000in}}%
\pgfpathlineto{\pgfqpoint{0.000000in}{-0.048611in}}%
\pgfusepath{stroke,fill}%
}%
\begin{pgfscope}%
\pgfsys@transformshift{1.521647in}{0.462654in}%
\pgfsys@useobject{currentmarker}{}%
\end{pgfscope}%
\end{pgfscope}%
\begin{pgfscope}%
\definecolor{textcolor}{rgb}{0.000000,0.000000,0.000000}%
\pgfsetstrokecolor{textcolor}%
\pgfsetfillcolor{textcolor}%
\pgftext[x=1.521647in,y=0.365432in,,top]{\color{textcolor}{\rmfamily\fontsize{8.000000}{9.600000}\selectfont\catcode`\^=\active\def^{\ifmmode\sp\else\^{}\fi}\catcode`\%=\active\def%{\%}$\mathdefault{2}$}}%
\end{pgfscope}%
\begin{pgfscope}%
\pgfpathrectangle{\pgfqpoint{0.573768in}{0.462654in}}{\pgfqpoint{2.720000in}{0.906250in}}%
\pgfusepath{clip}%
\pgfsetbuttcap%
\pgfsetroundjoin%
\pgfsetlinewidth{0.501875pt}%
\definecolor{currentstroke}{rgb}{0.690196,0.690196,0.690196}%
\pgfsetstrokecolor{currentstroke}%
\pgfsetstrokeopacity{0.250000}%
\pgfsetdash{{1.850000pt}{0.800000pt}}{0.000000pt}%
\pgfpathmoveto{\pgfqpoint{1.933768in}{0.462654in}}%
\pgfpathlineto{\pgfqpoint{1.933768in}{1.368904in}}%
\pgfusepath{stroke}%
\end{pgfscope}%
\begin{pgfscope}%
\pgfsetbuttcap%
\pgfsetroundjoin%
\definecolor{currentfill}{rgb}{0.000000,0.000000,0.000000}%
\pgfsetfillcolor{currentfill}%
\pgfsetlinewidth{0.803000pt}%
\definecolor{currentstroke}{rgb}{0.000000,0.000000,0.000000}%
\pgfsetstrokecolor{currentstroke}%
\pgfsetdash{}{0pt}%
\pgfsys@defobject{currentmarker}{\pgfqpoint{0.000000in}{-0.048611in}}{\pgfqpoint{0.000000in}{0.000000in}}{%
\pgfpathmoveto{\pgfqpoint{0.000000in}{0.000000in}}%
\pgfpathlineto{\pgfqpoint{0.000000in}{-0.048611in}}%
\pgfusepath{stroke,fill}%
}%
\begin{pgfscope}%
\pgfsys@transformshift{1.933768in}{0.462654in}%
\pgfsys@useobject{currentmarker}{}%
\end{pgfscope}%
\end{pgfscope}%
\begin{pgfscope}%
\definecolor{textcolor}{rgb}{0.000000,0.000000,0.000000}%
\pgfsetstrokecolor{textcolor}%
\pgfsetfillcolor{textcolor}%
\pgftext[x=1.933768in,y=0.365432in,,top]{\color{textcolor}{\rmfamily\fontsize{8.000000}{9.600000}\selectfont\catcode`\^=\active\def^{\ifmmode\sp\else\^{}\fi}\catcode`\%=\active\def%{\%}$\mathdefault{3}$}}%
\end{pgfscope}%
\begin{pgfscope}%
\pgfpathrectangle{\pgfqpoint{0.573768in}{0.462654in}}{\pgfqpoint{2.720000in}{0.906250in}}%
\pgfusepath{clip}%
\pgfsetbuttcap%
\pgfsetroundjoin%
\pgfsetlinewidth{0.501875pt}%
\definecolor{currentstroke}{rgb}{0.690196,0.690196,0.690196}%
\pgfsetstrokecolor{currentstroke}%
\pgfsetstrokeopacity{0.250000}%
\pgfsetdash{{1.850000pt}{0.800000pt}}{0.000000pt}%
\pgfpathmoveto{\pgfqpoint{2.345890in}{0.462654in}}%
\pgfpathlineto{\pgfqpoint{2.345890in}{1.368904in}}%
\pgfusepath{stroke}%
\end{pgfscope}%
\begin{pgfscope}%
\pgfsetbuttcap%
\pgfsetroundjoin%
\definecolor{currentfill}{rgb}{0.000000,0.000000,0.000000}%
\pgfsetfillcolor{currentfill}%
\pgfsetlinewidth{0.803000pt}%
\definecolor{currentstroke}{rgb}{0.000000,0.000000,0.000000}%
\pgfsetstrokecolor{currentstroke}%
\pgfsetdash{}{0pt}%
\pgfsys@defobject{currentmarker}{\pgfqpoint{0.000000in}{-0.048611in}}{\pgfqpoint{0.000000in}{0.000000in}}{%
\pgfpathmoveto{\pgfqpoint{0.000000in}{0.000000in}}%
\pgfpathlineto{\pgfqpoint{0.000000in}{-0.048611in}}%
\pgfusepath{stroke,fill}%
}%
\begin{pgfscope}%
\pgfsys@transformshift{2.345890in}{0.462654in}%
\pgfsys@useobject{currentmarker}{}%
\end{pgfscope}%
\end{pgfscope}%
\begin{pgfscope}%
\definecolor{textcolor}{rgb}{0.000000,0.000000,0.000000}%
\pgfsetstrokecolor{textcolor}%
\pgfsetfillcolor{textcolor}%
\pgftext[x=2.345890in,y=0.365432in,,top]{\color{textcolor}{\rmfamily\fontsize{8.000000}{9.600000}\selectfont\catcode`\^=\active\def^{\ifmmode\sp\else\^{}\fi}\catcode`\%=\active\def%{\%}$\mathdefault{4}$}}%
\end{pgfscope}%
\begin{pgfscope}%
\pgfpathrectangle{\pgfqpoint{0.573768in}{0.462654in}}{\pgfqpoint{2.720000in}{0.906250in}}%
\pgfusepath{clip}%
\pgfsetbuttcap%
\pgfsetroundjoin%
\pgfsetlinewidth{0.501875pt}%
\definecolor{currentstroke}{rgb}{0.690196,0.690196,0.690196}%
\pgfsetstrokecolor{currentstroke}%
\pgfsetstrokeopacity{0.250000}%
\pgfsetdash{{1.850000pt}{0.800000pt}}{0.000000pt}%
\pgfpathmoveto{\pgfqpoint{2.758011in}{0.462654in}}%
\pgfpathlineto{\pgfqpoint{2.758011in}{1.368904in}}%
\pgfusepath{stroke}%
\end{pgfscope}%
\begin{pgfscope}%
\pgfsetbuttcap%
\pgfsetroundjoin%
\definecolor{currentfill}{rgb}{0.000000,0.000000,0.000000}%
\pgfsetfillcolor{currentfill}%
\pgfsetlinewidth{0.803000pt}%
\definecolor{currentstroke}{rgb}{0.000000,0.000000,0.000000}%
\pgfsetstrokecolor{currentstroke}%
\pgfsetdash{}{0pt}%
\pgfsys@defobject{currentmarker}{\pgfqpoint{0.000000in}{-0.048611in}}{\pgfqpoint{0.000000in}{0.000000in}}{%
\pgfpathmoveto{\pgfqpoint{0.000000in}{0.000000in}}%
\pgfpathlineto{\pgfqpoint{0.000000in}{-0.048611in}}%
\pgfusepath{stroke,fill}%
}%
\begin{pgfscope}%
\pgfsys@transformshift{2.758011in}{0.462654in}%
\pgfsys@useobject{currentmarker}{}%
\end{pgfscope}%
\end{pgfscope}%
\begin{pgfscope}%
\definecolor{textcolor}{rgb}{0.000000,0.000000,0.000000}%
\pgfsetstrokecolor{textcolor}%
\pgfsetfillcolor{textcolor}%
\pgftext[x=2.758011in,y=0.365432in,,top]{\color{textcolor}{\rmfamily\fontsize{8.000000}{9.600000}\selectfont\catcode`\^=\active\def^{\ifmmode\sp\else\^{}\fi}\catcode`\%=\active\def%{\%}$\mathdefault{5}$}}%
\end{pgfscope}%
\begin{pgfscope}%
\pgfpathrectangle{\pgfqpoint{0.573768in}{0.462654in}}{\pgfqpoint{2.720000in}{0.906250in}}%
\pgfusepath{clip}%
\pgfsetbuttcap%
\pgfsetroundjoin%
\pgfsetlinewidth{0.501875pt}%
\definecolor{currentstroke}{rgb}{0.690196,0.690196,0.690196}%
\pgfsetstrokecolor{currentstroke}%
\pgfsetstrokeopacity{0.250000}%
\pgfsetdash{{1.850000pt}{0.800000pt}}{0.000000pt}%
\pgfpathmoveto{\pgfqpoint{3.170132in}{0.462654in}}%
\pgfpathlineto{\pgfqpoint{3.170132in}{1.368904in}}%
\pgfusepath{stroke}%
\end{pgfscope}%
\begin{pgfscope}%
\pgfsetbuttcap%
\pgfsetroundjoin%
\definecolor{currentfill}{rgb}{0.000000,0.000000,0.000000}%
\pgfsetfillcolor{currentfill}%
\pgfsetlinewidth{0.803000pt}%
\definecolor{currentstroke}{rgb}{0.000000,0.000000,0.000000}%
\pgfsetstrokecolor{currentstroke}%
\pgfsetdash{}{0pt}%
\pgfsys@defobject{currentmarker}{\pgfqpoint{0.000000in}{-0.048611in}}{\pgfqpoint{0.000000in}{0.000000in}}{%
\pgfpathmoveto{\pgfqpoint{0.000000in}{0.000000in}}%
\pgfpathlineto{\pgfqpoint{0.000000in}{-0.048611in}}%
\pgfusepath{stroke,fill}%
}%
\begin{pgfscope}%
\pgfsys@transformshift{3.170132in}{0.462654in}%
\pgfsys@useobject{currentmarker}{}%
\end{pgfscope}%
\end{pgfscope}%
\begin{pgfscope}%
\definecolor{textcolor}{rgb}{0.000000,0.000000,0.000000}%
\pgfsetstrokecolor{textcolor}%
\pgfsetfillcolor{textcolor}%
\pgftext[x=3.170132in,y=0.365432in,,top]{\color{textcolor}{\rmfamily\fontsize{8.000000}{9.600000}\selectfont\catcode`\^=\active\def^{\ifmmode\sp\else\^{}\fi}\catcode`\%=\active\def%{\%}$\mathdefault{6}$}}%
\end{pgfscope}%
\begin{pgfscope}%
\definecolor{textcolor}{rgb}{0.000000,0.000000,0.000000}%
\pgfsetstrokecolor{textcolor}%
\pgfsetfillcolor{textcolor}%
\pgftext[x=1.933768in,y=0.211111in,,top]{\color{textcolor}{\rmfamily\fontsize{8.000000}{9.600000}\selectfont\catcode`\^=\active\def^{\ifmmode\sp\else\^{}\fi}\catcode`\%=\active\def%{\%}Time (ms)}}%
\end{pgfscope}%
\begin{pgfscope}%
\pgfpathrectangle{\pgfqpoint{0.573768in}{0.462654in}}{\pgfqpoint{2.720000in}{0.906250in}}%
\pgfusepath{clip}%
\pgfsetbuttcap%
\pgfsetroundjoin%
\pgfsetlinewidth{0.501875pt}%
\definecolor{currentstroke}{rgb}{0.690196,0.690196,0.690196}%
\pgfsetstrokecolor{currentstroke}%
\pgfsetstrokeopacity{0.250000}%
\pgfsetdash{{1.850000pt}{0.800000pt}}{0.000000pt}%
\pgfpathmoveto{\pgfqpoint{0.573768in}{0.525661in}}%
\pgfpathlineto{\pgfqpoint{3.293768in}{0.525661in}}%
\pgfusepath{stroke}%
\end{pgfscope}%
\begin{pgfscope}%
\pgfsetbuttcap%
\pgfsetroundjoin%
\definecolor{currentfill}{rgb}{0.000000,0.000000,0.000000}%
\pgfsetfillcolor{currentfill}%
\pgfsetlinewidth{0.803000pt}%
\definecolor{currentstroke}{rgb}{0.000000,0.000000,0.000000}%
\pgfsetstrokecolor{currentstroke}%
\pgfsetdash{}{0pt}%
\pgfsys@defobject{currentmarker}{\pgfqpoint{-0.048611in}{0.000000in}}{\pgfqpoint{-0.000000in}{0.000000in}}{%
\pgfpathmoveto{\pgfqpoint{-0.000000in}{0.000000in}}%
\pgfpathlineto{\pgfqpoint{-0.048611in}{0.000000in}}%
\pgfusepath{stroke,fill}%
}%
\begin{pgfscope}%
\pgfsys@transformshift{0.573768in}{0.525661in}%
\pgfsys@useobject{currentmarker}{}%
\end{pgfscope}%
\end{pgfscope}%
\begin{pgfscope}%
\definecolor{textcolor}{rgb}{0.000000,0.000000,0.000000}%
\pgfsetstrokecolor{textcolor}%
\pgfsetfillcolor{textcolor}%
\pgftext[x=0.266667in, y=0.487081in, left, base]{\color{textcolor}{\rmfamily\fontsize{8.000000}{9.600000}\selectfont\catcode`\^=\active\def^{\ifmmode\sp\else\^{}\fi}\catcode`\%=\active\def%{\%}$\mathdefault{\ensuremath{-}20}$}}%
\end{pgfscope}%
\begin{pgfscope}%
\pgfpathrectangle{\pgfqpoint{0.573768in}{0.462654in}}{\pgfqpoint{2.720000in}{0.906250in}}%
\pgfusepath{clip}%
\pgfsetbuttcap%
\pgfsetroundjoin%
\pgfsetlinewidth{0.501875pt}%
\definecolor{currentstroke}{rgb}{0.690196,0.690196,0.690196}%
\pgfsetstrokecolor{currentstroke}%
\pgfsetstrokeopacity{0.250000}%
\pgfsetdash{{1.850000pt}{0.800000pt}}{0.000000pt}%
\pgfpathmoveto{\pgfqpoint{0.573768in}{0.921903in}}%
\pgfpathlineto{\pgfqpoint{3.293768in}{0.921903in}}%
\pgfusepath{stroke}%
\end{pgfscope}%
\begin{pgfscope}%
\pgfsetbuttcap%
\pgfsetroundjoin%
\definecolor{currentfill}{rgb}{0.000000,0.000000,0.000000}%
\pgfsetfillcolor{currentfill}%
\pgfsetlinewidth{0.803000pt}%
\definecolor{currentstroke}{rgb}{0.000000,0.000000,0.000000}%
\pgfsetstrokecolor{currentstroke}%
\pgfsetdash{}{0pt}%
\pgfsys@defobject{currentmarker}{\pgfqpoint{-0.048611in}{0.000000in}}{\pgfqpoint{-0.000000in}{0.000000in}}{%
\pgfpathmoveto{\pgfqpoint{-0.000000in}{0.000000in}}%
\pgfpathlineto{\pgfqpoint{-0.048611in}{0.000000in}}%
\pgfusepath{stroke,fill}%
}%
\begin{pgfscope}%
\pgfsys@transformshift{0.573768in}{0.921903in}%
\pgfsys@useobject{currentmarker}{}%
\end{pgfscope}%
\end{pgfscope}%
\begin{pgfscope}%
\definecolor{textcolor}{rgb}{0.000000,0.000000,0.000000}%
\pgfsetstrokecolor{textcolor}%
\pgfsetfillcolor{textcolor}%
\pgftext[x=0.266667in, y=0.883323in, left, base]{\color{textcolor}{\rmfamily\fontsize{8.000000}{9.600000}\selectfont\catcode`\^=\active\def^{\ifmmode\sp\else\^{}\fi}\catcode`\%=\active\def%{\%}$\mathdefault{\ensuremath{-}10}$}}%
\end{pgfscope}%
\begin{pgfscope}%
\pgfpathrectangle{\pgfqpoint{0.573768in}{0.462654in}}{\pgfqpoint{2.720000in}{0.906250in}}%
\pgfusepath{clip}%
\pgfsetbuttcap%
\pgfsetroundjoin%
\pgfsetlinewidth{0.501875pt}%
\definecolor{currentstroke}{rgb}{0.690196,0.690196,0.690196}%
\pgfsetstrokecolor{currentstroke}%
\pgfsetstrokeopacity{0.250000}%
\pgfsetdash{{1.850000pt}{0.800000pt}}{0.000000pt}%
\pgfpathmoveto{\pgfqpoint{0.573768in}{1.318145in}}%
\pgfpathlineto{\pgfqpoint{3.293768in}{1.318145in}}%
\pgfusepath{stroke}%
\end{pgfscope}%
\begin{pgfscope}%
\pgfsetbuttcap%
\pgfsetroundjoin%
\definecolor{currentfill}{rgb}{0.000000,0.000000,0.000000}%
\pgfsetfillcolor{currentfill}%
\pgfsetlinewidth{0.803000pt}%
\definecolor{currentstroke}{rgb}{0.000000,0.000000,0.000000}%
\pgfsetstrokecolor{currentstroke}%
\pgfsetdash{}{0pt}%
\pgfsys@defobject{currentmarker}{\pgfqpoint{-0.048611in}{0.000000in}}{\pgfqpoint{-0.000000in}{0.000000in}}{%
\pgfpathmoveto{\pgfqpoint{-0.000000in}{0.000000in}}%
\pgfpathlineto{\pgfqpoint{-0.048611in}{0.000000in}}%
\pgfusepath{stroke,fill}%
}%
\begin{pgfscope}%
\pgfsys@transformshift{0.573768in}{1.318145in}%
\pgfsys@useobject{currentmarker}{}%
\end{pgfscope}%
\end{pgfscope}%
\begin{pgfscope}%
\definecolor{textcolor}{rgb}{0.000000,0.000000,0.000000}%
\pgfsetstrokecolor{textcolor}%
\pgfsetfillcolor{textcolor}%
\pgftext[x=0.417518in, y=1.279565in, left, base]{\color{textcolor}{\rmfamily\fontsize{8.000000}{9.600000}\selectfont\catcode`\^=\active\def^{\ifmmode\sp\else\^{}\fi}\catcode`\%=\active\def%{\%}$\mathdefault{0}$}}%
\end{pgfscope}%
\begin{pgfscope}%
\definecolor{textcolor}{rgb}{0.000000,0.000000,0.000000}%
\pgfsetstrokecolor{textcolor}%
\pgfsetfillcolor{textcolor}%
\pgftext[x=0.211111in,y=0.915779in,,bottom,rotate=90.000000]{\color{textcolor}{\rmfamily\fontsize{8.000000}{9.600000}\selectfont\catcode`\^=\active\def^{\ifmmode\sp\else\^{}\fi}\catcode`\%=\active\def%{\%}$v_D$ (V)}}%
\end{pgfscope}%
\begin{pgfscope}%
\pgfpathrectangle{\pgfqpoint{0.573768in}{0.462654in}}{\pgfqpoint{2.720000in}{0.906250in}}%
\pgfusepath{clip}%
\pgfsetrectcap%
\pgfsetroundjoin%
\pgfsetlinewidth{1.505625pt}%
\definecolor{currentstroke}{rgb}{0.000000,0.619608,0.450980}%
\pgfsetstrokecolor{currentstroke}%
\pgfsetdash{}{0pt}%
\pgfpathmoveto{\pgfqpoint{0.697405in}{1.318145in}}%
\pgfpathlineto{\pgfqpoint{0.697448in}{1.321092in}}%
\pgfpathlineto{\pgfqpoint{0.697711in}{1.320212in}}%
\pgfpathlineto{\pgfqpoint{0.698119in}{0.525713in}}%
\pgfpathlineto{\pgfqpoint{0.698614in}{0.525733in}}%
\pgfpathlineto{\pgfqpoint{0.705819in}{0.526317in}}%
\pgfpathlineto{\pgfqpoint{0.705911in}{0.526851in}}%
\pgfpathlineto{\pgfqpoint{0.706084in}{0.899442in}}%
\pgfpathlineto{\pgfqpoint{0.706852in}{1.326683in}}%
\pgfpathlineto{\pgfqpoint{0.708665in}{1.326604in}}%
\pgfpathlineto{\pgfqpoint{0.710808in}{1.326654in}}%
\pgfpathlineto{\pgfqpoint{0.712951in}{1.326574in}}%
\pgfpathlineto{\pgfqpoint{0.714156in}{1.326599in}}%
\pgfpathlineto{\pgfqpoint{0.714206in}{1.326462in}}%
\pgfpathlineto{\pgfqpoint{0.715494in}{0.526278in}}%
\pgfpathlineto{\pgfqpoint{0.722217in}{0.526629in}}%
\pgfpathlineto{\pgfqpoint{0.722396in}{0.527957in}}%
\pgfpathlineto{\pgfqpoint{0.722571in}{0.921296in}}%
\pgfpathlineto{\pgfqpoint{0.722634in}{1.327301in}}%
\pgfpathlineto{\pgfqpoint{0.723312in}{1.327298in}}%
\pgfpathlineto{\pgfqpoint{0.728093in}{1.327219in}}%
\pgfpathlineto{\pgfqpoint{0.730691in}{1.327124in}}%
\pgfpathlineto{\pgfqpoint{0.731631in}{0.526764in}}%
\pgfpathlineto{\pgfqpoint{0.735257in}{0.526758in}}%
\pgfpathlineto{\pgfqpoint{0.738789in}{0.527376in}}%
\pgfpathlineto{\pgfqpoint{0.738881in}{0.528749in}}%
\pgfpathlineto{\pgfqpoint{0.739056in}{0.933360in}}%
\pgfpathlineto{\pgfqpoint{0.739119in}{1.327592in}}%
\pgfpathlineto{\pgfqpoint{0.739800in}{1.327584in}}%
\pgfpathlineto{\pgfqpoint{0.746229in}{1.327454in}}%
\pgfpathlineto{\pgfqpoint{0.747176in}{1.327387in}}%
\pgfpathlineto{\pgfqpoint{0.747189in}{1.324083in}}%
\pgfpathlineto{\pgfqpoint{0.747993in}{0.527129in}}%
\pgfpathlineto{\pgfqpoint{0.749642in}{0.526893in}}%
\pgfpathlineto{\pgfqpoint{0.755228in}{0.527297in}}%
\pgfpathlineto{\pgfqpoint{0.755366in}{0.529120in}}%
\pgfpathlineto{\pgfqpoint{0.755541in}{0.938455in}}%
\pgfpathlineto{\pgfqpoint{0.755604in}{1.327711in}}%
\pgfpathlineto{\pgfqpoint{0.756285in}{1.327690in}}%
\pgfpathlineto{\pgfqpoint{0.763661in}{1.327440in}}%
\pgfpathlineto{\pgfqpoint{0.763673in}{1.324889in}}%
\pgfpathlineto{\pgfqpoint{0.764481in}{0.527185in}}%
\pgfpathlineto{\pgfqpoint{0.766129in}{0.526921in}}%
\pgfpathlineto{\pgfqpoint{0.771712in}{0.527256in}}%
\pgfpathlineto{\pgfqpoint{0.771851in}{0.529029in}}%
\pgfpathlineto{\pgfqpoint{0.772026in}{0.937230in}}%
\pgfpathlineto{\pgfqpoint{0.772089in}{1.327683in}}%
\pgfpathlineto{\pgfqpoint{0.772783in}{1.327655in}}%
\pgfpathlineto{\pgfqpoint{0.779914in}{1.327390in}}%
\pgfpathlineto{\pgfqpoint{0.780146in}{1.327314in}}%
\pgfpathlineto{\pgfqpoint{0.780195in}{0.838868in}}%
\pgfpathlineto{\pgfqpoint{0.780951in}{0.527023in}}%
\pgfpathlineto{\pgfqpoint{0.782765in}{0.526754in}}%
\pgfpathlineto{\pgfqpoint{0.788243in}{0.527246in}}%
\pgfpathlineto{\pgfqpoint{0.788336in}{0.528508in}}%
\pgfpathlineto{\pgfqpoint{0.788511in}{0.929960in}}%
\pgfpathlineto{\pgfqpoint{0.788583in}{1.327507in}}%
\pgfpathlineto{\pgfqpoint{0.789308in}{1.327478in}}%
\pgfpathlineto{\pgfqpoint{0.794089in}{1.327199in}}%
\pgfpathlineto{\pgfqpoint{0.796631in}{1.326966in}}%
\pgfpathlineto{\pgfqpoint{0.797498in}{0.526608in}}%
\pgfpathlineto{\pgfqpoint{0.800795in}{0.526455in}}%
\pgfpathlineto{\pgfqpoint{0.804728in}{0.526773in}}%
\pgfpathlineto{\pgfqpoint{0.804821in}{0.527659in}}%
\pgfpathlineto{\pgfqpoint{0.804995in}{0.916517in}}%
\pgfpathlineto{\pgfqpoint{0.805066in}{1.327154in}}%
\pgfpathlineto{\pgfqpoint{0.805789in}{1.327108in}}%
\pgfpathlineto{\pgfqpoint{0.809251in}{1.326766in}}%
\pgfpathlineto{\pgfqpoint{0.811394in}{1.326590in}}%
\pgfpathlineto{\pgfqpoint{0.813115in}{1.326127in}}%
\pgfpathlineto{\pgfqpoint{0.813943in}{0.526128in}}%
\pgfpathlineto{\pgfqpoint{0.820867in}{0.526065in}}%
\pgfpathlineto{\pgfqpoint{0.821262in}{0.526348in}}%
\pgfpathlineto{\pgfqpoint{0.821353in}{0.527721in}}%
\pgfpathlineto{\pgfqpoint{0.821572in}{1.326426in}}%
\pgfpathlineto{\pgfqpoint{0.822283in}{1.326328in}}%
\pgfpathlineto{\pgfqpoint{0.825085in}{1.325628in}}%
\pgfpathlineto{\pgfqpoint{0.826239in}{1.325343in}}%
\pgfpathlineto{\pgfqpoint{0.828052in}{1.323214in}}%
\pgfpathlineto{\pgfqpoint{0.828217in}{1.323708in}}%
\pgfpathlineto{\pgfqpoint{0.828382in}{1.280542in}}%
\pgfpathlineto{\pgfqpoint{0.828547in}{1.322536in}}%
\pgfpathlineto{\pgfqpoint{0.828946in}{0.503848in}}%
\pgfpathlineto{\pgfqpoint{0.829413in}{0.527521in}}%
\pgfpathlineto{\pgfqpoint{0.829577in}{0.639618in}}%
\pgfpathlineto{\pgfqpoint{0.829763in}{0.525641in}}%
\pgfpathlineto{\pgfqpoint{0.830177in}{0.525678in}}%
\pgfpathlineto{\pgfqpoint{0.836936in}{0.525695in}}%
\pgfpathlineto{\pgfqpoint{0.837795in}{0.525735in}}%
\pgfpathlineto{\pgfqpoint{0.837880in}{0.542788in}}%
\pgfpathlineto{\pgfqpoint{0.838032in}{1.324308in}}%
\pgfpathlineto{\pgfqpoint{0.838727in}{1.322898in}}%
\pgfpathlineto{\pgfqpoint{0.839281in}{0.505031in}}%
\pgfpathlineto{\pgfqpoint{0.839607in}{0.547033in}}%
\pgfpathlineto{\pgfqpoint{0.840102in}{0.706973in}}%
\pgfpathlineto{\pgfqpoint{0.840417in}{0.589539in}}%
\pgfpathlineto{\pgfqpoint{0.840747in}{0.527945in}}%
\pgfpathlineto{\pgfqpoint{0.841076in}{0.613224in}}%
\pgfpathlineto{\pgfqpoint{0.841406in}{0.706448in}}%
\pgfpathlineto{\pgfqpoint{0.841722in}{0.588759in}}%
\pgfpathlineto{\pgfqpoint{0.842051in}{0.531535in}}%
\pgfpathlineto{\pgfqpoint{0.842216in}{0.557916in}}%
\pgfpathlineto{\pgfqpoint{0.842711in}{0.704976in}}%
\pgfpathlineto{\pgfqpoint{0.843027in}{0.585412in}}%
\pgfpathlineto{\pgfqpoint{0.843357in}{0.535777in}}%
\pgfpathlineto{\pgfqpoint{0.843522in}{0.566561in}}%
\pgfpathlineto{\pgfqpoint{0.844017in}{0.701813in}}%
\pgfpathlineto{\pgfqpoint{0.844181in}{0.644283in}}%
\pgfpathlineto{\pgfqpoint{0.844665in}{0.541227in}}%
\pgfpathlineto{\pgfqpoint{0.844830in}{0.578007in}}%
\pgfpathlineto{\pgfqpoint{0.845159in}{0.702616in}}%
\pgfpathlineto{\pgfqpoint{0.845489in}{0.633753in}}%
\pgfpathlineto{\pgfqpoint{0.845773in}{0.545588in}}%
\pgfpathlineto{\pgfqpoint{0.845991in}{0.698252in}}%
\pgfpathlineto{\pgfqpoint{0.846069in}{0.749600in}}%
\pgfpathlineto{\pgfqpoint{0.846413in}{0.525676in}}%
\pgfpathlineto{\pgfqpoint{0.854275in}{0.525772in}}%
\pgfpathlineto{\pgfqpoint{0.854362in}{0.539460in}}%
\pgfpathlineto{\pgfqpoint{0.854515in}{1.324649in}}%
\pgfpathlineto{\pgfqpoint{0.855213in}{1.323931in}}%
\pgfpathlineto{\pgfqpoint{0.856241in}{0.508909in}}%
\pgfpathlineto{\pgfqpoint{0.856570in}{0.583733in}}%
\pgfpathlineto{\pgfqpoint{0.856900in}{0.800178in}}%
\pgfpathlineto{\pgfqpoint{0.857226in}{0.626667in}}%
\pgfpathlineto{\pgfqpoint{0.857543in}{0.511372in}}%
\pgfpathlineto{\pgfqpoint{0.857873in}{0.609562in}}%
\pgfpathlineto{\pgfqpoint{0.858202in}{0.800316in}}%
\pgfpathlineto{\pgfqpoint{0.858517in}{0.611399in}}%
\pgfpathlineto{\pgfqpoint{0.858846in}{0.514703in}}%
\pgfpathlineto{\pgfqpoint{0.859176in}{0.642411in}}%
\pgfpathlineto{\pgfqpoint{0.859506in}{0.792523in}}%
\pgfpathlineto{\pgfqpoint{0.859796in}{0.600707in}}%
\pgfpathlineto{\pgfqpoint{0.860126in}{0.519290in}}%
\pgfpathlineto{\pgfqpoint{0.860291in}{0.559297in}}%
\pgfpathlineto{\pgfqpoint{0.860785in}{0.782609in}}%
\pgfpathlineto{\pgfqpoint{0.861079in}{0.585670in}}%
\pgfpathlineto{\pgfqpoint{0.861409in}{0.526157in}}%
\pgfpathlineto{\pgfqpoint{0.861574in}{0.579523in}}%
\pgfpathlineto{\pgfqpoint{0.861903in}{0.787706in}}%
\pgfpathlineto{\pgfqpoint{0.862384in}{0.643583in}}%
\pgfpathlineto{\pgfqpoint{0.862565in}{0.699463in}}%
\pgfpathlineto{\pgfqpoint{0.862709in}{0.526995in}}%
\pgfpathlineto{\pgfqpoint{0.862744in}{0.525605in}}%
\pgfpathlineto{\pgfqpoint{0.863485in}{0.525681in}}%
\pgfpathlineto{\pgfqpoint{0.869254in}{0.525714in}}%
\pgfpathlineto{\pgfqpoint{0.870767in}{0.525821in}}%
\pgfpathlineto{\pgfqpoint{0.870852in}{0.551824in}}%
\pgfpathlineto{\pgfqpoint{0.871003in}{1.324829in}}%
\pgfpathlineto{\pgfqpoint{0.871700in}{1.324297in}}%
\pgfpathlineto{\pgfqpoint{0.871865in}{1.322989in}}%
\pgfpathlineto{\pgfqpoint{0.872030in}{1.323877in}}%
\pgfpathlineto{\pgfqpoint{0.873020in}{0.554592in}}%
\pgfpathlineto{\pgfqpoint{0.873513in}{0.740266in}}%
\pgfpathlineto{\pgfqpoint{0.873677in}{0.778998in}}%
\pgfpathlineto{\pgfqpoint{0.873823in}{0.729424in}}%
\pgfpathlineto{\pgfqpoint{0.874297in}{0.559153in}}%
\pgfpathlineto{\pgfqpoint{0.874619in}{0.684745in}}%
\pgfpathlineto{\pgfqpoint{0.874784in}{0.766634in}}%
\pgfpathlineto{\pgfqpoint{0.875094in}{0.692854in}}%
\pgfpathlineto{\pgfqpoint{0.875412in}{0.564066in}}%
\pgfpathlineto{\pgfqpoint{0.875738in}{0.637546in}}%
\pgfpathlineto{\pgfqpoint{0.876067in}{0.777925in}}%
\pgfpathlineto{\pgfqpoint{0.876383in}{0.646041in}}%
\pgfpathlineto{\pgfqpoint{0.876710in}{0.566855in}}%
\pgfpathlineto{\pgfqpoint{0.877030in}{0.690553in}}%
\pgfpathlineto{\pgfqpoint{0.877195in}{0.767069in}}%
\pgfpathlineto{\pgfqpoint{0.877506in}{0.691638in}}%
\pgfpathlineto{\pgfqpoint{0.877824in}{0.570452in}}%
\pgfpathlineto{\pgfqpoint{0.878149in}{0.648245in}}%
\pgfpathlineto{\pgfqpoint{0.878477in}{0.776184in}}%
\pgfpathlineto{\pgfqpoint{0.879048in}{0.738775in}}%
\pgfpathlineto{\pgfqpoint{0.879286in}{0.525655in}}%
\pgfpathlineto{\pgfqpoint{0.879826in}{0.525679in}}%
\pgfpathlineto{\pgfqpoint{0.887190in}{0.525780in}}%
\pgfpathlineto{\pgfqpoint{0.887278in}{0.525951in}}%
\pgfpathlineto{\pgfqpoint{0.887516in}{1.324944in}}%
\pgfpathlineto{\pgfqpoint{0.888739in}{1.323115in}}%
\pgfpathlineto{\pgfqpoint{0.888904in}{1.323682in}}%
\pgfpathlineto{\pgfqpoint{0.889069in}{1.310742in}}%
\pgfpathlineto{\pgfqpoint{0.889233in}{1.322751in}}%
\pgfpathlineto{\pgfqpoint{0.889775in}{0.519433in}}%
\pgfpathlineto{\pgfqpoint{0.890086in}{0.687245in}}%
\pgfpathlineto{\pgfqpoint{0.890416in}{0.877722in}}%
\pgfpathlineto{\pgfqpoint{0.890695in}{0.617197in}}%
\pgfpathlineto{\pgfqpoint{0.891023in}{0.529491in}}%
\pgfpathlineto{\pgfqpoint{0.891181in}{0.599224in}}%
\pgfpathlineto{\pgfqpoint{0.891511in}{0.882165in}}%
\pgfpathlineto{\pgfqpoint{0.891817in}{0.691293in}}%
\pgfpathlineto{\pgfqpoint{0.892131in}{0.527657in}}%
\pgfpathlineto{\pgfqpoint{0.892457in}{0.674008in}}%
\pgfpathlineto{\pgfqpoint{0.892787in}{0.877604in}}%
\pgfpathlineto{\pgfqpoint{0.893067in}{0.636946in}}%
\pgfpathlineto{\pgfqpoint{0.893389in}{0.535615in}}%
\pgfpathlineto{\pgfqpoint{0.893550in}{0.598338in}}%
\pgfpathlineto{\pgfqpoint{0.893879in}{0.871509in}}%
\pgfpathlineto{\pgfqpoint{0.894183in}{0.707428in}}%
\pgfpathlineto{\pgfqpoint{0.894497in}{0.536885in}}%
\pgfpathlineto{\pgfqpoint{0.894821in}{0.671509in}}%
\pgfpathlineto{\pgfqpoint{0.895105in}{0.877047in}}%
\pgfpathlineto{\pgfqpoint{0.895529in}{0.731112in}}%
\pgfpathlineto{\pgfqpoint{0.895933in}{0.525667in}}%
\pgfpathlineto{\pgfqpoint{0.896427in}{0.525675in}}%
\pgfpathlineto{\pgfqpoint{0.903731in}{0.525879in}}%
\pgfpathlineto{\pgfqpoint{0.903819in}{0.548196in}}%
\pgfpathlineto{\pgfqpoint{0.903977in}{1.325143in}}%
\pgfpathlineto{\pgfqpoint{0.904697in}{1.324792in}}%
\pgfpathlineto{\pgfqpoint{0.905522in}{1.322453in}}%
\pgfpathlineto{\pgfqpoint{0.905686in}{1.323912in}}%
\pgfpathlineto{\pgfqpoint{0.906701in}{0.504005in}}%
\pgfpathlineto{\pgfqpoint{0.906825in}{0.543438in}}%
\pgfpathlineto{\pgfqpoint{0.906957in}{0.557064in}}%
\pgfpathlineto{\pgfqpoint{0.907451in}{0.901033in}}%
\pgfpathlineto{\pgfqpoint{0.907741in}{0.674249in}}%
\pgfpathlineto{\pgfqpoint{0.908053in}{0.548159in}}%
\pgfpathlineto{\pgfqpoint{0.908373in}{0.746377in}}%
\pgfpathlineto{\pgfqpoint{0.908538in}{0.882644in}}%
\pgfpathlineto{\pgfqpoint{0.908839in}{0.740784in}}%
\pgfpathlineto{\pgfqpoint{0.909151in}{0.552651in}}%
\pgfpathlineto{\pgfqpoint{0.909472in}{0.686269in}}%
\pgfpathlineto{\pgfqpoint{0.909802in}{0.889979in}}%
\pgfpathlineto{\pgfqpoint{0.910081in}{0.664350in}}%
\pgfpathlineto{\pgfqpoint{0.910399in}{0.560928in}}%
\pgfpathlineto{\pgfqpoint{0.910558in}{0.631062in}}%
\pgfpathlineto{\pgfqpoint{0.910879in}{0.885050in}}%
\pgfpathlineto{\pgfqpoint{0.911184in}{0.716529in}}%
\pgfpathlineto{\pgfqpoint{0.911497in}{0.560431in}}%
\pgfpathlineto{\pgfqpoint{0.911792in}{0.780041in}}%
\pgfpathlineto{\pgfqpoint{0.911980in}{1.024761in}}%
\pgfpathlineto{\pgfqpoint{0.912272in}{0.525666in}}%
\pgfpathlineto{\pgfqpoint{0.912434in}{0.525666in}}%
\pgfpathlineto{\pgfqpoint{0.920231in}{0.525965in}}%
\pgfpathlineto{\pgfqpoint{0.920406in}{0.927967in}}%
\pgfpathlineto{\pgfqpoint{0.920481in}{1.325232in}}%
\pgfpathlineto{\pgfqpoint{0.921218in}{1.324920in}}%
\pgfpathlineto{\pgfqpoint{0.922372in}{1.322900in}}%
\pgfpathlineto{\pgfqpoint{0.922537in}{1.323823in}}%
\pgfpathlineto{\pgfqpoint{0.923622in}{0.636920in}}%
\pgfpathlineto{\pgfqpoint{0.924088in}{0.796850in}}%
\pgfpathlineto{\pgfqpoint{0.924253in}{0.809820in}}%
\pgfpathlineto{\pgfqpoint{0.924392in}{0.761072in}}%
\pgfpathlineto{\pgfqpoint{0.924702in}{0.643981in}}%
\pgfpathlineto{\pgfqpoint{0.925022in}{0.710304in}}%
\pgfpathlineto{\pgfqpoint{0.925331in}{0.814200in}}%
\pgfpathlineto{\pgfqpoint{0.925618in}{0.714481in}}%
\pgfpathlineto{\pgfqpoint{0.925936in}{0.646668in}}%
\pgfpathlineto{\pgfqpoint{0.926221in}{0.748774in}}%
\pgfpathlineto{\pgfqpoint{0.926374in}{0.805869in}}%
\pgfpathlineto{\pgfqpoint{0.926696in}{0.738616in}}%
\pgfpathlineto{\pgfqpoint{0.927006in}{0.647144in}}%
\pgfpathlineto{\pgfqpoint{0.927319in}{0.736453in}}%
\pgfpathlineto{\pgfqpoint{0.927634in}{0.808060in}}%
\pgfpathlineto{\pgfqpoint{0.927923in}{0.694630in}}%
\pgfpathlineto{\pgfqpoint{0.928193in}{0.654831in}}%
\pgfpathlineto{\pgfqpoint{0.928277in}{0.742029in}}%
\pgfpathlineto{\pgfqpoint{0.928457in}{0.909960in}}%
\pgfpathlineto{\pgfqpoint{0.928757in}{0.525681in}}%
\pgfpathlineto{\pgfqpoint{0.928804in}{0.525690in}}%
\pgfpathlineto{\pgfqpoint{0.933183in}{0.525732in}}%
\pgfpathlineto{\pgfqpoint{0.936707in}{0.525967in}}%
\pgfpathlineto{\pgfqpoint{0.936874in}{0.876407in}}%
\pgfpathlineto{\pgfqpoint{0.936949in}{1.325351in}}%
\pgfpathlineto{\pgfqpoint{0.937670in}{1.325092in}}%
\pgfpathlineto{\pgfqpoint{0.938824in}{1.323655in}}%
\pgfpathlineto{\pgfqpoint{0.938989in}{1.324299in}}%
\pgfpathlineto{\pgfqpoint{0.939154in}{1.322848in}}%
\pgfpathlineto{\pgfqpoint{0.939319in}{1.323944in}}%
\pgfpathlineto{\pgfqpoint{0.940624in}{0.549030in}}%
\pgfpathlineto{\pgfqpoint{0.940779in}{0.637642in}}%
\pgfpathlineto{\pgfqpoint{0.941096in}{0.983624in}}%
\pgfpathlineto{\pgfqpoint{0.941399in}{0.737118in}}%
\pgfpathlineto{\pgfqpoint{0.941707in}{0.548078in}}%
\pgfpathlineto{\pgfqpoint{0.942025in}{0.773836in}}%
\pgfpathlineto{\pgfqpoint{0.942190in}{0.957897in}}%
\pgfpathlineto{\pgfqpoint{0.942628in}{0.633680in}}%
\pgfpathlineto{\pgfqpoint{0.942793in}{0.555430in}}%
\pgfpathlineto{\pgfqpoint{0.943110in}{0.719335in}}%
\pgfpathlineto{\pgfqpoint{0.943440in}{0.970104in}}%
\pgfpathlineto{\pgfqpoint{0.943712in}{0.680591in}}%
\pgfpathlineto{\pgfqpoint{0.944026in}{0.567846in}}%
\pgfpathlineto{\pgfqpoint{0.944176in}{0.659189in}}%
\pgfpathlineto{\pgfqpoint{0.944491in}{0.970797in}}%
\pgfpathlineto{\pgfqpoint{0.944986in}{0.733207in}}%
\pgfpathlineto{\pgfqpoint{0.945195in}{0.525654in}}%
\pgfpathlineto{\pgfqpoint{0.945782in}{0.525669in}}%
\pgfpathlineto{\pgfqpoint{0.953201in}{0.526032in}}%
\pgfpathlineto{\pgfqpoint{0.953377in}{0.931661in}}%
\pgfpathlineto{\pgfqpoint{0.953451in}{1.325401in}}%
\pgfpathlineto{\pgfqpoint{0.954188in}{1.325159in}}%
\pgfpathlineto{\pgfqpoint{0.955342in}{1.324072in}}%
\pgfpathlineto{\pgfqpoint{0.955506in}{1.324467in}}%
\pgfpathlineto{\pgfqpoint{0.955671in}{1.323624in}}%
\pgfpathlineto{\pgfqpoint{0.955836in}{1.324185in}}%
\pgfpathlineto{\pgfqpoint{0.956001in}{1.322815in}}%
\pgfpathlineto{\pgfqpoint{0.956166in}{1.323795in}}%
\pgfpathlineto{\pgfqpoint{0.957340in}{0.658192in}}%
\pgfpathlineto{\pgfqpoint{0.957627in}{0.778622in}}%
\pgfpathlineto{\pgfqpoint{0.957940in}{0.855141in}}%
\pgfpathlineto{\pgfqpoint{0.958223in}{0.706096in}}%
\pgfpathlineto{\pgfqpoint{0.958388in}{0.661494in}}%
\pgfpathlineto{\pgfqpoint{0.958701in}{0.763906in}}%
\pgfpathlineto{\pgfqpoint{0.959012in}{0.860997in}}%
\pgfpathlineto{\pgfqpoint{0.959291in}{0.723782in}}%
\pgfpathlineto{\pgfqpoint{0.959445in}{0.669584in}}%
\pgfpathlineto{\pgfqpoint{0.959760in}{0.743413in}}%
\pgfpathlineto{\pgfqpoint{0.960067in}{0.864828in}}%
\pgfpathlineto{\pgfqpoint{0.960348in}{0.747391in}}%
\pgfpathlineto{\pgfqpoint{0.960661in}{0.672388in}}%
\pgfpathlineto{\pgfqpoint{0.960938in}{0.793628in}}%
\pgfpathlineto{\pgfqpoint{0.961293in}{0.940417in}}%
\pgfpathlineto{\pgfqpoint{0.961502in}{0.742677in}}%
\pgfpathlineto{\pgfqpoint{0.961972in}{0.525671in}}%
\pgfpathlineto{\pgfqpoint{0.962302in}{0.525674in}}%
\pgfpathlineto{\pgfqpoint{0.969670in}{0.525995in}}%
\pgfpathlineto{\pgfqpoint{0.969758in}{0.551824in}}%
\pgfpathlineto{\pgfqpoint{0.969916in}{1.325494in}}%
\pgfpathlineto{\pgfqpoint{0.970638in}{1.325286in}}%
\pgfpathlineto{\pgfqpoint{0.971792in}{1.324282in}}%
\pgfpathlineto{\pgfqpoint{0.971956in}{1.324726in}}%
\pgfpathlineto{\pgfqpoint{0.972121in}{1.323944in}}%
\pgfpathlineto{\pgfqpoint{0.972286in}{1.324522in}}%
\pgfpathlineto{\pgfqpoint{0.972781in}{1.322410in}}%
\pgfpathlineto{\pgfqpoint{0.972946in}{1.323916in}}%
\pgfpathlineto{\pgfqpoint{0.974149in}{0.523253in}}%
\pgfpathlineto{\pgfqpoint{0.974617in}{1.100147in}}%
\pgfpathlineto{\pgfqpoint{0.974911in}{0.784490in}}%
\pgfpathlineto{\pgfqpoint{0.975216in}{0.524783in}}%
\pgfpathlineto{\pgfqpoint{0.975527in}{0.779192in}}%
\pgfpathlineto{\pgfqpoint{0.975856in}{1.068111in}}%
\pgfpathlineto{\pgfqpoint{0.976120in}{0.647322in}}%
\pgfpathlineto{\pgfqpoint{0.976276in}{0.537675in}}%
\pgfpathlineto{\pgfqpoint{0.976572in}{0.674865in}}%
\pgfpathlineto{\pgfqpoint{0.976894in}{1.096083in}}%
\pgfpathlineto{\pgfqpoint{0.977177in}{0.730172in}}%
\pgfpathlineto{\pgfqpoint{0.977476in}{0.538654in}}%
\pgfpathlineto{\pgfqpoint{0.977732in}{0.848519in}}%
\pgfpathlineto{\pgfqpoint{0.977914in}{1.215682in}}%
\pgfpathlineto{\pgfqpoint{0.978391in}{0.525671in}}%
\pgfpathlineto{\pgfqpoint{0.986155in}{0.526023in}}%
\pgfpathlineto{\pgfqpoint{0.986243in}{0.552658in}}%
\pgfpathlineto{\pgfqpoint{0.986401in}{1.325563in}}%
\pgfpathlineto{\pgfqpoint{0.987123in}{1.325375in}}%
\pgfpathlineto{\pgfqpoint{0.988277in}{1.324525in}}%
\pgfpathlineto{\pgfqpoint{0.988442in}{1.324897in}}%
\pgfpathlineto{\pgfqpoint{0.988607in}{1.324277in}}%
\pgfpathlineto{\pgfqpoint{0.988772in}{1.324724in}}%
\pgfpathlineto{\pgfqpoint{0.989596in}{1.322510in}}%
\pgfpathlineto{\pgfqpoint{0.989761in}{1.323943in}}%
\pgfpathlineto{\pgfqpoint{0.990988in}{0.522746in}}%
\pgfpathlineto{\pgfqpoint{0.991137in}{0.585528in}}%
\pgfpathlineto{\pgfqpoint{0.991457in}{1.098395in}}%
\pgfpathlineto{\pgfqpoint{0.991884in}{0.631882in}}%
\pgfpathlineto{\pgfqpoint{0.992049in}{0.531113in}}%
\pgfpathlineto{\pgfqpoint{0.992364in}{0.749871in}}%
\pgfpathlineto{\pgfqpoint{0.992694in}{1.112460in}}%
\pgfpathlineto{\pgfqpoint{0.992957in}{0.683832in}}%
\pgfpathlineto{\pgfqpoint{0.993244in}{0.541509in}}%
\pgfpathlineto{\pgfqpoint{0.993389in}{0.641455in}}%
\pgfpathlineto{\pgfqpoint{0.993701in}{1.110277in}}%
\pgfpathlineto{\pgfqpoint{0.994132in}{0.603508in}}%
\pgfpathlineto{\pgfqpoint{0.994687in}{0.525659in}}%
\pgfpathlineto{\pgfqpoint{0.994456in}{0.754923in}}%
\pgfpathlineto{\pgfqpoint{0.994905in}{0.525702in}}%
\pgfpathlineto{\pgfqpoint{0.998697in}{0.525742in}}%
\pgfpathlineto{\pgfqpoint{1.002659in}{0.526151in}}%
\pgfpathlineto{\pgfqpoint{1.002908in}{1.325601in}}%
\pgfpathlineto{\pgfqpoint{1.004305in}{1.325232in}}%
\pgfpathlineto{\pgfqpoint{1.005458in}{1.324352in}}%
\pgfpathlineto{\pgfqpoint{1.005623in}{1.324660in}}%
\pgfpathlineto{\pgfqpoint{1.005788in}{1.324060in}}%
\pgfpathlineto{\pgfqpoint{1.005953in}{1.324451in}}%
\pgfpathlineto{\pgfqpoint{1.006777in}{1.319638in}}%
\pgfpathlineto{\pgfqpoint{1.006942in}{1.323291in}}%
\pgfpathlineto{\pgfqpoint{1.007878in}{0.580554in}}%
\pgfpathlineto{\pgfqpoint{1.008492in}{1.037026in}}%
\pgfpathlineto{\pgfqpoint{1.008756in}{0.690551in}}%
\pgfpathlineto{\pgfqpoint{1.008911in}{0.591381in}}%
\pgfpathlineto{\pgfqpoint{1.009203in}{0.750437in}}%
\pgfpathlineto{\pgfqpoint{1.009519in}{1.058263in}}%
\pgfpathlineto{\pgfqpoint{1.009789in}{0.749256in}}%
\pgfpathlineto{\pgfqpoint{1.010076in}{0.597497in}}%
\pgfpathlineto{\pgfqpoint{1.010367in}{0.875726in}}%
\pgfpathlineto{\pgfqpoint{1.010678in}{1.042498in}}%
\pgfpathlineto{\pgfqpoint{1.010955in}{0.745180in}}%
\pgfpathlineto{\pgfqpoint{1.011243in}{0.525672in}}%
\pgfpathlineto{\pgfqpoint{1.011656in}{0.525681in}}%
\pgfpathlineto{\pgfqpoint{1.019124in}{0.526052in}}%
\pgfpathlineto{\pgfqpoint{1.019212in}{0.551385in}}%
\pgfpathlineto{\pgfqpoint{1.019370in}{1.325646in}}%
\pgfpathlineto{\pgfqpoint{1.020094in}{1.325482in}}%
\pgfpathlineto{\pgfqpoint{1.021248in}{1.324773in}}%
\pgfpathlineto{\pgfqpoint{1.021907in}{1.324375in}}%
\pgfpathlineto{\pgfqpoint{1.022072in}{1.324802in}}%
\pgfpathlineto{\pgfqpoint{1.022896in}{1.323098in}}%
\pgfpathlineto{\pgfqpoint{1.023061in}{1.324150in}}%
\pgfpathlineto{\pgfqpoint{1.023226in}{1.321350in}}%
\pgfpathlineto{\pgfqpoint{1.023391in}{1.323789in}}%
\pgfpathlineto{\pgfqpoint{1.024397in}{0.505207in}}%
\pgfpathlineto{\pgfqpoint{1.024693in}{0.757155in}}%
\pgfpathlineto{\pgfqpoint{1.025022in}{1.176477in}}%
\pgfpathlineto{\pgfqpoint{1.025284in}{0.695331in}}%
\pgfpathlineto{\pgfqpoint{1.025564in}{0.537427in}}%
\pgfpathlineto{\pgfqpoint{1.025707in}{0.641335in}}%
\pgfpathlineto{\pgfqpoint{1.026017in}{1.170158in}}%
\pgfpathlineto{\pgfqpoint{1.026448in}{0.603690in}}%
\pgfpathlineto{\pgfqpoint{1.026612in}{0.538446in}}%
\pgfpathlineto{\pgfqpoint{1.026765in}{0.612744in}}%
\pgfpathlineto{\pgfqpoint{1.027186in}{1.240681in}}%
\pgfpathlineto{\pgfqpoint{1.027494in}{0.634607in}}%
\pgfpathlineto{\pgfqpoint{1.027603in}{0.525661in}}%
\pgfpathlineto{\pgfqpoint{1.028197in}{0.525681in}}%
\pgfpathlineto{\pgfqpoint{1.035622in}{0.526137in}}%
\pgfpathlineto{\pgfqpoint{1.035800in}{0.934812in}}%
\pgfpathlineto{\pgfqpoint{1.035875in}{1.325653in}}%
\pgfpathlineto{\pgfqpoint{1.036614in}{1.325496in}}%
\pgfpathlineto{\pgfqpoint{1.037768in}{1.324902in}}%
\pgfpathlineto{\pgfqpoint{1.038592in}{1.324847in}}%
\pgfpathlineto{\pgfqpoint{1.039746in}{1.323041in}}%
\pgfpathlineto{\pgfqpoint{1.039911in}{1.323938in}}%
\pgfpathlineto{\pgfqpoint{1.040076in}{1.321123in}}%
\pgfpathlineto{\pgfqpoint{1.040240in}{1.323488in}}%
\pgfpathlineto{\pgfqpoint{1.041135in}{0.505711in}}%
\pgfpathlineto{\pgfqpoint{1.041285in}{0.535906in}}%
\pgfpathlineto{\pgfqpoint{1.041750in}{1.247028in}}%
\pgfpathlineto{\pgfqpoint{1.042019in}{0.741344in}}%
\pgfpathlineto{\pgfqpoint{1.042290in}{0.521249in}}%
\pgfpathlineto{\pgfqpoint{1.042583in}{0.851991in}}%
\pgfpathlineto{\pgfqpoint{1.042746in}{1.185399in}}%
\pgfpathlineto{\pgfqpoint{1.043168in}{0.640805in}}%
\pgfpathlineto{\pgfqpoint{1.043333in}{0.529218in}}%
\pgfpathlineto{\pgfqpoint{1.043637in}{0.831968in}}%
\pgfpathlineto{\pgfqpoint{1.043817in}{1.320651in}}%
\pgfpathlineto{\pgfqpoint{1.044269in}{0.525674in}}%
\pgfpathlineto{\pgfqpoint{1.052101in}{0.526130in}}%
\pgfpathlineto{\pgfqpoint{1.052270in}{0.890115in}}%
\pgfpathlineto{\pgfqpoint{1.052345in}{1.325721in}}%
\pgfpathlineto{\pgfqpoint{1.053072in}{1.325578in}}%
\pgfpathlineto{\pgfqpoint{1.054226in}{1.325016in}}%
\pgfpathlineto{\pgfqpoint{1.055050in}{1.325010in}}%
\pgfpathlineto{\pgfqpoint{1.055875in}{1.324039in}}%
\pgfpathlineto{\pgfqpoint{1.056040in}{1.324539in}}%
\pgfpathlineto{\pgfqpoint{1.056204in}{1.323652in}}%
\pgfpathlineto{\pgfqpoint{1.056369in}{1.324318in}}%
\pgfpathlineto{\pgfqpoint{1.056864in}{1.321126in}}%
\pgfpathlineto{\pgfqpoint{1.057029in}{1.323645in}}%
\pgfpathlineto{\pgfqpoint{1.058034in}{0.507587in}}%
\pgfpathlineto{\pgfqpoint{1.058333in}{0.900164in}}%
\pgfpathlineto{\pgfqpoint{1.058497in}{1.255918in}}%
\pgfpathlineto{\pgfqpoint{1.058918in}{0.606714in}}%
\pgfpathlineto{\pgfqpoint{1.059083in}{0.514547in}}%
\pgfpathlineto{\pgfqpoint{1.059232in}{0.560900in}}%
\pgfpathlineto{\pgfqpoint{1.059712in}{1.233831in}}%
\pgfpathlineto{\pgfqpoint{1.059953in}{0.697769in}}%
\pgfpathlineto{\pgfqpoint{1.060689in}{0.525690in}}%
\pgfpathlineto{\pgfqpoint{1.060367in}{0.821848in}}%
\pgfpathlineto{\pgfqpoint{1.060784in}{0.525699in}}%
\pgfpathlineto{\pgfqpoint{1.067873in}{0.525844in}}%
\pgfpathlineto{\pgfqpoint{1.068608in}{0.526342in}}%
\pgfpathlineto{\pgfqpoint{1.068851in}{1.325749in}}%
\pgfpathlineto{\pgfqpoint{1.070078in}{1.325364in}}%
\pgfpathlineto{\pgfqpoint{1.070902in}{1.325295in}}%
\pgfpathlineto{\pgfqpoint{1.072056in}{1.324627in}}%
\pgfpathlineto{\pgfqpoint{1.072551in}{1.324667in}}%
\pgfpathlineto{\pgfqpoint{1.073705in}{1.322525in}}%
\pgfpathlineto{\pgfqpoint{1.073870in}{1.323525in}}%
\pgfpathlineto{\pgfqpoint{1.074848in}{0.733187in}}%
\pgfpathlineto{\pgfqpoint{1.075111in}{0.883803in}}%
\pgfpathlineto{\pgfqpoint{1.075262in}{0.962288in}}%
\pgfpathlineto{\pgfqpoint{1.075549in}{0.855660in}}%
\pgfpathlineto{\pgfqpoint{1.075845in}{0.731559in}}%
\pgfpathlineto{\pgfqpoint{1.076122in}{0.862402in}}%
\pgfpathlineto{\pgfqpoint{1.076407in}{0.958577in}}%
\pgfpathlineto{\pgfqpoint{1.076867in}{0.874237in}}%
\pgfpathlineto{\pgfqpoint{1.077380in}{0.525678in}}%
\pgfpathlineto{\pgfqpoint{1.077710in}{0.525683in}}%
\pgfpathlineto{\pgfqpoint{1.085108in}{0.526586in}}%
\pgfpathlineto{\pgfqpoint{1.085346in}{1.325741in}}%
\pgfpathlineto{\pgfqpoint{1.086407in}{1.325537in}}%
\pgfpathlineto{\pgfqpoint{1.087891in}{1.324981in}}%
\pgfpathlineto{\pgfqpoint{1.088715in}{1.324833in}}%
\pgfpathlineto{\pgfqpoint{1.089869in}{1.323638in}}%
\pgfpathlineto{\pgfqpoint{1.090034in}{1.323987in}}%
\pgfpathlineto{\pgfqpoint{1.090199in}{1.323028in}}%
\pgfpathlineto{\pgfqpoint{1.090364in}{1.323592in}}%
\pgfpathlineto{\pgfqpoint{1.090528in}{1.321325in}}%
\pgfpathlineto{\pgfqpoint{1.090693in}{1.322939in}}%
\pgfpathlineto{\pgfqpoint{1.091578in}{0.702978in}}%
\pgfpathlineto{\pgfqpoint{1.091991in}{1.009874in}}%
\pgfpathlineto{\pgfqpoint{1.092394in}{0.767189in}}%
\pgfpathlineto{\pgfqpoint{1.092559in}{0.702616in}}%
\pgfpathlineto{\pgfqpoint{1.092842in}{0.860542in}}%
\pgfpathlineto{\pgfqpoint{1.093172in}{1.110779in}}%
\pgfpathlineto{\pgfqpoint{1.093408in}{0.673904in}}%
\pgfpathlineto{\pgfqpoint{1.093643in}{0.525674in}}%
\pgfpathlineto{\pgfqpoint{1.094188in}{0.525680in}}%
\pgfpathlineto{\pgfqpoint{1.101591in}{0.526549in}}%
\pgfpathlineto{\pgfqpoint{1.101827in}{1.325758in}}%
\pgfpathlineto{\pgfqpoint{1.102886in}{1.325558in}}%
\pgfpathlineto{\pgfqpoint{1.104699in}{1.324906in}}%
\pgfpathlineto{\pgfqpoint{1.105523in}{1.324747in}}%
\pgfpathlineto{\pgfqpoint{1.106677in}{1.323392in}}%
\pgfpathlineto{\pgfqpoint{1.106842in}{1.323806in}}%
\pgfpathlineto{\pgfqpoint{1.107007in}{1.322549in}}%
\pgfpathlineto{\pgfqpoint{1.107172in}{1.323323in}}%
\pgfpathlineto{\pgfqpoint{1.107934in}{0.673318in}}%
\pgfpathlineto{\pgfqpoint{1.108356in}{0.971424in}}%
\pgfpathlineto{\pgfqpoint{1.108511in}{1.082110in}}%
\pgfpathlineto{\pgfqpoint{1.108783in}{0.838728in}}%
\pgfpathlineto{\pgfqpoint{1.109079in}{0.672858in}}%
\pgfpathlineto{\pgfqpoint{1.109352in}{0.920238in}}%
\pgfpathlineto{\pgfqpoint{1.109657in}{1.143307in}}%
\pgfpathlineto{\pgfqpoint{1.109896in}{0.669923in}}%
\pgfpathlineto{\pgfqpoint{1.110026in}{0.525666in}}%
\pgfpathlineto{\pgfqpoint{1.110751in}{0.525681in}}%
\pgfpathlineto{\pgfqpoint{1.118050in}{0.526250in}}%
\pgfpathlineto{\pgfqpoint{1.118302in}{1.325800in}}%
\pgfpathlineto{\pgfqpoint{1.119700in}{1.325547in}}%
\pgfpathlineto{\pgfqpoint{1.120854in}{1.325035in}}%
\pgfpathlineto{\pgfqpoint{1.121679in}{1.325006in}}%
\pgfpathlineto{\pgfqpoint{1.122833in}{1.323895in}}%
\pgfpathlineto{\pgfqpoint{1.122997in}{1.324374in}}%
\pgfpathlineto{\pgfqpoint{1.123162in}{1.323471in}}%
\pgfpathlineto{\pgfqpoint{1.123327in}{1.324128in}}%
\pgfpathlineto{\pgfqpoint{1.123492in}{1.322738in}}%
\pgfpathlineto{\pgfqpoint{1.123657in}{1.323805in}}%
\pgfpathlineto{\pgfqpoint{1.123822in}{1.301530in}}%
\pgfpathlineto{\pgfqpoint{1.123986in}{1.323310in}}%
\pgfpathlineto{\pgfqpoint{1.124893in}{0.573528in}}%
\pgfpathlineto{\pgfqpoint{1.125182in}{0.834211in}}%
\pgfpathlineto{\pgfqpoint{1.125498in}{1.209267in}}%
\pgfpathlineto{\pgfqpoint{1.125754in}{0.742206in}}%
\pgfpathlineto{\pgfqpoint{1.126011in}{0.587102in}}%
\pgfpathlineto{\pgfqpoint{1.126241in}{1.052810in}}%
\pgfpathlineto{\pgfqpoint{1.126306in}{1.171931in}}%
\pgfpathlineto{\pgfqpoint{1.126717in}{0.525685in}}%
\pgfpathlineto{\pgfqpoint{1.134535in}{0.526277in}}%
\pgfpathlineto{\pgfqpoint{1.134787in}{1.325845in}}%
\pgfpathlineto{\pgfqpoint{1.136186in}{1.325608in}}%
\pgfpathlineto{\pgfqpoint{1.137340in}{1.325137in}}%
\pgfpathlineto{\pgfqpoint{1.138164in}{1.325116in}}%
\pgfpathlineto{\pgfqpoint{1.139318in}{1.324198in}}%
\pgfpathlineto{\pgfqpoint{1.139483in}{1.324577in}}%
\pgfpathlineto{\pgfqpoint{1.139647in}{1.323904in}}%
\pgfpathlineto{\pgfqpoint{1.139812in}{1.324383in}}%
\pgfpathlineto{\pgfqpoint{1.140636in}{1.319155in}}%
\pgfpathlineto{\pgfqpoint{1.140801in}{1.323376in}}%
\pgfpathlineto{\pgfqpoint{1.141683in}{0.546456in}}%
\pgfpathlineto{\pgfqpoint{1.141824in}{0.566985in}}%
\pgfpathlineto{\pgfqpoint{1.142278in}{1.301654in}}%
\pgfpathlineto{\pgfqpoint{1.142623in}{0.730199in}}%
\pgfpathlineto{\pgfqpoint{1.143008in}{0.525650in}}%
\pgfpathlineto{\pgfqpoint{1.143981in}{0.525686in}}%
\pgfpathlineto{\pgfqpoint{1.151017in}{0.526261in}}%
\pgfpathlineto{\pgfqpoint{1.151270in}{1.325855in}}%
\pgfpathlineto{\pgfqpoint{1.152669in}{1.325623in}}%
\pgfpathlineto{\pgfqpoint{1.153823in}{1.325159in}}%
\pgfpathlineto{\pgfqpoint{1.154647in}{1.325148in}}%
\pgfpathlineto{\pgfqpoint{1.155801in}{1.324268in}}%
\pgfpathlineto{\pgfqpoint{1.155966in}{1.324638in}}%
\pgfpathlineto{\pgfqpoint{1.156131in}{1.324001in}}%
\pgfpathlineto{\pgfqpoint{1.156296in}{1.324459in}}%
\pgfpathlineto{\pgfqpoint{1.157120in}{1.321716in}}%
\pgfpathlineto{\pgfqpoint{1.157285in}{1.323584in}}%
\pgfpathlineto{\pgfqpoint{1.158356in}{0.653735in}}%
\pgfpathlineto{\pgfqpoint{1.158780in}{1.149203in}}%
\pgfpathlineto{\pgfqpoint{1.159245in}{0.808877in}}%
\pgfpathlineto{\pgfqpoint{1.159402in}{0.590846in}}%
\pgfpathlineto{\pgfqpoint{1.159709in}{0.525675in}}%
\pgfpathlineto{\pgfqpoint{1.160204in}{0.525696in}}%
\pgfpathlineto{\pgfqpoint{1.167396in}{0.525941in}}%
\pgfpathlineto{\pgfqpoint{1.167489in}{0.526179in}}%
\pgfpathlineto{\pgfqpoint{1.167577in}{0.556612in}}%
\pgfpathlineto{\pgfqpoint{1.167735in}{1.325881in}}%
\pgfpathlineto{\pgfqpoint{1.168460in}{1.325772in}}%
\pgfpathlineto{\pgfqpoint{1.169614in}{1.325343in}}%
\pgfpathlineto{\pgfqpoint{1.170438in}{1.325380in}}%
\pgfpathlineto{\pgfqpoint{1.171592in}{1.324666in}}%
\pgfpathlineto{\pgfqpoint{1.172252in}{1.324287in}}%
\pgfpathlineto{\pgfqpoint{1.172417in}{1.324742in}}%
\pgfpathlineto{\pgfqpoint{1.173241in}{1.323152in}}%
\pgfpathlineto{\pgfqpoint{1.173406in}{1.324165in}}%
\pgfpathlineto{\pgfqpoint{1.173571in}{1.322012in}}%
\pgfpathlineto{\pgfqpoint{1.173735in}{1.323867in}}%
\pgfpathlineto{\pgfqpoint{1.174927in}{0.582738in}}%
\pgfpathlineto{\pgfqpoint{1.175465in}{1.233928in}}%
\pgfpathlineto{\pgfqpoint{1.175805in}{0.752747in}}%
\pgfpathlineto{\pgfqpoint{1.175984in}{0.525602in}}%
\pgfpathlineto{\pgfqpoint{1.176606in}{0.525673in}}%
\pgfpathlineto{\pgfqpoint{1.182375in}{0.525849in}}%
\pgfpathlineto{\pgfqpoint{1.183986in}{0.526260in}}%
\pgfpathlineto{\pgfqpoint{1.184164in}{0.939014in}}%
\pgfpathlineto{\pgfqpoint{1.184240in}{1.325865in}}%
\pgfpathlineto{\pgfqpoint{1.184979in}{1.325756in}}%
\pgfpathlineto{\pgfqpoint{1.186133in}{1.325367in}}%
\pgfpathlineto{\pgfqpoint{1.186957in}{1.325363in}}%
\pgfpathlineto{\pgfqpoint{1.188111in}{1.324725in}}%
\pgfpathlineto{\pgfqpoint{1.188936in}{1.324719in}}%
\pgfpathlineto{\pgfqpoint{1.189760in}{1.323416in}}%
\pgfpathlineto{\pgfqpoint{1.189925in}{1.324133in}}%
\pgfpathlineto{\pgfqpoint{1.190090in}{1.322675in}}%
\pgfpathlineto{\pgfqpoint{1.190254in}{1.323828in}}%
\pgfpathlineto{\pgfqpoint{1.190419in}{1.296031in}}%
\pgfpathlineto{\pgfqpoint{1.190584in}{1.323368in}}%
\pgfpathlineto{\pgfqpoint{1.191470in}{0.544865in}}%
\pgfpathlineto{\pgfqpoint{1.191612in}{0.576485in}}%
\pgfpathlineto{\pgfqpoint{1.191992in}{1.321466in}}%
\pgfpathlineto{\pgfqpoint{1.192371in}{0.603281in}}%
\pgfpathlineto{\pgfqpoint{1.192722in}{0.525679in}}%
\pgfpathlineto{\pgfqpoint{1.193051in}{0.525683in}}%
\pgfpathlineto{\pgfqpoint{1.200457in}{0.526177in}}%
\pgfpathlineto{\pgfqpoint{1.200545in}{0.554480in}}%
\pgfpathlineto{\pgfqpoint{1.200703in}{1.325892in}}%
\pgfpathlineto{\pgfqpoint{1.201429in}{1.325787in}}%
\pgfpathlineto{\pgfqpoint{1.202582in}{1.325366in}}%
\pgfpathlineto{\pgfqpoint{1.203407in}{1.325411in}}%
\pgfpathlineto{\pgfqpoint{1.204561in}{1.324730in}}%
\pgfpathlineto{\pgfqpoint{1.205220in}{1.324386in}}%
\pgfpathlineto{\pgfqpoint{1.205385in}{1.324816in}}%
\pgfpathlineto{\pgfqpoint{1.206209in}{1.323459in}}%
\pgfpathlineto{\pgfqpoint{1.206374in}{1.324305in}}%
\pgfpathlineto{\pgfqpoint{1.206539in}{1.322771in}}%
\pgfpathlineto{\pgfqpoint{1.206704in}{1.324056in}}%
\pgfpathlineto{\pgfqpoint{1.206869in}{1.319742in}}%
\pgfpathlineto{\pgfqpoint{1.207033in}{1.323726in}}%
\pgfpathlineto{\pgfqpoint{1.208042in}{0.573191in}}%
\pgfpathlineto{\pgfqpoint{1.208310in}{0.816814in}}%
\pgfpathlineto{\pgfqpoint{1.208525in}{1.321904in}}%
\pgfpathlineto{\pgfqpoint{1.208942in}{0.525727in}}%
\pgfpathlineto{\pgfqpoint{1.210142in}{0.525697in}}%
\pgfpathlineto{\pgfqpoint{1.216942in}{0.526185in}}%
\pgfpathlineto{\pgfqpoint{1.217030in}{0.554675in}}%
\pgfpathlineto{\pgfqpoint{1.217188in}{1.325906in}}%
\pgfpathlineto{\pgfqpoint{1.217914in}{1.325804in}}%
\pgfpathlineto{\pgfqpoint{1.219068in}{1.325394in}}%
\pgfpathlineto{\pgfqpoint{1.219892in}{1.325440in}}%
\pgfpathlineto{\pgfqpoint{1.221046in}{1.324789in}}%
\pgfpathlineto{\pgfqpoint{1.221870in}{1.324874in}}%
\pgfpathlineto{\pgfqpoint{1.222694in}{1.323670in}}%
\pgfpathlineto{\pgfqpoint{1.222859in}{1.324404in}}%
\pgfpathlineto{\pgfqpoint{1.223024in}{1.323155in}}%
\pgfpathlineto{\pgfqpoint{1.223189in}{1.324183in}}%
\pgfpathlineto{\pgfqpoint{1.223354in}{1.322079in}}%
\pgfpathlineto{\pgfqpoint{1.223518in}{1.323901in}}%
\pgfpathlineto{\pgfqpoint{1.224735in}{0.623464in}}%
\pgfpathlineto{\pgfqpoint{1.225004in}{1.005990in}}%
\pgfpathlineto{\pgfqpoint{1.225150in}{1.320364in}}%
\pgfpathlineto{\pgfqpoint{1.225521in}{0.525682in}}%
\pgfpathlineto{\pgfqpoint{1.225562in}{0.525701in}}%
\pgfpathlineto{\pgfqpoint{1.228282in}{0.525743in}}%
\pgfpathlineto{\pgfqpoint{1.233439in}{0.526278in}}%
\pgfpathlineto{\pgfqpoint{1.233604in}{0.894239in}}%
\pgfpathlineto{\pgfqpoint{1.233680in}{1.325914in}}%
\pgfpathlineto{\pgfqpoint{1.234406in}{1.325812in}}%
\pgfpathlineto{\pgfqpoint{1.235559in}{1.325451in}}%
\pgfpathlineto{\pgfqpoint{1.236384in}{1.325456in}}%
\pgfpathlineto{\pgfqpoint{1.237538in}{1.324897in}}%
\pgfpathlineto{\pgfqpoint{1.238362in}{1.324907in}}%
\pgfpathlineto{\pgfqpoint{1.239186in}{1.323973in}}%
\pgfpathlineto{\pgfqpoint{1.239351in}{1.324459in}}%
\pgfpathlineto{\pgfqpoint{1.239516in}{1.323619in}}%
\pgfpathlineto{\pgfqpoint{1.239681in}{1.324253in}}%
\pgfpathlineto{\pgfqpoint{1.240175in}{1.321857in}}%
\pgfpathlineto{\pgfqpoint{1.240340in}{1.323648in}}%
\pgfpathlineto{\pgfqpoint{1.241409in}{0.660952in}}%
\pgfpathlineto{\pgfqpoint{1.241466in}{0.708664in}}%
\pgfpathlineto{\pgfqpoint{1.241709in}{1.066127in}}%
\pgfpathlineto{\pgfqpoint{1.241974in}{0.525677in}}%
\pgfpathlineto{\pgfqpoint{1.242042in}{0.525721in}}%
\pgfpathlineto{\pgfqpoint{1.244485in}{0.525735in}}%
\pgfpathlineto{\pgfqpoint{1.249919in}{0.526260in}}%
\pgfpathlineto{\pgfqpoint{1.250086in}{0.888194in}}%
\pgfpathlineto{\pgfqpoint{1.250162in}{1.325940in}}%
\pgfpathlineto{\pgfqpoint{1.250888in}{1.325842in}}%
\pgfpathlineto{\pgfqpoint{1.252042in}{1.325480in}}%
\pgfpathlineto{\pgfqpoint{1.252866in}{1.325500in}}%
\pgfpathlineto{\pgfqpoint{1.254020in}{1.324950in}}%
\pgfpathlineto{\pgfqpoint{1.254844in}{1.324985in}}%
\pgfpathlineto{\pgfqpoint{1.255668in}{1.324106in}}%
\pgfpathlineto{\pgfqpoint{1.255833in}{1.324581in}}%
\pgfpathlineto{\pgfqpoint{1.255998in}{1.323805in}}%
\pgfpathlineto{\pgfqpoint{1.256163in}{1.324401in}}%
\pgfpathlineto{\pgfqpoint{1.256657in}{1.322631in}}%
\pgfpathlineto{\pgfqpoint{1.256822in}{1.323905in}}%
\pgfpathlineto{\pgfqpoint{1.256987in}{1.292348in}}%
\pgfpathlineto{\pgfqpoint{1.257152in}{1.323499in}}%
\pgfpathlineto{\pgfqpoint{1.257974in}{0.700458in}}%
\pgfpathlineto{\pgfqpoint{1.258133in}{0.762711in}}%
\pgfpathlineto{\pgfqpoint{1.258212in}{0.802799in}}%
\pgfpathlineto{\pgfqpoint{1.258561in}{0.525685in}}%
\pgfpathlineto{\pgfqpoint{1.266410in}{0.526305in}}%
\pgfpathlineto{\pgfqpoint{1.266588in}{0.940261in}}%
\pgfpathlineto{\pgfqpoint{1.266664in}{1.325938in}}%
\pgfpathlineto{\pgfqpoint{1.267404in}{1.325841in}}%
\pgfpathlineto{\pgfqpoint{1.268888in}{1.325427in}}%
\pgfpathlineto{\pgfqpoint{1.269712in}{1.325432in}}%
\pgfpathlineto{\pgfqpoint{1.270866in}{1.324863in}}%
\pgfpathlineto{\pgfqpoint{1.271690in}{1.324873in}}%
\pgfpathlineto{\pgfqpoint{1.272515in}{1.323905in}}%
\pgfpathlineto{\pgfqpoint{1.272679in}{1.324413in}}%
\pgfpathlineto{\pgfqpoint{1.272844in}{1.323528in}}%
\pgfpathlineto{\pgfqpoint{1.273009in}{1.324199in}}%
\pgfpathlineto{\pgfqpoint{1.273504in}{1.321296in}}%
\pgfpathlineto{\pgfqpoint{1.273668in}{1.323559in}}%
\pgfpathlineto{\pgfqpoint{1.274826in}{0.553615in}}%
\pgfpathlineto{\pgfqpoint{1.275316in}{0.525687in}}%
\pgfpathlineto{\pgfqpoint{1.275645in}{0.525692in}}%
\pgfpathlineto{\pgfqpoint{1.282876in}{0.526186in}}%
\pgfpathlineto{\pgfqpoint{1.282965in}{0.543889in}}%
\pgfpathlineto{\pgfqpoint{1.283122in}{1.325955in}}%
\pgfpathlineto{\pgfqpoint{1.283843in}{1.325862in}}%
\pgfpathlineto{\pgfqpoint{1.284997in}{1.325470in}}%
\pgfpathlineto{\pgfqpoint{1.285822in}{1.325530in}}%
\pgfpathlineto{\pgfqpoint{1.286975in}{1.324939in}}%
\pgfpathlineto{\pgfqpoint{1.287800in}{1.325040in}}%
\pgfpathlineto{\pgfqpoint{1.288624in}{1.324094in}}%
\pgfpathlineto{\pgfqpoint{1.288789in}{1.324666in}}%
\pgfpathlineto{\pgfqpoint{1.288954in}{1.323791in}}%
\pgfpathlineto{\pgfqpoint{1.289118in}{1.324504in}}%
\pgfpathlineto{\pgfqpoint{1.289613in}{1.322611in}}%
\pgfpathlineto{\pgfqpoint{1.289778in}{1.324072in}}%
\pgfpathlineto{\pgfqpoint{1.289943in}{1.288456in}}%
\pgfpathlineto{\pgfqpoint{1.290108in}{1.323738in}}%
\pgfpathlineto{\pgfqpoint{1.291101in}{0.740568in}}%
\pgfpathlineto{\pgfqpoint{1.291177in}{0.756942in}}%
\pgfpathlineto{\pgfqpoint{1.291434in}{0.525665in}}%
\pgfpathlineto{\pgfqpoint{1.292045in}{0.525703in}}%
\pgfpathlineto{\pgfqpoint{1.297815in}{0.525849in}}%
\pgfpathlineto{\pgfqpoint{1.299374in}{0.526266in}}%
\pgfpathlineto{\pgfqpoint{1.299540in}{0.883096in}}%
\pgfpathlineto{\pgfqpoint{1.299615in}{1.325949in}}%
\pgfpathlineto{\pgfqpoint{1.300338in}{1.325854in}}%
\pgfpathlineto{\pgfqpoint{1.301492in}{1.325503in}}%
\pgfpathlineto{\pgfqpoint{1.302316in}{1.325523in}}%
\pgfpathlineto{\pgfqpoint{1.303470in}{1.324998in}}%
\pgfpathlineto{\pgfqpoint{1.304294in}{1.325031in}}%
\pgfpathlineto{\pgfqpoint{1.305118in}{1.324229in}}%
\pgfpathlineto{\pgfqpoint{1.305283in}{1.324655in}}%
\pgfpathlineto{\pgfqpoint{1.305448in}{1.323970in}}%
\pgfpathlineto{\pgfqpoint{1.305613in}{1.324491in}}%
\pgfpathlineto{\pgfqpoint{1.306437in}{1.321972in}}%
\pgfpathlineto{\pgfqpoint{1.306602in}{1.323741in}}%
\pgfpathlineto{\pgfqpoint{1.307862in}{0.525681in}}%
\pgfpathlineto{\pgfqpoint{1.313636in}{0.525839in}}%
\pgfpathlineto{\pgfqpoint{1.315881in}{0.526480in}}%
\pgfpathlineto{\pgfqpoint{1.316124in}{1.325922in}}%
\pgfpathlineto{\pgfqpoint{1.317353in}{1.325641in}}%
\pgfpathlineto{\pgfqpoint{1.318507in}{1.325551in}}%
\pgfpathlineto{\pgfqpoint{1.319990in}{1.325027in}}%
\pgfpathlineto{\pgfqpoint{1.320814in}{1.324974in}}%
\pgfpathlineto{\pgfqpoint{1.321968in}{1.324056in}}%
\pgfpathlineto{\pgfqpoint{1.322133in}{1.324396in}}%
\pgfpathlineto{\pgfqpoint{1.322298in}{1.323743in}}%
\pgfpathlineto{\pgfqpoint{1.322463in}{1.324181in}}%
\pgfpathlineto{\pgfqpoint{1.322957in}{1.322473in}}%
\pgfpathlineto{\pgfqpoint{1.323122in}{1.323540in}}%
\pgfpathlineto{\pgfqpoint{1.324623in}{0.525680in}}%
\pgfpathlineto{\pgfqpoint{1.325447in}{0.525702in}}%
\pgfpathlineto{\pgfqpoint{1.332307in}{0.526093in}}%
\pgfpathlineto{\pgfqpoint{1.332396in}{0.527447in}}%
\pgfpathlineto{\pgfqpoint{1.332553in}{1.325913in}}%
\pgfpathlineto{\pgfqpoint{1.333419in}{1.325727in}}%
\pgfpathlineto{\pgfqpoint{1.335232in}{1.325482in}}%
\pgfpathlineto{\pgfqpoint{1.337045in}{1.324871in}}%
\pgfpathlineto{\pgfqpoint{1.337870in}{1.324723in}}%
\pgfpathlineto{\pgfqpoint{1.339353in}{1.323068in}}%
\pgfpathlineto{\pgfqpoint{1.339518in}{1.323545in}}%
\pgfpathlineto{\pgfqpoint{1.339683in}{1.321897in}}%
\pgfpathlineto{\pgfqpoint{1.339848in}{1.322968in}}%
\pgfpathlineto{\pgfqpoint{1.341017in}{0.525685in}}%
\pgfpathlineto{\pgfqpoint{1.348821in}{0.526220in}}%
\pgfpathlineto{\pgfqpoint{1.348909in}{0.555445in}}%
\pgfpathlineto{\pgfqpoint{1.349067in}{1.325964in}}%
\pgfpathlineto{\pgfqpoint{1.349793in}{1.325872in}}%
\pgfpathlineto{\pgfqpoint{1.350947in}{1.325504in}}%
\pgfpathlineto{\pgfqpoint{1.351771in}{1.325552in}}%
\pgfpathlineto{\pgfqpoint{1.352925in}{1.325006in}}%
\pgfpathlineto{\pgfqpoint{1.353749in}{1.325085in}}%
\pgfpathlineto{\pgfqpoint{1.354574in}{1.324261in}}%
\pgfpathlineto{\pgfqpoint{1.354739in}{1.324738in}}%
\pgfpathlineto{\pgfqpoint{1.354903in}{1.324014in}}%
\pgfpathlineto{\pgfqpoint{1.355068in}{1.324590in}}%
\pgfpathlineto{\pgfqpoint{1.355892in}{1.322304in}}%
\pgfpathlineto{\pgfqpoint{1.356057in}{1.323947in}}%
\pgfpathlineto{\pgfqpoint{1.356552in}{1.018060in}}%
\pgfpathlineto{\pgfqpoint{1.356826in}{1.320539in}}%
\pgfpathlineto{\pgfqpoint{1.356840in}{1.320723in}}%
\pgfpathlineto{\pgfqpoint{1.356883in}{1.318314in}}%
\pgfpathlineto{\pgfqpoint{1.357218in}{0.609378in}}%
\pgfpathlineto{\pgfqpoint{1.357505in}{0.525681in}}%
\pgfpathlineto{\pgfqpoint{1.357995in}{0.525687in}}%
\pgfpathlineto{\pgfqpoint{1.365305in}{0.526215in}}%
\pgfpathlineto{\pgfqpoint{1.365394in}{0.555321in}}%
\pgfpathlineto{\pgfqpoint{1.365552in}{1.325956in}}%
\pgfpathlineto{\pgfqpoint{1.366278in}{1.325863in}}%
\pgfpathlineto{\pgfqpoint{1.367432in}{1.325492in}}%
\pgfpathlineto{\pgfqpoint{1.368256in}{1.325541in}}%
\pgfpathlineto{\pgfqpoint{1.369410in}{1.324989in}}%
\pgfpathlineto{\pgfqpoint{1.370234in}{1.325071in}}%
\pgfpathlineto{\pgfqpoint{1.371058in}{1.324229in}}%
\pgfpathlineto{\pgfqpoint{1.371223in}{1.324719in}}%
\pgfpathlineto{\pgfqpoint{1.371388in}{1.323976in}}%
\pgfpathlineto{\pgfqpoint{1.371553in}{1.324568in}}%
\pgfpathlineto{\pgfqpoint{1.372377in}{1.322104in}}%
\pgfpathlineto{\pgfqpoint{1.372542in}{1.323910in}}%
\pgfpathlineto{\pgfqpoint{1.373037in}{1.068246in}}%
\pgfpathlineto{\pgfqpoint{1.373328in}{1.304398in}}%
\pgfpathlineto{\pgfqpoint{1.373345in}{1.308142in}}%
\pgfpathlineto{\pgfqpoint{1.373414in}{1.249274in}}%
\pgfpathlineto{\pgfqpoint{1.373959in}{0.525679in}}%
\pgfpathlineto{\pgfqpoint{1.374403in}{0.525688in}}%
\pgfpathlineto{\pgfqpoint{1.381773in}{0.526134in}}%
\pgfpathlineto{\pgfqpoint{1.381862in}{0.528981in}}%
\pgfpathlineto{\pgfqpoint{1.382018in}{1.325952in}}%
\pgfpathlineto{\pgfqpoint{1.382726in}{1.325865in}}%
\pgfpathlineto{\pgfqpoint{1.383880in}{1.325513in}}%
\pgfpathlineto{\pgfqpoint{1.384705in}{1.325545in}}%
\pgfpathlineto{\pgfqpoint{1.385859in}{1.325027in}}%
\pgfpathlineto{\pgfqpoint{1.386683in}{1.325080in}}%
\pgfpathlineto{\pgfqpoint{1.387507in}{1.324310in}}%
\pgfpathlineto{\pgfqpoint{1.387672in}{1.324733in}}%
\pgfpathlineto{\pgfqpoint{1.387837in}{1.324079in}}%
\pgfpathlineto{\pgfqpoint{1.388002in}{1.324585in}}%
\pgfpathlineto{\pgfqpoint{1.388826in}{1.322629in}}%
\pgfpathlineto{\pgfqpoint{1.388991in}{1.323945in}}%
\pgfpathlineto{\pgfqpoint{1.389156in}{1.307599in}}%
\pgfpathlineto{\pgfqpoint{1.389320in}{1.323586in}}%
\pgfpathlineto{\pgfqpoint{1.390371in}{0.525687in}}%
\pgfpathlineto{\pgfqpoint{1.396498in}{0.525850in}}%
\pgfpathlineto{\pgfqpoint{1.398275in}{0.526219in}}%
\pgfpathlineto{\pgfqpoint{1.398363in}{0.555428in}}%
\pgfpathlineto{\pgfqpoint{1.398522in}{1.325964in}}%
\pgfpathlineto{\pgfqpoint{1.399248in}{1.325873in}}%
\pgfpathlineto{\pgfqpoint{1.400402in}{1.325507in}}%
\pgfpathlineto{\pgfqpoint{1.401226in}{1.325557in}}%
\pgfpathlineto{\pgfqpoint{1.402380in}{1.325019in}}%
\pgfpathlineto{\pgfqpoint{1.403204in}{1.325099in}}%
\pgfpathlineto{\pgfqpoint{1.404028in}{1.324297in}}%
\pgfpathlineto{\pgfqpoint{1.404193in}{1.324761in}}%
\pgfpathlineto{\pgfqpoint{1.404358in}{1.324063in}}%
\pgfpathlineto{\pgfqpoint{1.404523in}{1.324619in}}%
\pgfpathlineto{\pgfqpoint{1.405347in}{1.322577in}}%
\pgfpathlineto{\pgfqpoint{1.405512in}{1.324010in}}%
\pgfpathlineto{\pgfqpoint{1.405677in}{1.287450in}}%
\pgfpathlineto{\pgfqpoint{1.405842in}{1.323662in}}%
\pgfpathlineto{\pgfqpoint{1.406779in}{0.525654in}}%
\pgfpathlineto{\pgfqpoint{1.406894in}{0.525702in}}%
\pgfpathlineto{\pgfqpoint{1.409649in}{0.525744in}}%
\pgfpathlineto{\pgfqpoint{1.414765in}{0.526256in}}%
\pgfpathlineto{\pgfqpoint{1.414951in}{0.942732in}}%
\pgfpathlineto{\pgfqpoint{1.415028in}{1.325962in}}%
\pgfpathlineto{\pgfqpoint{1.415778in}{1.325871in}}%
\pgfpathlineto{\pgfqpoint{1.416932in}{1.325515in}}%
\pgfpathlineto{\pgfqpoint{1.417756in}{1.325555in}}%
\pgfpathlineto{\pgfqpoint{1.418910in}{1.325033in}}%
\pgfpathlineto{\pgfqpoint{1.419734in}{1.325099in}}%
\pgfpathlineto{\pgfqpoint{1.420558in}{1.324329in}}%
\pgfpathlineto{\pgfqpoint{1.420723in}{1.324762in}}%
\pgfpathlineto{\pgfqpoint{1.420888in}{1.324104in}}%
\pgfpathlineto{\pgfqpoint{1.421053in}{1.324619in}}%
\pgfpathlineto{\pgfqpoint{1.421877in}{1.322746in}}%
\pgfpathlineto{\pgfqpoint{1.422042in}{1.324014in}}%
\pgfpathlineto{\pgfqpoint{1.422207in}{1.320251in}}%
\pgfpathlineto{\pgfqpoint{1.422372in}{1.323693in}}%
\pgfpathlineto{\pgfqpoint{1.423321in}{0.525672in}}%
\pgfpathlineto{\pgfqpoint{1.423476in}{0.525674in}}%
\pgfpathlineto{\pgfqpoint{1.431253in}{0.526276in}}%
\pgfpathlineto{\pgfqpoint{1.431418in}{0.883359in}}%
\pgfpathlineto{\pgfqpoint{1.431494in}{1.325965in}}%
\pgfpathlineto{\pgfqpoint{1.432217in}{1.325874in}}%
\pgfpathlineto{\pgfqpoint{1.433370in}{1.325538in}}%
\pgfpathlineto{\pgfqpoint{1.434195in}{1.325561in}}%
\pgfpathlineto{\pgfqpoint{1.435349in}{1.325072in}}%
\pgfpathlineto{\pgfqpoint{1.436173in}{1.325110in}}%
\pgfpathlineto{\pgfqpoint{1.436997in}{1.324407in}}%
\pgfpathlineto{\pgfqpoint{1.437327in}{1.324201in}}%
\pgfpathlineto{\pgfqpoint{1.437492in}{1.324638in}}%
\pgfpathlineto{\pgfqpoint{1.438316in}{1.323075in}}%
\pgfpathlineto{\pgfqpoint{1.438481in}{1.324049in}}%
\pgfpathlineto{\pgfqpoint{1.438646in}{1.321959in}}%
\pgfpathlineto{\pgfqpoint{1.438810in}{1.323742in}}%
\pgfpathlineto{\pgfqpoint{1.439947in}{0.525678in}}%
\pgfpathlineto{\pgfqpoint{1.440112in}{0.525689in}}%
\pgfpathlineto{\pgfqpoint{1.447675in}{0.526035in}}%
\pgfpathlineto{\pgfqpoint{1.447764in}{0.526561in}}%
\pgfpathlineto{\pgfqpoint{1.448005in}{1.325945in}}%
\pgfpathlineto{\pgfqpoint{1.449234in}{1.325682in}}%
\pgfpathlineto{\pgfqpoint{1.450718in}{1.325535in}}%
\pgfpathlineto{\pgfqpoint{1.452201in}{1.325025in}}%
\pgfpathlineto{\pgfqpoint{1.453026in}{1.324966in}}%
\pgfpathlineto{\pgfqpoint{1.454180in}{1.324089in}}%
\pgfpathlineto{\pgfqpoint{1.454344in}{1.324402in}}%
\pgfpathlineto{\pgfqpoint{1.454509in}{1.323797in}}%
\pgfpathlineto{\pgfqpoint{1.454674in}{1.324195in}}%
\pgfpathlineto{\pgfqpoint{1.455498in}{1.319676in}}%
\pgfpathlineto{\pgfqpoint{1.455603in}{1.323169in}}%
\pgfpathlineto{\pgfqpoint{1.456405in}{0.525682in}}%
\pgfpathlineto{\pgfqpoint{1.456726in}{0.525685in}}%
\pgfpathlineto{\pgfqpoint{1.464196in}{0.526142in}}%
\pgfpathlineto{\pgfqpoint{1.464285in}{0.528876in}}%
\pgfpathlineto{\pgfqpoint{1.464442in}{1.325971in}}%
\pgfpathlineto{\pgfqpoint{1.465150in}{1.325886in}}%
\pgfpathlineto{\pgfqpoint{1.466304in}{1.325550in}}%
\pgfpathlineto{\pgfqpoint{1.467128in}{1.325579in}}%
\pgfpathlineto{\pgfqpoint{1.468282in}{1.325094in}}%
\pgfpathlineto{\pgfqpoint{1.469106in}{1.325139in}}%
\pgfpathlineto{\pgfqpoint{1.469930in}{1.324451in}}%
\pgfpathlineto{\pgfqpoint{1.470260in}{1.324255in}}%
\pgfpathlineto{\pgfqpoint{1.470425in}{1.324687in}}%
\pgfpathlineto{\pgfqpoint{1.471249in}{1.323234in}}%
\pgfpathlineto{\pgfqpoint{1.471414in}{1.324137in}}%
\pgfpathlineto{\pgfqpoint{1.471579in}{1.322377in}}%
\pgfpathlineto{\pgfqpoint{1.471744in}{1.323859in}}%
\pgfpathlineto{\pgfqpoint{1.473099in}{0.525688in}}%
\pgfpathlineto{\pgfqpoint{1.473758in}{0.525700in}}%
\pgfpathlineto{\pgfqpoint{1.480699in}{0.526234in}}%
\pgfpathlineto{\pgfqpoint{1.480788in}{0.555743in}}%
\pgfpathlineto{\pgfqpoint{1.480946in}{1.325987in}}%
\pgfpathlineto{\pgfqpoint{1.481672in}{1.325899in}}%
\pgfpathlineto{\pgfqpoint{1.482826in}{1.325546in}}%
\pgfpathlineto{\pgfqpoint{1.483650in}{1.325595in}}%
\pgfpathlineto{\pgfqpoint{1.484804in}{1.325087in}}%
\pgfpathlineto{\pgfqpoint{1.485628in}{1.325165in}}%
\pgfpathlineto{\pgfqpoint{1.486453in}{1.324439in}}%
\pgfpathlineto{\pgfqpoint{1.486782in}{1.324241in}}%
\pgfpathlineto{\pgfqpoint{1.486947in}{1.324726in}}%
\pgfpathlineto{\pgfqpoint{1.487771in}{1.323199in}}%
\pgfpathlineto{\pgfqpoint{1.487936in}{1.324204in}}%
\pgfpathlineto{\pgfqpoint{1.488101in}{1.322296in}}%
\pgfpathlineto{\pgfqpoint{1.488266in}{1.323947in}}%
\pgfpathlineto{\pgfqpoint{1.489426in}{0.525679in}}%
\pgfpathlineto{\pgfqpoint{1.489921in}{0.525701in}}%
\pgfpathlineto{\pgfqpoint{1.497099in}{0.525987in}}%
\pgfpathlineto{\pgfqpoint{1.497197in}{0.526331in}}%
\pgfpathlineto{\pgfqpoint{1.497379in}{0.953786in}}%
\pgfpathlineto{\pgfqpoint{1.497455in}{1.325982in}}%
\pgfpathlineto{\pgfqpoint{1.498200in}{1.325893in}}%
\pgfpathlineto{\pgfqpoint{1.499354in}{1.325566in}}%
\pgfpathlineto{\pgfqpoint{1.500178in}{1.325589in}}%
\pgfpathlineto{\pgfqpoint{1.501332in}{1.325118in}}%
\pgfpathlineto{\pgfqpoint{1.502156in}{1.325155in}}%
\pgfpathlineto{\pgfqpoint{1.503310in}{1.324312in}}%
\pgfpathlineto{\pgfqpoint{1.503475in}{1.324712in}}%
\pgfpathlineto{\pgfqpoint{1.503640in}{1.324085in}}%
\pgfpathlineto{\pgfqpoint{1.503805in}{1.324563in}}%
\pgfpathlineto{\pgfqpoint{1.504629in}{1.322695in}}%
\pgfpathlineto{\pgfqpoint{1.504794in}{1.323917in}}%
\pgfpathlineto{\pgfqpoint{1.504959in}{1.319629in}}%
\pgfpathlineto{\pgfqpoint{1.505061in}{1.323636in}}%
\pgfpathlineto{\pgfqpoint{1.505964in}{0.525691in}}%
\pgfpathlineto{\pgfqpoint{1.506128in}{0.525692in}}%
\pgfpathlineto{\pgfqpoint{1.513669in}{0.526241in}}%
\pgfpathlineto{\pgfqpoint{1.513757in}{0.555902in}}%
\pgfpathlineto{\pgfqpoint{1.513916in}{1.325998in}}%
\pgfpathlineto{\pgfqpoint{1.514642in}{1.325911in}}%
\pgfpathlineto{\pgfqpoint{1.515796in}{1.325564in}}%
\pgfpathlineto{\pgfqpoint{1.516620in}{1.325612in}}%
\pgfpathlineto{\pgfqpoint{1.517774in}{1.325115in}}%
\pgfpathlineto{\pgfqpoint{1.518598in}{1.325190in}}%
\pgfpathlineto{\pgfqpoint{1.519423in}{1.324491in}}%
\pgfpathlineto{\pgfqpoint{1.519917in}{1.324766in}}%
\pgfpathlineto{\pgfqpoint{1.520082in}{1.324075in}}%
\pgfpathlineto{\pgfqpoint{1.520247in}{1.324626in}}%
\pgfpathlineto{\pgfqpoint{1.521071in}{1.322657in}}%
\pgfpathlineto{\pgfqpoint{1.521236in}{1.324031in}}%
\pgfpathlineto{\pgfqpoint{1.521401in}{1.318495in}}%
\pgfpathlineto{\pgfqpoint{1.521524in}{1.323761in}}%
\pgfpathlineto{\pgfqpoint{1.522435in}{0.525682in}}%
\pgfpathlineto{\pgfqpoint{1.522600in}{0.525704in}}%
\pgfpathlineto{\pgfqpoint{1.528534in}{0.525871in}}%
\pgfpathlineto{\pgfqpoint{1.530155in}{0.526253in}}%
\pgfpathlineto{\pgfqpoint{1.530243in}{0.558205in}}%
\pgfpathlineto{\pgfqpoint{1.530402in}{1.326002in}}%
\pgfpathlineto{\pgfqpoint{1.531128in}{1.325915in}}%
\pgfpathlineto{\pgfqpoint{1.532282in}{1.325573in}}%
\pgfpathlineto{\pgfqpoint{1.533106in}{1.325618in}}%
\pgfpathlineto{\pgfqpoint{1.534260in}{1.325129in}}%
\pgfpathlineto{\pgfqpoint{1.535084in}{1.325198in}}%
\pgfpathlineto{\pgfqpoint{1.535909in}{1.324517in}}%
\pgfpathlineto{\pgfqpoint{1.536568in}{1.324113in}}%
\pgfpathlineto{\pgfqpoint{1.536733in}{1.324639in}}%
\pgfpathlineto{\pgfqpoint{1.537557in}{1.322800in}}%
\pgfpathlineto{\pgfqpoint{1.537722in}{1.324054in}}%
\pgfpathlineto{\pgfqpoint{1.537887in}{1.320784in}}%
\pgfpathlineto{\pgfqpoint{1.538009in}{1.323794in}}%
\pgfpathlineto{\pgfqpoint{1.538923in}{0.525683in}}%
\pgfpathlineto{\pgfqpoint{1.539088in}{0.525705in}}%
\pgfpathlineto{\pgfqpoint{1.544857in}{0.525854in}}%
\pgfpathlineto{\pgfqpoint{1.546640in}{0.526255in}}%
\pgfpathlineto{\pgfqpoint{1.546728in}{0.558244in}}%
\pgfpathlineto{\pgfqpoint{1.546886in}{1.326005in}}%
\pgfpathlineto{\pgfqpoint{1.547613in}{1.325918in}}%
\pgfpathlineto{\pgfqpoint{1.548767in}{1.325577in}}%
\pgfpathlineto{\pgfqpoint{1.549591in}{1.325621in}}%
\pgfpathlineto{\pgfqpoint{1.550745in}{1.325135in}}%
\pgfpathlineto{\pgfqpoint{1.551569in}{1.325204in}}%
\pgfpathlineto{\pgfqpoint{1.552393in}{1.324527in}}%
\pgfpathlineto{\pgfqpoint{1.553053in}{1.324128in}}%
\pgfpathlineto{\pgfqpoint{1.553218in}{1.324648in}}%
\pgfpathlineto{\pgfqpoint{1.554042in}{1.322850in}}%
\pgfpathlineto{\pgfqpoint{1.554207in}{1.324070in}}%
\pgfpathlineto{\pgfqpoint{1.554372in}{1.321095in}}%
\pgfpathlineto{\pgfqpoint{1.554494in}{1.323813in}}%
\pgfpathlineto{\pgfqpoint{1.555430in}{0.525682in}}%
\pgfpathlineto{\pgfqpoint{1.555760in}{0.525688in}}%
\pgfpathlineto{\pgfqpoint{1.563132in}{0.526300in}}%
\pgfpathlineto{\pgfqpoint{1.563300in}{0.894273in}}%
\pgfpathlineto{\pgfqpoint{1.563375in}{1.326000in}}%
\pgfpathlineto{\pgfqpoint{1.564103in}{1.325912in}}%
\pgfpathlineto{\pgfqpoint{1.565257in}{1.325587in}}%
\pgfpathlineto{\pgfqpoint{1.566081in}{1.325614in}}%
\pgfpathlineto{\pgfqpoint{1.567235in}{1.325150in}}%
\pgfpathlineto{\pgfqpoint{1.568060in}{1.325192in}}%
\pgfpathlineto{\pgfqpoint{1.569214in}{1.324379in}}%
\pgfpathlineto{\pgfqpoint{1.569378in}{1.324767in}}%
\pgfpathlineto{\pgfqpoint{1.569543in}{1.324167in}}%
\pgfpathlineto{\pgfqpoint{1.569708in}{1.324627in}}%
\pgfpathlineto{\pgfqpoint{1.570532in}{1.322980in}}%
\pgfpathlineto{\pgfqpoint{1.570697in}{1.324032in}}%
\pgfpathlineto{\pgfqpoint{1.570862in}{1.321662in}}%
\pgfpathlineto{\pgfqpoint{1.570982in}{1.323769in}}%
\pgfpathlineto{\pgfqpoint{1.572053in}{0.525693in}}%
\pgfpathlineto{\pgfqpoint{1.572383in}{0.525698in}}%
\pgfpathlineto{\pgfqpoint{1.579602in}{0.526212in}}%
\pgfpathlineto{\pgfqpoint{1.579691in}{0.541373in}}%
\pgfpathlineto{\pgfqpoint{1.579848in}{1.326011in}}%
\pgfpathlineto{\pgfqpoint{1.580567in}{1.325927in}}%
\pgfpathlineto{\pgfqpoint{1.581721in}{1.325570in}}%
\pgfpathlineto{\pgfqpoint{1.582545in}{1.325632in}}%
\pgfpathlineto{\pgfqpoint{1.583699in}{1.325122in}}%
\pgfpathlineto{\pgfqpoint{1.584524in}{1.325220in}}%
\pgfpathlineto{\pgfqpoint{1.585348in}{1.324504in}}%
\pgfpathlineto{\pgfqpoint{1.585842in}{1.324809in}}%
\pgfpathlineto{\pgfqpoint{1.586007in}{1.324093in}}%
\pgfpathlineto{\pgfqpoint{1.586172in}{1.324674in}}%
\pgfpathlineto{\pgfqpoint{1.586996in}{1.322721in}}%
\pgfpathlineto{\pgfqpoint{1.587161in}{1.324114in}}%
\pgfpathlineto{\pgfqpoint{1.587326in}{1.320018in}}%
\pgfpathlineto{\pgfqpoint{1.587456in}{1.323862in}}%
\pgfpathlineto{\pgfqpoint{1.588376in}{0.525684in}}%
\pgfpathlineto{\pgfqpoint{1.588541in}{0.525706in}}%
\pgfpathlineto{\pgfqpoint{1.594311in}{0.525854in}}%
\pgfpathlineto{\pgfqpoint{1.596095in}{0.526256in}}%
\pgfpathlineto{\pgfqpoint{1.596183in}{0.558260in}}%
\pgfpathlineto{\pgfqpoint{1.596341in}{1.326006in}}%
\pgfpathlineto{\pgfqpoint{1.597067in}{1.325919in}}%
\pgfpathlineto{\pgfqpoint{1.598221in}{1.325578in}}%
\pgfpathlineto{\pgfqpoint{1.599046in}{1.325622in}}%
\pgfpathlineto{\pgfqpoint{1.600200in}{1.325135in}}%
\pgfpathlineto{\pgfqpoint{1.601024in}{1.325204in}}%
\pgfpathlineto{\pgfqpoint{1.601848in}{1.324527in}}%
\pgfpathlineto{\pgfqpoint{1.602507in}{1.324127in}}%
\pgfpathlineto{\pgfqpoint{1.602672in}{1.324647in}}%
\pgfpathlineto{\pgfqpoint{1.603497in}{1.322843in}}%
\pgfpathlineto{\pgfqpoint{1.603661in}{1.324067in}}%
\pgfpathlineto{\pgfqpoint{1.603826in}{1.321043in}}%
\pgfpathlineto{\pgfqpoint{1.603949in}{1.323809in}}%
\pgfpathlineto{\pgfqpoint{1.604883in}{0.525682in}}%
\pgfpathlineto{\pgfqpoint{1.605213in}{0.525688in}}%
\pgfpathlineto{\pgfqpoint{1.612586in}{0.526298in}}%
\pgfpathlineto{\pgfqpoint{1.612755in}{0.895248in}}%
\pgfpathlineto{\pgfqpoint{1.612830in}{1.325998in}}%
\pgfpathlineto{\pgfqpoint{1.613558in}{1.325909in}}%
\pgfpathlineto{\pgfqpoint{1.614712in}{1.325583in}}%
\pgfpathlineto{\pgfqpoint{1.615537in}{1.325610in}}%
\pgfpathlineto{\pgfqpoint{1.616691in}{1.325143in}}%
\pgfpathlineto{\pgfqpoint{1.617515in}{1.325186in}}%
\pgfpathlineto{\pgfqpoint{1.618339in}{1.324541in}}%
\pgfpathlineto{\pgfqpoint{1.618998in}{1.324147in}}%
\pgfpathlineto{\pgfqpoint{1.619163in}{1.324615in}}%
\pgfpathlineto{\pgfqpoint{1.619988in}{1.322911in}}%
\pgfpathlineto{\pgfqpoint{1.620152in}{1.324010in}}%
\pgfpathlineto{\pgfqpoint{1.620317in}{1.321382in}}%
\pgfpathlineto{\pgfqpoint{1.620437in}{1.323739in}}%
\pgfpathlineto{\pgfqpoint{1.621373in}{0.525681in}}%
\pgfpathlineto{\pgfqpoint{1.621702in}{0.525687in}}%
\pgfpathlineto{\pgfqpoint{1.629071in}{0.526296in}}%
\pgfpathlineto{\pgfqpoint{1.629238in}{0.890365in}}%
\pgfpathlineto{\pgfqpoint{1.629314in}{1.325994in}}%
\pgfpathlineto{\pgfqpoint{1.630040in}{1.325905in}}%
\pgfpathlineto{\pgfqpoint{1.631194in}{1.325579in}}%
\pgfpathlineto{\pgfqpoint{1.632018in}{1.325605in}}%
\pgfpathlineto{\pgfqpoint{1.633172in}{1.325138in}}%
\pgfpathlineto{\pgfqpoint{1.633997in}{1.325177in}}%
\pgfpathlineto{\pgfqpoint{1.635150in}{1.324351in}}%
\pgfpathlineto{\pgfqpoint{1.635315in}{1.324744in}}%
\pgfpathlineto{\pgfqpoint{1.635480in}{1.324132in}}%
\pgfpathlineto{\pgfqpoint{1.635645in}{1.324600in}}%
\pgfpathlineto{\pgfqpoint{1.636469in}{1.322859in}}%
\pgfpathlineto{\pgfqpoint{1.636634in}{1.323983in}}%
\pgfpathlineto{\pgfqpoint{1.636799in}{1.321132in}}%
\pgfpathlineto{\pgfqpoint{1.636920in}{1.323702in}}%
\pgfpathlineto{\pgfqpoint{1.637781in}{0.525690in}}%
\pgfpathlineto{\pgfqpoint{1.638111in}{0.525691in}}%
\pgfpathlineto{\pgfqpoint{1.645589in}{0.526732in}}%
\pgfpathlineto{\pgfqpoint{1.645826in}{1.325963in}}%
\pgfpathlineto{\pgfqpoint{1.646883in}{1.325826in}}%
\pgfpathlineto{\pgfqpoint{1.649026in}{1.325352in}}%
\pgfpathlineto{\pgfqpoint{1.650180in}{1.325199in}}%
\pgfpathlineto{\pgfqpoint{1.651664in}{1.324463in}}%
\pgfpathlineto{\pgfqpoint{1.652158in}{1.324477in}}%
\pgfpathlineto{\pgfqpoint{1.653312in}{1.322455in}}%
\pgfpathlineto{\pgfqpoint{1.653419in}{1.323401in}}%
\pgfpathlineto{\pgfqpoint{1.653526in}{1.320914in}}%
\pgfpathlineto{\pgfqpoint{1.653793in}{1.321307in}}%
\pgfpathlineto{\pgfqpoint{1.654401in}{0.525687in}}%
\pgfpathlineto{\pgfqpoint{1.654895in}{0.525701in}}%
\pgfpathlineto{\pgfqpoint{1.662004in}{0.526112in}}%
\pgfpathlineto{\pgfqpoint{1.662093in}{0.527555in}}%
\pgfpathlineto{\pgfqpoint{1.662251in}{1.325955in}}%
\pgfpathlineto{\pgfqpoint{1.662952in}{1.325869in}}%
\pgfpathlineto{\pgfqpoint{1.665095in}{1.325425in}}%
\pgfpathlineto{\pgfqpoint{1.666249in}{1.325277in}}%
\pgfpathlineto{\pgfqpoint{1.667732in}{1.324628in}}%
\pgfpathlineto{\pgfqpoint{1.668557in}{1.324467in}}%
\pgfpathlineto{\pgfqpoint{1.669711in}{1.322476in}}%
\pgfpathlineto{\pgfqpoint{1.669861in}{1.323307in}}%
\pgfpathlineto{\pgfqpoint{1.670328in}{1.230751in}}%
\pgfpathlineto{\pgfqpoint{1.670769in}{0.525687in}}%
\pgfpathlineto{\pgfqpoint{1.671099in}{0.525689in}}%
\pgfpathlineto{\pgfqpoint{1.678517in}{0.526233in}}%
\pgfpathlineto{\pgfqpoint{1.678606in}{0.555722in}}%
\pgfpathlineto{\pgfqpoint{1.678764in}{1.325985in}}%
\pgfpathlineto{\pgfqpoint{1.679490in}{1.325897in}}%
\pgfpathlineto{\pgfqpoint{1.680644in}{1.325543in}}%
\pgfpathlineto{\pgfqpoint{1.681468in}{1.325593in}}%
\pgfpathlineto{\pgfqpoint{1.682622in}{1.325082in}}%
\pgfpathlineto{\pgfqpoint{1.683447in}{1.325160in}}%
\pgfpathlineto{\pgfqpoint{1.684271in}{1.324429in}}%
\pgfpathlineto{\pgfqpoint{1.684601in}{1.324229in}}%
\pgfpathlineto{\pgfqpoint{1.684765in}{1.324719in}}%
\pgfpathlineto{\pgfqpoint{1.685590in}{1.323163in}}%
\pgfpathlineto{\pgfqpoint{1.685754in}{1.324190in}}%
\pgfpathlineto{\pgfqpoint{1.685919in}{1.322206in}}%
\pgfpathlineto{\pgfqpoint{1.686084in}{1.323929in}}%
\pgfpathlineto{\pgfqpoint{1.687242in}{0.525679in}}%
\pgfpathlineto{\pgfqpoint{1.687572in}{0.525684in}}%
\pgfpathlineto{\pgfqpoint{1.695017in}{0.526342in}}%
\pgfpathlineto{\pgfqpoint{1.695274in}{1.325978in}}%
\pgfpathlineto{\pgfqpoint{1.696679in}{1.325798in}}%
\pgfpathlineto{\pgfqpoint{1.697833in}{1.325434in}}%
\pgfpathlineto{\pgfqpoint{1.698657in}{1.325457in}}%
\pgfpathlineto{\pgfqpoint{1.699811in}{1.324906in}}%
\pgfpathlineto{\pgfqpoint{1.700635in}{1.324946in}}%
\pgfpathlineto{\pgfqpoint{1.701459in}{1.324068in}}%
\pgfpathlineto{\pgfqpoint{1.701624in}{1.324547in}}%
\pgfpathlineto{\pgfqpoint{1.701789in}{1.323771in}}%
\pgfpathlineto{\pgfqpoint{1.701954in}{1.324370in}}%
\pgfpathlineto{\pgfqpoint{1.702448in}{1.322628in}}%
\pgfpathlineto{\pgfqpoint{1.702613in}{1.323885in}}%
\pgfpathlineto{\pgfqpoint{1.702778in}{1.313231in}}%
\pgfpathlineto{\pgfqpoint{1.702879in}{1.323586in}}%
\pgfpathlineto{\pgfqpoint{1.703684in}{0.525688in}}%
\pgfpathlineto{\pgfqpoint{1.703982in}{0.525688in}}%
\pgfpathlineto{\pgfqpoint{1.711444in}{0.526069in}}%
\pgfpathlineto{\pgfqpoint{1.711534in}{0.526878in}}%
\pgfpathlineto{\pgfqpoint{1.711773in}{1.325960in}}%
\pgfpathlineto{\pgfqpoint{1.712675in}{1.325761in}}%
\pgfpathlineto{\pgfqpoint{1.714158in}{1.325616in}}%
\pgfpathlineto{\pgfqpoint{1.715642in}{1.325173in}}%
\pgfpathlineto{\pgfqpoint{1.716466in}{1.325106in}}%
\pgfpathlineto{\pgfqpoint{1.717620in}{1.324427in}}%
\pgfpathlineto{\pgfqpoint{1.718115in}{1.324474in}}%
\pgfpathlineto{\pgfqpoint{1.719268in}{1.322191in}}%
\pgfpathlineto{\pgfqpoint{1.719367in}{1.323411in}}%
\pgfpathlineto{\pgfqpoint{1.719465in}{1.319678in}}%
\pgfpathlineto{\pgfqpoint{1.719732in}{1.321145in}}%
\pgfpathlineto{\pgfqpoint{1.720424in}{0.525691in}}%
\pgfpathlineto{\pgfqpoint{1.720918in}{0.525700in}}%
\pgfpathlineto{\pgfqpoint{1.727972in}{0.526240in}}%
\pgfpathlineto{\pgfqpoint{1.728060in}{0.555873in}}%
\pgfpathlineto{\pgfqpoint{1.728219in}{1.325996in}}%
\pgfpathlineto{\pgfqpoint{1.728945in}{1.325909in}}%
\pgfpathlineto{\pgfqpoint{1.730099in}{1.325560in}}%
\pgfpathlineto{\pgfqpoint{1.730923in}{1.325609in}}%
\pgfpathlineto{\pgfqpoint{1.732077in}{1.325109in}}%
\pgfpathlineto{\pgfqpoint{1.732901in}{1.325186in}}%
\pgfpathlineto{\pgfqpoint{1.733726in}{1.324482in}}%
\pgfpathlineto{\pgfqpoint{1.734220in}{1.324759in}}%
\pgfpathlineto{\pgfqpoint{1.734385in}{1.324060in}}%
\pgfpathlineto{\pgfqpoint{1.734550in}{1.324617in}}%
\pgfpathlineto{\pgfqpoint{1.735374in}{1.322596in}}%
\pgfpathlineto{\pgfqpoint{1.735539in}{1.324015in}}%
\pgfpathlineto{\pgfqpoint{1.735704in}{1.300117in}}%
\pgfpathlineto{\pgfqpoint{1.735827in}{1.323723in}}%
\pgfpathlineto{\pgfqpoint{1.736725in}{0.525686in}}%
\pgfpathlineto{\pgfqpoint{1.736890in}{0.525699in}}%
\pgfpathlineto{\pgfqpoint{1.744387in}{0.526015in}}%
\pgfpathlineto{\pgfqpoint{1.744474in}{0.526368in}}%
\pgfpathlineto{\pgfqpoint{1.744726in}{1.325983in}}%
\pgfpathlineto{\pgfqpoint{1.746125in}{1.325804in}}%
\pgfpathlineto{\pgfqpoint{1.747609in}{1.325391in}}%
\pgfpathlineto{\pgfqpoint{1.748433in}{1.325398in}}%
\pgfpathlineto{\pgfqpoint{1.749587in}{1.324837in}}%
\pgfpathlineto{\pgfqpoint{1.750411in}{1.324847in}}%
\pgfpathlineto{\pgfqpoint{1.751236in}{1.323905in}}%
\pgfpathlineto{\pgfqpoint{1.751401in}{1.324398in}}%
\pgfpathlineto{\pgfqpoint{1.751565in}{1.323547in}}%
\pgfpathlineto{\pgfqpoint{1.751730in}{1.324190in}}%
\pgfpathlineto{\pgfqpoint{1.752225in}{1.321709in}}%
\pgfpathlineto{\pgfqpoint{1.752330in}{1.323651in}}%
\pgfpathlineto{\pgfqpoint{1.753432in}{0.525691in}}%
\pgfpathlineto{\pgfqpoint{1.754091in}{0.525704in}}%
\pgfpathlineto{\pgfqpoint{1.760908in}{0.526100in}}%
\pgfpathlineto{\pgfqpoint{1.760998in}{0.527262in}}%
\pgfpathlineto{\pgfqpoint{1.761155in}{1.325942in}}%
\pgfpathlineto{\pgfqpoint{1.762013in}{1.325808in}}%
\pgfpathlineto{\pgfqpoint{1.765475in}{1.325172in}}%
\pgfpathlineto{\pgfqpoint{1.768763in}{1.323163in}}%
\pgfpathlineto{\pgfqpoint{1.769187in}{1.321033in}}%
\pgfpathlineto{\pgfqpoint{1.769692in}{0.525683in}}%
\pgfpathlineto{\pgfqpoint{1.770351in}{0.525694in}}%
\pgfpathlineto{\pgfqpoint{1.777436in}{0.526304in}}%
\pgfpathlineto{\pgfqpoint{1.777602in}{0.887907in}}%
\pgfpathlineto{\pgfqpoint{1.777677in}{1.325985in}}%
\pgfpathlineto{\pgfqpoint{1.778401in}{1.325896in}}%
\pgfpathlineto{\pgfqpoint{1.779555in}{1.325574in}}%
\pgfpathlineto{\pgfqpoint{1.780379in}{1.325593in}}%
\pgfpathlineto{\pgfqpoint{1.781533in}{1.325132in}}%
\pgfpathlineto{\pgfqpoint{1.782358in}{1.325162in}}%
\pgfpathlineto{\pgfqpoint{1.783512in}{1.324341in}}%
\pgfpathlineto{\pgfqpoint{1.783676in}{1.324722in}}%
\pgfpathlineto{\pgfqpoint{1.783841in}{1.324121in}}%
\pgfpathlineto{\pgfqpoint{1.784006in}{1.324575in}}%
\pgfpathlineto{\pgfqpoint{1.784830in}{1.322827in}}%
\pgfpathlineto{\pgfqpoint{1.784995in}{1.323939in}}%
\pgfpathlineto{\pgfqpoint{1.785160in}{1.320959in}}%
\pgfpathlineto{\pgfqpoint{1.785282in}{1.323643in}}%
\pgfpathlineto{\pgfqpoint{1.786155in}{0.525692in}}%
\pgfpathlineto{\pgfqpoint{1.786485in}{0.525693in}}%
\pgfpathlineto{\pgfqpoint{1.793911in}{0.526241in}}%
\pgfpathlineto{\pgfqpoint{1.794000in}{0.555907in}}%
\pgfpathlineto{\pgfqpoint{1.794158in}{1.325998in}}%
\pgfpathlineto{\pgfqpoint{1.794884in}{1.325911in}}%
\pgfpathlineto{\pgfqpoint{1.796038in}{1.325564in}}%
\pgfpathlineto{\pgfqpoint{1.796863in}{1.325613in}}%
\pgfpathlineto{\pgfqpoint{1.798016in}{1.325115in}}%
\pgfpathlineto{\pgfqpoint{1.798841in}{1.325191in}}%
\pgfpathlineto{\pgfqpoint{1.799665in}{1.324492in}}%
\pgfpathlineto{\pgfqpoint{1.800159in}{1.324767in}}%
\pgfpathlineto{\pgfqpoint{1.800324in}{1.324076in}}%
\pgfpathlineto{\pgfqpoint{1.800489in}{1.324626in}}%
\pgfpathlineto{\pgfqpoint{1.801313in}{1.322662in}}%
\pgfpathlineto{\pgfqpoint{1.801478in}{1.324032in}}%
\pgfpathlineto{\pgfqpoint{1.801643in}{1.318715in}}%
\pgfpathlineto{\pgfqpoint{1.801766in}{1.323762in}}%
\pgfpathlineto{\pgfqpoint{1.802676in}{0.525681in}}%
\pgfpathlineto{\pgfqpoint{1.802841in}{0.525704in}}%
\pgfpathlineto{\pgfqpoint{1.808611in}{0.525853in}}%
\pgfpathlineto{\pgfqpoint{1.810398in}{0.526253in}}%
\pgfpathlineto{\pgfqpoint{1.810486in}{0.558206in}}%
\pgfpathlineto{\pgfqpoint{1.810644in}{1.326002in}}%
\pgfpathlineto{\pgfqpoint{1.811370in}{1.325915in}}%
\pgfpathlineto{\pgfqpoint{1.812524in}{1.325573in}}%
\pgfpathlineto{\pgfqpoint{1.813349in}{1.325618in}}%
\pgfpathlineto{\pgfqpoint{1.814503in}{1.325129in}}%
\pgfpathlineto{\pgfqpoint{1.815327in}{1.325198in}}%
\pgfpathlineto{\pgfqpoint{1.816151in}{1.324517in}}%
\pgfpathlineto{\pgfqpoint{1.816810in}{1.324112in}}%
\pgfpathlineto{\pgfqpoint{1.816975in}{1.324639in}}%
\pgfpathlineto{\pgfqpoint{1.817800in}{1.322798in}}%
\pgfpathlineto{\pgfqpoint{1.817964in}{1.324054in}}%
\pgfpathlineto{\pgfqpoint{1.818129in}{1.320769in}}%
\pgfpathlineto{\pgfqpoint{1.818252in}{1.323793in}}%
\pgfpathlineto{\pgfqpoint{1.819165in}{0.525683in}}%
\pgfpathlineto{\pgfqpoint{1.819329in}{0.525705in}}%
\pgfpathlineto{\pgfqpoint{1.825099in}{0.525854in}}%
\pgfpathlineto{\pgfqpoint{1.826883in}{0.526255in}}%
\pgfpathlineto{\pgfqpoint{1.826971in}{0.558241in}}%
\pgfpathlineto{\pgfqpoint{1.827129in}{1.326004in}}%
\pgfpathlineto{\pgfqpoint{1.827855in}{1.325917in}}%
\pgfpathlineto{\pgfqpoint{1.829009in}{1.325576in}}%
\pgfpathlineto{\pgfqpoint{1.829833in}{1.325621in}}%
\pgfpathlineto{\pgfqpoint{1.830987in}{1.325134in}}%
\pgfpathlineto{\pgfqpoint{1.831812in}{1.325203in}}%
\pgfpathlineto{\pgfqpoint{1.832636in}{1.324526in}}%
\pgfpathlineto{\pgfqpoint{1.833295in}{1.324126in}}%
\pgfpathlineto{\pgfqpoint{1.833460in}{1.324647in}}%
\pgfpathlineto{\pgfqpoint{1.834284in}{1.322844in}}%
\pgfpathlineto{\pgfqpoint{1.834449in}{1.324068in}}%
\pgfpathlineto{\pgfqpoint{1.834614in}{1.321058in}}%
\pgfpathlineto{\pgfqpoint{1.834737in}{1.323811in}}%
\pgfpathlineto{\pgfqpoint{1.835672in}{0.525682in}}%
\pgfpathlineto{\pgfqpoint{1.836001in}{0.525688in}}%
\pgfpathlineto{\pgfqpoint{1.843374in}{0.526300in}}%
\pgfpathlineto{\pgfqpoint{1.843543in}{0.895247in}}%
\pgfpathlineto{\pgfqpoint{1.843618in}{1.326000in}}%
\pgfpathlineto{\pgfqpoint{1.844346in}{1.325911in}}%
\pgfpathlineto{\pgfqpoint{1.845500in}{1.325586in}}%
\pgfpathlineto{\pgfqpoint{1.846325in}{1.325613in}}%
\pgfpathlineto{\pgfqpoint{1.847478in}{1.325149in}}%
\pgfpathlineto{\pgfqpoint{1.848303in}{1.325191in}}%
\pgfpathlineto{\pgfqpoint{1.849127in}{1.324552in}}%
\pgfpathlineto{\pgfqpoint{1.849786in}{1.324164in}}%
\pgfpathlineto{\pgfqpoint{1.849951in}{1.324626in}}%
\pgfpathlineto{\pgfqpoint{1.850775in}{1.322969in}}%
\pgfpathlineto{\pgfqpoint{1.850940in}{1.324030in}}%
\pgfpathlineto{\pgfqpoint{1.851105in}{1.321621in}}%
\pgfpathlineto{\pgfqpoint{1.851225in}{1.323766in}}%
\pgfpathlineto{\pgfqpoint{1.852086in}{0.525685in}}%
\pgfpathlineto{\pgfqpoint{1.852581in}{0.525706in}}%
\pgfpathlineto{\pgfqpoint{1.859764in}{0.525989in}}%
\pgfpathlineto{\pgfqpoint{1.859860in}{0.526311in}}%
\pgfpathlineto{\pgfqpoint{1.860044in}{0.949009in}}%
\pgfpathlineto{\pgfqpoint{1.860120in}{1.325997in}}%
\pgfpathlineto{\pgfqpoint{1.860866in}{1.325909in}}%
\pgfpathlineto{\pgfqpoint{1.862020in}{1.325581in}}%
\pgfpathlineto{\pgfqpoint{1.862844in}{1.325610in}}%
\pgfpathlineto{\pgfqpoint{1.863998in}{1.325141in}}%
\pgfpathlineto{\pgfqpoint{1.864823in}{1.325186in}}%
\pgfpathlineto{\pgfqpoint{1.865647in}{1.324537in}}%
\pgfpathlineto{\pgfqpoint{1.866306in}{1.324141in}}%
\pgfpathlineto{\pgfqpoint{1.866471in}{1.324616in}}%
\pgfpathlineto{\pgfqpoint{1.867295in}{1.322892in}}%
\pgfpathlineto{\pgfqpoint{1.867460in}{1.324012in}}%
\pgfpathlineto{\pgfqpoint{1.867625in}{1.321302in}}%
\pgfpathlineto{\pgfqpoint{1.867727in}{1.323760in}}%
\pgfpathlineto{\pgfqpoint{1.868605in}{0.525684in}}%
\pgfpathlineto{\pgfqpoint{1.869100in}{0.525707in}}%
\pgfpathlineto{\pgfqpoint{1.876245in}{0.525982in}}%
\pgfpathlineto{\pgfqpoint{1.876337in}{0.526254in}}%
\pgfpathlineto{\pgfqpoint{1.876425in}{0.558224in}}%
\pgfpathlineto{\pgfqpoint{1.876583in}{1.326003in}}%
\pgfpathlineto{\pgfqpoint{1.877310in}{1.325916in}}%
\pgfpathlineto{\pgfqpoint{1.878464in}{1.325574in}}%
\pgfpathlineto{\pgfqpoint{1.879288in}{1.325618in}}%
\pgfpathlineto{\pgfqpoint{1.880442in}{1.325129in}}%
\pgfpathlineto{\pgfqpoint{1.881266in}{1.325198in}}%
\pgfpathlineto{\pgfqpoint{1.882090in}{1.324516in}}%
\pgfpathlineto{\pgfqpoint{1.882750in}{1.324110in}}%
\pgfpathlineto{\pgfqpoint{1.882915in}{1.324637in}}%
\pgfpathlineto{\pgfqpoint{1.883739in}{1.322785in}}%
\pgfpathlineto{\pgfqpoint{1.883904in}{1.324050in}}%
\pgfpathlineto{\pgfqpoint{1.884069in}{1.320661in}}%
\pgfpathlineto{\pgfqpoint{1.884191in}{1.323787in}}%
\pgfpathlineto{\pgfqpoint{1.885102in}{0.525682in}}%
\pgfpathlineto{\pgfqpoint{1.885267in}{0.525705in}}%
\pgfpathlineto{\pgfqpoint{1.891037in}{0.525853in}}%
\pgfpathlineto{\pgfqpoint{1.892822in}{0.526253in}}%
\pgfpathlineto{\pgfqpoint{1.892910in}{0.558207in}}%
\pgfpathlineto{\pgfqpoint{1.893068in}{1.326002in}}%
\pgfpathlineto{\pgfqpoint{1.893795in}{1.325915in}}%
\pgfpathlineto{\pgfqpoint{1.894949in}{1.325572in}}%
\pgfpathlineto{\pgfqpoint{1.895773in}{1.325616in}}%
\pgfpathlineto{\pgfqpoint{1.896927in}{1.325126in}}%
\pgfpathlineto{\pgfqpoint{1.897751in}{1.325195in}}%
\pgfpathlineto{\pgfqpoint{1.898575in}{1.324510in}}%
\pgfpathlineto{\pgfqpoint{1.899235in}{1.324101in}}%
\pgfpathlineto{\pgfqpoint{1.899400in}{1.324632in}}%
\pgfpathlineto{\pgfqpoint{1.900224in}{1.322750in}}%
\pgfpathlineto{\pgfqpoint{1.900389in}{1.324040in}}%
\pgfpathlineto{\pgfqpoint{1.900553in}{1.320351in}}%
\pgfpathlineto{\pgfqpoint{1.900676in}{1.323774in}}%
\pgfpathlineto{\pgfqpoint{1.901581in}{0.525682in}}%
\pgfpathlineto{\pgfqpoint{1.901746in}{0.525704in}}%
\pgfpathlineto{\pgfqpoint{1.907516in}{0.525853in}}%
\pgfpathlineto{\pgfqpoint{1.909307in}{0.526252in}}%
\pgfpathlineto{\pgfqpoint{1.909395in}{0.558188in}}%
\pgfpathlineto{\pgfqpoint{1.909553in}{1.326001in}}%
\pgfpathlineto{\pgfqpoint{1.910280in}{1.325913in}}%
\pgfpathlineto{\pgfqpoint{1.911433in}{1.325570in}}%
\pgfpathlineto{\pgfqpoint{1.912258in}{1.325614in}}%
\pgfpathlineto{\pgfqpoint{1.913412in}{1.325123in}}%
\pgfpathlineto{\pgfqpoint{1.914236in}{1.325192in}}%
\pgfpathlineto{\pgfqpoint{1.915060in}{1.324504in}}%
\pgfpathlineto{\pgfqpoint{1.915555in}{1.324767in}}%
\pgfpathlineto{\pgfqpoint{1.915720in}{1.324092in}}%
\pgfpathlineto{\pgfqpoint{1.915884in}{1.324626in}}%
\pgfpathlineto{\pgfqpoint{1.916709in}{1.322716in}}%
\pgfpathlineto{\pgfqpoint{1.916873in}{1.324030in}}%
\pgfpathlineto{\pgfqpoint{1.917038in}{1.319936in}}%
\pgfpathlineto{\pgfqpoint{1.917161in}{1.323761in}}%
\pgfpathlineto{\pgfqpoint{1.918066in}{0.525681in}}%
\pgfpathlineto{\pgfqpoint{1.918231in}{0.525704in}}%
\pgfpathlineto{\pgfqpoint{1.924001in}{0.525852in}}%
\pgfpathlineto{\pgfqpoint{1.925792in}{0.526252in}}%
\pgfpathlineto{\pgfqpoint{1.925880in}{0.558173in}}%
\pgfpathlineto{\pgfqpoint{1.926038in}{1.325999in}}%
\pgfpathlineto{\pgfqpoint{1.926764in}{1.325912in}}%
\pgfpathlineto{\pgfqpoint{1.927918in}{1.325568in}}%
\pgfpathlineto{\pgfqpoint{1.928743in}{1.325613in}}%
\pgfpathlineto{\pgfqpoint{1.929896in}{1.325120in}}%
\pgfpathlineto{\pgfqpoint{1.930721in}{1.325190in}}%
\pgfpathlineto{\pgfqpoint{1.931545in}{1.324499in}}%
\pgfpathlineto{\pgfqpoint{1.932040in}{1.324763in}}%
\pgfpathlineto{\pgfqpoint{1.932204in}{1.324085in}}%
\pgfpathlineto{\pgfqpoint{1.932369in}{1.324622in}}%
\pgfpathlineto{\pgfqpoint{1.933193in}{1.322688in}}%
\pgfpathlineto{\pgfqpoint{1.933358in}{1.324022in}}%
\pgfpathlineto{\pgfqpoint{1.933523in}{1.319420in}}%
\pgfpathlineto{\pgfqpoint{1.933646in}{1.323750in}}%
\pgfpathlineto{\pgfqpoint{1.934533in}{0.525683in}}%
\pgfpathlineto{\pgfqpoint{1.934698in}{0.525702in}}%
\pgfpathlineto{\pgfqpoint{1.941127in}{0.525873in}}%
\pgfpathlineto{\pgfqpoint{1.942279in}{0.526268in}}%
\pgfpathlineto{\pgfqpoint{1.942451in}{0.896297in}}%
\pgfpathlineto{\pgfqpoint{1.942527in}{1.325997in}}%
\pgfpathlineto{\pgfqpoint{1.943256in}{1.325909in}}%
\pgfpathlineto{\pgfqpoint{1.944410in}{1.325568in}}%
\pgfpathlineto{\pgfqpoint{1.945235in}{1.325609in}}%
\pgfpathlineto{\pgfqpoint{1.946389in}{1.325120in}}%
\pgfpathlineto{\pgfqpoint{1.947213in}{1.325184in}}%
\pgfpathlineto{\pgfqpoint{1.948037in}{1.324499in}}%
\pgfpathlineto{\pgfqpoint{1.948696in}{1.324085in}}%
\pgfpathlineto{\pgfqpoint{1.948861in}{1.324611in}}%
\pgfpathlineto{\pgfqpoint{1.949686in}{1.322690in}}%
\pgfpathlineto{\pgfqpoint{1.949850in}{1.324003in}}%
\pgfpathlineto{\pgfqpoint{1.950015in}{1.319459in}}%
\pgfpathlineto{\pgfqpoint{1.950134in}{1.323729in}}%
\pgfpathlineto{\pgfqpoint{1.951012in}{0.525687in}}%
\pgfpathlineto{\pgfqpoint{1.951177in}{0.525698in}}%
\pgfpathlineto{\pgfqpoint{1.958696in}{0.526024in}}%
\pgfpathlineto{\pgfqpoint{1.958783in}{0.526427in}}%
\pgfpathlineto{\pgfqpoint{1.959029in}{1.325978in}}%
\pgfpathlineto{\pgfqpoint{1.960258in}{1.325704in}}%
\pgfpathlineto{\pgfqpoint{1.961412in}{1.325638in}}%
\pgfpathlineto{\pgfqpoint{1.962896in}{1.325159in}}%
\pgfpathlineto{\pgfqpoint{1.963720in}{1.325139in}}%
\pgfpathlineto{\pgfqpoint{1.964874in}{1.324397in}}%
\pgfpathlineto{\pgfqpoint{1.965204in}{1.324189in}}%
\pgfpathlineto{\pgfqpoint{1.965368in}{1.324533in}}%
\pgfpathlineto{\pgfqpoint{1.966193in}{1.323043in}}%
\pgfpathlineto{\pgfqpoint{1.966357in}{1.323857in}}%
\pgfpathlineto{\pgfqpoint{1.966522in}{1.321867in}}%
\pgfpathlineto{\pgfqpoint{1.966630in}{1.323551in}}%
\pgfpathlineto{\pgfqpoint{1.967613in}{0.525690in}}%
\pgfpathlineto{\pgfqpoint{1.968602in}{0.525710in}}%
\pgfpathlineto{\pgfqpoint{1.975245in}{0.526240in}}%
\pgfpathlineto{\pgfqpoint{1.975333in}{0.555869in}}%
\pgfpathlineto{\pgfqpoint{1.975492in}{1.325996in}}%
\pgfpathlineto{\pgfqpoint{1.976218in}{1.325908in}}%
\pgfpathlineto{\pgfqpoint{1.977372in}{1.325559in}}%
\pgfpathlineto{\pgfqpoint{1.978196in}{1.325607in}}%
\pgfpathlineto{\pgfqpoint{1.979350in}{1.325106in}}%
\pgfpathlineto{\pgfqpoint{1.980174in}{1.325182in}}%
\pgfpathlineto{\pgfqpoint{1.980998in}{1.324473in}}%
\pgfpathlineto{\pgfqpoint{1.981328in}{1.324281in}}%
\pgfpathlineto{\pgfqpoint{1.981493in}{1.324751in}}%
\pgfpathlineto{\pgfqpoint{1.982317in}{1.323302in}}%
\pgfpathlineto{\pgfqpoint{1.982482in}{1.324244in}}%
\pgfpathlineto{\pgfqpoint{1.982647in}{1.322528in}}%
\pgfpathlineto{\pgfqpoint{1.982812in}{1.323998in}}%
\pgfpathlineto{\pgfqpoint{1.983558in}{0.958274in}}%
\pgfpathlineto{\pgfqpoint{1.984184in}{0.525690in}}%
\pgfpathlineto{\pgfqpoint{1.984349in}{0.525697in}}%
\pgfpathlineto{\pgfqpoint{1.991724in}{0.526211in}}%
\pgfpathlineto{\pgfqpoint{1.991813in}{0.543650in}}%
\pgfpathlineto{\pgfqpoint{1.991971in}{1.326001in}}%
\pgfpathlineto{\pgfqpoint{1.992691in}{1.325915in}}%
\pgfpathlineto{\pgfqpoint{1.993845in}{1.325556in}}%
\pgfpathlineto{\pgfqpoint{1.994670in}{1.325617in}}%
\pgfpathlineto{\pgfqpoint{1.995823in}{1.325100in}}%
\pgfpathlineto{\pgfqpoint{1.996648in}{1.325196in}}%
\pgfpathlineto{\pgfqpoint{1.997472in}{1.324463in}}%
\pgfpathlineto{\pgfqpoint{1.997967in}{1.324774in}}%
\pgfpathlineto{\pgfqpoint{1.998131in}{1.324030in}}%
\pgfpathlineto{\pgfqpoint{1.998296in}{1.324634in}}%
\pgfpathlineto{\pgfqpoint{1.999120in}{1.322459in}}%
\pgfpathlineto{\pgfqpoint{1.999285in}{1.324044in}}%
\pgfpathlineto{\pgfqpoint{2.000184in}{0.550343in}}%
\pgfpathlineto{\pgfqpoint{2.000721in}{0.525689in}}%
\pgfpathlineto{\pgfqpoint{2.000886in}{0.525700in}}%
\pgfpathlineto{\pgfqpoint{2.008155in}{0.526033in}}%
\pgfpathlineto{\pgfqpoint{2.008242in}{0.526491in}}%
\pgfpathlineto{\pgfqpoint{2.008487in}{1.325972in}}%
\pgfpathlineto{\pgfqpoint{2.009715in}{1.325708in}}%
\pgfpathlineto{\pgfqpoint{2.010869in}{1.325634in}}%
\pgfpathlineto{\pgfqpoint{2.012353in}{1.325167in}}%
\pgfpathlineto{\pgfqpoint{2.013177in}{1.325133in}}%
\pgfpathlineto{\pgfqpoint{2.014331in}{1.324413in}}%
\pgfpathlineto{\pgfqpoint{2.014826in}{1.324522in}}%
\pgfpathlineto{\pgfqpoint{2.015980in}{1.322047in}}%
\pgfpathlineto{\pgfqpoint{2.016086in}{1.323527in}}%
\pgfpathlineto{\pgfqpoint{2.016525in}{1.004154in}}%
\pgfpathlineto{\pgfqpoint{2.016932in}{0.525682in}}%
\pgfpathlineto{\pgfqpoint{2.017261in}{0.525686in}}%
\pgfpathlineto{\pgfqpoint{2.024707in}{0.526292in}}%
\pgfpathlineto{\pgfqpoint{2.024876in}{0.894937in}}%
\pgfpathlineto{\pgfqpoint{2.024951in}{1.325989in}}%
\pgfpathlineto{\pgfqpoint{2.025680in}{1.325900in}}%
\pgfpathlineto{\pgfqpoint{2.026833in}{1.325570in}}%
\pgfpathlineto{\pgfqpoint{2.027658in}{1.325598in}}%
\pgfpathlineto{\pgfqpoint{2.028812in}{1.325124in}}%
\pgfpathlineto{\pgfqpoint{2.029636in}{1.325167in}}%
\pgfpathlineto{\pgfqpoint{2.030460in}{1.324507in}}%
\pgfpathlineto{\pgfqpoint{2.031120in}{1.324097in}}%
\pgfpathlineto{\pgfqpoint{2.031284in}{1.324584in}}%
\pgfpathlineto{\pgfqpoint{2.032109in}{1.322736in}}%
\pgfpathlineto{\pgfqpoint{2.032273in}{1.323953in}}%
\pgfpathlineto{\pgfqpoint{2.032438in}{1.320208in}}%
\pgfpathlineto{\pgfqpoint{2.032558in}{1.323664in}}%
\pgfpathlineto{\pgfqpoint{2.033442in}{0.525678in}}%
\pgfpathlineto{\pgfqpoint{2.033607in}{0.525703in}}%
\pgfpathlineto{\pgfqpoint{2.039212in}{0.525859in}}%
\pgfpathlineto{\pgfqpoint{2.041186in}{0.526249in}}%
\pgfpathlineto{\pgfqpoint{2.041274in}{0.558120in}}%
\pgfpathlineto{\pgfqpoint{2.041432in}{1.325996in}}%
\pgfpathlineto{\pgfqpoint{2.042158in}{1.325908in}}%
\pgfpathlineto{\pgfqpoint{2.043312in}{1.325563in}}%
\pgfpathlineto{\pgfqpoint{2.044136in}{1.325608in}}%
\pgfpathlineto{\pgfqpoint{2.045290in}{1.325113in}}%
\pgfpathlineto{\pgfqpoint{2.046115in}{1.325183in}}%
\pgfpathlineto{\pgfqpoint{2.046939in}{1.324486in}}%
\pgfpathlineto{\pgfqpoint{2.047433in}{1.324754in}}%
\pgfpathlineto{\pgfqpoint{2.047598in}{1.324067in}}%
\pgfpathlineto{\pgfqpoint{2.047763in}{1.324611in}}%
\pgfpathlineto{\pgfqpoint{2.048587in}{1.322619in}}%
\pgfpathlineto{\pgfqpoint{2.048752in}{1.324004in}}%
\pgfpathlineto{\pgfqpoint{2.048917in}{1.308834in}}%
\pgfpathlineto{\pgfqpoint{2.049040in}{1.323716in}}%
\pgfpathlineto{\pgfqpoint{2.049940in}{0.525679in}}%
\pgfpathlineto{\pgfqpoint{2.050105in}{0.525703in}}%
\pgfpathlineto{\pgfqpoint{2.055710in}{0.525860in}}%
\pgfpathlineto{\pgfqpoint{2.057670in}{0.526251in}}%
\pgfpathlineto{\pgfqpoint{2.057758in}{0.558150in}}%
\pgfpathlineto{\pgfqpoint{2.057917in}{1.325998in}}%
\pgfpathlineto{\pgfqpoint{2.058643in}{1.325910in}}%
\pgfpathlineto{\pgfqpoint{2.059797in}{1.325566in}}%
\pgfpathlineto{\pgfqpoint{2.060621in}{1.325611in}}%
\pgfpathlineto{\pgfqpoint{2.061775in}{1.325118in}}%
\pgfpathlineto{\pgfqpoint{2.062599in}{1.325188in}}%
\pgfpathlineto{\pgfqpoint{2.063424in}{1.324496in}}%
\pgfpathlineto{\pgfqpoint{2.063918in}{1.324761in}}%
\pgfpathlineto{\pgfqpoint{2.064083in}{1.324081in}}%
\pgfpathlineto{\pgfqpoint{2.064248in}{1.324620in}}%
\pgfpathlineto{\pgfqpoint{2.065072in}{1.322676in}}%
\pgfpathlineto{\pgfqpoint{2.065237in}{1.324019in}}%
\pgfpathlineto{\pgfqpoint{2.065402in}{1.319146in}}%
\pgfpathlineto{\pgfqpoint{2.065525in}{1.323746in}}%
\pgfpathlineto{\pgfqpoint{2.066669in}{0.525690in}}%
\pgfpathlineto{\pgfqpoint{2.074192in}{0.526669in}}%
\pgfpathlineto{\pgfqpoint{2.074427in}{1.325964in}}%
\pgfpathlineto{\pgfqpoint{2.075481in}{1.325828in}}%
\pgfpathlineto{\pgfqpoint{2.077624in}{1.325359in}}%
\pgfpathlineto{\pgfqpoint{2.078778in}{1.325204in}}%
\pgfpathlineto{\pgfqpoint{2.080262in}{1.324480in}}%
\pgfpathlineto{\pgfqpoint{2.081086in}{1.324300in}}%
\pgfpathlineto{\pgfqpoint{2.082132in}{1.321291in}}%
\pgfpathlineto{\pgfqpoint{2.082399in}{1.321537in}}%
\pgfpathlineto{\pgfqpoint{2.082830in}{0.525688in}}%
\pgfpathlineto{\pgfqpoint{2.083627in}{0.525706in}}%
\pgfpathlineto{\pgfqpoint{2.090578in}{0.526031in}}%
\pgfpathlineto{\pgfqpoint{2.090666in}{0.526479in}}%
\pgfpathlineto{\pgfqpoint{2.090910in}{1.325972in}}%
\pgfpathlineto{\pgfqpoint{2.092139in}{1.325707in}}%
\pgfpathlineto{\pgfqpoint{2.093293in}{1.325635in}}%
\pgfpathlineto{\pgfqpoint{2.094777in}{1.325166in}}%
\pgfpathlineto{\pgfqpoint{2.095601in}{1.325135in}}%
\pgfpathlineto{\pgfqpoint{2.096755in}{1.324411in}}%
\pgfpathlineto{\pgfqpoint{2.097085in}{1.324207in}}%
\pgfpathlineto{\pgfqpoint{2.097249in}{1.324525in}}%
\pgfpathlineto{\pgfqpoint{2.098403in}{1.322034in}}%
\pgfpathlineto{\pgfqpoint{2.098510in}{1.323535in}}%
\pgfpathlineto{\pgfqpoint{2.098617in}{1.289492in}}%
\pgfpathlineto{\pgfqpoint{2.098927in}{1.309341in}}%
\pgfpathlineto{\pgfqpoint{2.099627in}{0.525690in}}%
\pgfpathlineto{\pgfqpoint{2.107118in}{0.526208in}}%
\pgfpathlineto{\pgfqpoint{2.107206in}{0.542554in}}%
\pgfpathlineto{\pgfqpoint{2.107364in}{1.326001in}}%
\pgfpathlineto{\pgfqpoint{2.108085in}{1.325916in}}%
\pgfpathlineto{\pgfqpoint{2.109238in}{1.325555in}}%
\pgfpathlineto{\pgfqpoint{2.110063in}{1.325617in}}%
\pgfpathlineto{\pgfqpoint{2.111217in}{1.325099in}}%
\pgfpathlineto{\pgfqpoint{2.112041in}{1.325197in}}%
\pgfpathlineto{\pgfqpoint{2.112865in}{1.324460in}}%
\pgfpathlineto{\pgfqpoint{2.113360in}{1.324775in}}%
\pgfpathlineto{\pgfqpoint{2.113525in}{1.324027in}}%
\pgfpathlineto{\pgfqpoint{2.113689in}{1.324636in}}%
\pgfpathlineto{\pgfqpoint{2.114514in}{1.322448in}}%
\pgfpathlineto{\pgfqpoint{2.114678in}{1.324048in}}%
\pgfpathlineto{\pgfqpoint{2.115848in}{0.525689in}}%
\pgfpathlineto{\pgfqpoint{2.117002in}{0.525715in}}%
\pgfpathlineto{\pgfqpoint{2.123553in}{0.526043in}}%
\pgfpathlineto{\pgfqpoint{2.123641in}{0.526570in}}%
\pgfpathlineto{\pgfqpoint{2.123884in}{1.325970in}}%
\pgfpathlineto{\pgfqpoint{2.125113in}{1.325714in}}%
\pgfpathlineto{\pgfqpoint{2.126596in}{1.325573in}}%
\pgfpathlineto{\pgfqpoint{2.128080in}{1.325084in}}%
\pgfpathlineto{\pgfqpoint{2.128904in}{1.325033in}}%
\pgfpathlineto{\pgfqpoint{2.130058in}{1.324234in}}%
\pgfpathlineto{\pgfqpoint{2.130388in}{1.323987in}}%
\pgfpathlineto{\pgfqpoint{2.130553in}{1.324333in}}%
\pgfpathlineto{\pgfqpoint{2.131587in}{1.319528in}}%
\pgfpathlineto{\pgfqpoint{2.131853in}{1.321498in}}%
\pgfpathlineto{\pgfqpoint{2.132652in}{0.525693in}}%
\pgfpathlineto{\pgfqpoint{2.132982in}{0.525699in}}%
\pgfpathlineto{\pgfqpoint{2.140093in}{0.526241in}}%
\pgfpathlineto{\pgfqpoint{2.140182in}{0.555909in}}%
\pgfpathlineto{\pgfqpoint{2.140340in}{1.325998in}}%
\pgfpathlineto{\pgfqpoint{2.141066in}{1.325911in}}%
\pgfpathlineto{\pgfqpoint{2.142220in}{1.325564in}}%
\pgfpathlineto{\pgfqpoint{2.143044in}{1.325612in}}%
\pgfpathlineto{\pgfqpoint{2.144198in}{1.325114in}}%
\pgfpathlineto{\pgfqpoint{2.145023in}{1.325189in}}%
\pgfpathlineto{\pgfqpoint{2.145847in}{1.324488in}}%
\pgfpathlineto{\pgfqpoint{2.146341in}{1.324763in}}%
\pgfpathlineto{\pgfqpoint{2.146506in}{1.324069in}}%
\pgfpathlineto{\pgfqpoint{2.146671in}{1.324622in}}%
\pgfpathlineto{\pgfqpoint{2.147495in}{1.322631in}}%
\pgfpathlineto{\pgfqpoint{2.147660in}{1.324024in}}%
\pgfpathlineto{\pgfqpoint{2.147825in}{1.313812in}}%
\pgfpathlineto{\pgfqpoint{2.147948in}{1.323746in}}%
\pgfpathlineto{\pgfqpoint{2.149040in}{0.525691in}}%
\pgfpathlineto{\pgfqpoint{2.156568in}{0.526189in}}%
\pgfpathlineto{\pgfqpoint{2.156657in}{0.533923in}}%
\pgfpathlineto{\pgfqpoint{2.156813in}{1.326007in}}%
\pgfpathlineto{\pgfqpoint{2.157525in}{1.325923in}}%
\pgfpathlineto{\pgfqpoint{2.158679in}{1.325560in}}%
\pgfpathlineto{\pgfqpoint{2.159503in}{1.325627in}}%
\pgfpathlineto{\pgfqpoint{2.160657in}{1.325108in}}%
\pgfpathlineto{\pgfqpoint{2.161481in}{1.325212in}}%
\pgfpathlineto{\pgfqpoint{2.162306in}{1.324477in}}%
\pgfpathlineto{\pgfqpoint{2.162800in}{1.324797in}}%
\pgfpathlineto{\pgfqpoint{2.162965in}{1.324052in}}%
\pgfpathlineto{\pgfqpoint{2.163130in}{1.324661in}}%
\pgfpathlineto{\pgfqpoint{2.163954in}{1.322560in}}%
\pgfpathlineto{\pgfqpoint{2.164119in}{1.324092in}}%
\pgfpathlineto{\pgfqpoint{2.164556in}{1.040490in}}%
\pgfpathlineto{\pgfqpoint{2.164880in}{1.126165in}}%
\pgfpathlineto{\pgfqpoint{2.165346in}{0.525682in}}%
\pgfpathlineto{\pgfqpoint{2.165676in}{0.525687in}}%
\pgfpathlineto{\pgfqpoint{2.173064in}{0.526254in}}%
\pgfpathlineto{\pgfqpoint{2.173152in}{0.558215in}}%
\pgfpathlineto{\pgfqpoint{2.173311in}{1.326003in}}%
\pgfpathlineto{\pgfqpoint{2.174037in}{1.325915in}}%
\pgfpathlineto{\pgfqpoint{2.175191in}{1.325573in}}%
\pgfpathlineto{\pgfqpoint{2.176015in}{1.325618in}}%
\pgfpathlineto{\pgfqpoint{2.177169in}{1.325129in}}%
\pgfpathlineto{\pgfqpoint{2.177993in}{1.325198in}}%
\pgfpathlineto{\pgfqpoint{2.178818in}{1.324515in}}%
\pgfpathlineto{\pgfqpoint{2.179477in}{1.324109in}}%
\pgfpathlineto{\pgfqpoint{2.179642in}{1.324637in}}%
\pgfpathlineto{\pgfqpoint{2.180466in}{1.322782in}}%
\pgfpathlineto{\pgfqpoint{2.180631in}{1.324049in}}%
\pgfpathlineto{\pgfqpoint{2.180796in}{1.320645in}}%
\pgfpathlineto{\pgfqpoint{2.180919in}{1.323786in}}%
\pgfpathlineto{\pgfqpoint{2.181829in}{0.525682in}}%
\pgfpathlineto{\pgfqpoint{2.181994in}{0.525705in}}%
\pgfpathlineto{\pgfqpoint{2.187764in}{0.525853in}}%
\pgfpathlineto{\pgfqpoint{2.189549in}{0.526254in}}%
\pgfpathlineto{\pgfqpoint{2.189637in}{0.558214in}}%
\pgfpathlineto{\pgfqpoint{2.189796in}{1.326002in}}%
\pgfpathlineto{\pgfqpoint{2.190522in}{1.325915in}}%
\pgfpathlineto{\pgfqpoint{2.191676in}{1.325573in}}%
\pgfpathlineto{\pgfqpoint{2.192500in}{1.325617in}}%
\pgfpathlineto{\pgfqpoint{2.193654in}{1.325128in}}%
\pgfpathlineto{\pgfqpoint{2.194478in}{1.325197in}}%
\pgfpathlineto{\pgfqpoint{2.195303in}{1.324514in}}%
\pgfpathlineto{\pgfqpoint{2.195962in}{1.324107in}}%
\pgfpathlineto{\pgfqpoint{2.196127in}{1.324635in}}%
\pgfpathlineto{\pgfqpoint{2.196951in}{1.322774in}}%
\pgfpathlineto{\pgfqpoint{2.197116in}{1.324047in}}%
\pgfpathlineto{\pgfqpoint{2.197281in}{1.320573in}}%
\pgfpathlineto{\pgfqpoint{2.197403in}{1.323783in}}%
\pgfpathlineto{\pgfqpoint{2.198312in}{0.525682in}}%
\pgfpathlineto{\pgfqpoint{2.198477in}{0.525705in}}%
\pgfpathlineto{\pgfqpoint{2.204246in}{0.525853in}}%
\pgfpathlineto{\pgfqpoint{2.206034in}{0.526253in}}%
\pgfpathlineto{\pgfqpoint{2.206122in}{0.558206in}}%
\pgfpathlineto{\pgfqpoint{2.206280in}{1.326002in}}%
\pgfpathlineto{\pgfqpoint{2.207007in}{1.325915in}}%
\pgfpathlineto{\pgfqpoint{2.208161in}{1.325572in}}%
\pgfpathlineto{\pgfqpoint{2.208985in}{1.325616in}}%
\pgfpathlineto{\pgfqpoint{2.210139in}{1.325126in}}%
\pgfpathlineto{\pgfqpoint{2.210963in}{1.325196in}}%
\pgfpathlineto{\pgfqpoint{2.211787in}{1.324511in}}%
\pgfpathlineto{\pgfqpoint{2.212447in}{1.324102in}}%
\pgfpathlineto{\pgfqpoint{2.212612in}{1.324633in}}%
\pgfpathlineto{\pgfqpoint{2.213436in}{1.322755in}}%
\pgfpathlineto{\pgfqpoint{2.213601in}{1.324041in}}%
\pgfpathlineto{\pgfqpoint{2.213766in}{1.320406in}}%
\pgfpathlineto{\pgfqpoint{2.213888in}{1.323776in}}%
\pgfpathlineto{\pgfqpoint{2.214794in}{0.525682in}}%
\pgfpathlineto{\pgfqpoint{2.214958in}{0.525705in}}%
\pgfpathlineto{\pgfqpoint{2.220728in}{0.525853in}}%
\pgfpathlineto{\pgfqpoint{2.222519in}{0.526253in}}%
\pgfpathlineto{\pgfqpoint{2.222607in}{0.558195in}}%
\pgfpathlineto{\pgfqpoint{2.222765in}{1.326001in}}%
\pgfpathlineto{\pgfqpoint{2.223492in}{1.325914in}}%
\pgfpathlineto{\pgfqpoint{2.224646in}{1.325571in}}%
\pgfpathlineto{\pgfqpoint{2.225470in}{1.325615in}}%
\pgfpathlineto{\pgfqpoint{2.226624in}{1.325124in}}%
\pgfpathlineto{\pgfqpoint{2.227448in}{1.325193in}}%
\pgfpathlineto{\pgfqpoint{2.228272in}{1.324507in}}%
\pgfpathlineto{\pgfqpoint{2.228932in}{1.324096in}}%
\pgfpathlineto{\pgfqpoint{2.229096in}{1.324629in}}%
\pgfpathlineto{\pgfqpoint{2.229921in}{1.322732in}}%
\pgfpathlineto{\pgfqpoint{2.230086in}{1.324035in}}%
\pgfpathlineto{\pgfqpoint{2.230250in}{1.320152in}}%
\pgfpathlineto{\pgfqpoint{2.230373in}{1.323767in}}%
\pgfpathlineto{\pgfqpoint{2.231278in}{0.525681in}}%
\pgfpathlineto{\pgfqpoint{2.231443in}{0.525704in}}%
\pgfpathlineto{\pgfqpoint{2.237213in}{0.525853in}}%
\pgfpathlineto{\pgfqpoint{2.239004in}{0.526252in}}%
\pgfpathlineto{\pgfqpoint{2.239092in}{0.558183in}}%
\pgfpathlineto{\pgfqpoint{2.239250in}{1.326000in}}%
\pgfpathlineto{\pgfqpoint{2.239976in}{1.325913in}}%
\pgfpathlineto{\pgfqpoint{2.241130in}{1.325569in}}%
\pgfpathlineto{\pgfqpoint{2.241955in}{1.325614in}}%
\pgfpathlineto{\pgfqpoint{2.243109in}{1.325122in}}%
\pgfpathlineto{\pgfqpoint{2.243933in}{1.325192in}}%
\pgfpathlineto{\pgfqpoint{2.244757in}{1.324503in}}%
\pgfpathlineto{\pgfqpoint{2.245252in}{1.324766in}}%
\pgfpathlineto{\pgfqpoint{2.245416in}{1.324090in}}%
\pgfpathlineto{\pgfqpoint{2.245581in}{1.324625in}}%
\pgfpathlineto{\pgfqpoint{2.246406in}{1.322709in}}%
\pgfpathlineto{\pgfqpoint{2.246570in}{1.324028in}}%
\pgfpathlineto{\pgfqpoint{2.246735in}{1.319835in}}%
\pgfpathlineto{\pgfqpoint{2.246858in}{1.323758in}}%
\pgfpathlineto{\pgfqpoint{2.247764in}{0.525681in}}%
\pgfpathlineto{\pgfqpoint{2.247928in}{0.525704in}}%
\pgfpathlineto{\pgfqpoint{2.253698in}{0.525852in}}%
\pgfpathlineto{\pgfqpoint{2.255489in}{0.526252in}}%
\pgfpathlineto{\pgfqpoint{2.255577in}{0.558172in}}%
\pgfpathlineto{\pgfqpoint{2.255735in}{1.325999in}}%
\pgfpathlineto{\pgfqpoint{2.256461in}{1.325912in}}%
\pgfpathlineto{\pgfqpoint{2.257615in}{1.325568in}}%
\pgfpathlineto{\pgfqpoint{2.258440in}{1.325613in}}%
\pgfpathlineto{\pgfqpoint{2.259593in}{1.325120in}}%
\pgfpathlineto{\pgfqpoint{2.260418in}{1.325190in}}%
\pgfpathlineto{\pgfqpoint{2.261242in}{1.324499in}}%
\pgfpathlineto{\pgfqpoint{2.261736in}{1.324764in}}%
\pgfpathlineto{\pgfqpoint{2.261901in}{1.324085in}}%
\pgfpathlineto{\pgfqpoint{2.262066in}{1.324622in}}%
\pgfpathlineto{\pgfqpoint{2.262890in}{1.322690in}}%
\pgfpathlineto{\pgfqpoint{2.263055in}{1.324023in}}%
\pgfpathlineto{\pgfqpoint{2.263220in}{1.319474in}}%
\pgfpathlineto{\pgfqpoint{2.263343in}{1.323751in}}%
\pgfpathlineto{\pgfqpoint{2.264280in}{0.525680in}}%
\pgfpathlineto{\pgfqpoint{2.264444in}{0.525705in}}%
\pgfpathlineto{\pgfqpoint{2.269884in}{0.525842in}}%
\pgfpathlineto{\pgfqpoint{2.271980in}{0.526295in}}%
\pgfpathlineto{\pgfqpoint{2.272147in}{0.889249in}}%
\pgfpathlineto{\pgfqpoint{2.272223in}{1.325992in}}%
\pgfpathlineto{\pgfqpoint{2.272949in}{1.325904in}}%
\pgfpathlineto{\pgfqpoint{2.274102in}{1.325577in}}%
\pgfpathlineto{\pgfqpoint{2.274927in}{1.325602in}}%
\pgfpathlineto{\pgfqpoint{2.276081in}{1.325135in}}%
\pgfpathlineto{\pgfqpoint{2.276905in}{1.325174in}}%
\pgfpathlineto{\pgfqpoint{2.278059in}{1.324345in}}%
\pgfpathlineto{\pgfqpoint{2.278224in}{1.324739in}}%
\pgfpathlineto{\pgfqpoint{2.278389in}{1.324125in}}%
\pgfpathlineto{\pgfqpoint{2.278553in}{1.324595in}}%
\pgfpathlineto{\pgfqpoint{2.279378in}{1.322835in}}%
\pgfpathlineto{\pgfqpoint{2.279542in}{1.323973in}}%
\pgfpathlineto{\pgfqpoint{2.279707in}{1.320994in}}%
\pgfpathlineto{\pgfqpoint{2.279829in}{1.323688in}}%
\pgfpathlineto{\pgfqpoint{2.280725in}{0.525682in}}%
\pgfpathlineto{\pgfqpoint{2.281054in}{0.525687in}}%
\pgfpathlineto{\pgfqpoint{2.288463in}{0.526283in}}%
\pgfpathlineto{\pgfqpoint{2.288648in}{0.944748in}}%
\pgfpathlineto{\pgfqpoint{2.288725in}{1.325991in}}%
\pgfpathlineto{\pgfqpoint{2.289474in}{1.325902in}}%
\pgfpathlineto{\pgfqpoint{2.290628in}{1.325563in}}%
\pgfpathlineto{\pgfqpoint{2.291452in}{1.325600in}}%
\pgfpathlineto{\pgfqpoint{2.292606in}{1.325113in}}%
\pgfpathlineto{\pgfqpoint{2.293430in}{1.325171in}}%
\pgfpathlineto{\pgfqpoint{2.294255in}{1.324486in}}%
\pgfpathlineto{\pgfqpoint{2.294584in}{1.324298in}}%
\pgfpathlineto{\pgfqpoint{2.294749in}{1.324735in}}%
\pgfpathlineto{\pgfqpoint{2.295573in}{1.323345in}}%
\pgfpathlineto{\pgfqpoint{2.295738in}{1.324217in}}%
\pgfpathlineto{\pgfqpoint{2.295903in}{1.322614in}}%
\pgfpathlineto{\pgfqpoint{2.296068in}{1.323963in}}%
\pgfpathlineto{\pgfqpoint{2.296233in}{1.306249in}}%
\pgfpathlineto{\pgfqpoint{2.296334in}{1.323681in}}%
\pgfpathlineto{\pgfqpoint{2.297393in}{0.525688in}}%
\pgfpathlineto{\pgfqpoint{2.304986in}{0.526798in}}%
\pgfpathlineto{\pgfqpoint{2.305224in}{1.325960in}}%
\pgfpathlineto{\pgfqpoint{2.306286in}{1.325825in}}%
\pgfpathlineto{\pgfqpoint{2.308099in}{1.325419in}}%
\pgfpathlineto{\pgfqpoint{2.309253in}{1.325280in}}%
\pgfpathlineto{\pgfqpoint{2.310737in}{1.324613in}}%
\pgfpathlineto{\pgfqpoint{2.311561in}{1.324474in}}%
\pgfpathlineto{\pgfqpoint{2.312715in}{1.322352in}}%
\pgfpathlineto{\pgfqpoint{2.312817in}{1.323401in}}%
\pgfpathlineto{\pgfqpoint{2.312920in}{1.320586in}}%
\pgfpathlineto{\pgfqpoint{2.313187in}{1.321236in}}%
\pgfpathlineto{\pgfqpoint{2.313741in}{0.525690in}}%
\pgfpathlineto{\pgfqpoint{2.314400in}{0.525700in}}%
\pgfpathlineto{\pgfqpoint{2.321427in}{0.526239in}}%
\pgfpathlineto{\pgfqpoint{2.321515in}{0.555855in}}%
\pgfpathlineto{\pgfqpoint{2.321673in}{1.325995in}}%
\pgfpathlineto{\pgfqpoint{2.322399in}{1.325907in}}%
\pgfpathlineto{\pgfqpoint{2.323553in}{1.325558in}}%
\pgfpathlineto{\pgfqpoint{2.324378in}{1.325606in}}%
\pgfpathlineto{\pgfqpoint{2.325532in}{1.325104in}}%
\pgfpathlineto{\pgfqpoint{2.326356in}{1.325180in}}%
\pgfpathlineto{\pgfqpoint{2.327180in}{1.324470in}}%
\pgfpathlineto{\pgfqpoint{2.327510in}{1.324277in}}%
\pgfpathlineto{\pgfqpoint{2.327675in}{1.324749in}}%
\pgfpathlineto{\pgfqpoint{2.328499in}{1.323293in}}%
\pgfpathlineto{\pgfqpoint{2.328664in}{1.324240in}}%
\pgfpathlineto{\pgfqpoint{2.328829in}{1.322508in}}%
\pgfpathlineto{\pgfqpoint{2.328993in}{1.323993in}}%
\pgfpathlineto{\pgfqpoint{2.329757in}{0.789443in}}%
\pgfpathlineto{\pgfqpoint{2.330119in}{0.525681in}}%
\pgfpathlineto{\pgfqpoint{2.330449in}{0.525685in}}%
\pgfpathlineto{\pgfqpoint{2.337925in}{0.526336in}}%
\pgfpathlineto{\pgfqpoint{2.338106in}{0.954377in}}%
\pgfpathlineto{\pgfqpoint{2.338182in}{1.325985in}}%
\pgfpathlineto{\pgfqpoint{2.338928in}{1.325896in}}%
\pgfpathlineto{\pgfqpoint{2.340082in}{1.325570in}}%
\pgfpathlineto{\pgfqpoint{2.340906in}{1.325592in}}%
\pgfpathlineto{\pgfqpoint{2.342060in}{1.325124in}}%
\pgfpathlineto{\pgfqpoint{2.342884in}{1.325159in}}%
\pgfpathlineto{\pgfqpoint{2.344038in}{1.324321in}}%
\pgfpathlineto{\pgfqpoint{2.344203in}{1.324717in}}%
\pgfpathlineto{\pgfqpoint{2.344368in}{1.324096in}}%
\pgfpathlineto{\pgfqpoint{2.344533in}{1.324568in}}%
\pgfpathlineto{\pgfqpoint{2.345357in}{1.322732in}}%
\pgfpathlineto{\pgfqpoint{2.345522in}{1.323925in}}%
\pgfpathlineto{\pgfqpoint{2.345686in}{1.320151in}}%
\pgfpathlineto{\pgfqpoint{2.345788in}{1.323647in}}%
\pgfpathlineto{\pgfqpoint{2.346621in}{0.525686in}}%
\pgfpathlineto{\pgfqpoint{2.346951in}{0.525689in}}%
\pgfpathlineto{\pgfqpoint{2.354421in}{0.526454in}}%
\pgfpathlineto{\pgfqpoint{2.354667in}{1.325972in}}%
\pgfpathlineto{\pgfqpoint{2.355896in}{1.325704in}}%
\pgfpathlineto{\pgfqpoint{2.357050in}{1.325634in}}%
\pgfpathlineto{\pgfqpoint{2.358533in}{1.325160in}}%
\pgfpathlineto{\pgfqpoint{2.359357in}{1.325133in}}%
\pgfpathlineto{\pgfqpoint{2.360511in}{1.324398in}}%
\pgfpathlineto{\pgfqpoint{2.360841in}{1.324191in}}%
\pgfpathlineto{\pgfqpoint{2.361006in}{1.324523in}}%
\pgfpathlineto{\pgfqpoint{2.362160in}{1.321897in}}%
\pgfpathlineto{\pgfqpoint{2.362267in}{1.323528in}}%
\pgfpathlineto{\pgfqpoint{2.362796in}{0.619762in}}%
\pgfpathlineto{\pgfqpoint{2.363071in}{0.525687in}}%
\pgfpathlineto{\pgfqpoint{2.363560in}{0.525698in}}%
\pgfpathlineto{\pgfqpoint{2.370836in}{0.526066in}}%
\pgfpathlineto{\pgfqpoint{2.370927in}{0.526839in}}%
\pgfpathlineto{\pgfqpoint{2.371165in}{1.325959in}}%
\pgfpathlineto{\pgfqpoint{2.372065in}{1.325765in}}%
\pgfpathlineto{\pgfqpoint{2.373548in}{1.325616in}}%
\pgfpathlineto{\pgfqpoint{2.375032in}{1.325178in}}%
\pgfpathlineto{\pgfqpoint{2.375856in}{1.325105in}}%
\pgfpathlineto{\pgfqpoint{2.377010in}{1.324436in}}%
\pgfpathlineto{\pgfqpoint{2.377505in}{1.324471in}}%
\pgfpathlineto{\pgfqpoint{2.378658in}{1.322258in}}%
\pgfpathlineto{\pgfqpoint{2.378759in}{1.323396in}}%
\pgfpathlineto{\pgfqpoint{2.378859in}{1.320106in}}%
\pgfpathlineto{\pgfqpoint{2.379126in}{1.321123in}}%
\pgfpathlineto{\pgfqpoint{2.379647in}{0.525689in}}%
\pgfpathlineto{\pgfqpoint{2.380307in}{0.525699in}}%
\pgfpathlineto{\pgfqpoint{2.387366in}{0.526239in}}%
\pgfpathlineto{\pgfqpoint{2.387454in}{0.555856in}}%
\pgfpathlineto{\pgfqpoint{2.387613in}{1.325995in}}%
\pgfpathlineto{\pgfqpoint{2.388339in}{1.325907in}}%
\pgfpathlineto{\pgfqpoint{2.389493in}{1.325558in}}%
\pgfpathlineto{\pgfqpoint{2.390317in}{1.325607in}}%
\pgfpathlineto{\pgfqpoint{2.391471in}{1.325105in}}%
\pgfpathlineto{\pgfqpoint{2.392295in}{1.325181in}}%
\pgfpathlineto{\pgfqpoint{2.393119in}{1.324473in}}%
\pgfpathlineto{\pgfqpoint{2.393449in}{1.324281in}}%
\pgfpathlineto{\pgfqpoint{2.393614in}{1.324752in}}%
\pgfpathlineto{\pgfqpoint{2.394438in}{1.323305in}}%
\pgfpathlineto{\pgfqpoint{2.394603in}{1.324245in}}%
\pgfpathlineto{\pgfqpoint{2.394768in}{1.322535in}}%
\pgfpathlineto{\pgfqpoint{2.394933in}{1.324000in}}%
\pgfpathlineto{\pgfqpoint{2.395683in}{0.917386in}}%
\pgfpathlineto{\pgfqpoint{2.396109in}{0.525686in}}%
\pgfpathlineto{\pgfqpoint{2.396439in}{0.525689in}}%
\pgfpathlineto{\pgfqpoint{2.403867in}{0.526363in}}%
\pgfpathlineto{\pgfqpoint{2.404120in}{1.325981in}}%
\pgfpathlineto{\pgfqpoint{2.405520in}{1.325801in}}%
\pgfpathlineto{\pgfqpoint{2.407003in}{1.325385in}}%
\pgfpathlineto{\pgfqpoint{2.407828in}{1.325393in}}%
\pgfpathlineto{\pgfqpoint{2.408981in}{1.324826in}}%
\pgfpathlineto{\pgfqpoint{2.409806in}{1.324838in}}%
\pgfpathlineto{\pgfqpoint{2.410630in}{1.323877in}}%
\pgfpathlineto{\pgfqpoint{2.410795in}{1.324383in}}%
\pgfpathlineto{\pgfqpoint{2.410960in}{1.323506in}}%
\pgfpathlineto{\pgfqpoint{2.411125in}{1.324172in}}%
\pgfpathlineto{\pgfqpoint{2.411619in}{1.321433in}}%
\pgfpathlineto{\pgfqpoint{2.411724in}{1.323621in}}%
\pgfpathlineto{\pgfqpoint{2.412558in}{0.525684in}}%
\pgfpathlineto{\pgfqpoint{2.413218in}{0.525694in}}%
\pgfpathlineto{\pgfqpoint{2.420348in}{0.526329in}}%
\pgfpathlineto{\pgfqpoint{2.420514in}{0.896053in}}%
\pgfpathlineto{\pgfqpoint{2.420589in}{1.325986in}}%
\pgfpathlineto{\pgfqpoint{2.421316in}{1.325896in}}%
\pgfpathlineto{\pgfqpoint{2.422800in}{1.325520in}}%
\pgfpathlineto{\pgfqpoint{2.423624in}{1.325533in}}%
\pgfpathlineto{\pgfqpoint{2.424778in}{1.325047in}}%
\pgfpathlineto{\pgfqpoint{2.425602in}{1.325067in}}%
\pgfpathlineto{\pgfqpoint{2.426756in}{1.324147in}}%
\pgfpathlineto{\pgfqpoint{2.426921in}{1.324573in}}%
\pgfpathlineto{\pgfqpoint{2.427086in}{1.323875in}}%
\pgfpathlineto{\pgfqpoint{2.427251in}{1.324401in}}%
\pgfpathlineto{\pgfqpoint{2.428075in}{1.321410in}}%
\pgfpathlineto{\pgfqpoint{2.428194in}{1.323640in}}%
\pgfpathlineto{\pgfqpoint{2.429057in}{0.525688in}}%
\pgfpathlineto{\pgfqpoint{2.429551in}{0.525700in}}%
\pgfpathlineto{\pgfqpoint{2.436768in}{0.526049in}}%
\pgfpathlineto{\pgfqpoint{2.436857in}{0.526624in}}%
\pgfpathlineto{\pgfqpoint{2.437098in}{1.325970in}}%
\pgfpathlineto{\pgfqpoint{2.438163in}{1.325832in}}%
\pgfpathlineto{\pgfqpoint{2.439976in}{1.325413in}}%
\pgfpathlineto{\pgfqpoint{2.441130in}{1.325293in}}%
\pgfpathlineto{\pgfqpoint{2.442284in}{1.324746in}}%
\pgfpathlineto{\pgfqpoint{2.443108in}{1.324660in}}%
\pgfpathlineto{\pgfqpoint{2.444262in}{1.323184in}}%
\pgfpathlineto{\pgfqpoint{2.444427in}{1.323800in}}%
\pgfpathlineto{\pgfqpoint{2.444592in}{1.322259in}}%
\pgfpathlineto{\pgfqpoint{2.444695in}{1.323481in}}%
\pgfpathlineto{\pgfqpoint{2.444799in}{1.319888in}}%
\pgfpathlineto{\pgfqpoint{2.445065in}{1.321423in}}%
\pgfpathlineto{\pgfqpoint{2.445779in}{0.525691in}}%
\pgfpathlineto{\pgfqpoint{2.446273in}{0.525702in}}%
\pgfpathlineto{\pgfqpoint{2.453305in}{0.526241in}}%
\pgfpathlineto{\pgfqpoint{2.453394in}{0.555904in}}%
\pgfpathlineto{\pgfqpoint{2.453552in}{1.325998in}}%
\pgfpathlineto{\pgfqpoint{2.454278in}{1.325911in}}%
\pgfpathlineto{\pgfqpoint{2.455432in}{1.325563in}}%
\pgfpathlineto{\pgfqpoint{2.456256in}{1.325612in}}%
\pgfpathlineto{\pgfqpoint{2.457410in}{1.325114in}}%
\pgfpathlineto{\pgfqpoint{2.458235in}{1.325189in}}%
\pgfpathlineto{\pgfqpoint{2.459059in}{1.324489in}}%
\pgfpathlineto{\pgfqpoint{2.459553in}{1.324764in}}%
\pgfpathlineto{\pgfqpoint{2.459718in}{1.324070in}}%
\pgfpathlineto{\pgfqpoint{2.459883in}{1.324623in}}%
\pgfpathlineto{\pgfqpoint{2.460707in}{1.322635in}}%
\pgfpathlineto{\pgfqpoint{2.460872in}{1.324025in}}%
\pgfpathlineto{\pgfqpoint{2.461037in}{1.315579in}}%
\pgfpathlineto{\pgfqpoint{2.461160in}{1.323750in}}%
\pgfpathlineto{\pgfqpoint{2.462288in}{0.525690in}}%
\pgfpathlineto{\pgfqpoint{2.469748in}{0.526075in}}%
\pgfpathlineto{\pgfqpoint{2.469838in}{0.526919in}}%
\pgfpathlineto{\pgfqpoint{2.470077in}{1.325968in}}%
\pgfpathlineto{\pgfqpoint{2.470981in}{1.325767in}}%
\pgfpathlineto{\pgfqpoint{2.472464in}{1.325627in}}%
\pgfpathlineto{\pgfqpoint{2.473948in}{1.325181in}}%
\pgfpathlineto{\pgfqpoint{2.474772in}{1.325122in}}%
\pgfpathlineto{\pgfqpoint{2.475926in}{1.324444in}}%
\pgfpathlineto{\pgfqpoint{2.476421in}{1.324504in}}%
\pgfpathlineto{\pgfqpoint{2.477575in}{1.322327in}}%
\pgfpathlineto{\pgfqpoint{2.477672in}{1.323493in}}%
\pgfpathlineto{\pgfqpoint{2.477768in}{1.320713in}}%
\pgfpathlineto{\pgfqpoint{2.478035in}{1.321622in}}%
\pgfpathlineto{\pgfqpoint{2.478540in}{0.525686in}}%
\pgfpathlineto{\pgfqpoint{2.479200in}{0.525696in}}%
\pgfpathlineto{\pgfqpoint{2.486285in}{0.526313in}}%
\pgfpathlineto{\pgfqpoint{2.486451in}{0.890165in}}%
\pgfpathlineto{\pgfqpoint{2.486526in}{1.325990in}}%
\pgfpathlineto{\pgfqpoint{2.487251in}{1.325901in}}%
\pgfpathlineto{\pgfqpoint{2.488735in}{1.325522in}}%
\pgfpathlineto{\pgfqpoint{2.489559in}{1.325539in}}%
\pgfpathlineto{\pgfqpoint{2.490713in}{1.325049in}}%
\pgfpathlineto{\pgfqpoint{2.491537in}{1.325078in}}%
\pgfpathlineto{\pgfqpoint{2.492361in}{1.324366in}}%
\pgfpathlineto{\pgfqpoint{2.492691in}{1.324151in}}%
\pgfpathlineto{\pgfqpoint{2.492856in}{1.324590in}}%
\pgfpathlineto{\pgfqpoint{2.493680in}{1.322927in}}%
\pgfpathlineto{\pgfqpoint{2.493845in}{1.323965in}}%
\pgfpathlineto{\pgfqpoint{2.494010in}{1.321457in}}%
\pgfpathlineto{\pgfqpoint{2.494132in}{1.323678in}}%
\pgfpathlineto{\pgfqpoint{2.495220in}{0.525690in}}%
\pgfpathlineto{\pgfqpoint{2.495384in}{0.525697in}}%
\pgfpathlineto{\pgfqpoint{2.502736in}{0.526130in}}%
\pgfpathlineto{\pgfqpoint{2.502825in}{0.527973in}}%
\pgfpathlineto{\pgfqpoint{2.502982in}{1.325974in}}%
\pgfpathlineto{\pgfqpoint{2.503686in}{1.325889in}}%
\pgfpathlineto{\pgfqpoint{2.505169in}{1.325541in}}%
\pgfpathlineto{\pgfqpoint{2.506323in}{1.325459in}}%
\pgfpathlineto{\pgfqpoint{2.507477in}{1.324979in}}%
\pgfpathlineto{\pgfqpoint{2.508302in}{1.324948in}}%
\pgfpathlineto{\pgfqpoint{2.509455in}{1.323975in}}%
\pgfpathlineto{\pgfqpoint{2.509620in}{1.324373in}}%
\pgfpathlineto{\pgfqpoint{2.509785in}{1.323645in}}%
\pgfpathlineto{\pgfqpoint{2.509950in}{1.324160in}}%
\pgfpathlineto{\pgfqpoint{2.510445in}{1.322184in}}%
\pgfpathlineto{\pgfqpoint{2.510591in}{1.323548in}}%
\pgfpathlineto{\pgfqpoint{2.511673in}{0.525687in}}%
\pgfpathlineto{\pgfqpoint{2.512332in}{0.525701in}}%
\pgfpathlineto{\pgfqpoint{2.519210in}{0.526092in}}%
\pgfpathlineto{\pgfqpoint{2.519299in}{0.527174in}}%
\pgfpathlineto{\pgfqpoint{2.519457in}{1.325926in}}%
\pgfpathlineto{\pgfqpoint{2.520305in}{1.325809in}}%
\pgfpathlineto{\pgfqpoint{2.525251in}{1.324591in}}%
\pgfpathlineto{\pgfqpoint{2.526734in}{1.323491in}}%
\pgfpathlineto{\pgfqpoint{2.527465in}{1.320281in}}%
\pgfpathlineto{\pgfqpoint{2.527490in}{1.318808in}}%
\pgfpathlineto{\pgfqpoint{2.527750in}{0.525664in}}%
\pgfpathlineto{\pgfqpoint{2.528654in}{0.525689in}}%
\pgfpathlineto{\pgfqpoint{2.535731in}{0.526243in}}%
\pgfpathlineto{\pgfqpoint{2.535819in}{0.557990in}}%
\pgfpathlineto{\pgfqpoint{2.535977in}{1.325987in}}%
\pgfpathlineto{\pgfqpoint{2.536704in}{1.325898in}}%
\pgfpathlineto{\pgfqpoint{2.537858in}{1.325550in}}%
\pgfpathlineto{\pgfqpoint{2.538682in}{1.325595in}}%
\pgfpathlineto{\pgfqpoint{2.539836in}{1.325093in}}%
\pgfpathlineto{\pgfqpoint{2.540660in}{1.325165in}}%
\pgfpathlineto{\pgfqpoint{2.541484in}{1.324451in}}%
\pgfpathlineto{\pgfqpoint{2.541814in}{1.324256in}}%
\pgfpathlineto{\pgfqpoint{2.541979in}{1.324727in}}%
\pgfpathlineto{\pgfqpoint{2.542803in}{1.323239in}}%
\pgfpathlineto{\pgfqpoint{2.542968in}{1.324206in}}%
\pgfpathlineto{\pgfqpoint{2.543133in}{1.322393in}}%
\pgfpathlineto{\pgfqpoint{2.543298in}{1.323950in}}%
\pgfpathlineto{\pgfqpoint{2.544458in}{0.525680in}}%
\pgfpathlineto{\pgfqpoint{2.545282in}{0.525708in}}%
\pgfpathlineto{\pgfqpoint{2.552129in}{0.525986in}}%
\pgfpathlineto{\pgfqpoint{2.552226in}{0.526320in}}%
\pgfpathlineto{\pgfqpoint{2.552408in}{0.951702in}}%
\pgfpathlineto{\pgfqpoint{2.552484in}{1.325984in}}%
\pgfpathlineto{\pgfqpoint{2.553230in}{1.325896in}}%
\pgfpathlineto{\pgfqpoint{2.554384in}{1.325567in}}%
\pgfpathlineto{\pgfqpoint{2.555208in}{1.325593in}}%
\pgfpathlineto{\pgfqpoint{2.556362in}{1.325121in}}%
\pgfpathlineto{\pgfqpoint{2.557186in}{1.325161in}}%
\pgfpathlineto{\pgfqpoint{2.558010in}{1.324504in}}%
\pgfpathlineto{\pgfqpoint{2.558340in}{1.324318in}}%
\pgfpathlineto{\pgfqpoint{2.558505in}{1.324722in}}%
\pgfpathlineto{\pgfqpoint{2.559329in}{1.323400in}}%
\pgfpathlineto{\pgfqpoint{2.559494in}{1.324197in}}%
\pgfpathlineto{\pgfqpoint{2.559659in}{1.322727in}}%
\pgfpathlineto{\pgfqpoint{2.559824in}{1.323938in}}%
\pgfpathlineto{\pgfqpoint{2.559988in}{1.320127in}}%
\pgfpathlineto{\pgfqpoint{2.560091in}{1.323664in}}%
\pgfpathlineto{\pgfqpoint{2.560929in}{0.525683in}}%
\pgfpathlineto{\pgfqpoint{2.561259in}{0.525687in}}%
\pgfpathlineto{\pgfqpoint{2.568715in}{0.526360in}}%
\pgfpathlineto{\pgfqpoint{2.568972in}{1.325987in}}%
\pgfpathlineto{\pgfqpoint{2.570377in}{1.325808in}}%
\pgfpathlineto{\pgfqpoint{2.571531in}{1.325451in}}%
\pgfpathlineto{\pgfqpoint{2.572355in}{1.325472in}}%
\pgfpathlineto{\pgfqpoint{2.573509in}{1.324935in}}%
\pgfpathlineto{\pgfqpoint{2.574333in}{1.324971in}}%
\pgfpathlineto{\pgfqpoint{2.575158in}{1.324134in}}%
\pgfpathlineto{\pgfqpoint{2.575322in}{1.324585in}}%
\pgfpathlineto{\pgfqpoint{2.575487in}{1.323857in}}%
\pgfpathlineto{\pgfqpoint{2.575652in}{1.324415in}}%
\pgfpathlineto{\pgfqpoint{2.576476in}{1.321220in}}%
\pgfpathlineto{\pgfqpoint{2.576577in}{1.323691in}}%
\pgfpathlineto{\pgfqpoint{2.577448in}{0.525683in}}%
\pgfpathlineto{\pgfqpoint{2.577942in}{0.525705in}}%
\pgfpathlineto{\pgfqpoint{2.585095in}{0.525984in}}%
\pgfpathlineto{\pgfqpoint{2.585189in}{0.526272in}}%
\pgfpathlineto{\pgfqpoint{2.585360in}{0.897351in}}%
\pgfpathlineto{\pgfqpoint{2.585436in}{1.326001in}}%
\pgfpathlineto{\pgfqpoint{2.586166in}{1.325913in}}%
\pgfpathlineto{\pgfqpoint{2.587320in}{1.325575in}}%
\pgfpathlineto{\pgfqpoint{2.588144in}{1.325615in}}%
\pgfpathlineto{\pgfqpoint{2.589298in}{1.325132in}}%
\pgfpathlineto{\pgfqpoint{2.590123in}{1.325194in}}%
\pgfpathlineto{\pgfqpoint{2.590947in}{1.324523in}}%
\pgfpathlineto{\pgfqpoint{2.591606in}{1.324122in}}%
\pgfpathlineto{\pgfqpoint{2.591771in}{1.324632in}}%
\pgfpathlineto{\pgfqpoint{2.592595in}{1.322831in}}%
\pgfpathlineto{\pgfqpoint{2.592760in}{1.324043in}}%
\pgfpathlineto{\pgfqpoint{2.592925in}{1.320985in}}%
\pgfpathlineto{\pgfqpoint{2.593044in}{1.323783in}}%
\pgfpathlineto{\pgfqpoint{2.593974in}{0.525682in}}%
\pgfpathlineto{\pgfqpoint{2.594303in}{0.525687in}}%
\pgfpathlineto{\pgfqpoint{2.601677in}{0.526299in}}%
\pgfpathlineto{\pgfqpoint{2.601846in}{0.895481in}}%
\pgfpathlineto{\pgfqpoint{2.601921in}{1.325999in}}%
\pgfpathlineto{\pgfqpoint{2.602650in}{1.325911in}}%
\pgfpathlineto{\pgfqpoint{2.603803in}{1.325585in}}%
\pgfpathlineto{\pgfqpoint{2.604628in}{1.325613in}}%
\pgfpathlineto{\pgfqpoint{2.605782in}{1.325148in}}%
\pgfpathlineto{\pgfqpoint{2.606606in}{1.325191in}}%
\pgfpathlineto{\pgfqpoint{2.607430in}{1.324552in}}%
\pgfpathlineto{\pgfqpoint{2.608090in}{1.324164in}}%
\pgfpathlineto{\pgfqpoint{2.608254in}{1.324626in}}%
\pgfpathlineto{\pgfqpoint{2.609079in}{1.322970in}}%
\pgfpathlineto{\pgfqpoint{2.609243in}{1.324031in}}%
\pgfpathlineto{\pgfqpoint{2.609408in}{1.321626in}}%
\pgfpathlineto{\pgfqpoint{2.609528in}{1.323767in}}%
\pgfpathlineto{\pgfqpoint{2.610332in}{0.525688in}}%
\pgfpathlineto{\pgfqpoint{2.610973in}{0.525695in}}%
\pgfpathlineto{\pgfqpoint{2.618164in}{0.526316in}}%
\pgfpathlineto{\pgfqpoint{2.618329in}{0.889404in}}%
\pgfpathlineto{\pgfqpoint{2.618405in}{1.325998in}}%
\pgfpathlineto{\pgfqpoint{2.619129in}{1.325910in}}%
\pgfpathlineto{\pgfqpoint{2.620613in}{1.325533in}}%
\pgfpathlineto{\pgfqpoint{2.621437in}{1.325552in}}%
\pgfpathlineto{\pgfqpoint{2.622591in}{1.325067in}}%
\pgfpathlineto{\pgfqpoint{2.623416in}{1.325098in}}%
\pgfpathlineto{\pgfqpoint{2.624240in}{1.324402in}}%
\pgfpathlineto{\pgfqpoint{2.624569in}{1.324195in}}%
\pgfpathlineto{\pgfqpoint{2.624734in}{1.324622in}}%
\pgfpathlineto{\pgfqpoint{2.625559in}{1.323066in}}%
\pgfpathlineto{\pgfqpoint{2.625723in}{1.324023in}}%
\pgfpathlineto{\pgfqpoint{2.625888in}{1.321942in}}%
\pgfpathlineto{\pgfqpoint{2.626010in}{1.323754in}}%
\pgfpathlineto{\pgfqpoint{2.627180in}{0.525692in}}%
\pgfpathlineto{\pgfqpoint{2.627839in}{0.525706in}}%
\pgfpathlineto{\pgfqpoint{2.634666in}{0.526496in}}%
\pgfpathlineto{\pgfqpoint{2.634911in}{1.325979in}}%
\pgfpathlineto{\pgfqpoint{2.636140in}{1.325717in}}%
\pgfpathlineto{\pgfqpoint{2.637294in}{1.325644in}}%
\pgfpathlineto{\pgfqpoint{2.638777in}{1.325182in}}%
\pgfpathlineto{\pgfqpoint{2.639601in}{1.325149in}}%
\pgfpathlineto{\pgfqpoint{2.640755in}{1.324445in}}%
\pgfpathlineto{\pgfqpoint{2.641250in}{1.324552in}}%
\pgfpathlineto{\pgfqpoint{2.642404in}{1.322333in}}%
\pgfpathlineto{\pgfqpoint{2.642510in}{1.323604in}}%
\pgfpathlineto{\pgfqpoint{2.642617in}{1.320265in}}%
\pgfpathlineto{\pgfqpoint{2.642884in}{1.321841in}}%
\pgfpathlineto{\pgfqpoint{2.643373in}{0.525687in}}%
\pgfpathlineto{\pgfqpoint{2.644033in}{0.525696in}}%
\pgfpathlineto{\pgfqpoint{2.651131in}{0.526297in}}%
\pgfpathlineto{\pgfqpoint{2.651299in}{0.890363in}}%
\pgfpathlineto{\pgfqpoint{2.651374in}{1.325995in}}%
\pgfpathlineto{\pgfqpoint{2.652101in}{1.325907in}}%
\pgfpathlineto{\pgfqpoint{2.653255in}{1.325582in}}%
\pgfpathlineto{\pgfqpoint{2.654079in}{1.325607in}}%
\pgfpathlineto{\pgfqpoint{2.655233in}{1.325142in}}%
\pgfpathlineto{\pgfqpoint{2.656057in}{1.325181in}}%
\pgfpathlineto{\pgfqpoint{2.657211in}{1.324361in}}%
\pgfpathlineto{\pgfqpoint{2.657376in}{1.324751in}}%
\pgfpathlineto{\pgfqpoint{2.657541in}{1.324145in}}%
\pgfpathlineto{\pgfqpoint{2.657706in}{1.324608in}}%
\pgfpathlineto{\pgfqpoint{2.658530in}{1.322904in}}%
\pgfpathlineto{\pgfqpoint{2.658695in}{1.323998in}}%
\pgfpathlineto{\pgfqpoint{2.658860in}{1.321355in}}%
\pgfpathlineto{\pgfqpoint{2.658981in}{1.323722in}}%
\pgfpathlineto{\pgfqpoint{2.659912in}{0.525681in}}%
\pgfpathlineto{\pgfqpoint{2.660242in}{0.525687in}}%
\pgfpathlineto{\pgfqpoint{2.667616in}{0.526296in}}%
\pgfpathlineto{\pgfqpoint{2.667785in}{0.894940in}}%
\pgfpathlineto{\pgfqpoint{2.667861in}{1.325994in}}%
\pgfpathlineto{\pgfqpoint{2.668589in}{1.325905in}}%
\pgfpathlineto{\pgfqpoint{2.669743in}{1.325578in}}%
\pgfpathlineto{\pgfqpoint{2.670567in}{1.325605in}}%
\pgfpathlineto{\pgfqpoint{2.671721in}{1.325136in}}%
\pgfpathlineto{\pgfqpoint{2.672545in}{1.325178in}}%
\pgfpathlineto{\pgfqpoint{2.673369in}{1.324528in}}%
\pgfpathlineto{\pgfqpoint{2.674029in}{1.324127in}}%
\pgfpathlineto{\pgfqpoint{2.674193in}{1.324602in}}%
\pgfpathlineto{\pgfqpoint{2.675018in}{1.322844in}}%
\pgfpathlineto{\pgfqpoint{2.675183in}{1.323987in}}%
\pgfpathlineto{\pgfqpoint{2.675347in}{1.321049in}}%
\pgfpathlineto{\pgfqpoint{2.675467in}{1.323709in}}%
\pgfpathlineto{\pgfqpoint{2.676554in}{0.525691in}}%
\pgfpathlineto{\pgfqpoint{2.676719in}{0.525697in}}%
\pgfpathlineto{\pgfqpoint{2.684084in}{0.526194in}}%
\pgfpathlineto{\pgfqpoint{2.684173in}{0.535149in}}%
\pgfpathlineto{\pgfqpoint{2.684329in}{1.326005in}}%
\pgfpathlineto{\pgfqpoint{2.685041in}{1.325922in}}%
\pgfpathlineto{\pgfqpoint{2.686195in}{1.325559in}}%
\pgfpathlineto{\pgfqpoint{2.687019in}{1.325625in}}%
\pgfpathlineto{\pgfqpoint{2.688173in}{1.325105in}}%
\pgfpathlineto{\pgfqpoint{2.688998in}{1.325209in}}%
\pgfpathlineto{\pgfqpoint{2.689822in}{1.324472in}}%
\pgfpathlineto{\pgfqpoint{2.690316in}{1.324792in}}%
\pgfpathlineto{\pgfqpoint{2.690481in}{1.324044in}}%
\pgfpathlineto{\pgfqpoint{2.690646in}{1.324655in}}%
\pgfpathlineto{\pgfqpoint{2.691470in}{1.322524in}}%
\pgfpathlineto{\pgfqpoint{2.691635in}{1.324081in}}%
\pgfpathlineto{\pgfqpoint{2.692414in}{0.877447in}}%
\pgfpathlineto{\pgfqpoint{2.692854in}{0.525682in}}%
\pgfpathlineto{\pgfqpoint{2.693184in}{0.525687in}}%
\pgfpathlineto{\pgfqpoint{2.700580in}{0.526252in}}%
\pgfpathlineto{\pgfqpoint{2.700667in}{0.558189in}}%
\pgfpathlineto{\pgfqpoint{2.700826in}{1.326001in}}%
\pgfpathlineto{\pgfqpoint{2.701552in}{1.325913in}}%
\pgfpathlineto{\pgfqpoint{2.702706in}{1.325570in}}%
\pgfpathlineto{\pgfqpoint{2.703530in}{1.325615in}}%
\pgfpathlineto{\pgfqpoint{2.704684in}{1.325123in}}%
\pgfpathlineto{\pgfqpoint{2.705509in}{1.325193in}}%
\pgfpathlineto{\pgfqpoint{2.706333in}{1.324505in}}%
\pgfpathlineto{\pgfqpoint{2.706992in}{1.324094in}}%
\pgfpathlineto{\pgfqpoint{2.707157in}{1.324628in}}%
\pgfpathlineto{\pgfqpoint{2.707981in}{1.322726in}}%
\pgfpathlineto{\pgfqpoint{2.708146in}{1.324033in}}%
\pgfpathlineto{\pgfqpoint{2.708311in}{1.320069in}}%
\pgfpathlineto{\pgfqpoint{2.708434in}{1.323764in}}%
\pgfpathlineto{\pgfqpoint{2.709339in}{0.525681in}}%
\pgfpathlineto{\pgfqpoint{2.709504in}{0.525704in}}%
\pgfpathlineto{\pgfqpoint{2.715274in}{0.525853in}}%
\pgfpathlineto{\pgfqpoint{2.717064in}{0.526252in}}%
\pgfpathlineto{\pgfqpoint{2.717152in}{0.558182in}}%
\pgfpathlineto{\pgfqpoint{2.717311in}{1.326000in}}%
\pgfpathlineto{\pgfqpoint{2.718037in}{1.325913in}}%
\pgfpathlineto{\pgfqpoint{2.719191in}{1.325569in}}%
\pgfpathlineto{\pgfqpoint{2.720015in}{1.325614in}}%
\pgfpathlineto{\pgfqpoint{2.721169in}{1.325122in}}%
\pgfpathlineto{\pgfqpoint{2.721993in}{1.325192in}}%
\pgfpathlineto{\pgfqpoint{2.722818in}{1.324503in}}%
\pgfpathlineto{\pgfqpoint{2.723312in}{1.324766in}}%
\pgfpathlineto{\pgfqpoint{2.723477in}{1.324091in}}%
\pgfpathlineto{\pgfqpoint{2.723642in}{1.324626in}}%
\pgfpathlineto{\pgfqpoint{2.724466in}{1.322712in}}%
\pgfpathlineto{\pgfqpoint{2.724631in}{1.324029in}}%
\pgfpathlineto{\pgfqpoint{2.724796in}{1.319874in}}%
\pgfpathlineto{\pgfqpoint{2.724919in}{1.323759in}}%
\pgfpathlineto{\pgfqpoint{2.725824in}{0.525681in}}%
\pgfpathlineto{\pgfqpoint{2.725989in}{0.525704in}}%
\pgfpathlineto{\pgfqpoint{2.731759in}{0.525852in}}%
\pgfpathlineto{\pgfqpoint{2.733549in}{0.526252in}}%
\pgfpathlineto{\pgfqpoint{2.733637in}{0.558175in}}%
\pgfpathlineto{\pgfqpoint{2.733796in}{1.326000in}}%
\pgfpathlineto{\pgfqpoint{2.734522in}{1.325912in}}%
\pgfpathlineto{\pgfqpoint{2.735676in}{1.325569in}}%
\pgfpathlineto{\pgfqpoint{2.736500in}{1.325613in}}%
\pgfpathlineto{\pgfqpoint{2.737654in}{1.325121in}}%
\pgfpathlineto{\pgfqpoint{2.738478in}{1.325190in}}%
\pgfpathlineto{\pgfqpoint{2.739303in}{1.324501in}}%
\pgfpathlineto{\pgfqpoint{2.739797in}{1.324765in}}%
\pgfpathlineto{\pgfqpoint{2.739962in}{1.324087in}}%
\pgfpathlineto{\pgfqpoint{2.740127in}{1.324624in}}%
\pgfpathlineto{\pgfqpoint{2.740951in}{1.322698in}}%
\pgfpathlineto{\pgfqpoint{2.741116in}{1.324025in}}%
\pgfpathlineto{\pgfqpoint{2.741281in}{1.319645in}}%
\pgfpathlineto{\pgfqpoint{2.741403in}{1.323755in}}%
\pgfpathlineto{\pgfqpoint{2.742332in}{0.525681in}}%
\pgfpathlineto{\pgfqpoint{2.742497in}{0.525704in}}%
\pgfpathlineto{\pgfqpoint{2.748267in}{0.525853in}}%
\pgfpathlineto{\pgfqpoint{2.750034in}{0.526252in}}%
\pgfpathlineto{\pgfqpoint{2.750122in}{0.558169in}}%
\pgfpathlineto{\pgfqpoint{2.750280in}{1.325999in}}%
\pgfpathlineto{\pgfqpoint{2.751007in}{1.325912in}}%
\pgfpathlineto{\pgfqpoint{2.752161in}{1.325568in}}%
\pgfpathlineto{\pgfqpoint{2.752985in}{1.325612in}}%
\pgfpathlineto{\pgfqpoint{2.754139in}{1.325120in}}%
\pgfpathlineto{\pgfqpoint{2.754963in}{1.325189in}}%
\pgfpathlineto{\pgfqpoint{2.755787in}{1.324499in}}%
\pgfpathlineto{\pgfqpoint{2.756282in}{1.324763in}}%
\pgfpathlineto{\pgfqpoint{2.756447in}{1.324084in}}%
\pgfpathlineto{\pgfqpoint{2.756612in}{1.324622in}}%
\pgfpathlineto{\pgfqpoint{2.757436in}{1.322686in}}%
\pgfpathlineto{\pgfqpoint{2.757601in}{1.324022in}}%
\pgfpathlineto{\pgfqpoint{2.757766in}{1.319390in}}%
\pgfpathlineto{\pgfqpoint{2.757888in}{1.323749in}}%
\pgfpathlineto{\pgfqpoint{2.758780in}{0.525683in}}%
\pgfpathlineto{\pgfqpoint{2.758945in}{0.525702in}}%
\pgfpathlineto{\pgfqpoint{2.765374in}{0.525873in}}%
\pgfpathlineto{\pgfqpoint{2.766523in}{0.526276in}}%
\pgfpathlineto{\pgfqpoint{2.766708in}{0.943389in}}%
\pgfpathlineto{\pgfqpoint{2.766786in}{1.325993in}}%
\pgfpathlineto{\pgfqpoint{2.767535in}{1.325905in}}%
\pgfpathlineto{\pgfqpoint{2.768689in}{1.325565in}}%
\pgfpathlineto{\pgfqpoint{2.769513in}{1.325604in}}%
\pgfpathlineto{\pgfqpoint{2.770667in}{1.325115in}}%
\pgfpathlineto{\pgfqpoint{2.771492in}{1.325176in}}%
\pgfpathlineto{\pgfqpoint{2.772316in}{1.324489in}}%
\pgfpathlineto{\pgfqpoint{2.772810in}{1.324743in}}%
\pgfpathlineto{\pgfqpoint{2.772975in}{1.324071in}}%
\pgfpathlineto{\pgfqpoint{2.773140in}{1.324599in}}%
\pgfpathlineto{\pgfqpoint{2.773964in}{1.322632in}}%
\pgfpathlineto{\pgfqpoint{2.774129in}{1.323981in}}%
\pgfpathlineto{\pgfqpoint{2.774294in}{1.313973in}}%
\pgfpathlineto{\pgfqpoint{2.774395in}{1.323713in}}%
\pgfpathlineto{\pgfqpoint{2.775531in}{0.525689in}}%
\pgfpathlineto{\pgfqpoint{2.783023in}{0.526410in}}%
\pgfpathlineto{\pgfqpoint{2.783272in}{1.325976in}}%
\pgfpathlineto{\pgfqpoint{2.784667in}{1.325796in}}%
\pgfpathlineto{\pgfqpoint{2.786150in}{1.325395in}}%
\pgfpathlineto{\pgfqpoint{2.786975in}{1.325385in}}%
\pgfpathlineto{\pgfqpoint{2.788129in}{1.324841in}}%
\pgfpathlineto{\pgfqpoint{2.788953in}{1.324823in}}%
\pgfpathlineto{\pgfqpoint{2.790107in}{1.323557in}}%
\pgfpathlineto{\pgfqpoint{2.790272in}{1.324143in}}%
\pgfpathlineto{\pgfqpoint{2.790436in}{1.323009in}}%
\pgfpathlineto{\pgfqpoint{2.790601in}{1.323868in}}%
\pgfpathlineto{\pgfqpoint{2.790766in}{1.321757in}}%
\pgfpathlineto{\pgfqpoint{2.790873in}{1.323566in}}%
\pgfpathlineto{\pgfqpoint{2.791868in}{0.525690in}}%
\pgfpathlineto{\pgfqpoint{2.792693in}{0.525708in}}%
\pgfpathlineto{\pgfqpoint{2.799487in}{0.526240in}}%
\pgfpathlineto{\pgfqpoint{2.799576in}{0.555873in}}%
\pgfpathlineto{\pgfqpoint{2.799734in}{1.325996in}}%
\pgfpathlineto{\pgfqpoint{2.800460in}{1.325908in}}%
\pgfpathlineto{\pgfqpoint{2.801614in}{1.325559in}}%
\pgfpathlineto{\pgfqpoint{2.802438in}{1.325608in}}%
\pgfpathlineto{\pgfqpoint{2.803592in}{1.325107in}}%
\pgfpathlineto{\pgfqpoint{2.804416in}{1.325183in}}%
\pgfpathlineto{\pgfqpoint{2.805241in}{1.324475in}}%
\pgfpathlineto{\pgfqpoint{2.805735in}{1.324753in}}%
\pgfpathlineto{\pgfqpoint{2.805900in}{1.324048in}}%
\pgfpathlineto{\pgfqpoint{2.806065in}{1.324610in}}%
\pgfpathlineto{\pgfqpoint{2.806889in}{1.322541in}}%
\pgfpathlineto{\pgfqpoint{2.807054in}{1.324001in}}%
\pgfpathlineto{\pgfqpoint{2.807805in}{0.909341in}}%
\pgfpathlineto{\pgfqpoint{2.808209in}{0.525688in}}%
\pgfpathlineto{\pgfqpoint{2.808539in}{0.525690in}}%
\pgfpathlineto{\pgfqpoint{2.816008in}{0.526609in}}%
\pgfpathlineto{\pgfqpoint{2.816249in}{1.325965in}}%
\pgfpathlineto{\pgfqpoint{2.817313in}{1.325830in}}%
\pgfpathlineto{\pgfqpoint{2.819126in}{1.325410in}}%
\pgfpathlineto{\pgfqpoint{2.820280in}{1.325290in}}%
\pgfpathlineto{\pgfqpoint{2.821434in}{1.324739in}}%
\pgfpathlineto{\pgfqpoint{2.822258in}{1.324654in}}%
\pgfpathlineto{\pgfqpoint{2.823412in}{1.323150in}}%
\pgfpathlineto{\pgfqpoint{2.823577in}{1.323785in}}%
\pgfpathlineto{\pgfqpoint{2.823742in}{1.322171in}}%
\pgfpathlineto{\pgfqpoint{2.823846in}{1.323459in}}%
\pgfpathlineto{\pgfqpoint{2.823950in}{1.318074in}}%
\pgfpathlineto{\pgfqpoint{2.824217in}{1.321217in}}%
\pgfpathlineto{\pgfqpoint{2.824862in}{0.525687in}}%
\pgfpathlineto{\pgfqpoint{2.825356in}{0.525702in}}%
\pgfpathlineto{\pgfqpoint{2.832428in}{0.526111in}}%
\pgfpathlineto{\pgfqpoint{2.832517in}{0.527523in}}%
\pgfpathlineto{\pgfqpoint{2.832675in}{1.325953in}}%
\pgfpathlineto{\pgfqpoint{2.833375in}{1.325867in}}%
\pgfpathlineto{\pgfqpoint{2.835518in}{1.325428in}}%
\pgfpathlineto{\pgfqpoint{2.836672in}{1.325275in}}%
\pgfpathlineto{\pgfqpoint{2.838156in}{1.324635in}}%
\pgfpathlineto{\pgfqpoint{2.838980in}{1.324464in}}%
\pgfpathlineto{\pgfqpoint{2.840134in}{1.322538in}}%
\pgfpathlineto{\pgfqpoint{2.840285in}{1.323296in}}%
\pgfpathlineto{\pgfqpoint{2.840435in}{1.315064in}}%
\pgfpathlineto{\pgfqpoint{2.840702in}{1.318731in}}%
\pgfpathlineto{\pgfqpoint{2.840963in}{0.525666in}}%
\pgfpathlineto{\pgfqpoint{2.841863in}{0.525689in}}%
\pgfpathlineto{\pgfqpoint{2.848943in}{0.526243in}}%
\pgfpathlineto{\pgfqpoint{2.849031in}{0.557977in}}%
\pgfpathlineto{\pgfqpoint{2.849189in}{1.325986in}}%
\pgfpathlineto{\pgfqpoint{2.849916in}{1.325897in}}%
\pgfpathlineto{\pgfqpoint{2.851070in}{1.325548in}}%
\pgfpathlineto{\pgfqpoint{2.851894in}{1.325594in}}%
\pgfpathlineto{\pgfqpoint{2.853048in}{1.325090in}}%
\pgfpathlineto{\pgfqpoint{2.853872in}{1.325162in}}%
\pgfpathlineto{\pgfqpoint{2.854696in}{1.324444in}}%
\pgfpathlineto{\pgfqpoint{2.855026in}{1.324247in}}%
\pgfpathlineto{\pgfqpoint{2.855191in}{1.324722in}}%
\pgfpathlineto{\pgfqpoint{2.856015in}{1.323215in}}%
\pgfpathlineto{\pgfqpoint{2.856180in}{1.324196in}}%
\pgfpathlineto{\pgfqpoint{2.856345in}{1.322334in}}%
\pgfpathlineto{\pgfqpoint{2.856510in}{1.323937in}}%
\pgfpathlineto{\pgfqpoint{2.857670in}{0.525679in}}%
\pgfpathlineto{\pgfqpoint{2.858329in}{0.525692in}}%
\pgfpathlineto{\pgfqpoint{2.865438in}{0.526316in}}%
\pgfpathlineto{\pgfqpoint{2.865620in}{0.951276in}}%
\pgfpathlineto{\pgfqpoint{2.865696in}{1.325982in}}%
\pgfpathlineto{\pgfqpoint{2.866442in}{1.325893in}}%
\pgfpathlineto{\pgfqpoint{2.867596in}{1.325563in}}%
\pgfpathlineto{\pgfqpoint{2.868420in}{1.325589in}}%
\pgfpathlineto{\pgfqpoint{2.869574in}{1.325115in}}%
\pgfpathlineto{\pgfqpoint{2.870398in}{1.325156in}}%
\pgfpathlineto{\pgfqpoint{2.871222in}{1.324492in}}%
\pgfpathlineto{\pgfqpoint{2.871552in}{1.324304in}}%
\pgfpathlineto{\pgfqpoint{2.871717in}{1.324713in}}%
\pgfpathlineto{\pgfqpoint{2.872541in}{1.323364in}}%
\pgfpathlineto{\pgfqpoint{2.872706in}{1.324182in}}%
\pgfpathlineto{\pgfqpoint{2.872871in}{1.322656in}}%
\pgfpathlineto{\pgfqpoint{2.873036in}{1.323919in}}%
\pgfpathlineto{\pgfqpoint{2.873200in}{1.318422in}}%
\pgfpathlineto{\pgfqpoint{2.873303in}{1.323637in}}%
\pgfpathlineto{\pgfqpoint{2.874126in}{0.525685in}}%
\pgfpathlineto{\pgfqpoint{2.874455in}{0.525687in}}%
\pgfpathlineto{\pgfqpoint{2.881932in}{0.526399in}}%
\pgfpathlineto{\pgfqpoint{2.882181in}{1.325976in}}%
\pgfpathlineto{\pgfqpoint{2.883577in}{1.325796in}}%
\pgfpathlineto{\pgfqpoint{2.885060in}{1.325392in}}%
\pgfpathlineto{\pgfqpoint{2.885885in}{1.325386in}}%
\pgfpathlineto{\pgfqpoint{2.887038in}{1.324839in}}%
\pgfpathlineto{\pgfqpoint{2.887863in}{1.324827in}}%
\pgfpathlineto{\pgfqpoint{2.889017in}{1.323555in}}%
\pgfpathlineto{\pgfqpoint{2.889182in}{1.324153in}}%
\pgfpathlineto{\pgfqpoint{2.889346in}{1.323007in}}%
\pgfpathlineto{\pgfqpoint{2.889511in}{1.323882in}}%
\pgfpathlineto{\pgfqpoint{2.889676in}{1.321755in}}%
\pgfpathlineto{\pgfqpoint{2.889783in}{1.323586in}}%
\pgfpathlineto{\pgfqpoint{2.890654in}{0.525684in}}%
\pgfpathlineto{\pgfqpoint{2.891643in}{0.525703in}}%
\pgfpathlineto{\pgfqpoint{2.898404in}{0.526294in}}%
\pgfpathlineto{\pgfqpoint{2.898572in}{0.892756in}}%
\pgfpathlineto{\pgfqpoint{2.898648in}{1.325992in}}%
\pgfpathlineto{\pgfqpoint{2.899375in}{1.325903in}}%
\pgfpathlineto{\pgfqpoint{2.900529in}{1.325576in}}%
\pgfpathlineto{\pgfqpoint{2.901353in}{1.325603in}}%
\pgfpathlineto{\pgfqpoint{2.902507in}{1.325135in}}%
\pgfpathlineto{\pgfqpoint{2.903331in}{1.325176in}}%
\pgfpathlineto{\pgfqpoint{2.904156in}{1.324528in}}%
\pgfpathlineto{\pgfqpoint{2.904815in}{1.324129in}}%
\pgfpathlineto{\pgfqpoint{2.904980in}{1.324601in}}%
\pgfpathlineto{\pgfqpoint{2.905804in}{1.322855in}}%
\pgfpathlineto{\pgfqpoint{2.905969in}{1.323987in}}%
\pgfpathlineto{\pgfqpoint{2.906134in}{1.321120in}}%
\pgfpathlineto{\pgfqpoint{2.906254in}{1.323708in}}%
\pgfpathlineto{\pgfqpoint{2.907377in}{0.525691in}}%
\pgfpathlineto{\pgfqpoint{2.907542in}{0.525699in}}%
\pgfpathlineto{\pgfqpoint{2.914840in}{0.526079in}}%
\pgfpathlineto{\pgfqpoint{2.914930in}{0.526963in}}%
\pgfpathlineto{\pgfqpoint{2.915149in}{1.325982in}}%
\pgfpathlineto{\pgfqpoint{2.915862in}{1.325889in}}%
\pgfpathlineto{\pgfqpoint{2.917346in}{1.325524in}}%
\pgfpathlineto{\pgfqpoint{2.918500in}{1.325458in}}%
\pgfpathlineto{\pgfqpoint{2.919654in}{1.324950in}}%
\pgfpathlineto{\pgfqpoint{2.920478in}{1.324947in}}%
\pgfpathlineto{\pgfqpoint{2.921632in}{1.323899in}}%
\pgfpathlineto{\pgfqpoint{2.921797in}{1.324373in}}%
\pgfpathlineto{\pgfqpoint{2.921961in}{1.323537in}}%
\pgfpathlineto{\pgfqpoint{2.922126in}{1.324159in}}%
\pgfpathlineto{\pgfqpoint{2.922621in}{1.321650in}}%
\pgfpathlineto{\pgfqpoint{2.922740in}{1.323582in}}%
\pgfpathlineto{\pgfqpoint{2.923577in}{0.525685in}}%
\pgfpathlineto{\pgfqpoint{2.924236in}{0.525694in}}%
\pgfpathlineto{\pgfqpoint{2.931387in}{0.526406in}}%
\pgfpathlineto{\pgfqpoint{2.931635in}{1.325975in}}%
\pgfpathlineto{\pgfqpoint{2.933031in}{1.325795in}}%
\pgfpathlineto{\pgfqpoint{2.934514in}{1.325394in}}%
\pgfpathlineto{\pgfqpoint{2.935338in}{1.325385in}}%
\pgfpathlineto{\pgfqpoint{2.936492in}{1.324841in}}%
\pgfpathlineto{\pgfqpoint{2.937317in}{1.324825in}}%
\pgfpathlineto{\pgfqpoint{2.938471in}{1.323561in}}%
\pgfpathlineto{\pgfqpoint{2.938635in}{1.324148in}}%
\pgfpathlineto{\pgfqpoint{2.938800in}{1.323018in}}%
\pgfpathlineto{\pgfqpoint{2.938965in}{1.323876in}}%
\pgfpathlineto{\pgfqpoint{2.939130in}{1.321792in}}%
\pgfpathlineto{\pgfqpoint{2.939237in}{1.323578in}}%
\pgfpathlineto{\pgfqpoint{2.940115in}{0.525681in}}%
\pgfpathlineto{\pgfqpoint{2.941104in}{0.525702in}}%
\pgfpathlineto{\pgfqpoint{2.947852in}{0.526251in}}%
\pgfpathlineto{\pgfqpoint{2.947940in}{0.558151in}}%
\pgfpathlineto{\pgfqpoint{2.948099in}{1.325998in}}%
\pgfpathlineto{\pgfqpoint{2.948825in}{1.325911in}}%
\pgfpathlineto{\pgfqpoint{2.949979in}{1.325567in}}%
\pgfpathlineto{\pgfqpoint{2.950803in}{1.325612in}}%
\pgfpathlineto{\pgfqpoint{2.951957in}{1.325120in}}%
\pgfpathlineto{\pgfqpoint{2.952781in}{1.325190in}}%
\pgfpathlineto{\pgfqpoint{2.953606in}{1.324501in}}%
\pgfpathlineto{\pgfqpoint{2.954100in}{1.324765in}}%
\pgfpathlineto{\pgfqpoint{2.954265in}{1.324089in}}%
\pgfpathlineto{\pgfqpoint{2.954430in}{1.324625in}}%
\pgfpathlineto{\pgfqpoint{2.955254in}{1.322711in}}%
\pgfpathlineto{\pgfqpoint{2.955419in}{1.324029in}}%
\pgfpathlineto{\pgfqpoint{2.955584in}{1.319897in}}%
\pgfpathlineto{\pgfqpoint{2.955706in}{1.323760in}}%
\pgfpathlineto{\pgfqpoint{2.956612in}{0.525681in}}%
\pgfpathlineto{\pgfqpoint{2.956777in}{0.525704in}}%
\pgfpathlineto{\pgfqpoint{2.962547in}{0.525853in}}%
\pgfpathlineto{\pgfqpoint{2.964337in}{0.526253in}}%
\pgfpathlineto{\pgfqpoint{2.964425in}{0.558208in}}%
\pgfpathlineto{\pgfqpoint{2.964583in}{1.326002in}}%
\pgfpathlineto{\pgfqpoint{2.965310in}{1.325915in}}%
\pgfpathlineto{\pgfqpoint{2.966464in}{1.325573in}}%
\pgfpathlineto{\pgfqpoint{2.967288in}{1.325618in}}%
\pgfpathlineto{\pgfqpoint{2.968442in}{1.325129in}}%
\pgfpathlineto{\pgfqpoint{2.969266in}{1.325199in}}%
\pgfpathlineto{\pgfqpoint{2.970090in}{1.324517in}}%
\pgfpathlineto{\pgfqpoint{2.970750in}{1.324113in}}%
\pgfpathlineto{\pgfqpoint{2.970915in}{1.324639in}}%
\pgfpathlineto{\pgfqpoint{2.971739in}{1.322800in}}%
\pgfpathlineto{\pgfqpoint{2.971904in}{1.324055in}}%
\pgfpathlineto{\pgfqpoint{2.972069in}{1.320787in}}%
\pgfpathlineto{\pgfqpoint{2.972191in}{1.323794in}}%
\pgfpathlineto{\pgfqpoint{2.973105in}{0.525683in}}%
\pgfpathlineto{\pgfqpoint{2.973270in}{0.525705in}}%
\pgfpathlineto{\pgfqpoint{2.979039in}{0.525854in}}%
\pgfpathlineto{\pgfqpoint{2.980822in}{0.526255in}}%
\pgfpathlineto{\pgfqpoint{2.980910in}{0.558241in}}%
\pgfpathlineto{\pgfqpoint{2.981068in}{1.326004in}}%
\pgfpathlineto{\pgfqpoint{2.981795in}{1.325917in}}%
\pgfpathlineto{\pgfqpoint{2.982949in}{1.325576in}}%
\pgfpathlineto{\pgfqpoint{2.983773in}{1.325621in}}%
\pgfpathlineto{\pgfqpoint{2.984927in}{1.325134in}}%
\pgfpathlineto{\pgfqpoint{2.985751in}{1.325203in}}%
\pgfpathlineto{\pgfqpoint{2.986575in}{1.324526in}}%
\pgfpathlineto{\pgfqpoint{2.987235in}{1.324125in}}%
\pgfpathlineto{\pgfqpoint{2.987400in}{1.324646in}}%
\pgfpathlineto{\pgfqpoint{2.988224in}{1.322842in}}%
\pgfpathlineto{\pgfqpoint{2.988389in}{1.324067in}}%
\pgfpathlineto{\pgfqpoint{2.988553in}{1.321045in}}%
\pgfpathlineto{\pgfqpoint{2.988676in}{1.323810in}}%
\pgfpathlineto{\pgfqpoint{2.989610in}{0.525682in}}%
\pgfpathlineto{\pgfqpoint{2.989940in}{0.525688in}}%
\pgfpathlineto{\pgfqpoint{2.997307in}{0.526255in}}%
\pgfpathlineto{\pgfqpoint{2.997395in}{0.558253in}}%
\pgfpathlineto{\pgfqpoint{2.997553in}{1.326005in}}%
\pgfpathlineto{\pgfqpoint{2.998280in}{1.325918in}}%
\pgfpathlineto{\pgfqpoint{2.999433in}{1.325577in}}%
\pgfpathlineto{\pgfqpoint{3.000258in}{1.325622in}}%
\pgfpathlineto{\pgfqpoint{3.001412in}{1.325135in}}%
\pgfpathlineto{\pgfqpoint{3.002236in}{1.325204in}}%
\pgfpathlineto{\pgfqpoint{3.003060in}{1.324528in}}%
\pgfpathlineto{\pgfqpoint{3.003720in}{1.324128in}}%
\pgfpathlineto{\pgfqpoint{3.003884in}{1.324648in}}%
\pgfpathlineto{\pgfqpoint{3.004709in}{1.322849in}}%
\pgfpathlineto{\pgfqpoint{3.004873in}{1.324069in}}%
\pgfpathlineto{\pgfqpoint{3.005038in}{1.321084in}}%
\pgfpathlineto{\pgfqpoint{3.005161in}{1.323812in}}%
\pgfpathlineto{\pgfqpoint{3.006096in}{0.525682in}}%
\pgfpathlineto{\pgfqpoint{3.006426in}{0.525688in}}%
\pgfpathlineto{\pgfqpoint{3.013798in}{0.526299in}}%
\pgfpathlineto{\pgfqpoint{3.013967in}{0.894637in}}%
\pgfpathlineto{\pgfqpoint{3.014042in}{1.325999in}}%
\pgfpathlineto{\pgfqpoint{3.014770in}{1.325911in}}%
\pgfpathlineto{\pgfqpoint{3.015924in}{1.325585in}}%
\pgfpathlineto{\pgfqpoint{3.016748in}{1.325612in}}%
\pgfpathlineto{\pgfqpoint{3.017902in}{1.325147in}}%
\pgfpathlineto{\pgfqpoint{3.018727in}{1.325189in}}%
\pgfpathlineto{\pgfqpoint{3.019551in}{1.324549in}}%
\pgfpathlineto{\pgfqpoint{3.020210in}{1.324159in}}%
\pgfpathlineto{\pgfqpoint{3.020375in}{1.324622in}}%
\pgfpathlineto{\pgfqpoint{3.021199in}{1.322951in}}%
\pgfpathlineto{\pgfqpoint{3.021364in}{1.324023in}}%
\pgfpathlineto{\pgfqpoint{3.021529in}{1.321552in}}%
\pgfpathlineto{\pgfqpoint{3.021649in}{1.323756in}}%
\pgfpathlineto{\pgfqpoint{3.022541in}{0.525681in}}%
\pgfpathlineto{\pgfqpoint{3.023036in}{0.525707in}}%
\pgfpathlineto{\pgfqpoint{3.030098in}{0.525901in}}%
\pgfpathlineto{\pgfqpoint{3.030277in}{0.526253in}}%
\pgfpathlineto{\pgfqpoint{3.030364in}{0.558206in}}%
\pgfpathlineto{\pgfqpoint{3.030523in}{1.326002in}}%
\pgfpathlineto{\pgfqpoint{3.031249in}{1.325914in}}%
\pgfpathlineto{\pgfqpoint{3.032403in}{1.325572in}}%
\pgfpathlineto{\pgfqpoint{3.033227in}{1.325616in}}%
\pgfpathlineto{\pgfqpoint{3.034381in}{1.325126in}}%
\pgfpathlineto{\pgfqpoint{3.035206in}{1.325196in}}%
\pgfpathlineto{\pgfqpoint{3.036030in}{1.324511in}}%
\pgfpathlineto{\pgfqpoint{3.036689in}{1.324103in}}%
\pgfpathlineto{\pgfqpoint{3.036854in}{1.324633in}}%
\pgfpathlineto{\pgfqpoint{3.037678in}{1.322757in}}%
\pgfpathlineto{\pgfqpoint{3.037843in}{1.324042in}}%
\pgfpathlineto{\pgfqpoint{3.038008in}{1.320421in}}%
\pgfpathlineto{\pgfqpoint{3.038131in}{1.323776in}}%
\pgfpathlineto{\pgfqpoint{3.039036in}{0.525682in}}%
\pgfpathlineto{\pgfqpoint{3.039201in}{0.525705in}}%
\pgfpathlineto{\pgfqpoint{3.044971in}{0.525853in}}%
\pgfpathlineto{\pgfqpoint{3.046761in}{0.526253in}}%
\pgfpathlineto{\pgfqpoint{3.046849in}{0.558196in}}%
\pgfpathlineto{\pgfqpoint{3.047008in}{1.326001in}}%
\pgfpathlineto{\pgfqpoint{3.047734in}{1.325914in}}%
\pgfpathlineto{\pgfqpoint{3.048888in}{1.325571in}}%
\pgfpathlineto{\pgfqpoint{3.049712in}{1.325615in}}%
\pgfpathlineto{\pgfqpoint{3.050866in}{1.325125in}}%
\pgfpathlineto{\pgfqpoint{3.051690in}{1.325194in}}%
\pgfpathlineto{\pgfqpoint{3.052515in}{1.324507in}}%
\pgfpathlineto{\pgfqpoint{3.053174in}{1.324098in}}%
\pgfpathlineto{\pgfqpoint{3.053339in}{1.324630in}}%
\pgfpathlineto{\pgfqpoint{3.054163in}{1.322737in}}%
\pgfpathlineto{\pgfqpoint{3.054328in}{1.324036in}}%
\pgfpathlineto{\pgfqpoint{3.054493in}{1.320211in}}%
\pgfpathlineto{\pgfqpoint{3.054616in}{1.323769in}}%
\pgfpathlineto{\pgfqpoint{3.055521in}{0.525682in}}%
\pgfpathlineto{\pgfqpoint{3.055686in}{0.525704in}}%
\pgfpathlineto{\pgfqpoint{3.061455in}{0.525853in}}%
\pgfpathlineto{\pgfqpoint{3.063246in}{0.526252in}}%
\pgfpathlineto{\pgfqpoint{3.063334in}{0.558186in}}%
\pgfpathlineto{\pgfqpoint{3.063493in}{1.326000in}}%
\pgfpathlineto{\pgfqpoint{3.064219in}{1.325913in}}%
\pgfpathlineto{\pgfqpoint{3.065373in}{1.325570in}}%
\pgfpathlineto{\pgfqpoint{3.066197in}{1.325614in}}%
\pgfpathlineto{\pgfqpoint{3.067351in}{1.325123in}}%
\pgfpathlineto{\pgfqpoint{3.068175in}{1.325192in}}%
\pgfpathlineto{\pgfqpoint{3.069000in}{1.324504in}}%
\pgfpathlineto{\pgfqpoint{3.069494in}{1.324767in}}%
\pgfpathlineto{\pgfqpoint{3.069659in}{1.324092in}}%
\pgfpathlineto{\pgfqpoint{3.069824in}{1.324626in}}%
\pgfpathlineto{\pgfqpoint{3.070648in}{1.322716in}}%
\pgfpathlineto{\pgfqpoint{3.070813in}{1.324030in}}%
\pgfpathlineto{\pgfqpoint{3.070978in}{1.319938in}}%
\pgfpathlineto{\pgfqpoint{3.071100in}{1.323761in}}%
\pgfpathlineto{\pgfqpoint{3.072006in}{0.525681in}}%
\pgfpathlineto{\pgfqpoint{3.072171in}{0.525704in}}%
\pgfpathlineto{\pgfqpoint{3.077940in}{0.525852in}}%
\pgfpathlineto{\pgfqpoint{3.079731in}{0.526252in}}%
\pgfpathlineto{\pgfqpoint{3.079819in}{0.558176in}}%
\pgfpathlineto{\pgfqpoint{3.079977in}{1.326000in}}%
\pgfpathlineto{\pgfqpoint{3.080704in}{1.325912in}}%
\pgfpathlineto{\pgfqpoint{3.081858in}{1.325569in}}%
\pgfpathlineto{\pgfqpoint{3.082682in}{1.325613in}}%
\pgfpathlineto{\pgfqpoint{3.083836in}{1.325121in}}%
\pgfpathlineto{\pgfqpoint{3.084660in}{1.325190in}}%
\pgfpathlineto{\pgfqpoint{3.085484in}{1.324500in}}%
\pgfpathlineto{\pgfqpoint{3.085979in}{1.324765in}}%
\pgfpathlineto{\pgfqpoint{3.086144in}{1.324087in}}%
\pgfpathlineto{\pgfqpoint{3.086309in}{1.324624in}}%
\pgfpathlineto{\pgfqpoint{3.087133in}{1.322697in}}%
\pgfpathlineto{\pgfqpoint{3.087298in}{1.324025in}}%
\pgfpathlineto{\pgfqpoint{3.087463in}{1.319618in}}%
\pgfpathlineto{\pgfqpoint{3.087585in}{1.323755in}}%
\pgfpathlineto{\pgfqpoint{3.088515in}{0.525680in}}%
\pgfpathlineto{\pgfqpoint{3.088680in}{0.525704in}}%
\pgfpathlineto{\pgfqpoint{3.094284in}{0.525861in}}%
\pgfpathlineto{\pgfqpoint{3.096216in}{0.526251in}}%
\pgfpathlineto{\pgfqpoint{3.096304in}{0.558167in}}%
\pgfpathlineto{\pgfqpoint{3.096462in}{1.325999in}}%
\pgfpathlineto{\pgfqpoint{3.097189in}{1.325911in}}%
\pgfpathlineto{\pgfqpoint{3.098343in}{1.325568in}}%
\pgfpathlineto{\pgfqpoint{3.099167in}{1.325612in}}%
\pgfpathlineto{\pgfqpoint{3.100321in}{1.325119in}}%
\pgfpathlineto{\pgfqpoint{3.101145in}{1.325189in}}%
\pgfpathlineto{\pgfqpoint{3.101969in}{1.324498in}}%
\pgfpathlineto{\pgfqpoint{3.102464in}{1.324762in}}%
\pgfpathlineto{\pgfqpoint{3.102629in}{1.324083in}}%
\pgfpathlineto{\pgfqpoint{3.102793in}{1.324621in}}%
\pgfpathlineto{\pgfqpoint{3.103618in}{1.322681in}}%
\pgfpathlineto{\pgfqpoint{3.103783in}{1.324020in}}%
\pgfpathlineto{\pgfqpoint{3.103947in}{1.319256in}}%
\pgfpathlineto{\pgfqpoint{3.104070in}{1.323747in}}%
\pgfpathlineto{\pgfqpoint{3.104928in}{0.525689in}}%
\pgfpathlineto{\pgfqpoint{3.105089in}{0.525699in}}%
\pgfpathlineto{\pgfqpoint{3.112635in}{0.526025in}}%
\pgfpathlineto{\pgfqpoint{3.112722in}{0.526431in}}%
\pgfpathlineto{\pgfqpoint{3.112969in}{1.325978in}}%
\pgfpathlineto{\pgfqpoint{3.114198in}{1.325705in}}%
\pgfpathlineto{\pgfqpoint{3.115352in}{1.325639in}}%
\pgfpathlineto{\pgfqpoint{3.116835in}{1.325162in}}%
\pgfpathlineto{\pgfqpoint{3.117659in}{1.325141in}}%
\pgfpathlineto{\pgfqpoint{3.118813in}{1.324402in}}%
\pgfpathlineto{\pgfqpoint{3.119143in}{1.324195in}}%
\pgfpathlineto{\pgfqpoint{3.119308in}{1.324535in}}%
\pgfpathlineto{\pgfqpoint{3.120132in}{1.323061in}}%
\pgfpathlineto{\pgfqpoint{3.120297in}{1.323862in}}%
\pgfpathlineto{\pgfqpoint{3.120462in}{1.321923in}}%
\pgfpathlineto{\pgfqpoint{3.120570in}{1.323558in}}%
\pgfpathlineto{\pgfqpoint{3.121054in}{0.689589in}}%
\pgfpathlineto{\pgfqpoint{3.121554in}{0.525690in}}%
\pgfpathlineto{\pgfqpoint{3.121884in}{0.525694in}}%
\pgfpathlineto{\pgfqpoint{3.129184in}{0.526240in}}%
\pgfpathlineto{\pgfqpoint{3.129273in}{0.555880in}}%
\pgfpathlineto{\pgfqpoint{3.129431in}{1.325996in}}%
\pgfpathlineto{\pgfqpoint{3.130157in}{1.325909in}}%
\pgfpathlineto{\pgfqpoint{3.131311in}{1.325560in}}%
\pgfpathlineto{\pgfqpoint{3.132135in}{1.325609in}}%
\pgfpathlineto{\pgfqpoint{3.133289in}{1.325108in}}%
\pgfpathlineto{\pgfqpoint{3.134113in}{1.325184in}}%
\pgfpathlineto{\pgfqpoint{3.134938in}{1.324477in}}%
\pgfpathlineto{\pgfqpoint{3.135432in}{1.324754in}}%
\pgfpathlineto{\pgfqpoint{3.135597in}{1.324051in}}%
\pgfpathlineto{\pgfqpoint{3.135762in}{1.324612in}}%
\pgfpathlineto{\pgfqpoint{3.136586in}{1.322553in}}%
\pgfpathlineto{\pgfqpoint{3.136751in}{1.324004in}}%
\pgfpathlineto{\pgfqpoint{3.137490in}{1.074185in}}%
\pgfpathlineto{\pgfqpoint{3.137877in}{0.525686in}}%
\pgfpathlineto{\pgfqpoint{3.138196in}{0.525687in}}%
\pgfpathlineto{\pgfqpoint{3.145681in}{0.526328in}}%
\pgfpathlineto{\pgfqpoint{3.145847in}{0.895704in}}%
\pgfpathlineto{\pgfqpoint{3.145922in}{1.325985in}}%
\pgfpathlineto{\pgfqpoint{3.146648in}{1.325895in}}%
\pgfpathlineto{\pgfqpoint{3.148132in}{1.325519in}}%
\pgfpathlineto{\pgfqpoint{3.148956in}{1.325531in}}%
\pgfpathlineto{\pgfqpoint{3.150110in}{1.325044in}}%
\pgfpathlineto{\pgfqpoint{3.150934in}{1.325064in}}%
\pgfpathlineto{\pgfqpoint{3.152088in}{1.324138in}}%
\pgfpathlineto{\pgfqpoint{3.152253in}{1.324568in}}%
\pgfpathlineto{\pgfqpoint{3.152418in}{1.323863in}}%
\pgfpathlineto{\pgfqpoint{3.152583in}{1.324394in}}%
\pgfpathlineto{\pgfqpoint{3.153407in}{1.321253in}}%
\pgfpathlineto{\pgfqpoint{3.153527in}{1.323624in}}%
\pgfpathlineto{\pgfqpoint{3.154482in}{0.525690in}}%
\pgfpathlineto{\pgfqpoint{3.154812in}{0.525694in}}%
\pgfpathlineto{\pgfqpoint{3.162154in}{0.526240in}}%
\pgfpathlineto{\pgfqpoint{3.162242in}{0.555876in}}%
\pgfpathlineto{\pgfqpoint{3.162401in}{1.325996in}}%
\pgfpathlineto{\pgfqpoint{3.163127in}{1.325909in}}%
\pgfpathlineto{\pgfqpoint{3.164281in}{1.325560in}}%
\pgfpathlineto{\pgfqpoint{3.165105in}{1.325608in}}%
\pgfpathlineto{\pgfqpoint{3.166259in}{1.325108in}}%
\pgfpathlineto{\pgfqpoint{3.167083in}{1.325184in}}%
\pgfpathlineto{\pgfqpoint{3.167907in}{1.324477in}}%
\pgfpathlineto{\pgfqpoint{3.168402in}{1.324755in}}%
\pgfpathlineto{\pgfqpoint{3.168567in}{1.324052in}}%
\pgfpathlineto{\pgfqpoint{3.168732in}{1.324612in}}%
\pgfpathlineto{\pgfqpoint{3.169556in}{1.322557in}}%
\pgfpathlineto{\pgfqpoint{3.169721in}{1.324005in}}%
\pgfpathlineto{\pgfqpoint{3.170132in}{1.086804in}}%
\pgfpathlineto{\pgfqpoint{3.170132in}{1.086804in}}%
\pgfusepath{stroke}%
\end{pgfscope}%
\begin{pgfscope}%
\pgfsetrectcap%
\pgfsetmiterjoin%
\pgfsetlinewidth{0.803000pt}%
\definecolor{currentstroke}{rgb}{0.000000,0.000000,0.000000}%
\pgfsetstrokecolor{currentstroke}%
\pgfsetdash{}{0pt}%
\pgfpathmoveto{\pgfqpoint{0.573768in}{0.462654in}}%
\pgfpathlineto{\pgfqpoint{0.573768in}{1.368904in}}%
\pgfusepath{stroke}%
\end{pgfscope}%
\begin{pgfscope}%
\pgfsetrectcap%
\pgfsetmiterjoin%
\pgfsetlinewidth{0.803000pt}%
\definecolor{currentstroke}{rgb}{0.000000,0.000000,0.000000}%
\pgfsetstrokecolor{currentstroke}%
\pgfsetdash{}{0pt}%
\pgfpathmoveto{\pgfqpoint{3.293768in}{0.462654in}}%
\pgfpathlineto{\pgfqpoint{3.293768in}{1.368904in}}%
\pgfusepath{stroke}%
\end{pgfscope}%
\begin{pgfscope}%
\pgfsetrectcap%
\pgfsetmiterjoin%
\pgfsetlinewidth{0.803000pt}%
\definecolor{currentstroke}{rgb}{0.000000,0.000000,0.000000}%
\pgfsetstrokecolor{currentstroke}%
\pgfsetdash{}{0pt}%
\pgfpathmoveto{\pgfqpoint{0.573768in}{0.462654in}}%
\pgfpathlineto{\pgfqpoint{3.293768in}{0.462654in}}%
\pgfusepath{stroke}%
\end{pgfscope}%
\begin{pgfscope}%
\pgfsetrectcap%
\pgfsetmiterjoin%
\pgfsetlinewidth{0.803000pt}%
\definecolor{currentstroke}{rgb}{0.000000,0.000000,0.000000}%
\pgfsetstrokecolor{currentstroke}%
\pgfsetdash{}{0pt}%
\pgfpathmoveto{\pgfqpoint{0.573768in}{1.368904in}}%
\pgfpathlineto{\pgfqpoint{3.293768in}{1.368904in}}%
\pgfusepath{stroke}%
\end{pgfscope}%
\begin{pgfscope}%
\definecolor{textcolor}{rgb}{0.000000,0.000000,0.000000}%
\pgfsetstrokecolor{textcolor}%
\pgfsetfillcolor{textcolor}%
\pgftext[x=0.628168in,y=1.305467in,left,top]{\color{textcolor}{\rmfamily\fontsize{8.000000}{9.600000}\selectfont\catcode`\^=\active\def^{\ifmmode\sp\else\^{}\fi}\catcode`\%=\active\def%{\%}(c)}}%
\end{pgfscope}%
\begin{pgfscope}%
\pgfsetbuttcap%
\pgfsetmiterjoin%
\definecolor{currentfill}{rgb}{1.000000,1.000000,1.000000}%
\pgfsetfillcolor{currentfill}%
\pgfsetlinewidth{0.000000pt}%
\definecolor{currentstroke}{rgb}{0.000000,0.000000,0.000000}%
\pgfsetstrokecolor{currentstroke}%
\pgfsetstrokeopacity{0.000000}%
\pgfsetdash{}{0pt}%
\pgfpathmoveto{\pgfqpoint{2.259013in}{3.125849in}}%
\pgfpathlineto{\pgfqpoint{3.238213in}{3.125849in}}%
\pgfpathlineto{\pgfqpoint{3.238213in}{3.452099in}}%
\pgfpathlineto{\pgfqpoint{2.259013in}{3.452099in}}%
\pgfpathlineto{\pgfqpoint{2.259013in}{3.125849in}}%
\pgfpathclose%
\pgfusepath{fill}%
\end{pgfscope}%
\begin{pgfscope}%
\pgfpathrectangle{\pgfqpoint{2.259013in}{3.125849in}}{\pgfqpoint{0.979200in}{0.326250in}}%
\pgfusepath{clip}%
\pgfsetbuttcap%
\pgfsetroundjoin%
\pgfsetlinewidth{0.401500pt}%
\definecolor{currentstroke}{rgb}{0.690196,0.690196,0.690196}%
\pgfsetstrokecolor{currentstroke}%
\pgfsetstrokeopacity{0.250000}%
\pgfsetdash{{1.480000pt}{0.640000pt}}{0.000000pt}%
\pgfpathmoveto{\pgfqpoint{2.259013in}{3.125849in}}%
\pgfpathlineto{\pgfqpoint{2.259013in}{3.452099in}}%
\pgfusepath{stroke}%
\end{pgfscope}%
\begin{pgfscope}%
\pgfsetbuttcap%
\pgfsetroundjoin%
\definecolor{currentfill}{rgb}{0.000000,0.000000,0.000000}%
\pgfsetfillcolor{currentfill}%
\pgfsetlinewidth{0.803000pt}%
\definecolor{currentstroke}{rgb}{0.000000,0.000000,0.000000}%
\pgfsetstrokecolor{currentstroke}%
\pgfsetdash{}{0pt}%
\pgfsys@defobject{currentmarker}{\pgfqpoint{0.000000in}{0.000000in}}{\pgfqpoint{0.000000in}{0.027778in}}{%
\pgfpathmoveto{\pgfqpoint{0.000000in}{0.000000in}}%
\pgfpathlineto{\pgfqpoint{0.000000in}{0.027778in}}%
\pgfusepath{stroke,fill}%
}%
\begin{pgfscope}%
\pgfsys@transformshift{2.259013in}{3.125849in}%
\pgfsys@useobject{currentmarker}{}%
\end{pgfscope}%
\end{pgfscope}%
\begin{pgfscope}%
\definecolor{textcolor}{rgb}{0.000000,0.000000,0.000000}%
\pgfsetstrokecolor{textcolor}%
\pgfsetfillcolor{textcolor}%
\pgftext[x=2.259013in,y=3.111960in,,top]{\color{textcolor}{\rmfamily\fontsize{5.000000}{6.000000}\selectfont\catcode`\^=\active\def^{\ifmmode\sp\else\^{}\fi}\catcode`\%=\active\def%{\%}5.9}}%
\end{pgfscope}%
\begin{pgfscope}%
\pgfpathrectangle{\pgfqpoint{2.259013in}{3.125849in}}{\pgfqpoint{0.979200in}{0.326250in}}%
\pgfusepath{clip}%
\pgfsetbuttcap%
\pgfsetroundjoin%
\pgfsetlinewidth{0.401500pt}%
\definecolor{currentstroke}{rgb}{0.690196,0.690196,0.690196}%
\pgfsetstrokecolor{currentstroke}%
\pgfsetstrokeopacity{0.250000}%
\pgfsetdash{{1.480000pt}{0.640000pt}}{0.000000pt}%
\pgfpathmoveto{\pgfqpoint{2.503813in}{3.125849in}}%
\pgfpathlineto{\pgfqpoint{2.503813in}{3.452099in}}%
\pgfusepath{stroke}%
\end{pgfscope}%
\begin{pgfscope}%
\pgfsetbuttcap%
\pgfsetroundjoin%
\definecolor{currentfill}{rgb}{0.000000,0.000000,0.000000}%
\pgfsetfillcolor{currentfill}%
\pgfsetlinewidth{0.803000pt}%
\definecolor{currentstroke}{rgb}{0.000000,0.000000,0.000000}%
\pgfsetstrokecolor{currentstroke}%
\pgfsetdash{}{0pt}%
\pgfsys@defobject{currentmarker}{\pgfqpoint{0.000000in}{0.000000in}}{\pgfqpoint{0.000000in}{0.027778in}}{%
\pgfpathmoveto{\pgfqpoint{0.000000in}{0.000000in}}%
\pgfpathlineto{\pgfqpoint{0.000000in}{0.027778in}}%
\pgfusepath{stroke,fill}%
}%
\begin{pgfscope}%
\pgfsys@transformshift{2.503813in}{3.125849in}%
\pgfsys@useobject{currentmarker}{}%
\end{pgfscope}%
\end{pgfscope}%
\begin{pgfscope}%
\definecolor{textcolor}{rgb}{0.000000,0.000000,0.000000}%
\pgfsetstrokecolor{textcolor}%
\pgfsetfillcolor{textcolor}%
\pgftext[x=2.503813in,y=3.111960in,,top]{\color{textcolor}{\rmfamily\fontsize{5.000000}{6.000000}\selectfont\catcode`\^=\active\def^{\ifmmode\sp\else\^{}\fi}\catcode`\%=\active\def%{\%}5.9}}%
\end{pgfscope}%
\begin{pgfscope}%
\pgfpathrectangle{\pgfqpoint{2.259013in}{3.125849in}}{\pgfqpoint{0.979200in}{0.326250in}}%
\pgfusepath{clip}%
\pgfsetbuttcap%
\pgfsetroundjoin%
\pgfsetlinewidth{0.401500pt}%
\definecolor{currentstroke}{rgb}{0.690196,0.690196,0.690196}%
\pgfsetstrokecolor{currentstroke}%
\pgfsetstrokeopacity{0.250000}%
\pgfsetdash{{1.480000pt}{0.640000pt}}{0.000000pt}%
\pgfpathmoveto{\pgfqpoint{2.748613in}{3.125849in}}%
\pgfpathlineto{\pgfqpoint{2.748613in}{3.452099in}}%
\pgfusepath{stroke}%
\end{pgfscope}%
\begin{pgfscope}%
\pgfsetbuttcap%
\pgfsetroundjoin%
\definecolor{currentfill}{rgb}{0.000000,0.000000,0.000000}%
\pgfsetfillcolor{currentfill}%
\pgfsetlinewidth{0.803000pt}%
\definecolor{currentstroke}{rgb}{0.000000,0.000000,0.000000}%
\pgfsetstrokecolor{currentstroke}%
\pgfsetdash{}{0pt}%
\pgfsys@defobject{currentmarker}{\pgfqpoint{0.000000in}{0.000000in}}{\pgfqpoint{0.000000in}{0.027778in}}{%
\pgfpathmoveto{\pgfqpoint{0.000000in}{0.000000in}}%
\pgfpathlineto{\pgfqpoint{0.000000in}{0.027778in}}%
\pgfusepath{stroke,fill}%
}%
\begin{pgfscope}%
\pgfsys@transformshift{2.748613in}{3.125849in}%
\pgfsys@useobject{currentmarker}{}%
\end{pgfscope}%
\end{pgfscope}%
\begin{pgfscope}%
\definecolor{textcolor}{rgb}{0.000000,0.000000,0.000000}%
\pgfsetstrokecolor{textcolor}%
\pgfsetfillcolor{textcolor}%
\pgftext[x=2.748613in,y=3.111960in,,top]{\color{textcolor}{\rmfamily\fontsize{5.000000}{6.000000}\selectfont\catcode`\^=\active\def^{\ifmmode\sp\else\^{}\fi}\catcode`\%=\active\def%{\%}6.0}}%
\end{pgfscope}%
\begin{pgfscope}%
\pgfpathrectangle{\pgfqpoint{2.259013in}{3.125849in}}{\pgfqpoint{0.979200in}{0.326250in}}%
\pgfusepath{clip}%
\pgfsetbuttcap%
\pgfsetroundjoin%
\pgfsetlinewidth{0.401500pt}%
\definecolor{currentstroke}{rgb}{0.690196,0.690196,0.690196}%
\pgfsetstrokecolor{currentstroke}%
\pgfsetstrokeopacity{0.250000}%
\pgfsetdash{{1.480000pt}{0.640000pt}}{0.000000pt}%
\pgfpathmoveto{\pgfqpoint{2.993413in}{3.125849in}}%
\pgfpathlineto{\pgfqpoint{2.993413in}{3.452099in}}%
\pgfusepath{stroke}%
\end{pgfscope}%
\begin{pgfscope}%
\pgfsetbuttcap%
\pgfsetroundjoin%
\definecolor{currentfill}{rgb}{0.000000,0.000000,0.000000}%
\pgfsetfillcolor{currentfill}%
\pgfsetlinewidth{0.803000pt}%
\definecolor{currentstroke}{rgb}{0.000000,0.000000,0.000000}%
\pgfsetstrokecolor{currentstroke}%
\pgfsetdash{}{0pt}%
\pgfsys@defobject{currentmarker}{\pgfqpoint{0.000000in}{0.000000in}}{\pgfqpoint{0.000000in}{0.027778in}}{%
\pgfpathmoveto{\pgfqpoint{0.000000in}{0.000000in}}%
\pgfpathlineto{\pgfqpoint{0.000000in}{0.027778in}}%
\pgfusepath{stroke,fill}%
}%
\begin{pgfscope}%
\pgfsys@transformshift{2.993413in}{3.125849in}%
\pgfsys@useobject{currentmarker}{}%
\end{pgfscope}%
\end{pgfscope}%
\begin{pgfscope}%
\definecolor{textcolor}{rgb}{0.000000,0.000000,0.000000}%
\pgfsetstrokecolor{textcolor}%
\pgfsetfillcolor{textcolor}%
\pgftext[x=2.993413in,y=3.111960in,,top]{\color{textcolor}{\rmfamily\fontsize{5.000000}{6.000000}\selectfont\catcode`\^=\active\def^{\ifmmode\sp\else\^{}\fi}\catcode`\%=\active\def%{\%}6.0}}%
\end{pgfscope}%
\begin{pgfscope}%
\pgfpathrectangle{\pgfqpoint{2.259013in}{3.125849in}}{\pgfqpoint{0.979200in}{0.326250in}}%
\pgfusepath{clip}%
\pgfsetbuttcap%
\pgfsetroundjoin%
\pgfsetlinewidth{0.401500pt}%
\definecolor{currentstroke}{rgb}{0.690196,0.690196,0.690196}%
\pgfsetstrokecolor{currentstroke}%
\pgfsetstrokeopacity{0.250000}%
\pgfsetdash{{1.480000pt}{0.640000pt}}{0.000000pt}%
\pgfpathmoveto{\pgfqpoint{3.238213in}{3.125849in}}%
\pgfpathlineto{\pgfqpoint{3.238213in}{3.452099in}}%
\pgfusepath{stroke}%
\end{pgfscope}%
\begin{pgfscope}%
\pgfsetbuttcap%
\pgfsetroundjoin%
\definecolor{currentfill}{rgb}{0.000000,0.000000,0.000000}%
\pgfsetfillcolor{currentfill}%
\pgfsetlinewidth{0.803000pt}%
\definecolor{currentstroke}{rgb}{0.000000,0.000000,0.000000}%
\pgfsetstrokecolor{currentstroke}%
\pgfsetdash{}{0pt}%
\pgfsys@defobject{currentmarker}{\pgfqpoint{0.000000in}{0.000000in}}{\pgfqpoint{0.000000in}{0.027778in}}{%
\pgfpathmoveto{\pgfqpoint{0.000000in}{0.000000in}}%
\pgfpathlineto{\pgfqpoint{0.000000in}{0.027778in}}%
\pgfusepath{stroke,fill}%
}%
\begin{pgfscope}%
\pgfsys@transformshift{3.238213in}{3.125849in}%
\pgfsys@useobject{currentmarker}{}%
\end{pgfscope}%
\end{pgfscope}%
\begin{pgfscope}%
\definecolor{textcolor}{rgb}{0.000000,0.000000,0.000000}%
\pgfsetstrokecolor{textcolor}%
\pgfsetfillcolor{textcolor}%
\pgftext[x=3.238213in,y=3.111960in,,top]{\color{textcolor}{\rmfamily\fontsize{5.000000}{6.000000}\selectfont\catcode`\^=\active\def^{\ifmmode\sp\else\^{}\fi}\catcode`\%=\active\def%{\%}6.0}}%
\end{pgfscope}%
\begin{pgfscope}%
\pgfpathrectangle{\pgfqpoint{2.259013in}{3.125849in}}{\pgfqpoint{0.979200in}{0.326250in}}%
\pgfusepath{clip}%
\pgfsetbuttcap%
\pgfsetroundjoin%
\pgfsetlinewidth{0.401500pt}%
\definecolor{currentstroke}{rgb}{0.690196,0.690196,0.690196}%
\pgfsetstrokecolor{currentstroke}%
\pgfsetstrokeopacity{0.250000}%
\pgfsetdash{{1.480000pt}{0.640000pt}}{0.000000pt}%
\pgfpathmoveto{\pgfqpoint{2.259013in}{3.125849in}}%
\pgfpathlineto{\pgfqpoint{3.238213in}{3.125849in}}%
\pgfusepath{stroke}%
\end{pgfscope}%
\begin{pgfscope}%
\pgfsetbuttcap%
\pgfsetroundjoin%
\definecolor{currentfill}{rgb}{0.000000,0.000000,0.000000}%
\pgfsetfillcolor{currentfill}%
\pgfsetlinewidth{0.803000pt}%
\definecolor{currentstroke}{rgb}{0.000000,0.000000,0.000000}%
\pgfsetstrokecolor{currentstroke}%
\pgfsetdash{}{0pt}%
\pgfsys@defobject{currentmarker}{\pgfqpoint{0.000000in}{0.000000in}}{\pgfqpoint{0.027778in}{0.000000in}}{%
\pgfpathmoveto{\pgfqpoint{0.000000in}{0.000000in}}%
\pgfpathlineto{\pgfqpoint{0.027778in}{0.000000in}}%
\pgfusepath{stroke,fill}%
}%
\begin{pgfscope}%
\pgfsys@transformshift{2.259013in}{3.125849in}%
\pgfsys@useobject{currentmarker}{}%
\end{pgfscope}%
\end{pgfscope}%
\begin{pgfscope}%
\definecolor{textcolor}{rgb}{0.000000,0.000000,0.000000}%
\pgfsetstrokecolor{textcolor}%
\pgfsetfillcolor{textcolor}%
\pgftext[x=2.075368in, y=3.101736in, left, base]{\color{textcolor}{\rmfamily\fontsize{5.000000}{6.000000}\selectfont\catcode`\^=\active\def^{\ifmmode\sp\else\^{}\fi}\catcode`\%=\active\def%{\%}9.98}}%
\end{pgfscope}%
\begin{pgfscope}%
\pgfpathrectangle{\pgfqpoint{2.259013in}{3.125849in}}{\pgfqpoint{0.979200in}{0.326250in}}%
\pgfusepath{clip}%
\pgfsetbuttcap%
\pgfsetroundjoin%
\pgfsetlinewidth{0.401500pt}%
\definecolor{currentstroke}{rgb}{0.690196,0.690196,0.690196}%
\pgfsetstrokecolor{currentstroke}%
\pgfsetstrokeopacity{0.250000}%
\pgfsetdash{{1.480000pt}{0.640000pt}}{0.000000pt}%
\pgfpathmoveto{\pgfqpoint{2.259013in}{3.288974in}}%
\pgfpathlineto{\pgfqpoint{3.238213in}{3.288974in}}%
\pgfusepath{stroke}%
\end{pgfscope}%
\begin{pgfscope}%
\pgfsetbuttcap%
\pgfsetroundjoin%
\definecolor{currentfill}{rgb}{0.000000,0.000000,0.000000}%
\pgfsetfillcolor{currentfill}%
\pgfsetlinewidth{0.803000pt}%
\definecolor{currentstroke}{rgb}{0.000000,0.000000,0.000000}%
\pgfsetstrokecolor{currentstroke}%
\pgfsetdash{}{0pt}%
\pgfsys@defobject{currentmarker}{\pgfqpoint{0.000000in}{0.000000in}}{\pgfqpoint{0.027778in}{0.000000in}}{%
\pgfpathmoveto{\pgfqpoint{0.000000in}{0.000000in}}%
\pgfpathlineto{\pgfqpoint{0.027778in}{0.000000in}}%
\pgfusepath{stroke,fill}%
}%
\begin{pgfscope}%
\pgfsys@transformshift{2.259013in}{3.288974in}%
\pgfsys@useobject{currentmarker}{}%
\end{pgfscope}%
\end{pgfscope}%
\begin{pgfscope}%
\definecolor{textcolor}{rgb}{0.000000,0.000000,0.000000}%
\pgfsetstrokecolor{textcolor}%
\pgfsetfillcolor{textcolor}%
\pgftext[x=2.028107in, y=3.264861in, left, base]{\color{textcolor}{\rmfamily\fontsize{5.000000}{6.000000}\selectfont\catcode`\^=\active\def^{\ifmmode\sp\else\^{}\fi}\catcode`\%=\active\def%{\%}10.04}}%
\end{pgfscope}%
\begin{pgfscope}%
\pgfpathrectangle{\pgfqpoint{2.259013in}{3.125849in}}{\pgfqpoint{0.979200in}{0.326250in}}%
\pgfusepath{clip}%
\pgfsetbuttcap%
\pgfsetroundjoin%
\pgfsetlinewidth{0.401500pt}%
\definecolor{currentstroke}{rgb}{0.690196,0.690196,0.690196}%
\pgfsetstrokecolor{currentstroke}%
\pgfsetstrokeopacity{0.250000}%
\pgfsetdash{{1.480000pt}{0.640000pt}}{0.000000pt}%
\pgfpathmoveto{\pgfqpoint{2.259013in}{3.452099in}}%
\pgfpathlineto{\pgfqpoint{3.238213in}{3.452099in}}%
\pgfusepath{stroke}%
\end{pgfscope}%
\begin{pgfscope}%
\pgfsetbuttcap%
\pgfsetroundjoin%
\definecolor{currentfill}{rgb}{0.000000,0.000000,0.000000}%
\pgfsetfillcolor{currentfill}%
\pgfsetlinewidth{0.803000pt}%
\definecolor{currentstroke}{rgb}{0.000000,0.000000,0.000000}%
\pgfsetstrokecolor{currentstroke}%
\pgfsetdash{}{0pt}%
\pgfsys@defobject{currentmarker}{\pgfqpoint{0.000000in}{0.000000in}}{\pgfqpoint{0.027778in}{0.000000in}}{%
\pgfpathmoveto{\pgfqpoint{0.000000in}{0.000000in}}%
\pgfpathlineto{\pgfqpoint{0.027778in}{0.000000in}}%
\pgfusepath{stroke,fill}%
}%
\begin{pgfscope}%
\pgfsys@transformshift{2.259013in}{3.452099in}%
\pgfsys@useobject{currentmarker}{}%
\end{pgfscope}%
\end{pgfscope}%
\begin{pgfscope}%
\definecolor{textcolor}{rgb}{0.000000,0.000000,0.000000}%
\pgfsetstrokecolor{textcolor}%
\pgfsetfillcolor{textcolor}%
\pgftext[x=2.028107in, y=3.427986in, left, base]{\color{textcolor}{\rmfamily\fontsize{5.000000}{6.000000}\selectfont\catcode`\^=\active\def^{\ifmmode\sp\else\^{}\fi}\catcode`\%=\active\def%{\%}10.11}}%
\end{pgfscope}%
\begin{pgfscope}%
\pgfpathrectangle{\pgfqpoint{2.259013in}{3.125849in}}{\pgfqpoint{0.979200in}{0.326250in}}%
\pgfusepath{clip}%
\pgfsetrectcap%
\pgfsetroundjoin%
\pgfsetlinewidth{1.505625pt}%
\definecolor{currentstroke}{rgb}{0.000000,0.447059,0.698039}%
\pgfsetstrokecolor{currentstroke}%
\pgfsetdash{}{0pt}%
\pgfpathmoveto{\pgfqpoint{2.255351in}{3.318975in}}%
\pgfpathlineto{\pgfqpoint{2.272138in}{3.285909in}}%
\pgfpathlineto{\pgfqpoint{2.284802in}{3.261581in}}%
\pgfpathlineto{\pgfqpoint{2.299490in}{3.236746in}}%
\pgfpathlineto{\pgfqpoint{2.314178in}{3.215544in}}%
\pgfpathlineto{\pgfqpoint{2.328866in}{3.197972in}}%
\pgfpathlineto{\pgfqpoint{2.338658in}{3.188274in}}%
\pgfpathlineto{\pgfqpoint{2.348450in}{3.180188in}}%
\pgfpathlineto{\pgfqpoint{2.358242in}{3.173714in}}%
\pgfpathlineto{\pgfqpoint{2.368034in}{3.168849in}}%
\pgfpathlineto{\pgfqpoint{2.377826in}{3.165595in}}%
\pgfpathlineto{\pgfqpoint{2.387618in}{3.163948in}}%
\pgfpathlineto{\pgfqpoint{2.397410in}{3.163909in}}%
\pgfpathlineto{\pgfqpoint{2.407202in}{3.165476in}}%
\pgfpathlineto{\pgfqpoint{2.416994in}{3.168646in}}%
\pgfpathlineto{\pgfqpoint{2.426786in}{3.173420in}}%
\pgfpathlineto{\pgfqpoint{2.436578in}{3.179795in}}%
\pgfpathlineto{\pgfqpoint{2.446370in}{3.187769in}}%
\pgfpathlineto{\pgfqpoint{2.456162in}{3.197341in}}%
\pgfpathlineto{\pgfqpoint{2.465954in}{3.208509in}}%
\pgfpathlineto{\pgfqpoint{2.480642in}{3.228248in}}%
\pgfpathlineto{\pgfqpoint{2.495330in}{3.251565in}}%
\pgfpathlineto{\pgfqpoint{2.509480in}{3.277406in}}%
\pgfpathlineto{\pgfqpoint{2.535850in}{3.328210in}}%
\pgfpathlineto{\pgfqpoint{2.550538in}{3.351827in}}%
\pgfpathlineto{\pgfqpoint{2.565226in}{3.371709in}}%
\pgfpathlineto{\pgfqpoint{2.575018in}{3.382890in}}%
\pgfpathlineto{\pgfqpoint{2.584810in}{3.392412in}}%
\pgfpathlineto{\pgfqpoint{2.594602in}{3.400278in}}%
\pgfpathlineto{\pgfqpoint{2.604394in}{3.406488in}}%
\pgfpathlineto{\pgfqpoint{2.614186in}{3.411042in}}%
\pgfpathlineto{\pgfqpoint{2.623978in}{3.413943in}}%
\pgfpathlineto{\pgfqpoint{2.633770in}{3.415192in}}%
\pgfpathlineto{\pgfqpoint{2.643562in}{3.414791in}}%
\pgfpathlineto{\pgfqpoint{2.653354in}{3.412740in}}%
\pgfpathlineto{\pgfqpoint{2.663146in}{3.409043in}}%
\pgfpathlineto{\pgfqpoint{2.672938in}{3.403700in}}%
\pgfpathlineto{\pgfqpoint{2.682730in}{3.396714in}}%
\pgfpathlineto{\pgfqpoint{2.692522in}{3.388088in}}%
\pgfpathlineto{\pgfqpoint{2.702314in}{3.377824in}}%
\pgfpathlineto{\pgfqpoint{2.712106in}{3.365924in}}%
\pgfpathlineto{\pgfqpoint{2.726794in}{3.345014in}}%
\pgfpathlineto{\pgfqpoint{2.741482in}{3.320443in}}%
\pgfpathlineto{\pgfqpoint{2.756538in}{3.291618in}}%
\pgfpathlineto{\pgfqpoint{2.773400in}{3.258393in}}%
\pgfpathlineto{\pgfqpoint{2.788088in}{3.233228in}}%
\pgfpathlineto{\pgfqpoint{2.802776in}{3.211697in}}%
\pgfpathlineto{\pgfqpoint{2.817464in}{3.193798in}}%
\pgfpathlineto{\pgfqpoint{2.827256in}{3.183881in}}%
\pgfpathlineto{\pgfqpoint{2.837048in}{3.175578in}}%
\pgfpathlineto{\pgfqpoint{2.846840in}{3.168886in}}%
\pgfpathlineto{\pgfqpoint{2.856632in}{3.163806in}}%
\pgfpathlineto{\pgfqpoint{2.866424in}{3.160335in}}%
\pgfpathlineto{\pgfqpoint{2.876216in}{3.158473in}}%
\pgfpathlineto{\pgfqpoint{2.886008in}{3.158219in}}%
\pgfpathlineto{\pgfqpoint{2.895800in}{3.159571in}}%
\pgfpathlineto{\pgfqpoint{2.905592in}{3.162528in}}%
\pgfpathlineto{\pgfqpoint{2.915384in}{3.167088in}}%
\pgfpathlineto{\pgfqpoint{2.925176in}{3.173250in}}%
\pgfpathlineto{\pgfqpoint{2.934968in}{3.181012in}}%
\pgfpathlineto{\pgfqpoint{2.944760in}{3.190372in}}%
\pgfpathlineto{\pgfqpoint{2.954552in}{3.201329in}}%
\pgfpathlineto{\pgfqpoint{2.969240in}{3.220753in}}%
\pgfpathlineto{\pgfqpoint{2.983928in}{3.243756in}}%
\pgfpathlineto{\pgfqpoint{2.998517in}{3.270140in}}%
\pgfpathlineto{\pgfqpoint{3.025257in}{3.321511in}}%
\pgfpathlineto{\pgfqpoint{3.039945in}{3.345056in}}%
\pgfpathlineto{\pgfqpoint{3.054633in}{3.364867in}}%
\pgfpathlineto{\pgfqpoint{3.064425in}{3.376002in}}%
\pgfpathlineto{\pgfqpoint{3.074217in}{3.385479in}}%
\pgfpathlineto{\pgfqpoint{3.084009in}{3.393299in}}%
\pgfpathlineto{\pgfqpoint{3.093801in}{3.399463in}}%
\pgfpathlineto{\pgfqpoint{3.103593in}{3.403974in}}%
\pgfpathlineto{\pgfqpoint{3.113385in}{3.406831in}}%
\pgfpathlineto{\pgfqpoint{3.123177in}{3.408036in}}%
\pgfpathlineto{\pgfqpoint{3.132969in}{3.407592in}}%
\pgfpathlineto{\pgfqpoint{3.142761in}{3.405499in}}%
\pgfpathlineto{\pgfqpoint{3.152553in}{3.401759in}}%
\pgfpathlineto{\pgfqpoint{3.162345in}{3.396375in}}%
\pgfpathlineto{\pgfqpoint{3.172137in}{3.389348in}}%
\pgfpathlineto{\pgfqpoint{3.181929in}{3.380682in}}%
\pgfpathlineto{\pgfqpoint{3.191721in}{3.370378in}}%
\pgfpathlineto{\pgfqpoint{3.201513in}{3.358439in}}%
\pgfpathlineto{\pgfqpoint{3.216201in}{3.337471in}}%
\pgfpathlineto{\pgfqpoint{3.230889in}{3.312855in}}%
\pgfpathlineto{\pgfqpoint{3.238213in}{3.299291in}}%
\pgfpathlineto{\pgfqpoint{3.238213in}{3.299291in}}%
\pgfusepath{stroke}%
\end{pgfscope}%
\begin{pgfscope}%
\pgfsetrectcap%
\pgfsetmiterjoin%
\pgfsetlinewidth{0.803000pt}%
\definecolor{currentstroke}{rgb}{0.000000,0.000000,0.000000}%
\pgfsetstrokecolor{currentstroke}%
\pgfsetdash{}{0pt}%
\pgfpathmoveto{\pgfqpoint{2.259013in}{3.125849in}}%
\pgfpathlineto{\pgfqpoint{2.259013in}{3.452099in}}%
\pgfusepath{stroke}%
\end{pgfscope}%
\begin{pgfscope}%
\pgfsetrectcap%
\pgfsetmiterjoin%
\pgfsetlinewidth{0.803000pt}%
\definecolor{currentstroke}{rgb}{0.000000,0.000000,0.000000}%
\pgfsetstrokecolor{currentstroke}%
\pgfsetdash{}{0pt}%
\pgfpathmoveto{\pgfqpoint{3.238213in}{3.125849in}}%
\pgfpathlineto{\pgfqpoint{3.238213in}{3.452099in}}%
\pgfusepath{stroke}%
\end{pgfscope}%
\begin{pgfscope}%
\pgfsetrectcap%
\pgfsetmiterjoin%
\pgfsetlinewidth{0.803000pt}%
\definecolor{currentstroke}{rgb}{0.000000,0.000000,0.000000}%
\pgfsetstrokecolor{currentstroke}%
\pgfsetdash{}{0pt}%
\pgfpathmoveto{\pgfqpoint{2.259013in}{3.125849in}}%
\pgfpathlineto{\pgfqpoint{3.238213in}{3.125849in}}%
\pgfusepath{stroke}%
\end{pgfscope}%
\begin{pgfscope}%
\pgfsetrectcap%
\pgfsetmiterjoin%
\pgfsetlinewidth{0.803000pt}%
\definecolor{currentstroke}{rgb}{0.000000,0.000000,0.000000}%
\pgfsetstrokecolor{currentstroke}%
\pgfsetdash{}{0pt}%
\pgfpathmoveto{\pgfqpoint{2.259013in}{3.452099in}}%
\pgfpathlineto{\pgfqpoint{3.238213in}{3.452099in}}%
\pgfusepath{stroke}%
\end{pgfscope}%
\begin{pgfscope}%
\pgfsetbuttcap%
\pgfsetmiterjoin%
\definecolor{currentfill}{rgb}{1.000000,1.000000,1.000000}%
\pgfsetfillcolor{currentfill}%
\pgfsetlinewidth{0.000000pt}%
\definecolor{currentstroke}{rgb}{0.000000,0.000000,0.000000}%
\pgfsetstrokecolor{currentstroke}%
\pgfsetstrokeopacity{0.000000}%
\pgfsetdash{}{0pt}%
\pgfpathmoveto{\pgfqpoint{2.259013in}{2.056474in}}%
\pgfpathlineto{\pgfqpoint{3.238213in}{2.056474in}}%
\pgfpathlineto{\pgfqpoint{3.238213in}{2.382724in}}%
\pgfpathlineto{\pgfqpoint{2.259013in}{2.382724in}}%
\pgfpathlineto{\pgfqpoint{2.259013in}{2.056474in}}%
\pgfpathclose%
\pgfusepath{fill}%
\end{pgfscope}%
\begin{pgfscope}%
\pgfpathrectangle{\pgfqpoint{2.259013in}{2.056474in}}{\pgfqpoint{0.979200in}{0.326250in}}%
\pgfusepath{clip}%
\pgfsetbuttcap%
\pgfsetroundjoin%
\pgfsetlinewidth{0.401500pt}%
\definecolor{currentstroke}{rgb}{0.690196,0.690196,0.690196}%
\pgfsetstrokecolor{currentstroke}%
\pgfsetstrokeopacity{0.250000}%
\pgfsetdash{{1.480000pt}{0.640000pt}}{0.000000pt}%
\pgfpathmoveto{\pgfqpoint{2.259013in}{2.056474in}}%
\pgfpathlineto{\pgfqpoint{2.259013in}{2.382724in}}%
\pgfusepath{stroke}%
\end{pgfscope}%
\begin{pgfscope}%
\pgfsetbuttcap%
\pgfsetroundjoin%
\definecolor{currentfill}{rgb}{0.000000,0.000000,0.000000}%
\pgfsetfillcolor{currentfill}%
\pgfsetlinewidth{0.803000pt}%
\definecolor{currentstroke}{rgb}{0.000000,0.000000,0.000000}%
\pgfsetstrokecolor{currentstroke}%
\pgfsetdash{}{0pt}%
\pgfsys@defobject{currentmarker}{\pgfqpoint{0.000000in}{0.000000in}}{\pgfqpoint{0.000000in}{0.027778in}}{%
\pgfpathmoveto{\pgfqpoint{0.000000in}{0.000000in}}%
\pgfpathlineto{\pgfqpoint{0.000000in}{0.027778in}}%
\pgfusepath{stroke,fill}%
}%
\begin{pgfscope}%
\pgfsys@transformshift{2.259013in}{2.056474in}%
\pgfsys@useobject{currentmarker}{}%
\end{pgfscope}%
\end{pgfscope}%
\begin{pgfscope}%
\definecolor{textcolor}{rgb}{0.000000,0.000000,0.000000}%
\pgfsetstrokecolor{textcolor}%
\pgfsetfillcolor{textcolor}%
\pgftext[x=2.259013in,y=2.042585in,,top]{\color{textcolor}{\rmfamily\fontsize{5.000000}{6.000000}\selectfont\catcode`\^=\active\def^{\ifmmode\sp\else\^{}\fi}\catcode`\%=\active\def%{\%}5.9}}%
\end{pgfscope}%
\begin{pgfscope}%
\pgfpathrectangle{\pgfqpoint{2.259013in}{2.056474in}}{\pgfqpoint{0.979200in}{0.326250in}}%
\pgfusepath{clip}%
\pgfsetbuttcap%
\pgfsetroundjoin%
\pgfsetlinewidth{0.401500pt}%
\definecolor{currentstroke}{rgb}{0.690196,0.690196,0.690196}%
\pgfsetstrokecolor{currentstroke}%
\pgfsetstrokeopacity{0.250000}%
\pgfsetdash{{1.480000pt}{0.640000pt}}{0.000000pt}%
\pgfpathmoveto{\pgfqpoint{2.503813in}{2.056474in}}%
\pgfpathlineto{\pgfqpoint{2.503813in}{2.382724in}}%
\pgfusepath{stroke}%
\end{pgfscope}%
\begin{pgfscope}%
\pgfsetbuttcap%
\pgfsetroundjoin%
\definecolor{currentfill}{rgb}{0.000000,0.000000,0.000000}%
\pgfsetfillcolor{currentfill}%
\pgfsetlinewidth{0.803000pt}%
\definecolor{currentstroke}{rgb}{0.000000,0.000000,0.000000}%
\pgfsetstrokecolor{currentstroke}%
\pgfsetdash{}{0pt}%
\pgfsys@defobject{currentmarker}{\pgfqpoint{0.000000in}{0.000000in}}{\pgfqpoint{0.000000in}{0.027778in}}{%
\pgfpathmoveto{\pgfqpoint{0.000000in}{0.000000in}}%
\pgfpathlineto{\pgfqpoint{0.000000in}{0.027778in}}%
\pgfusepath{stroke,fill}%
}%
\begin{pgfscope}%
\pgfsys@transformshift{2.503813in}{2.056474in}%
\pgfsys@useobject{currentmarker}{}%
\end{pgfscope}%
\end{pgfscope}%
\begin{pgfscope}%
\definecolor{textcolor}{rgb}{0.000000,0.000000,0.000000}%
\pgfsetstrokecolor{textcolor}%
\pgfsetfillcolor{textcolor}%
\pgftext[x=2.503813in,y=2.042585in,,top]{\color{textcolor}{\rmfamily\fontsize{5.000000}{6.000000}\selectfont\catcode`\^=\active\def^{\ifmmode\sp\else\^{}\fi}\catcode`\%=\active\def%{\%}5.9}}%
\end{pgfscope}%
\begin{pgfscope}%
\pgfpathrectangle{\pgfqpoint{2.259013in}{2.056474in}}{\pgfqpoint{0.979200in}{0.326250in}}%
\pgfusepath{clip}%
\pgfsetbuttcap%
\pgfsetroundjoin%
\pgfsetlinewidth{0.401500pt}%
\definecolor{currentstroke}{rgb}{0.690196,0.690196,0.690196}%
\pgfsetstrokecolor{currentstroke}%
\pgfsetstrokeopacity{0.250000}%
\pgfsetdash{{1.480000pt}{0.640000pt}}{0.000000pt}%
\pgfpathmoveto{\pgfqpoint{2.748613in}{2.056474in}}%
\pgfpathlineto{\pgfqpoint{2.748613in}{2.382724in}}%
\pgfusepath{stroke}%
\end{pgfscope}%
\begin{pgfscope}%
\pgfsetbuttcap%
\pgfsetroundjoin%
\definecolor{currentfill}{rgb}{0.000000,0.000000,0.000000}%
\pgfsetfillcolor{currentfill}%
\pgfsetlinewidth{0.803000pt}%
\definecolor{currentstroke}{rgb}{0.000000,0.000000,0.000000}%
\pgfsetstrokecolor{currentstroke}%
\pgfsetdash{}{0pt}%
\pgfsys@defobject{currentmarker}{\pgfqpoint{0.000000in}{0.000000in}}{\pgfqpoint{0.000000in}{0.027778in}}{%
\pgfpathmoveto{\pgfqpoint{0.000000in}{0.000000in}}%
\pgfpathlineto{\pgfqpoint{0.000000in}{0.027778in}}%
\pgfusepath{stroke,fill}%
}%
\begin{pgfscope}%
\pgfsys@transformshift{2.748613in}{2.056474in}%
\pgfsys@useobject{currentmarker}{}%
\end{pgfscope}%
\end{pgfscope}%
\begin{pgfscope}%
\definecolor{textcolor}{rgb}{0.000000,0.000000,0.000000}%
\pgfsetstrokecolor{textcolor}%
\pgfsetfillcolor{textcolor}%
\pgftext[x=2.748613in,y=2.042585in,,top]{\color{textcolor}{\rmfamily\fontsize{5.000000}{6.000000}\selectfont\catcode`\^=\active\def^{\ifmmode\sp\else\^{}\fi}\catcode`\%=\active\def%{\%}6.0}}%
\end{pgfscope}%
\begin{pgfscope}%
\pgfpathrectangle{\pgfqpoint{2.259013in}{2.056474in}}{\pgfqpoint{0.979200in}{0.326250in}}%
\pgfusepath{clip}%
\pgfsetbuttcap%
\pgfsetroundjoin%
\pgfsetlinewidth{0.401500pt}%
\definecolor{currentstroke}{rgb}{0.690196,0.690196,0.690196}%
\pgfsetstrokecolor{currentstroke}%
\pgfsetstrokeopacity{0.250000}%
\pgfsetdash{{1.480000pt}{0.640000pt}}{0.000000pt}%
\pgfpathmoveto{\pgfqpoint{2.993413in}{2.056474in}}%
\pgfpathlineto{\pgfqpoint{2.993413in}{2.382724in}}%
\pgfusepath{stroke}%
\end{pgfscope}%
\begin{pgfscope}%
\pgfsetbuttcap%
\pgfsetroundjoin%
\definecolor{currentfill}{rgb}{0.000000,0.000000,0.000000}%
\pgfsetfillcolor{currentfill}%
\pgfsetlinewidth{0.803000pt}%
\definecolor{currentstroke}{rgb}{0.000000,0.000000,0.000000}%
\pgfsetstrokecolor{currentstroke}%
\pgfsetdash{}{0pt}%
\pgfsys@defobject{currentmarker}{\pgfqpoint{0.000000in}{0.000000in}}{\pgfqpoint{0.000000in}{0.027778in}}{%
\pgfpathmoveto{\pgfqpoint{0.000000in}{0.000000in}}%
\pgfpathlineto{\pgfqpoint{0.000000in}{0.027778in}}%
\pgfusepath{stroke,fill}%
}%
\begin{pgfscope}%
\pgfsys@transformshift{2.993413in}{2.056474in}%
\pgfsys@useobject{currentmarker}{}%
\end{pgfscope}%
\end{pgfscope}%
\begin{pgfscope}%
\definecolor{textcolor}{rgb}{0.000000,0.000000,0.000000}%
\pgfsetstrokecolor{textcolor}%
\pgfsetfillcolor{textcolor}%
\pgftext[x=2.993413in,y=2.042585in,,top]{\color{textcolor}{\rmfamily\fontsize{5.000000}{6.000000}\selectfont\catcode`\^=\active\def^{\ifmmode\sp\else\^{}\fi}\catcode`\%=\active\def%{\%}6.0}}%
\end{pgfscope}%
\begin{pgfscope}%
\pgfpathrectangle{\pgfqpoint{2.259013in}{2.056474in}}{\pgfqpoint{0.979200in}{0.326250in}}%
\pgfusepath{clip}%
\pgfsetbuttcap%
\pgfsetroundjoin%
\pgfsetlinewidth{0.401500pt}%
\definecolor{currentstroke}{rgb}{0.690196,0.690196,0.690196}%
\pgfsetstrokecolor{currentstroke}%
\pgfsetstrokeopacity{0.250000}%
\pgfsetdash{{1.480000pt}{0.640000pt}}{0.000000pt}%
\pgfpathmoveto{\pgfqpoint{3.238213in}{2.056474in}}%
\pgfpathlineto{\pgfqpoint{3.238213in}{2.382724in}}%
\pgfusepath{stroke}%
\end{pgfscope}%
\begin{pgfscope}%
\pgfsetbuttcap%
\pgfsetroundjoin%
\definecolor{currentfill}{rgb}{0.000000,0.000000,0.000000}%
\pgfsetfillcolor{currentfill}%
\pgfsetlinewidth{0.803000pt}%
\definecolor{currentstroke}{rgb}{0.000000,0.000000,0.000000}%
\pgfsetstrokecolor{currentstroke}%
\pgfsetdash{}{0pt}%
\pgfsys@defobject{currentmarker}{\pgfqpoint{0.000000in}{0.000000in}}{\pgfqpoint{0.000000in}{0.027778in}}{%
\pgfpathmoveto{\pgfqpoint{0.000000in}{0.000000in}}%
\pgfpathlineto{\pgfqpoint{0.000000in}{0.027778in}}%
\pgfusepath{stroke,fill}%
}%
\begin{pgfscope}%
\pgfsys@transformshift{3.238213in}{2.056474in}%
\pgfsys@useobject{currentmarker}{}%
\end{pgfscope}%
\end{pgfscope}%
\begin{pgfscope}%
\definecolor{textcolor}{rgb}{0.000000,0.000000,0.000000}%
\pgfsetstrokecolor{textcolor}%
\pgfsetfillcolor{textcolor}%
\pgftext[x=3.238213in,y=2.042585in,,top]{\color{textcolor}{\rmfamily\fontsize{5.000000}{6.000000}\selectfont\catcode`\^=\active\def^{\ifmmode\sp\else\^{}\fi}\catcode`\%=\active\def%{\%}6.0}}%
\end{pgfscope}%
\begin{pgfscope}%
\pgfpathrectangle{\pgfqpoint{2.259013in}{2.056474in}}{\pgfqpoint{0.979200in}{0.326250in}}%
\pgfusepath{clip}%
\pgfsetbuttcap%
\pgfsetroundjoin%
\pgfsetlinewidth{0.401500pt}%
\definecolor{currentstroke}{rgb}{0.690196,0.690196,0.690196}%
\pgfsetstrokecolor{currentstroke}%
\pgfsetstrokeopacity{0.250000}%
\pgfsetdash{{1.480000pt}{0.640000pt}}{0.000000pt}%
\pgfpathmoveto{\pgfqpoint{2.259013in}{2.056474in}}%
\pgfpathlineto{\pgfqpoint{3.238213in}{2.056474in}}%
\pgfusepath{stroke}%
\end{pgfscope}%
\begin{pgfscope}%
\pgfsetbuttcap%
\pgfsetroundjoin%
\definecolor{currentfill}{rgb}{0.000000,0.000000,0.000000}%
\pgfsetfillcolor{currentfill}%
\pgfsetlinewidth{0.803000pt}%
\definecolor{currentstroke}{rgb}{0.000000,0.000000,0.000000}%
\pgfsetstrokecolor{currentstroke}%
\pgfsetdash{}{0pt}%
\pgfsys@defobject{currentmarker}{\pgfqpoint{0.000000in}{0.000000in}}{\pgfqpoint{0.027778in}{0.000000in}}{%
\pgfpathmoveto{\pgfqpoint{0.000000in}{0.000000in}}%
\pgfpathlineto{\pgfqpoint{0.027778in}{0.000000in}}%
\pgfusepath{stroke,fill}%
}%
\begin{pgfscope}%
\pgfsys@transformshift{2.259013in}{2.056474in}%
\pgfsys@useobject{currentmarker}{}%
\end{pgfscope}%
\end{pgfscope}%
\begin{pgfscope}%
\definecolor{textcolor}{rgb}{0.000000,0.000000,0.000000}%
\pgfsetstrokecolor{textcolor}%
\pgfsetfillcolor{textcolor}%
\pgftext[x=2.075368in, y=2.032361in, left, base]{\color{textcolor}{\rmfamily\fontsize{5.000000}{6.000000}\selectfont\catcode`\^=\active\def^{\ifmmode\sp\else\^{}\fi}\catcode`\%=\active\def%{\%}0.00}}%
\end{pgfscope}%
\begin{pgfscope}%
\pgfpathrectangle{\pgfqpoint{2.259013in}{2.056474in}}{\pgfqpoint{0.979200in}{0.326250in}}%
\pgfusepath{clip}%
\pgfsetbuttcap%
\pgfsetroundjoin%
\pgfsetlinewidth{0.401500pt}%
\definecolor{currentstroke}{rgb}{0.690196,0.690196,0.690196}%
\pgfsetstrokecolor{currentstroke}%
\pgfsetstrokeopacity{0.250000}%
\pgfsetdash{{1.480000pt}{0.640000pt}}{0.000000pt}%
\pgfpathmoveto{\pgfqpoint{2.259013in}{2.219599in}}%
\pgfpathlineto{\pgfqpoint{3.238213in}{2.219599in}}%
\pgfusepath{stroke}%
\end{pgfscope}%
\begin{pgfscope}%
\pgfsetbuttcap%
\pgfsetroundjoin%
\definecolor{currentfill}{rgb}{0.000000,0.000000,0.000000}%
\pgfsetfillcolor{currentfill}%
\pgfsetlinewidth{0.803000pt}%
\definecolor{currentstroke}{rgb}{0.000000,0.000000,0.000000}%
\pgfsetstrokecolor{currentstroke}%
\pgfsetdash{}{0pt}%
\pgfsys@defobject{currentmarker}{\pgfqpoint{0.000000in}{0.000000in}}{\pgfqpoint{0.027778in}{0.000000in}}{%
\pgfpathmoveto{\pgfqpoint{0.000000in}{0.000000in}}%
\pgfpathlineto{\pgfqpoint{0.027778in}{0.000000in}}%
\pgfusepath{stroke,fill}%
}%
\begin{pgfscope}%
\pgfsys@transformshift{2.259013in}{2.219599in}%
\pgfsys@useobject{currentmarker}{}%
\end{pgfscope}%
\end{pgfscope}%
\begin{pgfscope}%
\definecolor{textcolor}{rgb}{0.000000,0.000000,0.000000}%
\pgfsetstrokecolor{textcolor}%
\pgfsetfillcolor{textcolor}%
\pgftext[x=2.075368in, y=2.195486in, left, base]{\color{textcolor}{\rmfamily\fontsize{5.000000}{6.000000}\selectfont\catcode`\^=\active\def^{\ifmmode\sp\else\^{}\fi}\catcode`\%=\active\def%{\%}1.05}}%
\end{pgfscope}%
\begin{pgfscope}%
\pgfpathrectangle{\pgfqpoint{2.259013in}{2.056474in}}{\pgfqpoint{0.979200in}{0.326250in}}%
\pgfusepath{clip}%
\pgfsetbuttcap%
\pgfsetroundjoin%
\pgfsetlinewidth{0.401500pt}%
\definecolor{currentstroke}{rgb}{0.690196,0.690196,0.690196}%
\pgfsetstrokecolor{currentstroke}%
\pgfsetstrokeopacity{0.250000}%
\pgfsetdash{{1.480000pt}{0.640000pt}}{0.000000pt}%
\pgfpathmoveto{\pgfqpoint{2.259013in}{2.382724in}}%
\pgfpathlineto{\pgfqpoint{3.238213in}{2.382724in}}%
\pgfusepath{stroke}%
\end{pgfscope}%
\begin{pgfscope}%
\pgfsetbuttcap%
\pgfsetroundjoin%
\definecolor{currentfill}{rgb}{0.000000,0.000000,0.000000}%
\pgfsetfillcolor{currentfill}%
\pgfsetlinewidth{0.803000pt}%
\definecolor{currentstroke}{rgb}{0.000000,0.000000,0.000000}%
\pgfsetstrokecolor{currentstroke}%
\pgfsetdash{}{0pt}%
\pgfsys@defobject{currentmarker}{\pgfqpoint{0.000000in}{0.000000in}}{\pgfqpoint{0.027778in}{0.000000in}}{%
\pgfpathmoveto{\pgfqpoint{0.000000in}{0.000000in}}%
\pgfpathlineto{\pgfqpoint{0.027778in}{0.000000in}}%
\pgfusepath{stroke,fill}%
}%
\begin{pgfscope}%
\pgfsys@transformshift{2.259013in}{2.382724in}%
\pgfsys@useobject{currentmarker}{}%
\end{pgfscope}%
\end{pgfscope}%
\begin{pgfscope}%
\definecolor{textcolor}{rgb}{0.000000,0.000000,0.000000}%
\pgfsetstrokecolor{textcolor}%
\pgfsetfillcolor{textcolor}%
\pgftext[x=2.075368in, y=2.358611in, left, base]{\color{textcolor}{\rmfamily\fontsize{5.000000}{6.000000}\selectfont\catcode`\^=\active\def^{\ifmmode\sp\else\^{}\fi}\catcode`\%=\active\def%{\%}2.10}}%
\end{pgfscope}%
\begin{pgfscope}%
\pgfpathrectangle{\pgfqpoint{2.259013in}{2.056474in}}{\pgfqpoint{0.979200in}{0.326250in}}%
\pgfusepath{clip}%
\pgfsetrectcap%
\pgfsetroundjoin%
\pgfsetlinewidth{1.505625pt}%
\definecolor{currentstroke}{rgb}{0.835294,0.368627,0.000000}%
\pgfsetstrokecolor{currentstroke}%
\pgfsetdash{}{0pt}%
\pgfpathmoveto{\pgfqpoint{2.255351in}{2.071988in}}%
\pgfpathlineto{\pgfqpoint{2.261522in}{2.065765in}}%
\pgfpathlineto{\pgfqpoint{2.269179in}{2.056818in}}%
\pgfpathlineto{\pgfqpoint{2.271233in}{2.058163in}}%
\pgfpathlineto{\pgfqpoint{2.275614in}{2.063257in}}%
\pgfpathlineto{\pgfqpoint{2.516943in}{2.367298in}}%
\pgfpathlineto{\pgfqpoint{2.519180in}{2.365931in}}%
\pgfpathlineto{\pgfqpoint{2.756538in}{2.058315in}}%
\pgfpathlineto{\pgfqpoint{2.758982in}{2.056370in}}%
\pgfpathlineto{\pgfqpoint{2.761430in}{2.058277in}}%
\pgfpathlineto{\pgfqpoint{2.768865in}{2.067415in}}%
\pgfpathlineto{\pgfqpoint{3.006364in}{2.366827in}}%
\pgfpathlineto{\pgfqpoint{3.008590in}{2.365507in}}%
\pgfpathlineto{\pgfqpoint{3.234551in}{2.072021in}}%
\pgfpathlineto{\pgfqpoint{3.238213in}{2.068674in}}%
\pgfpathlineto{\pgfqpoint{3.238213in}{2.068674in}}%
\pgfusepath{stroke}%
\end{pgfscope}%
\begin{pgfscope}%
\pgfsetrectcap%
\pgfsetmiterjoin%
\pgfsetlinewidth{0.803000pt}%
\definecolor{currentstroke}{rgb}{0.000000,0.000000,0.000000}%
\pgfsetstrokecolor{currentstroke}%
\pgfsetdash{}{0pt}%
\pgfpathmoveto{\pgfqpoint{2.259013in}{2.056474in}}%
\pgfpathlineto{\pgfqpoint{2.259013in}{2.382724in}}%
\pgfusepath{stroke}%
\end{pgfscope}%
\begin{pgfscope}%
\pgfsetrectcap%
\pgfsetmiterjoin%
\pgfsetlinewidth{0.803000pt}%
\definecolor{currentstroke}{rgb}{0.000000,0.000000,0.000000}%
\pgfsetstrokecolor{currentstroke}%
\pgfsetdash{}{0pt}%
\pgfpathmoveto{\pgfqpoint{3.238213in}{2.056474in}}%
\pgfpathlineto{\pgfqpoint{3.238213in}{2.382724in}}%
\pgfusepath{stroke}%
\end{pgfscope}%
\begin{pgfscope}%
\pgfsetrectcap%
\pgfsetmiterjoin%
\pgfsetlinewidth{0.803000pt}%
\definecolor{currentstroke}{rgb}{0.000000,0.000000,0.000000}%
\pgfsetstrokecolor{currentstroke}%
\pgfsetdash{}{0pt}%
\pgfpathmoveto{\pgfqpoint{2.259013in}{2.056474in}}%
\pgfpathlineto{\pgfqpoint{3.238213in}{2.056474in}}%
\pgfusepath{stroke}%
\end{pgfscope}%
\begin{pgfscope}%
\pgfsetrectcap%
\pgfsetmiterjoin%
\pgfsetlinewidth{0.803000pt}%
\definecolor{currentstroke}{rgb}{0.000000,0.000000,0.000000}%
\pgfsetstrokecolor{currentstroke}%
\pgfsetdash{}{0pt}%
\pgfpathmoveto{\pgfqpoint{2.259013in}{2.382724in}}%
\pgfpathlineto{\pgfqpoint{3.238213in}{2.382724in}}%
\pgfusepath{stroke}%
\end{pgfscope}%
\begin{pgfscope}%
\pgfsetbuttcap%
\pgfsetmiterjoin%
\definecolor{currentfill}{rgb}{1.000000,1.000000,1.000000}%
\pgfsetfillcolor{currentfill}%
\pgfsetlinewidth{0.000000pt}%
\definecolor{currentstroke}{rgb}{0.000000,0.000000,0.000000}%
\pgfsetstrokecolor{currentstroke}%
\pgfsetstrokeopacity{0.000000}%
\pgfsetdash{}{0pt}%
\pgfpathmoveto{\pgfqpoint{2.259013in}{0.987099in}}%
\pgfpathlineto{\pgfqpoint{3.238213in}{0.987099in}}%
\pgfpathlineto{\pgfqpoint{3.238213in}{1.313349in}}%
\pgfpathlineto{\pgfqpoint{2.259013in}{1.313349in}}%
\pgfpathlineto{\pgfqpoint{2.259013in}{0.987099in}}%
\pgfpathclose%
\pgfusepath{fill}%
\end{pgfscope}%
\begin{pgfscope}%
\pgfpathrectangle{\pgfqpoint{2.259013in}{0.987099in}}{\pgfqpoint{0.979200in}{0.326250in}}%
\pgfusepath{clip}%
\pgfsetbuttcap%
\pgfsetroundjoin%
\pgfsetlinewidth{0.401500pt}%
\definecolor{currentstroke}{rgb}{0.690196,0.690196,0.690196}%
\pgfsetstrokecolor{currentstroke}%
\pgfsetstrokeopacity{0.250000}%
\pgfsetdash{{1.480000pt}{0.640000pt}}{0.000000pt}%
\pgfpathmoveto{\pgfqpoint{2.259013in}{0.987099in}}%
\pgfpathlineto{\pgfqpoint{2.259013in}{1.313349in}}%
\pgfusepath{stroke}%
\end{pgfscope}%
\begin{pgfscope}%
\pgfsetbuttcap%
\pgfsetroundjoin%
\definecolor{currentfill}{rgb}{0.000000,0.000000,0.000000}%
\pgfsetfillcolor{currentfill}%
\pgfsetlinewidth{0.803000pt}%
\definecolor{currentstroke}{rgb}{0.000000,0.000000,0.000000}%
\pgfsetstrokecolor{currentstroke}%
\pgfsetdash{}{0pt}%
\pgfsys@defobject{currentmarker}{\pgfqpoint{0.000000in}{0.000000in}}{\pgfqpoint{0.000000in}{0.027778in}}{%
\pgfpathmoveto{\pgfqpoint{0.000000in}{0.000000in}}%
\pgfpathlineto{\pgfqpoint{0.000000in}{0.027778in}}%
\pgfusepath{stroke,fill}%
}%
\begin{pgfscope}%
\pgfsys@transformshift{2.259013in}{0.987099in}%
\pgfsys@useobject{currentmarker}{}%
\end{pgfscope}%
\end{pgfscope}%
\begin{pgfscope}%
\definecolor{textcolor}{rgb}{0.000000,0.000000,0.000000}%
\pgfsetstrokecolor{textcolor}%
\pgfsetfillcolor{textcolor}%
\pgftext[x=2.259013in,y=0.973210in,,top]{\color{textcolor}{\rmfamily\fontsize{5.000000}{6.000000}\selectfont\catcode`\^=\active\def^{\ifmmode\sp\else\^{}\fi}\catcode`\%=\active\def%{\%}5.9}}%
\end{pgfscope}%
\begin{pgfscope}%
\pgfpathrectangle{\pgfqpoint{2.259013in}{0.987099in}}{\pgfqpoint{0.979200in}{0.326250in}}%
\pgfusepath{clip}%
\pgfsetbuttcap%
\pgfsetroundjoin%
\pgfsetlinewidth{0.401500pt}%
\definecolor{currentstroke}{rgb}{0.690196,0.690196,0.690196}%
\pgfsetstrokecolor{currentstroke}%
\pgfsetstrokeopacity{0.250000}%
\pgfsetdash{{1.480000pt}{0.640000pt}}{0.000000pt}%
\pgfpathmoveto{\pgfqpoint{2.503813in}{0.987099in}}%
\pgfpathlineto{\pgfqpoint{2.503813in}{1.313349in}}%
\pgfusepath{stroke}%
\end{pgfscope}%
\begin{pgfscope}%
\pgfsetbuttcap%
\pgfsetroundjoin%
\definecolor{currentfill}{rgb}{0.000000,0.000000,0.000000}%
\pgfsetfillcolor{currentfill}%
\pgfsetlinewidth{0.803000pt}%
\definecolor{currentstroke}{rgb}{0.000000,0.000000,0.000000}%
\pgfsetstrokecolor{currentstroke}%
\pgfsetdash{}{0pt}%
\pgfsys@defobject{currentmarker}{\pgfqpoint{0.000000in}{0.000000in}}{\pgfqpoint{0.000000in}{0.027778in}}{%
\pgfpathmoveto{\pgfqpoint{0.000000in}{0.000000in}}%
\pgfpathlineto{\pgfqpoint{0.000000in}{0.027778in}}%
\pgfusepath{stroke,fill}%
}%
\begin{pgfscope}%
\pgfsys@transformshift{2.503813in}{0.987099in}%
\pgfsys@useobject{currentmarker}{}%
\end{pgfscope}%
\end{pgfscope}%
\begin{pgfscope}%
\definecolor{textcolor}{rgb}{0.000000,0.000000,0.000000}%
\pgfsetstrokecolor{textcolor}%
\pgfsetfillcolor{textcolor}%
\pgftext[x=2.503813in,y=0.973210in,,top]{\color{textcolor}{\rmfamily\fontsize{5.000000}{6.000000}\selectfont\catcode`\^=\active\def^{\ifmmode\sp\else\^{}\fi}\catcode`\%=\active\def%{\%}5.9}}%
\end{pgfscope}%
\begin{pgfscope}%
\pgfpathrectangle{\pgfqpoint{2.259013in}{0.987099in}}{\pgfqpoint{0.979200in}{0.326250in}}%
\pgfusepath{clip}%
\pgfsetbuttcap%
\pgfsetroundjoin%
\pgfsetlinewidth{0.401500pt}%
\definecolor{currentstroke}{rgb}{0.690196,0.690196,0.690196}%
\pgfsetstrokecolor{currentstroke}%
\pgfsetstrokeopacity{0.250000}%
\pgfsetdash{{1.480000pt}{0.640000pt}}{0.000000pt}%
\pgfpathmoveto{\pgfqpoint{2.748613in}{0.987099in}}%
\pgfpathlineto{\pgfqpoint{2.748613in}{1.313349in}}%
\pgfusepath{stroke}%
\end{pgfscope}%
\begin{pgfscope}%
\pgfsetbuttcap%
\pgfsetroundjoin%
\definecolor{currentfill}{rgb}{0.000000,0.000000,0.000000}%
\pgfsetfillcolor{currentfill}%
\pgfsetlinewidth{0.803000pt}%
\definecolor{currentstroke}{rgb}{0.000000,0.000000,0.000000}%
\pgfsetstrokecolor{currentstroke}%
\pgfsetdash{}{0pt}%
\pgfsys@defobject{currentmarker}{\pgfqpoint{0.000000in}{0.000000in}}{\pgfqpoint{0.000000in}{0.027778in}}{%
\pgfpathmoveto{\pgfqpoint{0.000000in}{0.000000in}}%
\pgfpathlineto{\pgfqpoint{0.000000in}{0.027778in}}%
\pgfusepath{stroke,fill}%
}%
\begin{pgfscope}%
\pgfsys@transformshift{2.748613in}{0.987099in}%
\pgfsys@useobject{currentmarker}{}%
\end{pgfscope}%
\end{pgfscope}%
\begin{pgfscope}%
\definecolor{textcolor}{rgb}{0.000000,0.000000,0.000000}%
\pgfsetstrokecolor{textcolor}%
\pgfsetfillcolor{textcolor}%
\pgftext[x=2.748613in,y=0.973210in,,top]{\color{textcolor}{\rmfamily\fontsize{5.000000}{6.000000}\selectfont\catcode`\^=\active\def^{\ifmmode\sp\else\^{}\fi}\catcode`\%=\active\def%{\%}6.0}}%
\end{pgfscope}%
\begin{pgfscope}%
\pgfpathrectangle{\pgfqpoint{2.259013in}{0.987099in}}{\pgfqpoint{0.979200in}{0.326250in}}%
\pgfusepath{clip}%
\pgfsetbuttcap%
\pgfsetroundjoin%
\pgfsetlinewidth{0.401500pt}%
\definecolor{currentstroke}{rgb}{0.690196,0.690196,0.690196}%
\pgfsetstrokecolor{currentstroke}%
\pgfsetstrokeopacity{0.250000}%
\pgfsetdash{{1.480000pt}{0.640000pt}}{0.000000pt}%
\pgfpathmoveto{\pgfqpoint{2.993413in}{0.987099in}}%
\pgfpathlineto{\pgfqpoint{2.993413in}{1.313349in}}%
\pgfusepath{stroke}%
\end{pgfscope}%
\begin{pgfscope}%
\pgfsetbuttcap%
\pgfsetroundjoin%
\definecolor{currentfill}{rgb}{0.000000,0.000000,0.000000}%
\pgfsetfillcolor{currentfill}%
\pgfsetlinewidth{0.803000pt}%
\definecolor{currentstroke}{rgb}{0.000000,0.000000,0.000000}%
\pgfsetstrokecolor{currentstroke}%
\pgfsetdash{}{0pt}%
\pgfsys@defobject{currentmarker}{\pgfqpoint{0.000000in}{0.000000in}}{\pgfqpoint{0.000000in}{0.027778in}}{%
\pgfpathmoveto{\pgfqpoint{0.000000in}{0.000000in}}%
\pgfpathlineto{\pgfqpoint{0.000000in}{0.027778in}}%
\pgfusepath{stroke,fill}%
}%
\begin{pgfscope}%
\pgfsys@transformshift{2.993413in}{0.987099in}%
\pgfsys@useobject{currentmarker}{}%
\end{pgfscope}%
\end{pgfscope}%
\begin{pgfscope}%
\definecolor{textcolor}{rgb}{0.000000,0.000000,0.000000}%
\pgfsetstrokecolor{textcolor}%
\pgfsetfillcolor{textcolor}%
\pgftext[x=2.993413in,y=0.973210in,,top]{\color{textcolor}{\rmfamily\fontsize{5.000000}{6.000000}\selectfont\catcode`\^=\active\def^{\ifmmode\sp\else\^{}\fi}\catcode`\%=\active\def%{\%}6.0}}%
\end{pgfscope}%
\begin{pgfscope}%
\pgfpathrectangle{\pgfqpoint{2.259013in}{0.987099in}}{\pgfqpoint{0.979200in}{0.326250in}}%
\pgfusepath{clip}%
\pgfsetbuttcap%
\pgfsetroundjoin%
\pgfsetlinewidth{0.401500pt}%
\definecolor{currentstroke}{rgb}{0.690196,0.690196,0.690196}%
\pgfsetstrokecolor{currentstroke}%
\pgfsetstrokeopacity{0.250000}%
\pgfsetdash{{1.480000pt}{0.640000pt}}{0.000000pt}%
\pgfpathmoveto{\pgfqpoint{3.238213in}{0.987099in}}%
\pgfpathlineto{\pgfqpoint{3.238213in}{1.313349in}}%
\pgfusepath{stroke}%
\end{pgfscope}%
\begin{pgfscope}%
\pgfsetbuttcap%
\pgfsetroundjoin%
\definecolor{currentfill}{rgb}{0.000000,0.000000,0.000000}%
\pgfsetfillcolor{currentfill}%
\pgfsetlinewidth{0.803000pt}%
\definecolor{currentstroke}{rgb}{0.000000,0.000000,0.000000}%
\pgfsetstrokecolor{currentstroke}%
\pgfsetdash{}{0pt}%
\pgfsys@defobject{currentmarker}{\pgfqpoint{0.000000in}{0.000000in}}{\pgfqpoint{0.000000in}{0.027778in}}{%
\pgfpathmoveto{\pgfqpoint{0.000000in}{0.000000in}}%
\pgfpathlineto{\pgfqpoint{0.000000in}{0.027778in}}%
\pgfusepath{stroke,fill}%
}%
\begin{pgfscope}%
\pgfsys@transformshift{3.238213in}{0.987099in}%
\pgfsys@useobject{currentmarker}{}%
\end{pgfscope}%
\end{pgfscope}%
\begin{pgfscope}%
\definecolor{textcolor}{rgb}{0.000000,0.000000,0.000000}%
\pgfsetstrokecolor{textcolor}%
\pgfsetfillcolor{textcolor}%
\pgftext[x=3.238213in,y=0.973210in,,top]{\color{textcolor}{\rmfamily\fontsize{5.000000}{6.000000}\selectfont\catcode`\^=\active\def^{\ifmmode\sp\else\^{}\fi}\catcode`\%=\active\def%{\%}6.0}}%
\end{pgfscope}%
\begin{pgfscope}%
\pgfpathrectangle{\pgfqpoint{2.259013in}{0.987099in}}{\pgfqpoint{0.979200in}{0.326250in}}%
\pgfusepath{clip}%
\pgfsetbuttcap%
\pgfsetroundjoin%
\pgfsetlinewidth{0.401500pt}%
\definecolor{currentstroke}{rgb}{0.690196,0.690196,0.690196}%
\pgfsetstrokecolor{currentstroke}%
\pgfsetstrokeopacity{0.250000}%
\pgfsetdash{{1.480000pt}{0.640000pt}}{0.000000pt}%
\pgfpathmoveto{\pgfqpoint{2.259013in}{0.987099in}}%
\pgfpathlineto{\pgfqpoint{3.238213in}{0.987099in}}%
\pgfusepath{stroke}%
\end{pgfscope}%
\begin{pgfscope}%
\pgfsetbuttcap%
\pgfsetroundjoin%
\definecolor{currentfill}{rgb}{0.000000,0.000000,0.000000}%
\pgfsetfillcolor{currentfill}%
\pgfsetlinewidth{0.803000pt}%
\definecolor{currentstroke}{rgb}{0.000000,0.000000,0.000000}%
\pgfsetstrokecolor{currentstroke}%
\pgfsetdash{}{0pt}%
\pgfsys@defobject{currentmarker}{\pgfqpoint{0.000000in}{0.000000in}}{\pgfqpoint{0.027778in}{0.000000in}}{%
\pgfpathmoveto{\pgfqpoint{0.000000in}{0.000000in}}%
\pgfpathlineto{\pgfqpoint{0.027778in}{0.000000in}}%
\pgfusepath{stroke,fill}%
}%
\begin{pgfscope}%
\pgfsys@transformshift{2.259013in}{0.987099in}%
\pgfsys@useobject{currentmarker}{}%
\end{pgfscope}%
\end{pgfscope}%
\begin{pgfscope}%
\definecolor{textcolor}{rgb}{0.000000,0.000000,0.000000}%
\pgfsetstrokecolor{textcolor}%
\pgfsetfillcolor{textcolor}%
\pgftext[x=2.042575in, y=0.962986in, left, base]{\color{textcolor}{\rmfamily\fontsize{5.000000}{6.000000}\selectfont\catcode`\^=\active\def^{\ifmmode\sp\else\^{}\fi}\catcode`\%=\active\def%{\%}-2.00}}%
\end{pgfscope}%
\begin{pgfscope}%
\pgfpathrectangle{\pgfqpoint{2.259013in}{0.987099in}}{\pgfqpoint{0.979200in}{0.326250in}}%
\pgfusepath{clip}%
\pgfsetbuttcap%
\pgfsetroundjoin%
\pgfsetlinewidth{0.401500pt}%
\definecolor{currentstroke}{rgb}{0.690196,0.690196,0.690196}%
\pgfsetstrokecolor{currentstroke}%
\pgfsetstrokeopacity{0.250000}%
\pgfsetdash{{1.480000pt}{0.640000pt}}{0.000000pt}%
\pgfpathmoveto{\pgfqpoint{2.259013in}{1.150224in}}%
\pgfpathlineto{\pgfqpoint{3.238213in}{1.150224in}}%
\pgfusepath{stroke}%
\end{pgfscope}%
\begin{pgfscope}%
\pgfsetbuttcap%
\pgfsetroundjoin%
\definecolor{currentfill}{rgb}{0.000000,0.000000,0.000000}%
\pgfsetfillcolor{currentfill}%
\pgfsetlinewidth{0.803000pt}%
\definecolor{currentstroke}{rgb}{0.000000,0.000000,0.000000}%
\pgfsetstrokecolor{currentstroke}%
\pgfsetdash{}{0pt}%
\pgfsys@defobject{currentmarker}{\pgfqpoint{0.000000in}{0.000000in}}{\pgfqpoint{0.027778in}{0.000000in}}{%
\pgfpathmoveto{\pgfqpoint{0.000000in}{0.000000in}}%
\pgfpathlineto{\pgfqpoint{0.027778in}{0.000000in}}%
\pgfusepath{stroke,fill}%
}%
\begin{pgfscope}%
\pgfsys@transformshift{2.259013in}{1.150224in}%
\pgfsys@useobject{currentmarker}{}%
\end{pgfscope}%
\end{pgfscope}%
\begin{pgfscope}%
\definecolor{textcolor}{rgb}{0.000000,0.000000,0.000000}%
\pgfsetstrokecolor{textcolor}%
\pgfsetfillcolor{textcolor}%
\pgftext[x=2.075368in, y=1.126111in, left, base]{\color{textcolor}{\rmfamily\fontsize{5.000000}{6.000000}\selectfont\catcode`\^=\active\def^{\ifmmode\sp\else\^{}\fi}\catcode`\%=\active\def%{\%}0.00}}%
\end{pgfscope}%
\begin{pgfscope}%
\pgfpathrectangle{\pgfqpoint{2.259013in}{0.987099in}}{\pgfqpoint{0.979200in}{0.326250in}}%
\pgfusepath{clip}%
\pgfsetbuttcap%
\pgfsetroundjoin%
\pgfsetlinewidth{0.401500pt}%
\definecolor{currentstroke}{rgb}{0.690196,0.690196,0.690196}%
\pgfsetstrokecolor{currentstroke}%
\pgfsetstrokeopacity{0.250000}%
\pgfsetdash{{1.480000pt}{0.640000pt}}{0.000000pt}%
\pgfpathmoveto{\pgfqpoint{2.259013in}{1.313349in}}%
\pgfpathlineto{\pgfqpoint{3.238213in}{1.313349in}}%
\pgfusepath{stroke}%
\end{pgfscope}%
\begin{pgfscope}%
\pgfsetbuttcap%
\pgfsetroundjoin%
\definecolor{currentfill}{rgb}{0.000000,0.000000,0.000000}%
\pgfsetfillcolor{currentfill}%
\pgfsetlinewidth{0.803000pt}%
\definecolor{currentstroke}{rgb}{0.000000,0.000000,0.000000}%
\pgfsetstrokecolor{currentstroke}%
\pgfsetdash{}{0pt}%
\pgfsys@defobject{currentmarker}{\pgfqpoint{0.000000in}{0.000000in}}{\pgfqpoint{0.027778in}{0.000000in}}{%
\pgfpathmoveto{\pgfqpoint{0.000000in}{0.000000in}}%
\pgfpathlineto{\pgfqpoint{0.027778in}{0.000000in}}%
\pgfusepath{stroke,fill}%
}%
\begin{pgfscope}%
\pgfsys@transformshift{2.259013in}{1.313349in}%
\pgfsys@useobject{currentmarker}{}%
\end{pgfscope}%
\end{pgfscope}%
\begin{pgfscope}%
\definecolor{textcolor}{rgb}{0.000000,0.000000,0.000000}%
\pgfsetstrokecolor{textcolor}%
\pgfsetfillcolor{textcolor}%
\pgftext[x=2.075368in, y=1.289236in, left, base]{\color{textcolor}{\rmfamily\fontsize{5.000000}{6.000000}\selectfont\catcode`\^=\active\def^{\ifmmode\sp\else\^{}\fi}\catcode`\%=\active\def%{\%}2.00}}%
\end{pgfscope}%
\begin{pgfscope}%
\pgfpathrectangle{\pgfqpoint{2.259013in}{0.987099in}}{\pgfqpoint{0.979200in}{0.326250in}}%
\pgfusepath{clip}%
\pgfsetrectcap%
\pgfsetroundjoin%
\pgfsetlinewidth{1.505625pt}%
\definecolor{currentstroke}{rgb}{0.000000,0.619608,0.450980}%
\pgfsetstrokecolor{currentstroke}%
\pgfsetdash{}{0pt}%
\pgfpathmoveto{\pgfqpoint{2.255351in}{1.161634in}}%
\pgfpathlineto{\pgfqpoint{2.256724in}{0.980154in}}%
\pgfpathmoveto{\pgfqpoint{2.259618in}{0.980154in}}%
\pgfpathlineto{\pgfqpoint{2.260237in}{1.160265in}}%
\pgfpathlineto{\pgfqpoint{2.260494in}{1.159595in}}%
\pgfpathlineto{\pgfqpoint{2.262893in}{1.158980in}}%
\pgfpathlineto{\pgfqpoint{2.266936in}{1.157023in}}%
\pgfpathlineto{\pgfqpoint{2.268157in}{0.980154in}}%
\pgfpathmoveto{\pgfqpoint{2.518709in}{0.980154in}}%
\pgfpathlineto{\pgfqpoint{2.519180in}{1.166361in}}%
\pgfpathlineto{\pgfqpoint{2.526058in}{1.166059in}}%
\pgfpathlineto{\pgfqpoint{2.540746in}{1.166177in}}%
\pgfpathlineto{\pgfqpoint{2.555434in}{1.165759in}}%
\pgfpathlineto{\pgfqpoint{2.570122in}{1.165890in}}%
\pgfpathlineto{\pgfqpoint{2.584810in}{1.165402in}}%
\pgfpathlineto{\pgfqpoint{2.599498in}{1.165552in}}%
\pgfpathlineto{\pgfqpoint{2.614186in}{1.164968in}}%
\pgfpathlineto{\pgfqpoint{2.619082in}{1.165294in}}%
\pgfpathlineto{\pgfqpoint{2.623978in}{1.164802in}}%
\pgfpathlineto{\pgfqpoint{2.628874in}{1.165153in}}%
\pgfpathlineto{\pgfqpoint{2.633770in}{1.164622in}}%
\pgfpathlineto{\pgfqpoint{2.638666in}{1.165001in}}%
\pgfpathlineto{\pgfqpoint{2.643562in}{1.164424in}}%
\pgfpathlineto{\pgfqpoint{2.648458in}{1.164837in}}%
\pgfpathlineto{\pgfqpoint{2.653354in}{1.164207in}}%
\pgfpathlineto{\pgfqpoint{2.658250in}{1.164660in}}%
\pgfpathlineto{\pgfqpoint{2.663146in}{1.163965in}}%
\pgfpathlineto{\pgfqpoint{2.668042in}{1.164466in}}%
\pgfpathlineto{\pgfqpoint{2.672938in}{1.163691in}}%
\pgfpathlineto{\pgfqpoint{2.677834in}{1.164253in}}%
\pgfpathlineto{\pgfqpoint{2.682730in}{1.163377in}}%
\pgfpathlineto{\pgfqpoint{2.687626in}{1.164016in}}%
\pgfpathlineto{\pgfqpoint{2.692522in}{1.163007in}}%
\pgfpathlineto{\pgfqpoint{2.697418in}{1.163749in}}%
\pgfpathlineto{\pgfqpoint{2.702314in}{1.162560in}}%
\pgfpathlineto{\pgfqpoint{2.707210in}{1.163444in}}%
\pgfpathlineto{\pgfqpoint{2.712106in}{1.161993in}}%
\pgfpathlineto{\pgfqpoint{2.717002in}{1.163087in}}%
\pgfpathlineto{\pgfqpoint{2.721898in}{1.161215in}}%
\pgfpathlineto{\pgfqpoint{2.726794in}{1.162657in}}%
\pgfpathlineto{\pgfqpoint{2.731690in}{1.159975in}}%
\pgfpathlineto{\pgfqpoint{2.736586in}{1.162117in}}%
\pgfpathlineto{\pgfqpoint{2.741482in}{1.156622in}}%
\pgfpathlineto{\pgfqpoint{2.745047in}{1.161502in}}%
\pgfpathlineto{\pgfqpoint{2.746879in}{0.980154in}}%
\pgfpathmoveto{\pgfqpoint{2.748967in}{0.980154in}}%
\pgfpathlineto{\pgfqpoint{2.749837in}{1.160468in}}%
\pgfpathlineto{\pgfqpoint{2.750094in}{1.159508in}}%
\pgfpathlineto{\pgfqpoint{2.752493in}{1.158891in}}%
\pgfpathlineto{\pgfqpoint{2.756538in}{1.156750in}}%
\pgfpathlineto{\pgfqpoint{2.757791in}{0.980154in}}%
\pgfpathmoveto{\pgfqpoint{3.008134in}{0.980154in}}%
\pgfpathlineto{\pgfqpoint{3.008590in}{1.166384in}}%
\pgfpathlineto{\pgfqpoint{3.015465in}{1.166027in}}%
\pgfpathlineto{\pgfqpoint{3.030153in}{1.166204in}}%
\pgfpathlineto{\pgfqpoint{3.044841in}{1.165722in}}%
\pgfpathlineto{\pgfqpoint{3.059529in}{1.165921in}}%
\pgfpathlineto{\pgfqpoint{3.074217in}{1.165359in}}%
\pgfpathlineto{\pgfqpoint{3.079113in}{1.165702in}}%
\pgfpathlineto{\pgfqpoint{3.084009in}{1.165217in}}%
\pgfpathlineto{\pgfqpoint{3.088905in}{1.165586in}}%
\pgfpathlineto{\pgfqpoint{3.093801in}{1.165070in}}%
\pgfpathlineto{\pgfqpoint{3.098697in}{1.165464in}}%
\pgfpathlineto{\pgfqpoint{3.103593in}{1.164912in}}%
\pgfpathlineto{\pgfqpoint{3.108489in}{1.165333in}}%
\pgfpathlineto{\pgfqpoint{3.113385in}{1.164741in}}%
\pgfpathlineto{\pgfqpoint{3.118281in}{1.165194in}}%
\pgfpathlineto{\pgfqpoint{3.123177in}{1.164556in}}%
\pgfpathlineto{\pgfqpoint{3.128073in}{1.165045in}}%
\pgfpathlineto{\pgfqpoint{3.132969in}{1.164352in}}%
\pgfpathlineto{\pgfqpoint{3.137865in}{1.164885in}}%
\pgfpathlineto{\pgfqpoint{3.142761in}{1.164127in}}%
\pgfpathlineto{\pgfqpoint{3.147657in}{1.164712in}}%
\pgfpathlineto{\pgfqpoint{3.152553in}{1.163875in}}%
\pgfpathlineto{\pgfqpoint{3.157449in}{1.164523in}}%
\pgfpathlineto{\pgfqpoint{3.162345in}{1.163589in}}%
\pgfpathlineto{\pgfqpoint{3.167241in}{1.164316in}}%
\pgfpathlineto{\pgfqpoint{3.172137in}{1.163257in}}%
\pgfpathlineto{\pgfqpoint{3.177033in}{1.164087in}}%
\pgfpathlineto{\pgfqpoint{3.181929in}{1.162865in}}%
\pgfpathlineto{\pgfqpoint{3.186825in}{1.163829in}}%
\pgfpathlineto{\pgfqpoint{3.191721in}{1.162383in}}%
\pgfpathlineto{\pgfqpoint{3.196617in}{1.163536in}}%
\pgfpathlineto{\pgfqpoint{3.201513in}{1.161758in}}%
\pgfpathlineto{\pgfqpoint{3.206409in}{1.163195in}}%
\pgfpathlineto{\pgfqpoint{3.211305in}{1.160868in}}%
\pgfpathlineto{\pgfqpoint{3.216201in}{1.162789in}}%
\pgfpathlineto{\pgfqpoint{3.221097in}{1.159306in}}%
\pgfpathlineto{\pgfqpoint{3.225993in}{1.162286in}}%
\pgfpathlineto{\pgfqpoint{3.230889in}{1.078920in}}%
\pgfpathlineto{\pgfqpoint{3.234551in}{1.161642in}}%
\pgfpathlineto{\pgfqpoint{3.235914in}{0.980154in}}%
\pgfpathlineto{\pgfqpoint{3.235914in}{0.980154in}}%
\pgfusepath{stroke}%
\end{pgfscope}%
\begin{pgfscope}%
\pgfsetrectcap%
\pgfsetmiterjoin%
\pgfsetlinewidth{0.803000pt}%
\definecolor{currentstroke}{rgb}{0.000000,0.000000,0.000000}%
\pgfsetstrokecolor{currentstroke}%
\pgfsetdash{}{0pt}%
\pgfpathmoveto{\pgfqpoint{2.259013in}{0.987099in}}%
\pgfpathlineto{\pgfqpoint{2.259013in}{1.313349in}}%
\pgfusepath{stroke}%
\end{pgfscope}%
\begin{pgfscope}%
\pgfsetrectcap%
\pgfsetmiterjoin%
\pgfsetlinewidth{0.803000pt}%
\definecolor{currentstroke}{rgb}{0.000000,0.000000,0.000000}%
\pgfsetstrokecolor{currentstroke}%
\pgfsetdash{}{0pt}%
\pgfpathmoveto{\pgfqpoint{3.238213in}{0.987099in}}%
\pgfpathlineto{\pgfqpoint{3.238213in}{1.313349in}}%
\pgfusepath{stroke}%
\end{pgfscope}%
\begin{pgfscope}%
\pgfsetrectcap%
\pgfsetmiterjoin%
\pgfsetlinewidth{0.803000pt}%
\definecolor{currentstroke}{rgb}{0.000000,0.000000,0.000000}%
\pgfsetstrokecolor{currentstroke}%
\pgfsetdash{}{0pt}%
\pgfpathmoveto{\pgfqpoint{2.259013in}{0.987099in}}%
\pgfpathlineto{\pgfqpoint{3.238213in}{0.987099in}}%
\pgfusepath{stroke}%
\end{pgfscope}%
\begin{pgfscope}%
\pgfsetrectcap%
\pgfsetmiterjoin%
\pgfsetlinewidth{0.803000pt}%
\definecolor{currentstroke}{rgb}{0.000000,0.000000,0.000000}%
\pgfsetstrokecolor{currentstroke}%
\pgfsetdash{}{0pt}%
\pgfpathmoveto{\pgfqpoint{2.259013in}{1.313349in}}%
\pgfpathlineto{\pgfqpoint{3.238213in}{1.313349in}}%
\pgfusepath{stroke}%
\end{pgfscope}%
\end{pgfpicture}%
\makeatother%
\endgroup%

	\caption{(a) tensão de saída, (b) corrente do indutor e (c) tensão no diodo do circuito Buck, com diodo ideal e carga de $1$ A.}
	\label{fig:ex4_3}
\end{figure}

Observando que a corrente no indutor se anula ao final do período $T_s$, é esperado que o conversor esteja operando no limiar entre o CCM e o DCM. Segundo~\cite{Eri:20}, a carga máxima (corrente mínima ou resistência máxima) pode ser calculada por
\begin{equation*}
	R_{o,max} = \frac{2L_{min}}{(1-D)T_s} = 10~\Omega.
\end{equation*}
Este é justamente o valor da carga utilizada na simulação, o que corrobora com o resultado apresentado.

\medskip
\textbf{d)} Calcule o valor da tensão de saída, caso se opere no MCD com corrente média de saída de $0{,}5$ A.

\medskip
\textbf{Resolução}: a característica do conversor Buck operando em MCD é
\begin{equation*}
	M(D,K) = \frac{2}{1 + \sqrt{1 + 4K/D^2}},
\end{equation*}
em que
\begin{equation*}
	K = \frac{2L}{RT_s}.
\end{equation*}

Como a resistência de carga pode ser escrita como
\begin{equation*}
	R = \frac{V_o}{I_o},
\end{equation*}
segue que
\begin{equation*}
	K = \frac{2LI_o}{V_oT_s}.
\end{equation*}

Substituindo-se esta expressão na característica do conversor, obtém-se
\begin{equation*}
	\frac{V_o}{E}
	=
	\frac{2}
	{1+\sqrt{1+\dfrac{8LI_o}{D^2T_sV_o}}}.
\end{equation*}

Multiplicando-se ambos os lados da equação por $E$ e rearranjando os termos, resulta
\begin{equation*}
	V_o\sqrt{1+\frac{8LI_o}{D^2T_sV_o}}
	=
	2E-V_o.
\end{equation*}

Elevando-se ambos os lados ao quadrado, obtém-se
\begin{equation*}
	V_o^2
	\left(
	1+\frac{8LI_o}{D^2T_sV_o}
	\right)
	=
	(2E-V_o)^2.
\end{equation*}

Expandindo ambos os lados e cancelando os termos em $V_o^2$, segue que
\begin{equation*}
	\frac{8LI_o}{D^2T_s}V_o
	=
	4E^2-4EV_o.
\end{equation*}

Isolando-se a tensão de saída, obtém-se
\begin{equation*}
	V_o
	=
	\frac{4E^2}
	{4E+\dfrac{8LI_o}{D^2T_s}}
	=
	\frac{E}
	{1+\dfrac{2LI_o}{ED^2T_s}}.
\end{equation*}

Substituindo-se os valores $E=20$ V, $D=0{,}5$, $L=100~\mu$H, $T_s=40~\mu$s e $I_o=0{,}5$ A, obtém-se
\begin{equation*}
	V_o \approx 13{,}33~\text{V}.
\end{equation*}